% Self-contained effective-topos manuscript. Compile with pdfLaTeX.
\documentclass[11pt,a4paper]{article}
\usepackage[T1]{fontenc}
\usepackage[utf8]{inputenc}
\usepackage{lmodern}
\usepackage[margin=29mm]{geometry}
\usepackage{amsmath,amssymb,amsthm,mathtools}
\usepackage{microtype}
\usepackage{enumitem,booktabs,array,longtable}
\usepackage{tikz-cd}
\usepackage[colorlinks=true,linkcolor=blue!45!black,citecolor=blue!45!black,urlcolor=blue!45!black]{hyperref}
\usepackage{aliascnt}
\usepackage[nameinlink,noabbrev,capitalise]{cleveref}
\hypersetup{pdftitle={Synthetic Computability and the Sheaf Theory of Oracles and Forcing},pdfauthor={Conor Sheridan}}
\setlist{itemsep=2pt,topsep=5pt}
\setlength{\emergencystretch}{2em}
\numberwithin{equation}{section}
\newtheorem{theorem}{Theorem}[section]
\newaliascnt{proposition}{theorem}
\newtheorem{proposition}[proposition]{Proposition}
\aliascntresetthe{proposition}
\crefname{proposition}{Proposition}{Propositions}
\Crefname{proposition}{Proposition}{Propositions}
\newaliascnt{lemma}{theorem}
\newtheorem{lemma}[lemma]{Lemma}
\aliascntresetthe{lemma}
\crefname{lemma}{Lemma}{Lemmas}
\Crefname{lemma}{Lemma}{Lemmas}
\newaliascnt{corollary}{theorem}
\newtheorem{corollary}[corollary]{Corollary}
\aliascntresetthe{corollary}
\crefname{corollary}{Corollary}{Corollaries}
\Crefname{corollary}{Corollary}{Corollaries}
\theoremstyle{definition}
\newaliascnt{definition}{theorem}
\newtheorem{definition}[definition]{Definition}
\aliascntresetthe{definition}
\crefname{definition}{Definition}{Definitions}
\Crefname{definition}{Definition}{Definitions}
\newaliascnt{construction}{theorem}
\newtheorem{construction}[construction]{Construction}
\aliascntresetthe{construction}
\crefname{construction}{Construction}{Constructions}
\Crefname{construction}{Construction}{Constructions}
\newaliascnt{example}{theorem}
\newtheorem{example}[example]{Example}
\aliascntresetthe{example}
\crefname{example}{Example}{Examples}
\Crefname{example}{Example}{Examples}
\theoremstyle{remark}
\newaliascnt{remark}{theorem}
\newtheorem{remark}[remark]{Remark}
\aliascntresetthe{remark}
\crefname{remark}{Remark}{Remarks}
\Crefname{remark}{Remark}{Remarks}
\newcommand{\E}{\Eff}
\newcommand{\N}{\mathbb N}
\newcommand{\Pow}{\mathcal P}
\newcommand{\OO}{\mathcal O}
\newcommand{\Idl}{\operatorname{Idl}}
\newcommand{\Up}{\operatorname{Up}}
\newcommand{\Sc}{\operatorname{Sc}}
\newcommand{\Pt}{\operatorname{Pt}_{\Sigma}}
\newcommand{\Pre}{\mathbf{Preord}}
\newcommand{\ObsPre}{\mathbf{OPre}_{\Sigma}}
\newcommand{\Dcpo}{\mathbf{Dcpo}}
\newcommand{\ED}{\mathbf{EDcpo}_{\Sigma}}
\newcommand{\Alg}{\operatorname{Alg}}
\newcommand{\Par}{\operatorname{Par}}
\newcommand{\Pres}{\mathbf{Pres}}
\newcommand{\PAsm}{\mathbf{PAsm}}
\newcommand{\Eff}{\mathbf{Eff}}
\newcommand{\Hit}{\operatorname{Hit}}
\newcommand{\dom}{\operatorname{dom}}
\newcommand{\im}{\operatorname{im}}
\newcommand{\id}{\operatorname{id}}
\newcommand{\ev}{\operatorname{ev}}
\newcommand{\lfp}{\operatorname{lfp}}
\newcommand{\Fin}{\operatorname{Fin}}
\newcommand{\PC}{\mathbf{PC}}
\newcommand{\FunDom}{\mathcal D_{\mathrm{fun}}}
\newcommand{\FG}{\mathbf G}
\newcommand{\way}{\mathrel{\ll}}
\newcommand{\onto}{\twoheadrightarrow}
\newcommand{\incl}{\hookrightarrow}
\newcommand{\pmap}{\rightharpoonup}
\newcommand{\T}{\mathsf T}
\newcommand{\Th}{\widehat{\mathsf T}}
\newcommand{\true}{\mathsf{true}}
\newcommand{\false}{\mathsf{false}}
\newcommand{\overt}{\operatorname{Overt}_{\Sigma}}
\newcommand{\cl}{\operatorname{cl}_{\Sigma}}
\newcommand{\st}{\mathrel{|}}
\newcommand{\llbracket}{\mathopen{[\![}}
\newcommand{\rrbracket}{\mathclose{]\!]}}
\DeclareMathOperator{\supp}{supp}
\DeclareMathOperator{\graph}{graph}
\title{Synthetic Computability and the Sheaf Theory of Oracles and Forcing}
\author{
  Conor Sheridan\\[0.5ex]
  \href{mailto:siridconor@gmail.com}
       {\texttt{siridconor@gmail.com}}
}
\date{\today}
\newcommand{\C}{\mathcal C}
\newcommand{\Cont}{\mathbf{Cont}_{\Sigma}}
\newcommand{\Search}{\mathbf{Search}_{\Sigma}}
\newcommand{\HH}{\underline{\mathcal C}}
\newcommand{\Sier}{\mathbb S}
\newcommand{\Lift}{\mathsf L}
\newcommand{\Kl}{\operatorname{Kl}}
\newcommand{\Kar}{\operatorname{Kar}}
\newcommand{\Chu}{\operatorname{Chu}}
\newcommand{\Bun}{\mathbf{Bun}}
\newcommand{\Tr}{\operatorname{Tr}}
\newcommand{\disc}{\operatorname{disc}}
\newcommand{\Hom}{\operatorname{Hom}}

\newcommand{\RR}{\mathbb R}
\newcommand{\QQ}{\mathbb Q}
\newcommand{\Baire}{\mathbb B}
\newcommand{\RO}{\operatorname{RO}_{\Sigma}}
\newcommand{\defin}{\operatorname{def}}
\newcommand{\J}{\mathsf J_\Sigma}
\newcommand{\F}{F_\Sigma}
\newcommand{\LL}{L_\Sigma}
\newcommand{\Pos}{\mathbf{Pos}}
\newcommand{\LPos}{\operatorname{Pos}}
\newcommand{\AS}{\mathcal A_\Sigma}
\newcommand{\CS}{\mathcal C_\Sigma}
\newcommand{\Spl}{\operatorname{Spl}}
\newcommand{\Spec}{\operatorname{Spec}}
\newcommand{\Fix}{\operatorname{Fix}}
\newcommand{\Eq}{\operatorname{Eq}}
\newcommand{\Code}{\operatorname{Code}}
\newcommand{\Sub}{\operatorname{Sub}}
\newcommand{\Ran}{\operatorname{Ran}}
\theoremstyle{plain}
\newtheorem{maintheorem}{Theorem}
\renewcommand{\themaintheorem}{\Alph{maintheorem}}
\crefname{maintheorem}{Theorem}{Theorems}
\Crefname{maintheorem}{Theorem}{Theorems}
\setcounter{tocdepth}{1}
\newcolumntype{L}[1]{>{\raggedright\arraybackslash}p{#1}}
\newcommand{\Char}{\operatorname{Char}}


\newcommand{\FF}{\mathcal F}
\newcommand{\GG}{\mathcal G}
\newcommand{\Loc}{\mathbf{Loc}}
\newcommand{\Frm}{\mathbf{Frm}}
\newcommand{\Sh}{\operatorname{Sh}}
\newcommand{\ScLoc}{\mathsf{Sc}}
\newcommand{\Alex}{\mathsf{Alex}}
\newcommand{\Stab}{\operatorname{Stab}}
\newcommand{\Adm}{\operatorname{Adm}}
\newcommand{\Res}{\operatorname{Res}}
\newcommand{\Nuc}{\operatorname{Nuc}}
\newcommand{\LT}{\operatorname{LT}}
\newcommand{\Obs}{\operatorname{Obs}}
\newcommand{\Mod}{\operatorname{Mod}}
\newcommand{\Must}{\operatorname{Must}}
\newcommand{\May}{\operatorname{May}}
\newcommand{\Hist}{\operatorname{Hist}}
\newcommand{\Desc}{\operatorname{Desc}}
\newcommand{\Coh}{\operatorname{Coh}}
\newcommand{\Ans}{\operatorname{Ans}}
\newcommand{\ContIdl}{\operatorname{CIdl}}
\newcommand{\eps}{\varepsilon}
\newcommand{\colim}{\operatorname{colim}}
\newcommand{\coeq}{\operatorname{coeq}}
\newcommand{\Span}{\operatorname{Span}}
\newcommand{\Meet}{\operatorname{Meet}}
\newcommand{\BigWitness}{\operatorname{Big}}
\newcommand{\Inc}{\operatorname{Inc}}
\newcommand{\Cand}{\operatorname{Cand}}
\newcommand{\Ord}{\operatorname{Ord}}
\newcommand{\Cov}{\operatorname{Cov}}
\newcommand{\pr}{\operatorname{pr}}
\newcommand{\Sat}{\operatorname{Sat}}
\newcommand{\at}[1]{\left[#1\right]}
\tikzcdset{row sep=large,column sep=large}
\hypersetup{pdfsubject={Presented locales, effective observations, oracle interfaces, and forcing}}


\newcommand{\Comp}{\mathbb T_{\mathrm{comp}}}
\newcommand{\Ocomp}{\mathcal O_{\mathrm{comp}}}

\newcommand{\PresPt}{\mathbf{Pres}_{\Pt}}
\newcommand{\relpair}[2]{\langle #1,#2\rangle}
\newcommand{\RealQ}{\mathsf Q}
\newcommand{\Syn}{\mathsf{Syn}}
\newcommand{\TrAsm}{\mathcal A^{\mathrm{tr}}}
\newcommand{\Prior}{\operatorname{Prior}}
\newcommand{\Tree}{\mathsf{Tree}}
\newcommand{\code}[1]{\ulcorner #1\urcorner}


% Presented arithmetic and computational forcing predicates.
\newcommand{\SigA}{\mathsf S}
\newcommand{\PiA}{\mathsf P}
\newcommand{\Dec}{\mathsf{Dec}_{\mathrm{prog}}}
\newcommand{\Small}{\operatorname{Small}}
\newcommand{\RegOpen}{\operatorname{RO}}
\newcommand{\Dens}{\mathsf d}
\newcommand{\Jdens}{J_{\mathrm{dens}}}
\newcommand{\DNC}{\operatorname{DNC}}
\newcommand{\rawforces}{\Vdash_{\mathrm{raw}}}
\newcommand{\val}[1]{\llbracket #1\rrbracket}
\newcommand{\converges}{\mathord\downarrow}
\newcommand{\diverges}{\mathord\uparrow}

\begin{document}
\maketitle
\begin{abstract}
Forcing conditions, oracle queries, and programs can describe the same
computed outcome while retaining different information about how that
outcome is obtained. Comparing them requires more than identifying their
truth values: a comparison must also account for semidecidable
observations, partial evaluation, local names, and the information lost
under localization. We develop these comparisons over Hyland's effective
topos, using domain completion and realizability to provide a common
computational interpretation.

A positive pairing with the Rosolini dominance selects an admitted-ideal
Kock--Z\"oberlein monad. Its algebra category is cartesian closed, and its
continuous algebras are characterized by adjoint splittings. An
observation-equipped completion and a realizability presentation identify
when these domains and their maps admit effective approximations. The
associated classified partiality supports least fixed points, complete
Elgot iteration, and feedback with exact finite-return domains. Tracked
presentations interpret these operations while retaining their numerical
witnesses and description syntax.

We then interpret forcing sites as localized ideal-domain sheaves and
oracle answer interfaces as modalities. Universal membership and
observation pullbacks specify what is preserved under change of base;
finite-prefix reconstruction distinguishes sequences formed in a forcing
topos from imported ground sequences. Artin gluing retains raw
computational descriptions alongside their sheafified semantics, while
region, history, and descent constructions describe information that a
realized oracle alone does not determine. Numerical oracle and fusion
constructions supply concrete instances of these comparisons. Finally,
finite naming sites and an explicit deformation calculation exhibit
obstructions to selecting global names and extending infinitesimal data.
The resulting framework separates agreement of computed values from
agreement of their observations, names, and execution witnesses.
\end{abstract}
\tableofcontents

\section*{Introduction}
\addcontentsline{toc}{section}{Introduction}
A forcing condition may specify finite oracle information, a region of possible outcomes, a promise governing future extensions, or a program subject to restraints. Such descriptions can determine the same realized oracle without carrying the same information. A condition retains its history; a query retains possible answers; and a program retains its finite execution evidence. Identifying only the resulting oracle therefore does not identify the computations and forcing presentations from which it arose.

The problem addressed here is to compare these descriptions without suppressing their computational content. A forcing semantics must interpret conditions and generic values, but it must also say which observations are effective, when an answer can be used by a continuation, and whether a description available locally can be recovered globally. We work over Hyland's effective topos $\Eff$ \cite{hyland}. Its realizability structure provides numerical implementations; its internal logic provides the images, dependent constructions, and sheaf semantics needed for comparisons that are uniform in parameters.

Two familiar examples exhibit different parts of this problem. For structural Cohen forcing, the two immediate extensions of a binary condition determine the next bit locally, and finite-level search computes saturation of positive stage predicates. The resulting generic sequence is formed from compatible finite prefixes in the forcing topos; it need not come from a local representative of a ground sequence. For Mathias forcing, a reservoir restricts future choices. Passing to the eventual increasing sequence forgets that restraint, so the realization fibre must retain histories not recoverable from the sequence alone. These examples require different comparison maps even though both are presented by forcing conditions.

We first construct the computational objects on which those maps act. Ordinary ideal completion provides directed suprema, while the Rosolini truth object $\Sigma\incl\Omega$ selects the weights and observations whose positive tests can be implemented. The order remains $\Omega$-valued: a positive approximation relation does not make inclusion of completed ideals semidecidable. The same distinction governs names. A regular-epimorphic presentation supplies names locally and permits tracked lifts on suitable covers, but it supplies no global selection of names. These are structural constraints on the comparison, not artifacts of a chosen notation.

The development consequently has three connected stages. Domain completion, realizability, and partiality establish what it means for a computation and its finite evidence to be effective. Forcing sites, oracle interfaces, change of base, and sheafification then compare the observations available under localization, distinguishing imported from native objects. Finally, regions, histories, descent, concrete oracle constructions, and naming obstructions determine what information is retained or lost. The main statements below describe these stages before their proofs and applications.

\paragraph{Computational foundations.}
The first three principal statements summarize the foundation and implementation comparisons. All constructions and their proofs appear below.

\begin{maintheorem}[Admitted completion and partiality]\label{main:A}
In $\Eff$, positive pairing selects a canonical locally monotone KZ submonad $\J$ of ordinary ideal completion. It is maximal among the parameter-stable submonads whose admitted suprema on $\Sigma$ are computed through $\Sigma\incl\Omega$. Its algebra category $\AS$ is cartesian closed. Free algebras are algebraic; continuous algebras are the split algebras characterized by $t_A\dashv a_A\dashv\eta_A$ and $\Spl(\AS)\simeq\Kar\Kl(\J)$. A distributive law $\J\LL\Rightarrow\LL\J$ lifts positive partiality. Exact finite returns construct complete Elgot iteration and coproduct trace; finite completion of guarded running processes gives the same operations and observation extents.
\end{maintheorem}
The completion statements are proved in \cref{jd:thm:J,jd:prop:maximal,jd:prop:continuous,jd:thm:ccc,jd:thm:distributive}; the execution comparisons are \cref{exec:thm:elgot,exec:thm:trace,exec:thm:comparison}.

\begin{maintheorem}[Semantic and realizability characterization]\label{main:B}
Observation-equipped free completions characterize algebraic effective dcpos; their retracts characterize continuous ones; observed dcpo presentations characterize all effective dcpos and their maps. Exact-completion presentations over $K_1$ realize these structures on arbitrary carriers. The enriched adjunction $\F\dashv\Pt$ induces $\J$, and its comparison restricts to $\CS\simeq\Spl(\AS)$. On overt modest approximation bases, positive rounded distributors characterize the searchable objects and maps. The comparisons preserve composition, products, observations, partiality, and parameter substitution. In the algebraic fragment the canonical compact part recognizes the presentation intrinsically.
\end{maintheorem}
The general-carrier characterization is \cref{thm:fulladequacy}. The split and searchable comparisons are \cref{jd:thm:splitcomparison,thm:adequacycomparison,cor:canonicalrecognition}. Positive spectra, intrinsic observation transposes, and their forcing interface are \cref{jd:thm:spectra,sec:intrinsic,thm:domain-marking-bridge}.

\begin{maintheorem}[Implementations and retained descriptions]\label{main:C}
Regular-epimorphic name presentations form a fibration. Projective covers implement the computational operations, including evaluation, specialization, least solutions, Elgot iteration, and feedback, with equality of domains and interpreted values. Quotienting implementations by denotation gives the effective semantic category on the presented sorts; this comparison preserves execution trace and the associated interaction construction. Tracked assembly presentations carry a cartesian comonadic
realizer exposure. Finite descriptions form a tracked substitution monad;
its algebra on partial values commutes exactly with name interpretation,
common refinement, and parameter reindexing. Numerical program exposure
realizes intensional fixed points and the uniform recursion compiler.
The finite-information operations retain their witness transformations.
\end{maintheorem}
The constructions are
\cref{int:thm:exposure,thm:generalpresentationenhancement,thm:canonicalcoding,int:thm:exactsyntax,int:thm:refinement,int:thm:ifp}.
Execution and its finite witnesses are
\cref{exec:thm:codefeedback,thm:operational-certificates}.
The temporal self-interpretation is \cref{int:thm:temporal}, and
exact transport and gluing are
\cref{int:thm:transport,struct:thm:modelgluing}.

\paragraph{Forcing and oracles as sheaf theory.}
The following statements identify the geometry, observation calculus, and retained implementations. Their assembly is proved in \cref{sec:assembly}.

\begin{maintheorem}[Presented forcing and generic interpretation]\label{main:presentation}
A condition site in $\Eff$ determines a forcing frame, a localic sheaf topos, and a generic continuous ideal. The frame isomorphism $\Up(P)\cong\Sc(\Idl P)$ transports its coverage nucleus and identifies its condition-site and localized ideal-domain sheaf presentations. Rounded ideals have a geometric classifier with frame $\Sc(D_C)$. Frame-valued models of their atomic relations give oracle realizations. The canonical basis site recovers the frame and generic model, and full Kripke--Joyal forcing is the subobject semantics of the sheaf topos. The contextual cover-answer interface identifies this same sheaf category and reflector with the oracle localization, including the generic membership model and retained computational presentations.
Artin gluing of the common adjunction constructs a localic sheaf topos
over the base, retaining raw presentations fully faithfully and
realizing their semantic reflection by an open modality. Its
parameter slices and code-transport equations are constructed in
\cref{sec:artin-gluing}; see also
\cref{thm:forcing-domain-oracle-identity}.
\end{maintheorem}

\begin{maintheorem}[Effective observations and localization]\label{main:effective}
Stable effective observations are the pullback of the marking along the fixed-frame inclusion. Maximal admissibility is an indexed right adjoint to composition with a marking and applies also to named computations. Effective restriction is normal lax under composition, with an explicit equality criterion on its domain. Positive quantifier domains and evaluator inverse image construct the effective computation observations. Answer maps generate least localizations in which their projections are epic; contextual cover answers recover the forcing localization. Localization towers and their individually effective fragments compare coherently.
\end{maintheorem}

\begin{maintheorem}[Oracle modalities and adaptive computation]\label{main:modalities}
Indexed propositional queries generate the Ahman--Bauer oracle modality, and every Lawvere--Tierney topology has a canonical query presentation. Container morphisms and instance reductions induce modal inequalities; regular-epimorphic query changes preserve the modality. Sheafification has its oracle-algebra and modal-singleton universal properties. Coherent well-founded protocols identify with equifoliate trees; their interpretation is a monad morphism and their modal kernel-pair quotient is sheafification. Numerical tagged realizers recover the extended-Weihrauch presentation of local operators on $\Eff$.
\end{maintheorem}

\begin{maintheorem}[Regions, histories, descent, and comparison]\label{main:descent}
Region data give realizations by saturated universal observation. Realization fibres have canonical history sites; least-history sections are characterized by an order-enriched adjunction. A local candidate span has a universal descended observation locale, with its maximal effective marking given by a pullback. Geometric presentations, universal membership relations, and domain--value spans construct coherent comparisons for classifiers, completion objects, and partial evaluation under change of base. Continuous ideals and stabilized finite observations describe the generic objects obtained by requirement and fusion constructions.
\end{maintheorem}

\begin{maintheorem}[Computationally structured forcing]
\label{main:intensional}
The tracked presentations, realizer exposure diagrams, description
algebras, exact partial evaluators, execution witnesses, and oracle
interfaces form a single regular computational structure.
Its marked sheaf realization on the Artin gluing recovers the raw
structure on the closed part and its sheafification on the semantic
open, with effective descent and coherent parameter reindexing.
A two-step marked flag retains raw, requirement and dense forcing
semantics in one structure and recovers their localization maps.
Native completions enter through their constructed comparison maps;
transported names cover exactly their represented native image.

Well-founded protocols have a canonical native comparison preserving
and reflecting finite successful transcripts and their supports.
It respects substitution and the coherent semantic quotient.
Numerical Boolean oracles give internal partial combinatory algebras
with native finite-hit application domains; explicit dialogue compilers
preserve definedness and values. Native sequences are
recovered from compatible prefixes; the Cohen generic has no local
imported representative. Maximal admissibility computes the effective
forcing observations on the retained named computations.
\end{maintheorem}
The structure and recovery are \cref{sec:computational-structure};
the protocol and numerical constructions are
\cref{struct:thm:Wcomparison,struct:thm:transcript,struct:cor:coherent,struct:thm:pca,cal:thm:dialogues,cal:thm:gluing}.


\begin{maintheorem}[Arithmetic oracles and constructed forcing requirements]\label{main:constructions}
Presented numerical predicates form an arithmetic hierarchy before any
quotient by mutual computability. Their first existential level is
exactly dominance-classified semidecidability. Dense decision
localizations construct successive jump oracles; no program in the
previous oracle interpreter decides its own jump. Closed-subobject
interpretation gives the exact arithmetic semantics of localization,
and the classical numerical interpretation satisfies the degree-free
hierarchy--jump theorem.

Finite bushy witnesses give forbidden-region normalization, smallness
preservation, and an explicit oracle-relative successor. A uniform
halting-oracle requirement compiler constructs a diagonally noncomputable (DNC) function computing
no computably bounded DNC function. Partition
promises have a constructed maximal promise and coherent restriction.
Specified dense decision queries construct prehomogeneous and
homogeneous sequences. An infinite computable binary tree has a native
path whose jump is computed by the ground halting oracle. Protected
fronts give a perfect oracle-relative fusion satisfying all diagonal
requirements. Finite restraints have explicit functional extension
operations, and their use-growth covers make generic finite-axiom
membership decidable. For the cited arithmetical partition forcing,
the retained raw predicates have exactly their semantic Boolean truth
values after regular-open saturation. The extension-density nucleus and
actual condition-code comparisons calibrate the named forcing sites;
finite execution witnesses give their numerical arithmetic semantics. The glued interpretation retains
each raw predicate together with that value and computes the
reflection by its open modality.
\end{maintheorem}
The constructions are \cref{arith:prop:first-level,arith:prop:jump-strict,arith:thm:post,bush:prop:normalization,bush:thm:successor,partition:prop:promises,ramsey:thm:homogeneous,closed:thm:low-basis,fusion:thm:diagonal,ks:thm:extension,ks:prop:generic,forcing:thm:raw-saturation,glue:cor:raw-forcing}.

\begin{maintheorem}[Naming descent, localization defects, and deformation]
\label{main:obstructions}
Regular-epimorphic naming systems and their implemented
diagrams satisfy effective descent. Their locally isomorphic
forms are classified by torsors for the automorphism group
of the displayed diagram. Every name cover has a canonical
augmentation torsor, whose nonvanishing obstructs a global
selector. A four-stage effective information site carries
locally isomorphic answer interfaces with the same
existence modality and strict container-morphism order.

The abelian Artin gluing admits an open--closed recollement,
exact kernel-and-cokernel localization-defect sequences,
and tensor projection formulas. On an explicitly constructed
31-stage site, a finite bit-naming module over the dual
numbers has a nonzero second-cohomology obstruction to
extension from $\epsilon^2=0$ to $\epsilon^3=0$. Both
finite naming obstructions persist after adjoining an
independent structural Cohen oracle.
\end{maintheorem}
The constructions and proofs are
\cref{obs:prop:descent,obs:prop:forms,obs:prop:augmentation,obs:prop:recollement,obs:cor:projection,obs:prop:resolution,obs:thm:crown,obs:thm:deformation,obs:prop:cohen-product}.


The computational base is $\Eff$; forcing topoi and their localizations are constructed over it. Carrier, cover, and answer presentations record the chosen computational problem, whereas their completions, observation domains, generated localizations, and protocol objects are constructed by universal properties; see \cref{jd:thm:J,thm:admissibility-coreflection,prop:generated-coverage,thm:refinement-quotient,def:well-founded-protocol}. In particular, a selected answer projection does not include a choice of its section. The comparisons distinguish arbitrary exact-completion carriers from modest presentations, and imported objects from objects formed natively after change of base. Attribution and the relationship to the literature are collected in \cref{subsec:literature}.

\paragraph{Organization.}
Parts~\ref{part:foundations}--\ref{part:completion} establish the computational foundation, positive partiality and enriched domain theory. Parts~\ref{part:relative}--\ref{part:coding} develop realizability, temporal self-interpretation, tracked presentations, description algebras and presented arithmetic. Parts~\ref{part:presentations}--\ref{part:effective} construct forcing semantics and effective observations. Part~\ref{part:basechange} supplies exact description transport, geometric comparisons and native sequences used by the oracle constructions in Part~\ref{part:modalities}. Regions, histories and descent occupy Parts~\ref{part:regions}--\ref{part:descent}; naming obstructions and transport of implementations follow in Parts~\ref{part:naming-obstructions} and \ref{part:presentation-transport}. Part~\ref{part:external} gives requirement constructions, fusion and the final compatibility synthesis. The guide in \cref{app:guide} indexes these constructions by their universal properties.

\part{The effective topos and positive computation}
\label{part:foundations}

The later forcing comparisons depend on a uniform account of positive truth and partial computation in the ground topos. We begin with the internal logic of regular images and dependent predicates, then identify finite-hit truth using numerical realizers. This supplies the distinction between a semidecidable extent and an arbitrary truth value that the domain and forcing constructions will preserve.


\section{Internal logic, classified families, and local data}
\label{sec:internalconventions}

The same claim may be represented by a predicate, a witness cover, or a family varying over a parameter object. Their comparison requires the image and descent operations of the effective topos. We fix these operations first, so that subsequent uses of existential witnesses do not silently become choices of representatives.


Throughout the ground development, $\E$ is Hyland's effective topos. We use its exact/regular presentation $\operatorname{Asm}(K_1)_{\mathrm{ex/reg}}$; in the external set-theoretic metatheory it is also the exact/lex completion $\PAsm(K_1)_{\mathrm{ex}}$ of partitioned assemblies over Kleene's first partial combinatory algebra \cite{vanoosten,frey}; the exact/lex description and its external-choice dependence are explicit in \cite[Theorem~4.2 and the following discussion]{freyassemblies}. Its natural numbers object is $\N$. Fixed-carrier projectivity for arbitrary partitioned assemblies uses external choice, as recorded in \cref{rem:external-projectivity}. The choice-free implementation calculus is instead by covered spans and descent, constructed in \cref{prop:refined-implementation,thm:refinement-quotient,thm:partial-refinement-quotient}; target-modest lifts are supplied by \cref{lem:modest-target-lift}. The realization and witness constructions in \cref{sec:realizability} prove the dominance and overtness properties used below; they are properties of this fixed model. Forcing sheaf topoi are subsequently constructed over $\Eff$, with their ground, imported, and native objects distinguished explicitly. The subobject classifier is written $\true:1\incl\Omega$. A predicate $R$ on $X$ means a subobject of $X$, or equivalently its characteristic map $X\to\Omega$. Substitution is pullback, and $\exists_f R$ is the regular image of $R\incl X\xrightarrow fY$. Thus existential elimination in an internal proof is an image argument. It does not select a witness as a function of every parameter.

A variable $x:X$ is allowed to occur over an arbitrary parameter object $S$. More explicitly, it may denote a morphism $S\to X$. A statement proved with unused variables is consequently a statement in the corresponding slice. This convention is important when a construction is required to be uniform: a map into a power object describes a family in every parameter context, including its generalized elements.

For a truth value $p:S\to\Omega$, the notation $u:p$ refers to the subobject $S_p\incl S$ it classifies. If $q$ is a predicate on $S_p$, then
\[
    \exists u:p\;q(u)
\]
is the classifier of the composite of its domain with $S_p\incl S$. A dependent expression is evaluated only over its stated domain. In particular a value $v:p\to Y$ is an arrow $S_p\to Y$; there is no specified value on the complement of $S_p$.

An elementary topos is locally cartesian closed and exact \cite{johnstone,maclanemoerdijk}. We use its dependent products and sums, power objects, regular images, and effective quotients. A total functional relation determines a unique arrow: its projection to the source is both monic and regular epic, hence an isomorphism. This form of unique choice will be used to assemble unique factorizations. It is different from selecting one member of an arbitrary inhabited fibre.

\begin{definition}[Local representatives]
\label{def:localrepresentative}
    Let $e:V\onto W$ be a regular epimorphism. A \textbf{local representative} of a morphism $f:S\to W$ along $e$ consists of a regular epimorphism $p:T\onto S$ and a morphism $\widetilde f:T\to V$ satisfying
    \[
        e\widetilde f=fp.
    \]
    Thus $f$ admits a lift after passage to a regular-epimorphic cover of its domain, without requiring a global choice of representatives.
\end{definition}

\begin{lemma}[Local representatives and descent]
\label{lem:localrepresentatives}
    Let $e:V\onto W$ be a regular epimorphism. For a map $f:S\to W$, its pullback $S\times_WV\onto S$ supplies a local representative of $f$. A map $g:V\to Z$ descends through $e$ precisely when its two composites with the kernel pair of $e$ agree. Equalities of maps, and inclusions of subobjects, can be checked after a regular-epimorphic cover.
\end{lemma}
\begin{proof}
\leavevmode
\begin{enumerate}[label=\textup{(\roman*)},leftmargin=*,itemsep=3pt,topsep=4pt]
\item \textbf{Local representatives and descent.} Regular epimorphisms are stable under pullback.
Since $\E$ is exact, $e$ is the coequalizer of its kernel pair; this gives the descent assertion and uniqueness.
Epimorphy reflects equality of maps.
If $R,S\incl W$ have $e^*R\subseteq e^*S$, apply the regular-image functor along $e$.
The images are $R$ and $S$, because the pullbacks of $e$ over these subobjects are regular epic.
Hence $R\subseteq S$.

The construction uses the cartesian square
\[
\begin{tikzcd}[column sep=large,row sep=large]
S\times_WV\ar[r,"\pi_V"]\ar[d,two heads,"\pi_S"']&V\ar[d,two heads,"e"]\\ S\ar[r,"f"']&W
\end{tikzcd}
\]
\end{enumerate}
\end{proof}

\subsection{Representability as an indexed coreflection}
\begin{theorem}[Indexed coreflection of admissible families]\label{thm:admissibility-coreflection}
For any mono $m:M\incl H$ in a topos, composition with $m$ and
pullback along $m$ give an adjunction
\[
 \Sigma_m:\E/M\rightleftarrows\E/H:m^*,\qquad
 \Sigma_m\dashv m^*.
\]
The left adjoint is fully faithful. Its counit at a family $f:T\to H$
is the mono $T\times_HM\incl T$, which is the largest subobject of
$T$ on which $f$ admits an $M$-representative. This construction is
stable under substitution, and iterated admissibility is identified
with the one-step pullback by canonical, coherently associative
pullback-pasting isomorphisms.
\end{theorem}
\begin{proof}
\leavevmode
\begin{enumerate}[label=\textup{(\roman*)},leftmargin=*,itemsep=3pt,topsep=4pt]
\item \textbf{Adjunction and full faithfulness.} A map over $H$ from $m u:S\to H$ to $f:T\to H$ is a map
$h:S\to T$ satisfying $fh=mu$.
The universal property of the cartesian square
\[
\begin{tikzcd}[column sep=large]
T\times_HM\ar[r]\ar[d,hook,"{\varepsilon_f}"']&M\ar[d,hook,"m"]\\
T\ar[r,"f"']&H
\end{tikzcd}
\]
identifies this with a map over $M$ from $u:S\to M$ to the
second projection.
Naturality of this
bijection is the adjunction.
Because $m$ is monic, an arrow between
two $M$-families commutes with their structure maps if and only if
its composite with $m$ does; hence $\Sigma_m$ is fully faithful.
\item \textbf{Maximality and substitution.} The counit is the first projection of the pullback, a mono.
If
$i:S\incl T$ and $fi$ factors through $m$, the pullback supplies a
unique factorization $S\to T\times_HM$ over $T$.
Thus no larger
admissible domain exists.
Reindexing a family and then taking this
pullback gives the same object as pulling back the original
admissibility mono.
For two or three consecutive markings, the
iterated and direct constructions are finite limits with the same
specified projections.
\item \textbf{Pullback-pasting coherence.} Their unique comparison maps commute, and
uniqueness also proves the unit and associativity coherence.
\end{enumerate}
\end{proof}


The pullback characterization of admissible families is now available in every parameter context. Later restrictions of forcing truths and named evaluations will use this same maximality construction; none of those restrictions requires a global choice of representatives.

\section{Finite-hit truth and the computational foundation}
\label{sec:dominance}

The order-theoretic constructions will use all truth values, but execution domains and positive observations need a smaller class with computational witnesses. The finite-hit image defines that class; its closure properties determine which partial maps and existential searches can be composed.


\begin{definition}[Finite-hit map and Rosolini object]
\label{def:finitehit}
    Put $2=1+1$. The \textbf{finite-hit map} is
    \[
        h:2^{\N}\longrightarrow\Omega,
        \qquad h(\alpha)\iff\exists n:\N\;[\alpha(n)=1].
    \]
    The \textbf{Rosolini object} $\Sigma$ is the regular image of $h$, given by its regular epi--mono factorization
    \[
        2^{\N}\xrightarrow{e}\Sigma\xrightarrow{i}\Omega.
    \]
    We identify $\Sigma$ with its image in $\Omega$ when writing formulas.
\end{definition}

\begin{lemma}
\label{lem:finiteclosure}
    The Rosolini object contains $\bot$ and $\top$ and is closed under binary meets and joins, uniformly in parameters.
\end{lemma}
\begin{proof}
\leavevmode
\begin{enumerate}[label=\textup{(\roman*)},leftmargin=*,itemsep=3pt,topsep=4pt]
\item \textbf{Construction of the comparison.} The constant zero and one streams represent $\bot$ and $\top$.
Interleave two streams to represent disjunction.
For conjunction use the stream whose $k$th value is one exactly when each input stream has a hit at a stage at most $k$.
The resulting hit propositions are respectively the join and meet of the input hit propositions.
Consequently the maps descend through the regular epimorphism $e\times e$, since their composites with $i$ depend only on the two classified truth values.
\item \textbf{Independence of stream names.} This proves the assertion without choosing a stream representative for every element of $\Sigma$.
\end{enumerate}
\end{proof}

\begin{definition}[Dominance and the Rosolini condition]
\label{def:dominance}
    A \textbf{dominance} is a subobject $j:\Delta\incl\Omega$ containing $\top$ and closed under dependent conjunction: in every parameter context, if $p\in\Delta$ and $q:p\to\Delta$, then $\exists u:p\,q(u)$ factors through $\Delta$. A mono is \textbf{$\Delta$-classified} when its characteristic map factors through $j$. For the finite-hit image $i:\Sigma\incl\Omega$, its classified monos are called \textbf{$\Sigma$-classified}. The realizability construction below proves dependent closure.
\end{definition}

\begin{proposition}[Composition criterion]
\label{prop:dominancecriterion}
    For a subobject $\Delta\incl\Omega$ containing $\top$, dependent closure is equivalent to closure of its classified monomorphisms under composition. Classified monomorphisms are stable under pullback and contain all identities.
\end{proposition}
\begin{proof}
\leavevmode
\begin{enumerate}[label=\textup{(\roman*)},leftmargin=*,itemsep=3pt,topsep=4pt]
\item \textbf{Composition criterion.} If $U\incl X$ has classifier $p:X\to\Delta$ and $V\incl U$ has classifier $q:U\to\Delta$, the composite has classifier $x\mapsto\exists u:p(x)\,q(x,u)$.
Dependent closure therefore gives composition closure.
\item \textbf{Converse construction.} Conversely, represent a dependent pair $p,q$ by its two successive classified subobjects in the given context and apply composition closure.
Pullback composes a classifier with the substitution map, and an identity is classified by the constant truth map.

The defining maps satisfy the following commutative diagram
\[
\begin{tikzcd}[column sep=large,row sep=large]
V\ar[r,hook]\ar[dr,hook]&U\ar[d,hook]\\ &X
\end{tikzcd}
\]
\end{enumerate}
\end{proof}

This is the form of dominance used by Rosolini \cite[Chapter 3, Proposition 3.1.6 and Corollary 3.1.7]{rosolini}; partial-map representation appears there as Theorem 3.1.4. A modern account of dominance-based partial elements and their constructive logical setting is \cite[\S\S2.5--2.8]{escardoknapp}. Finite meets of the canonical finite-hit image were proved above without dependent closure. 

\begin{definition}
\label{def:overt}
    An object $X$ is \textbf{$\Sigma$-overt} if existential quantification restricts to a map
    \[
        \exists_X^\Sigma:\Sigma^X\longrightarrow\Sigma.
    \]
    Equivalently, $\exists_X:\Omega^X\to\Omega$ carries $\Sigma$-valued predicates to $\Sigma$. The factorization is unique.
\end{definition}

The definition has a dependent form in context: $X\to S$ is \textbf{relatively overt} when the corresponding existential operation in the slice restricts to $\Sigma$. The canonical Rosolini object in a slice is the pullback of $\Sigma$, since its defining image factorization is stable under pullback.


\paragraph{Computability terminology.}
For numerical functions and sets we use \emph{computable}, \emph{partial computable}, and \emph{computably enumerable} (c.e.) \cite{soare-terminology}. The standard term \emph{primitive recursive} retains its more restrictive meaning; recursion names a defining construction or a fixed-point theorem. Internally, \emph{semidecidable} means $\Sigma$-classified. Its correspondence with c.e. predicates on numerical names is the finite-hit comparison below. The terms $A$-computable and $A$-c.e. mean computable and computably enumerable relative to the displayed Boolean oracle $A$.

\subsection{The numerical interpreter}\label{subsec:numerical-foundation}
Fix a finite register-program syntax with instructions for increment, zero test with decrement, jumps, and halt. Finite instruction lists have primitive recursive codes. Invalid codes are interpreted as a permanently looping program. Configurations, one-step execution, and bounded simulation are primitive recursive. A universal interpreter runs the coded configuration step by step. Thus there are primitive recursive predicates and output functions giving
\begin{equation}\label{eq:universal}
 \varphi_e(n)=m\iff\exists s\;T(e,n,s,m),
\end{equation}
where the finite witness records a halting computation. The graph is functional. Conversely the register-program constructions of composition, primitive recursion, and minimization are executable by register programs, and a register computation is expressible by primitive recursive stepping followed by minimization. This fixes the numbering of partial computable functions used to present $K_1$.


Finite-hit truth has therefore been defined by an image rather than by a choice of a deciding stream. The numerical presentation that follows will show how its witnesses can nevertheless be searched and transported when the appropriate names are available.

\section{Partitioned assemblies and finite-hit predicates}
\label{sec:realizability}

The finite-hit definition becomes computationally meaningful only when its classifiers are related to numerical names. We use the realizability presentation of $\Eff$ to extract such names, while keeping equality of represented elements separate from equality of their codes. The distinction is needed both for dominance and for later oracle implementations.


Kleene's first partial combinatory algebra $K_1$ has carrier $\N$ and application $en\simeq\varphi_e(n)$, for a fixed acceptable numbering of partial computable functions. The register interpreter above fixes this numbering; its specialization maps are developed in \cref{sec:numericalcodes}. Pairing, tags, numerals, primitive recursion and delayed abstraction are represented by numerical programs. Choose combinators $k,s$ with $kab=a$, $sab\downarrow$, and $sabc\simeq ac(bc)$; here Kleene equality includes equality of definedness. The index $sab$ is constructed without running either represented application. These are computational properties of $K_1$.

A partitioned assembly $P$ consists of a set $|P|$ and a naming map $\varepsilon_P:|P|\to\N$. A function $f:|P|\to|Q|$ is a morphism when some index $b$ satisfies $b\varepsilon_P(x)\downarrow$ and $b\varepsilon_P(x)=\varepsilon_Q(fx)$ for every $x\in|P|$. The naming map need not be injective. In particular arbitrary carriers are retained together with their naming maps. The effective topos is
\[
 \Eff=\PAsm(K_1)_{\mathrm{ex}}.
\]
Names and trackers use this fixed numerical application structure \cite{hyland,vanoosten,frey}.

The canonical embedding of partitioned assemblies in this exact completion is full; its objects are regular-projective and cover every object. Finite products, tagged finite sums, and carrier subsets with restricted naming maps remain partitioned assemblies. We use these exact-completion properties in the external metatheory \cite{vanoosten,frey}. Thus a projective cover means a partitioned-assembly cover or a partitioned refinement of a displayed pullback. Two selected covers are compared by a partitioned cover of their fibre product. The interpretation equations descend through that common refinement and do not choose one name for each semantic element.

\begin{remark}[External projectivity and canonical covers]
\label{rem:external-projectivity}
The assertion that every arbitrary-carrier partitioned assembly is
regular-projective uses choice in the external category of sets; it is
not an internal choice principle in $\Eff$. Indeed, give an external
surjection $u:X\onto Y$ the constant naming function $0$ on both
carriers. It is a regular epimorphism of assemblies, since a constant
realizer lifts names. Projectivity of the partitioned assembly on $Y$
would split this map, and the underlying function would split $u$.
Thus unrestricted partitioned-assembly projectivity entails external
choice. This is the dependence in the ex/lex projective presentation
by the ex/lex completion; the exact/regular completion of assemblies
and regular-epimorphic refinement do not invoke a fixed-carrier lift
\cite{hofstra,vanoosten,frey,freyassemblies}.

Canonical covering constructions themselves do not select witnesses.
For an assembly $X$ with realizing relation $n\Vdash_Xx$, the assembly
\[
 \widehat X=\{(x,n):n\Vdash_Xx\},\qquad
 \varepsilon_{\widehat X}(x,n)=n
\]
is partitioned and covers $X$ by projection. More generally an object
with realizability equality $E(x,y)\subseteq\N$ is covered by
$\{(x,n):n\in E(x,x)\}$ with the same naming map. No realizer is
chosen for each $x$. We use these covers in
\cref{lem:modest-target-lift,prop:refined-implementation}; projectivity
is retained only where a lift on a fixed arbitrary name carrier is
actually extracted.
\end{remark}

\begin{lemma}[Lifting to a modest name object]
\label{lem:modest-target-lift}
Let $C$ be a partitioned assembly, let $D$ be a modest assembly, and
let $r:D\onto Y$ be a regular epimorphism in $\Eff$. Every map
$f:C\to Y$ has a tracked lift $\widetilde f:C\to D$ through $r$.
This construction does not use external choice.
\end{lemma}
\begin{proof}
\leavevmode
\begin{enumerate}[label=\textup{(\roman*)},leftmargin=*,itemsep=3pt,topsep=4pt]
\item \textbf{Tracked implementation.} Regular epimorphy makes the realized relation $r(d)=f(c)$ total in
$c$. A realizer of this totality, applied to the distinguished name
$\varepsilon_C(c)$, yields a name $n_c$ realizing some $d\in D$ and
evidence of $r(d)=f(c)$.
Modesty says that a name realizes at most one
element.
Hence the relation $n_c\Vdash_Dd$ is total and functional;
unique-value comprehension defines $\widetilde f(c)=d$.
The program
extracting $n_c$ is its tracker, and the equality evidence gives
$r\widetilde f=f$.
No element of a multiply inhabited fibre is selected.

The construction uses the commutative lifting triangle
\[
\begin{tikzcd}[column sep=large,row sep=large]
C\ar[r,"\widetilde f"]\ar[dr,"f"']&D\ar[d,two heads,"r"]\\ &Y
\end{tikzcd}
\]
\end{enumerate}
\end{proof}


A partial equivalence relation $\rho$ on $K_1$ presents a modest object $X_\rho$. Its canonical cover is the partitioned assembly $P_\rho$ with carrier $|\rho|=\{a\mid\rho(a,a)\}$ and naming map the inclusion into $K_1$. Two names denote the same element exactly when they are $\rho$-related. Its kernel pair is the partitioned assembly of pairs $(a,a')$ satisfying $\rho(a,a')$.

\begin{proposition}
\label{prop:per}
    Modest objects are equivalent to PER presentations. A map $X_\rho\to X_\tau$ is represented by $b\in K_1$ that is defined on $|\rho|$, takes values in $|\tau|$, and satisfies
    \[
        \rho(a,a')\Longrightarrow\tau(ba,ba').
    \]
    Two representatives are equal precisely when $\tau(ba,ca)$ for every $a\in|\rho|$.
\end{proposition}
\begin{proof}
\leavevmode
\begin{enumerate}[label=\textup{(\roman*)},leftmargin=*,itemsep=3pt,topsep=4pt]
\item \textbf{Kernel-pair descent.} The support assembly $P_\tau$ is modest, so \cref{lem:modest-target-lift} lifts a map through $P_\tau\onto X_\tau$, giving its tracker without external choice.
Conversely an extensional tracker equalizes the kernel pair and descends.
Equality is tested after the epic cover. A modest assembly determines the PER relating names that realize the same element; modesty makes this a PER, and identity on names gives an isomorphism with its quotient presentation.
The converse quotient is modest because a common name forces equality of the two PER classes.

The interpretation maps form the commuting diagram
\[
\begin{tikzcd}[column sep=large,row sep=large]
P_\rho\ar[r,"\widehat f"]\ar[d,two heads,"q_\rho"']&P_\tau\ar[d,two heads,"q_\tau"]\\ X_\rho\ar[r,"f"']&X_\tau
\end{tikzcd}
\]
\end{enumerate}
\end{proof}

\begin{lemma}[Function names and subobjects]
\label{lem:modestclosure}
    Finite products, exponentials between modest objects, and subobjects of modest objects are modest in $\Eff$.
\end{lemma}
\begin{proof}
\leavevmode
\begin{enumerate}[label=\textup{(\roman*)},leftmargin=*,itemsep=3pt,topsep=4pt]
\item \textbf{Product structure.} Products use paired names and the product PER.
For $X_\rho$ and $X_\tau$, define the support of a function PER to consist of $a\in K_1$ such that $ax$ is defined and $\tau$-valid for every $\rho$-valid $x$, and $\rho(x,x')$ implies $\tau(ax,ax')$.
Relate two such names $a,a'$ when
    \[
        \forall x\in|\rho|\;\tau(ax,a'x).
    \]
    Symmetry and transitivity of $\tau$, together with extensionality on inputs, prove that this is a PER.
Evaluation is tracked by application.
\item \textbf{Adjunction and universal property.} To verify the exponential universal property for an arbitrary parameter object $Z$, take a canonical partitioned cover $P\onto Z$. A map $Z\times X_\rho\to X_\tau$ lifts on $P\times P_\rho$ to the modest target name cover by \cref{lem:modest-target-lift}.
If $t\in K_1$ tracks this lift, combinatory abstraction turns a name $z$ of $P$ into a name of the function $x\mapsto t\langle z,x\rangle$.
Equality of parameter denotations makes these function names equivalent.
They therefore descend to a map $Z\to X_{\tau^\rho}$.
Evaluation recovers the original map.
\item \textbf{Converse construction.} Conversely, equality after evaluation means precisely the displayed PER equality, so the transpose is unique.
This proves the exponential property in the whole topos.
\item \textbf{Functoriality and coherence.} For a mono $U\hookrightarrow X_\rho$, take a canonical partitioned cover $P\onto U$ and use \cref{lem:modest-target-lift} to lift $P\to X_\rho$ to the modest support assembly $P_\rho$.
If two elements of $P$ have the same name, their images under this tracked lift have the same name, hence the same $X_\rho$-denotation.
Monicity then makes their $U$-denotations equal.
Put a PER on the names occurring in $P$ by equality of these denotations.
\item \textbf{Kernel-pair descent.} Concretely, the equality can be tested at the lifted $\rho$-names: because $U\to X_\rho$ is monic, its kernel pair after pullback to $P$ is exactly the pullback of the $\rho$-equality relation.
The quotient of $P$ by this kernel pair is both $U$ and the modest assembly with the displayed PER.
Thus they are isomorphic.
This argument retains the names of the cover of $U$; replacing them by a bare subset of the carrier of $X_\rho$ would lose membership evidence.
\end{enumerate}
\end{proof}

A morphism $\theta:P\to2^{\N}$ is a family of total binary streams whose evaluation on $(\varepsilon_P(p),n)$ is uniformly tracked in $K_1$. Put
\[
    H_\theta=\{(p,n)\mid\theta_p(n)=1\},\qquad
    \Hit(\theta)=\{p\mid\exists n\;\theta_p(n)=1\}.
\]
The hit assembly is partitioned. Its projection to $P$ has regular image classified by $ie\theta$.

\begin{lemma}
\label{lem:hitsearch}
    There is a uniform tracker which, on a total binary stream with a hit, returns its first hit. On inputs with no hit, no definedness assertion is needed.
\end{lemma}
\begin{proof}
\leavevmode
\begin{enumerate}[label=\textup{(\roman*)},leftmargin=*,itemsep=3pt,topsep=4pt]
\item \textbf{Partiality and definedness.} Combinatory abstraction and a fixed-point combinator represent the guarded recursion which evaluates stage $n$, returns $n$ if it is a hit, and otherwise recurses at $n+1$.
If a least hit is $m$, downward induction from $m$ proves termination with output $m$.
All operations used belong to $K_1$.
\end{enumerate}
\end{proof}

\begin{lemma}[Support order]
\label{lem:support}
    For tracked streams on a partitioned assembly,
    \[
        e\theta\leq e\eta\quad\Longleftrightarrow\quad
        \Hit(\theta)\subseteq\Hit(\eta).
    \]
\end{lemma}
\begin{proof}
\leavevmode
\begin{enumerate}[label=\textup{(\roman*)},leftmargin=*,itemsep=3pt,topsep=4pt]
\item \textbf{Support order.} An inclusion of the classified subobjects puts each point of the image of $H_\theta$ in the image of $H_\eta$.
By the definition of that image there exists a hit of $\eta$ at the same parameter, proving support inclusion without selecting such a hit for every parameter.
\item \textbf{Converse construction.} Conversely, on $H_\theta$ search $\eta_p$ for its first hit.
Support inclusion guarantees termination, and \cref{lem:hitsearch} gives a tracked map $H_\theta\to H_\eta$ over $P$.
Its regular image gives the required subobject inclusion.
\end{enumerate}
\end{proof}


\subsection{Finite-hit witnesses and descent}
\label{subsec:finite-hit-witnesses-and-descent}

The support lemma is stronger than a comparison of subsets of an external carrier. Its reverse implication constructs a morphism of hit assemblies. If $\theta$ and $\eta$ have the same support, the two first-hit searches give arrows in both directions over the parameter object, and therefore the same regular image. This is what permits an external equality of semidecisions to become equality of internal classifiers.

For later use, finite logical operations can also be realized at the witness level. A disjunction stream interleaves $\theta$ and $\eta$. Its hit stage records which summand supplied a hit; parity and stage decoding give a tracked isomorphism between its hit assembly and the tagged sum of the two hit assemblies. A conjunction stream tests whether each stream has hit by stage $k$. A pair of input hits determines a conjunction hit at their maximum, while a conjunction hit allows bounded search up to $k$ to return one hit of each stream. These are maps in both directions over the parameter. They show that the regular images are the intended union and intersection and retain the witnesses used by later proof transformations.

For a PER presentation the relevant parameter equality is not equality of names. Given $\rho(a,a')$ and a hit of a saturated stream at $a$, saturation ensures a hit at $a'$, and first-hit search finds one. The output stage can differ arbitrarily from the input stage. Consequently no requirement that equal names have identical streams is imposed: equality concerns their classified predicates, and witness transport accounts for the change of name.


\begin{theorem}[Finite-hit classification]
\label{thm:finitehit}
    Maps $X_\rho\to\Sigma$ correspond naturally to tracked streams on $P_\rho$ whose hit support is invariant under $\rho$, modulo equality of supports.
\end{theorem}
\begin{proof}
\leavevmode
\begin{enumerate}[label=\textup{(\roman*)},leftmargin=*,itemsep=3pt,topsep=4pt]
\item \textbf{A tracked stream representation.} By \cref{lem:modestclosure}, the function PER for total binary-stream indices presents $2^{\N}$ as a modest assembly.
Apply \cref{lem:modest-target-lift} to lift the composite $P_\rho\to X_\rho\to\Sigma$ through $e:2^{\N}\onto\Sigma$.
Its two pullbacks along the kernel pair have equal classifiers, hence equal hit supports by \cref{lem:support}.
\item \textbf{Saturation along the kernel pair.} Conversely invariant support gives equality of those classifiers and therefore descent.
Equality of descended maps is equality after the epic cover, again equivalent to support equality.
Precomposition with a tracker substitutes its output name into the stream; extensionality ensures independence of the chosen representative.

The interpretation maps form the commuting diagram
\[
\begin{tikzcd}[column sep=large,row sep=large]
P_\rho\ar[r,"\theta"]\ar[d,two heads,"q_\rho"']&2^{\N}\ar[d,two heads,"e"]\\ X_\rho\ar[r,"\sigma"']&\Sigma
\end{tikzcd}
\]
\end{enumerate}
\end{proof}

\begin{theorem}
\label{thm:relativefoundation}
    The canonical Rosolini object in $\Eff$ is a dominance, and $\N$ is $\Sigma$-overt.
\end{theorem}
\begin{proof}
\leavevmode
\begin{enumerate}[label=\textup{(\roman*)},leftmargin=*,itemsep=3pt,topsep=4pt]
\item \textbf{Composition of positive domains.} For classified monos $V\incl U\incl X$, take the canonical partitioned cover $P\onto X$ of \cref{rem:external-projectivity}.
The modest-target lifting construction gives a stream $\alpha$ for the pullback of $U$.
On the partitioned hit assembly $H_\alpha$ it similarly lifts the classifier of $V$ to a stream $\beta$.
\item \textbf{Evaluation and transposition.} Define
    \[
        \gamma_p(\langle n,k\rangle)=
        \begin{cases}
            \beta_{(p,n)}(k),&\alpha_p(n)=1,\\
            0,&\alpha_p(n)=0.
        \end{cases}
    \]
    Guarded evaluation makes this a total tracked binary stream: the first branch is used only on a valid hit name.
Its hit image in $P$ is the pullback of $V$.
Hence $ie\gamma$ is the pullback of the characteristic map of $V\incl X$.
Monicity of $i$ implies that $e\gamma$ equalizes the kernel pair of $P\onto X$, so it descends to a $\Sigma$-classifier on $X$.
\item \textbf{Positive existential search over $\N$.} For overtness, take a canonical partitioned cover $P\onto\Sigma^{\N}$.
Apply \cref{lem:modest-target-lift} to the evaluation map on $P\times\N$ and the modest stream cover $e$, giving streams $\alpha_{p,n}$.
The stream $\gamma_p(\langle n,k\rangle)=\alpha_{p,n}(k)$ classifies existential quantification along $\N$ after pullback to $P$.
The same descent argument gives $\exists_\N^\Sigma$.
\end{enumerate}
\end{proof}

\subsection{Finite execution certificates}
Finite execution evidence can be retained directly as an object of the topos. For the register-program syntax of \cref{subsec:numerical-foundation}, let $\mathsf{Run}$ be the decidable subset of $\N$ coding finite lists of configurations. Define
\[
 V(e,x,y,t)=1
\]
when $t$ parses as a run of program $e$, its first configuration is the prescribed input $x$, each consecutive pair obeys the one-step rule, and its final configuration halts with output $y$. Invalid encodings are rejected. Parsing, bounded list traversal, and the one-step rule are primitive recursive, so $V:\N^4\to2$ is a total morphism in $\Eff$. It checks a supplied finite execution; it does not decide whether an arbitrary program eventually halts.

\begin{theorem}[Operational certificates and exact execution support]\label{thm:operational-certificates}
The accepted-run family
\[
 \operatorname{Cert}(e,x,y)=\{t:\N\mid V(e,x,y,t)=1\}
\]
has the exact support
\[
 \exists t\,V(e,x,y,t)=1\quad\Longleftrightarrow\quad
 \varphi_e(x)=y.
\]
Its existential support over $y,t$ is the $\Sigma$-classified domain of universal evaluation.
\end{theorem}
\begin{proof}
\leavevmode
\begin{enumerate}[label=\textup{(\roman*)},leftmargin=*,itemsep=3pt,topsep=4pt]
\item \textbf{Finite witness construction.} An accepted list is a computation of the deterministic register interpreter, by induction on its length.
Conversely a halting computation has a finite list of successive configurations, which the checker accepts.
Coding this list and decoding a bounded simulation witness give computable translations between these witnesses and those of \eqref{eq:universal}.
Pair $y$ and $t$ to obtain a binary stream whose $\langle y,t\rangle$th bit is $V(e,x,y,t)$.
Its finite-hit image is precisely the evaluation domain, uniformly in $e,x$.
\end{enumerate}
\end{proof}

Overtness concerns a uniform existential operation on the actual presentation; an externally countable carrier alone does not supply it.


The realizability calculations supply the two operations needed below: finite-hit witnesses compose, and numerical existential search over $\N$ is positive. The partial-map classifier can therefore retain exact domains of computation without making termination decidable.

\section{Classified partiality and its universal span}\label{sec:classified-partiality}

A positive truth value can specify when a computation returns without deciding whether it returns. We turn this extent into a partial-map classifier, so that composition is defined by pullback and retains the returned value on its exact domain. The same construction will later interpret numerical execution and finite-return feedback.


\begin{definition}
\label{def:partial-map}
    A \textbf{$\Sigma$-partial map} $X\pmap Y$ is an isomorphism class of spans
    \[
        X\xleftarrow{m}U\xrightarrow{f}Y
    \]
    with $m$ a $\Sigma$-classified mono. Composition is by pullback. Write $\Par_\Sigma(\E)$ for the resulting category.

    The \textbf{partial-map classifier} is the object
\begin{equation}\label{eq:lift}
    \LL Y=\sum_{p:\Sigma}Y^p.
\end{equation}
An element is a semidecidable \textbf{extent} $p$ together with a value defined under $p$. This dependent sum exists in the topos; equivalently one can use the object of subterminal subobjects of $Y$ whose support is $\Sigma$-classified.

\end{definition}

\subsection{The universal domain of definition}
\label{subsec:universal-domain-of-definition}

The description of $\LL$ can be expressed without generalized-element notation. Let
\[
    \mathsf D_Y=\{(z,y)\in \LL Y\times Y\mid y\in z\},
    \qquad d_Y:\mathsf D_Y\incl \LL Y,
\]
where $z$ is viewed as a subterminal subobject of $Y$. The projection $d_Y$ is monic because a subterminal subobject has at most one member, and its classifier is the support of $z$. The second projection $\ev_Y:\mathsf D_Y\to Y$ is the universal value map. For $f:X\to \LL Y$, pulling back this pair gives its domain and value:
\[
    \begin{tikzcd}
        X\times_{\LL Y}\mathsf D_Y\ar[r]\ar[d,hook]
        &\mathsf D_Y\ar[d,hook,"{d_Y}"]\ar[r]&Y\\
        X\ar[r,"{f}"']&\LL Y.&
    \end{tikzcd}
\]


\begin{proposition}
\label{prop:partialmonad}
    There are natural bijections
    \[
        \E(X,\LL Y)\cong\Par_\Sigma(\E)(X,Y).
    \]
    They identify partial-map composition with the Kleisli composition of a strong monad $\LL$. This is the dominance-classified partiality construction \cite{rosolini,escardoknapp}.
\end{proposition}
\begin{proof}
\leavevmode
\begin{enumerate}[label=\textup{(\roman*)},leftmargin=*,itemsep=3pt,topsep=4pt]
\item \textbf{Classification by the universal span.} Pullback of the universal span $(d_Y,\ev_Y)$ gives a natural map from $\E(X,\LL Y)$ to classified partial spans.
Conversely a span $(m:U\incl X,u:U\to Y)$ determines the functional relation
\[
 R(x,y)\iff\exists a:U\,[m(a)=x\wedge u(a)=y].
\]
Its fibres are subterminal and its support is the classifier of $m$.
It therefore determines a unique map $X\to \LL Y$.
The two constructions recover the same domain mono and value map; the pullback universal property proves that they are inverse and natural in parameters.
\item \textbf{Partiality and definedness.} If $f$ has extent $p(x)$ and value $u_x:p(x)\to Y$, and $g$ has extent $q(y)$ and value $v_y:q(y)\to Z$, their pullback composite has extent
\begin{equation}\label{eq:partial-composite-extent}
 \dom(gf)(x)=\exists a:p(x)\;q(u_x(a))
\end{equation}
and value $v_{u_x(a)}(b)$.
Dominance supplies the positive classifier, and subterminality makes the value independent of its witness notation.
Pullback pasting identifies either threefold composite with the same domain and value; its associativity coherence follows from uniqueness of the maps preserving the projections.
Total identity spans give the units.
\item \textbf{Monad compatibility.} Under the classification, the unit is $y\mapsto(\top,\lambda *.y)$ and multiplication is the same flattening of dependent extents.
Thus the span identities are the monad identities.
The strength is
\[
 (x,(p,v))\longmapsto(p,\lambda a.(x,v(a))).
\]
It classifies product with a total parameter.
Products preserve the displayed pullbacks, so retaining parameters commutes with units and with flattening; uniqueness of the classifying map gives the strength coherence equations.
\end{enumerate}
\end{proof}

\begin{remark}
    The defined part of $\LL Y$ is canonically $Y$, but the whole object is not generally $1+Y$. A partial element can have an extent for which neither truth nor falsity is known.
\end{remark}

\begin{proposition}[Restriction and substitution]\label{prop:restrictions}
    For a partial map $f=(m,u):X\pmap Y$, let $\overline f=(m,m):X\pmap X$. Then $\Par_\Sigma(\E)$ is a restriction category in the sense of \cite{cockettlack}. The same restriction equations hold in $\Par(\E)$ for arbitrary mono domains. Moreover the partial-map classification is stable under substitution of parameters, and the strength of $\LL$ classifies pairing a total parameter with a partial value.
\end{proposition}
\begin{proof}
\leavevmode
\begin{enumerate}[label=\textup{(\roman*)},leftmargin=*,itemsep=3pt,topsep=4pt]
\item \textbf{Unit and identity laws.} The equation $f\overline f=f$ follows from pulling the domain mono back along itself.
Restriction idempotents with a common source commute because their composite is the intersection of their domains with identity value.
Likewise $\overline{g\overline f}=\overline g\,\overline f$ for maps with a common source.
For composable $f,g$, the maps $\overline g f$ and $f\overline{gf}$ have the domain of $gf$ and the value of $f$ there.
These are the restriction identities.
\item \textbf{Restriction stability of domain monos.} They use only pullbacks of monos, hence apply to arbitrary mono domains; dominance closes the selected class of $\Sigma$-domains under composition.
Parameter substitution and strength are the corresponding pullback and product operations in \cref{prop:partialmonad}.
\end{enumerate}
\end{proof}


The universal span makes the extent and value of a partial computation independent of a chosen representative. Its composition equation will reappear in the numerical interpreter and, with finite returning histories, in Elgot iteration.

\section{Overtness and positive relational calculus}
\label{sec:overt}

Finite-hit predicates can be conjoined, but composing positive relations requires an existential search over their intermediate object. Overtness is exactly the property that makes this search positive in every parameter context. Its closure and Beck--Chevalley laws supply the relational calculus used for approximation matrices.


\begin{lemma}
\label{lem:overtclosure}
    Finite discrete objects are overt. Overt objects are closed under finite products, finite sums, and regular-epimorphic images. A dependent sum of a relatively overt family over an overt object is overt. A $\Sigma$-classified subobject of an overt object is overt.
\end{lemma}
\begin{proof}
\leavevmode
\begin{enumerate}[label=\textup{(\roman*)},leftmargin=*,itemsep=3pt,topsep=4pt]
\item \textbf{Product calculation.} For finite objects, existential quantification is a finite disjunction.
For products and dependent sums, use iterated restricted existential quantification.
For sums use disjunction.
If $q:X\onto Y$ and $X$ is overt, then
    \[
        \exists y:Y\,u(y)\iff\exists x:X\,u(qx),
    \]
    so the right side gives the required restriction.
For a classified subobject $U=\{x\mid p(x)\}$ and $v:U\to\Sigma$, dominance makes $\exists a:p(x)\,v(x,a)$ a $\Sigma$-predicate on $X$; then quantify over $X$.
\end{enumerate}
\end{proof}

A \textbf{$\Sigma$-relation} from $X$ to $Y$ is a map $X\times Y\to\Sigma$. If the middle object $Y$ is overt, relational composition restricts to
\begin{equation}\label{eq:sigmacomposition}
    (S R)(x,z)=\exists_Y^\Sigma y\,[R(x,y)\wedge S(y,z)].
\end{equation}
Associativity follows from the corresponding equality in $\Omega$ and monicity of $i$. The order and equality predicates continue to take values in $\Omega$. In particular, the appropriate identities in the rounded presentation category will be approximation relations.

\subsection{Dependent quantifiers and parameter coherence}\label{app:quantifiers}
The positive quantifier calculus is the restriction of the subobject adjunctions. All constructions below take place in arbitrary parameter contexts.

A map is \textbf{relatively $\Sigma$-compact} when its dependent universal quantifier restricts to positive predicates in every slice. For a map $p:E\to B$, pullback on subobjects has adjoints
\[
 \exists_p\dashv p^*\dashv\forall_p.
\]
The left adjoint is image of a composite. The right adjoint sends a predicate $R$ to the predicate that every element in the fibre satisfies $R$. These adjoints exist in a topos. In a pullback square
\[
\begin{tikzcd}
 E'\ar[r,"g'"]\ar[d,"p'"']&E\ar[d,"p"]\\
 B'\ar[r,"g"']&B
\end{tikzcd}
\]
one has
\begin{equation}\label{eq:bc-predicates}
 g^*\exists_p=\exists_{p'}(g')^*,
 \qquad
 g^*\forall_p=\forall_{p'}(g')^*.
\end{equation}

\begin{lemma}[Substitution and Frobenius]\label{lem:predicate-bc}
Equations \eqref{eq:bc-predicates} hold on all predicates. Also
\[
 \exists_p(R\wedge p^*S)=\exists_p(R)\wedge S.
\]
When the quantifiers restrict to the positive fragment, their restricted maps satisfy the same equations.
\end{lemma}
\begin{proof}
\leavevmode
\begin{enumerate}[label=\textup{(\roman*)},leftmargin=*,itemsep=3pt,topsep=4pt]
\item \textbf{Substitution and frobenius.} At $b'\in B'$, the two existential predicates say respectively that there is $e$ with $p(e)=g(b')$ satisfying $R$, or that there is $(b',e)\in E'$ satisfying its pulled-back predicate.
The pullback identifies these witnesses.
The two universal predicates quantify over exactly that same identified fibre, giving the second equation.
\item \textbf{Inverse comparison.} Frobenius is the equivalence
\[
 \exists e\,[p(e)=b\wedge R(e)\wedge S(b)]
 \iff(\exists e\,[p(e)=b\wedge R(e)])\wedge S(b).
\]
The forward implication projects witnesses; the reverse uses the supplied existential witness and the proof of $S(b)$.
These are internal equivalences, so identify subobjects.
\item \textbf{Reflection of positive equations.} If all predicates and restricted quantifier values lie in $\Sigma$, their composites with $i:\Sigma\incl\Omega$ are the full operations just calculated.
Monicity of $i$ therefore reflects the equations of their positive classifiers.
\end{enumerate}
\end{proof}

The square here is a pullback of objects \emph{within the current topos}. Preservation of a universal quantifier by an arbitrary geometric inverse-image functor is a different statement and is not implied by \eqref{eq:bc-predicates}.


\subsection{The laws of restricted existential quantification}
\label{subsec:laws-of-restricted-existential-quantification}

Overtness supplies an operation, not a selection of witnesses. Since its composite with $i$ is ordinary existential quantification, all equations between restricted existential expressions can be proved in $\Omega$ and reflected through the mono $i$. We record the laws used in relational composition and in observation closure.

\begin{proposition}
\label{prop:overtlaws}
    Restricted existential quantification satisfies substitution in parameters, Frobenius reciprocity, and Fubini's law whenever the indicated indexing objects are overt. In particular, for $u:S\times X\to\Sigma$, $v:S\to\Sigma$, and overt $X,Y$,
    \begin{align*}
        (\exists_X^\Sigma u)(f(t))
        &=\exists_X^\Sigma x\,u(f(t),x),\\
        v(s)\wedge\exists_X^\Sigma x\,u(s,x)
        &=\exists_X^\Sigma x\,[v(s)\wedge u(s,x)],\\
        \exists_X^\Sigma x\,\exists_Y^\Sigma y\,w(s,x,y)
        &=\exists_{X\times Y}^\Sigma(x,y)\,w(s,x,y).
    \end{align*}
    The same statements hold for relatively overt objects in a slice.
\end{proposition}
\begin{proof}
\leavevmode
\begin{enumerate}[label=\textup{(\roman*)},leftmargin=*,itemsep=3pt,topsep=4pt]
\item \textbf{Beck--Chevalley and Frobenius.} The image adjunction gives Frobenius and pullback Beck--Chevalley; the right adjoint also satisfies the latter, as proved in \cref{lem:predicate-bc}.
Composition of image adjoints gives Fubini.
Each restricted expression has its displayed positive classifier by overtness and finite closure, and monicity of $i$ reflects these equations.
The same adjunctions in a slice give the dependent assertions.
\end{enumerate}
\end{proof}

For example, if $R:X\times Y\to\Sigma$, $S:Y\times Z\to\Sigma$, and $T:Z\times W\to\Sigma$, with $Y,Z$ overt, both composites of three relations are the predicate
\[
    \exists_{Y\times Z}^\Sigma(y,z)\,
    [R(x,y)\wedge S(y,z)\wedge T(z,w)].
\]
The formula also shows that composition is monotone and commutes with substitution in every unused variable. If an identity relation is available in $\Sigma$, it acts as an identity by equality substitution. 

An overt object need not have decidable equality, and an overt subobject of an overt object need not be $\Sigma$-classified. The closure results concern the direction required by the formulas: a classified subobject of an overt object supports restricted existential quantification.

\begin{proposition}[Closure of positive quantifier maps]\label{prop:quantifier-closure}
Relatively overt maps are closed under composition and pullback. The same is true for relatively compact maps. A regular epimorphism with overt domain has overt codomain. A regular epimorphism with compact domain has compact codomain, with the given compact quantifier structure. These assertions hold relatively over a parameter base. When a restricted quantifier is not total, its maximal positive domain is constructed in \cref{thm:maximal-positive-quantifiers}; none of the present closure claims turns a partial restricted quantifier into a total one.
\end{proposition}
\begin{proof}
\leavevmode
\begin{enumerate}[label=\textup{(\roman*)},leftmargin=*,itemsep=3pt,topsep=4pt]
\item \textbf{Closure of positive quantifier maps.} For composable maps $p,q$, the adjunctions give
$\exists_{qp}=\exists_q\exists_p$ and
$\forall_{qp}=\forall_q\forall_p$.
Thus the composite restricted maps construct the positive quantifiers.
Stability under pullback is part of the relative definitions and agrees with \cref{lem:predicate-bc}.
\item \textbf{Existence and uniqueness.} If $e:A\onto B$, then
$\exists b\,R(b)\iff\exists a\,R(ea)$.
One implication substitutes a witness.
For the other, use epimorphy to lift the given $b$ locally and eliminate that existential lift.
Hence overtness of $A$ supplies the positive existential classifier on $B$.
Similarly
$\forall b\,R(b)\iff\forall a\,R(ea)$:
epimorphy makes a predicate containing all the images contain all of $B$.
Compactness of $A$ therefore supplies the positive universal classifier on $B$.
The maps are uniform in $R$, and uniqueness follows from the monic classifier.
Perform the same calculation in a slice for the relative assertion.
\end{enumerate}
\end{proof}

\begin{lemma}[Finite dependent conjunction]\label{lem:finite-compact}
For the family $\Fin(n)\to1$, and uniformly for $n$ varying in any context, universal quantification restricts to $\Sigma$. Finite decidable families are consequently $\Sigma$-compact.
\end{lemma}
\begin{proof}
\leavevmode
\begin{enumerate}[label=\textup{(\roman*)},leftmargin=*,itemsep=3pt,topsep=4pt]
\item \textbf{Finite conjunction with positive values.} Define the object of finite families as $\sum_{n:\N}\Sigma^{\Fin(n)}$.
Extend a finite family $a$ to a total stream of elements of $\Sigma$ by the value $\top$ outside its decidable domain.
Recursion defines
\[
 \operatorname{all}_\Sigma(0,a)=\top,\qquad
 \operatorname{all}_\Sigma(n+1,a)
 =\operatorname{all}_\Sigma(n,a|_n)\wedge a(n).
\]
Induction gives
$i(\operatorname{all}_\Sigma(n,a))\iff\forall k<n\,i(a(k))$.
The construction is in the slice of its parameters and uses only finite conjunction.
\item \textbf{Naturality and reindexing.} Pulling a displayed finite family back preserves its finite indexing and this formula defined by recursion.
It therefore supplies the required restricted universal quantifier after every parameter substitution.
\end{enumerate}
\end{proof}

\begin{proposition}[Finite positive quantification]\label{prop:finite-quantifiers}
For a displayed finite decidable family, both existential and universal quantification are positive, uniformly in its finite length. Finite products of overt objects are overt, and finite iterated dependent sums of relatively overt families are overt, by successive existential projection.
\end{proposition}
\begin{proof}
\leavevmode
\begin{enumerate}[label=\textup{(\roman*)},leftmargin=*,itemsep=3pt,topsep=4pt]
\item \textbf{Product calculation.} Finite existential quantification is recursion from $\bot$ using binary joins; the correctness proof is induction in the context of the finite family.
The dual recursion for universal quantification is \cref{lem:finite-compact}.
Finite products and iterated dependent sums follow by composing the restricted existential adjoints of \cref{prop:quantifier-closure}; their Fubini equations are \cref{prop:overtlaws}.
Thus all classifiers are uniform in the finite length and parameters.
\end{enumerate}
\end{proof}


The positive calculus now supports substitution, finite conjunction, and existential composition over overt intermediates. These are the operations needed to compose rounded approximation relations; the underlying order and full ideal construction remain $\Omega$-valued.

\part{Enriched category theory over the subobject classifier}
\label{part:enrichment}

Positive observations will eventually restrict which ideal suprema can be computed. Before imposing that restriction, we need the unrestricted order-theoretic construction: internal preorders, distributors, ideal completion, and approximation. Working over $\Omega$ makes their universal properties available without presupposing semidecidability of the order.


\section{Internal enrichment, variance, and flat weights}
\label{sec:omegainternal}

The forthcoming completion is best understood through its interaction with upper predicates. We therefore fix the variance of distributors and their existential composition before constructing ideal weights. This keeps the order-valued universal properties separate from the positive tests introduced earlier.


The internal truth values form a complete Heyting algebra \cite{johnstone}. Conjunction is a commutative tensor with unit $\top$, and implication satisfies $p\wedge q\leq r$ if and only if $p\leq(q\Rightarrow r)$. Existentially indexed joins distribute over finite meets. Completeness here refers to internally indexed families: for a predicate $S\hookrightarrow\Omega$, its join is $\exists p:S\,p$. We use this internal monoidal order as the enrichment base.

\begin{definition}[Internal preorders and separation]
\label{def:internalpreorder}
    An internal preorder is an object $P$ with a predicate $P(p,q)$ satisfying reflexivity and transitivity. It is \textbf{separated}, or an internal poset, when $P(p,q)\wedge P(q,p)$ implies $p=q$. A map $f:P\to Q$ is monotone when $P(p,q)\Rightarrow Q(fp,fq)$. The internal object of such maps is the corresponding subobject of $Q^P$ and has pointwise order.
\end{definition}

The enriched identity and composition conditions are exactly reflexivity and transitivity \cite{stubbe}. When we later say that a hom is \textbf{internally represented}, we mean that the relevant family of arrows is a subobject of an exponential in $\E$. 

The equivalence relation $p\sim q$ given by $p\leq q\leq p$ has an effective quotient $P_{\mathrm{sep}}$. Its order is induced by that of $P$, and monotone maps into posets factor uniquely through the quotient. Lower predicates, upper predicates, and ideals are invariant under $\sim$ and correspond to those of $P_{\mathrm{sep}}$. 

\begin{definition}[Distributors]
\label{def:distributor}
    A \textbf{distributor} $\Phi:P\nrightarrow Q$ is a predicate $\Phi(q,p)$ which is lower in $q$ and upper in $p$. Thus
    \[
        q'\leq q,\quad\Phi(q,p),\quad p\leq p'
        \quad\Longrightarrow\quad\Phi(q',p').
    \]
    The identity has matrix $1_P(p',p)=P(p',p)$. Composition and order are
    \[
        (\Psi\Phi)(r,p)=\exists q:Q\,[\Psi(r,q)\wedge\Phi(q,p)],
        \qquad\Phi\leq\Phi'\ \Longleftrightarrow\
        \forall q,p\,[\Phi(q,p)\Rightarrow\Phi'(q,p)].
    \]
\end{definition}

The identity equations follow from monotonicity and reflexivity; associativity is reassociation of the existential witnesses. Every hom-order of distributors is a complete lattice, with pointwise joins, and composition preserves those joins. For a monotone $f:P\to Q$, its graph and cograph distributors are
\[
    f_*(q,p)=Q(q,fp),\qquad f^*(p,q)=Q(fp,q).
\]
They satisfy $1_P\leq f^*f_*$ and $f_*f^*\leq1_Q$, so $f_*\dashv f^*$.

For $\Phi:P\nrightarrow Q$ and $\Psi:P\nrightarrow R$, the right extension has matrix
\[
    (\Ran_\Phi\Psi)(r,q)
    =\forall p:P\,[\Phi(q,p)\Rightarrow\Psi(r,p)].
\]
It is a distributor, by the two variance conditions, and satisfies
\[
    \Theta\Phi\leq\Psi\quad\Longleftrightarrow\quad
    \Theta\leq\Ran_\Phi\Psi.
\]
Indeed the left inequality says that each conjunction $\Theta(r,q)\wedge\Phi(q,p)$ implies $\Psi(r,p)$; implication and universal quantification give precisely the right inequality. This will type the way-below construction of \cref{jd:eq:wayuniv}.

A lower predicate on $P$ is a distributor $1\nrightarrow P$; an upper predicate is a distributor $P\nrightarrow1$. Their composite is the pairing $\exists p[I(p)\wedge U(p)]$ \cite{stubbe}. A lower weight is called \textbf{flat} here when this pairing, as a function of upper predicates, preserves finite meets, including the nullary meet $\top$. 

\begin{definition}[Internal frames]
\label{def:internalframe}
    An \textbf{internal frame} is an internal poset $L$ with joins of all internal subobjects $S\hookrightarrow L$ and finite meets, satisfying $a\wedge\bigvee S=\bigvee_{s:S}(a\wedge s)$. These operations are morphisms, including $\bigvee:\Pow L\to L$, and satisfy their universal properties internally. A frame homomorphism preserves finite meets and these joins. A character of $L$ is such a homomorphism $L\to\Omega$.
\end{definition}

For an indexed family $u:X\to L$, its join is the join of its image. Therefore preservation of joins of subobjects is equivalent to preservation of all internally indexed joins. 


The distributor conventions identify ideals as lower weights tested against upper predicates. We can therefore obtain completion from a character construction and subsequently restrict its weights by requiring only their positive pairings to remain semidecidable.

\section{The ideal-completion doctrine}
\label{sec:ideal}

A forcing frame will ultimately arise from upper predicates on a condition preorder. To compare it with a domain frame, we first need the unrestricted ideal completion and its trace on those predicates. The character adjunction supplies both the completion monad and the formula for the Scott opens of its free algebras.


An internal preorder $P$ is an internally $\Omega$-enriched category, where $\Omega$ is regarded as an internal monoidal poset with tensor $\wedge$ and unit $\top$. Its hom-predicate is $P(p,q)\equiv p\leq q$. The identity and composition inequalities are
\[
 \top\leq P(p,p),\qquad
 P(p,q)\wedge P(q,r)\leq P(p,r).
\]
An enriched functor is exactly a monotone map, since its hom condition is $P(p,q)\leq Q(f(p),f(q))$. The hom-preorders of such maps are ordered pointwise.  $\Omega$ is the internal frame of truth values. Lower and upper predicates are the two variances of enriched presheaf.

Write $\Up(P)\subseteq\Omega^P$ for the upper predicates. A lower predicate $I$ is an ideal when
\begin{align*}
    &\exists p\,I(p),\\
    &I(p)\wedge I(q)\Longrightarrow
    \exists r\,[I(r)\wedge p\leq r\wedge q\leq r].
\end{align*}
These conditions define $\Idl(P)\subseteq\Omega^P$, ordered by inclusion.

\subsection{Upper predicates and characters}
\label{subsec:upper-predicates-and-characters}

Write $\mathbf{Frm}(\E)$ for internal frames and internal frame homomorphisms. For a frame $L$, let
\[
    \operatorname{Char}(L)=\{\chi:L\to\Omega\mid
               \chi\text{ preserves finite meets and arbitrary joins}\},
\]
ordered pointwise. The preservation condition is internal: arbitrary joins include joins indexed by every internal subobject of $L$. Evaluation at $p\in P$ is a character of $\Up(P)$, and these evaluations give a monotone map from $P$ to its character object.

\begin{proposition}[The upper-predicate dual adjunction]
\label{prop:characters}
    There is an adjunction
    \[
        \Up:\Pre(\E)\ \rightleftarrows\ \mathbf{Frm}(\E)^{\mathrm{op}}
                  :\operatorname{Char}.
    \]
    Its hom-set bijection sends $h:L\to\Up(P)$ to $\widehat h:P\to\operatorname{Char}(L)$ with $\widehat h(p)(a)=h(a)(p)$.
\end{proposition}
\begin{proof}
\leavevmode
\begin{enumerate}[label=\textup{(\roman*)},leftmargin=*,itemsep=3pt,topsep=4pt]
\item \textbf{The upper-predicate dual adjunction.} Upper predicates form an internal frame with pointwise meets and joins.
For monotone $f:P\to Q$, inverse image $f^*:\Up(Q)\to\Up(P)$ preserves these operations. A frame homomorphism $u:L\to M$ acts on characters by precomposition, giving the stated variances.
\item \textbf{Evaluation and transposition.} For fixed $p$, evaluation at $p$ preserves all frame operations, so $a\mapsto h(a)(p)$ is a character.
For $p\leq q$, upperness of each $h(a)$ gives $\widehat h(p)\leq\widehat h(q)$.
\item \textbf{Converse construction.} Conversely, if $k:P\to\operatorname{Char}(L)$ is monotone, then $a\mapsto(p\mapsto k(p)(a))$ has upper values and preserves frame operations pointwise.
These two transpositions are inverse.
Precomposing in $P$ substitutes $f(p)$ for $p$; precomposing a frame homomorphism substitutes $u(a)$ for $a$.
Both commute with the displayed formula, proving naturality and the adjunction.

The universal property is represented by the commutative diagram
\[
\begin{tikzcd}[column sep=large,row sep=large]
P\ar[r,"\widehat h"]\ar[dr,"{p\mapsto h(a)(p)}"']&\operatorname{Char}(L)\ar[d,"{\operatorname{ev}_a}"]\\ &\Omega
\end{tikzcd}
\]
\end{enumerate}
\end{proof}

The unit and counit can be read off from transposition. The unit is
\[
    e_P:P\to\Char(\Up P),\qquad e_P(p)(U)=U(p),
\]
and the counit at $L$, expressed as a frame map because the frame category is opposite, is
\[
 \epsilon_L^{\mathrm{frm}}:L\to\Up(\Char L),
 \qquad\epsilon_L^{\mathrm{frm}}(a)(\chi)=\chi(a).
\]
To check the first triangle, evaluate a predicate $U$ at $p$: it is sent to $\chi\mapsto\chi(U)$ and then evaluated at $e_P(p)$, giving $U(p)$. For the other triangle, a character $\chi$ evaluated at $a$ returns $\chi(a)$ through either route. Extensionality gives the two map equations.

\begin{proposition}[Characters are ideals]
\label{prop:characterideals}
    There is a natural order isomorphism
    \[
        \operatorname{Char}(\Up(P))\cong\Idl(P),
        \qquad I_\chi(p)=\chi(\uparrow p),\qquad
        \chi_I(U)=\exists p\,[I(p)\wedge U(p)].
    \]
\end{proposition}
\begin{proof}
\leavevmode
\begin{enumerate}[label=\textup{(\roman*)},leftmargin=*,itemsep=3pt,topsep=4pt]
\item \textbf{Ideal and directedness calculation.} Since $\top=\bigvee_{p:P}\uparrow p$, preservation of joins and top makes $I_\chi$ inhabited.
If $q\leq p$, then $\uparrow p\subseteq\uparrow q$, so $I_\chi(p)$ implies $I_\chi(q)$.
For directedness use
    \[
        (\uparrow p)\cap(\uparrow q)
        =\bigvee_{r\,:,p\leq r\,\wedge\,q\leq r}\uparrow r.
    \]
    Preservation of this join and of the meet says that $I_\chi(p)\wedge I_\chi(q)$ implies the existence of $r\in I_\chi$ above both.
Thus $I_\chi$ is an ideal.
\item \textbf{Converse construction.} Conversely, inhabitedness gives $\chi_I(\top)=\top$, and existential quantification commutes with arbitrary unions of upper predicates.
If $\chi_I(U)$ and $\chi_I(V)$, choose their witnesses $p,q\in I$.
Directedness supplies $r\in I$ above both; upperness gives $U(r)\wedge V(r)$, proving preservation of binary meets.
The reverse implication follows by projection, so $\chi_I$ is a character.
\item \textbf{Existence and uniqueness.} Every upper predicate has its principal decomposition
    \[
        U=\bigvee_{p\,:,U(p)}\uparrow p.
    \]
    Applying $\chi$ proves $\chi_{I_\chi}(U)=\chi(U)$.
In the other direction, $I_{\chi_I}(p)\iff\exists q\,[I(q)\wedge p\leq q]\iff I(p)$.
The formulas preserve and reflect the pointwise orders.
\item \textbf{Naturality under monotone maps.} Finally, precomposition with $f^*$ gives at $\uparrow q$ the predicate $\exists p\,[I(p)\wedge q\leq f(p)]$, which is exactly the lower image along $f$.
This proves naturality.

The universal property is represented by the commutative diagram
\[
\begin{tikzcd}[column sep=large,row sep=large]
\Char(\Up P)\ar[r,"{\chi\mapsto I_\chi}"]\ar[dr,"\id"']&\Idl P\ar[d,"{I\mapsto\chi_I}"]\\ &\Char(\Up P)
\end{tikzcd}
\]
\end{enumerate}
\end{proof}

The unit of the dual adjunction sends $p$ to evaluation at $p$, hence to $\downarrow p$. Its counit, written in frame direction, sends $a\in L$ to the upper predicate $\chi\mapsto\chi(a)$. Consequently its induced multiplication sends an ideal $\mathcal I$ of ideals to the predicate
\[
 p\longmapsto\exists I\,[\mathcal I(I)\wedge I(p)].
\]

\subsection{The completion and its algebras}
\label{subsec:completion}

\begin{proposition}
\label{prop:idealmonad}
    Ideal completion defines a locally monotone monad $\T$ on $\Pre(\E)$, with
    \begin{align}
        \T P&=\Idl(P), & y_P(p)&=\downarrow p,\nonumber\\
        (\T f)(I)(q)&\iff\exists p\,[I(p)\wedge q\leq f(p)],\label{eq:idealmonad}\\
        \mu_P(\mathcal I)(p)&\iff\exists I\,[\mathcal I(I)\wedge I(p)].\nonumber
    \end{align}
    It is Kock--Z\"oberlein: $\T y_P\leq y_{\T P}$.
\end{proposition}
\begin{proof}
\leavevmode
\begin{enumerate}[label=\textup{(\roman*)},leftmargin=*,itemsep=3pt,topsep=4pt]
\item \textbf{Ideal and directedness calculation.} The lower image of an ideal is inhabited and directed.
The union of an ideal of ideals is inhabited; two of its elements lie in two member ideals, which have a common upper member, and directedness inside that member gives a common upper element.
Thus all displayed maps are defined.
\item \textbf{Unit and identity laws.} Lower closure gives $\T(\id)=\id$ and the two unit equations.
Associativity of $\mu$ expands on either side to the same iterated existential predicate.
The composition law for $\T$ follows by eliminating the intermediate lower-bound witness.
\item \textbf{The KZ inequality.} Finally $\T y_P(I)$ consists of ideals below a principal ideal $\downarrow p$ with $p\in I$.
Every such ideal is contained in $I$, which is precisely the KZ inequality.

Naturality is expressed by the commutative diagram
\[
\begin{tikzcd}[column sep=large,row sep=large]
P\ar[r,"y_P"]\ar[d,"f"']&\Idl P\ar[d,"\T f"]\\ Q\ar[r,"y_Q"']&\Idl Q
\end{tikzcd}
\]
\end{enumerate}
\end{proof}

\begin{theorem}
\label{thm:ordinaryalgebras}
    The Eilenberg--Moore category of $\T$ is the category of internal dcpos and Scott-continuous maps. Its algebra map is ideal supremum, and
    \[
        \sup_D\dashv y_D.
    \]
\end{theorem}
\begin{proof}
\leavevmode
\begin{enumerate}[label=\textup{(\roman*)},leftmargin=*,itemsep=3pt,topsep=4pt]
\item \textbf{Construction of the comparison.} Let $a:\Idl(D)\to D$ be an algebra.
From $a(\downarrow x)=x$, equivalent elements of the preorder are equal, so $D$ is separated.
If $x\in I$, then $\downarrow x\subseteq I$, hence $x\leq aI$.
If $x$ is an upper bound of $I$, then $I\subseteq\downarrow x$, hence $aI\leq x$.
Therefore $aI=\sup I$ and $a\dashv y_D$.
\item \textbf{Converse construction.} Conversely, ideal suprema give the unit equation.
Both $\sup\,\mu(\mathcal I)$ and $\sup\,\T(\sup)(\mathcal I)$ are the least upper bound of all elements of all ideals in $\mathcal I$, giving the algebra equation.
Inhabited directed families and ideals have the same suprema, by taking lower images.
Finally the homomorphism equation is exactly preservation of ideal suprema.

The comparison is witnessed by the retraction diagram
\[
\begin{tikzcd}[column sep=large,row sep=large]
D\ar[r,"y_D"]\ar[dr,"\id_D"']&\Idl D\ar[d,"\sup_D"]\\ &D
\end{tikzcd}
\]
\end{enumerate}
\end{proof}


\begin{proposition}[Free directed completion]
\label{prop:freeextension}
    For a preorder $P$ and a dcpo $D$, restriction along $y_P$ is a natural order isomorphism
    \[
        \Dcpo(\E)(\Idl P,D)\cong\Pre(\E)(P,D).
    \]
    The extension of $f:P\to D$ is $f^\#(I)=\sup_D\T f(I)$.

    Equivalently, $f^\#$ is the unique Scott-continuous filler in
    \[
        \begin{tikzcd}[column sep=large]
            P\ar[r,"{y_P}"]\ar[dr,"{f}"']&\Idl P\ar[d,dashed,"{f^\#}"]\\
            &D.
        \end{tikzcd}
    \]
\end{proposition}
\begin{proof}
\leavevmode
\begin{enumerate}[label=\textup{(\roman*)},leftmargin=*,itemsep=3pt,topsep=4pt]
\item \textbf{Free directed completion.} The lower image of $I$ along $f$ is an ideal, and its supremum is the supremum of the directed family $\{f(p)\mid p\in I\}$. A directed union of ideals gives the same family of values whether it is flattened first or its members are evaluated first.
Thus $f^\#$ preserves directed suprema.
On a principal ideal its value is $f(p)$, since $f(p)$ is the largest value being joined.
\item \textbf{Converse construction.} Conversely every ideal is the directed union of its principal subideals. A Scott-continuous extension must therefore have the displayed value on every ideal, proving uniqueness.
Pointwise inequalities between $f$ and $g$ give inequalities between these suprema, and inequalities between extensions restrict along $y_P$.
This proves the order isomorphism; naturality follows from uniqueness of extension.
\end{enumerate}
\end{proof}

\paragraph{Kleisli extension and convolution.}
A Kleisli arrow $m:B\to\Idl(C)$ is a monotone family of ideals. We write $m(b,c)$ for $c\in m(b)$. Its extension and Kleisli composite are
\begin{align}
    m^\#(I)(c)&\iff\exists b\,[I(b)\wedge m(b,c)],\label{eq:kleisli}\\
    (n\star m)(b,d)&\iff\exists c\,[m(b,c)\wedge n(c,d)].\nonumber
\end{align}
The relation is monotone in the source variable and lower in the target variable. By \cref{prop:freeextension}, extension preserves directed unions, $m^\#(\downarrow b)=m(b)$, and every Scott-continuous map between free completions is obtained this way. These formulas identify the Kleisli category with the full category of free ideal completions.


These are the Kock--Z\"oberlein properties of directed completion; see \cite{hofmannsousa,hofmannwaszkiewicz} for the general categorical setting.

\subsection{Function objects in the underlying domain category}
\label{subsec:function-objects-and-domain-category}

We will need pointwise directed suprema of maps. Let $[D,E]$ be the internal object of Scott-continuous maps from $D$ to $E$, with pointwise order. This is a subobject of $E^D$, since preservation of ideal suprema is an internal formula.

\begin{proposition}
\label{prop:ordinaryexponential}
    The category of internal dcpos and Scott-continuous maps is cartesian closed \cite{hofmannwaszkiewicz}. Its exponential is $[D,E]$, whose directed suprema are pointwise.
\end{proposition}
\begin{proof}
\leavevmode
\begin{enumerate}[label=\textup{(\roman*)},leftmargin=*,itemsep=3pt,topsep=4pt]
\item \textbf{Ideal and directedness calculation.} Given an inhabited directed family $f_j$, define $f(x)=\bigvee_j f_j(x)$.
For a directed family $x_k$, both iterated suprema
    \[
        \bigvee_j\bigvee_k f_j(x_k),\qquad
        \bigvee_k\bigvee_j f_j(x_k)
    \]
    are the supremum of the same inhabited directed family indexed by pairs $(j,k)$.
Hence $f$ is Scott-continuous.
This proves directed completeness of the hom-object.
\item \textbf{Scott continuity.} Evaluation is Scott-continuous jointly.
For a directed family $(f_j,x_j)$ in $[D,E]\times D$, evaluating its componentwise supremum gives the supremum of all $f_j(x_k)$.
For any two indices a common later member bounds both coordinates, so the diagonal family $f_l(x_l)$ is cofinal in this family.
Its supremum is therefore the same, as required.
\item \textbf{Converse construction.} Conversely, for Scott-continuous $h:C\times D\to E$, each $h(c,-)$ is Scott-continuous, and preservation of directed suprema in the first variable proves that $c\mapsto h(c,-)$ is a Scott-continuous map into $[D,E]$.
The evaluation and transposition equations hold on the underlying maps and give the exponential universal property.

The defining arrows form the commutative diagram
\[
\begin{tikzcd}[column sep=large,row sep=large]
C\times D\ar[r,"{\lambda h\times\id_D}"]\ar[dr,"h"']&{[D,E]\times D}\ar[d,"\ev"]\\ &E
\end{tikzcd}
\]
\end{enumerate}
\end{proof}


For $U\in\Up(P)$ put
\begin{equation}\label{eq:dagger}
    U^\dagger(I)\iff\exists p\,[I(p)\wedge U(p)].
\end{equation}
A Scott open is an upper predicate inaccessible by directed suprema \cite{hofmannwaszkiewicz}. Equivalently it is a Scott-continuous map to $\Omega$ with its directed-complete order.

\begin{lemma}[Trace and extension]
\label{lem:tracefree}
    Restriction along $y_P$ gives an isomorphism
    \[
        \Sc(\Idl P)\cong\Up(P),
    \]
    with inverse $U\mapsto U^\dagger$. Extension preserves finite meets and arbitrary joins of upper predicates.
\end{lemma}
\begin{proof}
\leavevmode
\begin{enumerate}[label=\textup{(\roman*)},leftmargin=*,itemsep=3pt,topsep=4pt]
\item \textbf{Trace and extension.} Every ideal has the directed-union expression $I=\bigcup_{p\in I}\downarrow p$. A Scott open is therefore recovered from its values on the principal ideals, giving the two inverse identities.
Formula \eqref{eq:dagger} preserves joins by interchange of existential quantifiers.
\item \textbf{Finite-meet preservation.} To prove preservation of binary meets, witnesses $p,q\in I$ for $U^\dagger(I)$ and $V^\dagger(I)$ have a common upper element $r\in I$, and upperness yields $U(r)\wedge V(r)$.
Inhabitedness handles the nullary meet.
\end{enumerate}
\end{proof}

The same calculation expresses ideals as flat weights: $I$ is an ideal exactly when $U\mapsto\exists p\,I(p)\wedge U(p)$ preserves finite meets, including the top predicate. The converse uses the principal upper predicates $\uparrow p$ and $\uparrow q$.


The unrestricted doctrine provides three reusable interfaces: principal ideals embed a preorder, union is ideal multiplication, and Scott observations of free completions are upper predicates on the generators. The positive doctrine will restrict the first two while the observed doctrine retains the third.

\section{Ordinary approximation, compactness, and splittings}
\label{sec:ordinaryapproximation}

Ideal completion is algebraic, whereas domains appearing as retracts of free completions need not be. The way-below relation identifies what remains approximable under such retracts. Its adjoint formulation is the form that can later be transferred to the admitted and observation-equipped doctrines.


\begin{definition}[Way-below and continuity]
\label{def:ordinarycontinuous}
    For $a,x:D$, define
    \[
        a\way x\quad\Longleftrightarrow\quad
        \forall I:\Idl D\,[x\leq\sup I\Rightarrow a\in I].
    \]
    Write $\Downarrow x=\{a\mid a\way x\}$. The dcpo is continuous when every $\Downarrow x$ is an inhabited directed ideal with supremum $x$. An element $k$ is compact when $k\way k$. A dcpo is algebraic when its compact elements below each $x$ form an inhabited directed family with supremum $x$.
\end{definition}

This says that for a directed family $(x_i)$ that $x\leq\bigvee_i x_i$ implies $a\leq x_i$ for some $i$. One direction applies the definition to the family's lower image. The other applies the family form to the identity family on an ideal and uses lower closure.

\begin{proposition}[Approximation adjunction]
\label{prop:ordinaryapproxadjunction}
    A dcpo is continuous if and only if its supremum map has a left adjoint. In that case the adjoint is $\Downarrow$ and
    \[
        \Downarrow\dashv\sup\dashv y_D,\qquad
        \sup\Downarrow=\id_D,\qquad\Downarrow\sup\leq\id_{\Idl D}.
    \]
    Moreover $\Downarrow$ is Scott-continuous and preserves ideal suprema.
\end{proposition}
\begin{proof}
\leavevmode
\begin{enumerate}[label=\textup{(\roman*)},leftmargin=*,itemsep=3pt,topsep=4pt]
\item \textbf{Approximation adjunction.} For continuous $D$, if $\Downarrow x\subseteq I$, then $x=\sup\Downarrow x\leq\sup I$.
Conversely $x\leq\sup I$ puts every $a\way x$ in $I$.
This proves $\Downarrow x\subseteq I$ if and only if $x\leq\sup I$, the first adjunction.
Its split identity follows from continuity; the second adjunction was proved for every dcpo.
\item \textbf{Converse construction.} Conversely, let $t\dashv\sup$.
Since $\sup yx=x$, one has $t(x)\subseteq yx$.
The adjunction gives $x\leq\sup t(x)$, and the containment gives the reverse inequality.
If $a\in t(x)$ and $x\leq\sup I$, then $t(x)\subseteq I$, so $a\in I$.
If $a\way x$, apply its defining test to $t(x)$, whose supremum is $x$.
Hence $t(x)=\Downarrow x$.
\item \textbf{Ideal and directedness calculation.} A left adjoint preserves the existing suprema in its source and target orders: for a directed family $x_i$ and an ideal $K$, the conditions $t(\bigvee_i x_i)\subseteq K$ and $t(x_i)\subseteq K$ for all $i$ are equivalent by adjunction.
The supremum in $\Idl D$ is union, proving the last assertion.

The comparison is witnessed by the retraction diagram
\[
\begin{tikzcd}[column sep=large,row sep=large]
D\ar[r,"\Downarrow_D"]\ar[dr,"\id_D"']&\Idl D\ar[d,"\sup_D"]\\ &D
\end{tikzcd}
\]
\end{enumerate}
\end{proof}

\begin{lemma}[Interpolation and basic opens]
\label{lem:ordinaryinterpolation}
    In a continuous dcpo, $a\way x$ implies $a\way b\way x$ for some $b$. Each predicate $a\way(-)$ is Scott open.
\end{lemma}
\begin{proof}
\leavevmode
\begin{enumerate}[label=\textup{(\roman*)},leftmargin=*,itemsep=3pt,topsep=4pt]
\item \textbf{Interpolation and basic opens.} The family of ideals $\Downarrow b$, for $b\way x$, is inhabited and directed because $\Downarrow x$ is an ideal and $\Downarrow$ is monotone.
Its union $K$ is an ideal.
Every member is below $x$; each $b\way x$ is the supremum of a subideal of $K$, so $\sup K=x$.
Applying $a\way x$ gives $a\in K$, hence $a\way b\way x$ for some $b$.
\item \textbf{Ideal and directedness calculation.} Upperness of $a\way(-)$ follows from its definition.
If $a\way\sup I$, interpolate to $a\way b\way\sup I$.
The ideal test gives $b\in I$; this is a witness in $I$ for the predicate $a\way(-)$.
Thus it is inaccessible by directed suprema.
\end{enumerate}
\end{proof}

\begin{proposition}[Free algebraic completion]
\label{prop:ordinarycompactrecognition}
    The compact elements of $\Idl P$ are exactly its principal ideals. A dcpo is algebraic if and only if its canonical map $\Idl K(D)\to D$, given by supremum of the included compact elements, is an isomorphism.
\end{proposition}
\begin{proof}
\leavevmode
\begin{enumerate}[label=\textup{(\roman*)},leftmargin=*,itemsep=3pt,topsep=4pt]
\item \textbf{Free algebraic completion.} If $\downarrow p\subseteq\bigcup_j I_j$, membership of $p$ gives one $I_j$ containing $\downarrow p$.
Thus principal ideals are compact.
Conversely $I=\bigcup_{p\in I}\downarrow p$ is a directed union; compactness of $I$ puts it below one of these principal ideals, and hence makes it equal to that principal ideal.
For a poset $P$, its generator is unique; for a preorder the compact part is its separated quotient.
\item \textbf{Unit and identity laws.} For algebraic $D$, send $x$ to the ideal of compacts below $x$.
Algebraicity gives one inverse identity with supremum.
If a compact $k$ lies below the supremum of an ideal of compacts, compactness puts $k$ below one of its members, so lower closure proves the other identity.
\item \textbf{Converse construction.} Conversely an isomorphism with an ideal completion supplies the compact generating family just described.
\end{enumerate}
\end{proof}

\begin{definition}[Kleisli arrows and the Karoubi envelope]
\label{def:karoubi}
    For a monad $M$, a Kleisli arrow $P\to Q$ is a map $P\to MQ$; composition is $n\star m=\mu\,Mn\,m$, and identity is the unit. An object of $\Kar\Kl(M)$ is $(P,r)$ with $r\star r=r$. An arrow $(P,r)\to(Q,s)$ is a Kleisli arrow $m$ satisfying $s\star m=m=m\star r$. Its identity is $r$. A splitting of an idempotent $e:X\to X$ consists of $j:D\to X$ and $q:X\to D$ with $jq=e$ and $qj=\id_D$.
\end{definition}

In the distributor calculus \cite{stubbe}, a Kleisli arrow for ordinary ideals has matrix $m(p,q)\equiv q\in m(p)$, upper in $p$ and lower in $q$. Its extension is $m^\#(I)(q)=\exists p[I(p)\wedge m(p,q)]$, and Kleisli composition is relational convolution. Restriction of a Scott-continuous map $\Idl P\to\Idl Q$ to principal ideals is its inverse description by \cref{prop:freeextension}. Thus the Kleisli category identifies fully faithfully with the category of free completions.

\begin{theorem}[Ordinary continuous completion]
\label{thm:ordinarykaroubi}
    The idempotent completion $\Kar\Kl(\T)$ is equivalent, as a locally ordered category, to the category of continuous dcpos and Scott-continuous maps. Every continuous dcpo is a retract of a free completion by a deflationary idempotent. Arbitrary retracts of free completions are also continuous.
\end{theorem}
\begin{proof}
\leavevmode
\begin{enumerate}[label=\textup{(\roman*)},leftmargin=*,itemsep=3pt,topsep=4pt]
\item \textbf{Ordinary continuous completion.} For continuous $D$, the maps $\Downarrow:D\to\Idl D$ and $\sup:\Idl D\to D$ are Scott-continuous maps, their composite on $D$ is the identity, and $\Downarrow\sup\leq\id$.
This gives the required free retraction.
\item \textbf{Converse construction.} Conversely let $s:D\to\Idl P$ and $q:\Idl P\to D$ be Scott-continuous maps with $qs=\id$.
For $p\in s(x)$, $q(\downarrow p)\way x$: if $x\leq\sup I$, apply $s$ and Scott continuity to the lower image of $I$; membership of $p$ in the resulting union gives $z\in I$ with $\downarrow p\subseteq s(z)$.
Applying $q$ gives $q(\downarrow p)\leq z$, so it belongs to $I$.
These basic approximants are inhabited and directed, and their supremum is $qs(x)=x$.
Their lower closure $K_x$ is therefore an ideal contained in $\Downarrow x$ with supremum $x$.
Applying the way-below test to $K_x$ gives the reverse containment.
Hence $D$ is continuous.
\item \textbf{Ideal and directedness calculation.} An idempotent Scott-continuous map splits by its fixed subobject; directed unions of fixed points are fixed.
Extension identifies Kleisli arrows with arrows between free objects, and the compatibility equations in \cref{def:karoubi} identify maps between their splittings.
This proves full faithfulness and the stated essential image.

The comparison is witnessed by the retraction diagram
\[
\begin{tikzcd}[column sep=large,row sep=large]
D\ar[r,"s"]\ar[dr,"\id_D"']&\Idl P\ar[d,"q"]\\ &D
\end{tikzcd}
\]
\end{enumerate}
\end{proof}


Continuity can now be recognized either through way-below ideals or through an adjoint splitting of supremum; algebraicity is recognized by compact generators. These categorical formulations are what allow the two effective completion doctrines to be compared without changing the ordinary meaning of approximation.

\part{Synthetic domain theory}\label{part:completion}
\label{part:searchable}

Ideal completion supplies all internal directed suprema, but these suprema need not preserve effective observations. This part separates two responses to that problem: select admitted weights using $\Sigma$, or retain all ideals and equip their completion with specified Scott observations. The effective-point comparison relates the two approaches, while searchable bases make the comparison executable.


We construct the admitted and observation-equipped doctrines separately, then compare them through the effective-point adjunction. The comparison is established first on free objects and subsequently on continuous retracts; searchable presentations make its operations explicit.

\section{A dictionary for the three completion categories}
\label{sec:categorydictionary}


    The following constructions use three different completeness operations. This dictionary fixes their object and morphism data; the definitions and comparisons follow in this part.
\begin{center}
\begin{tabular}{p{.19\linewidth}p{.34\linewidth}p{.36\linewidth}}
\toprule
    Category & Object data & Morphisms\\
    \midrule
    $\Alg(\T)$ & all-ideal dcpo & preservation of all directed suprema\\
    $\ED(\E)=\Alg(\Th)$ & all-ideal dcpo with specified observations & Scott-continuous maps preserving those observations\\
    $\AS=\Alg(\J)$ & poset with suprema of admitted ideals & preservation of admitted suprema\\
\bottomrule
\end{tabular}
\end{center}

The admitted doctrine is defined in \cref{jd:def:J}; the observation-equipped doctrine in \cref{def:obs,def:effective-dcpo}. The semantic retract comparison is established in \cref{jd:thm:splitcomparison}. A semantic dcpo retains all directed suprema, whereas its effective-point object retains the points with a jointly positive observation profile.


\section{A canonical restriction of ideal weights}
\label{sec:restriction-of-ideal-weights}

An unrestricted ideal pairs with every upper predicate to yield an arbitrary truth value. For effective computation we instead require its pairing with every positive upper predicate to remain in $\Sigma$. Expressing that requirement as a pullback gives the admitted-weight doctrine without choosing a basis or a family of observations.


Give $\Sigma$ the order induced by $i$. For an internal poset $P$, put
\[
     U_\Sigma P=\{u\in\Sigma^P\mid\forall p\leq q,\ u(p)\leq u(q)\},
\]
where the right side denotes the internal object of monotone maps. These are the positive upper predicates of $P$, determined by $P$ and $\Sigma$ with no selected observation family.

Every ideal $I\in\Idl P$ pairs with such predicates:
\[
    \chi_I(u)=\exists p:P\,[I(p)\wedge u(p)]\quad\in\Omega.
\]
Equivalently, it is the flat-weight pairing with upper predicates, restricted on the upper-predicate argument.

\begin{definition}
\label{jd:def:J}
    Define $\J P$ by the pullback
    \[
        \begin{tikzcd}[column sep=large,row sep=large]
            \J P\ar[r]\ar[d,hook,"j_P"']&\Sigma^{U_\Sigma P}\ar[d,hook,"i^{U_\Sigma P}"]\\
            \Idl P\ar[r,"\chi_P"']&\Omega^{U_\Sigma P}.
        \end{tikzcd}
    \]
    Call these the $\Sigma$-admissible ideals. They inherit inclusion order.
\end{definition}

\begin{theorem}[Canonical ideal doctrine]
\label{jd:thm:J}
    Ordinary ideal completion restricts to a locally monotone KZ monad $(\J,\eta,\mu)$ on $\Pos(\E)$. Its structural maps are
    \begin{align*}
        \eta_P(p)&=\downarrow p,\\
        (\J f)(I)(q)&\Longleftrightarrow\exists p\,[I(p)\wedge q\leq f(p)],\\
        \mu_P(\mathcal I)(p)&\Longleftrightarrow
        \exists I\in\J P\,[\mathcal I(I)\wedge I(p)].
    \end{align*}
    The inclusions $j_P:\J P\incl\Idl P$ form a monad morphism to ordinary ideal completion.
\end{theorem}
\begin{proof}
\leavevmode
\begin{enumerate}[label=\textup{(\roman*)},leftmargin=*,itemsep=3pt,topsep=4pt]
\item \textbf{Principal ideals and functoriality.} For a principal ideal, $\chi_{\downarrow p}(u)=u(p)$, so the unit factors through the pullback.
For monotone $f:P\to Q$ and $v\in U_\Sigma Q$,
    \[
        \chi_{\J f(I)}(v)=\chi_I(vf).
    \]
    The right side is in $\Sigma$, since $vf\in U_\Sigma P$.
This proves closure under lower image.
\item \textbf{Closure under ideal multiplication.} For $u\in U_\Sigma P$, the map
    \[
        \widehat u:\J P\to\Sigma,\qquad \widehat u(I)=\chi_I(u)
    \]
    is defined by the pullback and is monotone.
If $\mathcal I\in\J\J P$, its union is an ordinary ideal, and
    \begin{equation}
    \label{jd:eq:flatten}
        \chi_{\mu_P(\mathcal I)}(u)
        =\chi_{\mathcal I}(\widehat u)\quad\in\Sigma.
    \end{equation}
    This proves closure under flattening without asserting that $\J P$ is overt.
It is the main closure check.
\item \textbf{Monad and KZ laws.} The ordinary lower-image and union identities now prove functoriality, the unit equations, and associativity.
All restrictions are along monomorphisms, so the identities are reflected.
For the KZ inequality, if $K$ belongs to $(\J\eta_P)(I)$, there is $p\in I$ with $K\subseteq\downarrow p$, hence $K\subseteq I$.
Therefore
    \[
        \J\eta_P\leq\eta_{\J P}.
    \]
    These arguments are internal, so they also prove local monotonicity and parameter uniformity.
The same equations express compatibility of $j$ with the two monads.

Naturality is expressed by the commutative diagram
\[
\begin{tikzcd}[column sep=large,row sep=large]
\J\J P\ar[r,"{\mu_P^{\J}}"]\ar[d,"{\T(j_P)\,j_{\J P}}"']&\J P\ar[d,"j_P"]\\ \T\T P\ar[r,"{\mu_P^{\T}}"']&\T P
\end{tikzcd}
\]
\end{enumerate}
\end{proof}

\begin{lemma}[Preorder presentations of the same doctrine]
\label{lem:Jpreorders}
    The same pullback and formulas define $\J$ on internal preorders. Its algebra category is equivalent to $\AS=\Alg(\J|_{\Pos(\E)})$. If $P_{\mathrm{sep}}$ is the separated quotient of $P$, then $\J P\cong\J P_{\mathrm{sep}}$ canonically.
\end{lemma}
\begin{proof}
\leavevmode
\begin{enumerate}[label=\textup{(\roman*)},leftmargin=*,itemsep=3pt,topsep=4pt]
\item \textbf{Passage to separated quotients.} Upper predicates and lower ideals are invariant under $p\sim q\iff p\leq q\leq p$.
Their objects, pairings, and positive upper predicates therefore identify under the quotient.
The defining pullback gives the asserted isomorphism.
The closure calculations in \cref{jd:thm:J} use only the preorder laws, so they define the extended monad.
If $a:\J P\to P$ is an algebra, then $p\sim q$ gives $\downarrow p=\downarrow q$, and its unit equation gives $p=a(\downarrow p)=a(\downarrow q)=q$.
Thus every algebra is separated.
The inclusion of posets is full on the algebra maps, which proves the equivalence.
\end{enumerate}
\end{proof}

We use this extension when a basis has redundant names. In particular, a modest preorder may be used directly as a searchable basis.

\begin{proposition}[Universal property of the restriction]
\label{jd:prop:maximal}
    Among internal, parameter-stable submonads of ordinary ideal completion, $\J$ is the largest for which $\Sigma$ admits the prescribed ideal suprema computed through $i:\Sigma\incl\Omega$. Equivalently, $\J$ is terminal among such submonads equipped with an algebra on $\Sigma$ for which $i$ preserves the admitted suprema. Submonads here are considered with their internal parameter dependence, as is $\J$ itself.
\end{proposition}
\begin{proof}
\leavevmode
\begin{enumerate}[label=\textup{(\roman*)},leftmargin=*,itemsep=3pt,topsep=4pt]
\item \textbf{Universal property of the restriction.} For $I\in\J\Sigma$, the supremum computed in $\Omega$ is $\chi_I(\id_\Sigma)$ and belongs to $\Sigma$ by definition.
It is the least upper bound in $\Sigma$, and union gives its algebra laws.
\item \textbf{Converse construction.} Conversely, let $M\incl\Idl$ be a submonad with the stated algebra on $\Sigma$.
For $I\in MP$ and $u:P\to\Sigma$ monotone, $Mu(I)$ is an admitted ideal of $\Sigma$.
Its supremum, computed in $\Omega$, is $\chi_I(u)$, so this value belongs to $\Sigma$.
Thus $I\in\J P$, uniformly in the parameters.
The resulting factorization $M\to\J$ is unique because $\J\incl\Idl$ is monic; its unit and multiplication equations are reflected by that inclusion.

The universal property is represented by the commutative diagram
\[
\begin{tikzcd}[column sep=large,row sep=large]
MP\ar[r]\ar[dr,hook]&\J P\ar[d,hook,"j_P"]\\ &\Idl P
\end{tikzcd}
\]
\end{enumerate}
\end{proof}


Admissibility is thus characterized by a universal factorization of positive pairings, not by a semidecision of ideal inclusion. The algebra category inherits precisely these selected suprema, which can next be studied by the same adjoint methods used for ordinary ideal completion.

\section{Algebras, algebraicity, and approximation as adjunctions}
\label{sec:algebraicity-and-approx-as-adjunctions}

Once the admitted ideals form a monad, their algebras describe precisely which of those directed suprema exist. We next characterize the free and continuous cases intrinsically, through compact generators and adjoints to the supremum map. These criteria make later comparisons independent of a selected presentation.


\begin{definition}[Computational algebras and their maps]
\label{def:computationalalgebra}
    A $\J$-algebra is an internal poset $A$ with a monotone map $a:\J A\to A$ satisfying
    \[
        a\eta_A=\id_A,\qquad a\J a=a\mu_A.
    \]
    A morphism $(A,a)\to(B,b)$ is a monotone map $f:A\to B$ satisfying $f a=b\J f$. Write $\AS=\Alg(\J)$, with the pointwise order on its morphisms. These are the \textbf{computational domains}. An algebra is \textbf{algebraic} if it is isomorphic, as an algebra, to $(\J P,\mu_P)$ for some internal poset $P$. It is \textbf{continuous}, or \textbf{split}, if its structure map $a$ has an algebra-homomorphic section $t:A\to\J A$. The section is taken in $\Alg(\J)$: its equations are
    \[
        at=\id_A,\qquad ta=\mu_A\J t.
    \]
    In particular, a section of the underlying arrow of $\E$ alone is not the datum in this definition. The comparison with ordinary continuity and the semantic observation doctrine is proved below.
\end{definition}

For an algebra $a:\J A\to A$, the KZ property identifies $a$ with the supremum operation for the admitted ideals and gives $a\dashv\eta_A$. An algebra morphism is a monotone map preserving those weighted suprema. Effectivity is thus built into the chosen weights and the algebra-homomorphism equation.

\begin{proposition}[Categorical continuity]
\label{jd:prop:continuous}
    A $\J$-algebra $(A,a)$ is continuous precisely when $a$ admits a left adjoint:
    \[
        t_A\dashv a\dashv\eta_A,
        \qquad
        A\mathrel{\mathop{\rightleftarrows}^{t_A}_{a}}\J A.
    \]
    Then $at_A=\id_A$, $t_Aa\leq\id_{\J A}$, and $t_A\leq\eta_A$. Moreover $t_A$ is an algebra morphism and
    \[
        \mu_A\,\J t_A\,t_A=t_A.
    \]
    There is an equivalence
    \[
        \Spl(\AS)\simeq\Kar\Kl(\J).
    \]
\end{proposition}
\begin{proof}
\leavevmode
\begin{enumerate}[label=\textup{(\roman*)},leftmargin=*,itemsep=3pt,topsep=4pt]
\item \textbf{KZ-algebra characterization.} The split-algebra statements are the general KZ results, including the equivalence between a left adjoint to $a$ and an algebra-homomorphic section of $a$; see \cite[Proposition 3.8 and Theorem 3.10]{hofmannsousa}.
The required idempotent completeness of $\Pos(\E)$ is automatic: an idempotent $e:P\to P$ splits through its equalizer with $\id_P$, with the induced order.
Here the split-algebra assertions can also be checked directly.
\item \textbf{Retractions of free algebras.} The adjunctions give $t_A\leq\eta_A$ and $at_A=\id$. A left adjoint preserves every existing supremum, so $t_A$ preserves the admitted ideal suprema and is an algebra morphism.
Its algebra equation gives the displayed Kleisli idempotency.
\item \textbf{Adjunction from an algebra section.} For completeness, a retract of a free algebra is split: if $s:A\to\J B$ and $q:\J B\to A$ are algebra morphisms with $qs=\id$, set
    \[
        t_A=\J q\;\J\eta_B\;s.
    \]
    The two required calculations, with all algebra structures retained, are
    \begin{align*}
        at_A&=a\J q\,\J\eta_B\,s
        =q\mu_B\,\J\eta_B\,s=qs=\id_A,\\
        t_Aa&=\J q\,\J\eta_B\,sa
        =\J q\,\J\eta_B\,\mu_B\,\J s
        \leq\J q\,\J s=\id_{\J A}.
    \end{align*}
    Here $a\J q=q\mu_B$ and $sa=\mu_B\J s$ are the two homomorphism equations, and $\J\eta_B\,\mu_B\leq\id$ is the counit inequality of the free adjunction.
\item \textbf{Karoubi equivalence.} The equations prove $t_A\dashv a$.
Conversely, a left adjoint $t_A$ gives an algebra-homomorphic section as above and hence a retract of the free algebra on $A$.
\item \textbf{Recovering the adjoint from an algebra section.} Conversely, start with an algebra-homomorphic section of $a$.
Since $at_A=\id$, adjunction $a\dashv\eta_A$ gives $t_A\leq\eta_A$.
For $I\in\J A$, its homomorphism equation yields
    \[
        t_A(aI)=\mu_A\J t_A(I)
        =\bigcup_{x\in I}t_A(x)\subseteq I,
    \]
    where $t_A(x)\subseteq\downarrow x$ gives the inclusion.
Thus $t_A\dashv a$.
\item \textbf{Idempotent splitting.} Finally, split an idempotent algebra endomorphism by its fixed subobject: the image is closed under admitted suprema because the endomorphism preserves them.
This constructs the splittings used in $\Kar\Kl(\J)$ and proves the stated equivalence, with the same arrow formula as in \cref{def:karoubi}.

The comparison is witnessed by the retraction diagram
\[
\begin{tikzcd}[column sep=large,row sep=large]
A\ar[r,"t_A"]\ar[dr,"\id_A"']&\J A\ar[d,"a"]\\ &A
\end{tikzcd}
\]
\end{enumerate}
\end{proof}

\subsection{Way-below by a universal property}
\label{subsec:way-below}

The adjoint triple determines the approximation relation. Its membership formula is
\[
    b\way_\Sigma x\quad\Longleftrightarrow\quad b\in t_A(x).
\]
Equivalently, use $\Omega$-distributors with convention $\Phi:X\nrightarrow Y$ has values $\Phi(y,x)$. For $f:X\to Y$, put
\[
    f_*(y,x)=Y(y,fx),\qquad f^*(x,y)=Y(fx,y).
\]
Define $W_A:A\nrightarrow A$ by the right extension
\begin{equation}
\label{jd:eq:wayuniv}
     W_A=\Ran_{a_*}(\eta_A^*),\qquad
    \Phi a_*\leq\eta_A^*\quad\Longleftrightarrow\quad\Phi\leq W_A.
\end{equation}
Existence uses the complete hom-lattices of $\Omega$-distributors, not arbitrary joins in $\Sigma$. Expanding this universal property gives
\[
    W_A(b,x)\Longleftrightarrow
    \forall I\in\J A\,[x\leq aI\Rightarrow b\in I].
\]
When $t_A\dashv a$, this is equivalent to $b\in t_A(x)$: use the adjunction in one direction and the admitted ideal $t_A(x)$ in the other. For ordinary ideal completion the same construction gives the ordinary way-below distributor. The general enriched construction is due to \cite{hofmannwaszkiewicz}. Here the right extension is calculated for the dominance-selected ideal doctrine.

\subsection{Compactness and algebraicity}
\label{subsec:compactness-and-algebraicity}

For continuous $A$, define its compact part by the equalizer
\[
    \begin{tikzcd}[column sep=large]
        K(A)\ar[r,hook,"k"]&A\ar[r,shift left=.6ex,"t_A"]
        \ar[r,shift right=.6ex,"\eta_A"']&\J A.
    \end{tikzcd}
\]
The comparison from the free algebra on this compact part is
\[
    c_A=a\,\J k:\J K(A)\longrightarrow A.
\]
\begin{proposition}
\label{jd:prop:compact}
    For an internal poset $B$, $K(\J B)\cong B$, via $b\mapsto\downarrow b$. A continuous $\J$-algebra is algebraic if and only if $c_A$ is an isomorphism.
\end{proposition}
\begin{proof}
\leavevmode
\begin{enumerate}[label=\textup{(\roman*)},leftmargin=*,itemsep=3pt,topsep=4pt]
\item \textbf{Ideal and directedness calculation.} On $\J B$, the approximation map is $\J\eta_B$.
Its value at $I$ consists of admitted ideals below a principal ideal generated by a member of $I$.
If $\J\eta_B(I)=\eta_{\J B}(I)$, membership of $I$ in the right side gives $b\in I$ with $I\subseteq\downarrow b$.
Lower closure gives equality.
\item \textbf{Converse construction.} Conversely principal ideals satisfy the equalizer equation.
Antisymmetry makes the generating $b$ unique.
These constructions are internal inverse maps.
If $A$ is free, this identification makes $c_A$ an isomorphism.
The converse is immediate from its displayed source.
\end{enumerate}
\end{proof}

Thus algebraicity is recognized by a canonical comparison map, and continuity by an adjoint.

\section{Cartesian closure of the computational algebra category}
\label{sec:cartesian-closure}

Higher-order programs require function domains whose operations preserve admitted suprema. The pointwise order alone does not establish this: evaluation and currying must remain algebra morphisms. Product preservation of the admitted doctrine supplies the interchange of suprema needed to construct the exponential.


\begin{lemma}
\label{jd:lem:products}
    $\J$ preserves finite products. Under this identification its binary product operation sends $(I,J)$ to $I\times J$.
\end{lemma}
\begin{proof}
\leavevmode
\begin{enumerate}[label=\textup{(\roman*)},leftmargin=*,itemsep=3pt,topsep=4pt]
\item \textbf{Ideal and directedness calculation.} For $u:P\times Q\to\Sigma$ monotone, the function
    \[
        p\longmapsto\chi_J(u(p,-))
    \]
    is a positive upper predicate on $P$.
Thus
    \[
        \chi_{I\times J}(u)=\chi_I\bigl(p\mapsto\chi_J(u(p,-))\bigr)\in\Sigma.
    \]
    Conversely an ideal in $P\times Q$ is the product of its lower-image projections: directedness combines the witnesses for its two coordinates.
The projections of an admitted ideal are admitted by functoriality.
\item \textbf{The terminal product.} Finally the unique ideal of $1$ is admitted.
The formulas are natural and compatible with unit and multiplication.
\end{enumerate}
\end{proof}

\begin{theorem}
\label{jd:thm:ccc}
    $\AS$ is cartesian closed. Its exponential $[A,B]_\Sigma$ is the internal poset of $\J$-algebra morphisms, with pointwise order and pointwise admitted suprema.
\end{theorem}
\begin{proof}
We first form pointwise admitted suprema on the candidate function object. Product preservation then identifies the two possible orders of taking a supremum, which proves evaluation and transposition are algebra morphisms.

\leavevmode
\begin{enumerate}[label=\textup{(\roman*)},leftmargin=*,itemsep=3pt,topsep=4pt]
\item \textbf{Pointwise joins.} Let $H=[A,B]_\Sigma$ denote the stated subobject of $B^A$.
For $I\in\J H$, define
    \[
        f_I(x)=b\bigl(\J(\operatorname{ev}_x)(I)\bigr),
    \]
    where $b:\J B\to B$ is the target algebra.
This is the pointwise supremum of $I$ and is monotone.
\item \textbf{Preservation of admitted suprema.} To prove it is an algebra morphism, let $J\in\J A$.
The two iterated suprema
    \[
        \bigvee_{h\in I}\bigvee_{x\in J}h(x),
        \qquad
        \bigvee_{x\in J}\bigvee_{h\in I}h(x)
    \]
    exist: lower images under all evaluation and parameter maps are admitted, and \cref{jd:lem:products} admits $I\times J$.
Both expressions are the supremum of the lower image of $I\times J$ under evaluation, by the algebra multiplication equation.
\item \textbf{Algebra structure on the function object.} Since each $h$ preserves the supremum of $J$, this proves
    \[
        f_I(aJ)=b\bigl(\J f_I(J)\bigr).
    \]
    Thus $H$ has the stated admitted suprema, with algebra laws following from their least-upper-bound property.
\item \textbf{Joint continuity of evaluation.} For an admitted ideal in $H\times A$, \cref{jd:lem:products} identifies it with the product of its two projection ideals.
The preceding calculation proves that evaluation preserves its supremum.
\item \textbf{Currying along a principal ideal.} We construct the transpose uniformly in the first variable.
Let $h:C\times A\to B$ be an algebra morphism, let $c:C$ be a generalized element in an arbitrary context, and write $h_c(x)=h(c,x)$.
For $I\in\J A$, the product ideal $\downarrow c\times I$ is admitted and has supremum $(c,a_A I)$.
Hence
    \[
        h(c,a_A I)=b\bigl(\J h(\downarrow c\times I)\bigr).
    \]
    The lower image $\J h(\downarrow c\times I)$ is exactly $\J h_c(I)$.
\item \textbf{Identification of the lower image.} Indeed, a generator $h(c',x)$ with $c'\leq c$ and $x\in I$ lies below $h(c,x)$ by monotonicity, while the reverse inclusion uses the generator with $c'=c$.
Therefore
    \[
        h_c(a_A I)=b\bigl(\J h_c(I)\bigr),
    \]
    so every slice $h_c:A\to B$ is a $\J$-algebra morphism, uniformly in $c$.
Thus the ordinary transpose $\lambda h:C\to H$ is defined.
\item \textbf{Continuity of the transpose.} Now let $K\in\J C$ and fix $x:A$.
The admitted ideal $K\times\downarrow x$ has supremum $(a_CK,x)$.
Preservation by $h$ gives
    \[
        h(a_CK,x)=b\bigl(\J h(K\times\downarrow x)\bigr).
    \]
    As above, monotonicity in the second variable gives
    \[
        \J h(K\times\downarrow x)
        =\J(\operatorname{ev}_x\,\lambda h)(K).
    \]
    By the pointwise definition of the algebra structure on $H$, the right side is
    \[
        \operatorname{ev}_x\bigl(a_H(\J(\lambda h)(K))\bigr).
    \]
    Since this equality holds uniformly in $x$, extensionality in the exponential yields
    \[
        (\lambda h)(a_CK)=a_H\bigl(\J(\lambda h)(K)\bigr).
    \]
    Thus transposition is an algebra morphism.
The ordinary transpose and evaluation equations restrict, proving cartesian closure.

The universal property is represented by the commutative diagram
\[
\begin{tikzcd}[column sep=large,row sep=large]
C\times A\ar[r,"{\lambda h\times\id_A}"]\ar[dr,"h"']&{[A,B]_\Sigma\times A}\ar[d,"\ev"]\\ &B
\end{tikzcd}
\]
\end{enumerate}
\end{proof}

This supplies higher-order computation at the level of the whole algebra category. When a particular function domain is claimed to be continuous, its approximation structure still has to be exhibited.

\section{Observation structures and the lifted monad}
\label{sec:observations}

Restricting ideal weights is one way to model effectivity; another is to retain the full domain while specifying the observations that can be tested positively. We construct the second account as a lifting of ordinary ideal completion. Its algebra maps will preserve selected Scott opens, making semantic domains available before their effective points are extracted.


We now make the effective structure part of the base category on which ideal completion acts. This definition applies without a basis or a continuity condition.

\begin{definition}
\label{def:obs}
    An \textbf{observation structure} on $P$ is a subobject $\OO_P\incl\Up(P)$ with the following internal closure properties:
    \begin{enumerate}
        \item every constant predicate with value in $\Sigma$ belongs to $\OO_P$;

        \item $\OO_P$ is closed under finite meets;

        \item for every overt subobject $S\incl\OO_P$, its union $\bigvee S=\lambda p.\exists U:S\,U(p)$ belongs to $\OO_P$.
    \end{enumerate}
    A map $(P,\OO_P)\to(Q,\OO_Q)$ is a monotone map $f$ whose inverse-image map restricts to $f^*:\OO_Q\to\OO_P$. Write $\ObsPre(\E)$ for this locally preordered category.
\end{definition}

Closure under overt subfamilies is equivalent to closure under overt-indexed families: the image of an overt indexing object is overt by \cref{lem:overtclosure}. The subfamily formulation makes the definition internal, with quantification over $\Pow(\OO_P)$. In particular it defines a subobject of the power object of $\Up(P)$. All three clauses hold in context, so the restriction maps are uniform.

\paragraph{Internal construction of the observation structure.}
The quantification over overt subfamilies in \cref{def:obs} can be expressed entirely in higher-order internal logic. For a variable subobject $S\incl\Up(P)$, the formula
\[
 \overt(S)\;:\!\Longleftrightarrow\;
 \forall v\in\Sigma^S\;\bigl(\exists U:S\;i(v(U))\bigr)\in\Sigma
\]
asserts that the existential classifier factors through $i$. The unique factorization exists by unique choice in a topos, not by an axiom selecting stream representatives. For $O\subseteq\Up(P)$, closure under overt subfamilies is the bounded formula
\[
 \forall S\in\Pow(\Up(P))\;
 \left(S\subseteq O\wedge\overt(S)\Rightarrow\bigvee S\in O\right).
\]
Together with the constant and finite-meet clauses, this defines the object of observation structures used in \cref{lem:obshull}. 

If a family $u:I\to O$ is indexed by an overt object, its image $S$ is overt by the regular-image argument of \cref{lem:overtclosure}, and $\bigvee_Iu=\bigvee S$. Conversely every overt subfamily is an overt-indexed family by its inclusion map. This proves the equivalence of the two formulations with parameters. It is also the reason that passing to the image of the injective dagger map preserves the observation closure in \cref{thm:lifted}.


\begin{lemma}
\label{lem:obshull}
    Every family $A\subseteq\Up(P)$ is contained in a least observation structure $\cl(A)$. Inverse image commutes with the tests for membership in an observation structure: a monotone map preserving a generating family preserves its observation closure.
\end{lemma}
\begin{proof}
\leavevmode
\begin{enumerate}[label=\textup{(\roman*)},leftmargin=*,itemsep=3pt,topsep=4pt]
\item \textbf{Inverse comparison.} Intersect, internally, all subobjects $O\subseteq\Up(P)$ satisfying \cref{def:obs} and containing $A$.
This is a bounded power-object construction; the whole $\Up(P)$ is one such subobject.
The intersection contains the constants and finite meets.
An overt subfamily of the intersection maps isomorphically to a subfamily of each $O$, so its union lies in every $O$.
\item \textbf{Closure under inverse image.} For a monotone $f:P\to Q$, the collection of $V\in\Up(Q)$ with $f^*V\in\OO_P$ is an observation structure: inverse image preserves constants and finite meets, and the image under $f^*$ of an overt subfamily is overt.
This proves the last assertion by minimality.
\end{enumerate}
\end{proof}

\begin{theorem}[Lifted ideal completion]
\label{thm:lifted}
    The ideal-completion monad lifts to a KZ monad $\Th$ on $\ObsPre(\E)$ by
    \begin{equation}\label{eq:lifted}
        \Th(P,\OO_P)=\bigl(\Idl(P),\OO_P^\dagger\bigr),
        \qquad \OO_P^\dagger=\{U^\dagger\mid U\in\OO_P\}.
    \end{equation}
    The forgetful functor to $\Pre(\E)$ strictly preserves the monad.
\end{theorem}
\begin{proof}
\leavevmode
\begin{enumerate}[label=\textup{(\roman*)},leftmargin=*,itemsep=3pt,topsep=4pt]
\item \textbf{Lifted ideal completion.} By \cref{lem:tracefree}, dagger identifies $\OO_P$ with its image, preserves finite meets and joins, and sends constant predicates to the same constants.
In particular it carries overt subfamilies isomorphically to overt subfamilies.
Thus \eqref{eq:lifted} is an observation structure.
\item \textbf{Unit and identity laws.} For a monotone map $f$ and an upper predicate $V$ one has
    \begin{equation}\label{eq:liftchecks}
        (\T f)^*(V^\dagger)=(f^*V)^\dagger,
        \quad y_P^*(U^\dagger)=U,
        \quad \mu_P^*(U^\dagger)=(U^\dagger)^\dagger.
    \end{equation}
    The first identity uses upperness to eliminate a lower-image witness.
\item \textbf{Trace and multiplication identities.} The second is the trace identity.
The third expands to $\exists I,p\,[\mathcal I(I)\wedge I(p)\wedge U(p)]$.
These equations show that the functor, unit, and multiplication preserve the specified observations.
All monad equations and the KZ inequality are inherited through the faithful forgetful functor.
\end{enumerate}
\end{proof}

\begin{definition}
\label{def:effective-dcpo}
    An \textbf{effective dcpo} is a $\Th$-algebra. Its category is
    \[
        \ED(\E)=\Alg(\Th).
    \]
    We also write $\C=\ED(\E)$ when discussing its hom-objects or presentations. An \textbf{effective algebraic domain} and an \textbf{effective continuous domain} are objects of this category whose underlying dcpos are algebraic and continuous, respectively. The recognition theorems below identify these classes with free and split algebras for the lifted doctrine.
\end{definition}

\begin{theorem}[Algebra recognition]
\label{thm:recognition}
    An effective dcpo is equivalently an internal dcpo $D$ with an observation structure $\OO_D\subseteq\Sc(D)$. Its morphisms are Scott-continuous maps preserving these observations by inverse image.

    In particular the algebra laws are the commuting unit triangle and multiplication square in $\ObsPre(\E)$:
    \[
        \begin{tikzcd}[column sep=large]
            D\ar[r,"{y_D}"]\ar[dr,equal]&\Th D\ar[d,"{a}"]\\
            &D
        \end{tikzcd}
        \qquad
        \begin{tikzcd}[column sep=large]
            \Th\Th D\ar[r,"{\mu_D}"]\ar[d,"{\Th a}"']&\Th D\ar[d,"{a}"]\\
            \Th D\ar[r,"{a}"']&D.
        \end{tikzcd}
    \]
\end{theorem}
\begin{proof}
\leavevmode
\begin{enumerate}[label=\textup{(\roman*)},leftmargin=*,itemsep=3pt,topsep=4pt]
\item \textbf{Unit and identity laws.} The underlying algebra is a dcpo by \cref{thm:ordinaryalgebras}.
Effectiveness of $\sup_D$ says that for $U\in\OO_D$ there is $V\in\OO_D$ with
    \[
        U(\sup I)=V^\dagger(I).
    \]
    Restriction along $y_D$ gives $U=V$, by the algebra unit law and \eqref{eq:liftchecks}.
Hence
    \begin{equation}\label{eq:scottobs}
        U(\sup I)\iff\exists x\,[I(x)\wedge U(x)],
    \end{equation}
    which is Scott openness.
\item \textbf{Converse construction.} Conversely \eqref{eq:scottobs} makes the ordinary supremum map observation-preserving, hence a $\Th$-algebra.
The morphism statement follows from the underlying algebra equation and the definition of $\ObsPre(\E)$.

The defining arrows form the commutative diagram
\[
\begin{tikzcd}[column sep=large,row sep=large]
\Idl D\ar[r,"\sup_D"]\ar[dr,"U^\dagger"']&D\ar[d,"U"]\\ &\Omega
\end{tikzcd}
\]
\end{enumerate}
\end{proof}

\subsection{What factors through the dominance}
\label{subsec:what-factors-through-the-dominance}

The lifted construction resolves the variance issue in trying to make effectivity a change of enrichment. There are three evaluation maps with different types:
\[
    P\times P\longrightarrow\Omega,\qquad
    P\times\OO_P\longrightarrow\Omega,\qquad
    \Pt(P)\times\OO_P\longrightarrow\Sigma,
\]
where the final notation is defined by the same pullback as for dcpos below. The first map is the order; the second is evaluation of observations; the third is its restriction at effective points. It is the third map that factors through the dominance. On a searchable basis, the approximation matrix also takes values in $\Sigma$, and overtness supplies the existential compositions needed to pass between point and observation traces.

The $\Omega$-enriched structure still expresses all ideal joins, including ideals whose membership has no semidecision. Restricting the allowed ideals at this stage would change the directed-completeness notion. Instead, the equations in \eqref{eq:liftchecks} show that free completion transports the selected observation family while retaining the original ideals. This is why the algebra category has a general effective-dcpo notion before any continuity assumption is made.

\begin{proposition}[Effective free extension]
\label{prop:effectivefreeextension}
    For $(P,\OO_P)\in\ObsPre(\E)$ and an effective dcpo $(D,\OO_D)$, the free extension bijection restricts to
    \[
        \ED(\E)(\Th(P,\OO_P),D)
        \cong\ObsPre(\E)((P,\OO_P),(D,\OO_D)).
    \]
\end{proposition}
\begin{proof}
\leavevmode
\begin{enumerate}[label=\textup{(\roman*)},leftmargin=*,itemsep=3pt,topsep=4pt]
\item \textbf{Effective free extension.} If $f$ preserves observations, Scott openness of $U\in\OO_D$ gives
    \begin{align*}
        U(f^\#I)
        &\iff \exists p\,[I(p)\wedge U(f(p))]\\
        &\iff (f^*U)^\dagger(I).
    \end{align*}
    Thus the extension preserves the specified observations.
\item \textbf{Converse construction.} Conversely an effective Scott-continuous map restricts along the effective unit to an observation-preserving map.
The ordinary extension and restriction are inverse by \cref{prop:freeextension}, so these restricted assignments are inverse as well.

The defining arrows form the commutative diagram
\[
\begin{tikzcd}[column sep=large,row sep=large]
P\ar[r,"y_P"]\ar[dr,"f"']&\Idl P\ar[d,"f^\#"]\\ &D
\end{tikzcd}
\]
\end{enumerate}
\end{proof}

The proof supplies a computationally useful formula: an observation of a completed value is witnessed by one finite-stage input together with the pulled-back observation at that input. For general algebras this is the Scott-observation equation. For searchable rounded presentations it becomes a restricted existential quantifier with an explicit approximation witness.


The lifted doctrine has not restricted the ideals of its underlying dcpo. Instead it has made the chosen Scott observations part of the algebraic data. Effective points can now be defined by asking where those observations jointly have $\Sigma$-valued evaluations.

\section{Effective points and finite limits}
\label{sec:points}

An observation-equipped domain can contain points at which some selected observation is not semidecidable. The effective-point object retains exactly those points whose entire observation profile factors through $\Sigma$. We verify this factorization in families and show that it is compatible with finite-limit constructions.


Let $\epsilon_D:D\to\Omega^{\OO_D}$ be the transpose of evaluation.

\begin{definition}
\label{def:points}
    The object of effective points is the pullback
    \[
        \begin{tikzcd}
            \Pt(D)\ar[r]\ar[d,hook,"{j_D}"'] & \Sigma^{\OO_D}\ar[d,hook]\\
            D\ar[r,"{\epsilon_D}"'] & \Omega^{\OO_D}.
        \end{tikzcd}
    \]
    Its induced evaluation is denoted $\operatorname{Obs}_D:\Pt(D)\times\OO_D\to\Sigma$.
\end{definition}

\begin{proposition}
\label{prop:pointsfunctor}
    Effective points define a functor $\Pt:\ED(\E)\to\E$. For an effective map $f:D\to E$,
    \[
        \operatorname{Obs}_E(\Pt(f)(x),V)
        =\operatorname{Obs}_D(x,f^*V).
    \]
\end{proposition}
\begin{proof}
\leavevmode
\begin{enumerate}[label=\textup{(\roman*)},leftmargin=*,itemsep=3pt,topsep=4pt]
\item \textbf{Effective-point factorization.} The right-hand side is a $\Sigma$-valued evaluation profile for $f(j_Dx)$.
The pullback gives its unique factorization through $\Pt(E)$.
Identity and composition follow by uniqueness, and the equation is the evaluation form of that factorization.
\item \textbf{Comparison with the inverse image.} Naturality induces a canonical map
\[
 \Pt D\longrightarrow D\times_E\Pt E,
 \qquad x\longmapsto(j_Dx,\Pt(f)x).
\]
It is monic because its composite with the first projection is $j_D$.
Its construction proves commutativity, not cartesianness: effectiveness of $f(x)$ alone need not make $x$ effective.
The precise pullback statement for generated observations is proved in
\cref{prop:initial-observation-points}.

The construction uses the commutative square
\[
\begin{tikzcd}[column sep=large,row sep=large]
\Pt D\ar[r,"\Pt f"]\ar[d,hook,"j_D"']&\Pt E\ar[d,hook,"j_E"]\\ D\ar[r,"f"']&E
\end{tikzcd}
\]
\end{enumerate}
\end{proof}

There is a terminal effective dcpo $1$ with $\OO_1=\Sigma$. Its observation structure is closed under overt unions by definition of overtness. A point $1\to D$ is an effective morphism precisely when its evaluation profile factors through $\Sigma$, so global effective points agree with global elements of $\Pt(D)$.


\begin{lemma}[Testing generators in context]
\label{lem:testinggenerators}
    Suppose $\OO_D=\cl(G)$ and let $x:S\to D$ be any parameterized family. If evaluation restricts to a map $S\times G\to\Sigma$, then $x$ factors uniquely through $\Pt(D)$. Conversely any such factorization restricts to this evaluation on generators.
\end{lemma}
\begin{proof}
\leavevmode
\begin{enumerate}[label=\textup{(\roman*)},leftmargin=*,itemsep=3pt,topsep=4pt]
\item \textbf{Evaluation and transposition.} Work in the internal context $s:S$.
The collection of upper predicates $U$ for which $U(x(s))$ lies in $\Sigma$ contains the $\Sigma$-constants.
It is closed under finite meets by finite closure of $\Sigma$.
If an overt subfamily consists of such predicates, evaluation of its union is the restricted existential quantifier over that family, and therefore again lies in $\Sigma$.
The collection is thus an observation structure in this context.
It contains $G$ by hypothesis, so it contains $\cl(G)$ by minimality.
\item \textbf{Stability under parameter substitution.} The closure construction is internal and is preserved by substitution of the parameter, as its bounded formulation in \cref{def:obs} makes explicit.
\item \textbf{Joint factorization through $\Sigma$.} Consequently $U(x(s))$ lies in $\Sigma$ jointly in $(s,U)$.
Monicity of $i$ and unique choice give a single factorization $S\times\OO_D\to\Sigma$.
Transposing it supplies the map into the pullback defining $\Pt(D)$.
Uniqueness follows from monicity of $j_D$.
The converse is restriction of the resulting evaluation map.
\end{enumerate}
\end{proof}

This lemma keeps uniformity visible. Knowing a separate external semidecision for each observation is not its hypothesis; one needs a morphism on the parameterized generating family. Once that morphism is given, the internal observation closure constructs the full profile without selecting separate representatives for its values.

\begin{proposition}[Effective points of an initial observation structure]
\label{prop:initial-observation-points}
Given a finite family of Scott-continuous maps $f_i:D\to E_i$ into effective dcpos, equip $D$ with
\[
 \OO_D=\cl\{f_i^*V\mid i,\ V\in\OO_{E_i}\}.
\]
This constructs an effective dcpo and makes the following square cartesian:
\[
\begin{tikzcd}[column sep=large,row sep=large]
 \Pt D\ar[r,"{(\Pt f_i)_i}"]\ar[d,hook,"j_D"']
   &\prod_i\Pt E_i\ar[d,hook,"{\prod_i j_{E_i}}"]\\
 D\ar[r,"{(f_i)_i}"']&\prod_i E_i.
\end{tikzcd}
\]
\end{proposition}
\begin{proof}
\leavevmode
\begin{enumerate}[label=\textup{(\roman*)},leftmargin=*,itemsep=3pt,topsep=4pt]
\item \textbf{Construction of the observations.} Inverse images of the selected observations are Scott open. Scott opens are closed under the observation operations, so their generated structure is contained in $\Sc(D)$. Each $f_i$ is effective by construction.
\item \textbf{Parameterized universal property.} A map $x:T\to D$ together with lifts of all $f_ix$ supplies a $\Sigma$-valued evaluation on the tagged family $(i,V)$. Equal generators have equal truth values; monicity of $\Sigma\incl\Omega$ makes their classifiers equal. Thus this evaluation descends along the regular epimorphism from the tagged family onto its image of generators. By \cref{lem:testinggenerators}, $x$ factors uniquely through $\Pt D$. Conversely a lift of $x$ supplies all component lifts by \cref{prop:pointsfunctor}.
\item \textbf{Coherence.} These constructions are inverse in every parameter context and commute with substitution, since they are the unique factorizations through the displayed monomorphisms. They therefore prove the pullback universal property, including the empty-family case.
\end{enumerate}
\end{proof}

\begin{proposition}
\label{prop:finitelimits}
    The category $\ED(\E)$ has finite limits, created on underlying dcpos and equipped with the observation structures generated by the inverse images along their limit projections. The functor $\Pt$ preserves finite limits.
\end{proposition}
\begin{proof}
\leavevmode
\begin{enumerate}[label=\textup{(\roman*)},leftmargin=*,itemsep=3pt,topsep=4pt]
\item \textbf{Product structure.} Products of dcpos have componentwise order and directed suprema.
Equalizers of Scott-continuous maps inherit directed suprema, since the two images of any such supremum agree.
Equip a product with the closure of the inverse images of the factor observations, and an equalizer with the closure of restricted observations.
Scott opens themselves form an observation structure, so these closures consist of Scott opens.
The universal properties follow from \cref{lem:obshull}.
\item \textbf{Testing the limit generators.} Apply \cref{prop:initial-observation-points} to the constructed limit projections. For a product this gives the product of the effective-point objects. For an equalizer $K\incl D$, it gives $\Pt K=K\times_D\Pt D$; naturality and monicity of $j_E$ identify its equation with equality of the two maps into $\Pt E$. The empty-family case gives $\Pt1\cong1$. These are precisely the finite-limit comparisons.

Naturality is expressed by the commutative diagram
\[
\begin{tikzcd}[column sep=large,row sep=large]
\Pt(D\times E)\ar[r,"\cong"]\ar[d,"\Pt(\pi_D)"']&\Pt D\times\Pt E\ar[d,"\pi_D"]\\ \Pt D\ar[r,"\id"']&\Pt D
\end{tikzcd}
\]
\end{enumerate}
\end{proof}

For example, an arrow $x:S\to D\times E$ factors through the effective points precisely when each component does. In one direction compose with the effective projections. In the other, the two component profiles supply $\Sigma$-valued evaluation on the generating projection observations, and \cref{lem:testinggenerators} supplies the entire product profile. The resulting isomorphism
\[
    \Pt(D\times E)\cong\Pt(D)\times\Pt(E)
\]
is natural in $S$, hence is an isomorphism of objects. For an equalizer $K\incl D$ of effective maps $f,g$, its generating observations are the restrictions of $\OO_D$. The same argument gives
\[
    \Pt(K)\cong\{x\in\Pt(D)\mid\Pt(f)(x)=\Pt(g)(x)\}.
\]
The equation makes sense on denotations because $\Pt(E)\incl E$ is monic. This describes the full finite-limit comparison needed when presenting simultaneous computations and equations between them.


Finite-limit compatibility lets equations between computations be imposed before or after taking effective points. What remains is to identify which observed dcpos can be reconstructed from compact generators, from splittings, or from more general directed-cover relations.

\section{Algebraicity, continuity, and their categorical forms}\label{sec:continuity}

The lifted observation doctrine has the same algebraic and continuous shapes as ordinary domain theory, but the comparison maps must preserve observations as well as order. We identify free algebras by their compact parts, continuous ones by approximation adjoints, and arbitrary ones by the limits of those descriptions.


For the ordinary ideal doctrine, approximation is the distributor obtained by the right-extension construction of \cref{jd:eq:wayuniv}, with $\Idl$ in place of $\J$. A semantic dcpo is continuous when its supremum map admits a left adjoint in the underlying locally ordered category. We denote that adjoint by $\Downarrow_D$. Its calculation is
\[
    a\way x\quad\Longleftrightarrow\quad
    \forall I\in\Idl D\,[x\leq\sup I\Rightarrow a\in I].
\]
Continuity is equivalently the assertion that these approximation ideals have supremum $x$: the two implications are the adjunction calculation given in \cref{thm:continuous}. The compact part is the equalizer of $\Downarrow_D$ and $y_D$; a semantic domain is algebraic when the canonical ideal-completion map on its compact part is an isomorphism. The following recognition theorem identifies this categorical definition with generation by compact elements.

\begin{theorem}[Algebraic recognition]
\label{thm:algebraic}
    An effective dcpo is algebraic if and only if it is isomorphic to a free $\Th$-algebra. More precisely, for an algebraic $D$, let $K(D)$ be its compact subobject and let $\OO_K$ consist of the restrictions of $\OO_D$. Then
    \[
        (D,\OO_D)\cong\Th(K(D),\OO_K).
    \]
\end{theorem}
\begin{proof}
\leavevmode
\begin{enumerate}[label=\textup{(\roman*)},leftmargin=*,itemsep=3pt,topsep=4pt]
\item \textbf{Algebraic recognition.} In $\Idl(P)$, principal ideals are compact: if $\downarrow p$ is contained in a directed union, membership of $p$ puts it inside a member.
Conversely every compact ideal is principal, since $I=\bigcup_{p\in I}\downarrow p$.
Hence every free algebra is algebraic.
\item \textbf{Unit and identity laws.} For algebraic $D$, the maps
    \[
        x\longmapsto\{k\in K(D)\mid k\leq x\},
        \qquad I\longmapsto\sup_D I
    \]
    are inverse.
Compactness proves that a compact below $\sup I$ lies in $I$; algebraicity proves the other composite is the identity.
Every Scott open is determined by its restriction to $K(D)$, so restriction is monic on $\OO_D$.
Its image is an observation structure, since restriction preserves the defining operations and identifies the two families.
\item Finally
    \[
        U(x)\iff\exists k\in K(D)\,[k\leq x\wedge U(k)],
    \]
    so the two maps preserve the specified observations and are inverse effective maps.

The defining maps satisfy the following commutative diagram
\[
\begin{tikzcd}[column sep=large,row sep=large]
K(D)\ar[r,"y_{K(D)}"]\ar[d,hook,"k"']&\Idl K(D)\ar[d,"{\sup_D\Idl(k)}"]\\ D\ar[r,"\id_D"']&D
\end{tikzcd}
\]
\end{enumerate}
\end{proof}

\subsection{Algebraic domains and their maps}
\label{sec:algebraicpresentations}
Write $\mathbf{AlgED}_{\Sigma}$ for the full subcategory of $\ED$ whose underlying dcpos are algebraic. The observation structure on the compact part is part of the presentation data; it need not consist of positive predicates.

\begin{theorem}[Algebraic presentation equivalence]\label{thm:algebraicpresentation}
There is a locally ordered equivalence
\[
 \Kl(\Th)\simeq\mathbf{AlgED}_{\Sigma}.
\]
It sends $(P,\OO_P)$ to $(\Idl P,\OO_P^\dagger)$. Its canonical inverse on an algebraic effective dcpo is the observed compact part $(K(D),\OO_D|_{K(D)})$, up to the displayed completion equivalence. No search operation or positivity of the compact order is required.
\end{theorem}
\begin{proof}
\leavevmode
\begin{enumerate}[label=\textup{(\roman*)},leftmargin=*,itemsep=3pt,topsep=4pt]
\item \textbf{Image and quotient construction.} The hom-isomorphism of \cref{prop:effectivefreeextension} identifies Kleisli arrows with effective maps between free observed completions, including their pointwise order.
Essential surjectivity is \cref{thm:algebraic}.
The compact part of $\Idl P$ is the separated quotient of $P$: a compact ideal is principal because it is the directed union of its principal subideals.
The reconstruction and its observation equations are those in the proof of \cref{thm:algebraic}.
All constructions retain parameters.
\end{enumerate}
\end{proof}

Explicitly, an arrow is a monotone ideal-valued relation $m:P\to\Idl Q$ satisfying, uniformly for $V\in\OO_Q$,
\begin{equation}\label{eq:algebraicmatrix}
 m^*V(p)=\exists q\,[q\in m(p)\wedge V(q)]\in\OO_P.
\end{equation}
Its extension and the composite of two such matrices are
\[
 m^\#(I)(q)\iff\exists p\,[I(p)\wedge q\in m(p)],
 \qquad
 (n\star m)(p,r)\iff\exists q\,[q\in m(p)\wedge r\in n(q)].
\]
Identities are the ordinary order matrices. These formulas use regular images; the relations need not be $\Sigma$-classified. Observation preservation under composition follows from $(n\star m)^*=m^*n^*$, and ideality follows from directedness of the rows and monotonicity. On canonical compact presentations, an effective map $F:D\to E$ is represented by $m_F(k,l)\iff l\leq F(k)$ for $k\in K(D)$ and $l\in K(E)$. Algebraicity makes each row an ideal, and Scott openness gives \eqref{eq:algebraicmatrix}. Thus maps need not send compact elements to compact elements. In the searchable positive specialization, the coordinate observations recover the positive-matrix condition of \cref{thm:relations}.

For the computational doctrine the corresponding free-algebra equivalence is
\[
 \Kl(\J)\simeq\mathbf{FreeAlg}(\J),
\]
with the canonical compact comparison $\J K(A)\cong A$ of \cref{jd:prop:compact}. The observed and computational statements retain their different completion doctrines.

\begin{example}[Observation information beyond effective points]\label{ex:pointforgetting}
Give the discrete object $2$ either the family of $\Sigma$-constant observations or the family $\Sigma^2$ of all positive predicates. Both are observation structures: finite meets and overt joins are computed pointwise. A directed inhabited ideal of a discrete order is a singleton, so both observed domains are algebraic. Every point is effective in both, and both effective-point algebras are the same discrete algebra $2$. Nevertheless the identity from the constant-observation domain to the full positive-observation domain is not effective. Its inverse image of $b\mapsto[b=1]$ is not constant, as evaluation at $0$ and $1$ shows. Thus effective points do not recover arbitrary observation structures even in $\Eff$. The comparison with bare $\J$-algebras uses the canonical semantic objects of \cref{sec:semantic-comparison}.
\end{example}


\begin{lemma}[Interpolation]
\label{lem:interpolation}
    In a continuous dcpo, $a\way x$ implies $a\way b\way x$ for some $b$.
\end{lemma}
\begin{proof}
\leavevmode
\begin{enumerate}[label=\textup{(\roman*)},leftmargin=*,itemsep=3pt,topsep=4pt]
\item \textbf{Interpolation.} The assignment $x\mapsto\Downarrow x$ is monotone.
The union of the ideals $\Downarrow b$ for $b\way x$ is therefore an ideal $A_x$.
It has supremum $x$: every member is below $x$, while each $b\way x$ is the supremum of $\Downarrow b\subseteq A_x$.
Apply $a\way x$ to $A_x$ to obtain $d\in A_x$ with $a\leq d$.
Some $b\way x$ satisfies $d\way b$, whence $a\way b$.
\end{enumerate}
\end{proof}

\begin{theorem}[Continuous recognition]
\label{thm:continuous}
    For an effective dcpo $D$, the following are equivalent:
    \begin{enumerate}
        \item its underlying dcpo is continuous;

        \item its algebra map admits a left adjoint in $\ObsPre(\E)$, $\Downarrow_D\dashv\sup_D\dashv y_D$;

        \item it is effectively isomorphic to the splitting of a deflationary idempotent Scott-continuous endomorphism of a free $\Th$-algebra.
    \end{enumerate}
    For a continuous $D$ one has
    \begin{equation}
    \label{eq:approxobs}
        U^\dagger(\Downarrow x)=U(x)\qquad(U\in\OO_D).
    \end{equation}
\end{theorem}
\begin{proof}
\leavevmode
\begin{enumerate}[label=\textup{(\roman*)},leftmargin=*,itemsep=3pt,topsep=4pt]
\item \textbf{Scott continuity.} For continuous $D$, the equivalence
    \[
        \Downarrow x\subseteq I\quad\Longleftrightarrow\quad x\leq\sup I
    \]
    follows from continuity in one direction and the definition of $\way$ in the other.
Equation \eqref{eq:approxobs} is Scott openness applied to $x=\sup\Downarrow x$, so the left adjoint is observation-preserving.
Interpolation shows that $\Downarrow$ preserves directed suprema: if $a\way\sup_j x_j$, choose $a\way b\way\sup_j x_j$, then $b\leq x_j$ for some $j$, hence $a\way x_j$.
\item \textbf{The effective retraction.} Thus both adjoint maps in the retraction
    \[
        D\xrightarrow{\Downarrow}\Th D\xrightarrow{\sup}D
    \]
    are effective dcpo maps, and $\Downarrow\sup\leq\id$.
\item \textbf{Converse construction.} Conversely, if $t\dashv\sup$, then $t(x)\subseteq\downarrow x$ and $x\leq\sup t(x)$ give equality of this supremum with $x$.
If $a\in t(x)$ and $x\leq\sup I$, adjunction gives $a\in I$, so $a\way x$.
If $a\way x$, apply this definition to $t(x)$ and use lower closure to obtain $a\in t(x)$.
Hence $t(x)=\Downarrow x$.
\item \textbf{Continuity of retracts.} Finally let $D$ be a Scott retract of $\Idl(P)$, with section $j$ and retraction $q$.
For $p\in jx$, the element $q(\downarrow p)$ is way below $x$: apply $j$ to an ideal whose supremum is above $x$, and use compactness of $\downarrow p$.
These elements form an inhabited directed family with supremum $q(jx)=x$.
Their lower closure is contained in $\Downarrow x$ and is cofinal in it by the defining test for $\way$.
Thus $D$ is continuous.
This argument applies in particular to the displayed idempotent splittings.

The comparison is witnessed by the retraction diagram
\[
\begin{tikzcd}[column sep=large,row sep=large]
D\ar[r,"\Downarrow_D"]\ar[dr,"\id_D"']&\Th D\ar[d,"\sup_D"]\\ &D
\end{tikzcd}
\]
\end{enumerate}
\end{proof}

\begin{example}[A noncontinuous effective dcpo]\label{ex:noncontinuous}
The distinction between general and continuous observed domains occurs in $\Eff$ itself. Use its double-negation sheaf reflection $a:\Eff\to\mathbf{Set}$ and fully faithful inclusion $\nabla:\mathbf{Set}\to\Eff$ \cite{hyland,vanoosten}. Let $D_0$ be the ordered set with bottom $0$, top $1$, and two increasing chains $c_i(n)$, $i\in2$, $n\in\N$, with supremum $1$; elements on different chains have meet $0$. Form the internal ordered object $D=\nabla D_0$.

We first verify directed completeness in the original topos. For a parameterized inhabited directed lower predicate $I$ on $D$, apply $a$ to its membership relation and take its regular image in $aD=D_0$. The inhabitedness, directedness, and lowerness sequents are geometric and hence are preserved. This yields an ideal of $D_0$ in the reflected parameter context. Its ordinary supremum is a map into $D_0$. Adjunction transposes this map to a supremum in $D$. Indeed upper-bound inequalities have codomain the sheaf relation $\nabla(\leq_{D_0})$; they hold before reflection exactly when they hold after reflection, by the adjunction. This proves both upper-boundedness and leastness, uniformly in the parameter. Thus $D$ is an internal dcpo.

For each $i$, the lower ideal generated by the numerical sequence $c_i(n)$ has supremum $1$, by this calculation and preservation of the NNO by $a$. If $b\ll1$, the two ideal tests give numbers $n,m$ with $b\leq c_0(n)$ and $b\leq c_1(m)$. Their meet is $0$, so $b=0$. Conversely $0\ll1$ because every tested ideal is inhabited and lower. Therefore $\Downarrow1=\{0\}$, whose supremum is $0\ne1$, and $D$ is not continuous. Equip $D$ with the $\Sigma$-constant observation structure; it is closed under finite meets and overt joins, so this is an effective dcpo. This construction is an object of $\Eff$ throughout; the sheaf reflection is used to verify its directed-supremum map.
\end{example}


\subsection{Continuous domains as an idempotent completion}
\label{subsec:continuous-domains-as-an-idempoten-completion}

Write $\Cont(\E)$ for the full subcategory of $\C$ on continuous effective dcpos. The free-algebra functor induces a fully faithful comparison
\[
    J:\Kl(\Th)\longrightarrow\C,\qquad
    J(P)=\Th P,\quad J(m)=m^\#.
\]
The effective free-extension property identifies its hom-objects as well as its hom-sets. This gives the KZ-algebra formulation of continuous completion \cite{hofmannsousa} in which no basis is selected on an individual continuous domain.

\begin{theorem}[Idempotent-completion description]
\label{thm:karoubi}
    There is an equivalence of locally ordered categories
    \[
        \Kar\bigl(\Kl(\Th)\bigr)\simeq\Cont(\E).
    \]
    Every continuous algebra is the splitting of a deflationary idempotent on a free algebra. Allowing all idempotents gives the same essential image.
\end{theorem}
\begin{proof}
\leavevmode
\begin{enumerate}[label=\textup{(\roman*)},leftmargin=*,itemsep=3pt,topsep=4pt]
\item \textbf{Idempotent-completion description.} Finite limits in $\C$ split idempotents: if $e:A\to A$ satisfies $e^2=e$, the equalizer $j:D\incl A$ of $e$ and $\id_A$ receives a unique map $q:A\to D$ with $jq=e$, and $qj=\id_D$.
For a free $A=\Th P$, the retract argument in \cref{thm:continuous} proves continuity of $D$, even if $e$ is not deflationary.
\item \textbf{Converse construction.} Conversely, every continuous $D$ has the effective splitting
    \[
        \begin{tikzcd}[column sep=large]
            D\ar[r,shift left=.6ex,"{\Downarrow_D}"] &
            \Th D\ar[l,shift left=.6ex,"{\sup_D}"]
            \ar[loop right,"{e_D=\Downarrow_D\sup_D}"]
        \end{tikzcd}
        \qquad \sup_D\Downarrow_D=\id_D,
    \]
    and $e_D\leq\id_{\Th D}$.
Full faithfulness of $J$ transports each idempotent and each compatible map to the Kleisli category.
\item \textbf{Adjunction and universal property.} The universal property of the Karoubi envelope now gives a fully faithful functor with essential image exactly $\Cont(\E)$.
\end{enumerate}
\end{proof}

This is the free/split algebra comparison for the lifted KZ doctrine. The effective content lies in the fact that both arrows of the displayed splitting preserve the specified observations. It is this fact that permits approximation theory to remain inside $\C$.

\[
    \begin{tikzcd}[column sep=large]
        \ObsPre(\E)\ar[r,"{\Th}"]\ar[d,"{V}"'] &
        \ObsPre(\E)\ar[d,"{V}"]\\
        \Pre(\E)\ar[r,"{\T}"'] & \Pre(\E)
    \end{tikzcd}
    \qquad
    \begin{tikzcd}[column sep=large]
        \Kl(\Th)\ar[r,hook,"{J}"]\ar[d,hook] &\C\\
        \Kar(\Kl(\Th))\ar[ur,"{\simeq\ \Cont(\E)}"']&
    \end{tikzcd}
\]
The first square commutes with unit and multiplication. In the second diagram the diagonal denotes the equivalence followed by the full inclusion of continuous algebras.

\section{Dcpo presentations of general effective dcpos}
\label{sec:directedcovers}

Free completion and idempotent splitting describe algebraic and continuous domains, not every dcpo. To represent the remaining directed suprema we attach explicit cover inequalities to a preorder. The resulting universal construction gives a presentation of every observed dcpo without requiring a compact or searchable basis.


The algebraic and continuous characterizations have a common extension to all effective dcpos. The additional presentation data specify which directed suprema are to satisfy given inequalities. We internalize the directed-cover construction of Jung, Moshier, and Vickers \cite{jungmoshiervickers} and retain the chosen observations throughout.

\begin{definition}[Observed dcpo presentation]\label{def:directedcover}
A dcpo presentation consists of an internal preorder $P$, a cover-index object $C$, an endpoint map $\ell:C\to P$, and a family $I:C\to\Idl P$. The notation $\ell(c)\triangleleft I_c$ denotes the inequality to be imposed. A monotone map $f:P\to D$ into a dcpo respects the covers when
\[
 f(\ell(c))\leq\sup_{p\in I_c}f(p).
\]
A compatible upper predicate $u:P\to\Omega$ satisfies
\begin{equation}\label{eq:covercompatible}
 u(\ell(c))\Longrightarrow\exists p\,[I_c(p)\wedge u(p)].
\end{equation}
An observed presentation additionally specifies an observation structure $\OO_P$ consisting of compatible upper predicates. Equivalently, one may specify a family of compatible generators and form its observation closure as in \cref{def:obs}.
\end{definition}

\begin{lemma}[Internal directed-cover completion]\label{lem:covercompletion}
There is a dcpo $D_C(P)$ and a cover-preserving monotone map $\eta_C:P\to D_C(P)$ such that restriction induces a natural order isomorphism between Scott-continuous maps $D_C(P)\to D$ and cover-preserving monotone maps $P\to D$. This is the internal directed-cover construction corresponding to \cite[Theorem~2.7]{jungmoshiervickers}; its internal proof is recorded because the observation comparison below takes place in slices of $\Eff$.
\end{lemma}
\begin{proof}
\leavevmode
\begin{enumerate}[label=\textup{(\roman*)},leftmargin=*,itemsep=3pt,topsep=4pt]
\item \textbf{Cover-closed lower sets.} Let $\Lambda_C(P)$ be the object of lower subsets $K\subseteq P$ satisfying
\[
 I_c\subseteq K\Longrightarrow\ell(c)\in K.
\]
Arbitrary intersections preserve these conditions.
Thus $\Lambda_C(P)$ is a complete lattice, with joins given by closure of unions.
Write $\operatorname{cl}_C$ for this closure, and put $\eta_C(p)=\operatorname{cl}_C(\downarrow p)$.
Inside this lattice take the intersection of all sub-dcpos containing every $\eta_C(p)$; call it $D_C(P)$.
\item \textbf{Generation of the free dcpo.} The families of subsets and sub-dcpos used here are subobjects of power objects, so the intersections are internal constructions.
The rows $I_c$ are inhabited and directed, and closure gives
\[
 \eta_C(\ell(c))\subseteq\bigvee_{p\in I_c}\eta_C(p).
\]
This is a directed join in $D_C(P)$ as well as in the ambient lattice.
\item \textbf{Extension to a complete lattice.} First take a cover-preserving monotone $g:P\to L$ into a complete lattice.
For $z\in L$, the lower subset $\{p:g(p)\leq z\}$ is cover-closed.
It follows that
\[
 \widehat g(K)=\bigvee_{p\in K}g(p)
\]
is left adjoint to $z\mapsto\{p:g(p)\leq z\}$, preserves joins, and satisfies $\widehat g\eta_C=g$.
In particular closure adds no new upper-bound constraints beyond the stated covers.
\item \textbf{Scott-closed embedding.} For a dcpo $D$, let $\Gamma D$ be the complete lattice of Scott-closed lower subsets, where closed means closed under inhabited directed suprema.
The principal-lower-set map $j:D\to\Gamma D$ is an order embedding preserving directed suprema.
Indeed, a closed lower set contains a directed family exactly when it contains its supremum.
Hence $jf$ respects the covers and extends to a join-preserving $\widehat{jf}:\Lambda_C(P)\to\Gamma D$.
\item \textbf{Factorization through the generated dcpo.} The inverse image of the sub-dcpo $jD$ is a sub-dcpo containing $\eta_C P$, so restriction factors uniquely through a Scott-continuous map $\bar f:D_C(P)\to D$ with $\bar f\eta_C=f$.
\item \textbf{Uniqueness and order reflection.} The equalizer of two Scott-continuous maps is closed under directed suprema.
Consequently agreement on $\eta_C P$ forces agreement on $D_C(P)$.
The same argument applied to the subobject on which one map is below another proves order reflection.
This gives the asserted order isomorphism.
The constructions use power objects, images, and their slice versions, so the result is uniform in parameters.

The universal property is represented by the commutative diagram
\[
\begin{tikzcd}[column sep=large,row sep=large]
P\ar[r,"\eta_C"]\ar[dr,"f"']&D_C(P)\ar[d,"\bar f"]\\ &D
\end{tikzcd}
\]
\end{enumerate}
\end{proof}

Applying \cref{lem:covercompletion} to $\Omega$ gives the trace isomorphism
\begin{equation}\label{eq:covertrace}
 \Sc(D_C(P))\cong
 \{u\in\Up(P):u\text{ satisfies }\eqref{eq:covercompatible}\}.
\end{equation}
Write $\bar u$ for the unique extension. Compatible predicates are closed under the observation operations. For a binary meet, two witnesses in $I_c$ have a common upper bound there; upperness supplies both predicates at that bound. Inhabitedness handles constants, and existential interchange handles overt unions.

\begin{theorem}[Observed cover universal property]\label{thm:observedcovers}
Equip $D_C(P)$ with $\OO_{D_C(P)}=\{\bar u:u\in\OO_P\}$. This is an effective dcpo, and restriction induces the natural order isomorphism
\begin{align*}
 \ED(D_C(P),E)\cong\{f:P\to E\mid {}& f\text{ is monotone and respects covers},\\
                                  &Vf\in\OO_P\text{ uniformly for }V\in\OO_E\}.
\end{align*}
The comparison commutes with substitution in parameters.
\end{theorem}
\begin{proof}
\leavevmode
\begin{enumerate}[label=\textup{(\roman*)},leftmargin=*,itemsep=3pt,topsep=4pt]
\item \textbf{Observed cover universal property.} Extension preserves constants, finite meets, and overt unions.
For example $\bar u\wedge\bar v$ is Scott open and restricts to $u\wedge v$, so uniqueness in \eqref{eq:covertrace} gives $\overline{u\wedge v}=\bar u\wedge\bar v$.
The other operations are treated identically.
This proves that the displayed family is an observation structure.
\item \textbf{Scott continuity.} For a map on the right let $\bar f$ be its Scott extension.
For every $V\in\OO_E$, uniqueness of extension gives
\[
 V\bar f=\overline{Vf}\in\OO_{D_C(P)}.
\]
Thus $\bar f$ is effective.
\item \textbf{Converse construction.} Conversely, restriction of an effective map has precisely the displayed observation condition and respects the imposed directed joins.
The inverse equations and order structure follow from \cref{lem:covercompletion}.
All equalities hold in slices, proving the parameter assertion.

The universal property is represented by the commutative diagram
\[
\begin{tikzcd}[column sep=large,row sep=large]
P\ar[r,"\eta_C"]\ar[dr,"f"']&D_C(P)\ar[d,"\bar f"]\\ &E
\end{tikzcd}
\]
\end{enumerate}
\end{proof}

\begin{theorem}[Characterization of all effective dcpos]\label{thm:allcoverpresentations}
Every effective dcpo $D$ is canonically recovered from the observed presentation
\[
 P=D,\qquad C=\Idl D,\qquad I_c=c,\qquad
 \ell(c)=\sup_Dc,\qquad\OO_P=\OO_D.
\]
Let $\mathbf{Cov}_{\Sigma}$ have observed presentations as objects and cover-preserving observation maps $P\to D_{C'}(Q)$ as arrows, with composition by extension. Completion gives a locally ordered equivalence
\[
 \mathbf{Cov}_{\Sigma}\simeq\ED.
\]
\end{theorem}
\begin{proof}
\leavevmode
\begin{enumerate}[label=\textup{(\roman*)},leftmargin=*,itemsep=3pt,topsep=4pt]
\item \textbf{Characterization of all effective dcpos.} For the displayed presentation, a monotone map respects covers precisely when $f(\sup I)\leq\sup f[I]$.
The opposite inequality follows from monotonicity, so this is Scott continuity.
Compatible upper predicates are exactly Scott opens.
Thus \cref{thm:observedcovers} identifies the universal property with that of $D$ itself, including its observations.
\item \textbf{Identity law.} The identity of a presentation is $\eta_C$.
Extension composes arrows, and uniqueness of extension proves the category equations.
The same universal property proves full faithfulness; the displayed presentation proves essential surjectivity.
Notice that arrows take generators into the completed target. A requirement that they take generators to generators would not give this equivalence.

The universal property is represented by the commutative diagram
\[
\begin{tikzcd}[column sep=large,row sep=large]
D\ar[r,"\eta_C"]\ar[dr,"\id_D"']&D_C(D)\ar[d,"\overline{\id_D}"]\\ &D
\end{tikzcd}
\]
\end{enumerate}
\end{proof}

With no covers, $D_C(P)=\Idl P$, recovering \cref{thm:algebraicpresentation}. Retracts of these observed free objects give \cref{thm:karoubi}. General covers give all of $\ED$, including \cref{ex:noncontinuous}: the two covers putting the common top below each chain supremum express its two nontrivial directed joins. Thus algebraic completion, continuous retracts, and general directed-cover completion are three presentations within the same observed semantics.

\begin{proposition}[Canonical algebraic coequalizer]\label{prop:algebraiccoequalizer}
Every $\Th$-algebra $(D,a)$ is the reflexive coequalizer in $\ED$
\[
 \begin{tikzcd}[column sep=large]
 \Th\Th D\ar[r,shift left=.6ex,"\mu_D"]
              \ar[r,shift right=.6ex,"\Th a"']&
 \Th D\ar[r,"a"]&D.
 \end{tikzcd}
\]
Its common section is $\Th y_D$, and the two source objects are algebraic. Every $\J$-algebra similarly has the canonical reflexive coequalizer $\J\J A\rightrightarrows\J A\to A$ in $\AS$.
\end{proposition}
\begin{proof}
\leavevmode
\begin{enumerate}[label=\textup{(\roman*)},leftmargin=*,itemsep=3pt,topsep=4pt]
\item \textbf{Canonical algebraic coequalizer.} For either monad $M$, with unit $\eta$, let $h:MD\to(E,b)$ be an algebra map with $h\mu_D=hMa$.
Set $f=h\eta_D$.
Naturality and the monad laws give
\[
 fa=hMa\eta_{MD}=h\mu_D\eta_{MD}=h,
 \qquad bMf=bMhM\eta_D=h\mu_D M\eta_D=h.
\]
Thus $f$ is an algebra map and the unique factorization, since $a\eta_D=\id$.
Both arrows compose with $M\eta_D$ to the identity.
This proves the coequalizer and reflexivity statements.
\end{enumerate}
\end{proof}

The section $y_D$ of $a$ is an arrow of observed preorders and need not be Scott-continuous. The coequalizer therefore does not force every domain to be continuous. Nor does it describe an arbitrary directed-cover completion as a set-theoretic quotient of the ideal completion on its original generators: the generated sub-dcpo construction can introduce additional directed limits.


Observed covers therefore extend the free and split descriptions to arbitrary effective dcpos. They retain the inequalities that force directed suprema, while observation preservation is controlled by the compatible upper predicates. The following semantic comparison will identify the effective-point fragment of the canonical free case.

\section{The semantic comparison}
\label{sec:semantic-comparison}

The two domain doctrines have so far been developed separately. We now identify the precise bridge between them: effective points of a canonical observed completion are admitted ideals. The resulting adjunction compares free objects first; its restriction to splittings then compares continuous computational and semantic domains.


Recall the category $\ED$ of internal dcpos equipped with an observation family of Scott opens, closed under $\Sigma$-constants, finite meets, and overt-indexed joins. Let $\Pt D$ be its effective-point object, ordered as a subobject of $D$.

For a poset $P$ define a canonical semantic completion
\[
    \F P=\bigl(\Idl P,\{u^\dagger\mid u\in U_\Sigma P\}\bigr),
    \qquad u^\dagger(I)=\chi_I(u).
\]
This is the semantic lifted free completion applied to the canonical positive upper predicates. The family is determined before taking the ideal completion.

\begin{theorem}[Effective-point adjunction]
\label{jd:thm:adjunction}
    There is an internally enriched, locally ordered adjunction
    \[
        \F:\Pos(\E)\rightleftarrows\ED:\Pt,
        \qquad \Pt\F\cong\J.
    \]
    Its induced monad is the monad of \cref{jd:thm:J}. The canonical comparison functor
    \[
        K:\ED\longrightarrow\Alg(\J)
    \]
    equips $\Pt D$ with the supremum of its admitted ideals, calculated in $D$.
\end{theorem}
\begin{proof}
\leavevmode
\begin{enumerate}[label=\textup{(\roman*)},leftmargin=*,itemsep=3pt,topsep=4pt]
\item \textbf{Restricting effective maps to principal ideals.} An effective Scott-continuous map $\F P\to D$ restricts on principal ideals to a monotone map $f:P\to D$.
Effectiveness says exactly that every observation $U$ of $D$ has $Uf:P\to\Sigma$, hence that $f$ factors through $\Pt D$.
\item \textbf{Extending maps into effective points.} Conversely a monotone map into $\Pt D$ has a unique Scott extension, and
    \[
        U(f^\#I)=\chi_I(Uf)
    \]
    shows that its inverse image belongs to the specified observations of $\F P$.
This is the required hom-isomorphism, including its order and parameter dependence.
\item \textbf{The induced monad and counit.} The effective-point pullback for $\F P$ is exactly \cref{jd:def:J}.
The counit at $D$ sends an ideal of effective points to its supremum in $D$.
For $I\in\J(\Pt D)$ and $U\in\OO_D$, the predicate $U|_{\Pt D}$ is monotone and $\Sigma$-valued, so
    \[
        U\Bigl(\bigvee_{x\in I}x\Bigr)=\chi_I(U|_{\Pt D})\in\Sigma.
    \]
    The supremum is therefore an effective point.
On a free completion this operation is union, identifying the induced multiplication with $\mu$.

The universal property is represented by the commutative diagram
\[
\begin{tikzcd}[column sep=large,row sep=large]
P\ar[r,"\eta^F_P"]\ar[d,"g"']&\Pt\F P\ar[d,"\Pt(\bar g)"]\\ \Pt D\ar[r,"\id"']&\Pt D
\end{tikzcd}
\]
\end{enumerate}
\end{proof}

For later calculations, denote the adjunction unit by $\eta^{F}_P:P\to\Pt\F P$ and its counit by $\epsilon^{F}_D:\F\Pt D\to D$. Their formulas are
\[
    \eta^{F}_P(p)=\downarrow p,\qquad
    \epsilon^{F}_D(I)=\bigvee_{x\in I}j_D(x),
\]
where $j_D:\Pt D\hookrightarrow D$. The first triangle on $\F P$ says that an ideal is the union of its principal subideals. The second says that the supremum of the principal ideal of an effective point is that point. The hom-isomorphism preserves order because it and its inverse are given by restriction and directed supremum. It is internal: replacing an arrow by a family over a parameter gives the same two formulas. Thus the adjunction identifies the internal hom-objects and their external hom-sets. Its monad multiplication is $\Pt(\epsilon^F_{\F P})$, which is the union operation already computed for $\J$.

Let $\CS$ be the full subcategory of $\ED$ consisting of retracts of the canonical completions $\F P$. It includes the searchable rounded completions of \cref{dom:sec:rounded}: their approximation idempotents split effectively on $\F B$.

\begin{theorem}[Continuous comparison]
\label{jd:thm:splitcomparison}
    The comparison functor restricts to an equivalence
    \[
        \CS\simeq\Spl(\Alg(\J))
        \simeq\Kar\Kl(\J).
    \]
    Consequently a monotone map $f:\Pt D\to\Pt E$, for $D,E\in\CS$, extends to an effective Scott-continuous map $D\to E$ exactly when it is a $\J$-algebra morphism. Its extension is unique.
\end{theorem}
\begin{proof}
\leavevmode
\begin{enumerate}[label=\textup{(\roman*)},leftmargin=*,itemsep=3pt,topsep=4pt]
\item \textbf{Inverse comparison.} The adjunction gives
    \[
        \ED(\F P,\F Q)\cong\Pos(P,\J Q)
        \cong\Alg(\J)(\J P,\J Q),
    \]
    compatibly with identities, composition, order, and parameters.
Thus $K$ is an equivalence between the two full categories of free objects.
Taking idempotent completions preserves that equivalence.
On the semantic side the result is $\CS$; on the computational side it is the split algebras by \cref{jd:prop:continuous}.
 This establishes full faithfulness of the restricted comparison.

Naturality is expressed by the commutative diagram
\[
\begin{tikzcd}[column sep=large,row sep=large]
\ED(\F P,\F Q)\ar[r,"\cong"]\ar[d,"K"']&\Pos(P,\J Q)\ar[d,"(-)^\#"]\\ \AS(\J P,\J Q)\ar[r,"\id"']&\AS(\J P,\J Q)
\end{tikzcd}
\]
\end{enumerate}
\end{proof}

Thus $\CS$ has semantic retracts of free observed completions as objects and effective Scott-continuous maps as morphisms. The comparison $\Pt$ identifies it with the continuous $\J$-algebras. For searchable presentations, \cref{lem:roundedadmission,prop:roundedhomcomparison} calculate this comparison on ideals and arrows. This is the retract row of the dictionary in \cref{sec:categorydictionary}; it does not restrict all of $\ED(\E)$ to searchable objects.

Here split algebras are defined by an adjoint to their algebra map. It does not identify every arbitrary observation-equipped dcpo with a split algebra, nor does it identify all complete or well-complete modest objects of other accounts with this category.

\begin{proposition}[Comparison of approximation]
\label{jd:prop:waycomparison}
    For $D\in\CS$, the $\J$-way-below relation on $\Pt D$ is the restriction of the ordinary way-below relation of $D$.
\end{proposition}
\begin{proof}
\leavevmode
\begin{enumerate}[label=\textup{(\roman*)},leftmargin=*,itemsep=3pt,topsep=4pt]
\item \textbf{Adjunction and universal property.} Choose a semantic retraction $s:D\to\F B$, $q:\F B\to D$.
Both approximation relations, evaluated at effective $x,y$, have the description
    \[
        x\leq q(\downarrow b)\quad\text{for some }b\in s(y).
    \]
    On the computational side this follows from the formula for the left adjoint in the proof of \cref{jd:prop:continuous}, applied to the image retraction under $K$.
On the semantic side the points $q(\downarrow b)$ with $b\in s(y)$ form a directed family with supremum $y$.
\item \textbf{Ideal and directedness calculation.} Each is way below $y$: applying $s$ to a directed supremum above $y$ and using compactness of $\downarrow b$ supplies the required stage.
Conversely any point way below $y$ is below a member of this family.
This proves the same formula on both sides, independently of the chosen retraction.

The comparison is witnessed by the retraction diagram
\[
\begin{tikzcd}[column sep=large,row sep=large]
D\ar[r,"s"]\ar[dr,"\id_D"']&\F B\ar[d,"q"]\\ &D
\end{tikzcd}
\]
\end{enumerate}
\end{proof}

\paragraph{Observations become intrinsic morphisms.}
The dominance itself is a $\J$-algebra, with $a_\Sigma(I)=\chi_I(\id_\Sigma)$. It is also $\J 2$, where $2$ is the ordered two-element chain. On the semantic side $\F 2$ is the semantic Sierpi\'nski domain: its carrier is $\Omega$ and its effective points are $\Sigma$. Therefore \cref{jd:thm:splitcomparison} gives
\[
    \OO_D\cong\Alg(\J)(\Pt D,\Sigma)
    \qquad(D\in\CS).
\]
For searchable domains, the semantic observation-transpose criterion and preservation of $\Sigma$-admissible ideal suprema consequently characterize the same extensions. The Chu formulation is the point--observation representation of this comparison.


The comparison equivalence identifies continuous admitted algebras with a specific retract category of observed completions, not with every observed dcpo. This boundary matters when later changes of base act on full completion objects but preserve only their represented effective-point images.

\section{Searchable approximation presentations}
\label{dom:sec:rounded}

The abstract comparison becomes computational when a basis supports positive search and a semidecidable approximation relation. We encode that relation as a deflationary Kleisli idempotent, then identify its rounded fixed points and observations. Importantly, this requires no semidecision procedure for the completed order.


\begin{definition}
\label{def:approx}
    An \textbf{approximation presentation} is a preorder $B$ and a Kleisli endomorphism $r:B\to\Idl(B)$ such that $r\leq y_B$ and $r\star r=r$. Put
    \[
        \rho(a,b)\iff a\in r(b).
    \]
    The presentation is \textbf{searchable} when $B$ is $\Sigma$-overt and $\rho$ factors through $\Sigma$.
\end{definition}

The defining equations express refinement and interpolation:
\begin{equation}
\label{eq:rho}
    \rho(a,b)\Longrightarrow a\leq b,
    \qquad
    \rho(a,b)\iff\exists c\,[\rho(a,c)\wedge\rho(c,b)].
\end{equation}
Neither the order of $B$ nor inclusion of completed ideals is required to be $\Sigma$-classified.

\begin{theorem}[Rounded completion]
\label{thm:rounded}
    For an approximation presentation, the object
    \[
        D_r=\{I\in\Idl(B)\mid r^\# I=I\}
    \]
    is a continuous dcpo, and its directed suprema are unions. Its basic family is $\kappa(b)=r(b)$, and
    \begin{equation}\label{eq:roundedreconstruct}
        I=\bigvee_{b\in I}r(b),\qquad
        I(b)\Longrightarrow r(b)\way I.
    \end{equation}
    Every continuous dcpo has such a presentation, taking $B=D$ and $r(x)=\Downarrow x$.

    The construction is the idempotent splitting
    \[
        \begin{tikzcd}[column sep=large]
            D_r\ar[r,hook,shift left=.6ex,"{j_r}"]&
            \Idl B\ar[l,shift left=.6ex,"{q_r}"]\ar[loop right,"{r^\#}"]
        \end{tikzcd}
        \qquad j_rq_r=r^\#,\quad q_rj_r=\id.
    \]
\end{theorem}
\begin{proof}
\leavevmode
\begin{enumerate}[label=\textup{(\roman*)},leftmargin=*,itemsep=3pt,topsep=4pt]
\item \textbf{Rounded completion.} The extension $r^\#$ is a Scott-continuous idempotent below the identity.
Its fixed object is closed under directed unions.
Each row $r(b)$ is fixed, by $r\star r=r$.
The family in \eqref{eq:roundedreconstruct} is inhabited and directed by ideality of $I$ and monotonicity of $r$; its union is exactly $r^\#I=I$.
\item \textbf{Scott continuity.} If $I(b)$ and $I\subseteq\bigcup_jJ_j$, some $J_j$ contains $b$.
Lower closure and $r(b)\subseteq\downarrow b$ imply $r(b)\subseteq J_j$.
Hence $r(b)\way I$.
These approximants generate $I$, giving continuity as in \cref{thm:continuous}.
\item \textbf{Converse construction.} Conversely, take $B=D$ and $r(x)=\Downarrow x$.
Continuity makes each row an inhabited directed lower ideal;
    monotonicity of way-below makes $r$ monotone.
Testing against the
    principal ideal gives $r(x)\subseteq\downarrow x$, and the
    constructive interpolation of \cref{lem:ordinaryinterpolation}
    gives $r\star r=r$.
Define $a:D_r\to D$ by $a(I)=\sup I$ and
    $t:D\to D_r$ by $t(x)=\Downarrow x$.
Continuity gives $at=\id$.
If $b\in I$, roundedness gives $b\way c$ for some $c\in I$, so
    $b\way\sup I$.
\item \textbf{Further converse construction.} Conversely, $b\way\sup I$ puts $b$ in $I$ by
    the ideal test defining way-below.
Hence $ta=\id$.
This constructs
    the presentation and its comparison, including both inverse
    identities.

The defining arrows form the commutative diagram
\[
\begin{tikzcd}[column sep=large,row sep=large]
D_r\ar[r,hook,"j_r"]\ar[dr,"\id_{D_r}"']&\Idl B\ar[d,"q_r"]\\ &D_r
\end{tikzcd}
\]
\end{enumerate}
\end{proof}

\subsection{The Karoubi representation of rounded maps}\label{app:matrices}
\begin{theorem}[Rounded distributor representation]\label{thm:rounded-matrix}
Let $C=(B,r)$ and $C'=(B',r')$ be approximation presentations, with $D_C=D_r$ and $\kappa_C(b)=r(b)$. Scott-continuous maps $F:D_C\to D_{C'}$ correspond naturally to monotone ideal-valued maps $m:B\to\Idl(B')$ satisfying
\[
 r'\star m=m=m\star r.
\]
The formulas are
\[
 m_F(b)=F(\kappa_Cb),\qquad
 F_m(I)(c)\iff\exists b\,[I(b)\wedge c\in m(b)].
\]
The identity matrix is $r$, and composition is the Kleisli convolution \eqref{eq:kleisli}.
\end{theorem}
\begin{proof}
\leavevmode
\begin{enumerate}[label=\textup{(\roman*)},leftmargin=*,itemsep=3pt,topsep=4pt]
\item \textbf{Ideal and directedness calculation.} Write $j_C:D_C\incl\Idl(B)$ and $q_C:\Idl(B)\to D_C$ for the splitting of \cref{thm:rounded}, so $j_Cq_C=r^\#$ and $q_Cj_C=\id$, and similarly for $C'$.
The free-completion adjunction of \cref{prop:freeextension} identifies ideal-valued maps $m:B\to\Idl(B')$ with Scott-continuous maps $m^\#:\Idl(B)\to\Idl(B')$.
Under this identification, the rounded equations are exactly
\[
 (r')^\#m^\#=m^\#=m^\#r^\#.
\]
Thus they select the arrows between the two idempotents in the Karoubi envelope.
\item \textbf{Compatibility with the free splittings.} A map $F:D_C\to D_{C'}$ gives the compatible free map $j_{C'}Fq_C$, whose restriction to generators is
$j_{C'}Fq_C\eta_B=F\kappa_C=m_F$.
Conversely a compatible extension gives $q_{C'}m^\#j_C$.
The two composites are identities, since
\[
 q_{C'}(j_{C'}Fq_C)j_C=F,
 \qquad
 j_{C'}(q_{C'}m^\#j_C)q_C=(r')^\#m^\#r^\#=m^\#.
\]
Its membership formula is the Kleisli extension \eqref{eq:kleisli}, giving precisely $F_m$.
Identity and composition are inherited from the free-completion comparison, hence have matrices $r$ and existential convolution.
\item \textbf{Universal-property argument.} Restriction, extension and composition with the splitting maps are monotone, so the correspondence is natural and order-compatible.

The comparison is witnessed by the retraction diagram
\[
\begin{tikzcd}[column sep=large,row sep=large]
D_r\ar[r,hook,"j_r"]\ar[d,"F"']&\Idl B\ar[d,"m^\#"]\\ D_s\ar[r,hook,"j_s"']&\Idl C
\end{tikzcd}
\]
\end{enumerate}
\end{proof}

Thus these presented continuous domains form the indicated full subcategory of the Karoubi envelope of the ordinary ideal monad. The splitting retains the basis, its approximation relation and its interpretation maps. Searchability will select the effective arrows below.

\begin{example}[Continuity does not supply searchability]
\label{ex:continuous-not-searchable}
Let $K$ be diagonal halting and give
$A=\{e:\N\mid\neg K(e)\}$ the discrete order. An inhabited
directed subset of a discrete poset has a unique element, so its
supremum is that element and way-below is equality. Thus $A$ is a
continuous dcpo. It is not $\Sigma$-overt: otherwise the positive
predicate $(a,e)\mapsto[a=e]$ would give
\[
 e\longmapsto\exists a:A\,[a=e]\quad\in\Sigma^\N,
\]
classifying nonhalting, contrary to \cref{cor:nonsemidecidableorder}.

No alternative searchable approximation presentation of this dcpo
exists. In any rounded presentation of a discrete dcpo, every ideal
$I$ has a member $b$, and \eqref{eq:roundedreconstruct} gives
$\kappa(b)\leq I$. Discreteness forces equality, so the basic-family
map is regular epic. An overt basis would make its image overt by
\cref{lem:overtclosure}, contradicting the preceding argument.
The ordinary approximation equations are therefore constructed by
\cref{thm:rounded}, whereas searchability is a substantive property
of the resulting presentation.
\end{example}

For an upper predicate $S$ on $B$, define the upward rounding operator
\[
    (\rho^\uparrow S)(b)\iff\exists a\,[S(a)\wedge\rho(a,b)].
\]
A fixed point is automatically upper even if upperness was not required initially: use monotonicity of $r$ in its argument. If the presentation is searchable, this operation restricts to $\Sigma^B$ by overtness.

\begin{proposition}[Rounded trace representation]
\label{prop:roundedtrace}
    Scott opens of $D_r$ are in bijection with the fixed points of $\rho^\uparrow$ in $\Omega^B$. The inverse maps are
    \begin{equation}
    \label{eq:traces}
        \operatorname{tr}(U)(b)=U(r(b)),
        \qquad
        \langle S\rangle(I)\iff\exists b\,[I(b)\wedge S(b)].
    \end{equation}
    Equivalently, with $\delta_b(I)\iff b\in I$,
    \begin{equation}\label{eq:coordinate-expansion}
        U=\bigvee_{b:U(\kappa b)}\delta_b.
    \end{equation}
\end{proposition}
\begin{proof}
\leavevmode
\begin{enumerate}[label=\textup{(\roman*)},leftmargin=*,itemsep=3pt,topsep=4pt]
\item \textbf{Rounded trace representation.} Apply Scott openness to \eqref{eq:roundedreconstruct}; this recovers $U$ from its trace.
Apply the same equation to $r(b)$ to see that the trace is rounded.
Conversely a rounded predicate is upper, so its extension to $\Idl(B)$ is Scott open by \cref{lem:tracefree}.
Restrict to $D_r$.
Its trace is $\rho^\uparrow S=S$, which proves both inverse identities.
\end{enumerate}
\end{proof}

\begin{definition}
\label{def:O-serachable-presentation}
    For a searchable presentation set
    \[
        \OO_r=\{U\in\Sc(D_r)\mid \operatorname{tr}(U)\text{ factors through }\Sigma^B\}.
    \]
    Equivalently, $\OO_r$ is the object of $\Sigma$-valued rounded traces, interpreted by \eqref{eq:traces}.
\end{definition}

\begin{lemma}
\label{lem:roundedobs}
    The pair $(D_r,\OO_r)$ is an effective continuous dcpo. Each coordinate open
    \[
        \delta_a(I)=I(a)
    \]
    belongs to $\OO_r$, uniformly in $a$, with trace $\rho(a,-)$.
\end{lemma}
\begin{proof}
\leavevmode
\begin{enumerate}[label=\textup{(\roman*)},leftmargin=*,itemsep=3pt,topsep=4pt]
\item \textbf{Positive rounded traces.} Trace preserves constants, finite meets, and joins. A constant $p\in\Sigma$ is rounded because each row $r(b)$ is inhabited.
Finite meets have $\Sigma$-valued traces; an overt union of traces remains $\Sigma$-valued by restricted existential quantification.
Thus \cref{thm:recognition} applies.
Coordinate evaluation preserves directed unions, and its trace is the displayed semidecidable relation.
\end{enumerate}
\end{proof}

\begin{proposition}[Splitting in the lifted doctrine]
\label{prop:effectivesplitting}
    Put $\OO_B^\Sigma=\Up(B)\cap\Sigma^B$. A searchable approximation $r$ is a subidentity idempotent in the Kleisli category of $\Th$ on $(B,\OO_B^\Sigma)$. Its induced idempotent on the free algebra splits effectively as
    \[
        \Th(B,\OO_B^\Sigma)\xrightarrow{r^\#}(D_r,\OO_r)
        \xrightarrow{j}\Th(B,\OO_B^\Sigma),
    \]
    where $j$ is inclusion and the notation on the first arrow denotes the corestriction.
\end{proposition}
\begin{proof}
\leavevmode
\begin{enumerate}[label=\textup{(\roman*)},leftmargin=*,itemsep=3pt,topsep=4pt]
\item \textbf{Splitting in the lifted doctrine.} The $\Sigma$-valued upper predicates form an observation structure by finite closure and overt existential quantification.
For $U\in\OO_B^\Sigma$,
    \[
        r^*(U^\dagger)(b)=\exists_B^\Sigma a\,[\rho(a,b)\wedge U(a)],
    \]
    which is again a $\Sigma$-valued upper predicate.
Thus $r$ is an effective Kleisli arrow and retains its subidentity and idempotent equations in the lifted doctrine.
\item \textbf{Compatibility with the splitting.} Restriction of $U^\dagger$ along $j$ has trace $\rho^\uparrow U$, so $j$ is effective.
Conversely, for a rounded $\Sigma$-trace $S$,
    \begin{align*}
        \langle S\rangle(r^\#I)
        &\iff\exists a,b\,[I(b)\wedge\rho(a,b)\wedge S(a)]\\
        &\iff\exists b\,[I(b)\wedge S(b)]=S^\dagger(I).
    \end{align*}
    Hence the corestriction of $r^\#$ is effective.
Their composites are the identity on $D_r$ and the induced idempotent on the free algebra.
Moreover every observation of $D_r$ is the restriction of $S^\dagger$ for its own rounded trace $S$.
Thus the splitting carries precisely the observation structure $\OO_r$.
\end{enumerate}
\end{proof}

\begin{theorem}[Compatibility of effective points]
\label{thm:compatibility}
    There is a canonical isomorphism over $D_r$
    \begin{equation}
    \label{eq:pointcompat}
        \Pt(D_r)\cong D_r\times_{\Omega^B}\Sigma^B.
    \end{equation}
    Under it, observation is
    \[
        \operatorname{Obs}_r(I,S)=\exists_B^\Sigma b\,[I(b)\wedge S(b)].
    \]
    In particular each basic point $r(b)$ is effectively parametrized by $B$.
\end{theorem}
\begin{proof}
\leavevmode
\begin{enumerate}[label=\textup{(\roman*)},leftmargin=*,itemsep=3pt,topsep=4pt]
\item \textbf{Compatibility of effective points.} If $I$ is $\Sigma$-classified on $B$, the displayed expression lies in $\Sigma$ by overtness and equals every specified Scott observation by \eqref{eq:traces}.
This gives one map into the pullback of \cref{def:points}.
\item \textbf{Converse construction.} Conversely an effective point has a $\Sigma$-valued profile on all coordinate opens $\delta_a$, hence its underlying membership predicate $a\mapsto I(a)$ factors through $\Sigma$.
Both maps leave the underlying ideal unchanged, so they are inverse.
These are maps of objects, since all expressions are jointly uniform in $I$, $S$, and $a$.
The basic-family statement uses $r(b)(a)=\rho(a,b)$.
\end{enumerate}
\end{proof}

\begin{lemma}[Admission of rounded ideals]
\label{lem:roundedadmission}
    Let $(B,r)$ be searchable. For a rounded ideal $I\in D_r$, the following conditions are equivalent, uniformly in every parameter context:
    \[
        I\in\J B
        \quad\Longleftrightarrow\quad
        (b\mapsto I(b))\text{ factors through }\Sigma.
    \]
    Consequently the fixed subalgebra of $r^\#:\J B\to\J B$ is canonically $\Pt(D_r)$, as a subobject of $\Idl(B)$.
\end{lemma}
\begin{proof}
\leavevmode
\begin{enumerate}[label=\textup{(\roman*)},leftmargin=*,itemsep=3pt,topsep=4pt]
\item \textbf{Admission of rounded ideals.} Fix $a:B$.
The predicate $u_a(b)=\rho(a,b)$ is upper in $b$ and $\Sigma$-valued, jointly in $(a,b)$.
If $I$ is admitted, the pullback defining $\J B$ therefore supplies the $\Sigma$-valued map
    \[
        (I,a)\longmapsto\chi_I(u_a)
        =\exists b\,[I(b)\wedge\rho(a,b)]=I(a).
    \]
    The final equality is precisely rounding.
This proves the forward direction without making the order of $B$ semidecidable.
\item \textbf{Converse construction.} Conversely, if membership in $I$ is $\Sigma$-valued, then for every positive upper predicate $u$ the formula $\exists_B^\Sigma b\,[I(b)\wedge u(b)]$ is in $\Sigma$ by overtness.
Its uniformity gives the required factorization through the defining pullback, so $I\in\J B$.
\item \textbf{Positive rows and fixed points.} Every row $r(b)$ has positive membership and is therefore admitted.
Its extension restricts to $\J B$ by the monad laws.
The fixed object is closed under the admitted unions, since $r^\#$ is an algebra endomorphism.
The equivalence just proved identifies this fixed object with the pullback in \cref{thm:compatibility}; the algebra structures agree because both compute a supremum by union.
\end{enumerate}
\end{proof}


\subsection{Concrete bases and reconstruction}
\label{subsec:reconstruction}

We next connect the abstract idempotent presentation to a continuous dcpo that is initially given by a basis. A family $\kappa:B\to D$ is a \textbf{Scott basis} if every Scott open satisfies
\begin{equation}
\label{eq:scottbasis}
     U(x)\iff\exists b:B\,[U(\kappa b)\wedge\kappa b\way x].
\end{equation}
Give $B$ the induced preorder $a\leq b\iff\kappa a\leq\kappa b$. No injectivity of $\kappa$ is required. The basic predicates $x\mapsto d\way x$ are Scott open: interpolation followed by the defining directed-supremum test for $\way$ proves inaccessibility, and monotonicity gives upperness.

\begin{theorem}[Reconstruction from a Scott basis]
\label{thm:concretebasis}
    Put $\rho(a,b)\iff\kappa a\way\kappa b$. Then $r(b)(a)=\rho(a,b)$ is an approximation idempotent and
    \[
        D\cong D_r,\qquad
        x\longmapsto W_x=\{b\mid\kappa b\way x\},\qquad
        I\longmapsto\bigvee_{b\in I}\kappa b.
    \]
    If $B$ is overt and $\rho$ is $\Sigma$-classified, this becomes an effective isomorphism when $D$ is equipped with the Scott opens having $\Sigma$-valued traces along $\kappa$.
\end{theorem}
\begin{proof}
\leavevmode
\begin{enumerate}[label=\textup{(\roman*)},leftmargin=*,itemsep=3pt,topsep=4pt]
\item \textbf{Scott continuity.} Apply \eqref{eq:scottbasis} to top to see that $W_x$ is inhabited.
Its lower closure follows from $\kappa a\leq\kappa b\way x\Rightarrow\kappa a\way x$.
If $b,c\in W_x$, apply the basis equation to the Scott open
    \[
        \{z\mid\kappa b\way z\}\cap\{z\mid\kappa c\way z\}
    \]
    at $x$.
It gives $d\in W_x$ with $\kappa b,\kappa c\way\kappa d$, so $d$ bounds $b,c$.
Applying the same argument to a single basic open proves $W_x(a)\iff\exists b\,[W_x(b)\wedge\rho(a,b)]$.
\item \textbf{Reconstruction of rounded ideals.} Thus $W_x$ is a rounded ideal, and specializing $x=\kappa c$ gives both ideality of each row and interpolation of $\rho$.
Monotonicity and refinement follow from those of $\way$, proving that $r$ is the required idempotent.
\item \textbf{Universal-property argument.} Every $\kappa b$ with $b\in W_x$ lies below $x$.
Conversely, if $d\way x$, apply the basis equation to $z\mapsto d\way z$ to find $b\in W_x$ with $d\leq\kappa b$.
Continuity then gives $x=\bigvee_{b\in W_x}\kappa b$.
\item \textbf{Scott continuity.} For the other composite, let $x_I=\bigvee_{c\in I}\kappa c$.
If $b\in I$, rounding gives $c\in I$ with $\kappa b\way\kappa c\leq x_I$, so $b\in W_{x_I}$.
If $b\in W_{x_I}$, Scott openness of $z\mapsto\kappa b\way z$ gives $c\in I$ with $\kappa b\way\kappa c$.
Hence $b\leq c$, and lower closure yields $b\in I$.
The two maps are therefore inverse order maps and preserve all directed suprema.
\item \textbf{Identification of the basic families.} Finally the isomorphism sends $\kappa b$ to $r(b)$.
Restriction of a Scott open along these two families therefore has the same trace.
The specified observation structures correspond exactly, proving the effective assertion.

The comparison is witnessed by the retraction diagram
\[
\begin{tikzcd}[column sep=large,row sep=large]
D\ar[r,"{x\mapsto W_x}"]\ar[dr,"\id_D"']&D_r\ar[d,"{I\mapsto\sup\kappa[I]}"]\\ &D
\end{tikzcd}
\]
\end{enumerate}
\end{proof}

\begin{example}[An effective continuous domain that is not algebraic]
\label{ex:lowerreals}
    In $\Eff$ take the rational numbers $\mathbb Q$ with their ordered structure and let $\rho(a,b)$ mean $a<b$. The rationals have an effective enumeration, so they are overt; strict inequality is decidable, hence $\Sigma$-classified. Each strict lower ray is an inhabited directed lower set. Density of the rational order gives $\rho=\rho\circ\rho$. Thus this is a searchable approximation presentation.

    Its rounded ideals are inhabited lower cuts $I\subseteq\mathbb Q$ with no greatest element; the whole of $\mathbb Q$ is allowed. Directed joins are unions, and the basic point at $b$ is $(-\infty,b)$. Every such ideal has
    \[
        I=\bigvee_{b\in I}(-\infty,b).
    \]
    There are no compact elements. Indeed, compactness of $I$ applied to this directed expression would give $b\in I$ with $I\subseteq(-\infty,b)$, contradicting $b\not<b$. Since the domain has the point $(-\infty,0)$, it cannot be algebraic. Its effective points are exactly the semidecidable inhabited lower cuts, and $\delta_a(I)$ is the positive test that the cut contains $a$. The global \emph{effective} points, namely $\ED(\Eff)(1,D)\cong\Eff(1,\Pt D)$, are the lower semicomputable extended real values with nonempty rational lower cut. This does not classify all global elements of the underlying object $D$. This example separates effective continuity from effective algebraicity within the fixed computational model.
\end{example}


\begin{proposition}[Global effective lower reals]
\label{prop:effective-lower-reals}
For the domain $D$ of \cref{ex:lowerreals}, global elements of $\Pt D$ correspond exactly to inhabited rounded lower cuts whose rational membership is computably enumerable. The underlying object $D$ also has global elements with rational membership that is not computably enumerable.
\end{proposition}
\begin{proof}
\leavevmode
\begin{enumerate}[label=\textup{(\roman*)},leftmargin=*,itemsep=3pt,topsep=4pt]
\item \textbf{From positive membership to enumeration.} The coordinate observations of an effective point give a classifier $\chi:\QQ\to\Sigma$. The rational assembly is partitioned, and the stream object $2^\N$ is modest by \cref{lem:modestclosure}. Apply \cref{lem:modest-target-lift} to the regular epimorphism $2^\N\onto\Sigma$. The resulting tracked stream on rational names enumerates the cut by dovetailing rational names and finite hit stages.
\item \textbf{From enumeration to an effective point.} Given an inhabited rounded c.e. lower cut, bounded simulation of its enumeration gives a total binary stream whose finite-hit image is its membership predicate. Select one enumerated rational to witness inhabitedness. For lower closure, search the enumeration for the smaller rational; for directedness use the larger of the two named rationals; for rounding search for an enumerated rational strictly above the supplied member. The defining properties of the cut prove termination of these searches on their stated domains. These programs realize the ideal laws, and \cref{lem:testinggenerators} extends the positive membership classifier to the whole observation profile. Equality of cuts is equality of their support predicates, hence equality of the corresponding points. The whole rational cut gives $+\infty$; the empty cut is excluded by inhabitedness.
\item \textbf{Ordinary global elements.} For an external real $\alpha\in[0,1]$, form the partitioned assembly
\[
 C_\alpha=\{(q,r)\in\QQ^2\mid q<r<\alpha\},
\]
named by rational-pair codes, and take the regular image of its first projection:
\[
\begin{tikzcd}[column sep=large,row sep=large]
 C_\alpha\ar[r,two heads]\ar[dr,"\pi_1"']&I_\alpha\ar[d,hook]\\
 &\QQ.
\end{tikzcd}
\]
The pair $(-2,-1)$ witnesses inhabitedness. A witness $(q,r)$ transports to $(q',r)$ when $q'<q$. For two witnesses choose the one with larger first coordinate, giving directedness. Rounding replaces $q$ by $(q+r)/2$ and retains $r$. All these operations are tracked rational operations. Consequently $I_\alpha$ is an internal rounded ideal and determines $1\to D$, with external rational support $\{q:q<\alpha\}$. For any $\alpha$ whose lower cut is not c.e., part~(i) proves that this global element does not factor through $\Pt D$. Such reals exist because there are only countably many c.e. subsets of $\QQ$, whereas the real interval is uncountable in the external set-theoretic metatheory.
\end{enumerate}
\end{proof}

Rounded ideals provide semantic directed completeness; the positive approximation matrix makes their effective points and observations presentable. The next task is to characterize maps between these completions by the same finite-information data, rather than only to reconstruct their objects.

\section{Effective maps and rounded distributors}
\label{sec:relations}

Recovering objects from a basis is only half of a realizability comparison. We must also characterize their morphisms without requiring a map to send approximants to approximants. Rounded ideal-valued matrices supply exactly those arrows; positivity characterizes when their extensions preserve effective observations.


Let $(B,r)$ and $(C,s)$ be searchable presentations. For a Scott-continuous $F:D_r\to D_s$ put
\[
    m_F(b,c)\iff F(r(b))(c).
\]
This is an ideal-valued Kleisli arrow, and its rows already lie in $D_s$.

\begin{theorem}[Effective-map recognition]
\label{thm:relations}
    For such an $F$, the following are equivalent:
    \begin{enumerate}
        \item $F$ is an effective dcpo morphism;

        \item $m_F$ is $\Sigma$-classified;

        \item $F$ restricts to a morphism $\Pt(D_r)\to\Pt(D_s)$ over its underlying map.
    \end{enumerate}
    Effective maps are in bijection with $\Sigma$-classified ideal-valued relations $m:B\to\Idl(C)$ satisfying
    \begin{equation}
    \label{eq:roundedmap}
        m=s\star m\star r.
    \end{equation}
    The inverse correspondence is $F_m(I)=m^\#I$. Identities correspond to $r$, and composition corresponds to $\star$.

    Equivalently, the free extensions and their splittings fit into
    \[
        \begin{tikzcd}[column sep=large]
            \Idl B\ar[r,"{m^\#}"]\ar[d,"{q_r}"']&\Idl C\ar[d,"{q_s}"]\\
            D_r\ar[r,"{F_m}"']&D_s.
        \end{tikzcd}
    \]
    The roundedness equation is precisely what makes this factorization compatible with both splittings.
\end{theorem}
\begin{proof}
\leavevmode
\begin{enumerate}[label=\textup{(\roman*)},leftmargin=*,itemsep=3pt,topsep=4pt]
\item \textbf{Rounded distributor equations.} The ordinary map--matrix correspondence, its inverse and its composition law are \cref{thm:rounded-matrix}.
For idempotents $r,s$, the equation $m=s\star m\star r$ is equivalent to the two equations $s\star m=m=m\star r$: substitute the displayed equation and use idempotence.
Hence it is exactly the condition for an arrow between the two splittings.
\item \textbf{From observations to positive matrices.} If $F$ preserves observations, pull back the coordinate opens $\delta_c$ and evaluate their traces on the effective basic family.
This gives the jointly positive predicate $m_F(b,c)$.
\item \textbf{From positive matrices to effective maps.} Conversely a positive matrix gives, for every rounded effective trace $S$ on $C$,
\begin{equation}\label{eq:inverseobs}
 \operatorname{tr}(F^*\langle S\rangle)(b)
 =\exists_C^\Sigma c\,[m_F(b,c)\wedge S(c)].
\end{equation}
Overtness of $C$ constructs this classifier jointly in $(S,b)$, and \cref{prop:roundedtrace} identifies its underlying Scott open with $F^*\langle S\rangle$.
This proves observation preservation uniformly.
\item \textbf{Restriction to effective points.} For an effective point $I$, the extension has positive membership
\[
 c\longmapsto\exists_B^\Sigma b\,[I(b)\wedge m_F(b,c)].
\]
Thus \cref{thm:compatibility} gives the restriction to effective points.
\item \textbf{Converse construction.} Conversely such a restriction, composed with the effective basic family, makes $F(r(b))(c)$ jointly positive.
The splitting square is the extension of the two rounded equations, so the ordinary correspondence restricts to the asserted effective one.
\end{enumerate}
\end{proof}

\begin{proposition}[The morphisms in the comparison]
\label{prop:roundedhomcomparison}
    For searchable $(B,r)$ and $(C,s)$, the rounded $\J$-Kleisli arrows $m:B\to\J C$ are exactly the rounded ideal-valued relations whose membership matrix $m(b,c)$ is $\Sigma$-classified. Their extensions are exactly the algebra maps between the corresponding fixed algebras, and exactly the effective Scott-continuous maps $D_r\to D_s$ under passage to effective points.
\end{proposition}
\begin{proof}
\leavevmode
\begin{enumerate}[label=\textup{(\roman*)},leftmargin=*,itemsep=3pt,topsep=4pt]
\item \textbf{The morphisms in the comparison.} A rounded arrow satisfies $s\star m=m=m\star r$.
Thus each row $m(b)$ is fixed by $s^\#$.
Apply \cref{lem:roundedadmission} to these rows in the context $b:B$.
Admission implies that $m(b,c)$ is jointly $\Sigma$-valued.
\item \textbf{Converse construction.} Conversely positive rows are admitted by overtness of $C$, and the ordinary rounded equations are also the $\J$-Kleisli equations, since the inclusion in ordinary ideals reflects equality.
\item \textbf{Existence and uniqueness.} Write $j_r:A_r\to\J B$ and $q_r:\J B\to A_r$ for the splitting.
An algebra map $f:A_r\to A_s$ determines
    \[
        m=j_s f q_r\eta_B:B\longrightarrow\J C.
    \]
    Its extension is $j_s f q_r$, because the extension of a map from generators is the unique algebra map with those values.
The equations for the splittings give both rounded equations.
\item \textbf{Converse construction.} In the other direction a rounded $m$ determines $q_s m^\#j_r$; substitution of $j_rq_r=r^\#$ and $j_sq_s=s^\#$ shows these two constructions are inverse.
Now apply \cref{thm:relations} to the same positive matrix.
Composition is $\star$ throughout, and equality of maps is equality of these matrices.
This proves fullness and faithfulness on actual morphisms, not just agreement of the objects.
\end{enumerate}
\end{proof}

This theorem identifies the searchable presentation category with the full subcategory of $\ED(\E)$ on its generated objects, up to the chosen presentations. It also proves agreement between observation preservation and the matrix definition of an effective map. The matrix composition uses semidecidable approximation and search over the middle presentation; it does not require semidecidable inclusion of completed ideals.

For the algebraic specialization take $r=y_B$. Searchability then means that $B$ is overt and its order is $\Sigma$-classified. The completion is $\Idl(B)$, effective points are semidecidable ideals, and the effective observations are semidecidable upper traces. General effective algebraic objects, as described by \cref{thm:algebraic}, need not come with this particular kind of presentation.

\subsection{Products of searchable presentations}\label{app:products}
Let $(B,r)$ and $(C,s)$ be searchable presentations. On the product preorder put
\[
 (\rho\times\sigma)((b,c),(b',c'))
 =\rho(b,b')\wedge\sigma(c,c').
\]
The carrier is overt. Each row is a product of ideals, hence an ideal, and interpolation follows by pairing the two intermediate witnesses. This defines a searchable presentation.

\begin{proposition}\label{prop:productpresentation}
Its effective dcpo is canonically the categorical product of $(D_r,\OO_r)$ and $(D_s,\OO_s)$.
\end{proposition}
\begin{proof}
\leavevmode
\begin{enumerate}[label=\textup{(\roman*)},leftmargin=*,itemsep=3pt,topsep=4pt]
\item \textbf{Ideal and directedness calculation.} An ideal $K$ in $B\times C$ equals the product of its two lower projections.
One inclusion is immediate.
For the other, witnesses for $b$ and $c$ lie in two pairs of $K$; directedness gives a pair above both, and lower closure gives $(b,c)\in K$.
Inhabitedness ensures that the two projections recover the original factors.
The rounding equation is then equivalent to the two separate rounding equations.
This gives $D_{r\times s}\cong D_r\times D_s$ with coordinatewise directed suprema.
\item \textbf{Coherence equations.} Each coordinate observation of the product presentation is the conjunction of a coordinate observation of each factor.
For a general effective trace $S:B\times C\to\Sigma$, its associated open is
\[
 (I,J)\longmapsto\exists(b,c)\,[I(b)\wedge J(c)\wedge S(b,c)].
\]
This is an overt union of rectangle observations with $\Sigma$-constant weights.
\item \textbf{Converse construction.} Conversely the generating rectangle observations have $\Sigma$-valued traces, and the trace-defined observation family is closed under the prescribed operations.
Thus its observation structure is exactly the one of \cref{prop:finitelimits}.
\end{enumerate}
\end{proof}


\begin{remark}[Exact bases]
\label{rem:exact-bases}
    One always has $I(b)\Rightarrow r(b)\way I$. The converse need not hold for a redundant presentation. If the presentation is exact in the sense
    \[
        \rho(a,b)\iff r(a)\way r(b),
    \]
    then $I(b)\iff r(b)\way I$. Indeed, applying $r(b)\way I$ to \eqref{eq:roundedreconstruct} gives $r(b)\subseteq r(c)$ for some $c\in I$. Round $c$ once more to $d\in I$ with $\rho(c,d)$. Then $r(b)\way r(d)$, so exactness gives $\rho(b,d)$ and lower closure gives $b\in I$. Under exact presentations the matrix $m_F(b,c)$ can therefore be written $s(c)\way F(r(b))$.
\end{remark}


\subsection{Separation and determination by basic points}
\label{subsec:separation-and-determination-by-basic-points}

The specified observations of a searchable completion recover its order, even though that order is generally not itself semidecidable.

\begin{proposition}
\label{prop:observationseparation}
    For $I,J\in D_r$,
    \[
        I\subseteq J\quad\Longleftrightarrow\quad
        \forall U\in\OO_r\,[U(I)\Rightarrow U(J)].
    \]
    Moreover the effective-point functor is faithful on the full subcategory of searchable completions. Every effective map in this subcategory is determined both by its values on the basic family and by its inverse-image map on observations.
\end{proposition}
\begin{proof}
\leavevmode
\begin{enumerate}[label=\textup{(\roman*)},leftmargin=*,itemsep=3pt,topsep=4pt]
\item \textbf{Order reflection by coordinate tests.} Every specified observation is upper, which proves the forward implication.
Conversely take the coordinate observations $\delta_a$.
The assumed implication gives $I(a)\Rightarrow J(a)$ for every $a$, which is inclusion.
This argument is internal and uniform in $I,J$.
\item \textbf{Ideal and directedness calculation.} If $F,G:D_r\to D_s$ agree on effective points, compose that equality with $B\to\Pt(D_r)$ given by the basic family.
The resulting maps into $D_s$ agree at every $r(b)$.
Directed reconstruction then gives
    \[
        F(I)=\bigvee_{b\in I}F(r(b))
        =\bigvee_{b\in I}G(r(b))=G(I).
    \]
    If instead $F^*=G^*$ on observations, pull back every coordinate observation $\delta_c$ of the target.
One obtains $F(I)(c)=G(I)(c)$ for every $I,c$, hence equality of the maps.
The same reconstruction formula proves determination by the basic family directly.
\end{enumerate}
\end{proof}

The universal quantifier in the order-recovery formula remains $\Omega$-valued. The proposition therefore reconciles observational determination of an order with failure of a semidecision for that order. Faithfulness of $\Pt$ supplies uniqueness of extensions from effective points. The intrinsic extension theorem, \cref{thm:chu}, will characterize existence by an observation-transpose square. In the numerical model, the finite-information theorem makes this transpose available for every internal map of program objects, giving the Myhill--Shepherdson comparison.


An effective Scott-continuous map is now represented by a positive rounded distributor, and its composite is calculated by existential convolution over the intermediate overt basis. The spectrum construction will show that these matrices can also be recovered directly from the positive observation traces.

\section{Searchable bases and two positive spectra}
\label{sec:searchable-bases-and-two-positive-spectra}

The rounded trace object packages positive observations independently of any completed carrier. We test how much of the completed domain it determines by considering its $\Omega$-valued and $\Sigma$-valued characters. This separates reconstruction of all rounded ideals from reconstruction of their effective points.


The searchable presentations identify the admitted weights and determine a useful positive reconstruction.

\begin{proposition}
\label{jd:prop:basis}
    For an overt poset $B$ whose order is $\Sigma$-classified,
    \[
        \J B\cong\Idl(B)\times_{\Omega^B}\Sigma^B.
    \]
    In other words, its admitted ideals are exactly the semidecidable ideals.
\end{proposition}
\begin{proof}
\leavevmode
\begin{enumerate}[label=\textup{(\roman*)},leftmargin=*,itemsep=3pt,topsep=4pt]
\item \textbf{Positive membership of admitted ideals.} If $I\in\J B$, the principal upper predicate $b\leq-$ is $\Sigma$-valued and
    \[
        I(b)=\chi_I(b\leq-).
    \]
    This gives the membership classifier uniformly in $b$.
Conversely, if membership is $\Sigma$-valued, then $I(p)\wedge u(p)$ is a positive predicate on overt $B$, so its existential quantification is in $\Sigma$.
\end{enumerate}
\end{proof}

Combining \cref{jd:prop:compact,jd:prop:basis} gives a basis-independent criterion within the algebraic part: the algebras admitting an overt, positively ordered algebraic basis are exactly those whose \textbf{canonical compact part} $K(A)$ is overt and has $\Sigma$-classified order. Their presentation is the canonical map $\J K(A)\cong A$. Redundant preorder bases give the same result after their separated quotient; overtness and the classified order descend along its regular epimorphism. This is a useful concrete strengthening of the adequacy discussion.

\subsection{Rounded approximation without semidecidable order}
\label{subsec:rounded-approximation-without-semidecidable-order}

Now let $(B,r)$ be a searchable approximation presentation of \cref{dom:sec:rounded}: $B$ is overt, each $r(b)$ is an ideal, $r\leq\eta_B$, $r\star r=r$, and
\[
    \rho(a,b)\Longleftrightarrow a\in r(b)
\]
is $\Sigma$-classified. The order of $B$ can remain $\Omega$-valued. By overtness, every row $r(b)$ belongs to $\J B$, so $r:B\to\J B$ is an idempotent in $\Kl(\J)$ from the outset.

Define the positive trace object first:
\[
    \Gamma_r=\{u:B\to\Sigma\mid
    u(b)=\exists a:B\,[u(a)\wedge\rho(a,b)]\}.
\]
Its elements are upper predicates. It is closed under $\Sigma$-constants, finite meets, and overt-indexed joins. Directedness of the rows proves the finite-meet assertion; the remaining operations follow from row inhabitation and interchange of existential quantifiers. Put $\delta_a(b)=\rho(a,b)$.

For $V=\Sigma$ or $V=\Omega$, let $\Spec_V(\Gamma_r)$ be the object of monotone maps $\Gamma_r\to V$ preserving finite meets, overt-indexed joins, and $\Sigma$-constants, interpreted through $i$ for $V=\Omega$.

\begin{theorem}[Two-spectrum reconstruction]
\label{jd:thm:spectra}
    There are canonical isomorphisms
    \[
        \Spec_\Omega(\Gamma_r)\cong D_r,
        \qquad
        \Spec_\Sigma(\Gamma_r)\cong\Pt D_r,
    \]
    where $D_r$ is the full rounded-ideal completion. They identify postcomposition with $i$ with the effective-point inclusion:
    \[
        \begin{tikzcd}[column sep=large]
            \Spec_\Sigma(\Gamma_r)\ar[r,"\cong"]\ar[d,hook,"i\circ-"']
            &\Pt D_r\ar[d,hook]\\
            \Spec_\Omega(\Gamma_r)\ar[r,"\cong"']&D_r.
        \end{tikzcd}
    \]
    Both spectra are constructed from the positive trace object. No observation family is chosen after completing $B$.
\end{theorem}
\begin{proof}
The forward maps evaluate positive traces on rounded ideals. For the inverse, evaluations of the coordinate traces determine membership; the trace identities then prove ideality, rounding, and recovery of the original character.

\leavevmode
\begin{enumerate}[label=\textup{(\roman*)},leftmargin=*,itemsep=3pt,topsep=4pt]
\item \textbf{Two-spectrum reconstruction.} A rounded ideal $I$ gives
    \[
        \chi_I(u)=\exists b:B\,[I(b)\wedge u(b)].
    \]
    Ideality gives preservation of top and finite meets, and existential interchange gives preservation of overt joins.
Inhabitedness gives preservation of constants.
If $I$ has $\Sigma$-valued membership, overtness of $B$ makes $\chi_I$ $\Sigma$-valued.
\item \textbf{Converse construction.} Conversely, for a character $\chi$ define $I_\chi(a)=\chi(\delta_a)$.
This is a lower predicate because each row of $r$ is lower and $\chi$ is monotone.
The identities in $\Gamma_r$
    \begin{align}
        \top&=\bigvee_{b:B}\delta_b,\label{jd:eq:gen1}\\
        \delta_a&=\bigvee_{c:B}\rho(a,c)\wedge\delta_c,\label{jd:eq:gen2}\\
        \delta_a\wedge\delta_b
        &=\bigvee_{c:B}\rho(a,c)\wedge\rho(b,c)\wedge\delta_c\label{jd:eq:gen3}
    \end{align}
    show respectively that $I_\chi$ is inhabited, rounded, and directed.
\item \textbf{Ideal and directedness calculation.} For \eqref{jd:eq:gen3}, interpolate $\rho(a,x)$ and $\rho(b,x)$ once, then take a common upper member in the ideal $r(x)$; monotonicity gives the two relations to the intermediate $c$.
The reverse implication uses transitivity.
Thus $I_\chi\in D_r$.
In the $\Sigma$ case its membership is $\Sigma$-valued, giving an effective point.
\item \textbf{Recovery from positive traces.} Finally every trace has the positive expansion
    \[
        u=\bigvee_{b:B}u(b)\wedge\delta_b.
    \]
    Preservation of these operations yields $\chi_{I_\chi}=\chi$; rounding gives $I_{\chi_I}=I$.
These equations prove both isomorphisms and their compatibility with $i$.
\end{enumerate}
\end{proof}

\begin{corollary}[Positive observation duality]
\label{jd:cor:duality}
    For searchable presentations, effective Scott-continuous maps $D_r\to D_s$ are in bijection with maps $\Gamma_s\to\Gamma_r$ preserving finite meets, overt joins, and $\Sigma$-constants. This is contravariantly functorial.
\end{corollary}
\begin{proof}
\leavevmode
\begin{enumerate}[label=\textup{(\roman*)},leftmargin=*,itemsep=3pt,topsep=4pt]
\item \textbf{Scott continuity.} An effective Scott-continuous map acts by inverse image on observations.
Conversely let $h:\Gamma_s\to\Gamma_r$ preserve the stated operations.
Set
    \[
        m(b,c)=h(\delta_c)(b).
    \]
    Evaluation at $b$ composed with $h$ is a $\Sigma$-character, so \cref{jd:thm:spectra} makes each row a semidecidable rounded ideal.
Membership of $h(\delta_c)$ in $\Gamma_r$ gives rounding in the source; the character calculation gives rounding in the target.
Hence $m=s\star m\star r$.
Its ideal extension gives an effective Scott-continuous map.
\item \textbf{Trace calculation.} The trace expansions show that its inverse-image operation is $h$, and coordinate observations prove uniqueness.
Composition is either composition of inverse images or the corresponding overt relational convolution.
\end{enumerate}
\end{proof}


The two spectra recover the full rounded completion and its effective-point subobject from the same positive traces. Their difference is the codomain of evaluation, not a change in the generating tests. This distinction will persist when a partial computation is lifted or localized.

\section{Domain lifting and effective partiality}
\label{sec:domainlift}

Domain completion and partiality must agree about what it means for a value to be defined. The semantic lift allows arbitrary truth-valued extents; its effective points should recover exactly the $\Sigma$-classified partial elements. The comparison below proves this and gives the lifted observation formula.


The domain lift is defined for every effective dcpo. On its underlying object use all partial extents:
\[
    D_\bot=\sum_{p:\Omega}D^p.
\]
Order partial elements by extension and increase of value:
\[
    (p,x)\leq(q,y)\iff (p\Rightarrow q)\wedge
    (\forall a:p\;x(a)\leq y(a)).
\]
The least element has false extent and the unit $\eta_D(x)$ has true extent.

\begin{lemma}
\label{lem:ordinarylift}
    The object $D_\bot$ is a dcpo. It is the free pointed dcpo on $D$ for strict Scott-continuous maps out of it.
\end{lemma}
\begin{proof}
\leavevmode
\begin{enumerate}[label=\textup{(\roman*)},leftmargin=*,itemsep=3pt,topsep=4pt]
\item \textbf{Ideal and directedness calculation.} For an ideal $J$ of partial elements, take extent
    \[
        p_J\iff\exists z\in J\;\dom(z).
    \]
    Under this extent, the values of defined members of $J$ form an inhabited directed family in $D$: a common upper member of $J$ supplies a common upper value.
Its supremum defines the value of the required partial element.
Testing upper bounds proves that it is $\sup J$.
\item \textbf{Scott continuity.} For a Scott-continuous $f:D\to E$ into a pointed dcpo, define
    \[
        f_\bot(p,x)=\sup\bigl(\{\bot_E\}\cup\{f(x(a))\mid a:p\}\bigr).
    \]
    The family is inhabited and directed.
The formula preserves directed suprema by interchange of the two directed suprema and is strict.
Every $(p,x)$ is the directed supremum of the bottom partial element and the elements $\eta_D(x(a))$ for $a:p$, proving uniqueness.
\end{enumerate}
\end{proof}

For $U\in\OO_D$ define its positive extension by
\[
    U^+(p,x)\iff\exists a:p\;U(x(a)).
\]
It is Scott open by the supremum description above. Equip $D_\bot$ with the observation structure generated by these opens; $\Sigma$-constants are included by the closure operation. The open $\top^+$ is definedness.

\begin{theorem}
\label{thm:effectivelift}
    Domain lifting defines a monad on $\ED(\E)$. There is a natural isomorphism
    \begin{equation}
    \label{eq:liftcomparison}
        \Pt(D_\bot)\cong \LL\Pt(D)
    \end{equation}
    compatible with units and multiplication. For a searchable presentation $(B,r)$, this lift agrees with the rounded completion on $1+B$ obtained by adding a least basis element.

    Writing $\lambda_D$ for the isomorphism, naturality is the square
    \[
        \begin{tikzcd}[column sep=large]
            \Pt(D_\bot)\ar[r,"{\lambda_D}","{\cong}"']\ar[d,"{\Pt(F_\bot)}"']&
            \LL\Pt D\ar[d,"{\LL\Pt F}"]\\
            \Pt(E_\bot)\ar[r,"{\lambda_E}","{\cong}"']&\LL\Pt E.
        \end{tikzcd}
    \]
\end{theorem}
\begin{proof}
\leavevmode
\begin{enumerate}[label=\textup{(\roman*)},leftmargin=*,itemsep=3pt,topsep=4pt]
\item \textbf{Lifted monad on observations.} The ordinary free-pointed construction supplies the monad.
On observations,
    \[
        (f_\bot)^*(U^+)=(f^*U)^+,
        \qquad\eta_D^*(U^+)=U,
        \qquad\mu_D^*(U^+)=(U^+)^+.
    \]
    The formulas for constants are immediate.
The closure lemma then proves that the functor, unit, and multiplication are effective.
\item \textbf{Effective-point isomorphism.} If $(p,x)$ is an effective point of $D_\bot$, evaluating definedness shows $p\in\Sigma$.
Under $p$, evaluation of $U^+$ is evaluation of $U$ at $x$, uniformly in $U$.
Hence $x:p\to\Pt(D)$, giving a point of $\LL\Pt(D)$.
\item \textbf{A lifted searchable basis.} Conversely, from $p\in\Sigma$ and $x:p\to\Pt(D)$, each $U^+(p,x)$ lies in $\Sigma$ by dominance.
The observations taking a $\Sigma$-value at this partial element form an observation structure, so the statement extends from the generators to all of $\OO_{D_\bot}$.
These constructions are inverse.
They preserve the unit and multiplication because both sides use the same dependent extents and values.
\item \textbf{Compatibility of observations.} For the last assertion put $\widetilde B=1+B$, with new least element $*$, and define
    \[
        \widetilde\rho(*,z)=\top,
        \quad\widetilde\rho(\bar a,*)=\bot,
        \quad\widetilde\rho(\bar a,\bar b)=\rho(a,b).
    \]
    Every row is an ideal; interpolation follows by cases and \eqref{eq:rho}.
The carrier is overt and the relation is $\Sigma$-classified. A rounded ideal $J$ determines extent $p=\exists b\,J(\bar b)$ and, under $p$, value $\{b\mid J(\bar b)\}\in D_r$.
\item \textbf{Converse construction.} Conversely a partial ideal gives $J(*)=\top$ and $J(\bar b)=\exists a:p\,x(a)(b)$.
These are inverse order maps and preserve directed unions.
\item \textbf{Ideal and directedness calculation.} To compare observations, every Scott open $W$ on $D_\bot$ satisfies
    \[
        W(p,x)\iff W(\bot)\vee\exists a:p\,W(\eta_Dx(a)).
    \]
    This follows from the directed decomposition used in \cref{lem:ordinarylift}. A $\Sigma$-valued trace on $\widetilde B$ makes $W(\bot)$ a $\Sigma$-constant and makes $W\eta_D$ an observation on $D_r$.
Thus $W$ lies in the generated observation structure.
Conversely all its generators have $\Sigma$-valued traces.
The structures agree.
\end{enumerate}
\end{proof}

Consequently an effective map $D\to E_\bot$ induces a $\Sigma$-partial map $\Pt(D)\pmap\Pt(E)$, functorially for Kleisli composition. The comparison retains both the extent and the value. It does not replace a partial extent by a Boolean decision.

\begin{proposition}[Strength of domain lifting]
\label{prop:domainstrength}
    Domain lifting is a commutative strong monad on $\ED(\E)$ with respect to its finite products. Under the finite-product comparison for $\Pt$, the isomorphism $\lambda_D$ preserves strength.
\end{proposition}
\begin{proof}
\leavevmode
\begin{enumerate}[label=\textup{(\roman*)},leftmargin=*,itemsep=3pt,topsep=4pt]
\item \textbf{Strength of domain lifting.} On underlying partial elements, strength is
    \[
        \operatorname{st}_{D,E}(d,(p,x))=(p,\lambda a.(d,x(a))).
    \]
    It is Scott-continuous.
For a directed family, its extent at the supremum is the union of the component extents.
Under this extent, the defined members are cofinal among the relevant coordinates: once one defined member has been found, directedness bounds it together with any other member.
\item \textbf{Continuity of the strength map.} Consequently the supremum of the defined $D$-coordinates is the supremum of all $D$-coordinates, and the $E$-coordinate is the supremum used in the definition of the lift.
These are exactly the coordinates obtained by applying strength to the componentwise supremum.
\item \textbf{Effectivity of strength.} To prove effectivity, recall that product observations are generated by the two projection families.
The inverse images of their positive extensions are
    \begin{align*}
        \operatorname{st}^*((\pi_D^*U)^+)(d,z)
        &=U(d)\wedge\dom(z),\\
        \operatorname{st}^*((\pi_E^*V)^+)(d,z)
        &=V^+(z).
    \end{align*}
    Both are source observations.
Positive extension preserves binary meets and overt unions, while a $\Sigma$-constant extends to its conjunction with definedness.
\item \textbf{Observation closure for strength.} Thus preservation extends first over the generated product observations and then over the generated lift observations by \cref{lem:obshull}.
\item \textbf{Coherence equations.} The strength and monad equations hold on extents and values: keeping a parameter commutes with substituting a true extent and with flattening dependent extents.
The two induced maps $D_\bot\times E_\bot\to(D\times E)_\bot$ both have extent $p\wedge q$ and value $(x,y)$ under that extent.
They are equal, proving commutativity.
\item \textbf{Compatibility with classified strength.} Finally $\lambda$ retains these same extents and values, now with their $\Sigma$-classification at effective points.
The formula for strength therefore agrees with that of $\LL$ after the product comparison for $\Pt$.
\end{enumerate}
\end{proof}


\subsection{The monad comparison on partial computations}
\label{subsec:the-monad-comparison-on-partial-computations}

Let $\lambda_D:\Pt(D_\bot)\cong \LL\Pt(D)$ denote \eqref{eq:liftcomparison}. If $f:D\to E_\bot$ is effective, its interpretation on points is the partial map classified by
\[
    \Pt(D)\xrightarrow{\Pt(f)}\Pt(E_\bot)
    \xrightarrow{\lambda_E}\LL\Pt(E).
\]
For $g:E\to F_\bot$, the interpreted composite has extent ``$f$ is defined and its value is in the domain of $g$''. By the multiplication compatibility in \cref{thm:effectivelift}, this is also the interpretation of $\mu_F\,g_\bot f$. The comparison therefore defines a functor from the Kleisli category of domain lifting to $\Par_\Sigma(\E)$. It preserves restriction to definedness whenever that restriction is represented in the source, because the comparison includes the extent as well as the value.

For a searchable basis this statement can be checked at coordinates. A point of the lifted completion is a rounded ideal $J$ on $1+B$. It contains $*$, and its definedness is
\[
    \exists_B^\Sigma b\,J(\bar b).
\]
An observation of a defined value with rounded trace $S$ has the form $\exists_B^\Sigma b\,[J(\bar b)\wedge S(b)]$. Thus a semidecidable membership predicate for $J$ supplies semidecisions for both definedness and every observation made after definition. Conversely a semidecision of an extent, together with such a predicate under the extent, supplies all the lifted coordinates by dependent conjunction. These are exactly the data that the external partial-element presentation in Part~\ref{part:relative} will realize by guarded streams and hit witnesses.


Domain lifting has now been compared at the level of both extents and values. To lift partiality through the admitted-ideal algebra structure itself, we next construct the distributive law that orders those two operations.

\section{Partiality as a distributive law}
\label{sec:partiality-as-a-distributive law}

The effective-point comparison identifies individual partial elements, but composition with admitted suprema requires a stronger compatibility. A distributive law makes this interaction explicit: the extent of an admitted family of partial values is itself positive, and the defined values form an admitted ideal. This gives partiality a monad structure on computational algebras.


Equip $\LL P=\sum_{p:\Sigma}P^p$ with the extension order on partial elements. It is a monad on internal posets as well as on the underlying topos. Its extent map $\defin:\LL P\to\Sigma$ is monotone.

\begin{theorem}
\label{jd:thm:distributive}
    There is a distributive law
    \[
        \lambda:\J\LL\Longrightarrow\LL\J.
    \]
    For $I\in\J(\LL P)$ its extent is
    \[
        p_I=\chi_I(\defin)\in\Sigma.
    \]
    Under $p_I$, its value is the ideal generated by the values of the defined members of $I$.
\end{theorem}
\begin{proof}
The construction is determined by two components: the positive extent of a family of partial values and the ideal of values obtained under that extent. We establish admissibility of both components before checking the distributive-law equations.

\leavevmode
\begin{enumerate}[label=\textup{(\roman*)},leftmargin=*,itemsep=3pt,topsep=4pt]
\item \textbf{Defined-value ideal.} Under $p_I$ there is a defined member, so the proposed value ideal $V_I$ is inhabited.
If two defined members supply values, directedness of $I$ supplies a common upper partial element; it is defined and its value bounds both.
Taking lower closure gives an ideal of $P$.
\item \textbf{Admissibility of its positive pairing.} For $u\in U_\Sigma P$, extend it to the positive predicate
    \[
        \widetilde u(z)=\exists d:\defin(z)\;u(z(d))\quad\text{on }\LL P.
    \]
    Dominance makes it $\Sigma$-valued, and it is monotone.
Under $p_I$,
    \[
        \chi_{V_I}(u)=\chi_I(\widetilde u)\in\Sigma.
    \]
    Thus the value lies in $\J P$, and $\lambda$ is defined uniformly.
Lower image commutes with taking the extent and the generated value ideal, proving naturality.
\item \textbf{Naturality.} A partial element is determined by its classified extent and its value over that extent, so checking those two components suffices for equality.
The two unit equations have sources $\LL P$ and $\J P$ respectively; the multiplication equations have sources $\J\J\LL P$ and $\J\LL\LL P$.
Their common targets are $\LL\J P$.
The four distributive-law identities can be verified by these same two components.
For a principal ideal of partial elements, the extent and value are the original extent and the principal ideal of its value.
\item \textbf{Unit laws.} For an ideal of total elements, the extent is true and the value ideal is the original ideal.
These are the two unit laws.
\item \textbf{Multiplication coherence.} For multiplication of $\J$, both routes have extent
    \[
        \exists I\in\mathcal I\;\exists z\in I\;\defin(z)
    \]
    and, under that extent, the ideal generated by the defined values of these $z$.
Flattening the intermediate ideals adds only lower closures, which are already present.
\item \textbf{Multiplication law.} For multiplication of $\LL$, both routes have extent asserting that an outer defined member has a defined inner value; their values are the ideals generated by the resulting inner values.
This is reassociation of dependent existential quantification.
Equality of extents and value ideals proves the two multiplication laws.
\item Hence
    \begin{align*}
        \lambda\eta_{\LL}^{\J}&=\LL\eta^{\J},&
        \lambda\J\eta^{\LL}&=\eta^{\LL}_{\J},\\
        \lambda\mu^{\J}_{\LL}&=\LL\mu^{\J}\;\lambda_{\J}\;\J\lambda,&
        \lambda\J\mu^{\LL}&=\mu^{\LL}_{\J}\;\LL\lambda\;\lambda_{\LL}.
    \end{align*}

The defining maps satisfy the following commutative diagram
\[
\begin{tikzcd}[column sep=large,row sep=large]
\J P\ar[r,"\J\eta^{\LL}_P"]\ar[dr,"\eta^{\LL}_{\J P}"']&\J\LL P\ar[d,"\lambda_P"]\\ &\LL\J P
\end{tikzcd}
\]
\end{enumerate}
\end{proof}

\begin{corollary}
\label{jd:cor:lift}
    Partiality lifts to a monad on $\AS$. Its algebra map at $(A,a)$ is
    \[
        \J\LL A\xrightarrow{\lambda_A}\LL\J A
                      \xrightarrow{\LL a}\LL A.
    \]
    Under the semantic comparison $\Pt(D_\bot)\cong\LL\Pt D$, this is the algebra structure induced by semantic domain lifting.
\end{corollary}
\begin{proof}
\leavevmode
\begin{enumerate}[label=\textup{(\roman*)},leftmargin=*,itemsep=3pt,topsep=4pt]
\item \textbf{Ideal and directedness calculation.} The lifted algebra and its monad laws follow from the four equations above.
For the comparison, an admitted ideal of effective partial points has extent $\chi_I(\defin)$; the semantic supremum has this same extent and, when defined, the supremum of the same value ideal.
Thus the existing partiality isomorphism is also an isomorphism of $\J$-algebras.
Parameter strength agrees because retaining a total parameter changes neither calculation.
\end{enumerate}
\end{proof}

\paragraph{Iteration.}
In $\Eff$, where \cref{thm:relativefoundation} proves $\N$ overt, the lower ideal generated by an increasing sequence is $\Sigma$-admissible: pairing with $u$ is $\exists n:\N\,u(x_n)$. Thus pointed algebras admit the least-fixed-point construction. Cartesian closure and uniformity place parameterized iteration in the same category. \Cref{sec:intrinsicrecursion} proves that this operation equips the lifted partiality monad with complete Elgot iteration and that finite completion of its guarded dynamics yields coproduct feedback. The general Conway, Elgot, and trace formalisms are established theory; the comparison with the selected doctrine is made explicitly here.

\part{Realizability characterization in the effective topos}
\label{part:relative}

The preceding domain constructions are internal and do not assign codes to arbitrary completed points. To recover numerical meaning, we must represent the carriers, relation witnesses, searches, and maps without identifying a code with its denotation. Projective name covers and their kernel pairs provide that distinction; the resulting presentation theorems explain exactly which effective structures can be implemented.


We begin with modest objects, where PER names are available, and then pass to arbitrary carriers through exact-completion equality covers. The resulting tracked presentation calculus will be used for the intensional operations in Part~\ref{part:coding}.

\section{Modest objects and the names of completed points}
\label{sec:modestcompletion}

The positive truth object and searchable completions were defined internally. To assign numerical representations to their effective points, we first show that their carriers are modest and identify the corresponding PERs. The support covers then provide names on which later program and proof transformations can operate.


The modestness proofs use the closure properties of \cref{lem:modestclosure}; the resulting support covers are the code covers used in later implementation arguments.


\begin{lemma}[Finite-hit name presentations]
\label{lem:sigmamodest}
    In $\Eff$, $\Sigma$ is modest: its names are total binary-stream indices modulo equality of hit support.
\end{lemma}
\begin{proof}
\leavevmode
\begin{enumerate}[label=\textup{(\roman*)},leftmargin=*,itemsep=3pt,topsep=4pt]
\item \textbf{Finite-hit name presentations.} The numerical Boolean and numeral presentations in $K_1$ are modest.
Consequently the objects $2$ and $\N$ have their modest presentations, so \cref{lem:modestclosure} supplies a PER presentation of $2^{\N}$.
Let $P$ be its support cover.
On its valid stream names put the coarser PER identifying two streams when either both have a hit or neither has a hit; equivalently their predicates $\exists n\,[\theta(n)=1]$ have the same support.
This is an external definition of a PER, not a procedure deciding whether a stream hits.
Call its quotient $S$.
\item \textbf{Image and quotient construction.} The classifier $e:2^{\N}\to\Sigma$ descends through $S$.
Indeed, on the partitioned assembly of pairs of stream names with equal support, first-hit search gives maps between the two hit assemblies in both directions.
Their regular images, and hence their classifiers, are equal by \cref{lem:support}.
\item \textbf{Converse construction.} Conversely, equality of the two classifiers gives equality of these supports by the same lemma.
Thus the kernel relation on the stream-name cover induced by $S$ is exactly the kernel relation induced by $\Sigma$.
Both quotient maps are regular epic: for $S$ this is its support presentation, and for $\Sigma$ it is the composite of the stream-name cover and $e$.
Effectiveness of equivalence relations identifies the quotients.
Hence $S\cong\Sigma$.
\end{enumerate}
\end{proof}

\begin{theorem}[Modesty and covers of effective rounded ideals]
\label{thm:modestpoints}
    If $(P,r)$ is a searchable modest approximation presentation, then $\Pt(D_r)$ is modest. Its canonical PER can be described by stream names for membership, with evidence for ideality and rounding; equivalence is equality of the resulting membership predicates. Its support cover is regular epic, uniformly after passage to a projective cover of parameters.
\end{theorem}
\begin{proof}
\leavevmode
\begin{enumerate}[label=\textup{(\roman*)},leftmargin=*,itemsep=3pt,topsep=4pt]
\item \textbf{Modesty and covers of effective rounded ideals.} By \cref{lem:sigmamodest,lem:modestclosure}, $\Sigma^P$ is modest.
The identification
    \[
        \Pt(D_r)\cong D_r\times_{\Omega^P}\Sigma^P
    \]
    of \cref{thm:compatibility} exhibits the point object as the subobject of $\Sigma^P$ satisfying inhabitedness, lower closure, directedness, and rounding.
It is therefore modest by the subobject argument above.
\item \textbf{Tracked implementation.} Explicitly, for its resulting modest assembly $X$, define
    \[
      n\mathrel{R_X}m\iff
      \exists x:X\,[n\Vdash_Xx\wedge m\Vdash_Xx],\qquad
      D_X=\{n:R_X(n,n)\}.
    \]
    Modesty makes $R_X$ a PER.
Each $n\in D_X$ realizes a unique
    element, so unique-value comprehension gives the regular-epimorphic
    support cover $D_X\onto X$ with kernel pair $R_X$.
\item \textbf{Ideal and directedness calculation.} The realizing
    relation here is that of the constructed subobject assembly: it
    retains evidence for ideality and rounding in its subobject assembly.
\item \textbf{Observation preservation.} To calculate that cover, pull a family of these ideals back to a partitioned cover of its parameter object.
Lift its membership map, on the product with the basis-name cover, through the total-stream cover of $\Sigma$.
Combinatory abstraction packages the resulting tracker as a stream-family name.
Pull the proof objects for the four ideal conditions to the same projective context and lift their existential conclusions.
These lifts supply names of members, common upper bounds, and rounding witnesses; implications supply the corresponding witness transformations.
\item \textbf{Ideal and directedness calculation.} Equality transport along the basis PER is provided by first-hit search.
Their finite tuple is the asserted name of a rounded ideal.
\item \textbf{Converse construction.} Conversely a stream-family name with those witnesses gives a $\Sigma$-classified membership predicate satisfying exactly those four conditions, hence a point of the pullback.
Two such names denote the same point precisely when their membership predicates agree, since the pullback is a subobject of $\Sigma^P$.
The forward construction is available locally for every family of points, proving regular epicity in every parameter context.
Names and their uniform transformations are numerical in $K_1$.
Validity of a complete name is a mathematical support condition and need not be decidable.

The interpretation maps form the commuting diagram
\[
\begin{tikzcd}[column sep=large,row sep=large]
D_X\ar[r,two heads,"q_X"]\ar[dr,"\iota q_X"']&\Pt D_r\ar[d,hook,"\iota"]\\ &\Sigma^P
\end{tikzcd}
\]
\end{enumerate}
\end{proof}


Modesty identifies a valid name with at most one effective completed point, while the support cover supplies local names for every such point. Proof transformations must still respect PER equality; the next section extracts precisely those transformations from projective proof covers.

\section{Projective extraction of witnesses}
\label{sec:extraction}

An inclusion of interpreted predicates does not by itself display the program that converts one proof witness into another. Projective covers supply that conversion after choosing name presentations. We make the transport of equality evidence explicit so that the resulting witnesses remain extensional on quotient carriers.


The constructions in this section present semantic relations and their image inclusions. They supply witness extraction for the realizability comparison. The finite execution certificates used later are constructed from actual runs in \cref{thm:operational-certificates}.

The repeated realization arguments can be organized into a single lifting principle. Let $q:P\onto X$ be a selected partitioned cover, and let $R\incl X$ be any subobject. Choose a partitioned cover $P_R\onto P\times_X R$. The composite $P_R\to P$ is a \textbf{proof-name presentation} of $R$. It retains membership witnesses alongside the external support.

For a presented relation on a product of modest carriers, its proof names must transport along equality of carrier denotations. Such transport is extracted by projectivity: change the tuple of carrier names along its PER equality, map to the same internal relation, and lift back through the proof-name cover. Reflexivity and composition are respected after interpretation; proof-name equality is imposed only after interpretation.

\subsection{Extensional membership-witness covers}
\label{subsec:extensional proof-name covers}

Write $P$ for a product of selected PER covers, $X$ for the corresponding product of quotients, and $q:P\onto X$ for its interpretation. Its equality-name object is the kernel pair $E_P\rightrightarrows P$, whose names are tuples of PER-equivalent pairs. For a partitioned assembly $P_R$ and endpoint map $\partial:P_R\to P$, let $R\incl X$ be the regular image of $q\partial$. The canonical comparison is
\[
    \bar\partial:P_R\longrightarrow P\times_XR.
\]
There is a distinction between covering $R$ and covering its pullback at every selected carrier name. The following criterion supplies exactly the additional property.

\begin{proposition}[Extensional cover criterion]
\label{prop:extensionalcover}
The comparison $\bar\partial$ is regular epic if and only if there is a tracked equality-transport operation
\[
    \operatorname{tr}_R:
    P_R\times_{\partial,P,d_0}E_P\longrightarrow P_R,
    \qquad \partial\operatorname{tr}_R=d_1\pi_2.
\]
\end{proposition}
\begin{proof}
\leavevmode
\begin{enumerate}[label=\textup{(\roman*)},leftmargin=*,itemsep=3pt,topsep=4pt]
\item \textbf{Partiality and definedness.} Put $T=P_R\times_{\partial,P,d_0}E_P$.
There is a map $k:T\to P\times_XR$ taking a proof over a name $a$, together with equality of $a$ to $a'$, to the same relation instance at $a'$.
This map is regular epic.
Indeed, membership of a named point in $R$ means locally that it is represented by some proof in $P_R$; equality of the represented endpoints is precisely a point of $E_P$.
Thus every element of the target is locally in the image of $k$.
\item \textbf{Product calculation.} If transport exists, $k=\bar\partial\operatorname{tr}_R$.
Since $k$ is regular epic, so is $\bar\partial$.
Conversely, pull $\bar\partial$ back along $k$.
The resulting map is regular epic onto $T$.
This object is a carrier subassembly of a finite product of partitioned assemblies, hence is partitioned and projective. A section of the pullback followed by its other projection gives the required transport, with the stated endpoint equation.
Fullness of the partitioned-assembly embedding supplies one tracker for that map.
\end{enumerate}
\end{proof}

Thus a selected cover of an internal relation automatically has equality transport, and a proof-name presentation with transport covers the full pullback of its interpreted relation. Transport need not be unique or satisfy equations on proof names. Its endpoint equation is sufficient, since proofs are identified only after their regular image is taken. This separation lets ordinary relations retain the computational evidence that a bare support predicate would discard.


\begin{lemma}[Witness extraction]
\label{lem:extraction}
    Let $R,S$ be interpreted relations with selected partitioned proof-name covers. An internal implication $R\Rightarrow S$ admits a tracked transformation of proof names. For $R(\vec x)\Rightarrow\exists y\,S(\vec x,y)$, the transformation also returns a name of $y$ and a proof name of $S$. Conversely such transformations validate the corresponding sequents.
\end{lemma}
\begin{proof}
\leavevmode
\begin{enumerate}[label=\textup{(\roman*)},leftmargin=*,itemsep=3pt,topsep=4pt]
\item \textbf{Product calculation.} Pull the cover of the conclusion back along the map from the projective proof object of the premise to the interpreted conclusion.
The resulting regular epimorphism onto the premise proof object splits.
Its section provides the tracked transformation.
For an existential conclusion, use the projection of the proof object of $S$ to the context; its regular image is the existential interpretation.
The same lift supplies both the witness and its proof.
\item \textbf{Converse construction.} Conversely, taking regular images of a proof transformation gives precisely the claimed inclusion of subobjects.

The interpretation maps form the commuting diagram
\[
\begin{tikzcd}[column sep=large,row sep=large]
P_\phi\ar[r,"\widehat{\phi\Rightarrow\psi}"]\ar[d,two heads]&P_\psi\ar[d,two heads]\\ Q_\phi\ar[r,hook]&Q_\psi
\end{tikzcd}
\]
\end{enumerate}
\end{proof}


\subsection{Witness presentations of positive relations}
\label{proof-objects-for-formulas}

Fix a finite context $\vec x$ with name object $P_{\vec x}$, quotient $X_{\vec x}$, and cover $q_{\vec x}$. Terms have tracked evaluators by induction on composition of the selected operation trackers. For a positive-geometric formula $\phi$, put
\[
    Q_\phi=P_{\vec x}\times_{X_{\vec x}}\llbracket\phi\rrbracket.
\]
We construct a partitioned proof-name object $P_\phi$ over $P_{\vec x}$. For equality, use the subset on which the two evaluated names are PER-equivalent. For an ordinary relation, pull its proof-name endpoint map back along the tuple of term evaluators. For a positive relation, take the hit assembly of the substituted stream. Truth uses the context itself and falsity the empty assembly. The compound clauses are
\begin{align*}
    P_{\phi\wedge\psi}&=P_\phi\times_{P_{\vec x}}P_\psi,&
    P_{\phi\vee\psi}&=P_\phi+P_\psi,\\
    P_{\exists y\,\phi(\vec x,y)}&=P_{\phi(\vec x,y)}
       \quad\text{with the $y$-endpoint forgotten}.&&
\end{align*}
In the last clause the proof name still contains the witness name; only its projection to the context forgets it. These constructions use finite products, tagged sums, and carrier subsets, so each proof object remains partitioned.

\begin{lemma}[Formula covers]
\label{lem:formulacovers}
    The canonical map $P_\phi\to Q_\phi$ is regular epic for every such formula. It admits equality transport in the context, and the construction respects substitution up to maps between covers of the same interpreted relation.
\end{lemma}
\begin{proof}
\leavevmode
\begin{enumerate}[label=\textup{(\roman*)},leftmargin=*,itemsep=3pt,topsep=4pt]
\item \textbf{Formula covers.} Proceed by induction.
Equality is the pullback of the appropriate diagonal, represented by the PER kernel pair.
For an ordinary atom, \cref{prop:extensionalcover} and pullback stability supply the cover.
For a positive atom, the image of the substituted hit assembly has classifier equal to the composite of the interpreted positive relation with the term map, so it is precisely $Q_\phi$.
\item \textbf{Proof covers for conjunction.} Conjunction follows from pullback stability of regular epimorphisms.
For disjunction, the image of the tagged sum is the union of the two images.
Truth and falsity are immediate.
For the existential case let $R=\llbracket\phi(\vec x,y)\rrbracket$ and let $R\onto E\incl X_{\vec x}$ be the image of its projection.
\item \textbf{Universal-property argument.} The relevant composite is
    \[
        P_\phi\onto
        (P_{\vec x}\times P_y)\times_{X_{\vec x}\times X_y}R
        \onto P_{\vec x}\times_{X_{\vec x}}R
        \onto P_{\vec x}\times_{X_{\vec x}}E.
    \]
    The first arrow is the inductive cover; the second is a pullback of the carrier cover $P_y\onto X_y$; the third is a pullback of $R\onto E$.
All are regular epic.
This proves the existential clause without removing its witness from the proof name.
\item \textbf{Equality transport after formula induction.} The cover criterion now supplies context equality transport.
Under substitution the interpreted subobject is its pullback.
The substituted proof objects and selected proof objects both cover this pullback; projectivity gives comparison maps into either selected cover after the indicated refinement.
Their interpretations agree, which is the required substitution compatibility.

The construction uses the regular-image factorization triangle
\[
\begin{tikzcd}[column sep=large,row sep=large]
P_\phi\ar[r,two heads]\ar[dr]&Q_\phi\ar[d,hook]\\ &P_{\vec x}
\end{tikzcd}
\]
\end{enumerate}
\end{proof}

\begin{corollary}[Lifting inclusions of presented relations]
\label{cor:sequents}
    A sequent $\phi\vdash_{\vec x}\psi$ is valid internally if and only if there is a tracked map $P_\phi\to P_\psi$ over $P_{\vec x}$. The existence of such a map is independent of all selected proof covers.
\end{corollary}
\begin{proof}
\leavevmode
\begin{enumerate}[label=\textup{(\roman*)},leftmargin=*,itemsep=3pt,topsep=4pt]
\item \textbf{Lifting inclusions of presented relations.} For a valid sequent, compose $P_\phi\onto Q_\phi\incl Q_\psi$ and lift through $P_\psi\onto Q_\psi$ by projectivity.
Conversely a map over the context gives $Q_\phi\subseteq Q_\psi$ on regular images.
Pullback along the epic context cover reflects inclusion of subobjects, proving the original sequent.
Replacing proof covers leaves these images unchanged, so it leaves existence of a realization unchanged.
\end{enumerate}
\end{proof}

\begin{lemma}[Canonical modest proof-name covers]
\label{lem:modest-proof-covers}
On PER-presented carriers, every ordinary relation has a canonical modest partitioned proof-name cover. The formula constructions above then produce modest partitioned proof objects. Equality transport, substitution comparisons, and lifting a valid sequent between these covers use \cref{lem:modest-target-lift}, without external choice.
\end{lemma}
\begin{proof}
\leavevmode
\begin{enumerate}[label=\textup{(\roman*)},leftmargin=*,itemsep=3pt,topsep=4pt]
\item \textbf{Support cover of a relation.} For $R\incl X_{\vec x}$, the pullback $Q_R=P_{\vec x}\times_{X_{\vec x}}R$ is a subobject of the modest object $P_{\vec x}$. It is modest by \cref{lem:modestclosure}. Its PER support cover $P_R\onto Q_R$ is partitioned and modest; its endpoint is the composite into $P_{\vec x}$. This construction retains all support names and selects no name for a semantic element.
\item \textbf{Equality transport.} Let $T$ contain a proof name together with a context-equality name. Saturation of the interpreted relation supplies the map $T\to Q_R$ with the transported endpoint. The object $T$ is partitioned, so \cref{lem:modest-target-lift} constructs the upper arrow in
\[
\begin{tikzcd}[column sep=large,row sep=large]
 T\ar[r,"\operatorname{tr}"]\ar[d]&P_R\ar[d,two heads]\\
 Q_R\ar[r,"\id"']&Q_R.
\end{tikzcd}
\]
The endpoint equation follows after the cover, and transport equations hold after interpretation.
\item \textbf{Formula induction and sequents.} Products, tagged finite sums, and carrier subsets of modest partitioned assemblies remain modest and partitioned: their codes determine at most one element. The existential clause retains the witness code and only changes the endpoint map. Hence every formula proof object constructed from these atomic covers is modest. The epicity proof of \cref{lem:formulacovers} uses only images and pullbacks and applies unchanged. A valid sequent supplies $P_\phi\to Q_\psi$; lift through the modest target $P_\psi\onto Q_\psi$. The same argument in either direction compares substitution covers of the same relation. It constructs each requested witness transformation, uniformly in any tracked parameter family.
\end{enumerate}
\end{proof}

\subsection{Search as a witnessed existential operation}
\label{subsec:search-as-a-witnessed-existential-operation}

For a modest $X_\rho$, form a partitioned assembly $C_\rho$ of total extensional stream evaluators with parameters. An element is a triple $(\sigma,t,a)$, with $\sigma:|\rho|\times\N\to2$ and $t,a\in K_1$, satisfying
\[
    t\langle a,x,n\rangle=\sigma(x,n)\in\{0,1\}
    \quad(x\in|\rho|),
\]
whose finite-hit support is $\rho$-invariant. Name it by $\langle t,a\rangle$. Evaluation is tracked by a fixed combinator. Every tracked family with partitioned parameter $P$ specializes this evaluator by $p\mapsto(\sigma_p,t,\varepsilon_Pp)$.

Let $W_\rho\to C_\rho$ be the partitioned object of triples consisting of an evaluator, a valid $\rho$-name, and a hit stage. A \textbf{split search presentation} consists of a tracked stream $S:C_\rho\to2^{\N}$ and tracked maps over $C_\rho$
\begin{equation}
\label{eq:splitsearch}
    W_\rho\longrightarrow H_S,
    \qquad H_S\longrightarrow W_\rho.
\end{equation}
The maps need not be inverse; they establish equality of the two regular images and retain finite witnesses in each direction.

\begin{proposition}
\label{prop:searchpresentation}
    A modest object is $\Sigma$-overt if and only if it admits a split search presentation.
\end{proposition}
\begin{proof}
\leavevmode
\begin{enumerate}[label=\textup{(\roman*)},leftmargin=*,itemsep=3pt,topsep=4pt]
\item \textbf{Search witnesses from overtness.} A code $\langle t,a\rangle$ determines the entire function $\sigma$ on valid inputs, so $C_\rho$ is modest as well as partitioned. The retained evaluator, input, and stage codes likewise make $W_\rho$ and every $H_S$ modest and partitioned. If $X$ is overt, apply $\exists_X^\Sigma$ to the universal evaluator. Lift its classifier on $C_\rho$ through $2^\N\onto\Sigma$ by \cref{lem:modest-target-lift}. Both $W_\rho$ and $H_S$ cover the same existential subobject. Apply that lemma to each modest target cover to construct both maps in \eqref{eq:splitsearch}.
\item \textbf{Converse construction.} Conversely, take the canonical partitioned cover of $\Sigma^X$ and lift evaluation on its product with $P_\rho$ through the modest stream object by \cref{lem:modest-target-lift}.
Specialize the universal search presentation.
The two witness maps identify its hit image with the existential image of evaluation.
Therefore its $\Sigma$-classifier descends along the parameter cover and is the required $\exists_X^\Sigma$.
Equality of images also proves independence of evaluator names and naturality under substitution.
\end{enumerate}
\end{proof}

\subsection{Substitution and composition of search witnesses}
\label{subsec:substitution-and-composition-of-search-witnesses}

The universal evaluator is useful because it includes its own parameter name. If a family on $P\times P_\rho$ is tracked by $t$, specialization sends a name $a$ of $p$ to $\langle t,a\rangle$. A fixed combinator tracks this map; it need not inspect the external function component $\sigma$. Pulling the two maps in \eqref{eq:splitsearch} back along specialization gives, for every such family, both a witness-to-hit and a hit-to-witness transformation. This proves the uniformity required of an existential operation.

For example, suppose two positive relations $R(x,y)$ and $S(y,z)$ are presented by streams and the middle sort is searchable. Their conjunction has proof names $(x,y,z,n,k)$ containing a hit of each relation. Search over $y$ produces the stream for relational composition. Its decoder returns a name $y$ together with both hit witnesses; its encoder turns any such pair of hits into a hit of the composite. Hence the composite stream represents the actual regular image of the conjunction projection.

If a composite is formed three times, the two parenthesizations need not produce identical streams or identical search schedules. Their decoded witnesses both describe $R(x,y)\wedge S(y,z)\wedge T(z,w)$. The formula-cover theorem and \cref{lem:modest-proof-covers} give maps between their hit presentations over $(x,w)$; equivalently they have the same classifier. This is the external form of the associativity proved by restricted Fubini. It also explains the appropriate coherence condition: equations hold after interpretation, with explicit witness translations available on presentations.


The witness calculus retains more information than the regular image alone: it records a tracker converting each premise proof into a conclusion proof, together with any existential witness. That mechanism is now ready to be applied to whole signatures and their homomorphisms.

\section{The presentation theorem}
\label{sec:presentationtheorem}

The preceding witness lemmas concern individual predicates and searches. We now assemble them into an equivalence between tracked positive-geometric presentations and internal models. The general-carrier version uses exact-completion equality covers rather than relying on modestness of every sort.


Fix a finitary positive-geometric signature with specified carrier sorts, total operations, ordinary relations, and positive relations. Some carrier sorts may be designated searchable. A positive relation is interpreted by a $\Sigma$-classifier. A \textbf{split external model} consists of PER presentations of the carriers, tracked operations, proof-name presentations of ordinary relations with equality transport, saturated tracked streams for positive relations, split search presentations for searchable sorts, and tracked proof transformations for the axioms. An arrow consists of tracked extensional carrier maps preserving the operations and relations, with proof transformations witnessing ordinary relation preservation. Equality of arrows is equality of the induced carrier maps. Choices of witnesses are witness data, not additional components of arrow equality.


\subsection{Morphisms and their equality}
\label{subsec:morphisms-and-their-equality}

To check that these data form a category, one must verify that preservation is independent of the chosen carrier trackers. Suppose $f_S$ and $g_S$ have PER-equivalent values at every valid name of every sort. For an operation $\omega$, preservation by $g$ follows from preservation by $f$, extensionality of the target operation, and symmetry and transitivity of its PER. For an ordinary relation, let $T_f$ be a proof transformation over $f$. The pair of output names $(f(\vec a),g(\vec a))$ is a tracked equality name. Define
\[
    T_g(u)=\operatorname{tr}_{R'}
       \bigl(T_f(u),(f(\partial u),g(\partial u))\bigr).
\]
Its endpoint is $g(\partial u)$, as required. For a positive relation, saturation transports a target hit from the $f$-image to the $g$-image. Thus every preservation condition depends only on the interpreted carrier maps.

Identity carrier trackers have identity proof transformations. Composition uses composition of the carrier trackers and of the ordinary proof transformations; positive support inclusions compose. The preceding independence check makes this composition well-defined on arrow classes. Associativity and the identity equations hold in the PER quotients. Selected axiom realizers and search witnesses express properties of model objects; homomorphisms are not required to preserve their particular search schedules or proof names.


\begin{theorem}[Relative presentation theorem]
\label{thm:geometric}
    Internalization gives an equivalence between split external models and internal models on modest carriers, with positive relations $\Sigma$-classified and designated searchable carriers overt.
\end{theorem}
\begin{proof}
\leavevmode
\begin{enumerate}[label=\textup{(\roman*)},leftmargin=*,itemsep=3pt,topsep=4pt]
\item \textbf{Relative presentation theorem.} Internalize PERs by their quotients, operations by descent, ordinary relations by regular images, positive relations by \cref{thm:finitehit}, and search by \cref{prop:searchpresentation}.
Axiom transformations imply validity by \cref{lem:extraction}.
This defines the functor on objects and arrows.
\item \textbf{Essential surjectivity.} For essential surjectivity, select PER covers of the modest carriers.
Lift every operation through its target cover on the projective product of source covers.
Pull back every ordinary relation and select a projective proof-name cover; extract its equality transport as above.
Lift each positive classifier through the stream cover.
Theorems on finite hits and search give the corresponding external positive data.
Construct proof objects for the formulas and apply \cref{lem:extraction} to the valid axioms.
\item \textbf{Fullness.} For fullness, lift an internal homomorphism on carrier covers.
Operation preservation is equality after interpretation.
For an ordinary relation, pull back its target proof cover along the map from the source proof object and split the resulting epimorphism.
Positive relation preservation is support inclusion by \cref{lem:support}.
Faithfulness is \cref{prop:per} on each carrier.
These constructions prove the claimed equivalence.
\end{enumerate}
\end{proof}

\begin{remark}[Constructed witnesses in the modest presentation theorem]
\label{rem:modest-presentation-choice}
For the canonical PER support covers, operation lifts and carrier homomorphisms in \cref{thm:geometric} are constructed by \cref{lem:modest-target-lift}. Ordinary proof covers and their transformations are constructed by \cref{lem:modest-proof-covers}; positive streams and search witnesses are constructed by \cref{thm:finitehit,prop:searchpresentation}. Thus none of these individual lifts, or their uniformly tracked parameter versions, requires external choice. The theorem also allows arbitrary selected partitioned proof covers. Lifting into such a nonmodest fixed target still uses the external projectivity explicitly recorded in \cref{rem:external-projectivity}. The cover-span construction of \cref{thm:refinement-quotient,thm:partial-refinement-quotient} supplies the corresponding unrestricted semantic recovery without making those selections. In particular, the construction of a witness transformation is distinguished from a simultaneous selection of trackers for an arbitrary externally indexed collection of transformations.
\end{remark}

\subsection{Exact-completion presentations without modestness}
\label{sec:generalrealizability}
For an arbitrary carrier $X$, choose a partitioned cover $q_X:P_X\onto X$ and a partitioned proof cover
\[
 E_X\onto P_X\times_XP_X.
\]
Equality evidence retains both endpoints and a proof name. Reflexivity, symmetry, and transitivity are tracked lifts of the corresponding kernel-pair operations, taking a projective refinement of a premise object when needed. Their equations hold after interpretation. Ordinary relation presentations have equality transport using these equality proof names. Positive relations retain saturated finite-hit presentations on the carrier cover.

Conversely, such equality data on a partitioned carrier define an equivalence relation by regular image, and exactness supplies its quotient. A bare external equivalence relation on realizer labels is insufficient in general: different carrier elements can have the same label, and equality in the quotient has its own evidence. This is the exact-completion form of the effective-topos presentation \cite{birkedaloosten,stekelenburg,frey}.

\begin{lemma}[Search on a general carrier]\label{lem:generalsearchpresentation}
The witnessed-search characterization extends to any carrier $X$. Choose a projective cover $V\onto\Sigma^X$, a carrier cover $P_X\onto X$, and a stream $\alpha$ presenting their pulled-back evaluation. Put $W=\{(v,p,n):\alpha(v,p,n)=1\}$. Then $X$ is overt precisely when there is a stream $S$ on $V$ with tracked maps $W\to H_S$ and $H_S\to W$ over $V$.
\end{lemma}
\begin{proof}
\leavevmode
\begin{enumerate}[label=\textup{(\roman*)},leftmargin=*,itemsep=3pt,topsep=4pt]
\item \textbf{Search on a general carrier.} If $X$ is overt, its existential evaluation classifier pulls back to $V$ and lifts to a stream $S$.
The two displayed projective objects cover the same image, so projectivity supplies both maps.
\item \textbf{Converse construction.} Conversely these maps identify the hit support with the existential image of evaluation along the carrier cover.
That image is invariant under the kernel pair of $V\onto\Sigma^X$.
Since $\Sigma\incl\Omega$ is monic, its classifier descends to a map $\Sigma^X\to\Sigma$, whose composite with the inclusion is existential quantification.
This is overtness.
Pullback along a projectively presented family specializes the two maps and proves parameter uniformity.
\end{enumerate}
\end{proof}

\begin{theorem}[General relative presentation theorem]\label{thm:generalgeometric}
For a finitary positive-geometric signature, partitioned carrier presentations with equality proof covers, relation proof covers, tracked operations, and tracked axiom transformations, modulo equality of interpreted carrier maps, present the category of all internal models and homomorphisms. Specified positive relations are represented by finite-hit streams, and designated overt carriers by the witnessed search of \cref{lem:generalsearchpresentation}. This extends \cref{thm:geometric} beyond modest carriers.
\end{theorem}
\begin{proof}
\leavevmode
\begin{enumerate}[label=\textup{(\roman*)},leftmargin=*,itemsep=3pt,topsep=4pt]
\item \textbf{Image and quotient construction.} Take the regular image of each equality proof cover.
The tracked equality transformations prove reflexivity, symmetry, and transitivity of that image.
Form its effective quotient.
Equality transport makes the images of the other relation covers descend to these quotients, and extensionality makes the operation lifts descend.
Their equations and the interpreted axioms follow by taking images of the supplied proof transformations.
\item \textbf{Proof-object construction.} For the formula induction use a projective proof cover of equality, products over the context for conjunction, tagged sums for disjunction, and retention of the witness name for existential quantification.
If a resulting premise object is not already projective, cover it projectively.
Pullback stability and composition of regular epimorphisms prove exactly the image equations of \cref{lem:formulacovers}.
Thus witness extraction applies to every finite axiom sequent.
\item \textbf{Image and quotient construction.} Positive predicates use \cref{thm:finitehit}, which applies on the selected projective carrier cover without requiring a modest quotient.
\item \textbf{Converse construction.} Conversely an internal model supplies the displayed covers.
Projectivity lifts its operations and the maps expressing valid sequents.
It also supplies equality transport, while \cref{lem:generalsearchpresentation} treats any designated overt sorts.
An internal homomorphism lifts on the carrier covers, and its relation-preservation maps lift on the premise proof covers.
Hence interpretation is full.
It is faithful by the stated equality of arrows and epicity of the carrier covers, and essentially surjective by the construction just given.
\item \textbf{Identification of the resulting maps.} Common projective refinements compare all choices and make substitution coherent after interpretation.
\end{enumerate}
\end{proof}

\subsection{Tracked presentations of all algebraic and general domains}
\label{sec:trackedcovers}
The order, cover endpoints, row membership, and observation generators of \cref{def:directedcover} form finitary presentation data. For example, row ideality is expressed by
\begin{align*}
 \top&\vdash_c\exists p\,I_c(p),\\
 I_c(p)\wedge q\leq p&\vdash I_c(q),\\
 I_c(p)\wedge I_c(q)&\vdash
       \exists r\,[I_c(r)\wedge p\leq r\wedge q\leq r].
\end{align*}
For an observation-generator sort $G$ with evaluation $U(g,p)$, upperness and cover compatibility are the sequents
\[
 U(g,p)\wedge p\leq q\vdash U(g,q),\qquad
 U(g,\ell(c))\vdash\exists p\,[I_c(p)\wedge U(g,p)].
\]
These are ordinary relations with proof names; any independently specified positive relation additionally has its stream presentation. Form the observation closure and the directed-cover completion internally after interpreting these data. Consequently neither the completed carrier nor its full ideal object needs to be supplied as a generating sort.

\begin{theorem}[General and algebraic realizability characterizations]\label{thm:fulladequacy}
In $\Eff$, tracked observed dcpo presentations and their interpreted completion maps characterize all effective dcpos and effective Scott-continuous maps. Empty-cover presentations characterize all algebraic effective dcpos. These comparisons preserve and reflect equality of maps and commute with parameter substitution and refinement of presentations.
\end{theorem}
\begin{proof}
\leavevmode
\begin{enumerate}[label=\textup{(\roman*)},leftmargin=*,itemsep=3pt,topsep=4pt]
\item \textbf{Interpretation of presentation data.} Interpret the generating data by \cref{thm:generalgeometric}, take the internal observation closure, and apply \cref{thm:observedcovers}.
For a presentation arrow, present its map from the source generators into the completed target on projective covers.
Include the proof-name transformations for cover preservation and observation preservation; these express the actual relations in the completed target and its observation object.
Their regular images give exactly the conditions in the universal property.
\item \textbf{Lifting effective maps.} Conversely any effective map restricts to such a map on generators and lifts with these proofs by projectivity.
Thus extension and restriction give inverse correspondences on arrows.
\item \textbf{Canonical recovery.} Every effective dcpo occurs by the canonical presentation of \cref{thm:allcoverpresentations}, taking its given observations as generators.
With no covers, \cref{thm:algebraicpresentation} identifies the image with all algebraic effective dcpos, and \eqref{eq:algebraicmatrix} identifies the maps.
Common refinements preserve the interpreted generating data and the universal properties, proving independence and parameter compatibility.

The universal property is represented by the commutative diagram
\[
\begin{tikzcd}[column sep=large,row sep=large]
P\ar[r,"\eta_C"]\ar[dr,"f"']&D_C(P)\ar[d,"\bar f"]\\ &E
\end{tikzcd}
\]
\end{enumerate}
\end{proof}

The completion and observation closure are bounded higher-order constructions. The preceding theorem does not replace them by a finite list of first-order axioms, nor assert that a completed carrier is modest. Its finite sequents present the generators and cover laws; the universal property supplies the completed object. Morphism evidence refers to the completed object and its universal property. In the searchable continuous case, positive rounded distributors give the more explicit calculation of \cref{cor:approxcomparison,thm:adequacycomparison}.


\subsection{Application to approximation presentations}
\label{subsec:application-to-approximation-presentations}

For reference, the approximation signature has one carrier $B$, an ordinary relation $\leq$, and a positive relation $\rho$. In addition to reflexivity and transitivity of $\leq$, its finite axiom sequents are
\begin{align*}
    a'\leq a\wedge\rho(a,b)&\vdash\rho(a',b),\\
    \rho(a,b)\wedge b\leq b'&\vdash\rho(a,b'),\\
    \top&\vdash_b\exists a\,\rho(a,b),\\
    \rho(a,b)\wedge\rho(a',b)&\vdash
    \exists c\,[\rho(c,b)\wedge a\leq c\wedge a'\leq c],\\
    \rho(a,b)&\vdash a\leq b,\\
    \rho(a,b)&\dashv\vdash\exists c\,[\rho(a,c)\wedge\rho(c,b)].
\end{align*}
Each displayed equivalence abbreviates two sequents. These express, in order, lower closure, source monotonicity, inhabitedness, directedness, refinement, and interpolation. The specification includes searchability of the carrier separately through its split existential-search operation. 

For a searchable approximation presentation on a modest carrier, its ordinary order is an ordinary relation with proof names. The approximation predicate is a positive relation. Its structural sequents are reflexivity and transitivity of order, monotonicity and lower closure of approximation, inhabitedness and directedness of each row, refinement, and the two implications of interpolation \eqref{eq:rho}. Each is finitary positive-geometric. Thus \cref{thm:geometric} gives their external presentation by tracked witnesses. For example directedness transforms witnesses of $\rho(a,b)$ and $\rho(a',b)$ into a name $c$ and witnesses of
\[
    \rho(c,b),\qquad a\leq c,\qquad a'\leq c.
\]
Interpolation returns an intermediate approximation name. Relational composition uses the split search operation on the middle carrier. The same presentation method applies to ideal-valued rounded relations and their witness equations.

\begin{corollary}
\label{cor:approxcomparison}
    External split searchable approximation presentations and rounded relations present the full category of their internally generated effective continuous dcpos and effective maps, up to isomorphism. The comparison preserves and reflects equality of represented maps and is compatible with composition, finite products, and domain lifting.
\end{corollary}
\begin{proof}
\leavevmode
\begin{enumerate}[label=\textup{(\roman*)},leftmargin=*,itemsep=3pt,topsep=4pt]
\item \textbf{Product calculation.} Apply \cref{thm:geometric} to presentation objects and their relational diagrams, then \cref{thm:relations} to rounded completion.
Finite products use the product carrier and conjunction of approximation relations.
Lifting uses the tagged carrier and the case definition in \cref{thm:effectivelift}.
Both retain searchable carriers and tracked witnesses.
Composition is \eqref{eq:sigmacomposition}; the split search maps supply its witnesses.
Equality is equality of completed maps by \cref{thm:relations}.
\end{enumerate}
\end{proof}

\subsection{Witnesses for a rounded relation}
\label{subsec:the-witnesses-for-a-rounded-relation}

The application to maps uses a diagram theory with two carrier sorts $B,C$ and a positive relation $m(b,c)$. This is distinct from asking for an ordinary homomorphism of basis preorders: a domain map is represented by an ideal-valued relation. Its structural witnesses can all be read from the following finite sequents, together with the order and approximation sequents of its source and target:
\begin{align*}
    \top&\vdash_b\exists c\,m(b,c),\\
    m(b,c)\wedge c'\leq c&\vdash m(b,c'),\\
    m(b,c)\wedge m(b,c')&\vdash
    \exists d\,[m(b,d)\wedge c\leq d\wedge c'\leq d],\\
    b\leq b'\wedge m(b,c)&\vdash m(b',c),\\
    m(b,c)&\dashv\vdash
    \exists a,d\,[\rho(a,b)\wedge m(a,d)\wedge\sigma(c,d)].
\end{align*}
Here the last line denotes two sequents. Inhabitedness gives a selected target approximation from an input name; directedness combines two output approximation witnesses; roundedness returns the intermediate input and output approximants. Ordinary order conclusions carry proof names, while positive conclusions carry finite hits. The presentation theorem extracts exactly these different kinds of evidence.

To compose $m:B\to\Idl(C)$ and $n:C\to\Idl(D)$, the positive classifier is $\exists c\,[m(b,c)\wedge n(c,d)]$. Search on $C$ supplies its witness maps. A directedness witness for the composite first bounds the two intermediate $C$-names using directedness of the row of $m$, transports the two $n$-witnesses upward to that common name, and then uses directedness of the corresponding row of $n$. Lower closure and monotonicity are obtained by the same direct transport of witnesses. The roundedness equations follow by substituting the source and target roundedness equations and associating existential quantifiers. These constructions verify closure of the external relation category under composition, including its witnesses.

Finally its identity relation is $\rho$ at $(B,r)$. The left and right identity equations are the two roundedness equations for $m$, so they come with the already selected proof transformations. Internalization sends this external identity to the identity of $D_r$ by \cref{thm:relations}. This gives the explicit connection between the presentation theorem for finite sequents and the category of effective continuous domains.


\subsection{General algebras and higher-order structure}
\label{subsec:general-algebras-and-higher-order-structure}

Besides the generating presentation of \cref{thm:fulladequacy}, the algebra definition has a direct structural presentation. Choose projective covers of $D$, $\OO_D$, and $\Idl(D)$, and proof covers of their structure relations. Lift evaluation, the unit, and the supremum map on their projective sources. Present the algebra equations by maps into the appropriate kernel pairs. The lifted-monad identities \eqref{eq:liftchecks} then translate the effectiveness of supremum into the Scott-observation equation \eqref{eq:scottobs}.

 The universal quantifiers occurring in directed completeness are interpreted in that object; they are not inferred from a finite list of external order axioms. For arbitrary carriers their equality relations are represented by proof-name covers on the selected carrier presentations.

\begin{proposition}[Transport of general algebra presentations]
\label{prop:generalpresentation}
    Choose exact-completion presentations of the carrier, observation, and ideal objects of an effective dcpo. Tracked lifts of the structural maps, modulo equality after their quotient maps, present the same algebra category. Descent preserves and reflects the unit, multiplication, and observation equations, independently of the selected projective refinements.
\end{proposition}
\begin{proof}
\leavevmode
\begin{enumerate}[label=\textup{(\roman*)},leftmargin=*,itemsep=3pt,topsep=4pt]
\item \textbf{Kernel-pair descent.} Every structural arrow lifts on its projective source cover through its target cover.
Any two lifts have the same interpretation exactly when their pair maps into the target kernel pair.
Thus the presentation category, with arrows identified in this manner, has precisely the original hom-sets, and composition of lifts represents composition after descent.
This argument applies also to the maps on the actual ideal objects and their iterates.
\item \textbf{Kernel-pair characterization of algebra equations.} Consequently the two sides of each algebra equation have equal interpretations if and only if their lifted pair factors through the corresponding kernel pair.
Membership of an observation is represented by its selected proof cover, so the same argument and \cref{lem:extraction} identify effectiveness of the supremum map. A common projective refinement of two covers gives comparison maps in both directions over the represented object.
Their composites interpret the identity, proving independence of presentations.
\item \textbf{Ideal and directedness calculation.} The ideal object is the interpreted power-object subobject defined in \cref{sec:ideal}; transporting this object retains its universal property and therefore the quantification over all ideals.
\end{enumerate}
\end{proof}


\subsection{The algebra equations on names}
\label{subsec:the-algebra-equations-on-names}

The algebra laws admit the following kernel-pair presentation. Let
\[
    q_D:P_D\onto D,\qquad
    q_T:P_T\onto\Idl(D),\qquad
    q_{TT}:P_{TT}\onto\Idl(\Idl(D))
\]
be selected projective covers of the actual objects. Choose tracked lifts
\[
    \widehat y:P_D\to P_T,\qquad
    \widehat a:P_T\to P_D,\qquad
    \widehat\mu,\widehat{\T a}:P_{TT}\to P_T
\]
of the unit, algebra map, multiplication, and lower image of the algebra map. These lifts have interpretation equations, for example $q_D\widehat a=a q_T$ and $q_T\widehat\mu=\mu_Dq_{TT}$. The unit law is the statement that
\[
    (\widehat a\widehat y,\id_{P_D}):P_D\to P_D\times P_D
\]
factors through the kernel pair of $q_D$. The multiplication law is the corresponding factorization of
\[
    (\widehat a\widehat\mu,\widehat a\widehat{\T a})
                  :P_{TT}\to P_D\times P_D.
\]
If selected projective proof covers of these kernel pairs are used, projectivity supplies tracked proof-name maps for the factorizations. Conversely such proof maps give the two equations after descent. This is a direct presentation of the algebra laws, with no choice of a representative ideal for each set of ideal names.

For a homomorphism $f:D\to E$, choose lifts of $f$ and of its actual lower-image map $\T f$. Its algebra equation is represented by the pair of composites lifting $fa_D$ and $a_E\T f$, again with target the kernel pair of the selected cover of $E$. Identity and composition preserve these equations because they preserve the interpreted composites. The observation component is treated by covering $\OO_D$ and $\OO_E$, lifting inverse image, and requiring equality of the two evaluation maps. For supremum, that equality is exactly \eqref{eq:scottobs}. The existential side has proof names containing an ideal member and an observation witness, so \cref{cor:sequents} supplies the witness transformations once the actual ideal object is presented.

Refining a selected cover changes the lifts but leaves every displayed quotient map and equation unchanged after interpretation. A common projective refinement compares the two systems. Any two comparisons interpret the same maps and hence are equal modulo the corresponding kernel pairs. This gives the independence claimed in \cref{prop:generalpresentation} at the level of the algebra, its morphisms, and its observation structure. The higher-order content lies in presenting the actual ideal object and its defining universal property; the remaining algebra equations use the same projective lifting mechanism as the finitary presentation theorem.

This is the direct structural exact-completion description. The generating presentation in \cref{thm:fulladequacy} constructs the completed object from cover data; \cref{cor:approxcomparison} gives the positive-matrix presentation in the modest searchable fragment. The latter is what permits calculations on approximation names without presenting an entire higher-order ideal object separately.


The equivalence covers arbitrary interpreted carriers, including the higher-order ideal and observation objects when they are explicitly presented. The searchable rounded fragment has a simpler matrix calculus, which the adequacy theorem will recover by applying the same witness transformations to positive approximation data.

\section{Partial-map adequacy}
\label{sec:adequacy}

A partial program is not determined merely by where its tracker happens to halt: it has a specified domain predicate and a value defined on that predicate. We compare these data with dominance-classified spans, using hit assemblies to retain finite witnesses. The resulting restriction-category equivalence is the implementation bridge for later evaluators.


An external partial map from $X_\rho$ to $X_\tau$ is represented by a saturated stream $\theta$ on $P_\rho$ and a tracker $f\in K_1$ defined on $\Hit(\theta)$, with values in $|\tau|$ and extensional there. Its domain is specified by the stream, regardless of any additional applications at which $f$ happens to be defined. Two representatives are equal when their hit supports agree and their outputs are $\tau$-equivalent on that support.

\begin{theorem}[Partial-map adequacy]
\label{thm:partialadequacy}
    These external partial maps form a restriction category equivalent to the full subcategory of $\Par_\Sigma(\Eff)$ on modest objects. The equivalence is compatible with parameters and with the partial-map classifier.
\end{theorem}
\begin{proof}
\leavevmode
\begin{enumerate}[label=\textup{(\roman*)},leftmargin=*,itemsep=3pt,topsep=4pt]
\item \textbf{Composition.} For composition of $(\alpha,f)$ and $(\beta,g)$ use value tracker $g\circ f$ and domain stream
    \[
        \gamma_a(\langle n,k\rangle)=
        \begin{cases}
            \beta_{fa}(k),&\alpha_a(n)=1,\\
            0,&\alpha_a(n)=0.
        \end{cases}
    \]
    The guarded branch is total on its stated inputs.
Its support is exactly the sequential domain.
Identity uses the constantly-one stream and the identity tracker; restriction uses the original stream and identity on names.
Associativity and the restriction equations follow from equality of these domains and values.
\item \textbf{Stream representation of partial domains.} To externalize an internal partial map, lift its classifier to $\theta$ on $P_\rho$ through the modest stream object by \cref{lem:modest-target-lift}.
The partitioned hit assembly covers the pullback domain.
Apply the same lemma to lift its value map to the modest support assembly $P_\tau$, obtaining a tracked $v(a,n)$.
Search for a first hit $n_a$ and output $v(a,n_a)$.
On the stated domain search terminates and the output is extensional because it interprets the original value map.
\item \textbf{Converse construction.} Conversely, the classifier descends by saturation.
On $H_\theta$, send $(a,n)$ to $fa$.
Extensionality makes its interpretation constant on the kernel pair of the hit cover, so it descends to a value map on the classified domain.
The two constructions are inverse on represented partial maps.
Equality is reflected by support equality and PER equality of outputs.
All constructions can retain an additional parameter name; reindexing gives the same interpreted span.
\item \textbf{Monad compatibility.} The natural partial-map-classifier correspondence of \cref{prop:partialmonad} therefore identifies their unit and composition operations.

The interpretation maps form the commuting diagram
\[
\begin{tikzcd}[column sep=large,row sep=large]
H_\theta\ar[r,"v"]\ar[d,two heads,"\pi"']&P_\tau\ar[d,two heads,"q_\tau"]\\ D_\theta\ar[r,"u"']&X_\tau
\end{tikzcd}
\]
\end{enumerate}
\end{proof}

\begin{proposition}[Exact implementation with modest target names]
\label{prop:modest-partial-implementation}
Let $q:C\onto X$ be a partitioned name cover and $r:D\onto Y$ a modest assembly cover. Every $\Sigma$-partial map $f:X\pmap Y$ has a tracked partial implementation on $C$ with domain exactly $q^{-1}(\dom f)$ and output in $D$. No external choice is used. The interpreted domain--value spans commute with parameter substitution.
\end{proposition}
\begin{proof}
\leavevmode
\begin{enumerate}[label=\textup{(\roman*)},leftmargin=*,itemsep=3pt,topsep=4pt]
\item \textbf{Domain witnesses.} Put $A=C\times_X\dom f$. Lift its $\Sigma$-classifier through $2^\N\onto\Sigma$ using \cref{lem:modest-target-lift}. The resulting stream $\alpha$ has a partitioned hit assembly $H_\alpha\onto A$, with the exact image $A$ by \cref{thm:finitehit}.
\item \textbf{Values and the first hit.} Apply \cref{lem:modest-target-lift} once more:
\[
\begin{tikzcd}[column sep=large,row sep=large]
 H_\alpha\ar[r,"v"]\ar[d,two heads]&D\ar[d,two heads,"r"]\\
 A\ar[r,"{f\pi_{\dom f}}"']&Y.
\end{tikzcd}
\]
On a name of $c$, search the tracked stream for its first hit and run the tracker of $v$ on that hit name. The search halts exactly on $A$. Modesty of $D$ makes the resulting output name determine a unique element of $D$, so this specifies a tracked partial map with precisely the required interpretation.
\item \textbf{Parameters.} Pulling back the domain mono and the displayed value square gives the implementation in a new context. Any stream or value representatives obtained by repeating the construction have the same classified image and the same value after $r$. Thus their interpreted spans agree, and iterated substitutions agree by the pullback universal property.
\end{enumerate}
\end{proof}

\begin{corollary}[Finite execution witnesses on presented carriers]\label{cor:presented-execution-witnesses}
For an implemented partial operation on arbitrary presented carriers, pulling back its displayed domain stream and retaining the lifted output value gives the corresponding finite-hit witness cover, with exactly the semantic domain after interpretation. Carrier validity and equality evidence are retained in their proof covers.
\end{corollary}
\begin{proof}
\leavevmode
\begin{enumerate}[label=\textup{(\roman*)},leftmargin=*,itemsep=3pt,topsep=4pt]
\item \textbf{Partiality and definedness.} Pull the implemented operation's classifier back to its partitioned input cover and lift it through the modest stream object by \cref{lem:modest-target-lift}. Its supplied implementation gives the output on the hit assembly. Conversely, extracting a fixed output-name map for an arbitrary target uses the projectivity of \cref{rem:external-projectivity}; modest target names use \cref{prop:modest-partial-implementation}. For arbitrary target covers, \cref{thm:partial-refinement-quotient,prop:refined-finite-returns} constructs the corresponding domain and finite-return witnesses by covered spans without this selection.
Its hit assembly covers the pulled-back semantic domain.
The lifted value map gives the denotation on that cover.
Changing a stream representative transports hits by first-hit search, and changing an output name retains a witness in the output kernel pair.
Thus equality concerns domains and interpreted values.
The name presentation retains its validity predicate.
\end{enumerate}
\end{proof}

For any $z:P\to \LL X_\tau$ with $P$ partitioned, \cref{prop:modest-partial-implementation} supplies a stream for its extent and a tracked value on its hit assembly. Nested applications give the guarded stream for the monad multiplication. These are particularly useful presentations of partial elements because the value need only be computed after a finite witness of definedness has appeared.

\subsection{Names for partial elements and multiplication}
\label{subsec:names for partial elements and multiplication}

Over a partitioned parameter object $P$, a presentation of a map into $\LL X_\tau$ consists of a tracked stream $\theta$ and a tracked map $v:H_\theta\to P_\tau$ satisfying
\[
    \tau(v(p,n),v(p,m))\qquad
    \text{whenever $\theta_p(n)=\theta_p(m)=1$}.
\]
Its equality relation with another presentation $(\eta,w)$ requires equal supports and $\tau(v(p,n),w(p,m))$ on pairs of hits over the same parameter. The proof is direct descent: the kernel pair of the hit cover $H_\theta\onto D_\theta$ is $H_\theta\times_P H_\theta$, since $D_\theta\incl P$ is monic. The displayed condition therefore says precisely that $q_\tau v$ descends to a value on $D_\theta$. This supplies the natural bijection between these presentation classes and maps $P\to \LL X_\tau$.

The unit has constant-one stream and the given total value name. For multiplication, take an outer stream $\alpha:P\to2^\N$, an inner stream $\beta:H_\alpha\to2^\N$, and an inner value map $w:H_\beta\to P_\tau$. Guarded flattening gives
\[
    \gamma_p(\langle n,k\rangle)=
    \begin{cases}
        \beta_{(p,n)}(k),&\alpha_p(n)=1,\\
        0,&\alpha_p(n)=0.
    \end{cases}
\]
There is a tracked isomorphism of hit assemblies
\[
    H_\gamma\cong H_\beta,
    \qquad(p,\langle n,k\rangle)\longleftrightarrow((p,n),k),
\]
and the flattened value is $w$ after this isomorphism.

It remains to establish independence of the chosen outer hit. The outer value map $H_\alpha\to \LL X_\tau$ descends through its domain $D_\alpha$. Thus its two pullbacks to $H_\alpha\times_P H_\alpha$ are equal maps into the lift. Applying the preceding partial-element classification on this partitioned pullback shows that the two inner presentations have equal supports and equivalent outputs. Hence any two hits of $\gamma$ over $p$ give $\tau$-equivalent final values, even if they used different outer hits. This proves that guarded flattening represents multiplication, without assuming that the intermediate lift is modest.

For three nested partial elements, both multiplication orders decode to triples of hit witnesses and the same final value. Pairing and unpairing stage codes identify these witness objects; the same descent argument removes their choice. The unit cases remove a constantly true stage. In this form the monad equations have explicit presentation-level witnesses, while their equality remains equality of represented partial elements.


Partial-map adequacy identifies the numerical stream representing definedness with the semantic extent, and it identifies returned names only after quotient interpretation. Those are the precise conditions required for numerical program covers and their fixed-point compilers to implement semantic partial evaluation.

\section{Numerical semidecision and searchable carriers}
\label{sec:examples}


Application in $K_1$ is partial computable application on numerical names. A total tracked binary stream is a computable stream with parameters. Its hit support is c.e.; conversely bounded simulation turns any c.e. predicate into such a stream. Thus \cref{thm:finitehit} identifies $\Sigma$-predicates on a modest presentation with extensional c.e. predicates on its names. \Cref{thm:partialadequacy} identifies numerical partial maps with partial computable functions, retaining the chosen domain predicate.


\begin{example}[Countable carriers need not be searchable]
\label{ex:nonovert}
    Let $S\subseteq\N$ be an external non-c.e. subset and give it the partitioned naming map $s\mapsto s$ in $\Eff$. Its inclusion $j:S\to\N$ is tracked, and equality of a name $j(s)$ with a numerical parameter is decidable. If this modest object were overt, the predicate
    \[
        n\longmapsto\exists_S^\Sigma s\;[j(s)=n]
    \]
    would be a $\Sigma$-predicate on $\N$. Finite-hit classification would then make its support $S$ c.e., a contradiction. Thus an externally countable carrier with a perfectly valid numerical naming map need not support effective existential search. 
\end{example}

In the effective topos, the finite-graph basis of Part~\ref{part:synthetic} and the rational basis of \cref{ex:lowerreals} are searchable because they have effective enumerations. Their completion objects can have a much more complicated order. For instance, inclusion of arbitrary partial-function ideals tests all inputs; it is not obtained by a finite-hit procedure merely because finite graph inclusion is decidable. The difference between basis comparison and completed comparison is preserved by the observation construction and by its realizability presentation.

\section{Adequacy of completion and its computational structure}
\label{sec:adequacycomparison}


The general characterization is \cref{thm:fulladequacy}: observed directed covers characterize all of $\ED$, and empty covers characterize its algebraic part. We now calculate the canonical semantic--computational comparison explicitly in the searchable fragment, where overtness and positive approximation give stream matrices for maps.

Let $\mathbf{Ap}_\Sigma$ have searchable approximation objects $(P,r)$ and rounded ideal-valued relations, with the composition of \cref{thm:relations}. Its intrinsic interpretation is particularly simple. Positivity and overtness make every row of $r$ an admitted ideal, hence $r:P\to\J P$ is a Kleisli idempotent. Splitting it in $\AS$ gives a computational domain $A_r$. Splitting the corresponding ordinary ideal extension gives $D_r$. The comparison identifies $A_r$ with $\Pt D_r$.

More precisely, $\mathbf{Ap}_{\Sigma,\mathrm{mod}}$ retains a modest overt basis, its ordinary preorder, and its positive approximation idempotent; its arrows are the positive rounded distributors. The category $\mathbf{Ap}_{\mathrm{trk}}$ retains PER presentations of those carriers and proof-name presentations of the ordinary relations, together with stream names for positive relations and the witnesses of their axioms. Arrows are taken modulo equality of their interpreted matrices. The category $\mathbf{Search}_{\Sigma,\mathrm{mod}}$ is the full effective semantic category on the $D_r$ arising from these objects. Finally, $\mathcal A_{\Sigma,\mathrm{sp}}$ is the full subcategory of $\AS$ on the algebras isomorphic to
\[
    A_r=\{I\in\J B\mid r^\#I=I\}.
\]
These definitions make explicit the objects and the equality of arrows in all four corners. No numerical enumeration of a general modest basis is included: its overtness is represented by the search operation of \cref{sec:extraction}.

\begin{theorem}[Adequacy characterization]
\label{thm:adequacycomparison}
    In $\Eff$, internalization identifies tracked searchable modest approximation presentations with internal ones, including their rounded maps. Completion then identifies their category with the full category of $\J$-algebras admitting such an intrinsic presentation. The semantic and computational completions satisfy
    \[
        \begin{tikzcd}[column sep=large,row sep=large]
            \mathbf{Ap}_{\mathrm{trk}}\ar[r,"\simeq"]\ar[d,"D_{(-)}"']
            &\mathbf{Ap}_{\Sigma,\mathrm{mod}}\ar[d,"\text{Kleisli splitting}"]\\
            \mathbf{Search}_{\Sigma,\mathrm{mod}}\ar[r,"\Pt"',"\simeq"]
            &\mathcal A_{\Sigma,\mathrm{sp}}.
        \end{tikzcd}
    \]
    The square commutes up to its canonical comparison isomorphism. It preserves composition and finite products, identifies positive observations, and is compatible with domain lifting and substitution in parameters.
\end{theorem}
\begin{proof}
\leavevmode
\begin{enumerate}[label=\textup{(\roman*)},leftmargin=*,itemsep=3pt,topsep=4pt]
\item \textbf{Adjunction and universal property.} The upper equivalence is \cref{thm:geometric} applied to the approximation and rounded-map signatures.
The right-hand construction is defined before using realizers, by splitting $r$ in $\AS$.
Its arrows are the rounded $\J$-Kleisli maps by the universal property of the Karoubi envelope.
The nonformal identification with the arrows on the left is \cref{prop:roundedhomcomparison}: rounding recovers positive membership from admission, so these are exactly the positive rounded distributors of \cref{thm:relations}.
\item \textbf{Coherence equations.} On free objects the comparison is the hom-isomorphism of $\F\dashv\Pt$; on splittings it is its unique restriction.
This proves commutativity and full faithfulness.
Essential surjectivity is precisely the intrinsic splitting property defining $\mathcal A_{\Sigma,\mathrm{sp}}$.
\item \textbf{Composition.} For composition, both routes send matrices $m,n$ to the classifier $\exists c\,[m(b,c)\wedge n(c,d)]$.
Its tracked implementation is the split search on the middle object.
Associativity is interchange and reassociation of these existential composites.
Identity matrices are the approximation idempotents on both sides.
Products use conjunction of approximation matrices and pairing of witnesses; \cref{app:products} supplies the explicit presentation.
\item \textbf{Unit law.} Unit and ideal multiplication are principal ideals and flattening, so the comparison equations are the unit and union equations of \cref{jd:thm:J}.
\item \textbf{Comparison of observation pairings.} Observations are the rounded positive traces by \cref{prop:roundedtrace}; their evaluations coincide with the pairings in \cref{jd:thm:spectra}.
Lifting uses the classified extent and its value ideal.
The distributive law of \cref{jd:thm:distributive} and the guarded realizability composite of \cref{thm:partialadequacy} therefore give the same partial element.
Equality of the extent classifiers and of values on their common domain proves the partiality comparison, including unit and multiplication.
\item \textbf{Pullback construction.} All constructions retain a parameter.
Pullback carries a classifier and its value together and preserves the defining finite limits, images, and dependent constructions in slices.
Two successive substitutions have the canonical pullback comparison.
At the presentation level a projective refinement supplies trackers, and equality descends through the cover.
This proves coherent parameter compatibility.
\end{enumerate}
\end{proof}

\begin{corollary}[Canonical algebraic recognition]
\label{cor:canonicalrecognition}
    An algebraic computational domain has an overt modest positively ordered algebraic basis precisely when its canonical compact part has these properties. Its canonical presentation is $\J K(A)\cong A$. The corresponding effective-topos presentation is obtained by applying the finite-hit and search comparisons to $K(A)$.
\end{corollary}
\begin{proof}
\leavevmode
\begin{enumerate}[label=\textup{(\roman*)},leftmargin=*,itemsep=3pt,topsep=4pt]
\item \textbf{Canonical algebraic recognition.} For a free algebra, \cref{jd:prop:compact} identifies the compact part with its generating poset.
Conversely the same proposition identifies $A$ with $\J K(A)$ whenever the canonical algebraicity comparison is invertible.
Apply \cref{jd:prop:basis,thm:geometric} to this basis.
These constructions commute with isomorphism and therefore do not depend on a redundant choice of generators.
\end{enumerate}
\end{proof}

\begin{proposition}[Change of approximation presentation]
\label{prop:changebasis}
    Inverse rounded maps between two presentations induce inverse semantic and computational maps, inverse positive-observation maps, and equivalent fibres of the presentation fibration constructed in Part~\ref{part:coding}.
\end{proposition}
\begin{proof}
\leavevmode
\begin{enumerate}[label=\textup{(\roman*)},leftmargin=*,itemsep=3pt,topsep=4pt]
\item \textbf{Naturality and reindexing.} Apply ideal extension to the inverse equations in the rounded Kleisli category, then apply $\Pt$.
Inverse image supplies the inverse observation maps.
Reindexing a name cover along either resulting isomorphism gives inverse equivalences of presentation fibres.
Common refinements compare independently selected projective covers.
The equations are identities after interpretation, so the conclusion is unchanged by different tracker representatives.
\end{enumerate}
\end{proof}

This proves \cref{main:B}. In particular, the adequacy assertion compares both the completion objects and their induced operations. The external presentation still records the distinction between ordinary proof relations and positive relations: it is the approximation predicate that is positive, while the completed order is the order of the resulting algebra.


\part{Extensional synthetic computability}
\label{part:synthetic}

The semantic comparison is now available, but computational claims require objects on which evaluation, observations, and recursion can be calculated. We construct them from finite graphs and positive tests, and then prove intrinsic finite-information and fixed-point results. These statements concern denotations; the subsequent intensional part supplies their program indices and tracked implementations.


The computational category is $\AS$ with its lifted partiality monad. When using all-ideal semantic representatives we work in $\CS\subseteq\ED$ and the equivalence of \cref{jd:thm:splitcomparison}; for statements also valid on the larger semantic category we retain the notation $\C=\ED$. Write $\Search(\E)$ for the full subcategory of $\ED(\E)$ on the effective dcpos isomorphic to searchable rounded completions. Searchable representatives form this category. This part constructs evaluations, observations, decisions, and solutions of recursion equations before a program-name cover is chosen. Its final sections establish the stronger effective-topos statements about arbitrary internal maps on the computation objects.


\section{A typed function object}\label{sec:programdomain}

The abstract domain theory has supplied operations on effective points, but numerical computability needs an explicit object of partial functions. Finite graphs form an overt positively ordered basis for such an object. This presentation makes evaluation and currying calculable without assuming that every semantic function already has a program index.


The completion and comparison theorems give computational objects from positive approximation data. We begin the computability development with a typed example in which the natural numbers occur as input and output data. It is constructed internally in $\Eff$ from the overt structures proved in \cref{thm:relativefoundation}. Equip the discrete dcpo $\N$ with observations $\Sigma^{\N}$. Its effective-point object is $\N$ itself. Write $\N_\bot$ for its domain lift and $\LL\N$ for its effective-point object, as in \cref{thm:effectivelift}.

Let $\FG$ be the object of finite partial graphs from $\N$ to $\N$. A graph is represented canonically by a finite list of pairs, sorted by first coordinate, with no repeated first coordinate. Membership, inclusion, compatibility, and compatible union are decidable. Finite lists over $\N$ admit an internal enumeration. Normalization of a valid list and replacement of an invalid list by the empty graph give a regular epimorphism $\N\onto\FG$. Thus $\FG$ is overt and its inclusion order is $\Sigma$-classified. This use of finite combinatorics concerns approximations to typed data; no function has been assigned a program index.

Define the effective algebraic dcpo
\begin{equation}\label{eq:programdomain}
 \FunDom=\Idl(\FG),\qquad
 \OO_{\FunDom}=\{U\in\Sc(\FunDom)\mid U(\downarrow g)\in\Sigma
                       \text{ uniformly in }g\}.
\end{equation}
For $I\in\FunDom$, put
\[
 R_I(n,m)\iff\exists g\in I\;(n,m)\in g.
\]
Directedness of $I$ implies functionality of $R_I$: two witnesses lie in a common graph in $I$, which cannot give different values at the same input.

\begin{proposition}\label{prop:graphrepresentation}
The underlying dcpo $\FunDom$ is isomorphic to the object of $\Omega$-partial maps $\N\pmap\N$, ordered by graph inclusion. The inverse of $I\mapsto R_I$ sends a partial graph $R$ to its finite subgraphs. There is a canonical isomorphism
\begin{equation}\label{eq:programpts}
 \Pt(\FunDom)\cong(\LL\N)^{\N}.
\end{equation}
\end{proposition}
\begin{proof}
\leavevmode
\begin{enumerate}[label=\textup{(\roman*)},leftmargin=*,itemsep=3pt,topsep=4pt]
\item \textbf{Ideal and directedness calculation.} Finite subgraphs of a functional relation form an ideal: the empty graph supplies inhabitedness, lower closure is immediate, and compatible union supplies directedness.
Every member of $I$ is a finite subgraph of $R_I$.
\item \textbf{Converse construction.} Conversely, finitely many witnesses in $I$ can be bounded by one member using directedness; their finite subgraph therefore belongs to $I$.
These constructions are inverse and turn directed unions into graph unions.
\item \textbf{Positive graph membership.} If $I$ is $\Sigma$-classified on $\FG$, then $R_I$ is $\Sigma$-classified by overtness.
Its domain at $n$ is $\exists m\,R_I(n,m)\in\Sigma$ by overtness of $\N$, and functionality supplies the value.
\item \textbf{Converse construction.} Conversely, a family of $\Sigma$-partial numerical values has a $\Sigma$-classified graph: conjoin its extent with equality to the specified numeral, using dominance.
Membership of a finite graph is then a finite conjunction of such predicates.
This gives a $\Sigma$-classified ideal and proves \eqref{eq:programpts} uniformly.
\end{enumerate}
\end{proof}

For $g\in\FG$, write $\delta_g(I)$ for $g\in I$. Every observation has the formula
\begin{equation}\label{eq:graphopen}
 U(I)\iff\exists g\,[I(g)\wedge S_U(g)],
 \qquad S_U(g)=U(\downarrow g)\in\Sigma.
\end{equation}
Thus the observations are generated by the finite graph tests $\delta_g$ and $\Sigma$-constants under overt unions and finite meets.

\begin{theorem}[Effective numerical exponential]\label{thm:numericalexp}
Evaluation $\ev:\FunDom\times\N\to\N_\bot$ is effective, and for every effective dcpo $X$ there is a natural bijection
\[
 \ED(\E)(X,\FunDom)
 \cong\ED(\E)(X\times\N,\N_\bot).
\]
\end{theorem}
\begin{proof}
\leavevmode
\begin{enumerate}[label=\textup{(\roman*)},leftmargin=*,itemsep=3pt,topsep=4pt]
\item \textbf{Ideal and directedness calculation.} Evaluation is defined by the functional relation $R_I$.
It is Scott-continuous because evaluation of a directed union is the union of the evaluated graphs.
The inverse image of the observation ``defined with value $m$'' is
\[
 \exists n_0:\N\;[n=n_0\wedge\delta_{\{(n_0,m)\}}(I)],
\]
an overt union of rectangle observations.
For a general positive observation $V^+$ on $\N_\bot$, take the overt union over $m$ with weight $V(m)\in\Sigma$.
Constants and observation closure complete the proof of effectivity.
\item \textbf{Evaluation and transposition.} Given effective $h:X\times\N\to\N_\bot$, define $\lambda h(x)$ to be the ideal of finite graphs compatible with $n\mapsto h(x,n)$.
Its coordinate predicate at $g$ is the finite conjunction
\begin{equation}\label{eq:currying}
 \delta_g(\lambda h(x))
 \iff\bigwedge_{(n,m)\in g}[h(x,n)\text{ is defined with value }m].
\end{equation}
Each factor is an observation of $X$, by effectiveness of $h$ and insertion of the effective numeral.
Finite folding of conjunction is uniform in the finite graph.
\item \textbf{Pullback of finite-graph observations.} Formula \eqref{eq:graphopen} then shows that every observation pulls back to one of $X$, using overtness of $\FG$.
\item \textbf{Ideal and directedness calculation.} For Scott continuity, if a finite graph is contained in the value at a directed supremum, Scott continuity of each coordinate $h(-,n)$ supplies finitely many stages witnessing its entries.
Directedness supplies a common stage.
Hence $\lambda h$ preserves directed joins.
Evaluation and currying are inverse by functionality of the graphs, and all constructions commute with substitution in $X$.
\item \textbf{The evaluation universal property.} Thus the transpose is characterized by the evaluation triangle
\[
\begin{tikzcd}[column sep=large]
 X\times\N\ar[r,"{\lambda h\times\id}"]\ar[dr,"{h}"']&
 \FunDom\times\N\ar[d,"{\ev}"]\\
 &\N_\bot.
\end{tikzcd}
\]
\end{enumerate}
\end{proof}

This supplies the function object needed for numerical computation by a concrete construction. In particular, the program object
\[
 \PC=\Pt(\FunDom)
\]
has a semantic partial evaluator $U_{\mathrm{sem}}:\PC\times\N\pmap\N$ obtained from \cref{thm:effectivelift}. Its domain predicate is jointly $\Sigma$-classified.


\subsection{Typed composition and totality}
The exponential universal property gives a semantic composition operation
\[
 \operatorname{comp}:\FunDom\times\FunDom\longrightarrow\FunDom.
\]
It is obtained by currying sequential evaluation
\[
 (f,g,n)\longmapsto
 \mu_{\N}\bigl(\Lift(\ev_f)(\ev_g(n))\bigr),
\]
where strength supplies the retained parameter $f$. The monad laws give associativity and the total identity function gives the unit. At effective points this is composition of $\Sigma$-partial functions. Thus the operation later implemented by a numerical compiler already exists in the semantic category.

The construction can also be checked by its finite-information matrix: a finite output graph is produced precisely when a finite graph of the inner function and a finite graph of the outer function contain all the required intermediate values. The intermediate-value witnesses form an overt finite-data object. This verifies directly that composition is effective, without assigning a program index to either input function.

\begin{example}[Totality is not an observation]\label{ex:totality}
Let $\operatorname{Tot}(f)$ assert that the $\Omega$-partial function represented by $f\in\FunDom$ is defined at every natural number. For $k\in\N$, let $g_k$ be the finite graph of the identity on the initial segment below $k$. Then $\downarrow g_k$ is an increasing chain with supremum the total identity function. Each $g_k$ is undefined at $k$, so no member of the chain is total. Hence totality is not Scott open and does not belong to $\OO_{\FunDom}$. In the effective topos, \cref{thm:observationfullness} will moreover show that its restriction to $\Pt\FunDom$ is not $\Sigma$-classified. The first conclusion uses only approximation; the latter uses the additional model comparison.
\end{example}


\subsection{The subset domain and its effective points}\label{sec:enumerablesemantic}
The same ideal-completion construction \cite{jungmoshiervickers} applies to the overt poset $\Fin(\N)$ of finite decidable subsets, ordered by inclusion. Put
\[
 \mathcal W=\F\Fin(\N),\qquad W=\Pt\mathcal W=\J\Fin(\N).
\]
Union identifies the underlying semantic domain with $\Omega^\N$. An ideal is recovered as the family of finite subsets of its union. By \cref{jd:prop:basis}, the effective-point subobject is
\[
 W\cong\Sigma^\N.
\]
Indeed, membership of a finite set is a finite conjunction of its element-membership predicates; conversely singleton tests recover the latter. Observations are the upper positive traces on finite subsets. Thus both $P=\PC$ and $W$ are obtained by the same canonical doctrine, with different finite data as their bases.


The finite-graph presentation provides a concrete partial-function domain and a semantic evaluator with a jointly positive domain on effective points. It also shows why numerical indices are implementation data added later, not part of the definition of the semantic function space.

\section{Points, observations, and an intrinsic extension theorem}\label{sec:intrinsic}

An internal map between effective-point objects need not extend to an observation-preserving map of their full domains. The missing datum is its action on tests. Representing points and observations as hom-objects makes this condition a transpose square, whose solution can be reconstructed from a positive approximation matrix.


\subsection{The two representable functors}
For effective dcpos $D,E$, the conditions that a map $D\to E$ preserve ideal suprema and pull observations back to observations define a subobject
\[
 \underline{\C}(D,E)\incl E^D.
\]
The conditions are internal and hence include uniformity in a parameter. Identity and composition restrict to these subobjects. They make $\C$ an $\E$-enriched category in the specified sense of having hom-objects in $\E$. The order on the hom-objects is still the pointwise order inherited from the underlying dcpos.

The computational algebra category $\AS$ is cartesian closed by \cref{jd:thm:ccc}. The larger observation-equipped semantic category $\C=\ED$ has the hom-objects just defined; their existence alone does not supply semantic exponentials. The typed construction of \cref{thm:numericalexp} supplies the required semantic exponential and its observation calculation. On $\CS$, the comparison identifies effective semantic maps with algebra morphisms of computational points.

Put $\Sier=1_\bot$. Its underlying dcpo is $\Omega$, ordered by implication. Its observation structure is generated by the identity predicate $p\mapsto p$, together with the $\Sigma$-constants. In particular, $\Pt\Sier\cong\Sigma$ by \cref{thm:effectivelift}.

\begin{example}[The observation structure is genuine input]
\label{ex:observation-structure-input}
The same underlying dcpo $\Omega$ carries different observation structures, and the naturality square of \cref{prop:pointsfunctor} is not generally cartesian.
\end{example}
\begin{proof}
\leavevmode
\begin{enumerate}[label=\textup{(\roman*)},leftmargin=*,itemsep=3pt,topsep=4pt]
\item \textbf{Two constructed structures.} On $\Omega$ take only the constant predicates with values in $\Sigma$. Finite meets and overt unions remain constant, with values given by the corresponding operations in $\Sigma$, so this already is an observation structure. Call the resulting effective dcpo $D_0$. Its evaluation is everywhere $\Sigma$-valued, hence $\Pt D_0=\Omega$. In contrast, the structure on $\Sier$ contains the identity predicate and has $\Pt\Sier=\Sigma$. The identity predicate is not constant, as evaluation at bottom and top shows. Thus the two selected structures are distinct.
\item \textbf{The failed unrestricted pullback.} For the terminal map $\Sier\to1$, the lower pullback in the naturality square is $\Omega\times_1\Pt1\cong\Omega$, whereas its upper-left object is $\Sigma$. The inclusion $\Sigma\incl\Omega$ is not an isomorphism in $\Eff$: otherwise the subassembly of nonhalting indices in $\N$ would have a $\Sigma$-classifier. Lifting that classifier to binary streams on the partitioned natural-number assembly by \cref{lem:modest-target-lift} would enumerate the complement of the halting set. Together with the usual enumeration of the halting set this would decide halting, contradicting the diagonal argument. Therefore the naturality square is not a pullback. \Cref{prop:initial-observation-points} supplies exactly the pullback comparison for the observation structure it constructs.
\end{enumerate}
\end{proof}

\begin{proposition}[Representability of points and observations]\label{prop:representable}
There are natural isomorphisms in $\E$
\[
 \Pt D\cong\underline{\C}(1,D),\qquad
 \OO_D\cong\underline{\C}(D,\Sier).
\]
Under these identifications the observation pairing is enriched composition.
\end{proposition}
\begin{proof}
\leavevmode
\begin{enumerate}[label=\textup{(\roman*)},leftmargin=*,itemsep=3pt,topsep=4pt]
\item \textbf{Representability of points and observations.} A parameterized map $1\to D$ is effective exactly when its evaluation on all observations factors through $\Sigma$, which is the defining pullback of $\Pt D$. A Scott-continuous map $U:D\to\Omega$ is a Scott open.
If it is effective as a map to $\Sier$, pulling back the identity observation shows $U\in\OO_D$.
\item \textbf{Converse construction.} Conversely, if $U\in\OO_D$, the inverse images of the generators of $\OO_{\Sier}$ are $U$ and the $\Sigma$-constants.
The observation-closure lemma proves effectivity.
These arguments are internal and commute with parameters.
Composition evaluates $U$ at a point, so its value is precisely $\operatorname{Obs}_D$.

The defining maps satisfy the following commutative diagram
\[
\begin{tikzcd}[column sep=large,row sep=large]
\OO_D\times\Pt D\ar[r,"\operatorname{Obs}_D"]\ar[d,"\cong"']&\Sigma\ar[d,"\id"]\\ \underline{\C}(D,\Sier)\times\underline{\C}(1,D)\ar[r,"\circ"']&\Sigma
\end{tikzcd}
\]
\end{enumerate}
\end{proof}

Thus an admissible semidecision on $D$ is simply a morphism $D\to\Sier$. Define
\[
 R_D:\OO_D\longrightarrow\Sigma^{\Pt D}
\]
by restriction to effective points. Naturality is expressed by the square
\begin{equation}\label{eq:obsnaturality}
\begin{tikzcd}[column sep=large]
 \OO_E\ar[r,"{F^*}"]\ar[d,"{R_E}"']&\OO_D\ar[d,"{R_D}"]\\
 \Sigma^{\Pt E}\ar[r,"{\Sigma^{\Pt F}}"']&\Sigma^{\Pt D}.
\end{tikzcd}
\end{equation}
The problem of extending a map of effective points is therefore a problem of lifting its inverse-image operation through $R_D$ and $R_E$.

\subsection{Separation and the Chu representation}
For a searchable presentation $(B,r)$, the basic-point map
\[
 \beta_D:B\longrightarrow\Pt D,\qquad b\longmapsto r(b)
\]
is defined because the approximation relation is $\Sigma$-valued. The coordinate observations $\delta_b(I)=I(b)$ separate points of the underlying rounded-ideal domain. Every observation is determined by its composite with $\beta_D$ by \cref{prop:roundedtrace}. Consequently $R_D$ is monic. This is separation by a family of generalized effective points; it does not require enough global points of $\E$.

An object of $\Chu(\E,\Sigma)$ is a pairing $P\times O\to\Sigma$. A morphism from $(P,O,e)$ to $(Q,V,d)$ consists of $f:P\to Q$ and $g:V\to O$ satisfying
\[
 d(f(p),v)=e(p,g(v)).
\]
Equivalently, the following diagram commutes:
\[
\begin{tikzcd}[column sep=large]
 P\times V\ar[r,"{f\times\id}"]\ar[d,"{\id\times g}"']&
 Q\times V\ar[d,"{d}"]\\
 P\times O\ar[r,"{e}"']&\Sigma.
\end{tikzcd}
\]
This internal Chu construction follows \cite{pratt}. Here it is applied to internal effective points and observations.

\begin{theorem}[Intrinsic extension and full Chu representation]\label{thm:chu}
Let $D,E$ be searchable effective domains. A morphism $f:\Pt D\to\Pt E$ in $\E$ extends to an effective Scott-continuous map $F:D\to E$ if and only if there is a morphism $g:\OO_E\to\OO_D$ making
\begin{equation}\label{eq:extensioncriterion}
\begin{tikzcd}[column sep=large]
 \OO_E\ar[r,dashed,"{g}"]\ar[d,"{R_E}"']&\OO_D\ar[d,hook,"{R_D}"]\\
 \Sigma^{\Pt E}\ar[r,"{\Sigma^f}"']&\Sigma^{\Pt D}
\end{tikzcd}
\end{equation}
commute. Both $F$ and $g$ are unique, and $g=F^*$. Hence
\[
 D\longmapsto(\Pt D,\OO_D,\operatorname{Obs}_D)
\]
is a fully faithful $\E$-enriched functor from $\Search(\E)$ to $\Chu(\E,\Sigma)$.
\end{theorem}
\begin{proof}
\leavevmode
\begin{enumerate}[label=\textup{(\roman*)},leftmargin=*,itemsep=3pt,topsep=4pt]
\item \textbf{Construction of the matrix.} Necessity is \eqref{eq:obsnaturality}.
For sufficiency choose searchable presentations $(B,r)$ of $D$ and $(C,s)$ of $E$.
Define the $\Sigma$-valued matrix
\begin{equation}\label{eq:chumatrix}
 m(b,c)=\operatorname{Obs}_E(f\beta_D(b),\delta_c).
\end{equation}
For each $b$, this is the rounded ideal underlying the effective point $f\beta_D(b)$.
Thus its rows are ideals and are target-rounded.
By the commuting square it is also the trace of the observation $g(\delta_c)$ on $D$.
\item \textbf{Source and target roundedness.} Such a trace is upper and source-rounded:
\[
 m(b,c)\ \Longleftrightarrow\
 \exists a\,[\rho(a,b)\wedge m(a,c)].
\]
The row condition and this identity give $m=s\star m\star r$.
In particular $m$ is monotone in its source variable, so it is the rounded distributor required by \cref{thm:relations}.
\item \textbf{Extension and restriction.} That theorem constructs an effective map
\[
 F(I)(c)\iff\exists b\,[I(b)\wedge m(b,c)].
\]
For an effective point $I$, Scott openness of $g(\delta_c)$ and the reconstruction $I=\bigvee_{b\in I}r(b)$ give
\[
 F(I)(c)=g(\delta_c)(I)=\delta_c(f(I)).
\]
The coordinate observations separate $E$, so $\Pt F=f$.
\item \textbf{Uniqueness and enriched naturality.} An effective map with this restriction must agree with $F$ on every basic point $r(b)$ and hence on every directed reconstruction of an element of $D$.
This proves uniqueness of $F$.
Since $R_D$ is monic, \eqref{eq:extensioncriterion} also gives uniqueness of $g$, which must equal $F^*$.
All constructions can be performed in an arbitrary parameter context.
They therefore identify the internal hom-objects, proving enriched full faithfulness.
\end{enumerate}
\end{proof}

No algebraic preservation conditions on $g$ have to be imposed separately: the theorem identifies it with inverse image, so preservation of the available observation operations follows. The theorem also isolates the exact content of automatic continuity. If $R_D$ is an isomorphism, every map $\Pt D\to\Pt E$ has the required transpose and extends uniquely. This stronger hypothesis will be proved for the numerical function domain in $\Eff$.

\begin{remark}[The role of the observation structure]\label{rem:notautomatic}
For the discrete object $2$ with only $\Sigma$-constant observations from \cref{ex:pointforgetting}, the map $R_D:\OO_D\to\Sigma^{\Pt D}$ misses the predicate $b\mapsto[b=1]$. Thus observations cannot be recovered from effective points for arbitrary observed domains. On the searchable domains the extension theorem gives its precise observation-transpose criterion; the stronger fullness assertion for the numerical function domain is proved in \cref{thm:observationfullness}. These are comparisons for the specified observation structures.
\end{remark}


The extension criterion says that a map of effective points extends exactly when it transports the selected tests. On searchable domains this condition is expressed by a positive rounded distributor; it does not follow merely from the existence of an underlying set-theoretic function.

\section{Finite positive information and observable decisions}\label{sec:syntheticfinite}

The intrinsic extension theorem makes tests a way to recover semantic maps. We now identify the content of those tests: on a searchable domain every positive observation has a basic approximation witness. Combined with pointedness, this also limits which distinctions can be decided by an effective Boolean-valued map.


The finite-information statement below is internal: its witnesses are basic approximants in the topos, not program indices. The numerical interpretation of these witnesses will be established separately.

\begin{theorem}[Intrinsic finite-information theorem]\label{thm:intrinsicfinite}
For a searchable presentation $(B,r)$ of $D$, restriction to the basic-point family gives an isomorphism
\[
 \OO_D\cong\RO(B,r),
\]
where $\RO(B,r)$ is the object of upper predicates $S:B\to\Sigma$ satisfying
\[
 S(b)\iff\exists a\,[\rho(a,b)\wedge S(a)].
\]
The inverse sends $S$ to the observable predicate
\begin{equation}\label{eq:intrinsicfinite}
 U_S(I)\iff\exists b\,[I(b)\wedge S(b)].
\end{equation}
For every effective point $I$, the right side is a $\Sigma$-classified existential over an overt object.
\end{theorem}
\begin{proof}
\leavevmode
\begin{enumerate}[label=\textup{(\roman*)},leftmargin=*,itemsep=3pt,topsep=4pt]
\item \textbf{Intrinsic finite-information theorem.} This is the trace isomorphism of \cref{prop:roundedtrace}, now interpreted through \cref{prop:representable}.
Roundedness is restriction of Scott openness along $r(b)=\bigvee_{a\in r(b)}r(a)$.
\item \textbf{Converse construction.} Conversely the displayed existential is a Scott open, has trace $S$, and lies in $\OO_D$.
At an effective point $I$, both $I(b)$ and $S(b)$ are $\Sigma$-valued, and overtness of $B$ supplies the restricted existential.
The two inverse equations hold in every context.
\end{enumerate}
\end{proof}

Here ``finite information'' means a single basic approximation witnessing a positive observation. A basic approximation need not be a compact element or a literal finite set. The rational-cut example in \cref{ex:lowerreals} is continuous without being algebraic, and still satisfies the theorem. For finite-graph domains the basic approximations are literal finite graphs, recovering the finite-graph formulation.

The result can be expressed entirely by regular images. For an observation $U$, form its witness object
\[
 W_U=\{(I,b)\in\Pt D\times B\mid I(b)\wedge U(r(b))\}.
\]
The composite $W_U\to\Pt D$ has image exactly the classified subobject $R_D(U)$. Thus the finite-information statement asserts equality of two subobjects, with witnesses retained in the left-hand presentation:
\[
\begin{tikzcd}[column sep=large]
 W_U\ar[r,two heads]\ar[dr]&R_D(U)\ar[d,hook]\\
 &\Pt D.
\end{tikzcd}
\]
The image equation requires no section of the epimorphism. In realizability, projective extraction supplies explicit witness transformations after choosing a cover.

\begin{theorem}[Abstract Rice theorem]\label{thm:syntheticrice}
Let $D$ be an effective dcpo with an effective least element. If $K$ is an effective dcpo whose underlying order is discrete, precomposition with $!:D\to1$ gives an isomorphism
\[
 \underline{\C}(1,K)\cong\underline{\C}(D,K).
\]
In particular, every effective decision $D\to2$, with $2$ discrete and carrying observations $\Sigma^2$, is constant.
\end{theorem}
\begin{proof}
\leavevmode
\begin{enumerate}[label=\textup{(\roman*)},leftmargin=*,itemsep=3pt,topsep=4pt]
\item \textbf{Abstract Rice theorem.} The inequality $\bot_D!\leq\id_D$ gives, for every effective $d:D\to K$,
\[
 d\bot_D!\leq d.
\]
Discreteness of $K$ turns this inequality into equality.
Thus $d$ factors uniquely through $1$, with value $d\bot_D$.
Both assignments are effective and internal, so they give the asserted hom-object isomorphism.
\end{enumerate}
\end{proof}

For an observation $U\in\OO_D$, an effective decision means a complementary observation $V$ with $U\vee V=\top$ and $U\wedge V=\bot$. The corresponding partition defines $D\to2$; its inverse images of arbitrary $\Sigma$-valued predicates on $2$ are finite unions of the weighted observations $U,V$. The theorem therefore says that no nonconstant observational partition exists on a pointed domain. In $\Eff$, definedness on $D_\bot$ cannot have an effective Boolean decision whenever $D$ has an effective global point: it distinguishes bottom from that point's total lift. Definedness itself is nevertheless an observation.

This argument is independent of universal evaluation and self-reference. In $\Eff$, \cref{thm:intrinsichalting} proves constancy for every ambient decision on the function object, and \cref{thm:observationfullness} identifies all its semidecisions with observations. These extensional results precede the implementation and index-set versions.


Positive observations of searchable values are therefore witnessed by basic approximants, whereas effective Boolean decisions on pointed domains cannot distinguish those values nontrivially. The partial predicate-transformer calculus retains the termination observation without turning it into such a decision.

\section{Partial computation and positive predicate transformers}\label{sec:transformers}

For a partial computation, observing a result necessarily includes observing that the computation terminates. The positive-extension construction packages both facts into an inverse-image operator on observations. We check its contravariant composition and parameter laws before using it to interpret recursion and feedback.


Write $\Lift D=D_\bot$. A partial computation from $D$ to $E$ is a Kleisli arrow $f:D\to\Lift E$. By \cref{thm:effectivelift} it induces a $\Sigma$-partial map $\Pt D\pmap\Pt E$. For each observation $U$ of $E$, its positive extension $U^+$ observes both termination and the resulting value. Define
\[
 f^\diamond:\OO_E\longrightarrow\OO_D,
 \qquad f^\diamond(U)=f^*(U^+).
\]
In particular $f^\diamond(\top)$ is the domain observation of $f$. This formulation records definedness even when the partial value cannot be returned.

\subsection{Positive tests on partial values}\label{app:partial}
Let $V:Y\to\Sigma$ be a positive output predicate. Its test on a partial element is
\[
 V^+(p,u)=\exists a:[p]\;V(u(a)).
\]
The classifier exists by dominance. For $V=\top$, this is precisely the extent $p$. Thus output observation includes termination as part of its positive meaning.

\begin{proposition}[Substitution of partial observations]\label{prop:partial-observation}
For a $\Sigma$-partial map $f:X\pmap Y$, the predicate that it returns a value in $V$ is positive and has classifier
\[
 x\longmapsto\exists a:[p(x)]\;V(u_x(a)).
\]
For a composite $gf$, its output test is the substitution into $f$ of the output test of $g$, with exactly the composite extent \eqref{eq:partial-composite-extent}.
\end{proposition}
\begin{proof}
\leavevmode
\begin{enumerate}[label=\textup{(\roman*)},leftmargin=*,itemsep=3pt,topsep=4pt]
\item \textbf{Coherence equations.} The first formula is dependent closure of $\Sigma$.
For the composite, substitution gives
\[
 \exists a:[p(x)]\ \exists b:[q(u_x(a))]\;
 V(v_{u_x(a)}(b)).
\]
This is the output observation of the value of $gf$ on its displayed extent.
Reassociating the two existential witnesses identifies it with the direct lifted test.
For top output observation it reduces to the composite domain, so the equation includes definedness.
\end{enumerate}
\end{proof}

For semantic domains, a full element of $Y_\bot$ can have a general truth-valued extent. The positive lifted Scott open is constructed by its basis trace, as in \cref{lem:kernel-effective}; it need not evaluate positively on every full point. The effective-point comparison $\Pt(Y_\bot)\cong\LL\Pt Y$ is exactly the restriction where both its extent and its positive observations have the required computational type. This is why a forcing proof retains the full kernel and separately constructs its effective stable trace.


\begin{proposition}[Contravariant calculus of computation]\label{prop:transformers}
For Kleisli arrows $f:D\to\Lift E$ and $g:E\to\Lift F$,
\[
 (g\star f)^\diamond=f^\diamond g^\diamond,
 \qquad (\eta_D)^\diamond=\id_{\OO_D}.
\]
Each $f^\diamond$ preserves binary meets and overt-indexed joins. It preserves the top observation precisely when $f$ has everywhere true extent. If $D,E$ are searchable, $f\mapsto f^\diamond$ is faithful.
\end{proposition}
\begin{proof}
\leavevmode
\begin{enumerate}[label=\textup{(\roman*)},leftmargin=*,itemsep=3pt,topsep=4pt]
\item \textbf{Naturality and reindexing.} The positive-extension identities for unit and multiplication are
$\eta^*(U^+)=U$ and $\mu^*(U^+)=(U^+)^+$.
Together with naturality of positive extension, they give the two composition equations.
Two positive tests on a partial value hold exactly when its common extent holds and both tests hold of the value; this proves preservation of binary meets.
Existential interchange proves preservation of overt joins.
The assertion about top is the identity $f^\diamond(\top)=\defin f$.
\item \textbf{Observation preservation.} If $f^\diamond=g^\diamond$, the two maps into $\Lift E$ have equal inverse images of every positive generator.
They also agree on constant generators, and therefore have identical inverse images on $\OO_{\Lift E}$.
The searchable lift has separating coordinate observations by \cref{thm:effectivelift,prop:observationseparation}; consequently $f=g$.
\end{enumerate}
\end{proof}

Thus sequential computation has a contravariant observable semantics in positive predicate transformers. The loss of top preservation records partiality; no additional symbol for divergence is needed at the level of observations. Together with the Chu representation, this identifies a computation by its action on values and tests while allowing presentations of both to be supplied later.

\subsection{Contextuality and substitution}
The constructions above are stable under a change of parameter context. For $u:T\to S$, the pullback functor between slices is logical: it preserves finite limits, the slice subobject classifier, power objects, and the natural numbers object. It also preserves images. It therefore preserves the finite-hit factorization defining the slice Rosolini object and the bounded internal formulas defining ideals, observation structures, and their closure.

\begin{proposition}[Indexed semantics]\label{prop:indexed}
The completion, effective-point, observation, and lifting constructions give indexed categories and indexed functors over $\E$. There are canonical coherent isomorphisms
\[
\begin{gathered}
 u^*\mathsf J_{\Sigma,S}(P)\cong \mathsf J_{\Sigma,T}(u^*P),\qquad u^*\Idl_S(D)\cong\Idl_T(u^*D),\\
 u^*\operatorname{Pt}_{\Sigma,S} D\cong\operatorname{Pt}_{\Sigma,T}(u^*D),\qquad
 u^*\Lift_S D\cong\Lift_T(u^*D).
\end{gathered}
\]
The extension theorem, positive predicate transformers, and their comparison equations commute with these isomorphisms.
\end{proposition}
\begin{proof}
\leavevmode
\begin{enumerate}[label=\textup{(\roman*)},leftmargin=*,itemsep=3pt,topsep=4pt]
\item \textbf{Pullback construction.} These objects are defined by finite-limit, power-object, image, and internal bounded-predicate constructions. A logical pullback functor preserves their defining universal properties.
The least observation closure is the internal intersection of all observation structures containing its generators, so its bounded power-object definition is also preserved.
The resulting comparison isomorphisms commute with evaluation, unit, and multiplication by uniqueness.
Applying the same argument to the formulas for $m$ and $F$ in \cref{thm:chu} proves the extension comparison.
Identity and composition coherence follows from that of slice pullback.
\end{enumerate}
\end{proof}

For example, the effective-point comparison is the Beck--Chevalley square
\[
\begin{tikzcd}[column sep=large]
 \ED(\E/S)\ar[r,"{\operatorname{Pt}_{\Sigma,S}}"]\ar[d,"{u^*}"']&\E/S\ar[d,"{u^*}"]\\
 \ED(\E/T)\ar[r,"{\operatorname{Pt}_{\Sigma,T}}"']&\E/T,
\end{tikzcd}
\]
commuting up to the displayed canonical isomorphism. This indexed comparison concerns logical reindexing between slices of the fixed topos, rather than geometric inverse image between different topoi.


A positive transformer identifies what a partial computation can make observable, including whether it returns. Its substitution laws prepare the fixed-point constructions, where a successful observation will be witnessed by a finite iterate or by a finite returning execution path.

\section{Recursion, guarded execution, and interaction}\label{sec:intrinsicrecursion}

Several constructions produce behavior defined by recursion, but they solve different equations: least fixed points of domain maps, finite returns of partial transitions, and unique guarded solutions. We keep their interfaces distinct, then compare their finite witnesses. This makes the eventual feedback trace an equality of exact execution domains as well as of returned values.


The natural numbers object of $\Eff$ is overt by \cref{thm:relativefoundation}. Several operations in this part have different universal properties. We fix their notation before constructing them.
\begin{center}
\begin{tabular}{L{.19\linewidth}L{.32\linewidth}L{.38\linewidth}}
\toprule
Notation & Input & Selected solution or feedback\\
\midrule
$F^\dagger$ & $F:X\times D\to D$ & least fixed point in the domain order\\
$f^\ddagger$ & $f:X\to\LL(Y+X)$ & first finite return, with its exact domain\\
$\Tr_\times$ & cartesian feedback arrow & trace induced by Conway iteration\\
$\Tr_+$ & $f:X+U\to\LL(Y+U)$ & finite execution through interface $U$\\
$h^\sharp$ & $h:Z\times\triangleright A\to A$ & unique guarded solution in the temporal category\\
$\mathsf{rec}$ & numerical transformer index $t\in K_1$ & delayed self-reference of \cref{rel:thm:recursion}\\
\bottomrule
\end{tabular}
\end{center}
The first two are related by the partiality algebra, and finite completion relates the second to guarded execution. The comparisons are theorems below; the notation does not identify their different inputs or their different uniqueness properties.

A \textbf{pointed effective dcpo} has a least element which is an effective point. Given an effective map $F:X\times D\to D$ with $D$ pointed, define
\[
 a_0(x)=\bot_D,\qquad a_{n+1}(x)=F(x,a_n(x)),
 \qquad F^\dagger(x)=\bigvee_{n\in\N}a_n(x).
\]

\begin{theorem}\label{thm:lfp}
The map $F^\dagger:X\to D$ is effective and Scott-continuous. For each $x$ it is the least fixed point of $F(x,-)$, and for $U\in\OO_D$,
\begin{equation}\label{eq:lfpobs}
 U(F^\dagger(x))\iff\exists n\;U(a_n(x)).
\end{equation}
\end{theorem}
\begin{proof}
\leavevmode
\begin{enumerate}[label=\textup{(\roman*)},leftmargin=*,itemsep=3pt,topsep=4pt]
\item \textbf{Scott continuity.} Monotonicity makes the iteration increasing.
Scott continuity of $F(x,-)$ yields
\[
 F(x,F^\dagger(x))=\bigvee_n a_{n+1}(x)=F^\dagger(x).
\]
Induction bounds every iterate by any pre-fixed point, proving leastness.
Each $a_n$ is effective by induction, pairing, and composition; these inverse-image operations are uniform in $n$ by recursion.
Their observation preimages form an overt-indexed family since $\N$ is overt.
Equation \eqref{eq:lfpobs} therefore proves effectivity.
\item \textbf{Coherence equations.} Each iterate is Scott-continuous in $x$, and directed suprema in $x$ commute with the increasing supremum over $n$, proving Scott continuity of $F^\dagger$.
\end{enumerate}
\end{proof}


\begin{proposition}[Parameters and positive witnesses for recursion]\label{prop:lfpparameters}
For an effective map $h:Y\to X$, one has
\[
 \bigl(F\circ(h\times\id_D)\bigr)^\dagger=F^\dagger h.
\]
If $F\leq G$ pointwise, then $F^\dagger\leq G^\dagger$. In $\Eff$, under projective presentations of effective inputs and observations, a positive observation of $F^\dagger(x)$ is represented by a finite iteration number and a proof of the observation at that iterate.
\end{proposition}
\begin{proof}
\leavevmode
\begin{enumerate}[label=\textup{(\roman*)},leftmargin=*,itemsep=3pt,topsep=4pt]
\item \textbf{Parameters and positive witnesses for recursion.} The zeroth iterates commute with $h$ because they are constant bottom.
If the $n$th iterates commute, substitution into the recursion clause proves the same for the $(n+1)$st.
Taking their supremum proves the parameter equation.
For monotonicity, induction gives $a_n^F(x)\leq a_n^G(x)$: use the inductive inequality, monotonicity of $F$, and $F\leq G$.
Taking joins gives the claimed inequality.
\item \textbf{Parameter substitution.} For the witness assertion, work in $\Eff$ and restrict $x$ to $\Pt(X)$.
Cover $\Pt(X)\times\OO_D$ by a projective parameter object.
Evaluation of $U(a_n(x))$ is then $\Sigma$-valued, uniformly in $n$, since the iterates preserve effective points.
Lift this predicate to streams on the product of the parameter cover with $\N$.
Dovetail the stream index and $n$.
By \eqref{eq:lfpobs}, the resulting stream classifies $U(F^\dagger(x))$. A hit decodes to the selected iterate and a hit at that iterate, and a witness at any iterate encodes to a hit of the dovetailed stream.
\item \textbf{Scott continuity.} These are the split witness maps for the existential formula.
Thus the Scott continuity proof also determines the finite evidence used by the realizability presentation.
\end{enumerate}
\end{proof}


\subsection{Conway identities and simultaneous recursion}
Leastness gives more than individual solutions. It gives an algebra of constructions by recursion, including the compatibility needed to compose and substitute definitions by recursion. We state this on the full subcategory $\C_\bot$ of effectively pointed dcpos; its arrows need not preserve bottom.

\begin{theorem}[Conway recursion and trace]\label{thm:conway}
When $\N$ is overt, the assignment $F\mapsto F^\dagger$ of \cref{thm:lfp} is a Conway operator on $\C_\bot$. It satisfies the fixed-point and parameter equations, dinaturality, and the double-dagger equation. Consequently $\C_\bot$, with its cartesian product, is a traced category. For $h:X\times D\to Y\times D$, the trace is
\begin{equation}\label{eq:trace}
 \Tr^D_{X,Y}(h)
 =\pi_Yh\circ
 \left\langle\id_X,(\pi_Dh)^\dagger\right\rangle.
\end{equation}
These constructions commute with substitution in the ambient topos.
\end{theorem}
\begin{proof}
\leavevmode
\begin{enumerate}[label=\textup{(\roman*)},leftmargin=*,itemsep=3pt,topsep=4pt]
\item \textbf{Unit and identity laws.} The fixed-point equation and parameter identity are \cref{thm:lfp,prop:lfpparameters}.
For dinaturality, take effective maps $f:X\times D\to E$ and $g:X\times E\to D$.
Suppressing $x$, write $a=\mu d\,g(f(d))$ and $b=\mu e\,f(g(e))$.
The element $g(b)$ is a fixed point of $gf$, so $a\leq g(b)$.
Similarly $b\leq f(a)$, whence $g(b)\leq gf(a)=a$.
Thus
\begin{equation}\label{eq:dinatural}
 \bigl(g\circ\langle\pi_X,f\rangle\bigr)^\dagger
 =g\circ\left\langle\id_X,
 \bigl(f\circ\langle\pi_X,g\rangle\bigr)^\dagger\right\rangle.
\end{equation}
\item \textbf{Double-dagger equation.} For $h:X\times D\times D\to D$, let
$a=\mu d\,\mu e\,h(x,d,e)$ and $c=\mu d\,h(x,d,d)$.
The fixed-point equations imply $h(x,a,a)=a$, hence $c\leq a$.
Since $h(x,c,c)=c$, the inner least fixed point $\mu e\,h(x,c,e)$ is at most $c$.
Thus $c$ is pre-fixed for the outer operator, giving $a\leq c$.
This proves
\[
 (h^\dagger)^\dagger
 =\bigl(h\circ(\id_X\times\Delta_D)\bigr)^\dagger.
\]
All arrows involved are effective by \cref{thm:lfp}; the equations hold in the category $\C_\bot$ of effective pointed dcpos.
\item \textbf{Inverse comparison.} The equivalence between Conway operators and cartesian traces gives \eqref{eq:trace} and the trace identities \cite{hasegawa}; we use the traced-category terminology of \cite{jsv}.
Products of effectively pointed domains are effectively pointed, so this is a trace on the stated cartesian category.
The indexed assertion follows either from the iterates or from \cref{prop:indexed} and the defining leastness property.
\end{enumerate}
\end{proof}

The fixed-point equation itself is a commutative diagram in $\C$:
\[
\begin{tikzcd}[column sep=large]
 X\ar[r,"{\langle\id_X,F^\dagger\rangle}"]\ar[dr,"{F^\dagger}"']&
 X\times D\ar[d,"{F}"]\\
 &D.
\end{tikzcd}
\]
Leastness specifies this solution among the arrows satisfying the diagram. It is not replaced by the diagram alone.

\begin{proposition}[Beki\'c decomposition]\label{prop:bekic}
For effective $f:X\times D\times E\to D$ and $g:X\times D\times E\to E$, put
\[
 b(x,d)=\mu e\,g(x,d,e),\qquad
 a(x)=\mu d\,f(x,d,b(x,d)).
\]
Then the least simultaneous solution is
\[
 \langle f,g\rangle^\dagger(x)
 =\bigl(a(x),b(x,a(x))\bigr).
\]
\end{proposition}
\begin{proof}
\leavevmode
\begin{enumerate}[label=\textup{(\roman*)},leftmargin=*,itemsep=3pt,topsep=4pt]
\item \textbf{Bekic decomposition.} Both components satisfy their equations by construction.
If $(d,e)$ is a simultaneous pre-fixed point, then $b(x,d)\leq e$ by inner leastness.
Monotonicity gives $f(x,d,b(x,d))\leq d$, so $a(x)\leq d$.
Monotonicity of the parameterized least fixed point then gives $b(x,a(x))\leq b(x,d)\leq e$.
This proves leastness.
Every component is effective by closure under composition, products, and \cref{thm:lfp}.
\end{enumerate}
\end{proof}

\begin{proposition}[Strict simulation]\label{prop:strictsimulation}
Let $k:D\to E$ be effective, Scott-continuous, and strict. If $F:X\times D\to D$ and $G:X\times E\to E$ satisfy
\[
\begin{tikzcd}[column sep=large]
 X\times D\ar[r,"{F}"]\ar[d,"{\id_X\times k}"']&D\ar[d,"{k}"]\\
 X\times E\ar[r,"{G}"']&E,
\end{tikzcd}
\]
then $kF^\dagger=G^\dagger$.
\end{proposition}
\begin{proof}
\leavevmode
\begin{enumerate}[label=\textup{(\roman*)},leftmargin=*,itemsep=3pt,topsep=4pt]
\item \textbf{Strict simulation.} Strictness identifies the zeroth iterates, and the square identifies each successor iterate.
Scott continuity of $k$ then identifies their joins.
Without strictness the zeroth-iterate argument is unavailable, so the conclusion is asserted only for the stated class of simulations.

The defining arrows form the commutative diagram
\[
\begin{tikzcd}[column sep=large,row sep=large]
X\times D\ar[r,"F"]\ar[d,"\id_X\times k"']&D\ar[d,"k"]\\ X\times E\ar[r,"G"']&E
\end{tikzcd}
\]
\end{enumerate}
\end{proof}

The construction uses $\N$ to form a chain of iterates. It neither indexes all computations by $\N$ nor chooses names for the resulting fixed points. The later coding comparison lifts this already defined operation to presentations and thereby specifies which names denote the least semantic solution.


\subsection{Finite sums and the intrinsic solution object}
The execution construction uses the algebra category $\AS$ itself. Finite sums in this category are inherited from disjoint sums of internal posets.

\begin{lemma}[Disjoint sums of computational algebras]\label{exec:lem:sums}
There are natural isomorphisms
\[
 \J0\cong0,\qquad \J(P+Q)\cong\J P+\J Q.
\]
They commute with unit and multiplication. Consequently $\AS$ has finite coproducts created by its forgetful functor to internal posets.
\end{lemma}
\begin{proof}
\leavevmode
\begin{enumerate}[label=\textup{(\roman*)},leftmargin=*,itemsep=3pt,topsep=4pt]
\item \textbf{Disjoint sums of computational algebras.} An inhabited directed ideal in $P+Q$ cannot meet both summands: a common upper bound would have to belong to both.
Inhabitedness gives locally one summand, and uniqueness of that summand makes the choice descend.
Its restriction there is an ideal.
If it is admitted in $P+Q$, extend a positive upper predicate on the selected summand by false on the other; its pairing proves admission in the summand.
\item \textbf{Converse construction.} Conversely every positive upper predicate on $P+Q$ restricts to one on the selected summand, so admission there proves admission in the sum.
The constructions are inverse.
No ideal of the empty poset is inhabited.
Principal ideals and unions respect the selected summand, proving the monad equations and the coproduct assertion.
\end{enumerate}
\end{proof}

\begin{proposition}\label{prop:intrinsiclpf}
In the recursion setting of this section, every pointed $\J$-algebra $A$ has an algebra morphism
\[
 \lfp_A:[A,A]_\Sigma\longrightarrow A,
 \qquad
 \lfp_A(f)=\bigvee_{n:\N}f^n(\bot).
\]
It selects the least solution, is natural under strict intertwining maps, and yields parameterized Conway iteration. Under the semantic comparison it is the operator of \cref{thm:lfp,thm:conway}.
\end{proposition}
\begin{proof}
\leavevmode
\begin{enumerate}[label=\textup{(\roman*)},leftmargin=*,itemsep=3pt,topsep=4pt]
\item \textbf{Ideal and directedness calculation.} Overtness of $\N$ is available in the current setting, and is proved in $\Eff$ in \cref{thm:relativefoundation}.
The lower ideal generated by an increasing sequence is admitted because its pairing with $u$ is $\exists n\,u(x_n)$.
Thus the displayed supremum exists, and preservation of admitted suprema proves the fixed-point equation.
Induction bounds every finite iterate by any pre-fixed point, proving leastness.
\item \textbf{Observation preservation.} To prove the algebra equation for the operator, first note that $f\mapsto\lfp_A(f)$ is monotone by induction on its iterates.
Thus, for $I\in\J[A,A]_\Sigma$, its lower image under this map is already an admitted ideal: admission of lower images uses only monotonicity, not the algebra equation now being proved.
Put
\[
 g=\bigvee I,\qquad y=\bigvee_{h\in I}\lfp_A(h).
\]
For $f,h\in I$, directedness supplies $k\in I$ above both.
\item Therefore
\[
 f(\lfp_A(h))\leq k(\lfp_A(k))=\lfp_A(k)\leq y.
\]
Since $f$ preserves the admitted supremum defining $y$, this gives $f(y)\leq y$.
The supremum $g$ is pointwise by \cref{jd:thm:ccc}, so $g(y)\leq y$.
Leastness gives $\lfp_A(g)\leq y$.
\item \textbf{Converse construction.} Conversely $h\leq g$ gives $\lfp_A(h)\leq\lfp_A(g)$ for all $h\in I$, hence $y\leq\lfp_A(g)$.
This is precisely the homomorphism equation for $\lfp_A$.
\item \textbf{Scott continuity.} Strict intertwining follows by induction and preservation of the increasing supremum.
The parameter and Conway equations follow by the least-solution arguments of \cref{thm:conway}.
The comparison preserves the iterates and their suprema, proving the final assertion.
\end{enumerate}
\end{proof}

\subsection{Exact finite returns}
For this subsection an arrow $f:X\to\LL(Y+X)$ can be an arbitrary arrow of $\E$. Its partial-map representative and its two branches are
\[
 X\xleftarrow{\ m\ }P\xrightarrow{\ o\ }Y+X,
 \qquad P=P_Y+P_X,
\]
where $m$ is $\Sigma$-classified. Write $c_Y:P_Y\to X$, $c_X:P_X\to X$ for the source maps, $v:P_Y\to Y$ for the return, and $\nu:P_X\to X$ for the next state. Coproduct injections have decidable, hence $\Sigma$-classified, domains; composition therefore makes both branch domains $\Sigma$-classified.

\begin{definition}[Exact return object]
Define an internal family $\rho_n:R_n(f)\to X$ and output maps $v_n:R_n(f)\to Y$ by
\[
 R_0(f)=P_Y,\qquad
\begin{tikzcd}[column sep=large]
 R_{n+1}(f)\ar[r]\ar[d]&R_n(f)\ar[d,"\rho_n"]\\
 P_X\ar[r,"\nu"']&X.
\end{tikzcd}
\]
Here $\rho_0=c_Y$, $\rho_{n+1}$ is the left projection followed by $c_X$, and $v_{n+1}$ is the top projection followed by $v_n$. The family is constructed by recursion on subobjects of $X\times Y$, equivalently by finite lists of states and their transition witnesses.
\end{definition}

\begin{lemma}[Classified finite return]\label{exec:lem:return}
Every $\rho_n$ is a $\Sigma$-classified monomorphism, uniformly in $n$. Different return lengths have disjoint domains. The dependent sum has a $\Sigma$-classified monomorphism
\[
 D_f:=\sum_{n:\N}R_n(f)\ \incl\ X
\]
and an output map $D_f\to Y$. These constructions commute with parameter substitution.
\end{lemma}
\begin{proof}
\leavevmode
\begin{enumerate}[label=\textup{(\roman*)},leftmargin=*,itemsep=3pt,topsep=4pt]
\item \textbf{Induction on return domains.} The base case is $c_Y$.
Pullback and composition of classified monomorphisms prove the successor case.
Exit and continuation are disjoint, so $R_0\times_XR_{n+1}=0$.
Monicity of $c_X$ gives
\[
 R_{m+1}\times_XR_{n+1}
 \cong P_X\times_{\nu,X}(R_m\times_XR_n).
\]
Induction proves disjointness.
Within each length the source is monic, so the map from the dependent sum is monic as well.
Its classifier is
\[
 \defin_f(x)=\exists n:\N\;\rho_n(x),
\]
which belongs to $\Sigma$ by the established overtness of $\N$.
The $v_n$ give its output map.
\item \textbf{Coherence equations.} The sum denotes the dependent sum of one internally constructed family; no external countable coproduct is being invoked.
All the displayed pullbacks, recursions, and classified existential images commute with parameters.
\end{enumerate}
\end{proof}

Let $f^\ddagger:X\to\LL Y$ classify this partial map. If $\star$ denotes Kleisli extension, exposing the first step gives
\begin{equation}\label{exec:eq:unfold}
 f^\ddagger=[\eta_Y,f^\ddagger]^\star f.
\end{equation}
More explicitly, put $a_0=\bot$ and $a_{n+1}=[\eta_Y,a_n]^\star f$. The domain of $a_n$ is the sum of the $R_k$ for $k<n$, and its output is $v_k$ on that summand. Thus $f^\ddagger$ is their increasing union. Any solution of \eqref{exec:eq:unfold} agrees with $v_k$ on $R_k$, by induction. Hence $f^\ddagger$ is the least solution in the partial-map extension order.

\subsection{Complete Elgot partiality}
The following theorem places the preceding construction in the established theory of complete Elgot monads \cite{goncharovmiliusrauch}. For Kleisli arrows use $g\star f=g^\star f$, and write $g\oplus h$ for their coproduct.

\begin{theorem}[Complete Elgot structure]\label{exec:thm:elgot}
In the recursion setting of this section, $\LL$ is a complete Elgot monad on $\E$. Its lift to $\AS$ is also a complete Elgot monad. The same finite-return operation supplies both structures. For $f:X\to\LL(Y+X)$ it satisfies:
\begin{align}
 f^\ddagger&=[\eta_Y,f^\ddagger]\star f,
       \label{exec:eq:fix}\\
 g\star f^\ddagger
   &=\bigl((g\oplus\eta_X)\star f\bigr)^\ddagger,
       \label{exec:eq:nat}\\
 (h^\ddagger)^\ddagger
   &=\bigl(\LL[\id_{Y+X},\iota_X]h\bigr)^\ddagger,
       \label{exec:eq:codiag}
\end{align}
where $g:Y\to\LL Z$ and $h:X\to\LL((Y+X)+X)$. For $u:Z\to X$ total and $k:Z\to\LL(Y+Z)$, one also has
\begin{equation}\label{exec:eq:uniform}
 fu=\LL(\id_Y+u)k
 \quad\Longrightarrow\quad f^\ddagger u=k^\ddagger.
\end{equation}
The iteration operation
\[
 (-)^\ddagger:[X,\LL(Y+X)]_\Sigma\longrightarrow[X,\LL Y]_\Sigma
\]
is an algebra morphism. It agrees with intrinsic least-fixed-point iteration, commutes with parameter substitution, and agrees with semantic lifting wherever the computational--semantic comparison applies.
\end{theorem}
\begin{proof}
A finite successful history determines both a unique return value and a classified return domain. We compare these histories under naturality, codiagonal, and uniform simulation, and then identify the resulting operation with the least fixed point in the algebra category.

\leavevmode
\begin{enumerate}[label=\textup{(\roman*)},leftmargin=*,itemsep=3pt,topsep=4pt]
\item \textbf{Fixpoint and naturality.} The first law was proved by exposing the first transition.
For naturality, a successful run on the right follows exactly the continuations of $f$ and then performs $g$ on its return.
This identifies its domain and output with the left side, including the case in which $g$ is undefined.
\item \textbf{Codiagonal by finite histories.} For codiagonal, label a continuing $h$-transition by whether it uses the inner or outer occurrence of $X$. A finite successful history on the right is a finite word in those two labels followed by a $Y$-return.
Cut that word at every outer label.
The resulting finite list of inner blocks is precisely a history of $(h^\ddagger)^\ddagger$.
\item \textbf{Uniformity under simulation.} Conversely concatenate the blocks.
Both transformations preserve the initial state, all transition witnesses, and the final output; empty inner blocks account for consecutive outer steps.
They are inverse finite-list operations.
Equality of the classified return domains and output maps proves \eqref{exec:eq:codiag}.
\item \textbf{Lifting to computational algebras.} For uniformity, the commuting equation transports each $k$-transition along $u$, preserving its tag and returned value.
Inductively it identifies the domain of an $n$-step return from $z$ with the pullback of the corresponding domain from $u(z)$.
The same induction compares outputs.
Taking the classified existential image proves \eqref{exec:eq:uniform}.
These proofs are internal and hence work in every parameter context.
\item \textbf{Compatibility with least fixed points.} Now take $X,Y$ in $\AS$.
By \cref{exec:lem:sums}, copairing is available in $\AS$.
Cartesian closure and the lifted monad make
\[
 \Phi:[X,\LL(Y+X)]_\Sigma\times[X,\LL Y]_\Sigma
       \longrightarrow[X,\LL Y]_\Sigma,
 \qquad\Phi(f,a)=[\eta_Y,a]^\star f
\]
an algebra morphism.
\item \textbf{Scott continuity.} The second function object is pointed by the everywhere undefined map.
Apply \cref{prop:intrinsiclpf} to this parameterized endomorphism.
Its $n$th approximant is the $a_n$ above.
The algebra structure on $\LL Y$ computes the supremum of this chain by its classified union of extents and the corresponding values.
Once an input returns, subsequent approximants have the same value.
Its supremum is therefore exactly $f^\ddagger$.
\item \textbf{Algebraicity of the iteration operator.} This also proves that the iteration operation is an algebra morphism on the internal function objects.
The four equations in $\AS$ follow from those of underlying arrows.
The semantic comparison preserves lifting, the approximants, and their admitted suprema, giving the last comparison.
The parameter equation includes compatibility with the existing strength of $\LL$.

The interpretation maps form the commuting diagram
\[
\begin{tikzcd}[column sep=large,row sep=large]
X\ar[r,"f"]\ar[dr,"f^\ddagger"']&\LL(Y+X)\ar[d,"{[\eta_Y,f^\ddagger]^\star}"]\\ &\LL Y
\end{tikzcd}
\]
\end{enumerate}
\end{proof}

\begin{remark}
Complete Elgot iteration chooses coherent solutions to all equation morphisms $X\to\LL(Y+X)$. Elgot iteration selects the least coherent solution of each equation morphism. In particular, an equation which only continues admits more solutions than its least, everywhere undefined solution. Unique guarded solutions will occur in the temporal construction below. The internal argument also avoids replacing internally indexed chains by arbitrary external sequences of morphisms.
\end{remark}

\subsection{Coproduct feedback and its observations}
Write $\mathcal K_\Sigma=\Kl(\LL|_{\AS})$. Its finite coproducts are those of $\AS$, with the injections followed by the monad unit.

\begin{theorem}[Execution trace]\label{exec:thm:trace}
The symmetric monoidal category $(\mathcal K_\Sigma,+,0)$ has a trace. For $f:X+U\to\LL(Y+U)$ it is
\begin{equation}\label{exec:eq:trace}
 \Tr^{U}_{+}(f)
   =[\eta_Y,(f\iota_U)^\ddagger]^\star f\iota_X.
\end{equation}
It is an algebra morphism on the corresponding computational function objects and is stable under substitution in parameters. On the underlying partial maps, write $d,a,b,c$ for the branches $X\pmap Y$, $X\pmap U$, $U\pmap U$, and $U\pmap Y$. Then
\begin{equation}\label{exec:eq:paths}
 \Tr^{U}_{+}(f)
   =d\;\vee\!\bigvee_{n:\N}c\circ b^n\circ a.
\end{equation}
The summands in this formula have disjoint domains.
\end{theorem}
\begin{proof}
\leavevmode
\begin{enumerate}[label=\textup{(\roman*)},leftmargin=*,itemsep=3pt,topsep=4pt]
\item \textbf{Naturality and reindexing.} The formula first enters $f$ through $X$, then iterates its $U$-branch until it returns through $Y$.
The construction of exact returns proves \eqref{exec:eq:paths} and disjointness.
Naturality moves a map before entry or after exit.
Dinaturality identifies the finite paths obtained by cutting a composite interface on either side of its connecting map.
Vanishing at $0$ removes an unused interface; vanishing for $U+V$ cuts a finite path into its maximal visits to the respective interfaces, exactly as in the codiagonal proof.
\item \textbf{Observation preservation.} Superposing leaves the disjoint unaffected component unchanged.
For yanking, tracing the symmetry $U+U\to U+U$ sends an entering value once around the interface and then returns that same value.
These are the symmetric trace equations of \cite{jsv}.
Equivalently, they follow from the Conway--trace correspondence applied to the complete Elgot operation.
Algebra-morphism and parameter compatibility follow from \cref{exec:thm:elgot} and composition.

The interpretation maps form the commuting diagram
\[
\begin{tikzcd}[column sep=large,row sep=large]
X\ar[r,"\iota_X"]\ar[drr,"\Tr^U_+(f)"']&X+U\ar[r,"f"]&\LL(Y+U)\ar[d,"{[\eta_Y,(f\iota_U)^\ddagger]^\star}"]\\ &&\LL Y
\end{tikzcd}
\]
\end{enumerate}
\end{proof}

The union in \eqref{exec:eq:paths} concerns the complete entry--loop--exit paths. The powers $b^n$ alone need not be compatible partial functions, so no independent partial function $\bigvee_n b^n$ is required.

The interface feedback can be drawn in the symmetric monoidal graphical calculus as
\[
\begin{tikzpicture}[baseline=(current bounding box.center)]
 \node[draw,minimum width=2cm,minimum height=1cm] at (0,0) {$f$};
 \draw[->] (-2,0.25) node[left] {$X$} -- (-1,0.25);
 \draw[->] (1,0.25) -- (2,0.25) node[right] {$Y$};
 \draw[->] (1,-0.25) -- (1.55,-0.25) -- (1.55,-1.1)
    -- (-1.55,-1.1) node[midway,below] {$U$}
    -- (-1.55,-0.25) -- (-1,-0.25);
\end{tikzpicture}
\]
Its external arrow is $\Tr^U_+(f)$; its internal circuit is evaluated by finite return.

\begin{lemma}[Finite-history comparisons for the trace]\label{lem:tracehistories}
The trace equations identify classified domains as well as values. Their witness comparisons can be made by finite-list operations on complete returning histories.
\end{lemma}
\begin{proof}
\leavevmode
\begin{enumerate}[label=\textup{(\roman*)},leftmargin=*,itemsep=3pt,topsep=4pt]
\item \textbf{Finite-history comparisons for the trace.} A complete history consists of one entry transition, a finite list of interface transitions, and an exit transition, with the adjoining state equations.
Determinism makes its initial-state projection monic at each length; \cref{exec:lem:return} makes distinct lengths disjoint.
It therefore suffices to give witness maps over the initial state preserving the final output in both directions.
\item \textbf{Naturality.} For naturality with partial maps $a:X'\pmap X$ and $b:Y\pmap Y'$, retain the witness for $a$ as the initial item and the witness for $b$ as the final item.
The middle history is unchanged.
Deleting or adjoining those items gives the two descriptions of the same sequential domain.
In particular neither side is defined when its required initial or final partial call is undefined.
\item \textbf{Naturality and reindexing.} For dinaturality take $f:X+U\pmap Y+V$ and $g:V\pmap U$. A returning history of
$\Tr^U_+((\id_Y+g)f)$ has a sequence of alternating $f$- and $g$-calls after entry.
The history of $\Tr^V_+(f(\id_X+g))$ has the identical calls, but locates the interface cut immediately before each $g$-call; the first representation locates it immediately afterward.
Moving this cut changes no validity witness and no endpoint.
Moving it back is inverse.
This proves dinaturality for partial $g$, including its domain.
\item \textbf{Finite witness construction.} For vanishing at a sum $U+V$, label each visited state by its coproduct tag.
Cut the finite history at successive visits to $U$.
Each intervening finite block through $V$ is exactly an inner traced call; the remaining list is an outer history through $U$.
\item \textbf{Converse construction.} Conversely replace every inner call by its finite block.
Empty $V$-blocks and immediate returns are retained, so the constructions include all boundary cases.
Concatenation and cutting are inverse finite-list operations.
At the empty interface there is only the direct-return history.
For superposing, the unaffected summand never enters the interface, while histories in the other summand are unchanged.
For yanking, the symmetry provides exactly an entry and an exit carrying the same value.
\item \textbf{Coherence equations.} All these transformations use finite lists, coproduct cases, and composition of the displayed partial transitions.
They are internal and commute with parameters.
Their two directions identify the regular images of complete histories and preserve outputs, hence prove equality of the classified partial maps in every trace equation.
\end{enumerate}
\end{proof}

\begin{corollary}[Finite evidence for feedback]\label{exec:cor:observe}
For a positive observation $u:Y\to\Sigma$, the predicate transformer of feedback satisfies
\[
 (\Tr^U_+(f))^*u
 =d^*u\;\vee\!\bigvee_{n:\N}
       a^*(b^*)^n c^*u.
\]
Here the stars denote positive inverse-image transformers. Thus a positive observation of feedback is witnessed by a finite successful path and a positive observation of its returned value.
\end{corollary}
\begin{proof}
\leavevmode
\begin{enumerate}[label=\textup{(\roman*)},leftmargin=*,itemsep=3pt,topsep=4pt]
\item \textbf{Finite evidence for feedback.} Apply the transformer construction of \cref{sec:transformers} to the disjoint finite-return domains.
Composition reverses and classified existential images commute, giving the formula.
Overtness of $\N$ classifies the union.
In a realizability presentation, the witness translations of \cref{sec:extraction} retain the finite path and the terminal observation witness.
\end{enumerate}
\end{proof}

\subsection{Temporal self-interpretation and guarded recursion}
\label{exec:sec:guarded}\label{int:sec:temporal}
Let $\mathcal T=\operatorname{Seq}(\E)$ be the category of internal
inverse diagrams on the natural-number order: one object over $\N$
with restriction maps for $m\leq n$, satisfying identity and composition.
Finite limits are componentwise. This uses internal diagrams, without
external countable limits. We construct its prior--later interpretation
in the indexed terminology of \cite[Chapter~2]{ramesh}; the connection
with step-indexed domain theory is developed in \cite{birkedalguarded}.

\begin{construction}[The temporal indexed theory]\label{int:con:temporal}
Define
\[
 (\Prior A)_n=A_{n+1},\qquad
 (\triangleright A)_0=1,\quad(\triangleright A)_{n+1}=A_n.
\]
Adjacent restrictions give $p_A:\Prior A\to A$ and
$\operatorname{next}_A:A\to\triangleright A$, with the latter
terminal at zero. Put
\[
 \mathcal C(X)=\mathcal T/\Prior X,\qquad
 u^{\mathcal C*}=(\Prior u)^*,\qquad
 F_X=p_X^*:\mathcal T/X\to\mathcal C(X).
\]
A pre-introspective theory is a finite-limit category with an indexed
finite-limit category and such an indexed functor from its self-indexing.
It is locally introspective when the indexed hom-presheaves between
displayed objects are representable. These are the definitions of
\cite[Definitions~2.1 and~2.8]{ramesh}.
\end{construction}

\begin{theorem}[Local introspection and constructed guarded solutions]
\label{int:thm:temporal}\label{exec:prop:guardfix}
The displayed theory is locally introspective. Its global interpretation
is the identity, its indexed global-element object at $A$ is
$\triangleright A$, and its induced name map is $\operatorname{next}_A$.
For every $h:Z\times\triangleright A\to A$ there is a unique arrow
$h^\sharp:Z\to A$ satisfying
\[
\begin{tikzcd}[column sep=large]
 Z\ar[r,"{\langle\id,\operatorname{next}_Ah^\sharp\rangle}"]
   \ar[dr,"h^\sharp"']&Z\times\triangleright A\ar[d,"h"]\\
 &A.
\end{tikzcd}
\]
The construction is natural in $Z$, holds in slices over constant
diagrams, and supplies guarded Conway identities and guarded cartesian trace.
\end{theorem}
\begin{proof}
\leavevmode
\begin{enumerate}[label=\textup{(\roman*)},leftmargin=*,itemsep=3pt,topsep=4pt]
\item \textbf{Exponentials of inverse diagrams.} The finite-prefix formula
\[
 (B^A)_n=
 \left\{(u_k)_{k\leq n}\in\prod_{k\leq n}B_k^{A_k}
 \ \middle|\ r^B_{jk}u_k=u_jr^A_{jk}\ (j\leq k\leq n)\right\}
\]
constructs exponentials in $\mathcal T$. Restriction truncates;
evaluation takes $u_n$; currying restricts the parameter to each
earlier stage. The maps are inverse. In the slice over $X$, use at
$x\in X_n$ the same compatible family of maps between the fibres
over $r^X_{kn}x$. Dependent products and equalizers in $\E$ construct
these families uniformly, proving local cartesian closure.
\item \textbf{The prior–later adjunction.} There is an adjunction $\Prior\dashv\triangleright$:
$u:\Prior A\to B$ transposes to $\bar u_0=!$,
$\bar u_{n+1}=u_n$. Write its unit as
$\iota_A:A\to\triangleright\Prior A$. Naturality of $p$ gives
the indexed comparison
\[
\begin{tikzcd}[column sep=large]
 \mathcal T/X\ar[r,"p_X^*"]\ar[d,"u^*"']&
 \mathcal T/\Prior X\ar[d,"(\Prior u)^*"]\\
 \mathcal T/Y\ar[r,"p_Y^*"']&\mathcal T/\Prior Y.
\end{tikzcd}
\]
All functors preserve finite limits; the pullback comparison is
coherent by uniqueness of its projections.
\item \textbf{Representability of displayed homs.} The induced $\Prior_X:\mathcal T/X\to\mathcal T/\Prior X$
has right adjoint
$R_X(H)=X\times_{\triangleright\Prior X}\triangleright H$,
using $\iota_X$. This follows by restricting the adjunction to
commuting triangles. For $A,B\in\mathcal C(X)$ and $u:Y\to X$,
\[
 \Hom_{\mathcal T/X}(Y,R_X(B^A))
 \cong\Hom_{\mathcal T/\Prior Y}((\Prior u)^*A,(\Prior u)^*B).
\]
Thus $R_X(B^A)$ represents the indexed hom-presheaf. Since
$\Prior1=1$, $F_1=\id$. Indexed global elements of $A$ are
$\Hom_{\mathcal T}(\Prior X,A)$, represented by $\triangleright A$.
The transpose of $p_A$ is
$\triangleright p_A\,\iota_A=\operatorname{next}_A$.
\item \textbf{Unique guarded recursion.} Put $h^\sharp_0=h_0\langle\id,!\rangle$ and define the successor by recursion:
\[
 h^\sharp_{n+1}=h_{n+1}\langle\id,h^\sharp_nr^Z_n\rangle.
\]
Naturality of $h$ proves restriction compatibility by induction.
Every solution satisfies this recurrence, proving uniqueness.
Compatible finite prefixes form an internal object; induction gives
their unique existence, so their endpoints define the family without
choice. Parameter substitution commutes with each clause. In the
slice over $\Delta X$ replace the zeroth terminal object by $X$;
write the resulting shift as $\triangleright_X$. Unique guarded
solutions give the guarded Conway equations and guarded trace by \cite{miliuslitak}.
\end{enumerate}
\end{proof}

The higher-order rule is
\[
 \frac{z:Z,\ x:\triangleright A\ \vdash\ t:A}
 {z:Z\ \vdash\operatorname{fix}(x.t):A},\qquad
 \operatorname{fix}(x.t)=t[\operatorname{next}(\operatorname{fix}(x.t))/x].
\]
Domain algebras stage componentwise; the shift uses the terminal algebra
at zero. Local introspection here represents the displayed hom-objects.
Representing the whole indexed category is the stronger notion used
for the geminal construction in \cite[Construction~5.18]{ramesh};
the guarded solutions above follow from their explicit recurrence.
The realizer exposure and free term monad constructed below retain
their separate interpretation maps.


\subsection{Guarded running processes and finite completion}
Return to $f:X\to\LL(Y+X)$ and its two branch objects. In the temporal slice over $\Delta X$, let
\[
 \Psi_f(E)=\Delta P_Y+_{\Delta X}
       \Sigma_{\Delta c_X}(\Delta\nu)^*E.
\]
Here $\Sigma_{\Delta c_X}$ is dependent sum in the slice; the subscripted plus is its coproduct. Define $G_0=X$ and
\[
 G_{n+1}=P_Y+_X(P_X\times_{\nu,X}G_n).
\]
The first restriction is $[c_Y,c_X]:G_1\to X$ and each succeeding restriction is the identity on the return summand and the preceding restriction on the continuation summand. Write $\operatorname{Run}(f)$ for this inverse diagram over $X$.

\begin{proposition}[Guarded execution and its normal form]\label{exec:prop:running}
There is an isomorphism
\[
 \operatorname{Run}(f)\cong
      \triangleright_X\Psi_f(\operatorname{Run}(f)),
\]
unique up to the unique isomorphism respecting this fixed-object structure. If $H_n$ is the domain of $n$ successive continuations, with $H_0=X$, then
\begin{equation}\label{exec:eq:normal}
 G_n\cong\left(\sum_{k<n}R_k(f)\right)+_XH_n.
\end{equation}
All these isomorphisms are natural in parameter substitution.
\end{proposition}
\begin{proof}
\leavevmode
\begin{enumerate}[label=\textup{(\roman*)},leftmargin=*,itemsep=3pt,topsep=4pt]
\item \textbf{Finite witness construction.} The fixed-object equation is exactly the stage recurrence and its restriction maps.
Any other fixed object has zeroth stage isomorphic to $X$; its next-stage isomorphism is then forced by the preceding one.
Induction constructs the unique compatible isomorphism at every finite prefix, hence of the internal family.
For \eqref{exec:eq:normal}, the zeroth case is $H_0=X$.
Distribute the successor pullback over the finite coproduct, and use
\[
 P_Y=R_0,\qquad
 P_X\times_XR_k=R_{k+1},\qquad
 P_X\times_XH_n=H_{n+1},
\]
where the pullbacks use the continuation map.
\item \textbf{Pullback construction.} These identifications preserve the source and output maps.
The constructions use the same pullbacks in every parameter slice.
\end{enumerate}
\end{proof}

\begin{definition}[Finite completion of a running family]\label{exec:def:fincomp}
Let $(C_n\xrightarrow{\rho_n}X,w_n:C_n\to Y)_{n:\N}$ be an internally indexed family such that every $\rho_n$ is a $\Sigma$-classified mono, the domains are pairwise disjoint, and the structure is uniform in $n$. Its \textbf{finite completion} is the partial map represented by
\[
 X\xleftarrow{\ \rho\ }\sum_{n:\N}C_n
   \xrightarrow{\ [w_n]\ }Y,
\]
where $\rho$ is induced by the $\rho_n$. Pairwise disjointness makes $\rho$ monic, and its classifier is
\[
 x\longmapsto\exists n:\N\;[x\in C_n],
\]
which belongs to $\Sigma$ by overtness of $\N$. Thus finite completion is a constructed $\Sigma$-partial map, not an additional operation on arbitrary temporal objects.

For $\operatorname{Run}(f)$ we take $C_n=R_n(f)$ and $w_n=v_n$, and write the resulting partial map as $\operatorname{FinComp}(\operatorname{Run}(f))$.
\end{definition}

For feedback $h:X+U\to\LL(Y+U)$, let $D_d\incl X$ and $m_a:D_a\incl X$ be the domains of its direct-exit and entering branches, and let $v_a:D_a\to U$ be the entering value. Its guarded entering process is the following object over $\Delta X$:
\[
 \operatorname{GExec}^U(h)=\triangleright_X
 \left(\Delta D_d+_{\Delta X}
   \Sigma_{\Delta m_a}(\Delta v_a)^*
          \operatorname{Run}(h\iota_U)\right).
\]
Its first completed stratum is the direct exit, and its $(n+2)$nd completed stratum is the path $c b^n a$.

\begin{theorem}[Guarded execution, iteration, and feedback agree]\label{exec:thm:comparison}
Finite completion satisfies
\[
 \operatorname{FinComp}(\operatorname{Run}(f))=f^\ddagger.
\]
For $h:X+U\to\LL(Y+U)$, its guarded entering process has the direct-exit branch $d$ and the entering branch $a$ followed by $\operatorname{Run}(h\iota_U)$. Denote it by $\operatorname{GExec}^{U}(h)$. Then
\begin{equation}\label{exec:eq:guardtrace}
 \operatorname{FinComp}(\operatorname{GExec}^{U}(h))
       =\Tr^U_+(h).
\end{equation}
The equalities include classified termination domains, returned values, positive observations, and parameter substitution. For algebra morphisms, the completed operations belong to $\AS$.
\end{theorem}
\begin{proof}
\leavevmode
\begin{enumerate}[label=\textup{(\roman*)},leftmargin=*,itemsep=3pt,topsep=4pt]
\item \textbf{Guarded execution iteration and feedback agree.} By \eqref{exec:eq:normal}, the completed strata of $\operatorname{Run}(f)$ are exactly the $R_k(f)$, with the same source and output maps as in the construction of $f^\ddagger$.
Hence \cref{exec:def:fincomp} produces the very span
\[
 X\xleftarrow{}\sum_kR_k(f)\xrightarrow{[v_k]}Y
\]
used to define $f^\ddagger$.
This proves the first equality, including its classified domain and returned value.
\item \textbf{Partiality and definedness.} For entering feedback the completed strata are the direct-exit domain of $d$ and, for $n:\N$, the exact finite path $c\,b^n a$.
Determinism of the tagged partial computation and the disjointness argument of \cref{exec:lem:return} make these strata pairwise disjoint.
Their finite completion is therefore the partial map
\[
 d\;\vee\!\bigvee_{n:\N}c\circ b^n\circ a,
\]
with the same dependent-sum domain and output maps.
By \cref{exec:thm:trace} this is $\Tr^U_+(h)$, proving \eqref{exec:eq:guardtrace}.
\item \textbf{Pullback construction.} Apply \cref{exec:cor:observe} for the observation comparison and \cref{exec:thm:elgot} for the algebra-morphism assertion.
Every construction used in \cref{exec:def:fincomp} is stable under pullback, which gives the parameter statement.
\end{enumerate}
\end{proof}

Finite completion forgets elapsed time on successful histories. It is distinct from taking a limit of all running stages, which can also retain indefinitely continuing trajectories. The comparison is consequently an explicit map from the constructed execution data to partial-result semantics. Its relation to the general theory of guarded and unguarded iteration is discussed in \cite{goncharovetal}; foundational accounts connecting Elgot iteration, dominances, and restriction categories include \cite{goncharovfoundations}.

\begin{example}[A loop and its temporal information]\label{exec:ex:loop}
For $X=1$, the outcome $f=\eta\iota_X:1\to\LL(Y+1)$ always continues. Every $H_n$ is $1$, every $R_n$ is $0$, and $\operatorname{Run}(f)$ has every stage $1$. Its finite completion is the undefined partial result. If $Y=X=\N$ and the outcome returns $0$ at state $0$ and continues from $n+1$ to $n$, then $R_k$ consists of the initial state $k$ and the iteration returns $0$ on all inputs. The guarded diagram retains the number of transitions; finite completion retains their common returned value. Neither construction selects a program index.
\end{example}

\subsection{The interaction category}\label{exec:sec:interaction}
The coproduct trace gives a categorical interaction semantics by the $\operatorname{Int}$ construction of \cite{jsv}. Define
\[
 \mathcal I_\Sigma=\operatorname{Int}(\mathcal K_\Sigma).
\]
An object is a pair $(A^+,A^-)$ and its arrows are
\begin{equation}\label{exec:eq:inthom}
 \mathcal I_\Sigma((A^+,A^-),(B^+,B^-))
  =\AS(A^++B^-,\LL(B^++A^-)).
\end{equation}
Composition connects the $B^+$ and $B^-$ interfaces, using the symmetries of the coproduct, and traces out that shared interface. Thus composition performs precisely the finite-return execution already constructed.

\begin{corollary}[Compact-closed interaction semantics]\label{exec:cor:int}
The category $\mathcal I_\Sigma$ is compact closed, with dual $(A^+,A^-)^*=(A^-,A^+)$, tensor induced by coproduct on both components, and unit $(0,0)$. The functor $A\mapsto(A,0)$ embeds $\mathcal K_\Sigma$ fully faithfully and preserves its symmetric monoidal trace.
\end{corollary}
\begin{proof}
\leavevmode
\begin{enumerate}[label=\textup{(\roman*)},leftmargin=*,itemsep=3pt,topsep=4pt]
\item \textbf{Compact-closed interaction semantics.} Apply the symmetric $\operatorname{Int}$ theorem to \cref{exec:thm:trace}.
The evaluation and coevaluation are the identity wires with the indicated polarities; their triangle equations are yanking.
Associativity of interaction is vanishing for nested feedback.
The displayed hom formula proves full faithfulness at $(A,0)$.
\end{enumerate}
\end{proof}

This is the categorical interaction construction underlying the connection with Geometry of Interaction; see \cite{abramskyhaghverdisc,hasegawaclosed}. The result specifies a compact-closed category of effective partial interactions. The Cartesian trace of \cref{thm:conway} and the coproduct execution trace act on different categories and different tensors, while both are determined here by the same intrinsic least-solution construction. Code implementations of the latter comparison are constructed in \cref{exec:thm:codefeedback}.


\section{Lawvere fixed points and synthetic continuity}\label{sec:lawvere}


There are two useful categorical descriptions of solutions of recursion equations. Directed completion constructs a least solution; diagonal representation proves existence of a solution. We keep their solution objects and comparison maps explicit.

\subsection{The solution object and the diagonal}
For an object $D$ of a topos, define
\[
 \Fix(D)\incl D^D\times D
 \quad\text{as the equalizer of }\ev\text{ and }\pi_2,
 \qquad p_D:\Fix(D)\to D^D.
\]
The internal fixed-point property says that $p_D$ is epic. A chosen fixed-point operator is a section, and leastness is further structure on such a section.

\begin{theorem}[Lawvere, internal form]\label{thm:lawvere}
If $e:Q\onto D^Q$ is an epimorphism, then $p_D$ is epic.
\end{theorem}
\begin{proof}
\leavevmode
\begin{enumerate}[label=\textup{(\roman*)},leftmargin=*,itemsep=3pt,topsep=4pt]
\item \textbf{Existence and uniqueness.} For an endomorphism $f$ form the diagonal composite
\[
 Q\xrightarrow{\Delta}Q\times Q
 \xrightarrow{e\times\id}D^Q\times Q
 \xrightarrow{\ev}D\xrightarrow{f}D.
\]
Internally it is $x\mapsto f(e(x)(x))$.
Epicity gives a local representative $a$ of this map.
Then $e(a)(a)=f(e(a)(a))$.
The construction works in context $f$, which proves the asserted regular-epimorphic projection.
No section of $e$ has been chosen.
This is the internal epimorphic version of Lawvere's diagonal argument \cite{lawvere}.
\end{enumerate}
\end{proof}

\subsection{The effective-topos model theorem}
The following internal model theorem is used in the fixed topos $\Eff$. We use Bauer's terminology \cite[\S3]{bauerfixed}. A \textbf{countably generated $\omega$-cppo} is a pointed $\omega$-cpo equipped with a countable base whose induced order is decidable and such that every element is the supremum of an increasing sequence of basic elements. A \textbf{domain} is a countably generated $\omega$-cppo admitting such a base of compact elements. These are properties of the internal order; no program numbering is part of their definition.

\paragraph{Logical scope of the cited model theorem.}
Bauer's development is carried out in intuitionistic higher-order logic together with Countable Choice, Markov's principle, and the Enumeration axiom \cite[\S2]{bauerfixed}. The first supplies a choice function for a countable family of inhabited relations, Markov's principle turns nonconstancy of a binary sequence into the existence of a hit, and the Enumeration axiom provides an enumeration of the countable subsets of $\N$. Bauer records that these principles are validated in the effective topos and in realizability topoi based on Turing machines relative to a fixed oracle. This assertion concerns those models, not arbitrary localic forcing sheaf topoi. We use them only through the following cited model theorem. They are not hypotheses of the general dominance-selected theory developed in the preceding parts.

\begin{theorem}[Synthetic model theorem]\label{thm:syntheticmodel}
In $\Eff$, every countably generated $\omega$-cppo has the fixed-point property and every total relation on it has a fixed point. Every internal map from an $\omega$-cpo to a domain is monotone and preserves increasing sequential suprema. These conclusions hold internally, including parameters.
\end{theorem}
\begin{proof}
\leavevmode
\begin{enumerate}[label=\textup{(\roman*)},leftmargin=*,itemsep=3pt,topsep=4pt]
\item \textbf{Fixed-point theorem in the effective topos.} The fixed-point statement is \cite[Theorem 4.2 and Corollary 4.3]{bauerfixed}; the relational form is \cite[Theorem 5.2]{bauerfixed}; and the continuity statement is \cite[Theorem 6.4]{bauerfixed}.
Their proofs use precisely the model principles isolated above.
Since the results are internal universal statements in $\Eff$, they may be used in an arbitrary parameter context.
\end{enumerate}
\end{proof}

For a countable internal poset $B$ with bottom and decidable order, write $\omega\Idl(B)$ for the object of \textbf{countable ideals} in Bauer's sense. Concretely, an ideal $I\subseteq B$ is countable when internally there is a map
\[
 e:\N\longrightarrow 1+B
\]
whose non-dummy values have image exactly $I$. This is a property of the ideal, not a selected enumeration retained as part of the ideal datum.

\begin{lemma}[Rosolini ideals are the countable ideals]\label{lem:Jomega}
Let $B$ be an internal poset in $\Eff$ with a regular-epimorphic enumeration $\nu:\N\onto B$, decidable equality, and decidable order. Then there is a canonical equality of subobjects of $\Idl(B)$,
\[
 \J B=\omega\Idl(B).
\]
The equality is stable under parameters.
\end{lemma}
\begin{proof}
\leavevmode
\begin{enumerate}[label=\textup{(\roman*)},leftmargin=*,itemsep=3pt,topsep=4pt]
\item \textbf{Rosolini ideals are the countable ideals.} The enumeration makes $B$ overt: for $u:B\to\Sigma$,
\[
 \exists b:B\;u(b)
 \quad\Longleftrightarrow\quad
 \exists n:\N\;u(\nu n),
\]
and the right side belongs to $\Sigma$ by overtness of $\N$.
Decidable order is in particular $\Sigma$-classified.
Hence \cref{jd:prop:basis} identifies $\J B$ with the ideals whose membership predicate is $\Sigma$-classified.
\item \textbf{Ideal and directedness calculation.} Let such an ideal $I$ be given in a parameter context $S$.
Take a partitioned-assembly cover $p:P\onto S$.
The composite membership predicate
\[
 P\times\N\xrightarrow{\id\times\nu}P\times B
   \xrightarrow{I}\Sigma
\]
has projective source, so it lifts through $2^{\N}\onto\Sigma$ to a tracked binary-stream family $\theta(p,n,k)$.
\item \textbf{Product calculation.} Fix a primitive recursive pairing of two natural numbers and define, over $P$,
\[
 e_p(\langle n,k\rangle)=
 \begin{cases}
  \mathsf{inr}(\nu n),&\theta(p,n,k)=1,\\
  \mathsf{inl}(*),&\theta(p,n,k)=0.
 \end{cases}
\]
If a non-dummy value is emitted, its value belongs to $I$ by correctness of the lifted classifier.
\item \textbf{Converse construction.} Conversely, if $I(b)$ holds, regular epicity of $\nu$ gives locally an $n$ with $\nu n=b$, and the equality between the support of $\theta(p,n,-)$ and $I(\nu n)$ gives a hit $k$; hence $b$ occurs among the non-dummy values of $e_p$.
Thus the image of those values is exactly the pullback of $I$ to $P$.
Countability is an existential property, so it descends along the regular epimorphism $P\onto S$.
This proves $\J B\subseteq\omega\Idl(B)$ uniformly in $S$, without choosing a stream representative for every ideal.
\item \textbf{Further converse construction.} Conversely, let $I$ be countable and let $e:\N\to1+B$ enumerate it.
For $b:B$,
\[
 I(b)\quad\Longleftrightarrow\quad
 \exists n:\N\;[e(n)=\mathsf{inr}(b)].
\]
The bracketed predicate is decidable because the coproduct tag and equality in $B$ are decidable.
Overtness of $\N$ therefore makes the right side $\Sigma$-classified.
Hence $I\in\J B$ by \cref{jd:prop:basis}.
The two inclusions are inclusions of the same ideals inside $\Idl(B)$, so their comparison is canonical and stable under substitution.
\end{enumerate}
\end{proof}

\begin{lemma}[Cofinal chains from ideal enumerations]\label{lem:cofinalchain}
Let $B$ have decidable order and let $e:\N\to1+B$ enumerate an inhabited directed ideal $I$. Locally over the choice of one member of $I$, there is an increasing sequence $b_0\leq b_1\leq\cdots$ of members of $I$ which is cofinal in $I$. The construction is uniform in $e$ and the selected member.
\end{lemma}
\begin{proof}
\leavevmode
\begin{enumerate}[label=\textup{(\roman*)},leftmargin=*,itemsep=3pt,topsep=4pt]
\item \textbf{Ideal and directedness calculation.} Use the selected member as $b_0$.
Having constructed $b_n$, if $e(n)$ is dummy put $b_{n+1}=b_n$.
If $e(n)=\mathsf{inr}(c)$, search for the least $k$ such that $e(k)=\mathsf{inr}(d)$ and $b_n\leq d$, $c\leq d$.
This test is decidable.
Directedness gives a suitable $d\in I$, and enumeration gives an index for it, so the least such $k$ exists.
The least-index relation is total and functional in this context.
Unique choice and recursion on $\N$ consequently define the sequence.
\item \textbf{Termination of the unbounded search.} This is an unbounded search whose termination has been proved, not an invocation of countable choice.
\item \textbf{Induction.} Induction gives $b_n\in I$ and $b_n\leq b_{n+1}$.
Every non-dummy value $e(n)$ lies below $b_{n+1}$, so $I=\bigcup_n\downarrow b_n$.
The initial member exists locally because $I$ is inhabited.
In realizability, the decidable search is implemented by the same first-hit recursion used in \cref{lem:hitsearch}; hence the construction is tracked on its stated domain.
\end{enumerate}
\end{proof}

The lemma explains why a countable directed presentation and an increasing-chain presentation give the same ideal completion here. Equality of the resulting ideal is independent of the local initial member or the search enumeration. A chain supremum, when it is used below, is therefore the supremum of the already specified ideal.

\begin{proposition}[Identification of the intrinsic function domain]\label{prop:omegaidentification}
In $\Eff$, and in realizability topoi based on Turing machines relative to a fixed oracle, the computation object
\[
 P=\PC=\J\FG\cong(\LL\N)^{\N}
\]
is the countable-ideal completion $\omega\Idl(\FG)$ and is a domain in Bauer's sense. The same holds for
$W=\J\Fin(\N)\cong\Sigma^\N$. Also $\Sigma\cong\J2$ is a domain, where $2$ is ordered as a chain.
\end{proposition}
\begin{proof}
\leavevmode
\begin{enumerate}[label=\textup{(\roman*)},leftmargin=*,itemsep=3pt,topsep=4pt]
\item \textbf{Ideal and directedness calculation.} The finite-graph object has the empty graph as least element, its canonical finite-list coding gives a regular-epimorphic enumeration, and equality and inclusion are decidable.
Thus \cref{lem:Jomega} gives
\[
 \J\FG=\omega\Idl(\FG).
\]
Bauer's countable-ideal construction and its free-completion theorem \cite[\S3, Proposition 3.1]{bauerfixed} supply an $\omega$-cpo whose principal ideals form a countable base of compact elements; hence it is a domain.
\item \textbf{Partiality and definedness.} Under \cref{prop:graphrepresentation} this is the already constructed object of effective partial numerical functions.
\item \textbf{Ideal and directedness calculation.} The same argument applies to $\Fin(\N)$, whose bottom is the empty finite set and whose inclusion order is decidable.
For the two-element chain, \cref{lem:Jomega} identifies $\J2$ with its countable ideal completion, while the semantic comparison identifies this object with the Rosolini truth values $\Sigma$.
Bauer's proposition again makes it a domain.
All identifications are internal and therefore persist in parameter contexts.
\end{enumerate}
\end{proof}

No member of $P$ has been assigned a program number in this identification. Its basis consists of finite graphs, its elements are partial numerical functions, and its equality is extensional equality. The theorem uses the ambient realizability model to establish properties of this semantic object.

\section{Extensional halting, decisions, and finite information}\label{sec:extensionaltheorems}

Evaluation has a positive termination observation, but the existence of that observation does not make halting decidable. We distinguish these two assertions at the level of effective function objects. The subsequent observation-fullness result identifies precisely which extensional semidecisions come from finite information.


\subsection{Halting as observation of evaluation}
Whenever the indicated computational function object is used, termination is the composite
\[
 H_{X,Y}:[X,\LL Y]_\Sigma\times X
 \xrightarrow{\ev}\LL Y\xrightarrow{\defin}\Sigma.
\]
It is an observation by the lifted partiality construction. On the typed function object of \cref{sec:programdomain} this is the ambient map
\[
 H:P\times\N\to\Sigma,
 \qquad H(f,n)=\defin(f(n)).
\]
The two descriptions agree under \cref{eq:programpts}; they observe the same classified domain and value of evaluation.

\begin{lemma}[Fixed points and decidable partitions]\label{lem:fprice}
Internally, an inhabited object with the fixed-point property has only constant maps to $2$.
\end{lemma}
\begin{proof}
\leavevmode
\begin{enumerate}[label=\textup{(\roman*)},leftmargin=*,itemsep=3pt,topsep=4pt]
\item \textbf{Fixed points and decidable partitions.} For $d:D\to2$ and $x,y:D$, suppose $d(x)\ne d(y)$.
Boolean case distinction defines an endomap taking $z$ to $y$ when $d(z)=d(x)$ and to $x$ otherwise. A fixed point would have the opposite decision value to itself.
Thus $d(x)=d(y)$, since equality in $2$ is decidable.
The argument is in context $x,y$ and proves internal constancy.
This is the fixed-point form of the Rice argument.
\end{enumerate}
\end{proof}

\begin{theorem}[Intrinsic Rice and halting]\label{thm:intrinsichalting}
In $\Eff$ every internal decision $P\to2$ is constant, and likewise every decision $W\to2$. The termination observation $H:P\times\N\to\Sigma$ has no Boolean classifier and no complementary $\Sigma$-predicate. Definedness on $\LL\N$ has the same properties.
\end{theorem}
\begin{proof}
\leavevmode
\begin{enumerate}[label=\textup{(\roman*)},leftmargin=*,itemsep=3pt,topsep=4pt]
\item \textbf{Evaluation and transposition.} Apply \cref{thm:syntheticmodel,prop:omegaidentification,lem:fprice}. A classifier for $H$, restricted to any fixed numeral, would distinguish the nowhere-defined function from a total constant function, contradicting constancy on $P$. A classifier for definedness pulls back along evaluation at that numeral to such a classifier.
Complementary positive predicates induce a map to $2$: their disjoint union is the whole object, so the two constant Boolean maps descend to a decision.
This would contradict the statements just proved.
\end{enumerate}
\end{proof}

The preceding theorem concerns all ambient decision maps on these objects. The general observational decision theorem \cref{thm:syntheticrice} is valid in the constructed computational category over $\Eff$. The effective-topos fixed-point result supplies the stronger identification here.

\subsection{Observation fullness before program presentations}
\begin{theorem}[Intrinsic finite information]\label{thm:observationfullness}
In $\Eff$, restriction is an isomorphism
\[
 R_{\FunDom}:\OO_{\FunDom}\xrightarrow{\cong}\Sigma^P.
\]
For $u:P\to\Sigma$, its inverse has positive trace
$S_u(g)=u(\downarrow g)$ and is
\[
 U_u(I)=\exists g:\FG\,[I(g)\wedge S_u(g)].
\]
The isomorphism is internal and natural in every parameter. There is the analogous isomorphism $\OO_{\mathcal W}\cong\Sigma^W$ for the finite-subset domain.
\end{theorem}
\begin{proof}
\leavevmode
\begin{enumerate}[label=\textup{(\roman*)},leftmargin=*,itemsep=3pt,topsep=4pt]
\item \textbf{Synthetic Scott continuity.} By \cref{thm:syntheticmodel,prop:omegaidentification}, every $u:P\to\Sigma$ is monotone and preserves increasing sequential suprema.
Let $I\in P$.
By the countable-ideal presentation, the cofinal-chain construction in the proof of \cite[Proposition 3.1]{bauerfixed} supplies, internally in the present model, an increasing cofinal sequence $g_n\in I$ with
\[
 I=\bigvee_n\downarrow g_n.
\]
The supremum in $\Sigma$ of an increasing sequence of truth values is their overt existential, so continuity gives
\[
 u(I)=\bigvee_n u(\downarrow g_n)
      \iff\exists n:\N\;u(\downarrow g_n).
\]
Put $S_u(g)=u(\downarrow g)$.
\item \textbf{Cofinal finite approximations.} Cofinality and monotonicity give the basis-independent form
\[
 \exists n\;u(\downarrow g_n)
 \quad\Longleftrightarrow\quad
 \exists g:\FG\,[I(g)\wedge S_u(g)].
\]
Indeed, the forward implication uses $g_n\in I$; conversely, from $g\in I$ choose a cofinal stage with $g\leq g_n$ and use monotonicity of $u$ on principal ideals.
Thus the formula
\[
 U_u(J)=\exists g:\FG\,[J(g)\wedge S_u(g)]
\]
restricts to $u$ on every effective point $I$.
\item \textbf{Construction of the positive trace.} The trace $S_u$ is $\Sigma$-valued and upper, because $g\leq h$ gives $\downarrow g\subseteq\downarrow h$ and hence $u(\downarrow g)\leq u(\downarrow h)$.
Therefore its ordinary ideal extension $U_u$ is a Scott observation on all semantic ideals by \cref{eq:graphopen}.
\item \textbf{Inverse identity and parameter naturality.} Conversely, if $U\in\OO_{\FunDom}$, its restriction followed by the displayed construction has trace
\[
 g\longmapsto U(\downarrow g),
\]
which is the original trace of $U$; \cref{lem:tracefree} therefore recovers $U$ on every ideal.
The two constructions are inverse.
\item \textbf{Observation preservation.} The argument was carried out with arbitrary internal parameters: the cofinal-chain and continuity statements used above are the internal statements supplied by \cref{thm:syntheticmodel,prop:omegaidentification}, and every displayed formula is stable under substitution.
Hence the inverse is a morphism $\Sigma^P\to\OO_{\FunDom}$ and $R_{\FunDom}$ is an object isomorphism.
The finite-subset case is identical.

The comparison is witnessed by the retraction diagram
\[
\begin{tikzcd}[column sep=large,row sep=large]
\OO_{\FunDom}\ar[r,"R_{\FunDom}"]\ar[dr,"\id"']&\Sigma^{P}\ar[d,"U_{(-)}"]\\ &\OO_{\FunDom}
\end{tikzcd}
\]
\end{enumerate}
\end{proof}

This is the extensional Rice--Shapiro principle. A semidecision on partial functions has a positive finite-information witness because it is an internal map on the computational object. The result has been obtained from synthetic continuity and the canonical completion, without a predicate on program indices.

\begin{corollary}[Automatic effective extension]\label{thm:syntheticms}
Every internal map $P\to P$ extends uniquely to an effective Scott-continuous map $\FunDom\to\FunDom$. The same holds for $W$, and for maps between these two computation objects. The extension is natural in parameters and preserves composition.
\end{corollary}
\begin{proof}
\leavevmode
\begin{enumerate}[label=\textup{(\roman*)},leftmargin=*,itemsep=3pt,topsep=4pt]
\item \textbf{Automatic effective extension.} For $f:P\to P$, inverse image gives $\Sigma^f:\Sigma^P\to\Sigma^P$.
Conjugate by the observation isomorphisms of \cref{thm:observationfullness}.
The resulting observation transpose satisfies the evaluation square, so \cref{thm:chu} gives the unique extension.
The mixed and finite-subset cases use the corresponding source and target isomorphisms.
Uniqueness proves the composition and parameter equations.
\end{enumerate}
\end{proof}

\begin{corollary}\label{cor:intrinsicnonpositive}
Totality and the property of being the nowhere-defined element are not $\Sigma$-classified predicates on $P$.
\end{corollary}
\begin{proof}
\leavevmode
\begin{enumerate}[label=\textup{(\roman*)},leftmargin=*,itemsep=3pt,topsep=4pt]
\item \textbf{Unit and identity laws.} If totality were positive, \cref{thm:observationfullness} would extend it to an observation.
The finite approximations to the identity in \cref{ex:totality} contradict preservation of their increasing supremum.
If equality to bottom were positive, its extension would be an upper observation containing bottom and hence every point, contrary to its value at a total constant function.
\end{enumerate}
\end{proof}

\section{Relational recursion and transport to presentations}\label{sec:relationalrecursion}

The previous fixed-point arguments selected values or predicates. A total relation can express a solution condition without presenting a function that chooses a successor. We formulate its fixed-point object geometrically and then explain how such relations transport to the implementation structures constructed next.


A total relation on $D$ is a subobject $R\incl D\times D$ with epic first projection. A relational fixed point is a member of the pullback
\[
\begin{tikzcd}
 \Fix(R)\ar[r]\ar[d]&R\ar[d,hook]\\
 D\ar[r,"\Delta_D"']&D\times D.
\end{tikzcd}
\]
The multivalued conclusion of \cref{thm:syntheticmodel} says that this fixed-point object is inhabited for every total relation on the indicated domains.

\begin{proposition}[Relational transport]\label{prop:relationaltransport}
For an epimorphism $q:C\onto D$, let $h:C\to D$. If total relations on $D$ have fixed points, there is internally $c:C$ with $q(c)=h(c)$. In particular, for $t:C\to C$ there is internally $c$ with $q(c)=q(t(c))$.
\end{proposition}
\begin{proof}
\leavevmode
\begin{enumerate}[label=\textup{(\roman*)},leftmargin=*,itemsep=3pt,topsep=4pt]
\item \textbf{Relational transport.} Take the regular image of $(q,h):C\to D\times D$.
Its first projection is epic. A fixed point of this relation belongs to that image, hence locally comes from a $c$ with the required equation.
The assertion is existential and descends.
Equivalently, a fixed point exists for
$R(p,p')\iff\exists c\,[q(c)=p\wedge h(c)=p']$.

The defining maps satisfy the following commutative diagram
\[
\begin{tikzcd}[column sep=large,row sep=large]
C\ar[r,"{(q,h)}"]\ar[dr,"q"']&D\times D\ar[d,"\pi_1"]\\ &D
\end{tikzcd}
\]
\end{enumerate}
\end{proof}

The proposition is the representation-transport form of synthetic relational recursion; compare \cite[Proposition 5.4]{bauerfixed}. In $\Eff$, where \cref{thm:syntheticmodel} supplies its premise, it applies to every presentation of the indicated $P$, including nonnumerical presentations. Its use on another denotation object requires that object's fixed-point construction. The proposition gives an internal existence statement; a uniform compiler requires its own construction. The implementation of $\lfp$ selects the already constructed least denotation.

\paragraph{Dependency of the intrinsic development.}
The semantic function object was constructed from finite data, and its computational points were identified with the canonical admitted-ideal completion. The model theorem then supplied fixed points and continuity, from which observation fullness and effective extension followed. The forthcoming numerical calculations implement these results. None is used as a premise of \cref{thm:intrinsichalting,thm:observationfullness,thm:syntheticms}.


\part{Intensional structure and computational representation}\label{part:coding}

Extensional computability identifies what partial functions and observations mean. It does not specify how programs, proofs of definedness, or syntactic descriptions are transported. This part keeps those implementation data in regular-epimorphic presentations. Realizer exposure and description substitution provide their intensional operations, while exact interpretation relates them to the semantic maps already constructed.


The numerical program cover will also realize intensional fixed
points. The same exact interpretation equations will subsequently govern
oracle localization and the computational structure sheaf.

\section{Tracked presentations and realizer exposure}\label{sec:pres}

The extensional semantics specifies maps between completed computational objects, but it does not display a chosen implementation for each map. A name presentation records that additional data and its kernel pair records when implementations agree in denotation. We first establish how such presentations vary with the semantic object, then construct the realizer exposure that acts on their trackers.


\begin{definition}
Let $\Pres(\E)$ have as objects regular epimorphisms $q:C\onto X$ and as arrows commuting squares. Its projection to $\E$ sends $q$ to $X$. Elements of $C$ are \textbf{names} for elements of $X$; the kernel pair $C\times_X C$ records equality of denotations.
\end{definition}

\begin{proposition}\label{prop:fibration}
The projection $\Pres(\E)\to\E$ is a Grothendieck fibration. The cartesian lift of $f:Y\to X$ at $q$ is the pullback presentation
\[
 Y\times_X C\onto Y.
\]
\end{proposition}
\begin{proof}
\leavevmode
\begin{enumerate}[label=\textup{(\roman*)},leftmargin=*,itemsep=3pt,topsep=4pt]
\item \textbf{Pullback construction.} Regular epimorphisms are pullback-stable.
The universal property of the pullback is exactly the cartesian-arrow property.
The canonical comparison between an iterated pullback and the pullback along a composite satisfies the coherence equations because both sides are the unique map with the prescribed projections.

The construction uses the cartesian square
\[
\begin{tikzcd}[column sep=large,row sep=large]
Y\times_XC\ar[r,"\pi_C"]\ar[d,two heads,"\pi_Y"']&C\ar[d,two heads,"q"]\\ Y\ar[r,"f"']&X
\end{tikzcd}
\]
\end{enumerate}
\end{proof}

An arrow in the fibre over $X$ translates names while preserving their interpretation. A square over a map $f:X\to Y$ implements $f$ at the chosen presentations. The existence of such a square is additional information about those presentations; it does not follow from the existence of the two covers alone.

\begin{proposition}[Canonical implementation after refinement]
\label{prop:refined-implementation}
For $q:C\onto X$, $r:D\onto Y$ and $f:X\to Y$, take the
cartesian square
\[
\begin{tikzcd}[column sep=large,row sep=large]
 P_f=C\times_YD\ar[r,"{\pi_D}"]\ar[d,two heads,"{\pi_C}"']&D\ar[d,two heads,"r"]\\
 C\ar[r,"fq"']&Y.
\end{tikzcd}
\]
Thus $r\pi_D=fq\pi_C$, and the first projection is a regular
epimorphism. It is universal among implementations
over a refinement of $C$. If $C,D$ are assemblies, its canonical
partitioned cover gives a tracked implementation on a partitioned
refinement. These constructions do not require a section of either
name cover or external choice.
\end{proposition}
\begin{proof}
\leavevmode
\begin{enumerate}[label=\textup{(\roman*)},leftmargin=*,itemsep=3pt,topsep=4pt]
\item \textbf{Canonical implementation after refinement.} Pullback stability makes $\pi_C$ regular epic.
The pullback universal
property gives the displayed equality and the unique map from every
other pair of maps satisfying it.
When $C,D$ are assemblies, $P_f$
is a subobject of $C\times D$ and hence an assembly.
Its canonical
partitioned cover from \cref{rem:external-projectivity}, followed by
the two projections, supplies the tracked refinement.
\item \textbf{Pullback construction.} For a second presentation $s:E\onto Z$ and a map $g:Y\to Z$,
compose the spans by $P_f\times_D P_g$.
Its projection to $P_f$
is regular epic, and its two outside projections satisfy
$s\pi_E=gfq\pi_C$.
It therefore refines the direct pullback
$P_{gf}$. In the diagram below, $R_f=P_f$, $R_g=P_g$, $h$ is the output projection of $P_f$, and $p'_D$ is the source projection of $P_g$.
Associativity and parameter substitution are the canonical
pullback-pasting isomorphisms; their coherence follows from equality
of all projections.
\item Finally $q\pi_C$ is regular epic, so the
denotation is exactly $f$, and descent through its kernel pair is
unique by \cref{lem:localrepresentatives}.

The construction uses the cartesian square
\[
\begin{tikzcd}[column sep=large,row sep=large]
R_f\times_D R_g\ar[r,"\operatorname{pr}_2"]\ar[d,"\operatorname{pr}_1"']&R_g\ar[d,"p'_D"]\\ R_f\ar[r,"h"']&D
\end{tikzcd}
\]
\end{enumerate}
\end{proof}

\begin{theorem}[Choice-free semantic recovery by refinement spans]
\label{thm:refinement-quotient}
Let $\mathcal T$ be an exact category. Form a category whose objects
are regular epimorphisms $q:C\onto X$ and whose arrows from $q$ to
$r:D\onto Y$ are equivalence classes of spans
\[
\begin{tikzcd}[column sep=large,row sep=large]
 &R\ar[dl,two heads,"p"']\ar[dr,"h"]&\\
 C\ar[d,two heads,"q"']&&D\ar[d,two heads,"r"]\\
 X\ar[rr,"f"']&&Y
\end{tikzcd}
\]
where $rh$ descends uniquely
through the regular epimorphism $qp:R\onto X$ to a map $f:X\to Y$.
Two spans are equivalent when, after a common regular-epimorphic
refinement, their output maps agree after interpretation in $Y$.
Composition is pullback over the middle name object. The denotation
functor to $\mathcal T$ is an equivalence. It is compatible with
parameter pullback; for assemblies it admits tracked representatives
after a partitioned refinement. No external choice or fixed-carrier
projective lifting is used.
\end{theorem}
\begin{proof}
\leavevmode
\begin{enumerate}[label=\textup{(\roman*)},leftmargin=*,itemsep=3pt,topsep=4pt]
\item \textbf{Fullness by cartesian implementation.} For any $f:X\to Y$, the cartesian span $P_f$ of
\cref{prop:refined-implementation} represents $f$. This proves
fullness. The identity is the span with both legs $\id_C$.
\item \textbf{Faithfulness by kernel-pair descent.} If $(p,h)$ and $(p',h')$ both descend to $f$, their fibre
product $R\times_C R'$ is regular epic over each source. The two
resulting maps to $Y$ are both $fqp\pr_1=fqp'\pr_2$;
hence their pair factors through $D\times_YD$. Conversely, this
kernel-pair factorization on a common refinement, followed by
epicity over $X$, proves equality of the descended maps.
Thus denotation is faithful.
\item \textbf{Composition by pullback.} Composing $(p,h):C\leftarrow R\to D$ and
$(p',h'):D\leftarrow R'\to E$ gives
$C\leftarrow R\times_D R'\to E$. Its left leg is regular
epic by pullback stability and composition. The interpreted map is
the composite of the two denotations. Pullback associators and their
universal projections prove associativity and independence of common
refinements. The same argument in each slice proves reindexing
coherence.
\item \textbf{Essential surjectivity and tracked refinement.} Every object of $\mathcal T$ has its identity presentation,
so denotation is essentially surjective. In $\Eff$, the canonical
partitioned cover of an assembly fibre product turns any such span
into a tracked refinement, as in \cref{prop:refined-implementation}.
The construction is the regular-epimorphic instance of the
cover-span calculus, cf.\ \cite{makkai}.
\end{enumerate}
\end{proof}

\begin{theorem}[Exact partial maps by covered spans]
\label{thm:partial-refinement-quotient}
There is a restriction category of regular-epimorphic name presentations in $\Eff$ whose denotation functor is an equivalence with $\Par_\Sigma(\Eff)$, without external choice. Its objects are covers $q:C\onto X$. A representative from $q$ to $r:D\onto Y$ is a span
\[
 C\xleftarrow{p}R\xrightarrow{h}D
\]
with the following intrinsic properties: $A=\im(p)\incl C$ is $\Sigma$-classified and is saturated under $C\times_XC\rightrightarrows C$, and $rh$ equalizes the kernel pair of $qp$. Two representatives are identified when their source images agree and their interpreted outputs agree on $R\times_C R'$. Composition is pullback over the middle name object. Every semantic partial map has a canonical representative, and these constructions preserve exact domains and parameter substitution. For assembly name objects, every arrow class has a tracked representative after a canonical partitioned refinement.
\end{theorem}
\begin{proof}
\leavevmode
\begin{enumerate}[label=\textup{(\roman*)},leftmargin=*,itemsep=3pt,topsep=4pt]
\item \textbf{Canonical construction on the actual domain.} Write a semantic partial map as $X\xleftarrow{u}U\xrightarrow{v}Y$, with $u$ a $\Sigma$-classified mono. Put $A=C\times_XU$ and $R_f=A\times_YD$, using $v\pi_U$ for the map $A\to Y$. Thus
\[
\begin{tikzcd}[column sep=large,row sep=large]
 R_f\ar[r,"h"]\ar[d,two heads]&D\ar[d,two heads,"r"]
 &A\ar[r,"\pi_U"]\ar[d,hook]&U\ar[d,hook,"u"]\\
 A\ar[r,"{v\pi_U}"']&Y
 &C\ar[r,two heads,"q"']&X.
\end{tikzcd}
\]
Both squares are pullbacks. The first left leg is epic, so $\im(p)=A$ exactly. The second square proves saturation and positivity of $A$. Finally $rh=v\pi_U\pi_A$ depends only on the $U$-denotation, so it equalizes the kernel pair of $qp$.
\item \textbf{Intrinsic descent of the domain.} Conversely let a representative be given and put $U=\im(A\to X)$. In internal regular logic,
\[
 q(c)\in U\ \Longleftrightarrow\
 \exists a\in A\,[q(a)=q(c)]\ \Longleftrightarrow\ c\in A.
\]
The forward implication uses saturation and the reverse uses $a=c$. Hence $A=q^{-1}U$. The $\Sigma$-classifier of $A$ agrees on the kernel pair of $q$: the equality holds after $\Sigma\incl\Omega$, and this inclusion is monic. Effective descent along $q$ gives a classifier $X\to\Sigma$ for $U$. The composite $R\onto A\onto U$ is epic. Since $U\incl X$ is monic, its kernel pair is that of $qp$; the required equality of $rh$ therefore descends uniquely to a value map $U\to Y$.
\item \textbf{Equality of representatives.} Equal source images descend to equal subobjects of $X$, since pullback along $q$ reflects equality of subobjects. The projections from $R\times_C R'$ cover both representatives over their common image $A$. Agreement of their interpreted outputs there is consequently equivalent to equality of the descended maps $U\to Y$. Conversely equal partial denotations make those outputs agree on the displayed pullback. Thus the stated relation is an equivalence relation, and the arrow classes are in bijection with the actual partial maps, including their domains.
\item \textbf{Composition and coherence.} Suppose $S$ represents $g:Y\pmap Z$, with left leg $p':S\to D$. The composite representative is $R\times_D S$, formed using $h$ and $p'$. If $f$ has domain $U$ and $g$ has domain $V$, saturation of $\im(p')=r^{-1}V$ and epicity onto that image give
\[
 \im(R\times_D S\to C)
 =\{c\mid q(c)\in U\ \wedge\ v(q(c))\in V\}.
\]
Indeed, an element of the left-hand side supplies these two facts. Conversely, over such a $c$ the cover $R\onto A$ supplies a local output name $d$; its denotation lies in $V$, so $S\onto r^{-1}V$ supplies the next local witness. Their pullback gives the required image membership. Dominance makes this composite domain $\Sigma$-classified. The output descends to the usual composite value. Pullback associators preserve every projection; hence their pentagon and unit equations follow by uniqueness. The same argument proves independence of representatives and parameter coherence.
\item \textbf{Restriction and equivalence.} The identity is $C\xleftarrow{\id}C\xrightarrow{\id}C$. The restriction of $(p,h)$ is $(p,p):C\leftarrow R\to C$. It denotes the partial identity on $U$. The restriction equations therefore hold by the faithful domain--value correspondence just proved; this is the restriction-category structure of \cite{cockettlack}. The canonical representative in part~(i) proves fullness. Every $X$ has its identity presentation, and the total identity of $X$ has canonical representatives in both directions between $q$ and $\id_X$; their composites denote the identities. This gives the equivalence explicitly, without choosing representatives from arbitrary nonempty fibres. The total part recovers \cref{thm:refinement-quotient} and the cover-span viewpoint of \cite{makkai}.
\item \textbf{Tracked refinement.} When $C,D$ are assemblies, $R_f$ is a subobject of $C\times D$ and hence an assembly. Its canonical all-realizer cover is partitioned and the two projections are tracked. Precomposing by this cover preserves the source image and value, so it gives the asserted tracked representative. This construction uses neither a section of a name cover nor preservation of covers by the partial-map classifier.
\end{enumerate}
\end{proof}

\begin{proposition}[Finite execution by covered spans]
\label{prop:refined-finite-returns}
The covered-span presentation of \cref{thm:partial-refinement-quotient} constructs finite-return witnesses for every $f:X\pmap Y+X$ on arbitrary name covers. Their image is exactly the pullback of the domain of $f^\ddagger$, and their output descends to its value. The same construction presents coproduct feedback and commutes with parameters and common refinement, without external choice.
\end{proposition}
\begin{proof}
\leavevmode
\begin{enumerate}[label=\textup{(\roman*)},leftmargin=*,itemsep=3pt,topsep=4pt]
\item \textbf{Witness lists.} Use covers $q_X:C_X\onto X$ and $q_Y:C_Y\onto Y$ and construct the canonical representative $C_X\leftarrow R\to C_Y+C_X$. Pulling back the two coproduct injections splits it as $R=R_Y+R_X$. Retain the source map to $C_X$, and the output maps $R_Y\to C_Y$ and $R_X\to C_X$. Put
\[
 W_0=R_Y,\qquad W_{n+1}=R_X\times_{C_X}W_n,
\]
where the pullback matches the continuation output with the next source name. The output of $W_{n+1}$ is the output of its $W_n$ factor. This family is constructed internally as composable finite lists in $R$ whose last member lies in $R_Y$ and whose preceding members lie in $R_X$. Thus $\sum_{n:\N}W_n$ is the dependent sum of one internal family, not a choice of external witnesses at each length.
\item \textbf{Exact finite domains.} Let $R_n(f)\incl X$ and $v_n:R_n(f)\to Y$ be the finite-return objects of \cref{exec:lem:return}. Interpreting the witness list gives a commuting diagram
\[
\begin{tikzcd}[column sep=large,row sep=large]
 W_n\ar[r]\ar[d,two heads]&C_Y\ar[d,two heads,"q_Y"]\\
 C_X\times_XR_n(f)\ar[r,"{v_n\pi_2}"']&Y.
\end{tikzcd}
\]
The asserted left epimorphism is proved by induction. At $n=0$, it is the pullback of $q_Y$ in the return summand of the canonical representative. At the successor step a semantic continuation has, locally along $q_X$, a continuation name. The inductive cover supplies the rest of the list over that very name. Pullback and composition of these covers prove epicity, without selecting either witness globally. Consequently
\[
 \im(W_n\to C_X)=q_X^{-1}R_n(f).
\]
\item \textbf{The whole returning computation.} Write $W=\sum_{n:\N}W_n$. Images commute with this dependent existential, so
\[
 \im(W\to C_X)
 =q_X^{-1}\!\left(\bigvee_{n:\N}R_n(f)\right)
 =q_X^{-1}(\dom f^\ddagger).
\]
Each finite domain is $\Sigma$-classified, and $\N$-overtness classifies their union. The monicity and disjointness proved in \cref{exec:lem:return} show that equal initial semantic states in returning lists have the same length and the same terminal value. Hence the output $W\to C_Y\to Y$ equalizes the kernel pair of the initial denotation. By \cref{thm:partial-refinement-quotient}, the span $C_X\leftarrow W\to C_Y$ therefore denotes precisely $f^\ddagger$, including its exact domain.
\item \textbf{Feedback, parameters, and refinement.} Apply the same construction to the tagged transition defining $\Tr_+$, and restrict the initial state to its external-input summand. Its returning lists are exactly the feedback paths, so the same domain and value calculation proves the feedback assertion. Parameter pullback preserves the list construction, each finite pullback, and its image. Two representatives of the same transition have the same interpreted finite-return strata by induction; the preceding calculation therefore identifies their iteration and feedback classes. For assembly names the resulting witness objects admit canonical partitioned covers. This supplies the positive witnesses of \cref{exec:cor:observe} at arbitrary presentations and transports the Elgot and trace equations of \cref{exec:thm:elgot,exec:thm:trace} through the faithful equivalence.
\end{enumerate}
\end{proof}

\begin{lemma}\label{lem:projectivelift}
If $C$ is regular-projective, every $f:X\to Y$ admits an implementation between $q:C\onto X$ and any $r:D\onto Y$. Any two implementations have equal denotations, and hence their pair factors through $D\times_Y D$.
\end{lemma}
\begin{proof}
\leavevmode
\begin{enumerate}[label=\textup{(\roman*)},leftmargin=*,itemsep=3pt,topsep=4pt]
\item \textbf{Kernel-pair descent.} Lift $fq$ through the regular epimorphism $r$.
The second assertion is the defining property of the kernel pair of $r$.
\end{enumerate}
\end{proof}

Fixed-source lifting uses the projectivity specified in
\cref{rem:external-projectivity}; \cref{lem:modest-target-lift}
constructs it for modest target name assemblies without external choice.
The full fibration supplies cartesian reindexing. A projective cover of
a pulled-back presentation supplies implementation data when needed.
Composition retains the actual name map and tracker; its semantic
equations descend through the interpretation cover.

\subsection{Change of presentation and contextual codes}
The fibre $\Pres(\E)_X$ is the category of regular-epimorphic presentations of a fixed denotation object. If $q:C\onto X$ and $r:D\onto X$, their pullback
\[
 C\xleftarrow{\pi_C}C\times_XD\xrightarrow{\pi_D}D
\]
is a common refinement: both projections are regular epic and both interpretations into $X$ agree. When enough regular projectives are available, a projective cover of this pullback gives a common projective refinement. This is the useful comparison when neither of the original presentations carries a preferred translation to the other.

\begin{proposition}\label{prop:presentationdescent}
For presentations $q:C\onto X$ and $r:D\onto Y$, a map $u:C\to D$ descends to a map $X\to Y$ if and only if $ru$ equalizes the kernel pair of $q$. Its descent is unique. Two such maps $u,v$ have the same descent if and only if $(u,v):C\to D\times D$ factors through $D\times_YD$.
\end{proposition}
\begin{proof}
\leavevmode
\begin{enumerate}[label=\textup{(\roman*)},leftmargin=*,itemsep=3pt,topsep=4pt]
\item \textbf{Kernel-pair descent.} Regular epimorphisms in a topos are coequalizers of their kernel pairs.
Hence the equalizing condition is precisely the existence condition for a unique $f:X\to Y$ with $fq=ru$.
If $u,v$ descend to $f,g$, then $f=g$ if and only if $fq=gq$, since $q$ is epic.
This is the equation $ru=rv$, whose universal factorization is the indicated kernel pair.
No comparison of the names themselves is required.
\end{enumerate}
\end{proof}

A contextual name for $x:S\to X$ is a section of $x^*C\onto S$.
Such sections exist after projective refinement and agree in denotation
on common refinements. The kernel-pair criterion makes semantic equality
a congruence for implemented composition; projective lifting makes its
quotient full. Numerical covers below provide explicit compiler maps
within this indexed structure.

\subsection{Effective descent and twisted name presentations}
\label{obs:subsec:presentation-descent}
The same construction defines $\Pres(\mathcal T)$ in each of the
topoi constructed below. Its fibres admit effective descent along
epimorphisms. This strengthens parameter reindexing: compatible
local naming systems reconstruct an actual global naming system.

\begin{proposition}[Descent of names and implemented diagrams]
\label{obs:prop:descent}
For an epimorphism $e:U\onto X$ in a topos $\mathcal T$, pullback
identifies $\Pres(\mathcal T)_X$ with the category of presentations
$C_U\onto U$ equipped with an isomorphism between their two
pullbacks to $U\times_XU$, satisfying the cocycle equation.
Morphisms are the maps respecting that equation. Finite diagrams
of names, interpretation maps, and partial domain--value spans
descend with their equations. In particular, the groupoids of
presentations and invertible translations form a stack for the
epimorphism coverage.
\end{proposition}
\begin{proof}
\leavevmode
\begin{enumerate}[label=\textup{(\roman*)},leftmargin=*,itemsep=3pt,topsep=4pt]
\item \textbf{Unit and identity laws.} Write $\theta_{u,v}$ for the descent isomorphism, directed from
the fibre over $u$ to the fibre over $v$.
On $C_U$ define
\[
 c\sim d\quad\Longleftrightarrow\quad
 e(q_Uc)=e(q_Ud)\ \text{ and }\
 \theta_{q_Uc,q_Ud}(c)=d.
\]
The identity, inverse, and cocycle equations make this an
equivalence relation.
Exactness of a topos supplies its quotient
$C$.
The map $eq_U$ descends to $q:C\to X$, and $q$ is epic
because $eq_U$ is epic.
The canonical map
$C_U\to U\times_XC$ is monic by the identity equation.
\item \textbf{Inverse comparison.} It is epic:
locally represent a class by $d$ over $v$ and transport it to $u$
by $\theta_{v,u}$.
Thus it is an isomorphism.
This also proves
that the reconstructed descent isomorphism is the specified one.
\item \textbf{Image and quotient construction.} A compatible local map preserves the displayed equivalence
relations and therefore descends uniquely through their quotients.
This proves full faithfulness and essential surjectivity, without
selecting representatives.
Apply this construction to each object
and arrow of a diagram.
Composition and equality descend because
they can be checked after an epimorphism.
Finite limits and
monomorphisms are likewise checked after this pullback, so a
partial map retains its actual domain mono and value map.
\item \textbf{Domain calculation.} Domain
classification descends by descending the corresponding
classifying map.
Consequently implemented operations and their
domain-and-value equations descend as diagrams; a name cover alone
is not being assigned extra operations.
Restricting the resulting
equivalence to invertible arrows gives the asserted stack.
This is effective descent in the topos-theoretic sense
\cite{johnstone,maclanemoerdijk}, here applied to the actual
presentation diagrams.

The construction uses the cartesian square
\[
\begin{tikzcd}[column sep=large,row sep=large]
C_U\ar[r,"\cong"]\ar[d,two heads,"q_U"']&U\times_XC\ar[d,two heads,"\pi_1"]\\ U\ar[r,"\id_U"']&U
\end{tikzcd}
\]
\end{enumerate}
\end{proof}

\begin{proposition}[Forms of a naming system]
\label{obs:prop:forms}
Fix a presentation $q:C\onto X$ and work in $\mathcal T/X$.
The internal group $G=\operatorname{Aut}_X(C)$ classifies its
locally isomorphic forms: the groupoid of such presentations is
equivalent to the groupoid of right $G$-torsors. A torsor $T$ gives
the associated presentation
\[
 C_T=T\times^G C=(T\times_X C)/G\ \onto X,
 \qquad (t,c)g=(tg,g^{-1}c).
\]
For a displayed finite implemented signature, the same statement
holds with the automorphism group of its whole diagram, including
the partial domains and interpretation maps.
\end{proposition}
\begin{proof}
\leavevmode
\begin{enumerate}[label=\textup{(\roman*)},leftmargin=*,itemsep=3pt,topsep=4pt]
\item \textbf{Torsor comparison.} Construct $G$ inside the product of two endomorphism objects by
the equations $uv=vu=\id$; projection to either factor identifies
it with the invertible endomorphisms.
Composition and inversion
give its group structure. A right torsor, as in the descent formalism \cite{johnstone,maclanemoerdijk}, means an epimorphism
$T\onto X$ with
\[
 T\times_XG\longrightarrow T\times_XT,\qquad
 (t,g)\longmapsto(t,tg)
\]
an isomorphism.
Quotient the displayed diagonal action by its
orbit equivalence relation.
After pullback along $T\onto X$, a
chosen local torsor element identifies the quotient with $C$.
The associated map to $X$ is therefore epic and is locally the
original presentation.
\item \textbf{Converse construction.} Conversely, for a locally isomorphic presentation $C'\onto X$,
the object $T=\operatorname{Isom}_X(C,C')$ is constructed from
pairs of inverse maps in the slice.
It is locally inhabited and
has the right action $\phi g=\phi\circ g$.
Two isomorphisms differ
by a unique element of $G$, proving the torsor equation.
Evaluation descends to
$T\times^GC\to C'$ and is an isomorphism after the cover $T$,
hence globally.
\item \textbf{Further converse construction.} In the other direction,
$t\mapsto(c\mapsto[t,c])$ identifies a torsor with the
isomorphism torsor of its associated form.
These maps are
equivariant, natural, and inverse locally; descent proves their
inverse identities and coherence globally.
\item \textbf{Kernel-pair descent.} For the implemented signature take the subgroup of the product
of sort-automorphism groups that commutes with every displayed
arrow and preserves each domain subobject.
These are equalizer
conditions on actual diagrams.
Twisting every sort by this same
torsor and descending the equivariant arrows gives the required
operations.
Their interpretation and domain equations hold on
the trivializing cover and hence globally by
\cref{obs:prop:descent}.
\item \textbf{Identification of the resulting maps.} This classifies the forms of the
specified signature; it makes no identification of arbitrary
name covers with one another.

The defining maps satisfy the following commutative diagram
\[
\begin{tikzcd}[column sep=large,row sep=large]
T\times^G C\ar[r,"\cong"]\ar[d,two heads]&C'\ar[d,two heads]\\ X\ar[r,"\id_X"']&X
\end{tikzcd}
\]
\end{enumerate}
\end{proof}


\subsection{Tracked assemblies and realizer exposure}\label{int:sec:exposure}
The presentation fibration retains names over denotations. To retain the
implementation of a name transformation, we also keep its tracker.
We use the P-categorical formulation of exposures of
\cite[Sections~2--3 and~9]{kavvosarticle}. The choices of terminal and
product presentations below make the cartesian comparisons explicit.

\begin{definition}[The P-category of tracked assemblies]
\label{int:def:tracked}
Write $a\Vdash_X x$ for the realization relation of an assembly over
$K_1$. The P-category $\TrAsm$ has these assemblies as objects and
pairs $(f,r):X\to Y$ as arrows, where $f:|X|\to|Y|$ and
\[
 a\Vdash_Xx\quad\Longrightarrow\quad
 ra\downarrow\ \text{and}\ ra\Vdash_Yf(x).
\]
Its extensional equality is $(f,r)\sim(g,s)$ if $f=g$.
Identity is tracked by $\mathsf i$, and composition by
$(g,s)(f,r)=(gf,[a]s(ra))$. Thus the identity and associativity
equations hold up to $\sim$, and composition preserves $\sim$.
An exposure preserves identities and composition up to $\sim$ and
reflects $\sim$; preservation of $\sim$ is not part of its definition.
For an exposure $\RealQ$, write $u\equiv_{\RealQ}v$ when
$\RealQ u\sim\RealQ v$.
\end{definition}

\begin{construction}[Chosen cartesian presentations]
\label{int:con:cartesian}
Choose a computable bijective pairing of natural numbers with computable
projections, and make the following presentations.
\begin{enumerate}[label=\textup{(\roman*)}]
\item The terminal assembly $\mathbf1_*$ has carrier $\{*\}$ and
unique realizer $0$. The constant-zero program tracks its terminal map.
\item The carrier of $X\times Y$ is $|X|\times|Y|$, with
\[
 \langle a,b\rangle\Vdash_{X\times Y}(x,y)
 \quad\Longleftrightarrow\quad a\Vdash_Xx\ \wedge\ b\Vdash_Yy.
\]
Use the numerical projections as trackers of the projections, and use
$[c]\langle rc,sc\rangle$ for pairing tracked maps $(f,r)$ and $(g,s)$.
\item Product maps, associators, symmetries, and unitors use the
corresponding pairing, reassociation, exchange, and projection programs.
\end{enumerate}
The product universal property holds up to $\sim$, since its underlying
carrier is the ordinary product. The displayed trackers prove existence
of every map used in that universal property.
\end{construction}

\begin{definition}[Retained-realizer exposure]\label{int:def:exposure}
Define
\[
 |\RealQ X|=\{(x,a):a\Vdash_Xx\},\qquad
 b\Vdash_{\RealQ X}(x,a)\quad\Longleftrightarrow\quad b=a,
\]
and put $\RealQ(f,r)(x,a)=(f(x),ra)$, with retained tracker $r$.
The counit and duplication maps, both with identity tracker, are
\[
 \varepsilon^{\RealQ}_X(x,a)=x,\qquad
 \delta_X(x,a)=((x,a),a).
\]
\end{definition}

\begin{theorem}[Constructed cartesian comonadic exposure]
\label{int:thm:exposure}
The operation $\RealQ$ is a cartesian, product-preserving exposure on
$\TrAsm$ with the presentations of \cref{int:con:cartesian}.
The maps $\varepsilon^{\RealQ}$ and $\delta$ are P-natural and satisfy the
counit and coassociativity equations. There are P-natural isomorphisms
\[
 m_{X,Y}:\RealQ X\times\RealQ Y\xrightarrow{\cong}\RealQ(X\times Y),
 \qquad m_0:\mathbf1_*\xrightarrow{\cong}\RealQ\mathbf1_*.
\]
For every tracked point $u:\mathbf1_*\to\RealQ X$, the quotation
equation $\delta_Xu\sim\RealQ u\,m_0$ holds. In particular these data
form a comonadic exposure in the sense of \cite{kavvosarticle}.
\end{theorem}
\begin{proof}
\leavevmode
\begin{enumerate}[label=\textup{(\roman*)},leftmargin=*,itemsep=3pt,topsep=4pt]
\item \textbf{Exposure and reflection.} We verify the structure on the retained realizers.
Identity sends $(x,a)$ to itself. Both composition routes send it
to $(gf(x),s(ra))$, by delayed abstraction. If two exposed carrier maps
agree, their first coordinates agree on every realized $x$; every carrier
element is realized. Thus $\RealQ$ reflects $\sim$.
\item \textbf{P-naturality and comonad equations.} The naturality diagrams are
\[
\begin{tikzcd}[column sep=large]
 \RealQ X\ar[r,"{\RealQ(f,r)}"]\ar[d,"\varepsilon^{\RealQ}_X"']&
 \RealQ Y\ar[d,"\varepsilon^{\RealQ}_Y"]\\
 X\ar[r,"{(f,r)}"']&Y
\end{tikzcd}
\qquad
\begin{tikzcd}[column sep=large]
 \RealQ X\ar[r,"{\RealQ(f,r)}"]\ar[d,"\delta_X"']&
 \RealQ Y\ar[d,"\delta_Y"]\\
 \RealQ^2X\ar[r,"{\RealQ^2(f,r)}"']&\RealQ^2Y.
\end{tikzcd}
\]
Their values are respectively $f(x)$ and $((f(x),ra),ra)$.
Both counit composites on $(x,a)$ return $(x,a)$, and both
coassociativity composites return $(((x,a),a),a)$. Hence
\[
 \varepsilon^{\RealQ}_{\RealQ X}\delta_X\sim\id,
 \quad\RealQ\varepsilon^{\RealQ}_X\delta_X\sim\id,
 \quad\delta_{\RealQ X}\delta_X\sim\RealQ\delta_X\delta_X.
\]
\item \textbf{Cartesian comparison.} Set
\[
 m_{X,Y}((x,a),(y,b))=((x,y),\langle a,b\rangle).
\]
The inverse extracts the two numerical projections. Both composites
are carrier identities, since the pairing and projections are inverse
on the chosen product realizers. The map $m_0$ sends $*$ to $(*,0)$.
Naturality follows from the chosen tracker
$\langle a,b\rangle\mapsto\langle ra,sb\rangle$ of $f\times g$.
Reassociation sends the two paths on a triple to the same nested pair;
exchange swaps its coordinates; the unitors remove the terminal $0$.
These computations prove the monoidal coherence diagrams.
Furthermore projecting a paired tracker returns its specified
component, and pairing the projections of a product realizer returns
that same realizer. Thus the product beta and eta equations hold up
to $\equiv_{\RealQ}$. Every terminal arrow sends a valid realizer
to $0$, so terminal uniqueness also holds up to
$\equiv_{\RealQ}$. This proves cartesianness.
\item \textbf{Quotation equation.} If $u(*)=(x,a)$ and $r$ tracks $u$, then $r0=a$, because
$\RealQ X$ has the unique realizer $a$ at $(x,a)$. Both sides of
$\delta_Xu\sim\RealQ u\,m_0$ therefore return $((x,a),a)$.
\end{enumerate}
\end{proof}

\begin{example}[Exposure depends on retained presentations]
\label{int:ex:terminal}
Let $T$ have one carrier point realized by every natural number.
It too is terminal up to $\sim$, but $\RealQ T$ has carrier $\N$.
The constant-zero and constant-one trackers implement the same
$\id_T$ and have different exposed carrier maps. Thus $\RealQ$
does not preserve $\sim$, nor does it preserve every extensional
isomorphism between alternative terminal presentations. This is
compatible with \cref{int:thm:exposure}, which specifies its chosen
terminal presentation. The operation is an exposure on tracked arrows,
with intensional information retained in its tracked-arrow presentation.
\end{example}

\begin{proposition}[Exposed names and change of presentation]
\label{int:prop:presentations}
\begin{enumerate}[label=\textup{(\roman*)}]
\item For a modest assembly $X_\rho$, its exposed assembly
$\RealQ X_\rho$ is canonically isomorphic to the support presentation
$P_\rho$ of \cref{prop:per}; $\varepsilon^{\RealQ}_{X_\rho}$ is its name cover.
\item For an assembly presentation $q:C\onto\Pt D$, exposed names
interpret by $q\varepsilon^{\RealQ}_C$. A tracked implementation $(\widehat f,r)$
of $F:D\to E$ gives the commuting diagram
\[
\begin{tikzcd}[column sep=large]
 \RealQ C_D\ar[r,"{\RealQ(\widehat f,r)}"]
   \ar[d,"q_D\varepsilon^{\RealQ}_{C_D}"']&
 \RealQ C_E\ar[d,"q_E\varepsilon^{\RealQ}_{C_E}"]\\
 \Pt D\ar[r,"\Pt F"']&\Pt E.
\end{tikzcd}
\]
\item These comparisons respect implemented refinements and parameter
pullback followed by projective refinement. Their semantic quotient
recovers the category in \cref{thm:generalpresentationenhancement}.
\end{enumerate}
\end{proposition}
\begin{proof}
\leavevmode
\begin{enumerate}[label=\textup{(\roman*)},leftmargin=*,itemsep=3pt,topsep=4pt]
\item \textbf{Naturality and reindexing.} A realizer of a modest assembly determines its unique realized class.
Consequently $(x,a)\mapsto a$ and $a\mapsto([a],a)$ are inverse
tracked maps, proving (i).
For (ii), naturality of $\varepsilon^{\RealQ}$ gives
$q_E\varepsilon^{\RealQ}_{C_E}\RealQ(\widehat f,r)
=q_E\widehat f\varepsilon^{\RealQ}_{C_D}=\Pt Fq_D\varepsilon^{\RealQ}_{C_D}$.
For (iii), apply this equation to a refinement $u$ with $q'u=q$.
The exposure composition law compares successive refinements.
\item \textbf{Pullback construction.} Two
presentations admit a projective cover of their fibre product; the
same construction applies to pulled-back presentations.
Equalities
of interpretations descend along these covers.
Fullness of the
semantic quotient is projective implementation, and its faithfulness
is equality of the retained semantic arrow $F$.
On searchable domains
the point function detects this equality; on arbitrary observed domains
the arrow $F$ remains part of the presentation.
\item Finally,
$\varepsilon^{\RealQ}_C$ is the regular-epimorphic partitioned cover of
the assembly $C$, so passing to exposed names preserves the cover.
\end{enumerate}
\end{proof}

For a partitioned assembly $C$, $\RealQ C\cong C$: its name already
specifies its unique realizer. Nontrivial exposure on modest semantic
objects is supplied by (i). On arbitrary exact-completion objects the
construction acts on the displayed assembly presentations, retaining
their semantic arrows. This specifies the domain of the exposure
without treating every semantic object as an assembly.


\subsection{Domain presentations and extensional recovery}
Pull the name fibration back along the effective-point functor:
\begin{equation}\label{eq:presentationfibration}
\begin{tikzcd}[column sep=large]
 \PresPt(\C)\ar[r]\ar[d,"{\pi}"']&\Pres(\E)\ar[d,"{\operatorname{cod}}"]\\
 \C\ar[r,"{\Pt}"']&\E.
\end{tikzcd}
\end{equation}
A presented domain is $(D,q:C\onto\Pt D)$. An implementation arrow over $F:D\to E$ is an implementation $\widehat F:C\to C'$ satisfying
\begin{equation}\label{eq:implement}
\begin{tikzcd}[column sep=large]
 C\ar[r,"{\widehat F}"]\ar[d,two heads,"{q_D}"']&C'\ar[d,two heads,"{q_E}"]\\
 \Pt D\ar[r,"{\Pt F}"']&\Pt E.
\end{tikzcd}
\end{equation}
Cartesian lifts are the pullbacks along $\Pt F$. Projective implementations use projective refinements of those lifts.

\begin{definition}
An \textbf{implemented presentation of a computational signature}
consists of name presentations for its semantic sorts and selected tracked
implementations of its operations. A partial implementation has exactly
the pulled-back semantic domain and the specified interpreted value.
Its carrier map and tracker are retained in $\TrAsm$. Semantic equality
is equality after interpretation, through the target kernel pair;
it is the quotient used in the following theorem.
\end{definition}

\begin{theorem}[Extensional recovery from tracked presentations]
\label{thm:presentationconservative}\label{thm:generalpresentationenhancement}
For a full effective-domain category $\mathcal D\subseteq\ED$, use
projective presentations of $\Pt D$ and retain arrows $(F,\widehat F)$
satisfying \eqref{eq:implement}. Then:
\begin{enumerate}[label=\textup{(\roman*)}]
\item Quotienting by equality of the semantic component $F$ gives
an equivalence with $\mathcal D$.
\item On searchable domains equality of point denotations already
detects this quotient, so
$\PresPt^{\mathrm{proj}}(\Search)/{\sim}\simeq\Search$.
\item Specified semantic operations lift on projective presentations
of their actual input and output objects. In particular,
\begin{equation}\label{eq:lfpimplementation}
\begin{tikzcd}[column sep=large]
 C_X\ar[r,"\widehat{F^\dagger}"]\ar[d,two heads,"q_X"']&
 C_D\ar[d,two heads,"q_D"]\\
 \Pt X\ar[r,"\Pt(F^\dagger)"']&\Pt D.
\end{tikzcd}
\end{equation}
The equations are invariant under common refinement.
\item Without fixed-source projectivity, all of these semantic maps
have choice-free implementations after a regular-epimorphic
refinement, with the same semantic quotient.
\end{enumerate}
\end{theorem}
\begin{proof}
\leavevmode
\begin{enumerate}[label=\textup{(\roman*)},leftmargin=*,itemsep=3pt,topsep=4pt]
\item \textbf{Image and quotient construction.} The choice-free assertion (iv) is \cref{thm:refinement-quotient}
applied to the actual input and output point objects.
For the
fixed-source formulation (i)--(iii), the selected projective covers
are used as follows.
Every point object has a projective cover, proving essential
surjectivity.
Lift $\Pt Fq_D$ through $q_E$ to prove fullness.
Equality of retained semantic components makes the quotient faithful.
On searchable domains, equal name denotations give equal point maps
by epicity of $q_D$, and \cref{prop:observationseparation} gives
$F=G$.
\item \textbf{Pullback construction.} Finite products of projective input covers are projective;
alternatively cover the actual input object.
Applying the same lift
to the already constructed least-solution map proves
\eqref{eq:lfpimplementation}.
Two lifts pair into the target kernel
pair. A projective cover of the fibre product compares two systems
of names, where their interpreted equations agree and descend.
\item \textbf{Universal-property argument.} Fixed-source
lifting of arbitrary partitioned assemblies here uses the external
metatheory recorded in \cref{rem:external-projectivity}; it is not
needed for the covered-span form of the result.
\end{enumerate}
\end{proof}

The general presentation retains $F$ and its observation-preservation
data, as required by \cref{ex:pointforgetting}. For a lifted type,
cover its actual point object and use
$\Pt(\Lift D)\cong\LL\Pt D$; no preservation of covers by $\LL$
is needed. Leastness in \eqref{eq:lfpimplementation} comes from the
semantic least-solution construction. Its positive observations have
the finite-iterate witnesses of \eqref{eq:lfpobs}, implemented by
the dovetailing construction of \cref{prop:lfpparameters}.


\subsection{Exact implementation in context}
For program and data covers $q_{\mathrm{prog}},q_X,q_Y$, the
domain--value construction of \cref{thm:partialadequacy} gives
\begin{equation}\label{eq:adequacysquare}
 q_YU_{\mathrm{code}}
 =U_{\mathrm{sem}}(q_{\mathrm{prog}}\times q_X)
 \quad\text{in }\Par_\Sigma(\Eff).
\end{equation}
Thus the implementation domain is the pullback of the semantic domain.
On a partitioned source with modest target names, \cref{prop:modest-partial-implementation} lifts the classifier and value and constructs first-hit search. For arbitrary selected target names the fixed-carrier extraction uses the projectivity of \cref{rem:external-projectivity}; \cref{thm:partial-refinement-quotient} constructs the exact covered-span implementation without it. This constructs the
evaluator and its exact extent. For a parameter map, pull back this
span and refine its name cover projectively. Iterated parameter changes
have the canonical pullback-pasting comparison. On a common refinement
their classifiers and values agree, so the equation descends.
\Cref{int:thm:exactsyntax,int:thm:refinement} extends this single
primitive adequacy square to the entire free term monad.

\section{Exact implementations of the computational signature}\label{sec:canonicalcoding}

Realizer exposure becomes useful only when its action is compatible with the operations of the semantic category. We build names for these operations and compare their domain--value spans with the corresponding semantic maps. Exactness here means equality after interpretation, not an identification of distinct program codes.


Let $\mathcal C_{\mathrm{ad}}=\mathbf{Search}_{\Sigma,\mathrm{mod}}$
be the characterized semantic category. For its modest point assembly
$X_\rho$, \cref{int:prop:presentations} identifies the canonical
support cover with $\varepsilon^{\RealQ}_{X_\rho}:\RealQ X_\rho\onto X_\rho$.
For a completion $(B,r)$, a supported name comprises a saturated
membership stream and realizers of inhabitedness, lower closure,
directedness, and rounding. Equality is equality of its hit support.
\Cref{thm:modestpoints} identifies this quotient with $\Pt D_r$.
Projective stream lifts and witness extraction supply local names;
validity is the supported-name predicate on the numerical presentation.

\begin{theorem}[Canonical coding enhancement]\label{thm:canonicalcoding}
The effective-topos realizability characterization constructs the following data.
\begin{enumerate}
\item For each characterized domain a name cover $q_D:C_D\onto\Pt D$, canonical from its PER presentation. Together with all their reindexings these belong to a presentation fibration over $\mathcal C_{\mathrm{ad}}$.
\item For a specified computational signature interpreted by these domains, implementations of its total operations and its $\Sigma$-partial operations. The signature may include typed evaluation, composition, specialization, observations, the least-solution operators, Elgot iteration, and the feedback maps constructed in Part~\ref{part:synthetic}.
\item An equivalence
\[
 \PresPt^{\mathrm{proj}}(\mathcal C_{\mathrm{ad}})/{\simeq_{\mathrm{den}}}
 \simeq\mathcal C_{\mathrm{ad}},
\]
where parallel implementations are equivalent exactly when they have the same interpreted map. Adequacy includes equality of classified domains for partial operations. These comparisons are stable under parameter substitution and common refinement of presentations.
\item The same interpreted operations have implementations by covered
spans on arbitrary regular-epimorphic name presentations; their
refinement quotient is the semantic category without external choice.
For canonical modest target name covers, the fixed-carrier compiler maps in (2) are constructed without external choice by \cref{prop:modest-partial-implementation}. For arbitrary nonmodest target name covers their fixed-carrier extraction uses the explicitly selected projective covers.
\end{enumerate}
\end{theorem}
\begin{proof}
\leavevmode
\begin{enumerate}[label=\textup{(\roman*)},leftmargin=*,itemsep=3pt,topsep=4pt]
\item \textbf{Inverse comparison.} The preceding support construction gives the covers on the modest sorts.
For parameter sorts of the computational signature, take a partitioned-assembly cover of the actual interpretation object in $\E$.
The displayed category equivalence concerns the characterized domain sorts; parameter presentations also cover the hom-objects used by iteration. For a total operation $f$ with modest target name cover, \cref{lem:modest-target-lift} constructs the lift on the partitioned product of the input covers, giving $q_Y\widehat f=fq_X$. An arbitrary nonmodest target cover uses the stated external projectivity instead. No modestness of an arbitrary parameter or hom-object is inferred merely from its being a parameter sort.
\item \textbf{Partiality and definedness.} For a partial operation $f:X\pmap Y$, pull its classified domain back along $q_X$.
Lift that classifier to a stream on $C_X$ through the modest stream object. If the output name assembly is modest, \cref{prop:modest-partial-implementation} lifts the value and constructs first-hit execution without external choice. For an arbitrary selected target cover, projectivity of the hit assembly supplies the fixed-target value lift. First-hit search then gives the partial implementation on names.
Its extent is exactly the pulled-back extent of $f$, and its value denotes that of $f$.
This is the construction of \cref{thm:partialadequacy}, applied to each operation and to every parameter family.
In particular it applies to evaluation and to the existing least-fixed-point map.
\item \textbf{Unit and identity laws.} Identity and composition preserve the interpretation equations.
The equality relation on implementations is a congruence because it is equality after the target quotient.
Fullness of the quotient functor is projective lifting; faithfulness is equality after the source cover followed by faithfulness of $\Pt$ in \cref{jd:thm:splitcomparison}.
Every characterized object has its constructed cover, proving essential surjectivity.
\item \textbf{Pullback construction.} Pullback of a cover supplies cartesian substitution. A projective cover of that pullback supplies implementation data when required.
Two refinements have a common projective refinement; on it their classifiers and interpreted values agree.
Descent proves coherence.
Change of basis is handled by \cref{prop:changebasis}.
For (4), apply \cref{thm:refinement-quotient} to total maps and \cref{thm:partial-refinement-quotient} to partial maps. The latter constructs the exact source image, descends its classifier, and proves compatibility with composition. \Cref{prop:refined-finite-returns} constructs their iteration and feedback witnesses at arbitrary covers.
\item \textbf{Independence of name representatives.} Thus the entire construction, including its equations, is independent of the selected representatives in the stated sense.
Only fixed-source lifts into arbitrary nonmodest target name covers retain the external choice qualification of \cref{rem:external-projectivity}; the modest-target and covered-span constructions just cited do not.
\end{enumerate}
\end{proof}

\subsection{Evaluation and specialization}
For a computational function object $\mathrm{Prog}_{X,Y}$, the interpretation equation is the square of partial maps
\[
\begin{tikzcd}[column sep=large,row sep=large]
 C_{\mathrm{Prog}}\times C_X\ar[r,dashed,"\widehat{\ev}"]
       \ar[d,two heads,"q_{\mathrm{Prog}}\times q_X"']
 &C_Y\ar[d,two heads,"q_Y"]\\
 \mathrm{Prog}_{X,Y}\times\Pt X\ar[r,dashed,"\ev"']&\Pt Y.
\end{tikzcd}
\]
Commutativity means equality in $\Par_\Sigma(\E)$, including the domains. Typed specialization is already a total semantic operation: curry evaluation with one argument held fixed. Apply \cref{thm:canonicalcoding} to obtain its compiler. The evaluation--specialization equation holds after interpretation because it is the exponential evaluation equation below the covers.

The same construction implements a name as constant data. Regard a name object $C$ as a discrete $\J$-algebra; every ideal in a discrete poset is principal, so this algebra is canonical. The constant-function operation $C\to[1,C]_\Sigma$ lifts on its projective source to a presentation of that function sort. Evaluation of the resulting program denotes the original name. This operation works on names that have already been supplied; it does not split $q_D$ on semantic values. The relative finite-description construction in \cref{sec:relativecodes} retains the supplied carrier label and its ambient realizer. Its quotation operation uses tree data constructors; it does not convert that parameter into a numeral.

\subsection{Implementing iteration and interaction}
The parameter sorts in \cref{thm:canonicalcoding} include the actual computational function objects. This gives an implementation theorem for the operations constructed in \cref{sec:intrinsicrecursion}, including their dependence on the equation being solved.

\begin{theorem}[Coding preserves iteration and feedback]\label{exec:thm:codefeedback}
Choose projective name covers for
\[
 E_{X,Y}=[X,\LL(Y+X)]_\Sigma,\qquad
 P_{X,Y}=[X,\LL Y]_\Sigma,
\]
and for the analogous feedback parameter sort
$F_{X,Y,U}=[X+U,\LL(Y+U)]_\Sigma$.
There are total implementation maps $\operatorname{iterCode}$ and $\operatorname{feedbackCode}$ making the squares
\[
\begin{tikzcd}[column sep=large,row sep=large]
 C_E\ar[r,"\operatorname{iterCode}"]\ar[d,two heads,"q_E"']
   &C_P\ar[d,two heads,"q_P"]\\
 E_{X,Y}\ar[r,"(-)^\ddagger"']&P_{X,Y}
\end{tikzcd}
\qquad
\begin{tikzcd}[column sep=large,row sep=large]
 C_F\ar[r,"\operatorname{feedbackCode}"]\ar[d,two heads,"q_F"']
   &C_P\ar[d,two heads,"q_P"]\\
 F_{X,Y,U}\ar[r,"\Tr^U_+"']&P_{X,Y}
\end{tikzcd}
\]
commute. Evaluation of either square is equality of partial maps, including their $\Sigma$-classified extents. The equalities commute with parameters and common refinement of covers. On the quotient of presented partial computations by denotational equality, the interpretation equivalence preserves coproduct trace and induces an equivalence of their $\operatorname{Int}$ interaction categories.
\end{theorem}
\begin{proof}
\leavevmode
\begin{enumerate}[label=\textup{(\roman*)},leftmargin=*,itemsep=3pt,topsep=4pt]
\item \textbf{Observation preservation.} Both bottom arrows are algebra morphisms by \cref{exec:thm:elgot,exec:thm:trace}, hence are arrows of $\E$.
Apply projective lifting through $q_P$ on $C_E$ and $C_F$ to obtain the upper arrows. When $C_P$ is modest, \cref{lem:modest-target-lift} constructs these same lifts without external choice. For arbitrary covers, \cref{prop:refined-finite-returns} gives the exact iteration and feedback implementations as covered spans, and proves their finite-witness comparison without requiring a fixed-target lift.
Their totality concerns construction of an implementation, not termination of the implemented loop.
Compose with the evaluation operation and use its adequacy square.
Equality of partial elements gives equality of extents and values, and \cref{exec:cor:observe} identifies the positive witnesses.
\item \textbf{Image and quotient construction.} Finite sums of projective presentations are presented by the coproduct of their covers.
Implemented composition descends to composition of partial denotations, and the two displayed squares supply iteration and trace on this quotient.
Any two choices of an upper arrow have the same denotation; thus the trace is independent of that choice.
Projective lifting supplies every partial semantic arrow, by the domain-and-value construction in \cref{thm:partialadequacy}, and the imposed congruence makes interpretation faithful.
\item \textbf{Inverse comparison.} This gives the equivalence on the presented computational sorts and their finite-sum closure.
Refinement of covers preserves its interpretation equations by descent.
Finally the $\operatorname{Int}$ composition is formed from coproduct, symmetry, and trace; applying the equivalence to these operations yields the stated interaction equivalence.
\end{enumerate}
\end{proof}

\paragraph{Finite witness transformations.}
The iteration wrapper maintains a state, runs the supplied equation,
returns on its $Y$ tag, and updates on its $X$ tag. Feedback first
tags the external input and uses the same loop on the interface.
Specialization constructs these wrappers without executing them.
Their accepted runs split into finite lists of state-update runs and
a final return; conversely concatenation and the marked wrapper states
construct such runs. Register renaming, splitting, and concatenation
are total computable on the supplied finite data, with checker equations
verified one transition at a time. Induction on the list identifies
the support with $R_n$ of \cref{exec:lem:return}. Thus the interpreted
trace and guarded finite-completion equations retain exact domains
and values, while distinct wrappers can retain different syntax and timing.

\subsection{Conservativity of extensional decisions}
\begin{proposition}\label{prop:decisiondescent}
For a regular epimorphism $q:C\onto P$, a predicate on $C$ descends to $P$ precisely when its classifier equalizes the kernel pair of $q$. This holds for both $2$-classifiers and $\Sigma$-classifiers. Pullback and descent are inverse on these predicates, uniformly in parameters.
\end{proposition}
\begin{proof}
\leavevmode
\begin{enumerate}[label=\textup{(\roman*)},leftmargin=*,itemsep=3pt,topsep=4pt]
\item \textbf{Kernel-pair descent.} Apply the coequalizer universal property of $q$ to the classifier.
Its codomain is already $2$ or $\Sigma$, so descent retains that classification.
Epicity proves the inverse identities.
Pullback preserves the kernel-pair diagram and its regular-epimorphic quotient in each slice, proving parameter compatibility.
\end{enumerate}
\end{proof}

Thus a decidable property of names that respects denotation gives a decision on the computation object. Applying \cref{thm:intrinsichalting} gives the index form of Rice for every presentation of $P$ in $\Eff$. A halting decider on names similarly descends on the program argument. Positive saturated predicates descend to $\Sigma^P$ and hence, by \cref{thm:observationfullness}, to the finite-information observations. The numerical section makes these translations explicit.

For a name transformation $t:C\to C$, the image of $(q,qt)$ is
a total relation on its denotation object. On the computational objects
of \cref{thm:syntheticmodel}, \cref{prop:relationaltransport} gives its
internal fixed points. The exposed program cover below makes the
numerical case uniform; \eqref{eq:lfpimplementation} separately
implements the constructed least solution.


The implementation squares identify the domains and outputs of the chosen primitive operations without requiring identical names for identical denotations. The free term monad can therefore compose those implementations while keeping substitution of parameter names explicit.

\section{Description syntax, substitution, and interpretation}\label{sec:relativecodes}

Individual trackers can implement a primitive, but composite programs require a syntax whose substitution respects interpretation. Finite descriptions provide that syntax while retaining arbitrary parameter names. The main task is to prove that syntactic substitution and partial semantic evaluation give the same extent and returned value.


Fix the implemented signature of \cref{thm:canonicalcoding}.
Its descriptions form a substitution monad on the supplied name
carriers, interpreted exactly in $\LL$. Numerical trackers implement
its structural recursions; the tree carriers retain opaque parameter
labels. We first construct the delayed numerical operations used by
those recursions.

\subsection{Delayed abstraction and uniform compilers}
The fixed $K_1$ combinators satisfy $\mathsf s abc\simeq ac(bc)$, where $\simeq$ includes equality of definedness. A compiler constructs a program index before any application represented by its body is executed.

\begin{lemma}[Delayed abstraction]\label{rel:lem:abstraction}
For an applicative expression $t(x,\vec a)$, use bracket abstraction with
the clauses
\[
 [x]x=\mathsf i,\qquad [x]a=\mathsf k a\quad(a\text{ atomic},\ a\ne x),\qquad
 [x](uv)=\mathsf s([x]u)([x]v).
\]
Then $[x]t$ is defined for every assignment of ambient elements to the
remaining variables. In $K_1$,
$([x]t)b\simeq t[b/x]$. Abstracting over the remaining parameters gives a total computable compiler of indices in $K_1$.
\end{lemma}
\begin{proof}
\leavevmode
\begin{enumerate}[label=\textup{(\roman*)},leftmargin=*,itemsep=3pt,topsep=4pt]
\item \textbf{Delayed abstraction.} Induct on the expression.
The atomic clauses are defined by the pca laws.
At an application node both abstracted subexpressions are defined by
induction, and $\mathsf s uv$ is defined for every pair of elements $u,v$.
The application equation is the corresponding $\mathsf s$ equation and the
inductive equations.
The finite bracket algorithm and the total index constructors give the compiler assertion.
\item \textbf{Verification of the equations.} It is essential that the constant clause is restricted to
atomic expressions: replacing an unevaluated composite $uv$ by $\mathsf k(uv)$
would require $uv$ to be defined before the argument arrives.
\end{enumerate}
\end{proof}

\begin{proposition}[Uniform numerical compilers]\label{rel:prop:compilers}
There are total computable compilers for constant data, composition and
specialization, with the equations
\begin{align*}
 \mathsf{const}(a)&=\mathsf k a,& (\mathsf{const}(a))x&=a,\\
 \mathsf{comp}(b,a)x&\simeq b(ax),\\
 \mathsf{spec}(t,a)x&\simeq t\relpair{a}{x}.
\end{align*}
They return numerical program indices before executing their bodies.
\end{proposition}
\begin{proof}
\leavevmode
\begin{enumerate}[label=\textup{(\roman*)},leftmargin=*,itemsep=3pt,topsep=4pt]
\item \textbf{Uniform numerical compilers.} Use $\mathsf k$, $[b][a][x]b(ax)$ and $[t][a][x]t\relpair{a}{x}$, respectively,
with the abstraction of Lemma~\ref{rel:lem:abstraction}.
The abstractions over
the final input delay the applications.
Pairing is represented in $K_1$.
Substitution proves the equations, and the delayed bracket algorithm gives totality of the index constructors.
\end{enumerate}
\end{proof}

For presentations $q_X:P_X\onto X$ and $q_Y:P_Y\onto Y$, an implemented
map retains both $\widehat f:P_X\to P_Y$ and its tracker $t\in K_1$:
\[
 q_Y\widehat f=fq_X,
 \qquad t\varepsilon_X(p)=\varepsilon_Y(\widehat f(p)).
\]
These equations hold in the fixed $K_1$ interpretation. They retain the actual carrier labels as well as their numerical names. Composition of implementations uses
Proposition~\ref{rel:prop:compilers}; equality of their denotations is witnessed
through the target kernel pair.


\subsection{The free term monad}
\label{int:sec:syntax}
Fix the implemented finitary many-sorted
signature, with sort family $S$, supplied literal carriers, and supplied
arity and sorting data for its constructor names. In comparing two
implementations we keep these literal and operation symbols fixed;
their interpretations are the displayed maps of the respective models.
Higher-order operations, including specialization and iteration, occur
as typed primitive symbols with their already constructed implementations.

\begin{construction}[Finite descriptions with variable carriers]
\label{int:con:syntax}
For a family $V=(V_s)_{s\in S}$, let $\Syn(V)$ be the family of finite
well-sorted trees with variable leaves in $V$, the fixed supplied literal
leaves, and the specified constructor nodes. A map $u:V\to W$ relabels
variable leaves and fixes literals and operation symbols. There are maps
\[
 \eta^{\Syn}_V:V\longrightarrow\Syn(V),\qquad
 \mu^{\Syn}_V:\Syn(\Syn(V))\longrightarrow\Syn(V),
\]
given by variable insertion and substitution of the inner tree at an
outer variable leaf. For partitioned variable and literal carriers, a
tree's name is the tagged finite tree of their supplied names.
Carrier labels remain in the tree itself.
\end{construction}

\begin{proposition}[Constructed substitution monad]
\label{int:prop:syntax}\label{rel:prop:trees}
The construction gives a monad on the ordinary category of partitioned
assembly families and tracked carrier maps, with equality of carrier
functions. It also gives a monad on families in $\Eff$, agreeing on the
presented trees. Relabelling, insertion, and substitution have uniform
trackers on their supplied partitioned presentations. The raw tree
carrier is also partitioned. Constructors, finite variable shifts, and
quotation of a supplied tree into its tree-data extension are tracked;
the quoted expression evaluates to that tree.
\end{proposition}
\begin{proof}
\leavevmode
\begin{enumerate}[label=\textup{(\roman*)},leftmargin=*,itemsep=3pt,topsep=4pt]
\item \textbf{Finite witness construction.} Finite labelled trees are encoded by finite node lists, their incidence
data, and their node labels. The finite equalities specifying the root,
the ordered children, and their sorts form the well-sorted subobject.
The construction uses the NNO and finite lists. Give a literal
$p\in C_i$ the name $\langle\mathsf{lit}_i,\varepsilon_{C_i}(p)\rangle$,
a variable its tag and supplied name, and a compound node its
constructor name paired with the finite list of child names. The
carrier retains all labels, even when their realizers coincide.
Sorting and rule evidence use the supplied proof-name presentations;
no decidable equality on arbitrary labels is required. Relabelling
and substitution copy literals and recurse on children, supplying
their trackers. Finite shifts act computably on variable indices.
\item \textbf{Parameter substitution.} For quotation use the data extension with the literal carriers, finite
lists, trees, and their constructors. A variable becomes variable data,
a literal copies its supplied parameter, and a compound node applies
the tree-data constructor to the quotations of its children, obtained by structural recursion.
Structural recursion tracks this operation, and induction proves that
its evaluation is the original tree. It makes no enumeration or
selection of arbitrary semantic values.
\item \textbf{Unit and identity laws.} For the monad equations, proceed by the three tree clauses.
At a variable, either unit equation inserts and then removes the
variable node, while associativity substitutes the same specified inner
tree. At a literal all operations copy that literal. At a constructor,
the equations follow from the child equations and the unchanged
constructor. This proves naturality, both unit equations, and
associativity as equalities of carrier maps. Retaining a chosen tracker
is additional P-categorical data and does not alter this equality.
The internal construction has the same clauses and equations.
Encoding and decoding are inverse on each constructor, hence inverse
on trees by induction.
\end{enumerate}
\end{proof}

For finite binding contexts, weakening a substitution into a context
extended by $\Delta$ fixes the fresh variables and weakens the old
replacements. The old-variable and fresh-variable clauses give
$(\tau^\#\sigma)^{+\Delta}=(\tau^{+\Delta})^\#\sigma^{+\Delta}$;
induction through binding nodes gives the context-indexed substitution
equations. This is the algebraic organization of binding developed in
\cite{fioreplotkinturi}. The interpretation below concerns the supplied
typed primitive signature, so its constructor step uses the exact
primitive operations already constructed.

\begin{construction}[Interpretation on partial values]
\label{int:con:interpretation}
Let $X=(X_s)_{s\in S}$ be the semantic sorts. Write $L=\LL$ in this
construction. For each primitive $o:\prod_{k=1}^nX_{s_k}\to LX_t$, put
\begin{equation}\label{int:eq:liftedoperation}
 o^\sharp:
 \prod_{k=1}^nLX_{s_k}
 \xrightarrow{\operatorname{seq}}L\Bigl(\prod_{k=1}^nX_{s_k}\Bigr)
 \xrightarrow{L(o)}L^2X_t\xrightarrow{\mu^L}LX_t.
\end{equation}
Here $\operatorname{seq}$ is the finite product operation supplied by
the commutative strength of $L$; its extent is the conjunction of the
argument extents. The nullary case uses the unit at the terminal object.
Interpret each literal by its supplied partial value, using $\eta^L$
for a total value. Tree recursion defines
\[
 \alpha_X:\Syn(LX)\longrightarrow LX:
\]
a variable returns its partial value, a literal its supplied
interpretation, and a constructor applies $o^\sharp$ to the results
of its children.
\end{construction}

\begin{proposition}[Interpretation and substitution]
\label{int:prop:substitution}
The map $\alpha_X$ is a $\Syn$-algebra:
\begin{equation}\label{int:eq:treealgebra}
 \alpha_X\eta^{\Syn}_{LX}=\id_{LX},\qquad
 \alpha_X\mu^{\Syn}_{LX}=\alpha_X\Syn(\alpha_X).
\end{equation}
For $v:W\to LX$, put $\operatorname{eval}_v=\alpha_X\Syn(v)$.
If $\sigma:V\to\Syn(W)$ and
$\sigma^\#=\mu^{\Syn}_W\Syn(\sigma)$, then
\begin{equation}\label{int:eq:substitution}
 \operatorname{eval}_v\sigma^\#
   =\operatorname{eval}_{\operatorname{eval}_v\sigma}.
\end{equation}
These equations are equalities of partial results, including their extents.
\end{proposition}
\begin{proof}
\leavevmode
\begin{enumerate}[label=\textup{(\roman*)},leftmargin=*,itemsep=3pt,topsep=4pt]
\item \textbf{Interpretation and substitution.} The unit equation is the variable clause.
For the multiplication
equation, a variable of the outer tree contains an inner tree; both
routes evaluate it. A literal is unchanged.
At a constructor the
induction equations on its children give identical arguments to the
same $o^\sharp$.
This proves \eqref{int:eq:treealgebra}.
\item \textbf{Naturality.} Naturality of $\mu^{\Syn}$ and the algebra equation give
\begin{align*}
 \alpha_X\Syn(v)\mu^{\Syn}_W\Syn(\sigma)
 &=\alpha_X\mu^{\Syn}_{LX}\Syn(\Syn(v))\Syn(\sigma)\\
 &=\alpha_X\Syn(\alpha_X\Syn(v)\sigma),
\end{align*}
which is \eqref{int:eq:substitution}.
Equality in $LX$ identifies the
extent classifiers and the values on the common extent.
\end{enumerate}
\end{proof}

The environment in \eqref{int:eq:substitution} consists of partial
values; it does not execute every replacement in advance. Substituting
a divergent expression for an unused variable in the literal $0$ still
returns $0$. The strict expression $\mathsf{let}\ x=u\ \mathsf{in}\ 0$
instead has the domain of $u$. Its sequential law is Kleisli associativity.
Literal quotation is the tree-data construction of \cref{rel:prop:trees},
with its proved evaluation equation. It is distinct from $\delta$ of
$\RealQ$. The modal interpretation of staged code originates in
\cite{daviespfenning}; the distinction between retained realizers and
literal program text, including failure of strict naturality for
syntactic quotation, is developed in \cite[Remark~4.5]{chenko}.

\begin{theorem}[Exact interpretation of descriptions]
\label{int:thm:exactsyntax}
For the supplied presentations $q_s:C_s\onto X_s$, the implementation
and semantic tree algebras satisfy the commuting square
\begin{equation}\label{int:eq:exactalgebra}
\begin{tikzcd}[column sep=large]
 \Syn(\LL C)\ar[r,"\alpha_C"]\ar[d,"\Syn(\LL q)"']&
 \LL C\ar[d,"\LL q"]\\
 \Syn(\LL X)\ar[r,"\alpha_X"']&\LL X.
\end{tikzcd}
\end{equation}
Consequently every finite typed expression has equal implementation
and semantic domains and equal interpreted values there. The equality
respects substitution and change of parameter.
\end{theorem}
\begin{proof}
\leavevmode
\begin{enumerate}[label=\textup{(\roman*)},leftmargin=*,itemsep=3pt,topsep=4pt]
\item \textbf{Exact interpretation of descriptions.} For a primitive, exact adequacy of \cref{thm:partialadequacy} is
\[
 \LL(q_t)\widehat o=o\Bigl(\prod_kq_{s_k}\Bigr).
\]
Naturality of strength and multiplication extends it to
\[
 \LL(q_t)\widehat o^\sharp
   =o^\sharp\Bigl(\prod_k\LL(q_{s_k})\Bigr).
\]
\item \textbf{Induction.} Induct on the tree. The variable case is $\LL q$; the literal case
is its supplied adequacy equation; the displayed primitive equation
proves the constructor case. This gives \eqref{int:eq:exactalgebra}.
For $w:V\to\LL C$ it yields
$\LL q\operatorname{eval}^{C}_w
=\operatorname{eval}^{X}_{\LL q\,w}$.
The classifier of a partial element is unchanged by $\LL q$, so the
domain equality is exact. On that domain the equation identifies
the interpreted values.
\item \textbf{Pullback construction.} Apply \cref{int:prop:substitution} to obtain
the same equality after substitution. In a context $\Gamma$ perform
the construction in the slice, using its relative partial-element
monad. Pullback preserves the finite-tree clauses and the primitive
domain--value squares, giving the parameter statement.
\end{enumerate}
\end{proof}

\begin{theorem}[Common-refinement invariance of descriptions]
\label{int:thm:refinement}
Two projective presentations $q_s:C_s\onto X_s$ and
$q'_s:C'_s\onto X_s$, with exact implementations of the same signature
and the same interpreted literal symbols, admit partitioned projective
common refinements
\[
 P_s\onto C_s\times_{X_s}C'_s,
 \qquad p_s:P_s\to C_s,\quad p'_s:P_s\to C'_s,
 \qquad \bar q_s=q_sp_s=q'_sp'_s.
\]
For every partial-value environment $v:V\to\LL P$ one has
\begin{equation}\label{int:eq:refinement}
 \LL q\operatorname{eval}^{C}_{\LL p\,v}
 =\operatorname{eval}^{X}_{\LL\bar q\,v}
 =\LL q'\operatorname{eval}^{C'}_{\LL p'\,v}.
\end{equation}
The equalities include domains, respect expression substitution, and
are coherent after interpretation under further refinement and parameter
change. Paired primitive implementations lift to the common refinements,
using first-hit search for their partial domains.
\end{theorem}
\begin{proof}
\leavevmode
\begin{enumerate}[label=\textup{(\roman*)},leftmargin=*,itemsep=3pt,topsep=4pt]
\item \textbf{Common-refinement invariance of descriptions.} Pull back the two covers and take a partitioned projective cover.
Their projections cover $X_s$ and give the displayed family. A total
primitive on a product $P$ of common input presentations gives a pair
of output names with the same denotation. It therefore factors through
$C_t\times_{X_t}C'_t$ and lifts to $P_t$ by projectivity of $P$.
\item \textbf{Observation preservation.} For a partial primitive with semantic domain $U\incl Z$, let
$a:P\to Z$ interpret the common inputs. Exact adequacy identifies both
implementation domains with $U_P=P\times_ZU$. Their output pair is a
map $U_P\to C_t\times_{X_t}C'_t$. Lift the classifier $P\to\Sigma$
through the binary-stream cover, obtaining a total stream $\theta_p$.
The partitioned hit presentation
$H_\theta=\{(p,n):\theta_p(n)=1\}$ covers $U_P$. Projectivity of
$H_\theta$ lifts the output pair to $P_t$. For $p\in U_P$, search for
the least hit of $\theta_p$ and apply this lift. Its domain is exactly
$U_P$, and its two projections are the original output maps. The
search uses the supplied stream and makes no global selection of
names for arbitrary semantic values.
\item \textbf{Partiality and definedness.} Apply \cref{int:thm:exactsyntax} to each presentation and use
$\LL q\,\LL p=\LL\bar q=\LL q'\,\LL p'$.
This proves \eqref{int:eq:refinement}; taking definedness gives the
domain equation, and restriction to the common domain gives the value
equation. The substitution equation applies to each tree algebra.
\item \textbf{Pullback construction.} Further refinements give the same middle semantic evaluation.
Two choices admit a projective cover of their fibre product, where
these interpretations agree, and the equality descends by epicity.
After parameter pullback, refine the pulled-back covers projectively
and repeat (i)--(iii). Uniqueness of the interpretation and the
canonical pullback isomorphisms give identity and composition coherence.
\end{enumerate}
\end{proof}


\subsection{Delayed numerical recursion}
The delayed index constructors give the following uniform self-reference equation in $K_1$.

\begin{theorem}[Delayed numerical recursion]\label{rel:thm:recursion}
There is a total computable compiler $\mathsf{rec}:\N\to\N$ such that, for $t\in K_1$, $r=\mathsf{rec}(t)$ is defined and
\[
                    r x\simeq (t r)x.
\]
With a supplied parameter $z$, there is a uniform total computable compiler producing $r_z$ with
\[
 r_zx\simeq (h\relpair{z}{r_z})x.
\]
\end{theorem}
\begin{proof}
\leavevmode
\begin{enumerate}[label=\textup{(\roman*)},leftmargin=*,itemsep=3pt,topsep=4pt]
\item \textbf{Tracked implementation.} Using delayed abstraction, put
\[
 d=[z][x](zz)x,
 \qquad g_t=[z]t(dz),
 \qquad r=d g_t.
\]
The element $d$ is a fixed $K_1$ index.
For every $z$, $dz$ is defined because
it is the abstracted program waiting for $x$, even if $zz$ is undefined.
The element $g_t$ is defined for every $t$, and hence so is $d g_t$.
Calculation gives
\[
 (d g_t)x\simeq (g_tg_t)x
             \simeq (t(dg_t))x=(tr)x.
\]
Abstracting the expression $d([z]t(dz))$ over $t$ gives the required
compiler.
\item \textbf{Totality of the numerical compiler.} The finite bracket construction and delayed index application make this compiler total computable.
For the parameterized form apply the same construction to
$t_z=[r]h\relpair{z}{r}$; the parameter remains ambient data.
\end{enumerate}
\end{proof}

\begin{remark}[What the equation proves]
The theorem establishes equality of one-argument behaviors. It does not
establish the raw-name equation $r=tr$. More importantly, it does not
establish that $r$ belongs to any previously specified typed name
presentation $P_D$. A proof that a tracker operates correctly on valid
names cannot be applied to $r$ until its validity has been established.
Using that application to establish validity would be circular.
\end{remark}

\begin{proposition}[Application and represented partial domains]\label{rel:sec:weak}
Raw application $K_1\times K_1\pmap K_1$ has a $\Sigma$-classified domain in $\Eff$: bounded universal simulation gives its finite-hit stream. For a presented semantic operation, its implemented domain is the pullback of that operation's domain, and its value must satisfy the target interpretation equation.
\end{proposition}
\begin{proof}
\leavevmode
\begin{enumerate}[label=\textup{(\roman*)},leftmargin=*,itemsep=3pt,topsep=4pt]
\item \textbf{Application and represented partial domains.} For indices $e,n$, let the $s$th bit record whether universal execution of $e$ on $n$ halts within $s$ steps.
Bounded simulation is primitive recursive and its hit support is exactly defined application.
\item \textbf{Partiality and definedness.} The second assertion follows by pulling back the domain mono and lifting the value map on the hit assembly, as in \cref{thm:partialadequacy}. A raw fixed-point index therefore supplies the displayed application equation, while validity as a name in a particular presentation is established by its interpretation and domain equations.
\end{enumerate}
\end{proof}

\section{Numerical codes and specialization}\label{sec:numericalcodes}


Use the register enumeration and finite execution checker of \cref{subsec:numerical-foundation,thm:operational-certificates}. We now identify its program-cover semantics and specialization maps.
Bounded simulation of all pairs $(n,s)$ gives, uniformly in $e$, an increasing sequence of finite graphs $g_{e,s}$ whose union is $\graph(\varphi_e)$. The sequence can be chosen total and primitive recursive in $(e,s)$ by restricting inputs and running times to $s$. It yields a map
\begin{equation}\label{eq:qcodes}
 q:\N\longrightarrow\PC,
 \qquad q(e)=\{g\in\FG\mid g\subseteq\graph(\varphi_e)\}.
\end{equation}
Finite graph containment is a finite conjunction of convergence tests, so $q(e)$ is an effective point, uniformly in $e$.

\begin{proposition}\label{prop:codecover}
The map $q$ is a regular epimorphism. The object $\PC$ is the modest object of partial computable functions with program indices as names, and the kernel pair of $q$ is
\[
 e\sim d\iff\varphi_e=\varphi_d.
\]
\end{proposition}
\begin{proof}
\leavevmode
\begin{enumerate}[label=\textup{(\roman*)},leftmargin=*,itemsep=3pt,topsep=4pt]
\item \textbf{Observation preservation.} Cover $\PC$ by a partitioned assembly $P$.
Its point-membership predicate on $P\times\FG$ is $\Sigma$-classified.
Lift it through the stream cover on this projective source, obtaining a total computable stream evaluator on names.
Given a parameter name $a$ and input $n$, search the enumerated finite graphs and their hit stages until a graph containing a pair $(n,m)$ is found, then return $m$.
The ideal properties ensure functionality; the resulting function has exactly the graph represented by the parameter point.
\item \textbf{Existence and uniqueness.} This search is compiled uniformly from $a$, giving a tracked map $P\to\N$ whose composite with $q$ is the cover $P\onto\PC$.
Hence $q$ is epic and therefore regular epic in the topos.
\item \textbf{Ideal and directedness calculation.} Two codes have equal denotations precisely when their finite subgraph ideals coincide, equivalently when their partial functions coincide.
This equality has a uniform proof-name interpretation: given two extensionally equal codes and a hit witnessing finite graph containment for either one, search the other program for the same finitely many outputs.
The search terminates and transforms the witness.
The two transformations realize equality of the membership predicates uniformly in the pair of codes.
\item Hence the map from codes descends through the PER of extensional equality.
The two translations just constructed, between point names and program indices, induce inverse maps between $\PC$ and this PER quotient.
It is modest by \cref{prop:per}.
\end{enumerate}
\end{proof}


\begin{proposition}[The program cover is an exposure counit]
\label{int:prop:programcover}
For the numerical assembly presentation of $\PC$ in
\cref{prop:codecover}, there is a tracked isomorphism
\[
 \nu:\N\xrightarrow{\cong}\RealQ\PC,
 \qquad \nu(e)=(q(e),e),\qquad \varepsilon^{\RealQ}_{\PC}\nu=q.
\]
The counit has no section in $\Eff$. Two tracked implementations of
$\id_{\PC}$ have distinct exposures.
\end{proposition}
\begin{proof}
\leavevmode
\begin{enumerate}[label=\textup{(\roman*)},leftmargin=*,itemsep=3pt,topsep=4pt]
\item \textbf{The program cover is an exposure counit.} The realizer $e$ of a partial computable function determines precisely
the denotation $q(e)$. Projection to the retained index is therefore
inverse to $\nu$, and both maps have identity tracker.
\item \textbf{Tracked implementation.} For intensional separation, construct two distinct finite wrappers
containing the literal $e$, each executing the universal interpreter
on $e$ and its input, with an extra idle instruction in the second.
Their register-syntax codes are distinct. Choose the first unless its
code equals $e$, and otherwise choose the second. This is a total
computable padding map $b$ with $b(e)\ne e$ and $q(b(e))=q(e)$.
The identity and padding programs track $\id_{\PC}$, while their
exposures send the retained index $e$ to $e$ and $b(e)$.
\item \textbf{Universal-property argument.} If $\varepsilon^{\RealQ}_{\PC}$ had a section, its composite with
$\nu^{-1}$ would give $j:\PC\to\N$ with $qj=\id$.
The mono $j$ would make equality on $\PC$ decidable by numerical
comparison. Comparison with bottom would then give a nonconstant
decision $\PC\to2$, since bottom differs from the constant-zero
function. This contradicts the intrinsic Rice theorem
\cref{thm:intrinsichalting}. Thus no section exists.
\end{enumerate}
\end{proof}


\subsection{The name translations in the program cover}
It is useful to describe the two translations in \cref{prop:codecover} at the level of witnesses. Fix a computable enumeration $\nu:\N\onto\FG$. If a point name $a$ supplies a stream $\theta(a,g,k)$ for finite graph membership, the program produced from $a$ operates on input $n$ as follows. At stage $s$ it checks all pairs $(l,k)$ with $l,k\leq s$. When
\[
 \theta(a,\nu(l),k)=1
 \quad\text{and}\quad(n,m)\in\nu(l),
\]
it returns $m$. Every stage makes only finitely many calls to the total stream evaluator, so this is a partial computable procedure. If the represented partial function is defined at $n$, some member graph contains its pair and some membership stream has a hit, so a sufficiently large stage finds it. Conversely every returned pair lies in a member graph and therefore in the represented function. Directedness guarantees that different discoveries cannot return different values.

In the other direction a program index $e$ supplies the stream
\[
 \theta_e(g,s)=1\quad\Longleftrightarrow\quad
 \text{every pair $(n,m)\in g$ is verified by an $s$-step run of $e$ on $n$}.
\]
The empty graph is accepted immediately. This stream is total and computable jointly in $(e,g,s)$. Its hit support is the ideal of finite subgraphs of $\varphi_e$. These two procedures are uniform syntax constructions, so they operate on projective families of point names, not just on individual global elements.

For completeness, equality of the ideals produced in the first composite also has explicit witness translations. Given a hit of $\theta(a,g,-)$, run the translated program on the finitely many inputs in $g$ until their outputs are seen. It accepts the finite graph $g$. Conversely, a finite computation witness for every pair in $g$ records the member graphs and stream hits used by the search. Combine these finitely many ideal members by directedness, take their finite subgraph $g$, and obtain the corresponding membership of the original ideal. In a selected presentation, \cref{lem:extraction} turns the combination and lower-closure proofs into transformations of the selected proof names. Equivalently, since the membership stream is total and its support contains $g$, a direct search for its hit gives the same equality of classifiers. This proves the inverse equation uniformly without selecting a unique program for a function.

The numerical cover therefore has two separate features. Every denotation has a program name locally, expressed by regular epicity. Equality of two names is equality of all their finite graph tests, expressed by its kernel pair. Neither feature gives a decision procedure for that equality. The compiler operations below respect this kernel pair even though they act on the unreduced numerical syntax.


\begin{proposition}[Compilers]\label{prop:compilers}
There are total computable operations $s(e,a)$ and $c(e,d)$ satisfying
\begin{align}
 \varphi_{s(e,a)}(x)&\simeq\varphi_e(\langle a,x\rangle),\label{eq:smn}\\
 \varphi_{c(e,d)}(x)&\simeq\varphi_e(\varphi_d(x)).\label{eq:compilercomposition}
\end{align}
There is also a total computable $j:\FG\to\N$ coding each finite partial graph exactly.
\end{proposition}
\begin{proof}
\leavevmode
\begin{enumerate}[label=\textup{(\roman*)},leftmargin=*,itemsep=3pt,topsep=4pt]
\item \textbf{Compilers.} The operations $\mathsf{spec}$ and $\mathsf{comp}$ of \cref{rel:prop:compilers} already construct the required behaviors.
Choose their register representatives by the wrappers that, respectively, insert the literals $(e,a)$ and pair the input, or run $d$ and then $e$.
Writing these finite wrappers is total computable; the second interpreter is entered only after the first returns.
\item \textbf{Finite witness construction.} The defining equations therefore give
\[
 q(s(e,a))=q(\mathsf{spec}(e,a)),\qquad
 q(c(e,d))=q(\mathsf{comp}(e,d)).
\]
This is a comparison of implementations through the program cover, without identifying their raw codes.
For $j(g)$, output the finite list of input-equality tests returning the indicated value, with a loop if no test succeeds.
Finite syntax recursion gives totality, and iterated specialization handles any finite tuple of parameters.
\end{enumerate}
\end{proof}

Universal numerical evaluation $U(e,n)\simeq\varphi_e(n)$ and the semantic evaluator satisfy
\[
 U=U_{\mathrm{sem}}(q\times\id_\N)
 \quad\text{in }\Par_\Sigma(\Eff).
\]
This is \eqref{eq:adequacysquare} for the concrete program presentation. The compiler $s$ implements semantic parameterization, and $c$ implements partial-map composition. In particular all indices are valid names, including indices of everywhere divergent programs.


\subsection{Composition as a map of presented domains}
Composition of partial numerical functions is an effective map
\[
 \operatorname{comp}:\FunDom\times\FunDom\longrightarrow\FunDom.
\]
One can see this directly from its finite information. For finite graphs $a,b,c$, the matrix
\begin{equation}\label{eq:compositionmatrix}
 M((a,b),c)\iff c\subseteq a\circ b
\end{equation}
is decidable. Each row is a principal ideal of finite graphs, and the matrix is monotone in both input graphs. The product presentation of \cref{prop:productpresentation} and \cref{thm:relations} therefore give an effective Scott-continuous map. Its value on ideals $I,J$ is
\[
 c\in\operatorname{comp}(I,J)
 \iff\exists a\in I\;\exists b\in J\;c\subseteq a\circ b.
\]
For a finite graph $c$ contained in the composite of two arbitrary partial functions, choose the finitely many intermediate values used by its pairs. They give finite subgraphs $a,b$ as in this formula. The reverse implication follows by composition of graph inclusions. Thus this effective extension is exactly ordinary partial-function composition, including for arbitrary internal ideals.

On effective points, the compiler of \cref{prop:compilers} supplies the commuting square
\[
\begin{tikzcd}
 \N\times\N\ar[r,"{c}"]\ar[d,"{q\times q}"']&\N\ar[d,"{q}"]\\
 \PC\times\PC\ar[r,"{\Pt(\operatorname{comp})}"']&\PC.
\end{tikzcd}
\]
The upper arrow compiles a sequential interpreter; the lower arrow is characterized by finite graph composition. Adequacy identifies their domains and outputs. Associativity of composition gives equality of the two three-program compilations after $q$, while their numerical codes may differ. This is a concrete instance of composition in the presentation fibration and of why its kernel-pair equality is the relevant equation for programs.


Specialization and composition are now both semantic operations and implemented numerical compilers. Their agreement is expressed by the commutative interpretation squares, allowing the halting and recursion arguments to use real program indices without redefining the extensional function domain.

\section{Halting and recursion}\label{sec:recursion}

Once evaluation has numerical codes, diagonal halting and self-reference can be expressed as statements about the chosen numbering. The same interpreter that makes halting semidecidable prevents a total halting decision. Delayed specialization then constructs the recursion compiler, without identifying syntactic fixed points with least semantic fixed points.


The universal convergence predicate
\[
 H(e,n)\iff\exists s,m\;T(e,n,s,m)
\]
is $\Sigma$-classified. Its diagonal $K(e)=H(e,e)$ is therefore semidecidable.

\begin{theorem}[Halting]\label{thm:halting}
There is no decision morphism $d:\N\to2$ in $\Eff$ classifying $K$. Externally, the diagonal halting set is c.e. and not computable.
\end{theorem}
\begin{proof}
\leavevmode
\begin{enumerate}[label=\textup{(\roman*)},leftmargin=*,itemsep=3pt,topsep=4pt]
\item \textbf{Halting.} A decision morphism is tracked by a total computable Boolean function.
Compile the program that on input $n$ returns $0$ when $d(n)=0$ and loops when $d(n)=1$.
If $a$ is its code, adequacy gives
\[
 H(a,a)\iff d(a)=0,
\]
whereas the decision specification gives $H(a,a)\iff d(a)=1$.
The Boolean alternatives contradict each other.
The finite computation witness proves computable enumerability.
\end{enumerate}
\end{proof}


\begin{corollary}\label{cor:nonsemidecidableorder}
The complement of diagonal halting is not $\Sigma$-classified in $\Eff$. The order relation on the completed function domain $\FunDom$ is not $\Sigma$-classified either.
\end{corollary}
\begin{proof}
\leavevmode
\begin{enumerate}[label=\textup{(\roman*)},leftmargin=*,itemsep=3pt,topsep=4pt]
\item \textbf{Construction of the comparison.} If both diagonal halting and its complement had computable membership streams, dovetail the two streams.
Exactly one eventually hits; return the corresponding Boolean.
This would be a total computable decision procedure, contradicting \cref{thm:halting}.
\item \textbf{Ideal and directedness calculation.} Uniformly in $e$, compile the partial function $f_e$ which is defined with value $0$ at input $0$ exactly when the computation $\varphi_e(e)$ halts, and is undefined at every other input.
Its finite graph ideal is an effective point of $\FunDom$, uniformly in $e$.
If the completed order were $\Sigma$-classified, pulling it back along $(f_e,\bot)$ would classify
\[
 f_e\leq\bot\quad\Longleftrightarrow\quad\neg H(e,e),
\]
contradicting the first assertion.
Notice that the finite graph basis still has decidable inclusion.
\item \textbf{Universal-property argument.} The obstruction concerns the completed order, not the approximation matrix from which it was built.
\end{enumerate}
\end{proof}

This counterexample gives a computational reason for the placement of $\Sigma$ in Part~\ref{part:completion}. Requiring the completed hom-predicate to factor through $\Sigma$ would exclude the principal program domain of this paper. Factoring its effective observation evaluation, and the approximation relations on searchable presentations, retains that domain and gives the operations needed for computation.


\begin{theorem}[Intensional recursion and numerical specialization]
\label{thm:recursion}\label{int:thm:ifp}
Every tracked $h:\RealQ\PC\to\PC$ has a point
$y:\mathbf1_*\to\PC$ satisfying, up to $\sim$,
\begin{equation}\label{int:eq:ifp}
\begin{tikzcd}[column sep=large]
 \mathbf1_*\ar[r,"m_0"]\ar[dr,"y"']&
 \RealQ\mathbf1_*\ar[r,"\RealQ y"]&\RealQ\PC\ar[dl,"h"]\\
 &\PC.
\end{tikzcd}
\end{equation}
The construction is uniform in the retained tracker. In particular,
every total computable $t:\N\to\N$ has an effectively obtained $e$
with $q(e)=q(t(e))$, equivalently $\varphi_e=\varphi_{t(e)}$.
For any partial computable $h(e,z,x)$ there is a total computable $k(z)$
such that
\begin{equation}\label{eq:paramrecursion}
 \varphi_{k(z)}(x)\simeq h(k(z),z,x).
\end{equation}
These constructions are uniform on supplied partitioned parameter names.
\end{theorem}
\begin{proof}
\leavevmode
\begin{enumerate}[label=\textup{(\roman*)},leftmargin=*,itemsep=3pt,topsep=4pt]
\item \textbf{Tracked implementation.} If $u$ tracks $h$, then $ue\downarrow$ and
$q(ue)=h(\nu(e))$ for every $e$, by \cref{int:prop:programcover}.
Put $r=\mathsf{rec}(u)$ using \cref{rel:thm:recursion}. Then
\[
 \varphi_r(x)\simeq\varphi_{ur}(x),\qquad
 q(r)=q(ur)=h(q(r),r).
\]
Let $y(*)=q(r)$ with constant-$r$ tracker. The exposed point
$\RealQ y\,m_0$ has value $(q(r),r)$, proving \eqref{int:eq:ifp}.
Every index is a valid program name, so this does not use a circular
typed-validity argument. Taking $h(q(e),e)=q(t(e))$ proves the
total-transformer assertion.
\item \textbf{Partiality and definedness.} For partial computable $h(e,z,x)$, specialization gives a total
index builder $H(e,z)$ for $x\mapsto h(e,z,x)$. Delayed compilation
produces a program $u_z$ for $e\mapsto H(e,z)$, totally in $z$.
Then $k(z)=\mathsf{rec}(u_z)$ is total and satisfies
$\varphi_{k(z)}(x)\simeq\varphi_{H(k(z),z)}(x)\simeq h(k(z),z,x)$.
Specializing at the supplied name of a partitioned parameter gives
the same construction in that context.
\end{enumerate}
\end{proof}

This is the exposure formulation of the second recursion theorem
\cite[Sections~4 and~9]{kavvosarticle}. Its equation is after program
interpretation, and does not identify $e$ with $t(e)$. The uniform
choice depends on the retained tracker and parameter names. Leastness
belongs to the semantic construction of \cref{thm:lfp} and its
implementation square \eqref{eq:lfpimplementation}; guarded solutions
belong to the temporal indexed theory of \cref{int:thm:temporal}.


The halting predicate, diagonal obstruction, and self-reference compiler have different logical roles: positive termination, failure of decision, and equality after code interpretation. Keeping them distinct permits the finite-information index theorems to descend from the semantic ones without treating an intensional fixed point as an order-theoretic least solution.

\section{Numerical implementation of finite information}\label{sec:finiteinformation}


The intrinsic finite-information theorem is already established. Descent along the program cover turns it into the classical index theorem. We also calculate explicit transformations of positive witnesses, using the constructed compilers.

An extensional c.e. set of program indices is a c.e. set $W\subseteq\N$ satisfying
\[
 \varphi_e=\varphi_d\Longrightarrow(e\in W\Longleftrightarrow d\in W).
\]
Let $\mathcal A$ be its induced class of partial computable functions. Membership in $W$ has a fixed staged semidecision, uniformly in any additional numerical parameters. The two implementation lemmas below produce programs by \cref{prop:compilers,thm:recursion}; neither is a premise of the intrinsic theorem.

\paragraph{Bounded stage tests.}
Let $A(w,e,s)$ be a decidable stage test saying that a specified semidecision with code $w$ has accepted $e$ by stage $s$. Let $g_{a,s}$ be the increasing primitive recursive finite approximation to $\varphi_a$ used in \cref{sec:numericalcodes}.

The test $A$ includes bounded simulation, so no stage waits for the semidecision to halt. Parameters may include names from any partitioned parameter assembly.

\begin{lemma}[Monotonicity from extensionality]\label{lem:automaticmonotone}
If $f,g$ are partial computable, $f\subseteq g$, and $f\in\mathcal A$, then $g\in\mathcal A$. A witness of membership can be transformed uniformly using indices of $f,g$ and the given semidecision.
\end{lemma}
\begin{proof}
\leavevmode
\begin{enumerate}[label=\textup{(\roman*)},leftmargin=*,itemsep=3pt,topsep=4pt]
\item \textbf{Monotonicity from extensionality.} Given indices $a,b$ of $f,g$, parameterized recursion constructs $e=e(w,a,b)$.
At stage $s$, test $A(w,e,s)$.
Before its first hit emit $g_{a,s}$; from that hit onward emit $g_{b,s}$.
Retain all previously emitted pairs.
On input $n$, search this enumeration for $(n,m)$ and return $m$.
The graph is functional under $f\subseteq g$.
The index constructor is total on every numerical tuple; the inclusion is used only to establish its interpreted behavior.
\item \textbf{The no-hit case.} If the hit never occurs, $\varphi_e=f$.
Extensionality and $f\in\mathcal A$ then imply $e\in W$, contradicting absence of a hit.
Hence the search terminates.
After the hit, the full emitted graph is $g$, so $e\in W$ implies $g\in\mathcal A$ by extensionality.
This reasoning is in the external metatheory.
The witness transformation itself is the explicit search for the hit, followed, if a hit witness at a specified index of $g$ is required, by search for that index in $W$.
\item \textbf{Tracked implementation.} The proof establishes termination on every input satisfying the stated premise, and thus gives the required realizer without invoking an internal Markov principle.
\end{enumerate}
\end{proof}

\begin{lemma}[Finite witness extraction]\label{lem:finitewitness}
If $f\in\mathcal A$, some finite subgraph $g\subseteq f$, regarded as a partial function, belongs to $\mathcal A$. Such a graph and its membership witness can be found uniformly from an index of $f$, the semidecision, and the membership premise.
\end{lemma}
\begin{proof}
\leavevmode
\begin{enumerate}[label=\textup{(\roman*)},leftmargin=*,itemsep=3pt,topsep=4pt]
\item \textbf{Finite witness extraction.} Parameterized recursion constructs $e=e(w,a)$.
Maintain a finite graph, initially empty.
At stage $s$, first test $A(w,e,s)$; on a hit freeze the maintained graph permanently, and otherwise replace it by $g_{a,s}$.
The program on input $n$ searches the emitted graph for its value.
Every stage is bounded, and the first-hit stage determines the frozen graph effectively.
On a run with no hit the union is $f$.
\item \textbf{Verification of the equations.} If there is no hit, $\varphi_e=f$, contradicting extensionality and $f\in\mathcal A$.
Thus there is a hit and the final graph is a finite $g\subseteq f$.
Since $e\in W$ and $\varphi_e=g$, extensionality gives $j(g)\in W$.
The procedure searches for the first hit, returns the recorded finite graph, and searches for a hit at its canonical index $j(g)$ if an explicit witness is required.
All of these searches terminate by the preceding argument and are uniform in the original parameters.
\end{enumerate}
\end{proof}

\begin{theorem}[Finite-information theorem]\label{thm:finiteinformation}
For an extensional c.e. index set $W$, define the c.e. set of finite graphs
\[
 S_W=\{g\in\FG\mid j(g)\in W\}.
\]
For every partial computable $f$,
\begin{equation}\label{eq:rice-shapiro}
 f\in\mathcal A\quad\Longleftrightarrow\quad
 \exists g\in S_W\;(g\subseteq f).
\end{equation}
The translations of positive witnesses in both directions are uniform. Consequently restriction to effective points gives an isomorphism in $\Eff$
\begin{equation}\label{eq:observationiso}
 \OO_{\FunDom}\cong\Sigma^{\PC}.
\end{equation}
\end{theorem}
\begin{proof}
\leavevmode
\begin{enumerate}[label=\textup{(\roman*)},leftmargin=*,itemsep=3pt,topsep=4pt]
\item \textbf{Finite-information theorem.} By \cref{prop:decisiondescent}, the saturated classifier $\N\to\Sigma$ descends to a map $u:\PC\to\Sigma$.
Apply \cref{thm:observationfullness}.
Its finite trace is
\[
 S_u(g)=u(\downarrow g)=u(q(j(g))),
\]
which is exactly the predicate $j(g)\in W$.
The trace-extension formula gives \eqref{eq:rice-shapiro}, and \eqref{eq:observationiso} is the intrinsic isomorphism already proved there.
\item \textbf{Converse construction.} Conversely an upper positive trace pulls back along $q$ to an extensional c.e. predicate, by joint search for its finite graph and the corresponding computation witnesses.
The concrete witness transformations are \cref{lem:automaticmonotone,lem:finitewitness}; they realize the two directions of this equality uniformly.
\end{enumerate}
\end{proof}

The resulting index theorem implements the intrinsic one. Scott continuity is a consequence of the internal semantics, and no continuity premise is imposed on an index predicate. The following parameter calculation records the realizers of this already established object-level correspondence.

\subsection{Parameterized witness adequacy}
\begin{proposition}[Uniform implementation of the trace isomorphism]\label{prop:finite-information-parameters}
Let $P$ be partitioned, with naming map $\varepsilon$, and let
$\theta(\varepsilon(p),e,s)$ be a tracked total stream for the pullback of $u:P\times\PC\to\Sigma$. The trace
\[
 S_p(g)\iff\exists s\;\theta(\varepsilon(p),j(g),s)=1
\]
and its extension implement the parameterized instance of \cref{thm:observationfullness}. Their positive witnesses are transformed uniformly, and the construction descends along parameter covers.
\end{proposition}
\begin{proof}
\leavevmode
\begin{enumerate}[label=\textup{(\roman*)},leftmargin=*,itemsep=3pt,topsep=4pt]
\item \textbf{Uniform implementation of the trace isomorphism.} For each valid parameter name, the index support is extensional.
Switching transforms a hit at $j(g)$ along $g\subseteq h$, proving upperness of $S_p$.
Freezing, applied to a hit at an input index $e$, returns a finite $g\subseteq\varphi_e$ and a hit at $j(g)$; finite simulation supplies the containment witnesses.
\item \textbf{Converse construction.} Conversely those witnesses and a hit for $S_p(g)$ allow switching to find a hit at $e$.
The bounded algorithms above are uniform in the parameter name and terminate on the stated proof inputs, so they define maps of proof-name objects.
\item \textbf{Inverse comparison.} Equal parameter denotations give equal predicates $u$, hence equal traces by restriction along $qj$.
Extension respects those equalities by the witness transformations.
Kernel-pair descent therefore gives the construction on every parameter object.
The inverse identities are those of the intrinsic trace isomorphism: restriction recovers $S_p$, and directed reconstruction recovers the observation on every ideal, including those without a numerical global name.
\end{enumerate}
\end{proof}

\subsection{Computably enumerable subsets}
For the semantic finite-subset domain $\mathcal W$ of \cref{sec:enumerablesemantic}, we have $\Pt\mathcal W\cong\Sigma^\N$.
Program domains give the numerical enumeration $e\mapsto W_e=\dom(\varphi_e)$ of effective points. To present any semidecidable subset, search its membership stream on input $n$ and return $0$ on a hit; thus this enumeration is a cover by the same projective argument as \cref{prop:codecover}.

\begin{corollary}[Rice--Shapiro for computably enumerable sets]\label{cor:rssets}
An extensional c.e. property of indices for c.e. subsets of $\N$ has the form
\[
 W_e\in\mathcal A
 \quad\Longleftrightarrow\quad
 \exists F\in\mathcal S\;(F\subseteq W_e)
\]
for a c.e. family $\mathcal S$ of finite subsets. Equivalently, in $\Eff$,
\[
 \OO_{\mathcal W}\cong\Sigma^{\Pt(\mathcal W)}.
\]
\end{corollary}
\begin{proof}
\leavevmode
\begin{enumerate}[label=\textup{(\roman*)},leftmargin=*,itemsep=3pt,topsep=4pt]
\item \textbf{Observation preservation.} Descend the saturated index predicate to $W=\Sigma^\N$ and use the finite-subset instance of \cref{thm:observationfullness}.
Its trace consists of the finite subsets satisfying the descended predicate.
For explicit witness transformations, use increasing finite enumerations in the switching and freezing constructions.
Inclusion ensures correctness of switching, and freezing returns the required finite subset.
These operations are uniform on presented parameters.
\end{enumerate}
\end{proof}

\section{Rice's theorem and effective continuous operators}\label{sec:rice-ms}

The numerical coding theorem permits the intrinsic results to be read as theorems about partial computable functions. We first descend decidable extensional predicates to obtain Rice's theorem, then compare effective operators with the finite-information maps arising from the semantic function domain. The point is to relate index transformations to continuity without replacing one notion by the other.


\begin{theorem}[Rice]\label{thm:rice}
Every decidable extensional property of partial computable functions is constant. Hence the index set of a nontrivial extensional property is not computable. The analogous statement holds for c.e. subsets of $\N$.
\end{theorem}
\begin{proof}
\leavevmode
\begin{enumerate}[label=\textup{(\roman*)},leftmargin=*,itemsep=3pt,topsep=4pt]
\item \textbf{Rice.} The classifier of an extensional decidable index predicate equalizes the kernel pair of the program cover.
By \cref{prop:decisiondescent}, it descends to a decision $P\to2$, which is constant by \cref{thm:intrinsichalting}.
The finite-subset statement uses the same descent to $W$.
\end{enumerate}
\end{proof}

This is the presentation form of the intrinsic decision theorem. The observation-theoretic argument of \cref{thm:syntheticrice} gives its order-theoretic explanation once observation fullness has identified all ambient decisions.

\begin{theorem}[Myhill--Shepherdson]\label{thm:ms}
Let $t:\N\to\N$ be total computable and extensional in the sense that
\[
 \varphi_e=\varphi_d\Longrightarrow\varphi_{t(e)}=\varphi_{t(d)}.
\]
The induced operation on partial computable functions extends uniquely to an effective Scott-continuous operation on the dcpo of all partial functions $\N\pmap\N$. Categorically, restriction induces a bijection
\begin{equation}\label{eq:operatorbijection}
 \ED(\Eff)(\FunDom,\FunDom)
 \cong\Eff(\PC,\PC).
\end{equation}
\end{theorem}
\begin{proof}
\leavevmode
\begin{enumerate}[label=\textup{(\roman*)},leftmargin=*,itemsep=3pt,topsep=4pt]
\item \textbf{Myhill--shepherdson.} The bijection \eqref{eq:operatorbijection} is \cref{thm:syntheticms}.
The extensional map $qt$ descends through the kernel pair of $q$ to $f:\PC\to\PC$; conversely a map of effective points lifts through the modest program-name cover by \cref{lem:modest-target-lift}.
We calculate the matrix of its unique effective extension.
\item \textbf{Partiality and definedness.} Write $F(\varphi_e)=\varphi_{t(e)}$.
For each finite output graph $c$, the predicate
\[
 c\subseteq F(\varphi_e)
\]
is c.e., uniformly in $(e,c)$, and extensional in $e$.
\Cref{thm:finiteinformation} applies.
In particular these predicates imply that $F$ is monotone on partial computable functions: every finite piece of $F(f)$ is a finite positive observation and persists under $f\subseteq g$.
\item \textbf{Ideal and directedness calculation.} Define the semidecidable relation on finite graphs
\begin{equation}\label{eq:operatorrelation}
 M(b,c)\iff c\subseteq F(b),
\end{equation}
where $b$ is interpreted by its canonical index $j(b)$.
Each row is the ideal of finite subgraphs of the partial function $F(b)$, and monotonicity of $F$ makes the family monotone in $b$.
Thus $M$ is a $\Sigma$-classified Kleisli arrow on $\FG$.
\item \textbf{Scott continuity.} By \cref{thm:relations} it extends to an effective Scott-continuous map
\begin{equation}\label{eq:msextension}
 \overline F(I)(c)\iff\exists b\,[I(b)\wedge M(b,c)].
\end{equation}
On an effective program $f$, the finite-information theorem for the predicate $c\subseteq F(-)$ gives
\[
 c\subseteq F(f)\iff\exists b\subseteq_{\mathrm{fin}}f\;(c\subseteq F(b)),
\]
so \eqref{eq:msextension} agrees with $F$.
Uniqueness follows because every ideal is the directed union of its finite principal ideals.
\item \textbf{Identification of the resulting maps.} Externally the extension has the equivalent formula
\[
 \overline F(f)=\bigcup_{b\subseteq_{\mathrm{fin}}f}F(b).
\]
Its values are functional even for noncomputable $f$: two finite input subgraphs have a finite compatible union $b$, and monotonicity places both corresponding output pieces in $F(b)$.
\item \textbf{Tracked implementation.} The matrix construction is uniform in the parameters, and the displayed agreement on effective points identifies it with the intrinsic extension.
Thus the categorical bijection and its numerical realization are the same map.
\end{enumerate}
\end{proof}

\begin{corollary}\label{cor:mssets}
Every tracked extensional operation on c.e. subsets of $\N$ has a unique effective Scott-continuous extension to $\Pow(\N)$. It is given by the union of its values on finite subsets of the argument.
\end{corollary}
\begin{proof}
\leavevmode
\begin{enumerate}[label=\textup{(\roman*)},leftmargin=*,itemsep=3pt,topsep=4pt]
\item \textbf{Ideal and directedness calculation.} Apply \cref{cor:rssets} to each finite output observation, and repeat the ideal-valued relation construction \eqref{eq:operatorrelation}.
Finite unions of finite input subsets supply directedness.
The effective map is its Kleisli extension, and finite subset density proves uniqueness.
\end{enumerate}
\end{proof}


\subsection{Worked observations and operators}
Several examples make the finite-information theorems concrete. For fixed numerals $n,m$, the observation
\[
 V_{n,m}(I)=I(\{(n,m)\})
\]
tests whether a partial function returns $m$ at $n$. Its code predicate is the c.e. set of indices for which that computation terminates with that value. Its upper finite trace is the decidable family of graphs containing $(n,m)$. The observation of being defined at $n$ is $\bigvee_{m:\N}V_{n,m}$; overtness of $\N$ supplies its semidecision. A conjunction of finitely many such tests records one finite graph, provided their outputs are consistent.

Undefinedness at a fixed input fails upper closure: the empty function satisfies it, while its extension defined at that input does not. Consequently the predicate is not upward closed. Nor is totality an extensional c.e. property. If it were, \cref{lem:finitewitness} applied to any total computable function would produce a finite partial function that was total, which is impossible on $\N$. These are direct uses of the finite-witness conclusion, not additional undecidability assumptions.

On the subset domain $\mathcal W$, the basic test at a finite set $F$ is $F\subseteq A$. A semidecision runs the finitely many membership searches in parallel until each has a hit. A general observation searches for one finite $F$ in its c.e. trace and simultaneously verifies $F\subseteq A$. The witness is the finite set together with its membership hits, exactly as predicted by \cref{cor:rssets}.

\begin{example}[A direct-image operator]\label{ex:directimage}
Let $R\subseteq\N\times\N$ be c.e., and define on arbitrary subsets
\[
 T_R(A)=\{m\mid\exists n\,[n\in A\wedge R(n,m)]\}.
\]
This is monotone and preserves directed unions by interchange of existential quantifiers. Its finite trace matrix is
\[
 M_R(F,G)\iff
   \forall m\in G\;\exists n\in F\;R(n,m).
\]
Both quantifiers displayed here range over finite sets, so this is a c.e. predicate of finite codes. Its rows are ideals of finite subsets of $T_R(F)$, and its dependence on $F$ is monotone. Hence it is an effective map by the algebraic instance of \cref{thm:relations}. On a name enumerating $A$, a program dovetails that enumeration with the enumeration of $R$, outputting $m$ when a matching pair is discovered. The resulting compiler is an implementation of $\Pt(T_R)$ over the enumeration cover of $\Pt(\mathcal W)$. Its adequacy follows from the same existential formula.
\end{example}

Finite-information operators need not preserve arbitrary unions. For example, define
\[
 T(A)=\{n\in\N\mid n=0\wedge0\in A\wedge1\in A\}.
\]
This positive formula defines an effective operator. It preserves directed unions: two membership witnesses in a directed union occur together at one stage. It does not preserve the union of $\{0\}$ and $\{1\}$, since each separate output is empty while the output at their union contains $0$. Its single nontrivial finite input trigger is $\{0,1\}$. Thus the relevant categorical preservation law is directed continuity, encoded by ideal suprema.

\begin{example}[A nonalgebraic observation]
For the rounded rational domain of \cref{ex:lowerreals}, a basic test is $a\in I$. Its finite information is a rational lower bound witnessed in the cut; the completed domain has no compact elements. The relation $a<b$ makes the coordinate trace rounded: $a<b$ if and only if some rational $c$ satisfies $a<c<b$. A semidecidable cut can therefore be observed and transported by the same matrix calculus although this domain has no compact elements. This is a concrete use of the continuous part of the theory which the algebraic partial-function and subset-domain examples alone would not supply.
\end{example}


\section{Interpretation of the numerical implementation}\label{sec:recovery}


The abstract and numerical constructions now match at every stage. In $\Eff$, the presentation $q:\N\onto\PC$ has program indices as names and partial computable functions as denotations. The universal evaluator is the pullback of semantic evaluation; specialization and composition are arrows between the corresponding presentations. The recursion theorem produces equality after interpretation. Effective Scott observations coincide with extensional semidecidable properties of codes, and effective Scott-continuous operators coincide with tracked extensional operations.

\begin{center}
\begin{tabular}{@{}L{0.43\textwidth}L{0.49\textwidth}@{}}
\toprule
Categorical construction & Numerical interpretation\\
\midrule
$\Sigma$-classified subobject & Semidecidable predicate on names\\
$\Sigma$-partial map & Partial computable function on numerical data\\
Presentation $q:\N\onto\PC$ & Numbering of partial computable functions\\
Pullback of semantic evaluation & Universal partial computable function\\
Specialization of program presentations & $s$-$m$-$n$ compiler\\
Fixed point modulo the kernel pair of $q$ & Kleene's recursion theorem\\
Complete Elgot iteration & Compiled state-update loop with finite-return extent\\
Coproduct execution trace & Feedback wrapper on tagged interfaces\\
Finite execution certificates (\cref{thm:operational-certificates}) & The same semidecidable termination predicate\\
Observation isomorphism \eqref{eq:observationiso} & Finite-information / Rice--Shapiro theorem\\
Extension bijection \eqref{eq:operatorbijection} & Myhill--Shepherdson theorem\\
\bottomrule
\end{tabular}
\end{center}

These identifications are backed by the constructions and proofs above; they do not identify arbitrary global points of the underlying $\Omega$-dcpos with computable elements. Effectivity is expressed by the specified observation structure, its effective-point object, and the chosen code presentation.

\begin{proposition}[Ground implementations and oracle extension]\label{prop:relativization}
The numerical interpretation equations are equations in $\Eff$. Their transport to a forcing sheaf topos preserves their actual domains and values. Adaptive oracle operations are supplied by the query and protocol constructions of \cref{part:modalities}; their name comparisons are proved in \cref{sec:intensional-basechange}.
\end{proposition}
\begin{proof}
\leavevmode
\begin{enumerate}[label=\textup{(\roman*)},leftmargin=*,itemsep=3pt,topsep=4pt]
\item \textbf{Pullback construction.} Apply the regular-image and pullback description of an implementation to its domain--value span. A geometric inverse image preserves these finite-limit and regular-image equations, as proved in \cref{thm:partial-transport,thm:intensional-pres-fibration}.
For an oracle query, the continuation retains its answer input and its code; \cref{thm:intensional-oracle-coding} gives the resulting sheafified denotation.
These constructions compare the same ground code with its transported execution and the additional oracle protocol.
\end{enumerate}
\end{proof}


\section{Presented arithmetic before degree theory}\label{sec:presented-arithmetic}

The oracle constructions will refer to arithmetical predicates as well as to programs. To keep these notions distinct, we present the arithmetical hierarchy directly in the subobject fibration on numerical contexts. Quantifier operations then have explicit interpretations before any quotient by Turing equivalence is considered.


Use the subobject fibration on numerical contexts $\N^r$ in $\Eff$.
For a coordinate projection $\pi$, its adjoints are
\[
 \exists_\pi\dashv\pi^*\dashv\forall_\pi.
\]
Let $\Dec(\N^r)$ consist of predicates supplied by actual total Boolean
programs. Define syntactic classes, with their interpretations, by
\[
 \SigA_0=\PiA_0=\Dec,\qquad
 \SigA_{n+1}=\{\exists z\,R:R\in\PiA_n\},\qquad
 \PiA_{n+1}=\{\forall z\,R:R\in\SigA_n\}.
\]
Finite tuples of quantified numbers are coded by one number. These are
classes of presented predicates: the program and formula are retained.
The level-zero convention uses computable matrices. For positive levels it gives the
arithmetical classes in the classical numerical interpretation.

\begin{proposition}[The first existential level]\label{arith:prop:first-level}
On the ground numerical contexts, interpreted $\SigA_1$ predicates are
exactly the $\Sigma$-classified predicates of \cref{def:finitehit}, uniformly on
program presentations.
\end{proposition}
\begin{proof}
\leavevmode
\begin{enumerate}[label=\textup{(\roman*)},leftmargin=*,itemsep=3pt,topsep=4pt]
\item \textbf{The first existential level.} For $R(x)=\exists s\,D(x,s)$, run the total Boolean program for $D$ on
successive $s$ and halt on a positive answer.
This is a semidecision
program, and its domain is the indicated existential image.
\item \textbf{Converse construction.} Conversely, the finite-hit epimorphism $2^\N\twoheadrightarrow\Sigma$
and \cref{lem:modest-target-lift} lift a map
$\N^r\to\Sigma$ to a tracked stream family
$\alpha:\N^r\times\N\to2$.
Its hit support is
$\exists s\,[\alpha(x,s)=1]$.
This is the required computable matrix.
The use of the lifting lemma is on the partitioned assembly $\N^r$;
no lift for arbitrary parameter objects is being claimed.
\end{enumerate}
\end{proof}

Substitution and quantification have Beck--Chevalley equations
inside this subobject fibration. This constructs the hierarchy without
forming any quotient by mutual computability. It does not assert
intuitionistic prenex or complement identities that are false without
further interpretation. The closed interpretation and the hierarchy--jump theorem are constructed in \cref{sec:arithmetic-oracles}. The constructive analysis in \cite{forster-kirst-mueck} is a useful
reference for this precise distinction.


The hierarchy is now defined by actual predicate presentations and numerical quantifiers, independently of degree equivalence. This makes it possible to compare a forcing localization with a specific oracle interpreter, rather than using a degree label as a substitute for a construction.

\section{From effective domains to oracle observations}\label{sec:domain-forcing-bridge}


The preceding development supplies the computational objects used in the forcing constructions. We now make their interface explicit. The notation $\Sc(D)$ denotes the frame of Scott opens, and $\ScLoc(D)$ the locale with that frame. The forcing topos is written $\FF$; the semantic free-domain functor remains $\F$.

\begin{proposition}[Overtness of numerically presented bases]\label{prop:enumerated-overt}
A regular-epimorphic image of $\N$, and a $\Sigma$-classified subobject of $\N$, is $\Sigma$-overt. The corresponding relative assertions hold in parameter slices.
\end{proposition}
\begin{proof}
\leavevmode
\begin{enumerate}[label=\textup{(\roman*)},leftmargin=*,itemsep=3pt,topsep=4pt]
\item \textbf{Observation preservation.} Overtness of $\N$ is \cref{thm:relativefoundation}.
If $e:\N\onto B$, define
$\exists_B^\Sigma u=\exists_\N^\Sigma(u e)$. Regular
epimorphy gives equality of its image in $\Omega$ with
$\exists_Bu$, and monicity of $\Sigma\incl\Omega$ gives its
unique factorization, uniformly in parameters. On a partitioned
parameter cover, lift the pulled-back predicate $u(e(n))$ to a
tracked stream $\theta(n,k)$. The dovetailed stream
$\gamma(\langle n,k\rangle)=\theta(n,k)$ has a hit exactly
when $u$ holds at some $b:B$. Decoding a hit returns $e(n)$ and
the positive witness; encoding a supplied witness uses a local
preimage under $e$ and a hit for its stream. Their regular images
agree, so the resulting classifier descends, without a selected
section of $e$.
\item \textbf{Existence and uniqueness.} If $A\incl\N$ is classified by $p:\N\to\Sigma$,
first form the dependent conjunction
$\exists a:p(n)\,u(n,a)$ and then quantify over $n$.
Dominance and overtness of $\N$ give the positive classifier;
Beck--Chevalley and uniqueness of its factorization prove the
relative assertions.
\end{enumerate}
\end{proof}

\begin{theorem}[Enriched observations and forcing markings]\label{thm:domain-marking-bridge}
Let $C=(B,r)$ be a searchable approximation presentation. Its canonical positive trace marking, used in \cref{prop:canonical-positive-marking}, is exactly its effective observation object:
\[
 M_C=\Sc(D_C)\times_{\Omega^B}\Sigma^B=\OO_C.
\]
Writing $A_C=\Pt D_C$, the semantic comparison identifies these observations with the $\J$-algebra maps $A_C\to\Sigma$. Evaluation is
\[
 \Obs_C(I,U)=\exists_B^\Sigma b\,[I(b)\wedge U(r(b))].
\]
For an effective map $F:D_C\to D_{C'}$, inverse image on markings is its observation transpose, and
\[
 \Obs_C(I,F^*U)=\Obs_{C'}(\Pt F(I),U).
\]
These equations commute with composition, products, parameters, and the comparison $\Pt(D_\bot)\cong\LL\Pt D$.
\end{theorem}
\begin{proof}
\leavevmode
\begin{enumerate}[label=\textup{(\roman*)},leftmargin=*,itemsep=3pt,topsep=4pt]
\item \textbf{Pullback construction.} The first equality is the defining pullback of \cref{def:O-serachable-presentation}.
Rounded trace reconstruction, \cref{prop:roundedtrace}, gives the displayed existential formula.
Its positive classifier is furnished by overtness, and \cref{thm:compatibility,lem:roundedadmission} identify its effective points with the split $\J$-algebra.
The observation identification is the intrinsic observation comparison of \cref{jd:thm:splitcomparison}; equivalently it is the positive character reconstruction of \cref{jd:thm:spectra}.
\item \textbf{Evaluation and transposition.} The transpose equation is ordinary evaluation after inverse image, restricted through the monomorphism $\Sigma\incl\Omega$.
For positive rounded distributors it is also Fubini's equation for relational composition, so it is joint in all parameters.
Product compatibility is \cref{prop:productpresentation}.
For a lifted point $(p,v)$, both sides evaluate an output observation as $\exists a:p\,U(v(a))$, by \cref{thm:effectivelift}.
Hence the partial comparison includes the same extent as well as the same value.
\end{enumerate}
\end{proof}

The later construction of $M_C$ also applies when a presentation is not searchable. The theorem identifies the entire searchable computational interface with that construction; it does not turn arbitrary refinement or arbitrary frame truth into a positive predicate.

\begin{proposition}[Refinement and answer selection in $\Eff$]\label{prop:nonautomatic-effective-data}
There is no uniform positive classifier for every displayed preorder relation, and not every regular-epimorphic answer map in $\Eff$ has a section.
\end{proposition}
\begin{proof}
\leavevmode
\begin{enumerate}[label=\textup{(\roman*)},leftmargin=*,itemsep=3pt,topsep=4pt]
\item \textbf{Refinement and answer selection in.} For a generalized truth $u:\Omega$, put on $2$ the reflexive order with $0\leq1$ precisely when $u$, and with $1\leq0$ false. A positive classifier for these relations jointly in $u$ would give a retraction of $i:\Sigma\incl\Omega$, making $i$ an isomorphism.
That would classify the complement of halting positively, contrary to \cref{thm:intrinsichalting}.
\item \textbf{Tracked implementation.} For the second assertion use the program cover $q:\N\onto\PC$ constructed in \cref{sec:numericalcodes}. A section $s$ would be monic.
Since equality of numbers is decidable, the predicate $f\mapsto[s(f)=s(\bot)]$ would be a map $\PC\to2$ taking different values on the undefined function and the constant-zero function.
This contradicts \cref{thm:intrinsichalting}.
Thus the answer cover $q$ has no global selector.
\end{enumerate}
\end{proof}

\begin{center}
\begin{tabular}{L{.28\textwidth}L{.64\textwidth}}
\toprule Notation & Meaning\\\midrule
$\J,\AS$ & Admitted-ideal monad and its algebra category\\
$\F,\Pt,\LL$ & Semantic free domain, effective points, positive partiality\\
$\FunDom,\PC$ & Full numerical partial-function domain and its effective points\\
$\mathcal W,W$ & Full subset domain and its effective points\\
$\Sc(D),\ScLoc(D)$ & Scott frame and its locale\\
$\OO_D,\Pt D$ & Positive observations and their effective points\\
$j,\Omega_j,S_j$ & Nucleus, fixed frame, stable effective marking\\
$\FF_j,\pi_j$ & Forcing sheaf topos and projection to $\Eff$\\
$\rho:L_j\to X$ & Realization in the oracle locale\\
$f^*,f^\#$ & Inverse image of objects and of internal locales\\
$\Theta_f$ & Imported semidecision truth in the forcing topos\\
\bottomrule
\end{tabular}
\end{center}

\part{Presentations of forcing spaces}\label{part:presentations}

We now turn from computational domains to the sites that present forcing information. The same order and ideal constructions have a geometric interpretation: upper predicates form a frame, covers act by a nucleus, and continuous ideals give generic models. The purpose is to derive the forcing locale from the condition data and to compare it with the Scott locale of a domain completion.


\section{Conditions, covers, and basis reflection}\label{sec:basis}

A condition order records refinement, while a coverage specifies which collections of refinements jointly suffice. Neither alone should be confused with the eventual forcing truth value. We use the upper-predicate frame to construct the coverage nucleus and recover the forcing frame from the presented condition data.


A condition preorder is an internal preorder $P$ with the convention that $p\leq q$ means that $q$ strengthens $p$. A predicate $U:P\to\Omega$ is upper when $U(p)$ and $p\leq q$ imply $U(q)$. The object $\Up(P)$ is a frame: joins are pointwise existential disjunctions and finite meets are pointwise conjunctions. Write $\uparrow p=\{q\mid p\leq q\}$. Every upper predicate has the expansion
\begin{equation}\label{eq:upper-expansion}
 U=\bigvee_{p:U(p)}\uparrow p.
\end{equation}
The implication from right to left is upperness; the other uses the witness $p=q$ at $q\in U$.

\begin{definition}[Refinement coverage]\label{def:site}
A sieve at $p$ is an upper subobject $S\subseteq\uparrow p$. A refinement-covariant Grothendieck topology $J$ assigns covering sieves and satisfies: $\uparrow p$ covers $p$; a cover $S$ of $p$ restricts to a cover $S\cap\uparrow q$ of $q$ for $q\geq p$; and if $S$ covers $p$ and $R\cap\uparrow q$ covers $q$ for every $q\in S$, then $R$ covers $p$. Here $R$ is a sieve at $p$.
\end{definition}

Cover enlargement follows from these clauses. If $S\subseteq R$ and $S$ covers $p$, every $q\in S$ has $\uparrow q\subseteq R$, so local character makes $R$ cover $p$. Covers are closed under finite intersection: restrict one cover to each member of the other and apply local character. These are internal proofs and apply to families of covers.

\begin{proposition}[Coverage saturation]\label{prop:site-nucleus}
The operation
\[
 j_J(U)(p)\iff U\cap\uparrow p\in J(p)
\]
is a nucleus on $\Up(P)$.
\end{proposition}
\begin{proof}
\leavevmode
\begin{enumerate}[label=\textup{(\roman*)},leftmargin=*,itemsep=3pt,topsep=4pt]
\item \textbf{Coverage saturation.} If $p\leq q$, pullback stability makes a covering $U\cap\uparrow p$ restrict to a covering $U\cap\uparrow q$; hence $j_J(U)$ is upper.
Cover enlargement proves monotonicity.
If $U(p)$, upperness gives $\uparrow p\subseteq U$, and the maximal cover gives $j_J(U)(p)$.
\item \textbf{Localization and coverage.} Finite intersection of covers gives $j_J(U)\cap j_J(V)\subseteq j_J(U\cap V)$.
Monotonicity gives the reverse inclusion.
The top predicate is fixed.
Finally, if $j_Jj_J(U)(p)$, the sieve $j_J(U)\cap\uparrow p$ covers $p$, and each of its members has its restriction of $U$ covering.
Local character gives $j_J(U)(p)$.
Inflationarity supplies the reverse inequality.
Thus saturation is idempotent and finite-meet preserving.
\end{enumerate}
\end{proof}

\begin{definition}[Nucleus and forced frame]\label{def:nucleus}
A nucleus on a frame $H$ is a monotone, inflationary, idempotent map $j:H\to H$ preserving finite meets. Put $\Omega_j=\Fix(j)$, let $i_j:\Omega_j\incl H$ be the inclusion, and let $q_j:H\to\Omega_j$ have underlying map $j$. The locale with frame $\Omega_j$ is denoted $L_j$.
\end{definition}

\begin{proposition}[The fixed-frame reflection]\label{prop:fixed-frame}
The fixed points form a frame, with meets inherited from $H$ and joins
\[
 \bigvee\nolimits^{\Omega_j}_t a_t=j\Bigl(\bigvee\nolimits^H_t a_t\Bigr).
\]
Its bottom is $j0$ and its top is $1$. The map $q_j$ is a surjective frame homomorphism, $q_j\dashv i_j$, $q_ji_j=\id$, and $i_jq_j=j$. For fixed $a,b$ the implication in the fixed frame is the ambient implication $a\Rightarrow_H b$.
\end{proposition}
\begin{proof}
\leavevmode
\begin{enumerate}[label=\textup{(\roman*)},leftmargin=*,itemsep=3pt,topsep=4pt]
\item \textbf{The fixed-frame reflection.} If all $a_t$ are fixed, then $j(\bigwedge_ta_t)\leq a_t$ for every $t$, so the meet is fixed.
The displayed join is fixed, is an upper bound, and is below every fixed upper bound by monotonicity.
If $a$ is fixed, finite-meet preservation gives
\[
 a\wedge j\Bigl(\bigvee_t b_t\Bigr)
 =j\Bigl(a\wedge\bigvee_t b_t\Bigr)
 =j\Bigl(\bigvee_t(a\wedge b_t)\Bigr),
\]
which is the frame distributive law.
The same calculation proves finite-meet preservation of $q_j$; preservation of joins follows from $j(\bigvee_tj x_t)=j(\bigvee_t x_t)$.
\item \textbf{Equality of the two traces.} To check the latter equality, use $x_t\leq jx_t$ in one direction, and $jx_t\leq j(\bigvee_t x_t)$ in the other.
\item \textbf{Adjunction and universal property.} For fixed $a$, $jx\leq a$ is equivalent to $x\leq a$, proving the adjunction.
The two composite equations follow immediately.
Finally
\[
 j(a\Rightarrow b)\wedge a
 =j((a\Rightarrow b)\wedge a)\leq jb=b.
\]
Thus $j(a\Rightarrow b)\leq a\Rightarrow b$, so this implication is fixed.
Its universal property is inherited from $H$.

The defining maps satisfy the following commutative diagram
\[
\begin{tikzcd}[column sep=large,row sep=large]
H\ar[r,"q_j"]\ar[dr,"j"']&H_j\ar[d,hook,"i_j"]\\ &H
\end{tikzcd}
\]
\end{enumerate}
\end{proof}

\subsection{Extension-density and its localization}\label{sec:extension-density}
Put $H=\Up(P)$. Define
\begin{equation}\label{cal:eq:density}
 \Dens U(p)\quad\Longleftrightarrow\quad
 \forall q\geq p\;\exists r\geq q\; U(r).
\end{equation}
All quantifiers in this formula are internal. In particular, the existential
quantifier is not replaced by its double negation.

\begin{theorem}\label{cal:thm:density}
The operation $\Dens$ is a dense nucleus on $H$. It fixes each constant truth
predicate $U_v(p)=v$, for $v:\Omega_{\E}$. Its associated covering sieves are
exactly the upper sieves $S\subseteq\uparrow p$ satisfying
\[
 \forall q\geq p\;\exists r\geq q\;(r\in S).
\]
Moreover
\begin{equation}\label{cal:eq:nn}
 (\neg_H\neg_H U)(p)
 \quad\Longleftrightarrow\quad
 \forall q\geq p\;\neg\neg\exists r\geq q\;U(r),
 \qquad \Dens\leq\neg_H\neg_H.
\end{equation}
In the classical interpretation $\Dens=\neg_H\neg_H$ and
$\Fix(\Dens)=\RegOpen(P)$.
\end{theorem}
\begin{proof}
\leavevmode
\begin{enumerate}[label=\textup{(\roman*)},leftmargin=*,itemsep=3pt,topsep=4pt]
\item \textbf{Construction of the comparison.} Persistence and monotonicity follow directly from the formula.
If $U(p)$ and
$q\geq p$, use $r=q$; hence $U\leq \Dens U$.
For finite meets, start at $q\geq p$,
refine to $r\in U$, then to $s\geq r$ in $V$.
Persistence gives
$s\in U\cap V$.
Thus $\Dens U\cap \Dens V\leq \Dens(U\cap V)$; the reverse inclusion is
monotonicity.
Also $\Dens1=1$.
\item \textbf{Idempotence of density saturation.} For idempotence, if $\Dens\Dens U(p)$ and $q\geq p$, first obtain $r\geq q$ with
$\Dens U(r)$, and then obtain $s\geq r$ in $U$.
This proves $\Dens\Dens U\leq \Dens U$.
Taking $q=p$ shows $\Dens0=0$.
For a constant predicate, the same specialization
gives $\Dens U_v(p)\Rightarrow v$, and the converse uses $r=q$.
\item \textbf{Localization and coverage.} The displayed sieves satisfy the topology rules directly. A restriction to
$\uparrow q$ is dense by the same formula.
For transitivity, refine first into the
first dense sieve and then into the sieve covering that refinement.
Alternatively
these are precisely the sieves whose nucleus closure contains $p$.
\item \textbf{Existence and uniqueness.} In $H$, $\neg_H U(p)$ says $\forall q\geq p\,\neg U(q)$.
Applying this twice
and using $\neg\exists r\,U(r)\Leftrightarrow\forall r\,\neg U(r)$ yields
\eqref{cal:eq:nn}.
Classical logic removes the inner double negation, giving the
last assertion.
\end{enumerate}
\end{proof}

\begin{remark}[Why the distinction matters]
For the one-element preorder, $\Dens$ is the identity on $\Omega_{\E}$, whereas
double negation need not be the identity. Consequently extension-density cannot
be identified internally with Booleanization. This is already a counterexample
in the simplest possible condition presentation. Formula \eqref{cal:eq:density}
does not claim stability under arbitrary inverse image: it contains a universal
quantifier. When changing base, one forms the native condition site there and
uses the comparison maps of \cref{sec:geometricchange,sec:objectchange}.
\end{remark}


Write $\Jdens$ for this extension-dense topology. A requirement cover is
extension-dense when its proved refinement operation applies below every
stronger condition. The topology rules in \cref{cal:thm:density} then carry
this property to the generated coverage.

\begin{proposition}[Localization of dense requirements]\label{cal:thm:localize}
For a coverage $J_0$ generated by extension-dense requirements, its nucleus
satisfies $j_0\leq\Dens$. The restriction of $\Dens$ to
$H_{j_0}=\Fix(j_0)$ is a nucleus, with fixed frame $\Fix(\Dens)$.
Consequently there is a canonical subtopos inclusion
\[
 \Sh_\E(P,\Jdens)\longrightarrow\Sh_\E(P,J_0).
\]
Classically this is the regular-open forcing completion of the full
condition order.
\end{proposition}
\begin{proof}
\leavevmode
\begin{enumerate}[label=\textup{(\roman*)},leftmargin=*,itemsep=3pt,topsep=4pt]
\item \textbf{Localization of dense requirements.} Each generating sieve is a $\Jdens$-cover.
Closure under the topology rules
gives $J_0\subseteq\Jdens$ and hence $j_0\leq\Dens$.
Monotonicity,
inflationarity, and idempotence give
\[
 \Dens j_0=\Dens=j_0\Dens.
\]
Thus $\Dens$ acts on $H_{j_0}$ and has precisely its original fixed points.
The resulting fixed-frame quotient gives the stated localic inclusion.
The classical assertion is \cref{cal:thm:density}.
\end{enumerate}
\end{proof}

The case-specific refinement proofs below verify this construction for the
structural Cohen, Mathias, perfect-tree, closed-class, enumerated bushy, and
restrained-functional sites in the bases where their operations are defined;
their common calibration is collected in \cref{cal:cor:cases}.

\subsection{Basis reflection}
\begin{lemma}[Factorization through a forcing quotient]\label{lem:quotient-factor}
For a frame homomorphism $h:H\to A$, a factorization $h=\bar hq_j$ exists precisely when $hj=h$. It is unique.
\end{lemma}
\begin{proof}
\leavevmode
\begin{enumerate}[label=\textup{(\roman*)},leftmargin=*,itemsep=3pt,topsep=4pt]
\item \textbf{Factorization through a forcing quotient.} A factorization implies $hj=\bar hq_ji_jq_j=h$.
Conversely set $\bar h(a)=h(i_ja)$.
It preserves finite meets.
For fixed $a_t$,
\[
 \bar h\Bigl(j\bigvee_ta_t\Bigr)=h\Bigl(j\bigvee_ta_t\Bigr)
 =h\Bigl(\bigvee_ta_t\Bigr)=\bigvee_t\bar h(a_t).
\]
Hence it is a frame map and $\bar hq_j=h$.
Surjectivity of $q_j$ proves uniqueness.
\end{enumerate}
\end{proof}

\begin{theorem}[Basis reflection]\label{thm:basis}
Let $A$ be an internal frame and $d:P\to A$ satisfy
\begin{align}
 p\leq q&\ \Longrightarrow\ d_q\leq d_p,\label{eq:basis-order}\\
 1&=\bigvee_p d_p,\qquad
 d_p\wedge d_q=\bigvee_{r\geq p,q}d_r,\label{eq:basis-meet}\\
 a&=\bigvee_{d_p\leq a}d_p.\label{eq:basis-dense}
\end{align}
Define
\[
 q(U)=\bigvee_{p:U(p)}d_p,\qquad
 i(a)(p)\iff d_p\leq a.
\]
Then $q:\Up(P)\to A$ is a frame homomorphism, $q\dashv i$, $qi=\id$, and $iq$ is a nucleus. The topology
\[
 S\in J_d(p)\iff d_p=\bigvee_{q\in S}d_q
\]
has saturation $iq$, and $q,i$ induce inverse frame isomorphisms $\Fix(iq)\cong A$.
\end{theorem}
\begin{proof}
\leavevmode
\begin{enumerate}[label=\textup{(\roman*)},leftmargin=*,itemsep=3pt,topsep=4pt]
\item \textbf{Existence and uniqueness.} The order law makes $i(a)$ upper.
The definitions give
\[
 q(U)\leq a\iff \forall p\,[U(p)\Rightarrow d_p\leq a]
 \iff U\subseteq i(a).
\]
Thus $q\dashv i$, and join density gives $qi=\id$.
Preservation of joins follows by exchanging the two existential indices.
The top law is \eqref{eq:basis-meet}.
For finite meets,
\[
 q(U)\wedge q(V)
 =\bigvee_{p:U(p),\,q:V(q)}\bigvee_{r\geq p,q}d_r
 =\bigvee_{r:U(r)\wedge V(r)}d_r.
\]
Upperness proves the first containment implicit in the last equality; the choice of indices $p=q=r$ proves the other.
\item \textbf{Adjunction and universal property.} Since $i$ preserves meets as a right adjoint, $iq$ preserves finite meets.
Its inflationarity is the unit of the adjunction, and its idempotence uses $qi=\id$.
\item \textbf{Universal-property argument.} We verify the cover description.
If $d_p\leq q(U)$, distribute $d_p\wedge-$ and use the meet equation:
\[
 d_p=\bigvee_{q:U(q)}(d_p\wedge d_q)
     =\bigvee_{r\geq p,\,U(r)}d_r.
\]
The reverse implication follows because this last join is below $q(U)$.
Hence
\[
 iq(U)(p)\iff U\cap\uparrow p\in J_d(p).
\]
The maximal cover follows from $p\in\uparrow p$.
If $S$ covers $p$ and $r\geq p$, intersect its equation with $d_r$.
The meet equation expresses the result as the join of $d_t$ for $t\geq r$ in $S$, proving stability.
\item \textbf{Functoriality and coherence.} For local character, substitute the cover equation for each $q\in S$ into the cover equation for $p$.
Thus $J_d$ is a topology and has the asserted saturation.
\item \textbf{The two fixed-point identities.} Finally $iq$ fixes $i(a)$, while a fixed $U$ equals $iq(U)$.
Consequently the restrictions of $q,i$ are inverse.
They preserve the appropriate joins by the fixed-frame formula in \cref{prop:fixed-frame}.
\end{enumerate}
\end{proof}

\begin{corollary}[Canonical basis site]\label{cor:canonical-basis}
For any join-dense family $d:B\to A$, put $b\preceq c$ when $d_c\leq d_b$. This preorder satisfies the basis equations and its semantic cover topology presents $A$.
\end{corollary}
\begin{proof}
\leavevmode
\begin{enumerate}[label=\textup{(\roman*)},leftmargin=*,itemsep=3pt,topsep=4pt]
\item \textbf{Canonical basis site.} The top equation is density at $1$.
Expand $d_b\wedge d_c$ by density.
The indexing condition $d_a\leq d_b\wedge d_c$ is precisely $a\succeq b,c$.
The remaining equations are the definition and the given density.
Apply \cref{thm:basis}.
\end{enumerate}
\end{proof}

The order in this corollary is constructed from the frame. A displayed condition order can be compared with it through the basis map, but remains part of the original presentation. In particular, the formula $d_c\leq d_b$ contains a full frame inequality and does not itself produce an effective refinement test.

\begin{proposition}[Displayed basis comparison]\label{prop:displayed-basis}
Let $(P,J)$ carry a basis family $d:P\to A$ satisfying \eqref{eq:basis-order}--\eqref{eq:basis-dense}. If its covering sieves are exactly the semantic covers $J_d$, then
\[
 \Omega_J(P)=\Fix(j_J)\cong A.
\]
More generally, it suffices that every displayed cover is semantic and that the topology generated by the displayed covers contains every semantic cover.
\end{proposition}
\begin{proof}
\leavevmode
\begin{enumerate}[label=\textup{(\roman*)},leftmargin=*,itemsep=3pt,topsep=4pt]
\item \textbf{Displayed basis comparison.} Soundness gives $J\subseteq J_d$.
Detection gives the opposite inclusion, since $J$ is already a topology.
Their saturation formulas coincide.
The result is therefore the fixed-frame isomorphism of \cref{thm:basis}.
Its maps are explicitly
\[
 U\longmapsto\bigvee_{p:U(p)}d_p,\qquad
 a\longmapsto(p\mapsto d_p\leq a),
\]
so the comparison also identifies each distinguished stage.
\end{enumerate}
\end{proof}

\begin{proposition}[Requirement-generated condition sites]
\label{prop:generated-coverage}
For an internal family of proposed covers $(p_i,S_i)_{i:I}$,
where $S_i\subseteq\uparrow p_i$ is an upper sieve, there is a
least refinement-covariant Grothendieck topology containing them.
It is obtained without a chosen saturation schedule: let $j_*$
be the pointwise meet of the nuclei $k$ on $\Up(P)$ satisfying
$\uparrow p_i\leq k(S_i)$ for every $i$, and put
\[
 J_*(p)=\{S\subseteq\uparrow p\mid p\in j_*(S)\}.
\]
Then $j_{J_*}=j_*$, and $J_*$ is the least topology containing
all $(p_i,S_i)$.
\end{proposition}
\begin{proof}
\leavevmode
\begin{enumerate}[label=\textup{(\roman*)},leftmargin=*,itemsep=3pt,topsep=4pt]
\item \textbf{Localization and coverage.} The constant-top nucleus validates every proposed cover. The
internal family of validators is therefore inhabited, and
\cref{lem:generated-nucleus} constructs its least member $j_*$.
Inflationarity makes $\uparrow p$ cover $p$. If
$p\leq q$ and $p\in j_*(S)$, then $q\in j_*(S)$.
Since $j_*$ preserves finite meets and
$q\in j_*(\uparrow q)$, it follows that
$q\in j_*(S\cap\uparrow q)$. Thus restriction preserves covers.
\item \textbf{Transitivity of coverage.} If $S$ covers $p$ and $R\cap\uparrow q$ covers $q$
for every $q\in S$, then $S\subseteq j_*(R)$. Applying $j_*$
and idempotence gives
$p\in j_*(S)\subseteq j_*j_*(R)=j_*(R)$, proving local character.
\item \textbf{Localization and coverage.} By finite-meet preservation,
$j_{J_*}(U)(p)$ is equivalent to
$p\in j_*(U\cap\uparrow p)$, hence to $p\in j_*(U)$.
Every given $(p_i,S_i)$ lies in $J_*$ by its defining inequality.
Conversely, the nucleus of any topology containing these covers
validates the corresponding inequalities, so $j_*$ lies below it.
The preceding equivalence of saturations reflects inclusion of
coverages, proving leastness.
This is the internal nucleus--coverage construction; see also
\cite{johnstone,caramello}.
\end{enumerate}
\end{proof}

Thus specified requirement families generate their condition sites;
no independent coverage choice or transfinite saturation procedure
is used.


The coverage has become a nucleus on the condition frame, and the resulting fixed frame is the forcing locale. The next comparison identifies this same geometry in the Scott opens of an ideal completion, so that forcing and domain observations can be transported along a common frame isomorphism.

\section{Ideal completion and rounded classifiers}\label{sec:rounded}

The condition-site frame can also be read from a completed domain of ideals. To establish that these are two presentations of the same forcing geometry, we compare Alexandrov predicates on conditions with Scott opens on ideals. Rounded presentations then admit a corresponding geometric classifier.


The ordinary ideal doctrine of \cref{subsec:completion} supplies the following locale comparison.

\begin{proposition}[Scott--Alexandrov comparison]\label{prop:scott-alex}
The maps
\[
 \Phi_P(U)(I)\iff\exists p\in I\,U(p),
 \qquad
 \Psi_P(V)(p)\iff V(\downarrow p)
\]
are inverse frame isomorphisms $\Up(P)\cong\Sc(\Idl P)$.
\end{proposition}
\begin{proof}
\leavevmode
\begin{enumerate}[label=\textup{(\roman*)},leftmargin=*,itemsep=3pt,topsep=4pt]
\item \textbf{Scott--Alexandrov comparison.} Apply \cref{lem:tracefree}: $\Phi_P$ is its extension map and $\Psi_P$ its restriction along $y_P$.
That lemma proves both inverse identities and preservation of the frame operations.
In particular $\Phi_P(\uparrow p)(I)\iff p\in I$, by lowerness of $I$.
\end{enumerate}
\end{proof}

The locale with frame $\Up(P)$ is written $\Alex^\uparrow(P)$. The isomorphism identifies it with $\ScLoc(\Idl P)$. It identifies the basic open $\uparrow p$ with $\delta_p(I)\iff p\in I$.

For an approximation presentation $C=(B,r)$, write $D_C=D_r$ and $\kappa_C(b)=r(b)$. Its splitting and directed reconstruction are \cref{thm:rounded}; its frame of opens has the trace representation and coordinate expansion of \cref{prop:roundedtrace,eq:coordinate-expansion}. For searchable $C$, \cref{lem:roundedobs,thm:compatibility} supply the observation object and its effective-point pairing. We now use these constructions to prove the geometric classifier property.

\begin{definition}[Rounded-ideal theory]\label{def:rounded-theory}
The geometric presentation associated with $C$ has propositions $\mathsf I(b)$, indexed by $B$, and relations
\begin{align}
 1&=\bigvee_b\mathsf I(b),\label{eq:atomic-top}\\
 \mathsf I(a)\wedge\mathsf I(c)
   &=\bigvee_{d\geq a,c}\mathsf I(d),\label{eq:atomic-meet}\\
 \mathsf I(a)&=\bigvee_{\rho(a,b)}\mathsf I(b).\label{eq:atomic-round}
\end{align}
A model in a frame $A$ is a family $\alpha:B\to A$ satisfying these relations.
\end{definition}

The meet equation implies lowerness. If $a\leq c$, the index $d=c$ occurs on the right, so $\alpha_c\leq\alpha_a\wedge\alpha_c$. Inhabitedness and directedness are the first two equations. The last equation is the fixed-point condition for $r^\#$. Thus models in $\Omega$, in every parameter context, are precisely rounded ideals.

\begin{theorem}[Explicit rounded classifier]\label{thm:rounded-classifier}
The frame of the rounded-ideal presentation is $\Sc(D_C)$. More explicitly, for every frame $A$, the assignments
\[
 h\longmapsto(h(\delta_b))_{b\in B},\qquad
 \alpha\longmapsto
 \left(U\longmapsto\bigvee_{b:U(\kappa_Cb)}\alpha_b\right)
\]
give inverse, natural bijections between frame homomorphisms $\Sc(D_C)\to A$ and models of \eqref{eq:atomic-top}--\eqref{eq:atomic-round}.
\end{theorem}
\begin{proof}
\leavevmode
\begin{enumerate}[label=\textup{(\roman*)},leftmargin=*,itemsep=3pt,topsep=4pt]
\item \textbf{Explicit rounded classifier.} The coordinate opens satisfy the equations by ideality and roundedness. A frame map therefore gives a model.
Conversely let $\alpha$ be a model and define $h$ by the displayed formula.
Joins are preserved because a join of Scott opens is evaluated pointwise:
\[
 h\Bigl(\bigvee_tU_t\Bigr)
 =\bigvee_{b:\exists t\,U_t(\kappa b)}\alpha_b
 =\bigvee_t h(U_t).
\]
The first atomic equation gives $h(1)=1$.
\item \textbf{Product calculation.} For a binary meet, expand:
\[
 h(U)\wedge h(V)
 =\bigvee_{a:U(\kappa a),\,c:V(\kappa c)}
       \bigvee_{d\geq a,c}\alpha_d.
\]
Monotonicity of $\kappa$ makes every displayed $\alpha_d$ occur in $h(U\wedge V)$.
The reverse inequality follows from monotonicity of $h$ applied to the two projections $U\wedge V\leq U,V$.
Thus $h$ is a frame map.
Its value on a coordinate is
\[
 h(\delta_a)=\bigvee_{\rho(a,b)}\alpha_b=\alpha_a.
\]
Finally \eqref{eq:coordinate-expansion} forces any frame map with these coordinate values to equal $h$.
\item \textbf{Naturality.} Naturality in the target frame follows by applying a frame map to this join formula.
All equations are internal, proving naturality in parameters as well.
\end{enumerate}
\end{proof}

This theorem presents the space by its geometric universal property against internal frames. It proves the universal property against every internal frame $A$. Hence it will remain applicable when the target is a forcing frame, a relative frame in a slice, or the frame interpreting a further localization. Its change-of-base form is proved in Part~\ref{part:basechange}.

\begin{corollary}[Atomic realization]\label{cor:atomic-realization}
For a locale $L$ with frame $A$, models $\alpha:B\to A$ of the rounded-ideal equations correspond to locale maps $\rho:L\to\ScLoc(D_C)$. The inverse-image formula is
\[
 \rho^*(U)=\bigvee_{b:U(\kappa_Cb)}\alpha_b.
\]
\end{corollary}
\begin{proof}
\leavevmode
\begin{enumerate}[label=\textup{(\roman*)},leftmargin=*,itemsep=3pt,topsep=4pt]
\item \textbf{Atomic realization.} A locale map is, by definition, a frame homomorphism in the reverse direction.
Apply \cref{thm:rounded-classifier}.
\end{enumerate}
\end{proof}


The Scott--Alexandrov comparison expresses the condition frame and the ideal-domain frame through the same basic membership tests. A canonical site can now be reconstructed from the actual frame generators, providing a common reference for different presentations of the same forcing locale.

\section{Canonical sites and presentation comparisons}\label{sec:presentations}

An arbitrary basis presentation may contain redundant conditions. The completed forcing frame determines a canonical basis site whose generators have their actual frame values as meanings. Comparing the original and canonical sites makes the site invariance precise without identifying conditions merely because their written descriptions coincide.


For a Scott forcing datum $(C,j)$ put $H_C=\Sc(D_C)$ and
\[
 \beta_b=q_j(\delta_b)\in\Omega_j(C).
\]
These elements are join-dense by \eqref{eq:coordinate-expansion} and preservation of joins by $q_j$. The canonical condition order is
\[
 b\preceq_j c\iff \beta_c\leq\beta_b.
\]
It is the forced basis order. The original order on $B$, inclusion among the basic ideals $\kappa_Cb$, and this forced order are all retained.

\begin{theorem}[Recovery of a condition presentation]\label{thm:canonical-site}
Let $J_{C,j}$ be the semantic basis topology on $(B,\preceq_j)$. Then
\[
 \Omega_{J_{C,j}}(B)\cong\Omega_j(C),
\]
with $b$ corresponding to the stage $\beta_b$. Under this isomorphism the generic basis-membership subobjects agree.
\end{theorem}
\begin{proof}
\leavevmode
\begin{enumerate}[label=\textup{(\roman*)},leftmargin=*,itemsep=3pt,topsep=4pt]
\item \textbf{Scott continuity.} Apply \cref{cor:canonical-basis} to the join-dense family $\beta$.
Explicitly, a stable upper predicate $U$ is sent to $\bigvee_{b:U(b)}\beta_b$, and a fixed Scott open $a$ is sent to $b\mapsto\beta_b\leq a$.
The proof of \cref{thm:basis} shows these are inverse.
The membership subobject of the generic basis element $b$ is classified by its stage open.
Since both maps send that stage to $\beta_b$, they classify the same subobject in every parameter context.
\end{enumerate}
\end{proof}

\begin{theorem}[Site-to-Scott transport]\label{thm:site-scott}
For any condition site $(P,J)$, the Scott--Alexandrov isomorphism transports $j_J$ to
\[
 \bar j=\Phi_Pj_J\Phi_P^{-1}
\]
on $\Sc(\Idl P)$. It induces an isomorphism of forcing frames and of their localic sheaf topoi over $\E$. If $P$ is an effective condition preorder, it also identifies $\Up_\Sigma(P)$ with $\OO_{(P,\eta_P)}$, and identifies the individually effective stable fragments.
\end{theorem}
\begin{proof}
\leavevmode
\begin{enumerate}[label=\textup{(\roman*)},leftmargin=*,itemsep=3pt,topsep=4pt]
\item \textbf{Site-to-Scott transport.} Conjugation by a frame isomorphism preserves the nucleus laws.
The same isomorphism and its inverse restrict to the fixed points. A locale is determined by its frame, and the corresponding internal sheaf constructions are transported along this isomorphism.
\item \textbf{Ideal and directedness calculation.} For the effective assertion, a Scott open of $\Idl P$ is effective precisely when its principal-ideal trace is $\Sigma$-valued.
The maps $\Phi_P,\Psi_P$ of \cref{prop:scott-alex} therefore restrict to inverse maps of the effective observation objects.
\item \textbf{Observation preservation.} The squares
\[
\begin{tikzcd}
 \Up_\Sigma(P)\ar[r,"\Phi_{P,\Sigma}"]\ar[d,hook]&
 \OO_{(P,\eta_P)}\ar[d,hook]\\
 \Up(P)\ar[r,"\Phi_P"']&\Sc(\Idl P)
\end{tikzcd}
\qquad
\begin{tikzcd}
 \Omega_J(P)\ar[r,"\Phi_P^{\mathrm{fix}}"]\ar[d,hook]&
 \Omega_{\bar j}(P,\eta_P)\ar[d,hook]\\
 \Up(P)\ar[r,"\Phi_P"']&\Sc(\Idl P)
\end{tikzcd}
\]
commute.
Pulling the first along the fixed-frame inclusion gives the stable-fragment comparison.
This argument does not require saturation to preserve every effective predicate.

Naturality is expressed by the commutative diagram
\[
\begin{tikzcd}[column sep=large,row sep=large]
\Up(P)\ar[r,"j_J"]\ar[d,"\Phi_P"']&\Up(P)\ar[d,"\Phi_P"]\\ \Sc(\Idl P)\ar[r,"\bar j"']&\Sc(\Idl P)
\end{tikzcd}
\]
\end{enumerate}
\end{proof}

The last sentence is useful: presentation invariance of an individually effective fixed truth is a pullback statement. Global effective saturation is an additional property of the quotient map, treated in Part~\ref{part:effective}.

\subsection{Dense comparison of condition codes}
\begin{lemma}[Refinement lifting and fibre contraction]\label{cal:lem:codes}
For a monotone regular epimorphism $f:Q\to P$, suppose the following operations
are available locally:
\begin{enumerate}
\item from $q\in Q$ and $p'\geq f(q)$, produce $q'\geq q$ with $f(q')=p'$;
\item from $f(q)=f(q')$, produce $s\geq q,q'$ with $f(s)=f(q)$.
\end{enumerate}
Then inverse image induces a canonical frame isomorphism
\[
 f^{-1}:\Fix(\Dens_P)\xrightarrow{\ \cong\ }\Fix(\Dens_Q).
\]
Thus their dense sheaf topoi are equivalent. The generic denotation is
preserved.
\end{lemma}
\begin{proof}
\leavevmode
\begin{enumerate}[label=\textup{(\roman*)},leftmargin=*,itemsep=3pt,topsep=4pt]
\item \textbf{Refinement lifting and fibre contraction.} Refinement lifting proves
$f^{-1}\Dens_PV=\Dens_Qf^{-1}V$.
In one direction, refine $f(q')$ into $V$ and lift
that refinement from $q'$.
In the other, lift a proposed refinement of $f(q)$,
then use density in $Q$ and apply $f$.
\item \textbf{Existence and uniqueness.} If $U=\Dens_QU$, $q\in U$, and $f(q)=f(q')$, take any $r\geq q'$.
Lift $f(r)$
from $q$ to $t\geq q$.
Contract the equal fibre of $r,t$ to a common
strengthening $s$.
Then $s\in U$ and $s\geq r$.
Hence $q'\in \Dens_QU=U$.
Thus $U$ is constant on the fibres.
Set
\[
 V(p)\Longleftrightarrow\exists q\;(f(q)=p\ \wedge\ U(q)).
\]
Lifting makes $V$ upper; fibre constancy and epimorphy give $f^{-1}V=U$.
The commuting equation and epimorphy imply $\Dens_PV=V$.
\item \textbf{Coherence equations.} These constructions are
inverse, preserve finite meets, and preserve fixed-frame joins because closure
commutes with inverse image.
Their action on principal truth values identifies
the generic denotations.
\end{enumerate}
\end{proof}

\subsection{A complete small comparison: the two Cohen presentations}
Work in $\Eff$ for the effective claims, and internally in any base for the geometric construction. Let $P_{\mathrm{str}}=2^{<\N}$, ordered by prefix extension. Generate the topology by $\sigma\triangleleft\{\sigma0,\sigma1\}$. Define the Cantor locale $\mathbb C$ by cylinder generators $[\sigma]$, with $[\varnothing]=1$, binary splitting, and
\[
 [\sigma]\wedge[\tau]=
 \begin{cases}
 [\tau],&\sigma\preceq\tau,\\
 [\sigma],&\tau\preceq\sigma,\\
 0,&\sigma\perp\tau.
 \end{cases}
\]
These decidable cases exhaust the possibilities for two finite strings. The frame presented by these relations is precisely the site frame: \eqref{eq:upper-expansion} supplies the free upper-predicate frame, its meets give the displayed intersection equation, and its quotient imposes binary splitting.

Let $P_{\mathrm{pf}}$ be coded finite partial maps $\N\rightharpoonup2$, ordered by graph extension. Its generators are covered by assigning either value at each coordinate outside the domain. For $p\in P_{\mathrm{pf}}$, define the clopen
\[
 [p]=\bigwedge_{n\in\dom(p)}[x(n)=p(n)].
\]
It is the finite join of the string cylinders of any length exceeding every element of $\dom(p)$ that agree with $p$. Finite binary splitting proves independence of that length.

\begin{proposition}[Cohen frame comparison]\label{prop:cohen-comparison}
The string and finite-partial-map structural sites both present $\mathbb C$. Their cylinder interpretations identify their generic total binary oracles.
\end{proposition}
\begin{proof}
\leavevmode
\begin{enumerate}[label=\textup{(\roman*)},leftmargin=*,itemsep=3pt,topsep=4pt]
\item \textbf{Cohen frame comparison.} The finite-partial-map cylinders cover top, reverse the extension order, and have intersection $[p\cup q]$ when compatible and zero otherwise.
They generate the Cantor frame because every string is a finite partial map.
Their coordinate-decision covers hold by the binary partition at that coordinate.
\item \textbf{Partiality and definedness.} For completeness, reduce a finite partial map $p$ to its finite cover by all strings on an initial segment containing $\dom(p)$.
Any relation between its cylinders is thereby checked on string cylinders.
\item \textbf{Converse construction.} Conversely each string split is a coordinate-decision cover of its finite graph.
The two assignments on generators extend to frame maps in both directions, and their composites fix all string and finite-map cylinders.
Since these generate, the composites are identities. A coordinate observation is the join of precisely the cylinders assigning that coordinate its indicated value, so it is preserved.
\end{enumerate}
\end{proof}

\begin{proposition}[The information domains differ]\label{prop:cohen-domains}
Internally in $\Eff$, $\Idl(P_{\mathrm{str}})$ and $\Idl(P_{\mathrm{pf}})$ are not order-isomorphic.
\end{proposition}
\begin{proof}
\leavevmode
\begin{enumerate}[label=\textup{(\roman*)},leftmargin=*,itemsep=3pt,topsep=4pt]
\item \textbf{The information domains differ.} Every principal ideal of a string has its compact elements linearly ordered: they are its finite prefixes.
In fact all compact elements below any string ideal are linearly ordered, since an ideal contains only compatible strings and compatible strings are comparable.
\item \textbf{Ideal and directedness calculation.} For finite partial maps, the principal ideal of the assignment $\{(0,0),(1,0)\}$ has two incomparable compact elements, the principal ideals of its one-coordinate subassignments.
Principal ideals are compact: if $\downarrow p$ lies below a directed union, its member $p$ belongs to one ideal, which then contains $\downarrow p$.
\item \textbf{Converse construction.} Conversely a compact ideal is principal by its directed principal-ideal expansion.
Compactness and order below an element are preserved by an order isomorphism of dcpos.
The displayed difference rules out such an isomorphism.
\end{enumerate}
\end{proof}

Thus the forcing-locale comparison has an exact scope. It preserves the total oracle and the corresponding forcing truth. The information domain remains a separate object of the presentation. The effective comparisons in this example will be calculated by saturation in \cref{sec:quantifiers}.


The structural coordinate covers are extension-dense: any finite condition
can be extended to decide another coordinate. Hence
\cref{cal:thm:localize} constructs their dense forcing localization.
Classically, the frame isomorphism of \cref{prop:cohen-comparison} identifies
the two double-negation completions and the same Cohen generic. Equivalently,
strings are dense among finite partial maps: fill missing coordinates with
zero up to a length past the finite domain. This forcing comparison leaves the
distinct information domains of \cref{prop:cohen-domains} intact.

\subsection{The universal continuous-ideal presentation}
The site frame also has a useful universal property which will construct maps from spaces of candidate oracles.

\begin{proposition}[Frame-valued continuous ideals]\label{prop:site-classifier}
For an internal frame $A$, frame homomorphisms $\Omega_J(P)\to A$ correspond to families $\alpha:P\to A$ satisfying the ideal basis equations
\[
 \bigvee_p\alpha_p=1,\quad
 p\leq q\Rightarrow\alpha_q\leq\alpha_p,\quad
 \alpha_p\wedge\alpha_q=\bigvee_{r\geq p,q}\alpha_r,
\]
and the cover equations $\alpha_p=\bigvee_{q\in S}\alpha_q$ for $S\in J(p)$. It suffices to verify the cover equations on a generating family and its refinement pullbacks.
\end{proposition}
\begin{proof}
\leavevmode
\begin{enumerate}[label=\textup{(\roman*)},leftmargin=*,itemsep=3pt,topsep=4pt]
\item \textbf{Ideal and directedness calculation.} A frame homomorphism gives these relations by applying it to the stage basis.
Conversely the ideal equations make
$h(U)=\bigvee_{p:U(p)}\alpha_p$ a frame map on $\Up(P)$, by the meet calculation of \cref{thm:basis}.
If $j_J(U)(p)$, the restriction $U\cap\uparrow p$ covers $p$, so the cover equation gives
$\alpha_p\leq h(U)$.
Consequently $h(j_JU)\leq h(U)$; inflationarity gives equality.
\Cref{lem:quotient-factor} supplies the factorization through $\Omega_J(P)$.
It is unique since the stages generate.
\item \textbf{Pullback construction.} For the last assertion, let $J_\alpha(p)$ consist of sieves $S$ with $\alpha_p=\bigvee_{q\in S}\alpha_q$.
The maximal sieve is a cover.
Intersect its equation with $\alpha_r$ and use common refinements to prove pullback stability.
For transitivity substitute the cover equation at each member of a cover.
Thus $J_\alpha$ is a topology.
If it contains the given generators and their required refinements, it contains the generated topology.
\end{enumerate}
\end{proof}

\begin{theorem}[Points and continuous ideals]\label{thm:points-ideals}
Points of $L_{P,J}$ in any parameter context correspond to inhabited lower directed ideals $I\subseteq P$ satisfying
\[
 I(p)\ \wedge\ S\in J(p)\quad\Longrightarrow\quad
 \exists q\,[q\in S\wedge I(q)].
\]
For such an ideal, the inverse image of an upper predicate is
\[
 \alpha_I(U)=\exists p\,[I(p)\wedge U(p)].
\]
This frame map factors through $\Omega_J$ and is the inverse image of the corresponding point. The correspondence is natural under parameter substitution and geometric base change of the displayed presentation.
\end{theorem}
\begin{proof}
\leavevmode
\begin{enumerate}[label=\textup{(\roman*)},leftmargin=*,itemsep=3pt,topsep=4pt]
\item \textbf{Ideal and directedness calculation.} Specialize \cref{prop:site-classifier} to the truth-value frame $\Omega$ in the parameter slice.
Its ideal and cover equations are exactly the stated conditions, and its inverse construction is $\alpha_I$.
The generating-stage equations prove both inverse identities.
These are geometric membership and cover equations, so their pulled-back models give the same correspondence after geometric base change; \cref{thm:classifier-basechange} identifies the resulting classifier comparison.
\end{enumerate}
\end{proof}

\begin{lemma}[The topology detected by an ideal]\label{lem:ideal-topology}
For an inhabited lower directed ideal $I$, let $J_I(p)$ consist of the sieves $S$ such that
\[
 I(p)\Longrightarrow\exists q\,[I(q)\wedge q\in S].
\]
Then $J_I$ is a Grothendieck topology. Consequently it suffices to verify continuity on any generating coverage.
\end{lemma}
\begin{proof}
\leavevmode
\begin{enumerate}[label=\textup{(\roman*)},leftmargin=*,itemsep=3pt,topsep=4pt]
\item \textbf{The topology detected by an ideal.} For the $\Omega$-valued ideal family $\alpha_p=I(p)$, the topology $J_\alpha$ constructed in the proof of \cref{prop:site-classifier} has exactly this definition: lowerness gives $\bigvee_{q\in S}I(q)\leq I(p)$, so its cover equation is the displayed implication.
Hence that construction proves the topology axioms and the generating-coverage criterion.
\end{enumerate}
\end{proof}


\part{Generic oracles and full forcing semantics}\label{part:semantics}

A forcing frame determines a locale, but forcing assertions involve names and generic objects in its sheaf topos. We pass from condition coverage to that internal semantics, retaining the stage parameter in every forcing judgment. The resulting generic observation equations will supply the full truths from which effective observations are later selected.


\section{The forcing topos and its generic information}\label{sec:generic}

The forcing frame now determines a localic sheaf topos. Its generic object must be constructed inside that topos rather than treated as an externally selected filter. We first calculate the relative truth object and the generic ideal, so that subsequent statements about names are assertions in the correct context.


For a forcing locale $L$ write
\[
 \pi:\FF=\Sh_\E(L)\longrightarrow\E.
\]
This is the localic sheaf topos internal to $\E$. On a condition presentation it is equivalently the topos of internal sheaves on $(P,J)$. We use the internal site construction and the internal locale--localic-topos correspondence \cite{maclanemoerdijk,johnstone}. The following calculation specifies the relative truth object needed below.

\begin{proposition}[Relative truth classification]\label{prop:relative-truth}
For every $T\in\E$ there is a natural bijection
\[
 \E(T,\OO(L))\cong\Sub_\FF(\pi^*T).
\]
Consequently $\pi_*\Omega_\FF\cong\OO(L)$. A family $u:T\to\OO(L)$ classifies a subobject denoted $\widetilde u\incl\pi^*T$.
\end{proposition}
\begin{proof}
\leavevmode
\begin{enumerate}[label=\textup{(\roman*)},leftmargin=*,itemsep=3pt,topsep=4pt]
\item \textbf{Coherence equations.} A subterminal sheaf on the internal locale is an ideal of opens closed under the covering joins.
It is determined by their join, and the sheaf associated with an open $u$ consists of the opens below $u$.
This gives $\Sub_\FF(1)\cong\OO(L)$ as an internal statement over $\E$.
Applying the same construction in the slice $\E/T$ gives the asserted classification of subobjects of $\pi^*T$.
Restriction of a family is substitution of $u$, so the bijection is natural in $T$.
\item \textbf{Inverse comparison.} By the adjunction $\pi^*\dashv\pi_*$, its left side is naturally
$\E(T,\pi_*\Omega_\FF)$; Yoneda gives the object isomorphism.
\end{enumerate}
\end{proof}

This is a relative statement. An element of $\OO(L)$ varying over $T$ represents truth in the pulled-back context $\pi^*T$. For a general sheaf context $A$, predicates are subobjects of $A$, and are not replaced by base-indexed elements of $\OO(L)$.

Let $(P,J)$ present $L$. Its stage opens are
\[
 [p]_J=q_J(\uparrow p).
\]
They satisfy
\[
 1=\bigvee_p[p]_J,\qquad
 [p]_J\wedge[q]_J=\bigvee_{r\geq p,q}[r]_J.
\]
Let $\ell_p\incl1_\FF$ be the associated relative subterminal.

\begin{theorem}[Generic continuous condition ideal]\label{thm:generic-ideal}
The family $p\mapsto[p]_J$ classifies a subobject
\[
 \dot I_J\incl\pi^*P
\]
which is internally an inhabited lower directed ideal and meets every $J$-cover of each of its members.
\end{theorem}
\begin{proof}
\leavevmode
\begin{enumerate}[label=\textup{(\roman*)},leftmargin=*,itemsep=3pt,topsep=4pt]
\item \textbf{Localization and coverage.} Inhabitedness is the top equation for the stage basis.
If $p\leq q$, then $[q]_J\leq[p]_J$, which gives lowerness.
The meet equation gives
\[
 \dot I_J(p)\wedge\dot I_J(q)
 \ \Longrightarrow\ \exists r\geq p,q\;\dot I_J(r).
\]
For a covering sieve $S$ at $p$, the defining equation is
$[p]_J=\bigvee_{q\in S}[q]_J$.
Relative truth classification interprets this as
$\dot I_J(p)\Rightarrow\exists q\in S\,\dot I_J(q)$.
Every implication is geometric in the pulled-back relations, so it is a statement in $\FF$, with all cover parameters present.
\end{enumerate}
\end{proof}

In the Scott presentation $(C,j)$ the same construction uses $\beta_b=q_j\delta_b$. It gives the generic rounded ideal
\[
 \dot I_C\incl\pi^*B.
\]
The proof is the interpretation of the three equations of \cref{def:rounded-theory}. It defines an element of the rounded completion formed \emph{inside} $\FF$ from the pulled-back basis relations. The comparison with $\pi^*D_C$ is constructed in \cref{sec:objectchange}; no equality of these objects is needed here.

\begin{proposition}[Generic localic point]\label{prop:generic-localic}
There is a canonical point of $\pi^\#L$ in $\FF$. Under the external-relative locale correspondence it is represented by the diagonal section
$L\to L\times L$ over the first factor. Its observation on $u\in\OO(L)$ is the subterminal classified by $u$.
\end{proposition}
\begin{proof}
\leavevmode
\begin{enumerate}[label=\textup{(\roman*)},leftmargin=*,itemsep=3pt,topsep=4pt]
\item \textbf{Generic localic point.} An internal locale in $\Sh_\E(L)$ can be represented as a locale over $L$.
The pullback $\pi^\#L$ is represented by the second projection viewed in the product $L\times L$ over $L$. A point is a section, and the diagonal is such a section.
Pulling back the open from the second factor along the diagonal returns the original open on $L$.
This is exactly the relative truth statement of \cref{prop:relative-truth}.
\end{enumerate}
\end{proof}

For a forcing sublocale $L_j\incl\ScLoc(D_C)$, compose the generic localic point with this inclusion to obtain the generic information point. Its observation on $U\in\Sc(D_C)$ has truth $q_jU$. Thus its basis-membership relation is the generic rounded ideal just constructed. This identifies the ideal and localic constructions on every coordinate; \cref{eq:coordinate-expansion} then identifies all their Scott observations.

\begin{proposition}[Presentation invariance of the generic model]\label{prop:generic-invariance}
A frame isomorphism intertwining distinguished basis opens transports the associated sheaf topoi over $\E$ and their generic membership subobjects. It also transports every model whose atomic interpretations are intertwined.
\end{proposition}
\begin{proof}
\leavevmode
\begin{enumerate}[label=\textup{(\roman*)},leftmargin=*,itemsep=3pt,topsep=4pt]
\item \textbf{Presentation invariance of the generic model.} The inverse frame maps induce inverse locale maps and hence equivalent internal sheaf categories.
Under relative truth classification, the image of a membership subobject is classified by the image of its coordinate open.
The stated intertwining therefore identifies the relations.
An atomic model is the family of these relations with its geometric equations, so it is transported by the same argument.
Equality on arbitrary observations follows from the generator expansions.
\end{enumerate}
\end{proof}

\subsection{Stage representables and the relative Yoneda calculation}
The stage notation can also be derived from the site itself. Regard $P$ as the category with arrows $q\to p$ when $p\leq q$. The representable presheaf $y(p)$ is subterminal: at $q$ it is inhabited precisely when $p\leq q$. It therefore corresponds to the upper predicate $\uparrow p$.

For an arbitrary upper predicate $U$, the associated subterminal presheaf has a section at $p$ precisely when $U(p)$. Its sheafification has a section at $p$ precisely when $U$ holds on a covering sieve at $p$. To prove this assertion, a matching family of sections of a subterminal is just a cover on which the membership predicates hold; compatibility is automatic because there is at most one section. The first plus construction therefore gives the predicate $j_JU$. This predicate is already a sheaf: locality is precisely the fixed-point equation $j_J(j_JU)=j_JU$, and uniqueness of gluing is automatic. Consequently
\[
 a_J(U)=j_JU,\qquad a_Jy(p)=\ell_p.
\]

\begin{proposition}[Stage Yoneda]\label{prop:stage-yoneda}
For every sheaf $A$, the relative section object $A(p)\in\E$
has a naturally typed identification
\[
 A(p)\cong\pi_*[\ell_p,A]_{\FF},
\]
where $[-,-]_{\FF}$ is the internal exponential and $p$ may vary
in its full base parameter context. In particular,
$\E(T,A(p))\cong\FF(\pi^*T\times\ell_p,A)$ for every $T\in\E$.
These relative stages form a generating family. Hence the internal
site topos is the localic topos with frame $\Omega_J(P)$.
\end{proposition}
\begin{proof}
\leavevmode
\begin{enumerate}[label=\textup{(\roman*)},leftmargin=*,itemsep=3pt,topsep=4pt]
\item \textbf{Parameter substitution.} For each base parameter $T$, sheafification and internal Yoneda give
\[
 \begin{aligned}
 \FF(\pi^*T\times a_Jy(p),A)
 &\cong\widehat P(\pi_0^*T\times y(p),A)\\
 &\cong\E(T,A(p)).
 \end{aligned}
\]
The first bijection is the sheafification adjunction and the second
evaluates an internal natural transformation at the identity of $p$.
The inverse restricts the chosen section along all arrows into $p$.
Since $a_Jy(p)=\ell_p$, the exponential adjunction identifies
the left side with $\E(T,\pi_*[\ell_p,A]_{\FF})$.
Yoneda in
$\E$ gives the stated isomorphism of objects, including its
parameter uniformity.
\item \textbf{Evaluation and transposition.} Every presheaf, and hence every sheaf viewed as a presheaf, is covered by the evaluation family of its sections on representables.
Applying sheafification gives an epic family of stage sections.
Thus the relative subterminals $\ell_p$ generate.
The frame of all relative subterminals consists of the stable upper predicates by the preceding calculation.
The localic criterion for a topos generated by relative subterminals now identifies it with $\Sh_\E(L_{P,J})$ \cite{johnstone,maclanemoerdijk}.
\item \textbf{Universal-property argument.} This last use of the localic criterion depends on the constructed generating family, not on an assertion that a poset site has no nonconstant sheaves.
\end{enumerate}
\end{proof}

For the canonical forced basis site, representables are already sheaves. Indeed its order is $b\preceq c$ exactly when $\beta_c\leq\beta_b$. If a covering sieve at $c$ consists of stages in $\uparrow b$, then
\[
 \beta_c=\bigvee_{d\in S}\beta_d\leq\beta_b,
\]
so $c\in\uparrow b$. Thus $\uparrow b$ is stable. This proves subcanonicity and identifies its representable directly with the selected basic open. For a general site the sheafification step is essential.

\section{Names, contexts, and intuitionistic forcing}\label{sec:forcing}

A generic ideal by itself interprets atomic membership observations. Full forcing also requires predicates and quantifiers on arbitrary sheaf objects. We therefore define forcing in a named context as factorization through a subobject, and derive the logical clauses from that interpretation rather than from an external Boolean valuation.


An object $A\in\FF$ is a possible type of names. A name at stage $p$ is an arrow $a:\ell_p\to A$. In an external reading, the stage and the arrow are allowed to vary over arbitrary base parameter objects. This convention is essential for an internal site over a nontrivial topos.

\begin{definition}[Kripke--Joyal forcing]\label{def:forcing}
For a predicate $M\incl A$ and a name $a:\ell_p\to A$, write
\[
 p\Vdash M(a)
\]
when $a$ factors through $M$. For a formula $\varphi$ in context $A$, use its subobject interpretation $\llbracket\varphi\rrbracket\incl A$.
\end{definition}

\begin{lemma}[Persistence and locality]\label{lem:forcing-locality}
Forcing persists under strengthening. If a covering sieve $S$ at $p$ consists of stages forcing the restriction of $M(a)$, then $p\Vdash M(a)$.
\end{lemma}
\begin{proof}
\leavevmode
\begin{enumerate}[label=\textup{(\roman*)},leftmargin=*,itemsep=3pt,topsep=4pt]
\item \textbf{Persistence and locality.} For $q\geq p$, one has $\ell_q\incl\ell_p$, and a factorization restricts along this mono.
If all restrictions over $S$ factor, pull back $M$ along $a$.
This gives a subobject of $\ell_p$ containing all $\ell_q$ for $q\in S$.
Their join is $\ell_p$ by the cover equation, so the pulled-back subobject is all of $\ell_p$.
This is the required factorization.
\end{enumerate}
\end{proof}

\begin{lemma}[Local generation in a context]\label{lem:local-generation}
Internally over $\E$, local sections over basis stages generate every sheaf object. An inclusion or equality of subobjects can therefore be checked on such sections, allowing all parameter contexts.
\end{lemma}
\begin{proof}
\leavevmode
\begin{enumerate}[label=\textup{(\roman*)},leftmargin=*,itemsep=3pt,topsep=4pt]
\item \textbf{Local generation in a context.} For a sheaf $A$, its sections over an open $u$ are the maps from the represented subterminal $u$ to $A$.
The canonical family of all such sections has jointly epic evaluation: over any open, the section itself is one of the evaluation terms. A basis cover of $u$ restricts each section to basis opens, so basis sections still generate.
\item \textbf{Image and quotient construction.} For an internal locale the section objects and evaluation families are internal to $\E$; taking arbitrary generalized parameters gives the external formulation. A jointly epic family reflects subobject inclusion by regular image, proving the last assertion.
\end{enumerate}
\end{proof}

\begin{theorem}[Kripke--Joyal clauses]\label{thm:kj}
The Kripke--Joyal forcing relation satisfies the following clauses, with all names and stages understood in their current parameter contexts.
\begin{enumerate}
\item $p\Vdash\top$, and $p\Vdash\bot$ precisely when $[p]_J=0$.
\item $p\Vdash\varphi\wedge\psi$ precisely when both conjuncts are forced.
\item $p\Vdash\bigvee_t\varphi_t$ precisely when there is a covering family of restrictions on which a local index $t$ and a proof of $\varphi_t$ exist.
\item $p\Vdash\varphi\Rightarrow\psi$ precisely when every restriction of the context to $q\geq p$ forcing $\varphi$ also forces $\psi$.
\item $p\Vdash\exists x:B\,\theta(x)$ precisely when a cover admits local names $b$ forcing $\theta(b)$.
\item $p\Vdash\forall x:B\,\theta(x)$ precisely when every local name $b$ over every refinement $q\geq p$ forces $\theta(b)$.
\end{enumerate}
Equality of names is equality of their arrows on the stage. Negation is implication into the empty predicate. The clauses apply also to dependent types and power-object contexts.
\end{theorem}
\begin{proof}
\leavevmode
\begin{enumerate}[label=\textup{(\roman*)},leftmargin=*,itemsep=3pt,topsep=4pt]
\item \textbf{Kripke--Joyal clauses.} Top, bottom, and conjunction are immediate from the terminal, empty, and intersection subobjects.
For disjunction, the interpreting subobject is the image of the coproduct of the disjunct subobjects. A name factors through that image precisely when it locally lifts to this coproduct.
To justify the local lifting, pull the coproduct-to-image epimorphism back along the name; the result is an epimorphism over $\ell_p$.
\item \textbf{Verification of the equations.} The generating family of local sections in \cref{lem:local-generation} covers its target. A section of a coproduct is locally in a specified summand, giving both the local index and the forced disjunct.
\item \textbf{Converse construction.} Conversely those local witnesses factor through the image and glue by \cref{lem:forcing-locality}.
Existential quantification is the same regular-image argument for the projection of the interpreted relation.
\item \textbf{Functoriality and coherence.} For implication, let $R,S\incl A$ interpret $\varphi,\psi$.
The Heyting implication is characterized by $T\cap R\subseteq S$ if and only if $T\subseteq(R\Rightarrow S)$.
Pull this relation back along the name.
If $a$ factors through the implication, every restriction on which it lies in $R$ lies in $S$.
\item \textbf{Converse construction.} Conversely, suppose the stated restriction property holds.
Every basis section of $a^*R$ is a restriction satisfying $R$, hence satisfies $S$.
By local generation, $a^*R\subseteq a^*S$.
This is precisely the condition that $a$ factor through the implication.
\item \textbf{Pullback construction.} For a dependent projection $f:B\to A$, universal quantification is the right adjoint $\forall_f$ to pullback on subobjects. A name $a$ lies in $\forall_fN$ exactly when the pullback $\ell_p\times_AB$ is contained in the pulled-back $N$.
Test this inclusion on all its local sections.
Such a section is exactly a refinement of $a$ together with a local name in its dependent fibre, giving the universal clause.
The converse uses local generation again.
\item \textbf{Equality of interpreted names.} Equality is interpreted by the diagonal and hence factorization means equality of the two names.
Power objects and dependent types are objects of the topos, so the same argument applies to their names and predicates.
No restriction to the geometric fragment occurs in these clauses.
\end{enumerate}
\end{proof}

The words local index in the disjunction clause include a generalized index in the appropriate pulled-back object. Requiring one external index to work on all of $\ell_p$ would replace the image condition by a stronger global choice statement. A universal statement ranges over names in every parameter context.

\begin{corollary}[Closed and atomic stage truth]\label{cor:stage-truth}
If a closed predicate has truth $u\in\OO(L)$, then
\[
 p\Vdash u\iff[p]_J\leq u.
\]
If $u$ is represented by a $J$-stable upper predicate, this is equivalent to $u(p)$.
\end{corollary}
\begin{proof}
\leavevmode
\begin{enumerate}[label=\textup{(\roman*)},leftmargin=*,itemsep=3pt,topsep=4pt]
\item \textbf{Closed and atomic stage truth.} Factorization of one subterminal through another is inclusion of their opens.
For stable $u$, the adjunction gives
$q_J(\uparrow p)\leq u\iff\uparrow p\subseteq i_Ju$, which is $u(p)$ by upperness.
\end{enumerate}
\end{proof}

\begin{corollary}[Finite constant witnesses]\label{cor:finite-witnesses}
For a displayed finite constant type $\pi^*\Fin(n)$, a universal statement at a stage is equivalent to its $n$ constant instances there. An existential statement is equivalent to a cover on which one of those constant instances holds; the chosen index may vary across the cover.
\end{corollary}
\begin{proof}
\leavevmode
\begin{enumerate}[label=\textup{(\roman*)},leftmargin=*,itemsep=3pt,topsep=4pt]
\item \textbf{Forcing construction.} The finite coproduct decomposition $\Fin(n)=\coprod_{k<n}1$ is preserved by $\pi^*$. A local name in that object pulls the summands back to a disjoint cover of its stage, on each member of which the name is a constant $k$.
Thus forcing all the constant instances implies forcing every local instance by persistence and locality, proving the universal assertion; the reverse follows by testing constants.
For the existential assertion apply the same finite decomposition to each local witness in \cref{thm:kj}.
Conversely the displayed constant witnesses on a cover satisfy that existential clause.
\end{enumerate}
\end{proof}

\begin{example}[Local binary decisions]\label{ex:cohen-logic}
Let $U_{n,i}$ be the Cohen open saying that the generic bit at $n$ is $i$. Binary splitting gives $U_{n,0}\vee U_{n,1}=1$ and their intersection is zero. A string $\sigma$ forces $U_{n,i}$ precisely when $n<|\sigma|$ and $\sigma(n)=i$. Indeed this condition makes its cylinder lie in $U_{n,i}$. If $n\geq|\sigma|$, extend $\sigma$ to a string assigning the opposite bit at $n$; its nonempty cylinder refutes the containment. Thus a short string forces the disjunction while neither disjunct is forced there. The finite cover deciding coordinate $n$ supplies the local witnesses required by \cref{thm:kj}.
\end{example}

\subsection{Logical operations in a forcing frame}
For stable upper predicates $u,v$, conjunction is intersection, disjunction is saturation of union, and
\[
 (u\Rightarrow v)(p)
 \iff\forall q\geq p\,[u(q)\Rightarrow v(q)].
\]
This is first the implication in $\Up(P)$: intersecting an upper predicate $w$ with $u$ lies in $v$ precisely when every member of $w$ satisfies the displayed condition. Since $v$ is fixed, \cref{prop:fixed-frame} shows that the implication is already stable. In particular,
\[
 (\neg u)(p)\iff\forall q\geq p\,[u(q)\Rightarrow0_J(q)].
\]
Here $0_J$ may contain inconsistent stages. Omitting it and replacing the conclusion by ambient falsity would give the wrong clause at such stages.

For an $\E$-indexed family of closed stable truths, its existential disjunction is its join in $\Omega_J(P)$. For quantification over a general sheaf object, the context and local-name construction above is required. This distinguishes a convenient frame calculation from a general quantifier semantics.


Kripke--Joyal forcing is therefore subobject semantics in the sheaf topos, not an independently imposed list of Boolean clauses. The generic realization equations next express oracle observations in this logic, after which the positive marking will select their effective fragment.

\section{Realizations and generic oracle observations}\label{sec:realization}

The generic forcing ideal can be mapped into a separate oracle locale. This realization expresses what the generic oracle reveals while retaining the possibility that its forcing history contains more information. Pulling back oracle opens along the realization gives the exact generic observation equations used by later effective restrictions.


An oracle locale $X$ may be a Scott locale or a sublocale of one. A realization is a locale map $\rho:L\to X$. The generic realized oracle is the point of $\pi^\#X$ obtained by applying $\rho$ to the generic point of $L$.

\begin{theorem}[Generic observation equation]\label{thm:generic-observation}
For $U\in\OO(X)$, the truth of the generic realized oracle lying in $U$ is $\rho^*U$. Consequently
\[
 p\Vdash\dot G\in U\iff[p]_J\leq\rho^*U.
\]
These equations hold in families and for geometric combinations of the observations.
\end{theorem}
\begin{proof}
\leavevmode
\begin{enumerate}[label=\textup{(\roman*)},leftmargin=*,itemsep=3pt,topsep=4pt]
\item \textbf{Pullback construction.} The generic point evaluates an open of $L$ as that open by \cref{prop:generic-localic}.
The realized point first pulls $U$ back to $\rho^*U$; its evaluation is therefore the subterminal corresponding to $\rho^*U$.
Apply \cref{cor:stage-truth}.
Substitution and finite conjunction are pullbacks, and existential disjunction is image and join; inverse image of a locale preserves the corresponding operations.
This proves the family and geometric assertions.
\end{enumerate}
\end{proof}

\begin{proposition}[Direct realization]\label{prop:direct}
A Scott-continuous map $f:D_C\to D_X$ gives a realization
\[
 L_j\longrightarrow\ScLoc(D_C)\xrightarrow{\ScLoc(f)}\ScLoc(D_X)
\]
whose inverse image is $q_jf^*$. For a target sublocale $X_k\incl\ScLoc(D_X)$ this map factors through $X_k$ precisely when
\[
 q_jf^*k=q_jf^*,
\]
equivalently $f^*k(U)\leq j(f^*U)$ for every $U$.
\end{proposition}
\begin{proof}
\leavevmode
\begin{enumerate}[label=\textup{(\roman*)},leftmargin=*,itemsep=3pt,topsep=4pt]
\item \textbf{Direct realization.} Inverse image of the displayed composite is $q_jf^*$.
Factorization through the target quotient frame is \cref{lem:quotient-factor}.
Since $U\leq kU$, one direction of the required equality is automatic.
The reverse direction, after applying the fixed-frame adjunction, is exactly the displayed inequality.
The induced map on $\Fix(k)$ is the unique factorization.
\end{enumerate}
\end{proof}

For a realization given by atomic values $\alpha_b$, \cref{cor:atomic-realization} supplies the same construction without requiring an underlying map of completed objects. This is useful when conditions describe constraints on the eventual oracle.

\begin{proposition}[Common-approximation form]\label{prop:atomic-approximation}
The atomic meet equation can equivalently be written
\[
 \alpha_a\wedge\alpha_c
 =\bigvee_{\rho(a,b)\wedge\rho(c,b)}\alpha_b
\]
when the atomic rounding equations hold.
\end{proposition}
\begin{proof}
\leavevmode
\begin{enumerate}[label=\textup{(\roman*)},leftmargin=*,itemsep=3pt,topsep=4pt]
\item \textbf{Common-approximation form.} The right side is below the left because $\rho(a,b)$ and $\rho(c,b)$ imply $a,c\leq b$.
Conversely expand each factor using rounding:
\[
 \alpha_a\wedge\alpha_c
 =\bigvee_{\rho(a,u),\,\rho(c,v)}(\alpha_u\wedge\alpha_v)
 =\bigvee_{\rho(a,u),\,\rho(c,v),\,b\geq u,v}\alpha_b.
\]
Monotonicity of the approximation rows gives $\rho(a,b)$ and $\rho(c,b)$.
This proves the reverse inequality.
Starting from the common-approximation equation gives the ordinary meet equation because its indexing set is contained in $b\geq a,c$, while lowerness gives the converse inequality.
\end{enumerate}
\end{proof}

\begin{proposition}[Interpretation uniqueness]\label{prop:interpretation-unique}
Two realizations into a rounded Scott locale agreeing on every coordinate open are equal. If two constructions have the same atomic interpretation, their generic oracle observations agree on all Scott opens.
\end{proposition}
\begin{proof}
\leavevmode
\begin{enumerate}[label=\textup{(\roman*)},leftmargin=*,itemsep=3pt,topsep=4pt]
\item \textbf{Interpretation uniqueness.} Coordinate opens generate by \eqref{eq:coordinate-expansion}.
Frame maps preserve their joins, so agreement on coordinates implies agreement on all opens.
Locale maps are determined by their inverse-image maps.
The generic equation then identifies the observations.
\end{enumerate}
\end{proof}

\subsection{Object-level reconstruction by coordinates}
Suppose a locale is presented as a space of total functional relations between specified parameter objects $Q,R$. Let $E(q,r)$ be the interpreted relation in $\FF$. If
\[
 \forall q\;\exists r\;E(q,r),\qquad
 E(q,r)\wedge E(q,r')\Rightarrow r=r',
\]
then its projection to $Q$ is both regular epic and monic. It is therefore an isomorphism, and its inverse followed by the $R$-projection is the reconstructed function. These equations are geometric when functionality is expressed by equality. They can be checked on the coordinate opens of the presentation.

For the Cohen locale, take $Q=\pi^*\N$, $R=2$, and $E(n,i)$ classified by $U_{n,i}$. The partition equations of \cref{ex:cohen-logic} give a function
\[
 \dot g:\pi^*\N\longrightarrow2
\]
in $\FF$. Its transpose is an element of the binary-stream object formed in $\FF$. Every finite coordinate observation agrees with the localic generic by construction.
\Cref{seq:thm:prefixes,seq:prop:cohen} identifies this native stream
with its compatible prefixes and total prefix ideal.
\Cref{seq:prop:transport} constructs its comparison with
$\pi^*(2^\N)$, and \cref{prop:intensional-new-streams} proves the
generic stream has no local lift through that map.

For a general oracle realization, one retains the localic point and all its observations whether or not such a total functional coordinate presentation is present. Part~\ref{part:modalities} makes this reconstruction into a theorem about answer interfaces, and Part~\ref{part:basechange} compares it with pulled-back computational objects.


\part{Effective observations and oracle computation}\label{part:effective}

The generic forcing semantics classifies arbitrary truths; an implemented computation can test only a selected positive part. The task here is therefore not to redefine forcing truth, but to find its maximal effective domain and to check how such domains behave under restriction, quantification, and evaluation. These constructions will distinguish semidecision of execution from effective observation of that semidecision at a forcing stage.


\section{Stable effective truth and maximal admissibility}\label{sec:markings}


The full forcing frame interprets arbitrary predicates. An effective statement is a factorization through a specified positive observation object. It is therefore important to construct the factorization at the level of families. Separate existence of a representative for each external element would not supply the required internal map.

\begin{proposition}[Canonical positive-trace markings without searchability]
\label{prop:canonical-positive-marking}
For every approximation presentation $C=(B,r)$ there is a canonical
subobject of Scott opens defined by the pullback
\[
\begin{tikzcd}
 M_C\ar[r]\ar[d,hook,"m_C"']&\Sigma^B\ar[d,hook,"i^B"]\\
 \Sc(D_C)\ar[r,"\operatorname{tr}_C"']&\Omega^B,
\end{tikzcd}
\qquad
 \operatorname{tr}_C(U)(b)=U(\kappa_Cb).
\]
For any condition preorder $P$, whether or not its refinement relation
is positive, the same construction gives
\[
 M_P=\Up(P)\times_{\Omega^P}\Sigma^P.
\]
These are the largest subobjects admitting jointly positive basis,
respectively stage, traces; their classifiers are unique in every
parameter context. They contain top and are closed under finite
meets. For an $I$-indexed family of marked opens they preserve its
ambient join whenever the corresponding positive existential over
$I$ has been constructed, in particular whenever $I$ is
$\Sigma$-overt. This closure does not require overtness of $B$ or
$P$. If the approximation relation $\rho_C(a,b)$ has a jointly
$\Sigma$-valued classifier, every coordinate open $\delta_a$ is
marked. Overtness of $B$ is used separately for basis search and
the effective-point observation pairing.
\end{proposition}
\begin{proof}
\leavevmode
\begin{enumerate}[label=\textup{(\roman*)},leftmargin=*,itemsep=3pt,topsep=4pt]
\item \textbf{Canonical positive-trace markings without searchability.} The two displayed objects are pullbacks of the monomorphisms
$i^B$ and $i^P$, so their first projections are monic. A family
$U:T\to\Sc(D_C)$ lifts precisely when its trace
$T\times B\to\Omega$ factors through $i$; the pullback supplies
its unique positive classifier.
The argument for $P$ is identical.
The trace of top is constantly top, and traces commute with finite
meets.
\item \textbf{Scott continuity.} For $U_t$, the Scott frame joins are pointwise, hence
\[
 \operatorname{tr}_C\Bigl(\bigvee_{t:I}U_t\Bigr)(b)
 \iff\exists t:I\,\operatorname{tr}_C(U_t)(b).
\]
A supplied positive existential on the indexing object gives a
classifier for this trace, uniformly in $b$.
The same calculation
with evaluation at $p$ proves the claim for $M_P$.
\item Finally
$\operatorname{tr}_C(\delta_a)(b)$ is $a\in r(b)$, so a positive
approximation matrix marks all coordinates without using a
quantifier over $B$.
Overtness of $B$ is needed only when a later
formula explicitly quantifies over that basis.
\end{enumerate}
\end{proof}

For this part let $H$ be a frame and let $m:M\incl H$ be an observation marking. In applications $H=\Sc(D_C)$ and $M=\OO_C$, or $H=\Up(P)$ and $M=\Up_\Sigma(P)$. The order on $M$ is induced from $H$. We use the abstract notation to expose a pullback common to both applications; the positive meaning of the marking is supplied by these concrete observation objects.

\begin{definition}[Individually effective fixed truth]\label{def:stable-mark}
For a nucleus $j:H\to H$, define $S_j$ by
\[
\begin{tikzcd}
 S_j\ar[r,"e_j"]\ar[d,hook,"s_j"']&
 M\ar[d,hook,"m"]\\
 \Omega_j\ar[r,hook,"i_j"']&H.
\end{tikzcd}
\]
In the Scott case it is also written $\Stab_\Sigma(C,j)$; in the site case it is $\Stab_{J,\Sigma}(P)$.
\end{definition}
An element consists of a fixed truth and its effective ambient representative. Both projections are monic, so the representative is unique when it exists. In particular the definition uses the fixed-frame inclusion $i_j$.
Equivalently, its projection to $M$ is the equalizer
$\Eq(jm,m)\incl M$. Thus stability is constructed by finite limits,
not supplied as an extra existence principle.


The stable marking $S_j$, the realization domain $\Adm(\rho)$,
the positive quantifier domains of \cref{thm:maximal-positive-quantifiers},
and the history output domain of \cref{thm:maximal-history-observation}
are instances of this indexed construction. The sequential
restriction law of \cref{thm:lax} arises because imposing an
intermediate marking adds one pullback constraint absent from direct
restriction. The coreflection classifies \emph{representability},
not termination, and its counit need not be $\Sigma$-classified.

\begin{proposition}[Familywise universal property]\label{prop:stable-universal}
A family $u:T\to\Omega_j$ lifts to $S_j$ precisely when $i_ju$ factors through $M$. The lift is unique and is stable under substitution of $T$. If $M$ contains top and is closed under finite meets in $H$, then so is $S_j$, through its inclusion in $\Omega_j$.
\end{proposition}
\begin{proof}
\leavevmode
\begin{enumerate}[label=\textup{(\roman*)},leftmargin=*,itemsep=3pt,topsep=4pt]
\item \textbf{Familywise universal property.} The lifting and substitution assertions are \cref{thm:admissibility-coreflection} for $m$ and $i_j$.
Top is fixed by $j$ and marked by the stated operation.
For $a,b\in S_j$, the meet $i_js_j(a)\wedge i_js_j(b)$ is marked, and it equals the image of their fixed-frame meet.
The pullback therefore supplies the required meet operation.
Equality of the semilattice laws is reflected by the monic inclusion.
\end{enumerate}
\end{proof}

No arbitrary-join operation is implicit in this proposition. The join in the fixed frame is obtained by saturation, and its ambient representative must still be shown effective.

\begin{definition}[Global effective saturation]\label{def:global-effective}
The nucleus is globally effective for $M$ when there is a map $j_M:M\to M$ such that $mj_M=jm$.
\end{definition}

\begin{theorem}[Global restriction and individual fixed points]\label{thm:global-fixed}
The restricted map $j_M$, when it exists, is unique, monotone, inflationary, and idempotent. It preserves the finite meets inherited by $M$. There is a canonical isomorphism
\[
 \Fix(j_M)\cong S_j.
\]
For positive Scott or upper-predicate markings, the fixed object has the overt-indexed saturated joins whenever the corresponding overt joins of $M$ are available.
\end{theorem}
\begin{proof}
\leavevmode
\begin{enumerate}[label=\textup{(\roman*)},leftmargin=*,itemsep=3pt,topsep=4pt]
\item \textbf{Global restriction and individual fixed points.} Uniqueness and equality of maps are reflected by $m$.
The induced order on $M$ makes it an order embedding, so it also reflects inequalities.
Apply $m$ to each desired nucleus law.
For example,
\[
 mj_M^2=jmj_M=j^2m=jm=mj_M,
\]
and finite-meet preservation is reflected from $j(m(a)\wedge m(b))=jm(a)\wedge jm(b)$.
\item \textbf{Pullback construction.} A point $a$ of $\Fix(j_M)$ has $jm(a)=m(a)$, so it gives the pair $(m(a),a)$ in the defining pullback of $S_j$.
Conversely a pair in that pullback has $jm(a)=m(a)$, and monicity gives $j_Ma=a$.
These assignments are inverse and parameterized.
\item \textbf{Existence and uniqueness.} For an overt family $a_t$ of fixed marked elements, form its join $a$ in $M$ by positive existential quantification and then $j_Ma$.
Its ambient value is $j(\bigvee_t m(a_t))$, the fixed-frame join.
This proves the assertion and its uniformity.
The empty family gives the marked fixed bottom when the marking contains bottom.
\end{enumerate}
\end{proof}

A global restriction provides effective saturation for every raw effective input. The pullback $S_j$ is meaningful independently of that property. This separation allows the theory to retain useful computation in examples with non-effective full saturation.

Let an oracle locale $X$ have marking $M_X\incl\OO(X)$, and let $\rho:L_j\to X$ be a realization.

\begin{definition}[Admissible observation object]\label{def:admissibility}
Define
\[
\begin{tikzcd}
 \Adm(\rho)\ar[r]\ar[d,hook]&S_j\ar[d,hook,"s_j"]\\
 M_X\ar[r,"\rho^*m_X"']&\Omega_j.
\end{tikzcd}
\]
A family $K:T\to M_X$ is forcing-effective when it factors through $\Adm(\rho)$. The realization is observationally effective when $\Adm(\rho)\to M_X$ is an isomorphism.
\end{definition}

\begin{theorem}[Maximal admissibility]\label{thm:admissibility}
The object $\Adm(\rho)$ is the largest subobject of $M_X$ on which the realized observation has an effective stable representative. For every selected family its lift is unique. For observation markings closed under top and finite meets, admissibility has the same closure. A globally effective source nucleus and an effective pre-realization give observational effectivity.
\end{theorem}
\begin{proof}
\leavevmode
\begin{enumerate}[label=\textup{(\roman*)},leftmargin=*,itemsep=3pt,topsep=4pt]
\item \textbf{Maximal admissibility.} Apply \cref{thm:admissibility-coreflection} to $s_j$ and the family $\rho^*m_X$.
Its counit is the displayed admissibility mono, giving maximality, uniqueness and parameter coherence.
Inverse image preserves finite meets and top, and \cref{prop:stable-universal} provides the corresponding operations in $S_j$.
Pullback therefore gives them in $\Adm(\rho)$.
\item \textbf{Verification of the equations.} For the last assertion let $\rho^*=q_jh$ with $h$ restricting to $h_M:M_X\to M$.
The ambient representative of $\rho^*m_X$ is $jh\,m_X=mj_Mh_M$.
This defines its lift by \cref{prop:stable-universal}, for the whole family $M_X$.
\end{enumerate}
\end{proof}

The admissibility pullback constructs its domain and proves its
maximality, but does not make that domain the whole observation
object or make its inclusion $\Sigma$-classified. These are separate
properties of the already constructed mono. The Mathias containment
calculation in \cref{thm:mathias-unbounded} proves that total effective
admissibility fails even for a specific positive oracle observation.

\begin{corollary}[Global realization terminology]\label{cor:global-realization}
When $j_M$ exists, a realization has a lift of its full effective inverse-image family to $\Fix(j_M)$ if and only if it is observationally effective. Thus the globally effective realization formulation and the individually stable formulation select the same observations in that case.
\end{corollary}
\begin{proof}
\leavevmode
\begin{enumerate}[label=\textup{(\roman*)},leftmargin=*,itemsep=3pt,topsep=4pt]
\item \textbf{Global realization terminology.} Compose a lift with $\Fix(j_M)\cong S_j$ from \cref{thm:global-fixed}, or compose its inverse with the lift supplied by observational effectivity.
The two composites preserve the full truth and its effective representative, so are inverse by the pullback uniqueness.
\end{enumerate}
\end{proof}

\begin{proposition}[Effective postcomposition]\label{prop:admissibility-post}
If $t:X\to Y$ has an effective inverse-image restriction $t_M^*:M_Y\to M_X$, then
\[
 \Adm(t\rho)\cong M_Y\times_{M_X}\Adm(\rho).
\]
\end{proposition}
\begin{proof}
\leavevmode
\begin{enumerate}[label=\textup{(\roman*)},leftmargin=*,itemsep=3pt,topsep=4pt]
\item \textbf{Effective postcomposition.} Both sides are the pullback of $s_j:S_j\incl\Omega_j$ along
$M_Y\xrightarrow{t_M^*}M_X\xrightarrow{\rho^*m_X}\Omega_j$.
Pullback pasting gives the isomorphism, and uniqueness with the same projections gives its coherence under further postcomposition.
\end{enumerate}
\end{proof}

\subsection{Actual semidecisions}
In the site case the trace projection is $e_J:S_J\to\Up_\Sigma(P)$. Evaluation gives
\[
 \operatorname{StageObs}:P\times S_J\longrightarrow\Sigma,\qquad
 (p,a)\longmapsto e_J(a)(p).
\]
Its underlying proposition is
\[
 [p]_J\leq s_J(a)
\]
by \cref{cor:stage-truth}. Thus any selected forcing-effective family $T\to S_J$ gives a jointly uniform stage semidecision on $P\times T$. Neither overtness of the whole condition object nor an effective refinement relation is needed for this evaluation.

For a searchable Scott presentation, effective information states instead give
\[
 \operatorname{InfoObs}:\Pt D_C\times S_j\to\Sigma,\qquad
 (c,a)\longmapsto\Obs_C(c,e_j(a)).
\]
Its underlying proposition is $(i_js_j(a))(c)$. These two formulas are not identical in a general rounded presentation.

\begin{proposition}[Free information states and stages]\label{prop:stage-information}
In the free ideal presentation of an effective condition preorder, evaluation at $\downarrow p$ of a stable Scott open agrees with forcing at stage $p$. In a rounded presentation, stage truth is $q_j\delta_b\leq a$, while information-state evaluation is $(i_ja)(\kappa_Cb)$.
\end{proposition}
\begin{proof}
\leavevmode
\begin{enumerate}[label=\textup{(\roman*)},leftmargin=*,itemsep=3pt,topsep=4pt]
\item \textbf{Observation preservation.} For the free presentation, the trace at $\downarrow p$ is the upper predicate's value at $p$.
Under the fixed-frame comparison this is equivalent to $[p]_J\leq a$.
In the rounded presentation the two stated expressions follow respectively from the stage basis and the definition of the observation pairing.
They agree whenever $\delta_b$ is the principal Scott neighborhood of $\kappa_Cb$, since then $U(\kappa_Cb)\iff\delta_b\leq U$; the fixed-frame adjunction completes the comparison.
\end{enumerate}
\end{proof}

This supplies the exact comparison used in the free case and identifies the equation to calculate in a displayed rounded basis. The notation for a condition is therefore not used to suppress the difference between its stage and its basic information state.

\begin{lemma}[Self-approximation and stage coincidence]\label{lem:stage-compact-criterion}
For an approximation presentation $C=(B,r)$ and $b:B$, the following are equivalent:
\begin{equation}\label{eq:compact-coordinate-criterion}
 b\in r(b)\quad\Longleftrightarrow\quad \delta_b=\{I:D_C\mid\kappa_C(b)\subseteq I\}.
\end{equation}
Here the right-hand equality is an equality of Scott opens. Consequently, if $b\in r(b)$, then for every nucleus $j$ on $\Sc(D_C)$ and every $V\in\Omega_j(C)$,
\begin{equation}\label{eq:compact-stage-evaluation}
 q_j(\delta_b)\leq V\quad\Longleftrightarrow\quad
 (i_jV)(\kappa_Cb).
\end{equation}
For a basis-exact presentation, the first assertion follows when $\kappa_Cb\way\kappa_Cb$.
\end{lemma}
\begin{proof}
\leavevmode
\begin{enumerate}[label=\textup{(\roman*)},leftmargin=*,itemsep=3pt,topsep=4pt]
\item \textbf{Self-approximation and stage coincidence.} For a rounded ideal $I$, the implication $I(b)\Rightarrow r(b)\subseteq I$ follows from lowerness and the fact that $r(b)\subseteq\downarrow b$: if $a\in r(b)$ then $a\leq b$, hence $I(a)$.
More directly it follows from the ideal-extension reconstruction.
If $b\in r(b)$ and $r(b)\subseteq I$, then $I(b)$.
Thus the two Scott predicates agree.
\item \textbf{Converse construction.} Conversely, evaluate the asserted equality at $I=r(b)$: its right side is true, so $\delta_b(r(b))$ gives $b\in r(b)$.
For fixed $V$, the adjunction $q_j\dashv i_j$ yields
\[
 q_j(\delta_b)\leq V\ \Longleftrightarrow\ \delta_b\leq i_jV.
\]
Under the established equality $\delta_b=\uparrow\kappa_Cb$, upperness of $i_jV$ identifies this with $(i_jV)(\kappa_Cb)$.
In a basis-exact presentation, $b\in r(b)$ is the displayed instance $\kappa_Cb\way\kappa_Cb$ of the approximation relation.
\end{enumerate}
\end{proof}

\begin{remark}[Noncompact coordinates]\label{rem:noncompact-stage}
The criterion must not be silently applied to a rounded lower-cut basis, where the basic opens $\delta_b$ generally do not contain the corresponding basic cut $r(b)$. In the free ideal completion, in contrast, $r(b)=\downarrow b$ contains $b$, so stage forcing and principal-information-state evaluation coincide. The distinction is between the actual stage open $q_j\delta_b$ and evaluation at $\kappa_Cb$, not between two interpretations of the same Boolean predicate.
\end{remark}

\subsection{Detecting effectivity from the observation pairing}
The observation pairing also gives a converse criterion. This is useful when a program or a region calculation directly supplies a semidecision directly on effective information states.

For a searchable presentation $C$, the basis map $\kappa_C:B\to D_C$ lifts to a map $\operatorname{bas}_C:B\to\Pt D_C$. One can construct this lift directly. The observation profile of $\kappa_C(b)$ at $U\in\OO_C$ is the effective trace $U(\kappa_Cb)$. This is $\Sigma$-valued jointly in $(b,U)$ by the definition of $\OO_C$. Transposition and the defining pullback for $\Pt D_C$ give $\operatorname{bas}_C$, with
$j_D\operatorname{bas}_C=\kappa_C$.

\begin{theorem}[Detection by effective points]\label{thm:effective-detection}
For a stable family $u:T\to\Omega_j(C)$, the following are equivalent:
\begin{enumerate}
\item $u$ factors through $S_j$;
\item there is a map $v:\Pt D_C\times T\to\Sigma$ with
\[
 i(v(c,t))=(i_ju(t))(j_Dc).
\]
\end{enumerate}
Both assertions are about one uniform map, not separate choices for external points.
\end{theorem}
\begin{proof}
\leavevmode
\begin{enumerate}[label=\textup{(\roman*)},leftmargin=*,itemsep=3pt,topsep=4pt]
\item \textbf{Detection by effective points.} A lift $\widehat u:T\to S_j$ gives
$v(c,t)=\Obs_C(c,e_j\widehat u(t))$.
The two pullback equations for $\Pt D_C$ and $S_j$ prove the displayed equality.
\item \textbf{Converse construction.} Conversely define the trace
$\tau(t,b)=v(\operatorname{bas}_C(b),t)$.
Its underlying predicate is
$(i_ju(t))(\kappa_Cb)$, the basis trace of the full Scott open $i_ju(t)$.
The pullback defining $\OO_C$ therefore gives a unique family $w:T\to\OO_C$ with full open $i_ju$ and trace $\tau$.
The pair $(u,w)$ lies in the pullback defining $S_j$, giving the required lift.
Every construction uses maps in the full parameter context, which proves uniformity.
\end{enumerate}
\end{proof}

Thus a semidecision on all effective points detects precisely the effective stable family. Restricting that semidecision to the effective basis points is enough to reconstruct its trace. An assertion about only the external global points would not provide the map $\tau$ and is not used in this theorem.

\subsection{Effective reconstruction from atomic coordinates}
Let $X$ be a searchable oracle presentation. Suppose
$\alpha:B_X\to S_j$ has stable components satisfying the rounded-ideal equations. By \cref{cor:atomic-realization} it defines $\rho:L_j\to\ScLoc(D_X)$, and by construction its coordinate family is admissible. The following calculation describes exactly how that coordinate information extends to general effective observations.

Write $a_b=e_j\alpha(b)\in\OO_C$. For $U\in\OO_X$ form the full Scott open
\[
 W(U)=\bigvee_{b:U(\kappa_Xb)}m_C(a_b)\in\Sc(D_C).
\]
Here $m_C:\OO_C\incl\Sc(D_C)$ is the observation inclusion. Its basis trace is represented by
\begin{equation}\label{eq:atomic-positive-trace}
 \tau(U,c)=\exists^\Sigma_{b:B_X}
   \bigl(\operatorname{tr}_X(U,b)\wedge
         \operatorname{tr}_C(a_b,c)\bigr).
\end{equation}
The positive existential exists because $B_X$ is overt. Its underlying predicate is $W(U)(\kappa_Cc)$, by evaluation of the join. Hence the pair of the full open $W(U)$ and this trace constructs a map
$W_\Sigma:\OO_X\to\OO_C$.

\begin{proposition}[Atomic effective extension]\label{prop:atomic-effective}
The realized inverse image satisfies
\[
 \rho^*m_X=q_jm_CW_\Sigma.
\]
Consequently its maximal admissible domain is exactly the domain on which the family $j\,m_CW_\Sigma$ has a positive trace. If the source saturation has a global restriction $j_M$, the map $j_MW_\Sigma$ factors through $\Fix(j_M)\cong S_j$ and gives observational effectivity of $\rho$.
\end{proposition}
\begin{proof}
\leavevmode
\begin{enumerate}[label=\textup{(\roman*)},leftmargin=*,itemsep=3pt,topsep=4pt]
\item \textbf{Pullback construction.} The atomic reconstruction formula gives
\[
 \rho^*(U)=
 \bigvee_{b:U(\kappa_Xb)}^{\Omega_j}s_j\alpha(b).
\]
By the fixed-frame join formula this is $q_jW(U)$, proving the first equation.
The pullback definition of $\Adm(\rho)$ now gives the exact domain assertion.
Under global effective saturation,
$m_Cj_MW_\Sigma=j\,m_CW_\Sigma$,
and idempotence makes $j_MW_\Sigma$ fixed.
The resulting stable lift has full component $\rho^*m_X$ and is therefore the required observationally effective map.
\end{enumerate}
\end{proof}

This proof exposes the only saturation step. Overt reconstruction first builds an ordinary effective open by \eqref{eq:atomic-positive-trace}; the forcing frame then takes its saturated join. Individual coordinate effectivity proves neither that last operation effective on all inputs nor that $S_j$ is a frame.

\subsection{Target sublocales and selected families}
If the oracle itself is a sublocale $X_k\incl\ScLoc(D_X)$, its effective stable marking is
\[
 S_k^X=\Fix(k)\times_{\Sc(D_X)}\OO_X.
\]
Suppose a full direct realization factors through $X_k$ as in \cref{prop:direct}. Its inverse image on a marked target truth $U\in S_k^X$ is computed from the effective representative $e_k(U)$ before applying source saturation. For a selected family $K:T\to S_k^X$, any positive family $B:T\to\OO_C$ satisfying
\[
 m_CB=j\,F^{-1}m_Xe_kK
\]
therefore gives its effective stable lift. Conversely such a lift supplies precisely this $B$, by its trace projection. This is the pullback criterion with the target fixed fragment retained; it applies without extending $K$ to all target opens.


The maximal admissibility pullback selects the greatest domain on which a full forcing observation has the required positive representative. Since the construction concerns actual families, it can be applied to evaluator kernels and transported names without choosing independent representatives point by point.

\section{Effective transport as a partial restriction}\label{sec:lax}

A map of forcing frames always transports full truths. It need not transport every marked positive truth into the marked part of its target. Its largest valid domain is a pullback, so effective transport is generally a partial operation even when the underlying frame map is total.


Let $(H,M_H)$ and $(K,M_K)$ be marked frames. For a frame homomorphism $h:H\to K$, define
\[
 D_h=M_H\times_K M_K,\qquad
 \Res(h):M_H\pmap M_K.
\]
The maps to $K$ are $h\,m_H$ and $m_K$. Since $m_K$ is monic, $D_h\to M_H$ is monic. The other projection is the value map of the restriction.

These are partial maps in the ordinary $\Par(\E)$, with arbitrary mono domains. Their domains express the existence of a marked output. They are different from the $\Sigma$-classified domains of computations represented by $\LL$. In particular the construction does not claim to semidecide whether an arbitrary full truth has an effective representative.

\begin{theorem}[Maximal restriction and lax composition]\label{thm:lax}
Restriction is maximal among partial maps with graph lying over $h$. It satisfies
\[
 \Res(\id)=\id,\qquad
 \Res(g)\Res(h)\leq\Res(gh),
\]
where the order is restriction extension. The sequential domain is
\[
 \{a\in M_H\mid h(a)\in M_K,\ gh(a)\in M_J\},
\]
and the direct domain is
\[
 \{a\in M_H\mid gh(a)\in M_J\}.
\]
Equality holds precisely when the direct domain maps into $M_K$ under $h$. The comparison inclusions satisfy the unit and associativity coherence equations.
\end{theorem}
\begin{proof}
\leavevmode
\begin{enumerate}[label=\textup{(\roman*)},leftmargin=*,itemsep=3pt,topsep=4pt]
\item \textbf{Maximal restriction and lax composition.} Apply \cref{thm:admissibility-coreflection} to $m_K$ and $hm_H$.
Its counit is $D_h\incl M_H$, hence every partial restriction factors uniquely through this domain.
\item \textbf{Composition.} Composition of partial maps is pullback of the second domain along the first value map.
An element of its domain is therefore an $a\in M_H$ for which $h(a)$ is marked in $K$ and $g(h(a))$ is marked in $J$.
The direct domain omits the intermediate condition.
Inclusion gives the comparison, and both value maps have the same composite into $J$, so monicity of its marking identifies them.
The equality criterion follows by comparing these two subobjects.
\item \textbf{Identity law.} For the identity, the pullback of $M_H\incl H$ along itself is $M_H$, with identity value.
For three maps, either sequence of comparison inclusions sends the domain requiring all intermediate markings into the domain requiring only the final marking.
Both are inclusions of the same subobject into the same target.
They are therefore equal.
The same uniqueness proves the unit equations.
\end{enumerate}
\end{proof}

Thus $\Res$ is a normal lax functor into the locally ordered category of partial maps. The assertion compares the actual domains of definition of successive effective representatives.

\begin{example}[Strictness under localization in $\Eff$]\label{ex:lax-strict}
Take the internal frame $H=\Omega^\N$ with positive marking $M=\Sigma^\N$. Let $B$ be the complement of the halting predicate, viewed as a member of $H$, and set
\[
 j(U)=U\vee B,\qquad k(U)=\top.
\]
Distributivity proves that $j$ preserves binary meets; both maps are inflationary and idempotent and $j\leq k$. The positive input $\bot$ is not in the domain of $\Res(q_j)$, since $j(\bot)=B$ has no positive classifier by \cref{thm:intrinsichalting}. But $q_k(\bot)=\top$ is marked. Thus the direct effective restriction of the composite quotient is defined at $\bot$, while the sequential restriction through $\Fix(j)$ is not. This proves strictness of the lax comparison using actual internal frames and markings in the fixed ground topos.
\end{example}


\begin{proposition}[Transport along a marked isomorphism]\label{prop:marked-isomorphism}
A frame isomorphism $\phi:H\to H'$ together with an isomorphism $\phi_M:M\to M'$ over $\phi$ transports nuclei, stable markings, and admissibility objects. Global restrictions are transported by conjugation, and these comparisons are coherent under composition.
\end{proposition}
\begin{proof}
\leavevmode
\begin{enumerate}[label=\textup{(\roman*)},leftmargin=*,itemsep=3pt,topsep=4pt]
\item \textbf{Pullback construction.} Put $j'=\phi j\phi^{-1}$.
The fixed points are carried isomorphically to the fixed points of $j'$.
The two marking squares commute, so the universal property of the pullbacks gives the stable-marking isomorphism.
Applying the same argument to the realization pullback gives admissibility.
If $j_M$ exists, $\phi_Mj_M\phi_M^{-1}$ satisfies the required square and is the unique restricted map.
For two isomorphisms, both comparison composites preserve all the displayed projections; uniqueness gives coherence.
\end{enumerate}
\end{proof}

\begin{proposition}[Transport and reflection of effective truth]\label{prop:effective-reflection}
Let $j\leq k$ be nuclei on $H$, and write $q_{jk}:\Omega_j\to\Omega_k$ for their quotient frame map. Retain one observation marking $M\incl H$ and the individually effective stable subobjects $s_j:S_j\incl\Omega_j$, $s_k:S_k\incl\Omega_k$. Form
\[
 \begin{tikzcd}
 T_{jk}\ar[r]\ar[d,hook]&S_k\ar[d,hook,"s_k"]\\
 S_j\ar[r,"q_{jk}s_j"']&\Omega_k.
 \end{tikzcd}
\]
The left projection classifies exactly the $j$-stable effective truths whose localized images possess $k$-stable effective representatives. A total effective transport $q_{jk,\Sigma}:S_j\to S_k$ exists exactly when $q_{jk}s_j$ factors through $s_k$, and is then unique. When this factorization exists, it reflects effective representability of \emph{all} $j$-stable truths precisely when
\[
 \begin{tikzcd}
 S_j\ar[r,"q_{jk,\Sigma}"]\ar[d,hook,"s_j"']&S_k\ar[d,hook,"s_k"]\\
 \Omega_j\ar[r,"q_{jk}"']&\Omega_k
 \end{tikzcd}
\]
is a pullback.
\end{proposition}
\begin{proof}
\leavevmode
\begin{enumerate}[label=\textup{(\roman*)},leftmargin=*,itemsep=3pt,topsep=4pt]
\item \textbf{Pullback construction.} The first statement is the defining property of the displayed pullback: a map $U:T\to S_j$ factors through $T_{jk}$ if and only if the image $q_{jk}s_jU$ factors through $s_k$.
Monicity of $s_k$ proves the stated uniqueness.
For reflection, let $a:T\to\Omega_j$ be any generalized stable truth.
The lower square being cartesian says precisely that a lift of $q_{jk}a$ to $S_k$ has a unique corresponding lift of $a$ to $S_j$.
\item \textbf{Converse construction.} Conversely, if this equivalence of factorizations holds for every $T$ and every $a$, apply it to the universal pullback of $s_k$ along $q_{jk}$; Yoneda gives the displayed pullback square.
The argument is internal and stable in all parameter contexts.
\end{enumerate}
\end{proof}

\begin{example}[Neither preservation nor reflection is automatic]\label{ex:effective-reflection-fails}
In $\Eff$ use $H=\Omega^\N$, $M=\Sigma^\N$, and let $B$ be the nonhalting predicate. The nucleus $j(U)=U\vee B$ fails effective preservation because it sends the positive input $\bot$ to the nonpositive predicate $B$. The indiscrete nucleus $k(U)=\top$ fails reflection because it sends $B$ to the marked output $\top$. Moreover $j\leq k$ and $k$ is globally effective, although $j$ is not. Thus preservation, reflection, and monotonicity of global effectiveness all fail in the fixed ground model. Nonpositivity of $B$ is exactly the nonhalting conclusion of \cref{thm:intrinsichalting}.
\end{example}


\subsection{Two kinds of partial domain}
In \cref{thm:lax}, the domain of $\Res(h)$ is
\[
 \{a\in M_H\mid h(m_Ha)\text{ belongs to }M_K\}.
\]
It is a pullback of a mono and can be used in $\Par(\E)$. Its definition does not give a $\Sigma$-classifier. By contrast, \eqref{eq:partial-composite-extent} has a $\Sigma$-classifier supplied by dominance. The two partial-map constructions have the same categorical composition operation, but different domain classes.

This also explains their different roles in a proof. A computational domain describes when execution returns. An admissibility domain describes where a full semantic transformer has an effective representative. The latter may be used to state the exact scope of a theorem even when no procedure decides membership in that scope. Once a selected family factors through it, its resulting observation map is an actual semidecision.


\section{Quantification, certificates, and effective saturation}\label{sec:quantifiers}


For an object $Z$, write $\Pow_\Sigma(Z)$ for the subobject of $\Omega^Z$ consisting of maps through $\Sigma$. Pullback of predicates always restricts to this object. An overt or compact map provides a total restricted existential or universal quantifier, respectively; the following construction supplies the exact positive domain without taking either totality property as input.

\begin{theorem}[Maximal positive quantifier domains]\label{thm:maximal-positive-quantifiers}
For any map $p:A\to B$, form the two pullbacks
\[
\begin{tikzcd}[column sep=large]
 Q_p^\exists\ar[r,"\exists_p^\Sigma"]\ar[d,hook]&\Sigma^B\ar[d,hook,"i^B"]\\
 \Sigma^A\ar[r,"\exists_p i^A"']&\Omega^B
\end{tikzcd}
\qquad
\begin{tikzcd}[column sep=large]
 Q_p^\forall\ar[r,"\forall_p^\Sigma"]\ar[d,hook]&\Sigma^B\ar[d,hook,"i^B"]\\
 \Sigma^A\ar[r,"\forall_p i^A"']&\Omega^B.
\end{tikzcd}
\]
Their left legs are the maximal subobjects of positive input predicates on which the respective quantified output is positive. The restricted operations are unique and are stable under substitution in parameter contexts. A relatively overt, respectively compact, map is exactly one for which the corresponding left leg is an isomorphism in all slices.
\end{theorem}
\begin{proof}
\leavevmode
\begin{enumerate}[label=\textup{(\roman*)},leftmargin=*,itemsep=3pt,topsep=4pt]
\item \textbf{Observation preservation.} Apply \cref{thm:admissibility-coreflection} to $i^B:\Sigma^B\incl\Omega^B$ and to the families $\exists_pi^A$ and $\forall_pi^A$.
Its counits are the displayed domain monos, so it gives maximality, uniqueness and parameter coherence.
The full Beck--Chevalley equations of \cref{lem:predicate-bc} identify reindexing with the quantifier construction in the new context.
The counit is invertible precisely when the quantified family has a total positive representative, which is the stated overt or compact property in all slices.
\end{enumerate}
\end{proof}

For a map $s:K\to P$, relative $\Sigma$-overtness means that $Q_s^\exists\to\Sigma^K$ is invertible in every slice, and relative $\Sigma$-compactness of $q:E\to K$ means the analogous assertion for $Q_q^\forall$. These are properties of the quantified maps, now expressed by explicit comparisons.


\begin{theorem}[Effective quantified normal form]\label{thm:normal-form}
For a diagram
\[
 P\xleftarrow{r}E\xrightarrow{q}K\xrightarrow{s}P,
\]
the composite $N=\exists_s\forall_qr^*$ restricts to positive predicates when $q$ carries the compact quantifier structure and $s$ the overt quantifier structure. If, for a selected family $A:T\to\Up_\Sigma(P)$,
\[
 j_J(A(t))(p)\iff
 \exists k\,[s(k)=p\ \wedge\
       \forall e\,(q(e)=k\Rightarrow A(t)(r(e)))],
\]
then its saturated family has a unique lift $T\to S_J$. If the equation holds for the entire family $\Up_\Sigma(P)$, it constructs the global restriction $j_{J,\Sigma}$.
\end{theorem}
\begin{proof}
\leavevmode
\begin{enumerate}[label=\textup{(\roman*)},leftmargin=*,itemsep=3pt,topsep=4pt]
\item \textbf{Observation preservation.} Pulling $A(t)$ back along $r$ gives a positive predicate on $E$, uniformly in $t$.
Compact universal projection gives a positive predicate on $K$, and overt existential projection gives one on $P$.
The displayed equation identifies its underlying predicate with the saturation, hence also proves that it is upper.
It therefore defines a map $B:T\to\Up_\Sigma(P)$ satisfying $mB=j_JmA$.
The pair $(q_JmA,B)$ gives the lift by \cref{prop:stable-universal}.
Taking $T=\Up_\Sigma(P)$ and $A$ the identity gives the global restriction.
Its laws follow from \cref{thm:global-fixed}.
\end{enumerate}
\end{proof}

\subsection{Certificates and finite witness assembly}
A useful way to establish a normal-form equation is to exhibit certificates. For a condition $p$, a certificate $c$ has a displayed finite set of leaves, each labelled by a strengthening of $p$. A positive predicate holds throughout the certificate when it holds at every labelled leaf. Checking this is the finite universal operation just constructed. Searching for the certificate is an overt existential operation on its code object. Thus the formula
\[
 \exists c:\operatorname{Cert}(p)\;
 \forall \ell:\operatorname{Leaf}(c)\;
 U(\operatorname{label}(c,\ell))
\]
has a positive classifier whenever its displayed certificate search is overt.

The geometric proof that this formula equals the forcing saturation has two directions. Soundness verifies that every certificate leaf family covers its root, so truth on all leaves entails the saturated truth. Completeness converts a cover proof into one of the displayed certificates. For binary splitting this conversion is the finite-level uniformization in \cref{thm:cohen-saturation}. For a finite bushy witness, finite grafting composes the certificates as in \cref{prop:bushy-positive}.

When a certificate has finitely many leaves and each leaf has a further certificate, the witnesses for those finitely many certificates can be assembled by induction on the finite leaf list. At an empty list there is one empty assembly. At a successor step, eliminate the witness for the head and then the inductively assembled witnesses for the tail. This constructs one finite tuple of certificates without an axiom of choice. Grafting that tuple gives the flattened certificate. The same finite-witness argument underlies the idempotence proof for Cohen saturation.

The semantic equality and its positive classifier are therefore separate parts of one construction. The certificate proof establishes the full formula, and the quantifier operations compile that formula into a semidecision. Their conjunction supplies the stable lift. This is the concrete derivation behind the short categorical expression $\exists_s\forall_qr^*$.


This theorem calculates a representative. The geometric normal-form equality itself has to be established from the cover presentation; the next example gives that derivation.

\begin{theorem}[Constructed positive domain of a quantified forcing normal form]\label{thm:normal-form-domain}
For the diagram of \cref{thm:normal-form}, there is a canonical maximal domain $D_{\rm NF}\incl\Up_\Sigma(P)$ on which the \emph{sequential positive execution} of $\exists_s\forall_qr^*$ (positive intermediate universal quantification followed by positive existential quantification) is defined and agrees with $j_J$. It has a unique map $D_{\rm NF}\to S_J$ lifting its saturated truths. The result identifies the positive domain constructed by the displayed normal form.
\end{theorem}
\begin{proof}
\leavevmode
\begin{enumerate}[label=\textup{(\roman*)},leftmargin=*,itemsep=3pt,topsep=4pt]
\item \textbf{Constructed positive domain of a quantified forcing.} Pull $Q_q^\forall\incl\Sigma^E$ back along
$\Up_\Sigma(P)\xrightarrow{r^*}\Sigma^E$ and call the resulting subobject $D_1$.
Its universal projection supplies $v:D_1\to\Sigma^K$ classifying $\forall_qr^*A$.
Pull $Q_s^\exists\incl\Sigma^K$ back along $v$, obtaining $D_2\incl D_1$ and a map $w:D_2\to\Sigma^P$ representing $\exists_s\forall_qr^*A$.
\item Finally let $D_{\rm NF}$ be the equalizer, in $\Omega^P$, of $i^Pw$ and $j_JmA$ on $D_2$.
Equality forces the positive predicate $w$ to be upper and stable, so it defines a map to $\Up_\Sigma(P)$ and, together with the fixed quotient truth, a unique map to $S_J$ by its pullback property.
Every family on which these same two quantifier operations are positively defined and yield the saturation factors through $D_1$, then $D_2$, then the equalizer.
This proves maximality for the displayed normal form.
Each step is a pullback or equalizer and commutes with parameter substitution.
\end{enumerate}
\end{proof}

For Cohen saturation, finite compactness makes the first restricted quantifier total, $\N$-overtness makes the second total, and the explicit finite-level saturation equation makes the final equalizer all of $\Up_\Sigma(P_{\mathrm{str}})$. Thus \cref{thm:cohen-saturation} proves a genuinely global effective restriction. The domain $D_{\rm NF}$ need not exhaust $\Adm(\rho)$ in other forcing examples, since the latter also recognizes positive representatives obtained by different expressions.

\begin{lemma}[Finite uniformization]\label{lem:uniformization}
If $P(u,n)$ is a predicate indexed by $u\in2^m$ and
$\forall u\,\exists n\,P(u,n)$, then
\[
 \exists k\;\forall u\;\exists n\leq k\,P(u,n).
\]
If $P$ is monotone in $n$, the last clause can be replaced by $P(u,k)$.
\end{lemma}
\begin{proof}
\leavevmode
\begin{enumerate}[label=\textup{(\roman*)},leftmargin=*,itemsep=3pt,topsep=4pt]
\item \textbf{Finite uniformization.} Induct on $m$.
There is one index at $m=0$, so its existential witness suffices.
Split $2^{m+1}$ into the two copies of $2^m$ according to the first bit.
The induction hypothesis gives bounds for each half; their maximum bounds both.
These are two successive existential eliminations, so no choice function over a varying family is introduced.
Monotonicity promotes each bounded witness to the common bound.
\end{enumerate}
\end{proof}

\begin{theorem}[Cohen saturation]\label{thm:cohen-saturation}
For the binary-splitting topology on strings,
\[
 j_C(U)(\sigma)\iff
 \exists m:\N\;\forall u\in2^m\;U(\sigma^\frown u).
\]
In $\Eff$ this formula restricts uniformly to $\Up_\Sigma(P_{\mathrm{str}})$.
\end{theorem}
\begin{proof}
\leavevmode
\begin{enumerate}[label=\textup{(\roman*)},leftmargin=*,itemsep=3pt,topsep=4pt]
\item \textbf{Cohen saturation.} Call the displayed operation $N$.
It is upper in $\sigma$: if a strengthening already has length at least $|\sigma|+m$, upperness of $U$ gives the result with depth zero; otherwise use the remaining depth.
It is inflationary with witness $m=0$ and monotone.
If depths $m,n$ witness $NU,NV$, increase both to $\max(m,n)$ using upperness, proving finite-meet preservation.
\item \textbf{Observation preservation.} For idempotence, a witness for $NN(U)(\sigma)$ gives a finite level $2^m$, and each of its nodes has a further depth at which $U$ holds everywhere.
\Cref{lem:uniformization} gives a common further depth $k$.
Hence depth $m+k$ witnesses $N(U)(\sigma)$.
Inflationarity gives the other direction.
Thus $N$ is a nucleus.
\item \textbf{Localization and coverage.} It validates every binary cover, witnessed by depth one.
Conversely let $k$ be any nucleus validating the binary covers.
In $\Fix(k)$, repeated splitting gives
\[
 q_k(\uparrow\sigma)
   =\bigvee_{u\in2^m}q_k(\uparrow(\sigma^\frown u))
\]
by induction on $m$.
If all these extensions belong to $U$, the right side is below $q_kU$, so $\sigma\in kU$.
Therefore $N\leq k$.
The least-cover construction identifies $N=j_C$.
\item \textbf{Naturality.} For effectivity, the inner universal quantifier has a displayed finite decidable fibre and is positive by \cref{lem:finite-compact}.
The outer existential uses the overt natural numbers of $\Eff$.
Concatenation is primitive recursive, so all substitutions are effective.
\Cref{thm:normal-form} supplies the uniform restricted nucleus.
\end{enumerate}
\end{proof}

The finite-partial-map version is obtained by testing all assignments on a finite set of undecided coordinates:
\[
 j_{\mathrm{pf}}(U)(p)\iff
 \exists F\subseteq_{\mathrm{fin}}\N\setminus\dom(p)\;
     \forall v:F\to2\;U(p\cup v).
\]
To verify idempotence, collect the finitely many witness sets for the finitely many assignments to the first $F$ and take their union. Previously decided coordinates are removed; upperness transports each witness to the common assignment. The same finite witness argument proves meet preservation by union of two witness sets. Leastness follows by finite iteration of coordinate-decision covers. Coding finite sets gives an overt outer search, and the assignment object is finite, proving global effectivity. Under \cref{prop:cohen-comparison}, these computations agree on the common effective cylinder observations.

\begin{proposition}[Positive logical closure]\label{prop:positive-logic}
A family built from positive predicates by substitution, finite conjunction, and the displayed overt existential quantifiers is positive, uniformly in all free variables. The assertion holds relative to any dominance used for partial computation.
\end{proposition}
\begin{proof}
\leavevmode
\begin{enumerate}[label=\textup{(\roman*)},leftmargin=*,itemsep=3pt,topsep=4pt]
\item \textbf{Positive logical closure.} Substitution composes the classifier with a map.
Finite conjunction uses the dominance's conjunction operation.
Overt projection supplies the classifier at the existential step.
Induction on the formation of the expression gives a single classifier of the entire family.
Its equations are reflected through the inclusion into the full truth object, so the construction is stable under parameters.
\end{enumerate}
\end{proof}


The restricted quantifiers and certificate operations specify how positive assertions are built from finite evidence. Evaluator kernels can now be tested against these exact admissibility domains; the result may be positive as a computation without being positively observable in every forcing presentation.

\section{Evaluator kernels and semidecidable oracle observations}\label{sec:kernels}

For an oracle computation, a positive observation of its returned value determines a predicate on the oracle input. We construct this evaluator kernel from the actual partial evaluator, so it records termination together with the observed output. Its factorization through the chosen marking will determine whether that predicate can be tested effectively in a forcing presentation.


Let $E,X,A,Y$ be the program, oracle, input, and output domains of an implemented effective evaluator as in \cref{thm:canonicalcoding,thm:partialadequacy}:
\[
 \operatorname{Eval}:E\times X\times A\longrightarrow Y_\bot.
\]
The same notation denotes its Scott-continuous semantic map. For effective parameters $e,a$, its oracle-coordinate section is $v_{e,a}:X\to Y_\bot$.

For $V\in\OO_Y$, define the positive lifted observation
\[
 V^+(p,y)\iff\exists u:p\,V(y(u)).
\]
Its extent is evaluated before the value. Definedness is the special case $V=\top_Y$.

\begin{lemma}[Effective lifted observations and kernels]\label{lem:kernel-effective}
The observations $V^+$ are effective uniformly in $V$, and
\[
 K_{e,a,V}=v_{e,a}^*(V^+)\in\OO_X
\]
is effective uniformly in effective program, input, and output-observation parameters.
\end{lemma}
\begin{proof}
\leavevmode
\begin{enumerate}[label=\textup{(\roman*)},leftmargin=*,itemsep=3pt,topsep=4pt]
\item \textbf{Observation preservation.} Use the specified lift presentation with basis $1+B_Y$.
The trace of $V^+$ at the bottom basis element is $\bot$, and at a defined basis element $\bar b$ it is the trace of $V$ at $b$.
This trace is rounded: at the bottom tag both sides of the rounded equation are false; at a defined tag the bottom contribution is false and the remaining equation is exactly the rounding equation for the trace of $V$.
By \cref{prop:roundedtrace} it defines a Scott open on $Y_\bot$.
\item \textbf{Partiality and definedness.} Its value at a partial element is the displayed dependent existential, since the defined basis tags reconstruct precisely its defined value.
The trace is $\Sigma$-valued uniformly in $V$, so this constructs the effective lifted observation uniformly in its parameters.
\item \textbf{Observation preservation.} Let $Z=E\times X\times A$.
Effectivity of the evaluator gives the observation
$W_V=\operatorname{Eval}^{-1}(V^+)\in\OO_Z$, uniformly in $V$.
For $b\in B_X$, use the effective product point
$(e,\operatorname{bas}_Xb,a)\in\Pt Z$ and define
\[
 \tau(e,a,V,b)=
 \Obs_Z\bigl((e,\operatorname{bas}_Xb,a),W_V\bigr).
\]
Its underlying predicate is
\[
 V^+\bigl(\operatorname{Eval}(j_Ee,\kappa_Xb,j_Aa)\bigr).
\]
This is exactly the basis trace at $b$ of the full Scott open
$v_{e,a}^{-1}(V^+)$.
\item \textbf{Pullback construction.} The effective-open pullback therefore gives a unique map
\[
 \Pt E\times\Pt A\times\OO_Y\longrightarrow\OO_X
\]
with that full open and trace $\tau$.
This is the desired jointly uniform kernel map.
The argument constructs the trace in all parameters simultaneously and does not invoke a selection of representatives for individual kernels.
\end{enumerate}
\end{proof}

The kernel is an oracle observation: it says that evaluation terminates with output in $V$. It is not yet a semidecision of its forcing at every condition. That conclusion is the admissibility calculation below.

\subsection{Full observational evaluation at a localic generic oracle}
The effective kernels above are positive tests on outputs, whereas the
underlying evaluator also acts on arbitrary full Scott opens. This latter
interpretation remains available when the generic oracle exists only
as a point of its Scott locale, not as an element of a pulled-back dcpo.

\begin{theorem}[Full localic observational evaluation]\label{thm:full-observational-evaluator}
Let $\operatorname{Eval}:E\times X\times A\to Y_\bot$ be a
semantic Scott-continuous evaluator, and let
$e:T\to\Pt E$ and $a:T\to\Pt A$ be effective program and input
families over any parameter object $T$. Define
\[
 \iota_{e,a}:T\times X\longrightarrow E\times X\times A,
 \qquad(t,x)\longmapsto(j_Ee(t),x,j_Aa(t)),
\]
and $v_{e,a}=\operatorname{Eval}\,\iota_{e,a}$. In the slice
$\E/T$, every full output Scott open $W$ has a relative oracle
inverse image
\[
 H^\Omega_{e,a,W}=v_{e,a}^{-1}(W),
\]
which is a relative Scott-open family on $X$. For every realization
$\rho:L_j\to\ScLoc(X)$, inverse image along its pullback to
$\E/T$ yields a frame homomorphism on full output opens. On the
pulled-back old Scott observations it is given by
\[
 \operatorname{Eval}^{*}_{\rho,e,a}(W)
       =\rho_T^*H^\Omega_{e,a,W}.
\]
For each positive output observation $V\in\OO_Y$, the
underlying relative open of the effective computation kernel obeys
\begin{equation}\label{eq:full-evaluator-positive-kernel}
 H^\Omega_{e,a,V^+}=m_X(K_{e,a,V}),\qquad
 \operatorname{Eval}^{*}_{\rho,e,a}(V^+)
       =\rho_T^*m_X(K_{e,a,V}),
\end{equation}
with all terms interpreted in $\E/T$. If a generic object-valued
oracle $\dot g:1\to\pi^*X$ induces $\rho$ by its Scott
observations, substituting it into the pulled-back old evaluator
has this same full observational inverse-image map. The assertions
are uniform in $T$ and do not require a chosen global program or
input parameter.
\end{theorem}
\begin{proof}
\leavevmode
\begin{enumerate}[label=\textup{(\roman*)},leftmargin=*,itemsep=3pt,topsep=4pt]
\item \textbf{Ideal and directedness calculation.} The insertion is a map over $T$ after pulling the semantic domains
back to $\E/T$.
Constants preserve inhabited directed suprema, and
products compute those suprema componentwise; thus the insertion
and $v_{e,a}$ are Scott-continuous in that slice.
Inverse image
of Scott opens along $v_{e,a}$ and then along $\rho_T$ preserves
finite meets and arbitrary joins.
Their composite is therefore a
frame homomorphism, also on any newly formed relative output opens.
\item \textbf{Scott continuity.} For a pulled-back positive lifted open, compute in the context
$(t,x)$:
\[
 V^+\bigl(\operatorname{Eval}(j_Ee(t),x,j_Aa(t))\bigr)
 \iff K_{e(t),a(t),V}(x).
\]
The jointly uniform kernel construction of
\cref{lem:kernel-effective} classifies this full Scott open by the
marked family on the right, proving
\eqref{eq:full-evaluator-positive-kernel}.
When $\dot g$ exists,
object-level evaluation satisfies $W$ exactly when $\dot g$
belongs to $H^\Omega_{e,a,W}$.
By the defining Scott-observation
equation of $\dot g$ this truth is $\rho_T^*H^\Omega_{e,a,W}$.
\item \textbf{Naturality and reindexing.} This proves equality for every displayed full output observation.
All constructions are in their stated slice, so reindexing the
program and input families preserves the equations.
\end{enumerate}
\end{proof}

\begin{theorem}[Kernel interpretation and forcing effectivity]\label{thm:kernel-forcing}
Let $\rho:L_j\to\ScLoc(X)$ be a realization. The exact truth of the localic generic computation satisfying $V$ is
\[
 \rho^*K_{e,a,V}.
\]
A selected family of these kernels is forcing-effective precisely when its map into $\OO_X$ factors through $\Adm(\rho)$. Its lift $\widehat K:T\to S_j$ gives the semidecision
\[
 \Pt D_C\times T\longrightarrow\Sigma,\qquad
 (c,t)\longmapsto\Obs_C(c,e_j\widehat K(t)).
\]
For a condition-site presentation the corresponding stage semidecision is
$(p,t)\mapsto e_J\widehat K(t)(p)$.
\end{theorem}
\begin{proof}
\leavevmode
\begin{enumerate}[label=\textup{(\roman*)},leftmargin=*,itemsep=3pt,topsep=4pt]
\item \textbf{Pullback construction.} The induced Scott-locale map of $v_{e,a}$ composed with $\rho$ has inverse image
$\rho^*v_{e,a}^*$.
Applying it to $V^+$ gives the displayed truth.
By \cref{lem:kernel-effective}, the entire kernel family lands in $\OO_X$.
The pullback defining $\Adm(\rho)$ says exactly when its realized truths lift through $S_j$, proving the criterion.
Evaluation of the resulting effective trace gives the first semidecision.
In the site presentation, \cref{cor:stage-truth} identifies trace evaluation with the forced closed truth.
\end{enumerate}
\end{proof}

\begin{corollary}[Relative halting observation]\label{cor:halting}
For $V=\top_Y$, the kernel is the definedness observation. Its selected forcing-effective lift semidecides the corresponding relative halting truth. If $\rho$ is observationally effective, this construction applies to the full effective family of kernels.
\end{corollary}
\begin{proof}
\leavevmode
\begin{enumerate}[label=\textup{(\roman*)},leftmargin=*,itemsep=3pt,topsep=4pt]
\item \textbf{Relative halting observation.} The positive lift of top has precisely the computation's extent as its underlying predicate.
Apply \cref{thm:kernel-forcing}; observational effectivity makes the admissibility projection an isomorphism.
\end{enumerate}
\end{proof}

\subsection{Naturality of computational observations}
Let $t:Z\to X$ be effective. Pulling the oracle parameter back along $t$ replaces $K_{e,a,V}$ by $t^*K_{e,a,V}$, since both formulas are the inverse image of $V^+$ under $v_{e,a}t$. If $w:Y\to Y'$ is effective, observing $V'\in\OO_{Y'}$ after applying the lifted $w$ gives the kernel for $w^*V'$. To see this, on the defined extent the value is $w(y)$, and
$V'(w(y))$ is exactly $(w^*V')(y)$. Thus the comparison includes both termination and the output observation.

For a composite partial computation, the kernel formula uses the dependent extent of the composite. The equality follows by associativity of pullback of domain--value spans, or directly by reassociating the two existential extents. This is why the domain comparison of \cref{thm:partialadequacy} is retained. Part~\ref{part:basechange} proves the analogous statement after change of ambient topos.

\subsection{Finite transcripts and the Mathias distinction}
A certified finite oracle transcript records a program, input, finite answer assignment, terminating trace, and output. Syntactic validity is a finite decidable check in the realizability examples. Its oracle observation is the cylinder of answers used by the trace. A valid transcript proves that every oracle in that cylinder yields the recorded computation and output, by induction on the finite run.

For a fixed program and input, its full positive kernel is the join of the cylinders of all valid terminating transcripts with the chosen output observation. This is an exact geometric equation: termination supplies a finite run and hence a transcript; a valid transcript supplies termination on its cylinder. If a uniform use bound is available, all relevant answer assignments form one finite decidable object, so the compact universal method applies to the forcing question. Without that uniform bound the full join may require a different saturation calculation.

The Mathias construction in \cref{sec:mathias-calculation} makes this distinction explicit. Each finite answer cylinder has an effective stable trace; a uniformly bounded family of computations has one by finite conjunction. Nevertheless the kernel that searches for the least selected integer and halts precisely when it is even has no uniformly effective forced trace. The proof constructs a uniform family of reservoirs reducing its stage truth to nonhalting. Thus the general admissibility object retains exactly the selected computational observations proved effective, while the full frame continues to interpret all kernels.


\section{Intensional enhancement: named computation kernels}\label{sec:intensional-effective}


The name-presentation fibration of \cref{sec:pres} separates a semantic effective
point from its potentially many implementation names. Let
$q_E:C_E\onto\Pt E$ and $q_A:C_A\onto\Pt A$ be its chosen regular-epimorphic
program and input presentations. Compose the kernel construction of
\cref{lem:kernel-effective} with their interpretation maps to obtain
\[
 K^{\rm name}:T=C_E\times C_A\times\OO_Y\longrightarrow\OO_X,
 \quad K^{\rm name}(c,a,V)=K_{q_E(c),q_A(a),V}.
\]
Names are retained; the kernel depends only on their denotations.

\begin{theorem}[Universal effective forcing domain of named programs]
\label{thm:named-kernel-domain}
For a realization $\rho:L_j\to X$, form
\[
 T_\rho^{\rm adm}
   =T\times_{\OO_X}\Adm(\rho)
   \cong T\times_{\Omega_j}S_j.
\]
The first projection is the largest subobject of named program,
input and output-observation tuples for which the forced computation
has a jointly effective stable representative. Its canonical lift
$k_\rho:T_\rho^{\rm adm}\to S_j$ gives, for a condition site, the
uniform semidecision
\[
 (p,z)\longmapsto e_J(k_\rho(z))(p)\in\Sigma,
 \qquad
 i(e_J(k_\rho(z))(p))\iff
 [p]_J\leq\rho^*m_XK^{\rm name}(z).
\]
For a searchable Scott presentation, the analogous map is
$(c,z)\mapsto\Obs_C(c,e_jk_\rho(z))$. Name translations commuting
with interpretations preserve these constructions by pullback.
\end{theorem}
\begin{proof}
\leavevmode
\begin{enumerate}[label=\textup{(\roman*)},leftmargin=*,itemsep=3pt,topsep=4pt]
\item \textbf{Pullback construction.} The kernel is one map on the full parameter object $T$, by
\cref{lem:kernel-effective}.
Pullback pasting in
\cref{thm:admissibility-coreflection} identifies the two displayed
domains. A family $z:Z\to T$ has an effective forcing representative
precisely when its composite $\rho^*m_XK^{\rm name}z$ lifts through
$s_j:S_j\incl\Omega_j$.
The displayed pullback classifies that lift
uniquely, giving maximality.
Its evaluation is the stage formula from
\cref{sec:markings}; the information-state formula follows from the
effective-point pairing.
\item \textbf{Observation preservation.} If $u:C'_E\to C_E$ preserves program denotation,
then $K^{\rm name}(u(c),a,V)=K^{\rm name}(c,a,V)$; the same holds for
inputs.
Pullback uniqueness identifies the translated domains and
their effective traces.
No arbitrary choice of a name for a semantic
point is used.
\end{enumerate}
\end{proof}

The extent of an implemented $\Sigma$-partial program and the mono
$T_\rho^{\rm adm}\incl T$ have different meanings. The former
classifies termination on an input; the latter classifies whether a
\emph{forcing observation about that termination} has the selected
positive representative. The bounded and unbounded Mathias
calculations in \cref{thm:mathias-bounded,thm:mathias-unbounded}
show that the latter condition cannot be suppressed for all programs.

\part{Change of base and comparison}\label{part:basechange}

Forcing replaces the ground topos by a sheaf topos. An object imported by inverse image need not coincide with the corresponding object newly constructed there, especially when power objects, function spaces, or truth classifiers are involved. This part builds the comparison maps from universal membership and domain--value spans, then uses finite prefixes to exhibit the distinction for generic sequences.


\section{Geometric presentations under change of base}\label{sec:geometricchange}


There are three changes of context in this paper. A parameter object $T$ gives the slice $\E/T$. A geometric morphism $f:\FF\to\E$ gives inverse image $f^*$ on objects. It also gives inverse image $f^\#$ on internal locales, characterized by
\[
 \Sh_\FF(f^\#L)\simeq\FF\times_\E\Sh_\E(L)
\]
over $\FF$. The last expression is a pullback of geometric morphisms. It is not the assertion that the frame object of $f^\#L$ is $f^*\OO(L)$.

Write $\tau_f:f^*\Omega_\E\to\Omega_\FF$ for the characteristic map of $f^*\true$. If $R\incl X$ has characteristic map $\chi_R$, the characteristic map of the pulled-back subobject is
\begin{equation}\label{eq:truth-comparison}
 \chi_{f^*R}=\tau_f f^*\chi_R.
\end{equation}
In formulas, $(f^*R)(x)$ always means membership in this subobject. This convention is essential when a relation is originally given as a map to $\Omega_\E$.

\begin{lemma}[Geometric transport]\label{lem:geometric-transport}
Inverse image preserves finite limits, finite unions, existential direct images, and the natural numbers object. Consequently it preserves the validity of geometric sequents, with the displayed indexing objects and relations pulled back.
\end{lemma}
\begin{proof}
\leavevmode
\begin{enumerate}[label=\textup{(\roman*)},leftmargin=*,itemsep=3pt,topsep=4pt]
\item \textbf{Kernel-pair descent.} Finite limits are part of the definition.
For a family of subobjects $R_i\incl X$, their union is the image of the copairing $\coprod_iR_i\to X$.
Inverse image preserves coproducts.
It preserves regular epimorphisms because it preserves their kernel-pair coequalizers, and it preserves monos by left exactness.
It therefore preserves images and unions.
The same image argument applied to $R\incl X\to Y$ proves preservation of $\exists_{X\to Y}R$.
Finite conjunction and equality are finite-limit constructions.
These facts prove preservation of geometric sequents by induction.
\item \textbf{Verification of the equations.} Here is an explicit argument for the natural numbers object that does not use a countable-coproduct description.
Give $A$ in $\FF$ a point $z:1\to A$ and an endomorphism $s:A\to A$.
By adjunction, $z$ gives $\bar z:1\to f_*A$, and $f_*s$ is an endomorphism of $f_*A$.
The recursion property in $\E$ gives a unique $h:\N_\E\to f_*A$ with $h0=\bar z$ and $h\,\mathrm{succ}=f_*s\,h$.
Its adjunct $\bar h:f^*\N_\E\to A$ satisfies the required recursion equations by naturality of the adjunction.
\item \textbf{Converse construction.} Conversely any solution transposes to such an $h$, proving uniqueness.
Thus $f^*\N_\E$ is a natural numbers object in $\FF$.
\end{enumerate}
\end{proof}

The ideal and roundedness laws are geometric when membership is taken as a predicate. For example, directedness has the form
\[
 I(b)\wedge I(c)\ \vdash\
 \exists d\,\bigl(I(d)\wedge b\leq d\wedge c\leq d\bigr),
\]
and roundedness has the form
\[
 I(a)\ \dashv\vdash\ \exists b\,[I(b)\wedge\rho_C(a,b)].
\]
This makes it possible to transport the theory even when the object of all its models is not preserved.


The distinction between the three forms of change of context is now fixed. Universal membership is the next step: it constructs the comparison into native completions using geometric data actually preserved by inverse image, rather than assuming that internal power-object definitions commute with change of base.

\section{Universal membership and completion objects}\label{sec:objectchange}


We first construct the comparison between pulled-back ideal objects and ideals formed in the new topos; the geometric classifier comparison will then follow. If $P$ is a preorder, let
\[
 M_P\incl\Idl_\E(P)\times P
\]
be universal ideal membership. On $f^*\Idl_\E(P)$ retain the \emph{pulled-back inclusion order}. This is a specified relation; it is not defined afresh by a universal quantifier in $\FF$.

\begin{lemma}[Ideal-completion comparison]\label{lem:ideal-beta}
There is a canonical monotone map
\[
 \beta^{\Idl}_{f,P}:f^*\Idl_\E(P)\longrightarrow\Idl_\FF(f^*P)
\]
characterized by
\[
 b\in\beta^{\Idl}_{f,P}(I)\iff(f^*M_P)(I,b).
\]
It is natural in monotone maps, preserves principal ideals, and satisfies
\begin{equation}\label{eq:kleisli-beta}
 \beta^{\Idl}_{f,Q}\,f^*(m^\#)
 =(m^f)^\#\beta^{\Idl}_{f,P},
 \qquad m^f=\beta^{\Idl}_{f,Q}f^*m
\end{equation}
for every Kleisli arrow $m:P\to\Idl_\E(Q)$.
\end{lemma}
\begin{proof}
\leavevmode
\begin{enumerate}[label=\textup{(\roman*)},leftmargin=*,itemsep=3pt,topsep=4pt]
\item \textbf{Ideal and directedness calculation.} The projection $M_P\to\Idl_\E(P)$ is epic because ideals are inhabited.
Lowerness is the inclusion of the subobject defined by $M_P(I,q)\wedge p\leq q$ in the subobject defined by $M_P(I,p)$.
Directedness is epimorphy of the projection
\[
 \{(I,p,q,d)\mid I(p)\wedge I(q)\wedge I(d)\wedge p\leq d\wedge q\leq d\}
 \longrightarrow\{(I,p,q)\mid I(p)\wedge I(q)\}.
\]
Each statement is preserved by $f^*$.
Thus $f^*M_P$ is a family of ideals, and the object of ideals in $\FF$ classifies it, giving $\beta^{\Idl}_{f,P}$.
\item \textbf{Verification of the equations.} The membership implication $I\subseteq J,\ I(p)\vdash J(p)$ is also preserved.
Consequently the constructed map is monotone for the specified pulled-back order.
This proves monotonicity without replacing that order by a new implication-valued relation.
\item \textbf{Naturality and reindexing.} For a monotone $g:P\to Q$, the defining formula
\[
 q\in\Idl(g)(I)\iff\exists p\,[I(p)\wedge q\leq g(p)]
\]
is geometric.
Pulling it back gives exactly the formula for
$\Idl(f^*g)(\beta^{\Idl}_{f,P}I)$.
Extensionality proves naturality.
\item \textbf{Unit law.} For the unit, membership in either $\beta^{\Idl}_{f,P}(f^*\eta_P(p))$ or $\eta_{f^*P}(p)$ is $q\leq p$.
\item Finally,
\[
 q\in m^\#(I)\iff\exists p\,[I(p)\wedge q\in m(p)].
\]
After inverse image, the right side becomes
$\exists p\,[(f^*M_P)(I,p)\wedge q\in m^f(p)]$.
This is the defining membership formula on the right of \eqref{eq:kleisli-beta}.
Equality of membership predicates gives equality of maps.
\end{enumerate}
\end{proof}

\begin{corollary}[Composition of relation presentations]\label{cor:kleisli-basechange}
For composable Kleisli arrows $m,n$,
\[
 (n\star m)^f=n^f\star m^f.
\]
For composable geometric morphisms $g:\GG\to\FF$ and $f:\FF\to\E$,
\[
 \beta^{\Idl}_{fg,P}
 =\beta^{\Idl}_{g,f^*P}\,g^*\beta^{\Idl}_{f,P},
\]
under the canonical inverse-image identifications.
\end{corollary}
\begin{proof}
\leavevmode
\begin{enumerate}[label=\textup{(\roman*)},leftmargin=*,itemsep=3pt,topsep=4pt]
\item \textbf{Composition of relation presentations.} For the first equation, substitute $n$ into \eqref{eq:kleisli-beta} and precompose with $f^*m$.
For the second, both sides classify $(fg)^*M_P$: pulling back the defining membership equation for $\beta^{\Idl}_{f,P}$ and then the equation for $\beta^{\Idl}_{g,f^*P}$ gives this predicate.
The uniqueness of a classifying map proves equality.
\end{enumerate}
\end{proof}

\begin{proposition}[Pulled-back approximation presentations]\label{prop:base-approximation}
Put
\[
 f^\dagger C=(f^*B,r^f),\qquad
 r^f=\beta^{\Idl}_{f,B}f^*r.
\]
Then $r^f$ is an approximation idempotent and its transpose relation is $f^*\rho_C$. The comparison restricts to
\[
 \beta_{f,C}:f^*D_C\longrightarrow D_{f^\dagger C}.
\]
If $m:B\to\Idl(B')$ represents a Scott-continuous map $F:D_C\to D_{C'}$ by
$m=r'\star m\star r$, then $m^f$ represents a Scott-continuous map $F^f$ and
\begin{equation}\label{eq:rounded-naturality}
 \beta_{f,C'}f^*F=F^f\beta_{f,C}.
\end{equation}
Identities and composites are preserved.
\end{proposition}
\begin{proof}
\leavevmode
\begin{enumerate}[label=\textup{(\roman*)},leftmargin=*,itemsep=3pt,topsep=4pt]
\item \textbf{Pulled-back approximation presentations.} Monotonicity follows from \cref{lem:ideal-beta}.
The relation
$\rho_C(a,b)\vdash a\leq b$ pulls back, so $r^f\leq\eta$.
The Kleisli calculation gives
\[
 r^f\star r^f=(r\star r)^f=r^f.
\]
By its definition, $a\in r^f(b)$ is precisely $(f^*\rho_C)(a,b)$.
\item \textbf{Kernel-pair descent.} The object $D_C$ is the equalizer of $r^\#$ and the identity on $\Idl(B)$.
If $I$ belongs to its inverse image, then $f^*(r^\#)(I)=I$.
Equation \eqref{eq:kleisli-beta} gives
\[
 (r^f)^\#\bigl(\beta^{\Idl}_{f,B}I\bigr)
 =\beta^{\Idl}_{f,B}\bigl(f^*(r^\#)I\bigr)
 =\beta^{\Idl}_{f,B}I.
\]
Thus the image is rounded, defining $\beta_{f,C}$.
\item \textbf{Unit and identity laws.} The equation for $m$ is preserved by \cref{cor:kleisli-basechange}, so $m^f$ acts on rounded ideals by
\[
 F^f(I)(c)\iff\exists b\,[I(b)\wedge c\in m^f(b)].
\]
Applying \eqref{eq:kleisli-beta} to $m$ and restricting to rounded ideals proves \eqref{eq:rounded-naturality}.
The identity is represented by $r$, and composite maps by Kleisli composite relations.
The already proved equations therefore give functoriality.
\end{enumerate}
\end{proof}

In particular the naturality square is a square of \emph{constructed} maps:
\[
\begin{tikzcd}
 f^*D_C\ar[r,"\beta_{f,C}"]\ar[d,"f^*F"']&
 D_{f^\dagger C}\ar[d,"F^f"]\\
 f^*D_{C'}\ar[r,"\beta_{f,C'}"']&
 D_{f^\dagger C'}.
\end{tikzcd}
\]
The comparison need not be invertible.

\subsection{Classifier comparison}
\begin{theorem}[Classifier base change]\label{thm:classifier-basechange}
Let $C=(B,r)$ be an approximation presentation. Pull back its basis preorder and approximation relation and form the corresponding presentation $f^\dagger C$. There is a canonical isomorphism
\[
 f^\#\ScLoc(D_C)\cong\ScLoc(D_{f^\dagger C}).
\]
For a condition site $(P,J)$, pull back its displayed cover relations, including their parameters, and generate the topology $J_f$ from those relations. Then
\[
 f^\#L_{P,J}\cong L_{f^*P,J_f}.
\]
These isomorphisms carry the respective generic membership predicates to each other and satisfy the identity and composition coherence equations.
\end{theorem}
\begin{proof}
\leavevmode
\begin{enumerate}[label=\textup{(\roman*)},leftmargin=*,itemsep=3pt,topsep=4pt]
\item \textbf{Ideal and directedness calculation.} The existence of $f^\dagger C$ as an approximation presentation is proved directly in \cref{prop:base-approximation}; that proof uses only membership relations.
By \cref{thm:rounded-classifier}, $\ScLoc(D_C)$ presents precisely the geometric theory of rounded ideals on $B$.
Its inverse-image classifier presents the pulled-back theory.
The equations of that theory are the rounded-ideal equations for $f^\dagger C$, since its order and approximation predicates are exactly the pulled-back predicates.
Apply \cref{thm:rounded-classifier} in $\FF$.
\item \textbf{Inverse comparison.} Both locales therefore classify the same models, naturally in every test locale over $\FF$.
The representing isomorphism preserves the universal model.
\item \textbf{Pullback construction.} For conditions, use the explicit frame presentation in \cref{prop:site-classifier}.
It consists of the ideal equations and the equations
\[
 [p]=\bigvee_{q\in S}[q]
\]
for the displayed covers.
Pullback of this geometric presentation gives the same equations in $\FF$.
Passing from a generating coverage to the generated topology does not alter its models: the topology of covers valid in any one ideal contains the generators and hence their closure.
Thus the pulled-back classifier is $L_{f^*P,J_f}$.
\item \textbf{Localization and coverage.} This construction pulls back the relation of membership in each parameterized covering family; it does not pull back a power object and call it the new object of all sieves.
\item \textbf{Unit and identity laws.} For composable geometric morphisms, the two resulting isomorphisms preserve exactly the same pulled-back generic predicates. A map out of either classifier is determined by those predicates.
Hence the two composites agree.
The identity comparison follows by the same uniqueness argument.
\end{enumerate}
\end{proof}

This is the useful base-independent statement. Its language is geometric presentation, not preservation of arbitrary power objects or exponentials. It is an instance of the classifier methods of \cite{vickersau}. The direct relation with synthetic domain constructions belongs to the broader setting discussed in \cite{sterlingye}; the concrete comparisons needed here are proved below.


\begin{proposition}[Failure of completion preservation]\label{prop:beta-not-epi}
Let $\pi:\Sh_{\Eff}(\mathbb C)\to\Eff$ be the Cohen projection and let $B=2^{<\N}$ have prefix order. Then $\beta^{\Idl}_{\pi,B}:\pi^*\Idl(B)\to\Idl_{\Sh_{\Eff}(\mathbb C)}(\pi^*B)$ is not epic.
\end{proposition}
\begin{proof}
\leavevmode
\begin{enumerate}[label=\textup{(\roman*)},leftmargin=*,itemsep=3pt,topsep=4pt]
\item \textbf{Failure of completion preservation.} The generic binary stream $g$ defines the native ideal $I_g$ of its finite prefixes.
We prove that it has no local lift.
In a parameter context $T$, a section of a pulled-back ground object is locally represented by ground data, and the Cohen locale has the cylinder basis in that context.
Thus a lift would, after the corresponding covers, give a ground ideal $I$ and a cylinder $[\sigma]$ on which $\beta I=I_g$.
\item \textbf{Pullback construction.} Use the two subcylinders $[\sigma0]$ and $[\sigma1]$.
Membership of the ground string $\sigma0$ in $\beta I$ is the pullback of its ground membership truth in $I$, by \cref{lem:ideal-beta}.
On the first subcylinder equality requires that truth, and on the second it requires its negation.
Both subcylinders have sections over the parameter context, given by the displayed prefix followed by zeros.
Pulling back along these sections forces both the ground truth and its negation in the same context.
That context is therefore initial.
\item \textbf{Forcing construction.} All pieces of a covering family on which a lift could exist are initial, which is impossible since $I_g$ is defined on the whole nondegenerate forcing topos.
Hence the comparison is not epic.
The argument uses only finite strings, their decidable incompatibility, and the defining membership comparison.
\end{enumerate}
\end{proof}


\subsection{Products and the defined part of a lift}
The product presentation $C\boxtimes C'$ has basis $B\times B'$ and approximation relation
\[
 \rho_{C\boxtimes C'}((a,a'),(b,b'))
 \iff\rho_C(a,b)\wedge\rho_{C'}(a',b').
\]
Its completed domain is $D_C\times D_{C'}$. To see the ideal statement directly, project an ideal $I$ on $B\times B'$ to its two lower images. If $b$ and $b'$ belong to those images, choose their witnessing pairs in $I$; directedness supplies one pair above both, so lowerness gives $(b,b')\in I$. Hence $I$ is the product of its projections. Roundedness is then coordinatewise. The argument uses existential witnesses only locally, so it is internal.

The lift presentation has basis $1+B$, with a bottom tag $*$ and defined tags $\bar b$. Its order makes $*$ least and retains the order on $B$. The approximation relation is
\[
 *\ \widetilde\rho\ y,\qquad
 \bar a\ \widetilde\rho\ \bar b\iff\rho_C(a,b),\qquad
 \neg(\bar a\ \widetilde\rho\ *).
\]
Every lifted rounded ideal contains $*$. Its definedness predicate is
\[
 \defin_C(I)\iff\exists b\,I(\bar b).
\]
On this extent the defined tags form an inhabited rounded ideal on $B$, and conversely adjoining $*$ to such an ideal gives a defined element. These constructions are inverse in the slice over the extent.

\begin{proposition}[Product and lift comparisons]\label{prop:base-lift}
There are canonical presentation isomorphisms
\[
 f^\dagger(C\boxtimes C')\cong f^\dagger C\boxtimes f^\dagger C',
 \qquad
 \lambda_{f,C}:f^\dagger(C_\bot)\cong(f^\dagger C)_\bot.
\]
Let $\chi_{f,C}=|\lambda_{f,C}|\beta_{f,C_\bot}$. Then
\begin{align}
 \chi_{f,C}f^*\eta_C&=\eta_{f^\dagger C}\beta_{f,C},\label{eq:lift-unit-base}\\
 \defin_{f^\dagger C}\chi_{f,C}
   &=\tau_f f^*\defin_C,\label{eq:lift-def-base}\\
 \chi_{f,C'}f^*(F_\bot)&=(F^f)_\bot\chi_{f,C}.\label{eq:lift-map-base}
\end{align}
Equation \eqref{eq:lift-def-base} is an equality of $\Omega_\FF$-valued predicates. On defined parts $\chi_{f,C}$ restricts to $\beta_{f,C}$.
\end{proposition}
\begin{proof}
\leavevmode
\begin{enumerate}[label=\textup{(\roman*)},leftmargin=*,itemsep=3pt,topsep=4pt]
\item \textbf{Product and lift comparisons.} Inverse image preserves the product basis and the conjunction defining its relation.
It preserves the coproduct $1+B$, the two tags, and the cases defining the lifted relation.
These give the stated presentation isomorphisms.
\item \textbf{Ideal and directedness calculation.} For \eqref{eq:lift-unit-base}, both lifted ideals contain $*$.
Membership of $\bar b$ in either side is $(f^*M_C)(I,b)$, by the defining property of $\beta$.
Thus they are equal.
For \eqref{eq:lift-def-base}, membership in the defined part is the geometric formula $\exists b\,I(\bar b)$.
Pulling it back gives exactly the formula for definedness after $\chi$.
The characteristic-map convention \eqref{eq:truth-comparison} supplies the displayed $\tau_f$.
\item \textbf{Partiality and definedness.} Restricting to this extent, membership of $b$ on both sides is again $(f^*M_C)(I,b)$, proving the assertion about defined parts.
\item \textbf{Coherence equations.} The relation presenting $F_\bot$ is obtained from the relation presenting $F$ by adjoining the bottom case.
This finite coproduct construction commutes with inverse image.
Apply \eqref{eq:rounded-naturality} to the lifted relation and transport along $\lambda_{f,C},\lambda_{f,C'}$ to obtain \eqref{eq:lift-map-base}.
\end{enumerate}
\end{proof}

\begin{corollary}[Product compatibility of completion]\label{cor:beta-product}
Under the displayed product identifications,
\[
 \beta_{f,C\boxtimes C'}=\beta_{f,C}\times\beta_{f,C'}.
\]
\end{corollary}
\begin{proof}
\leavevmode
\begin{enumerate}[label=\textup{(\roman*)},leftmargin=*,itemsep=3pt,topsep=4pt]
\item \textbf{Product compatibility of completion.} Membership of $(b,b')$ in the product of the two compared ideals is the conjunction of their membership predicates.
By the defining equation for $\beta$, this is the inverse image of $I(b)\wedge I'(b')$, which is membership in the original product ideal.
Thus the two maps classify the same relation on the product basis.
Extensionality identifies them.
\end{enumerate}
\end{proof}

\begin{definition}[Completion and lift exactness]\label{def:completion-exact}
A geometric morphism is completion-exact at $C$ when $\beta_{f,C}$ is an isomorphism, and lift-exact at $C$ when $\chi_{f,C}$ is an isomorphism. These describe the specific comparison maps just constructed.
\end{definition}
Since $\lambda_{f,C}$ is an isomorphism, lift exactness at $C$ is exactly completion exactness at $C_\bot$. An object isomorphism here identifies the displayed models. Any order or effective structure transported through it is the transported structure; a new inclusion order in the target is compared by the membership equations, not silently substituted on the source.

\subsection{Object representatives of generic information}
Let $\pi:\FF_j\to\E$ be the forcing morphism for a nucleus on $\Sc(D_C)$. The generic rounded ideal gives a native element
\[
 \dot h^\dagger:1\longrightarrow D_{\pi^\dagger C}
\]
without any object exactness assertion. An element $x:1\to\pi^*D$ of a pulled-back dcpo, when available, gives a relative Scott point by
\[
 \theta_x(U)=\langle x\in\pi^*U\rangle\in\Omega_j.
\]
This is a frame map: top and finite meets follow from left exactness, while an internally indexed union of Scott opens is an existential image and is preserved by \cref{lem:geometric-transport}. This proof uses the universal membership relation and therefore holds uniformly in $U$, not just for global opens.


\begin{theorem}[Unconditional native generic computation]
\label{thm:native-generic-computation}
Let $C=(B,r)$ be a rounded information presentation, let $j$ be a
nucleus on $\Sc(D_C)$, and let $F:D_C\to D_{C'}$ be a presented
Scott-continuous map, with rounded distributor
$m_F(b)(c)\iff c\in F(\kappa_Cb)$. In the forcing sheaf topos,
let $\dot I^\dagger:1\to D_{\pi^\dagger C}$ be the native generic
rounded ideal and put
$\dot G^\dagger=F^\pi(\dot I^\dagger)$. Then, for each target
coordinate $c:B'$, its relative truth is exactly
\begin{equation}\label{eq:native-generic-coordinate}
 \llbracket c\in\dot G^\dagger\rrbracket
     =q_j\bigl(F^{-1}\delta_c\bigr).
\end{equation}
The equation extends by the Scott coordinate expansion to all full
observations transported by the geometric classifier comparison.
Consequently the localic point of the native generic oracle is the
direct realization with frame map $q_jF^{-1}$.

For every presented partial evaluator and effective program/input
parameters transported through the product and lift comparison maps,
its native generic-relative computation satisfies
\begin{equation}\label{eq:native-generic-kernel}
 \llbracket
 V^+\bigl(\operatorname{Eval}^\pi(e^\dagger,
                \dot G^\dagger,a^\dagger)\bigr)
 \rrbracket
   =\rho^*K_{e,a,V}
\end{equation}
for each transported positive output observation $V$, with
$V=\top$ giving the entire definedness predicate. These are
unconditional native domain-and-observation equations; no inverse to
$\beta_{\pi,C}$ is needed.
\end{theorem}
\begin{proof}
\leavevmode
\begin{enumerate}[label=\textup{(\roman*)},leftmargin=*,itemsep=3pt,topsep=4pt]
\item \textbf{Unconditional native generic computation.} The rounded distributor extension and universal native membership give
\[
 c\in\dot G^\dagger
 \iff\exists b:B\,
       [\dot I^\dagger(b)\wedge m_F^\pi(b)(c)].
\]
The generic membership truth at $b$ is $q_j\delta_b$.
\item \textbf{Ideal and directedness calculation.} Relative
truth classification and preservation of existential joins therefore
turn the displayed membership into
\[
 \bigvee_{b:m_F(b)(c)}q_j\delta_b
  =q_j\left(\bigvee_{b:m_F(b)(c)}\delta_b\right).
\]
The matrix reconstruction formula, evaluated at an arbitrary
rounded ideal $I$, says
\[
  (F^{-1}\delta_c)(I)
   \iff\exists b\in I\;m_F(b)(c).
\]
By the coordinate expansion this is precisely
$F^{-1}\delta_c=\bigvee_{b:m_F(b)(c)}\delta_b$,
proving \eqref{eq:native-generic-coordinate}.
\item \textbf{Scott continuity.} All full Scott opens
are generated by the target coordinates, and the two inverse-image
maps preserve joins and finite meets.
Hence their equality on
coordinates extends to all full observations and identifies the
localic realization.
\item \textbf{Partiality and definedness.} Apply the same construction to the evaluator matrix into the lifted
output presentation.
The product and lift relations are obtained by
pulling back their geometric membership relations, so
\cref{prop:base-lift,thm:partial-transport} identify both the
native evaluator relation and its defined tag.
Taking the inverse
image of the positive lifted output open $V^+$ gives exactly the
kernel $K_{e,a,V}$ before realization.
The first part of the theorem
therefore identifies the truth of that native observation with
$\rho^*K_{e,a,V}$.
\item \textbf{Observation preservation.} For $V=\top$, the lifted observation is the
extent of the partial value, by \cref{prop:partial-observation}.
This proves the definedness equation as well as the output-test
equations, jointly in all parameters.
\end{enumerate}
\end{proof}


\begin{theorem}[Reconstruction and its observation equation]\label{thm:object-generic}
For the additional comparison with an old-object representative, when $\pi$ is completion-exact at $C$, put
$\dot h=\beta_{\pi,C}^{-1}\dot h^\dagger$. Then
\[
 \theta_{\dot h}=q_j:\Sc(D_C)\longrightarrow\Omega_j.
\]
For a presented Scott-continuous map $F:D_C\to D_{C'}$, the element
$\dot g=\pi^*F(\dot h)$ induces the direct realization with frame map $q_jF^{-1}$. Moreover
$\beta_{\pi,C'}\dot g=F^\pi(\dot h^\dagger)$.
\end{theorem}
\begin{proof}
\leavevmode
\begin{enumerate}[label=\textup{(\roman*)},leftmargin=*,itemsep=3pt,topsep=4pt]
\item \textbf{Ideal and directedness calculation.} The membership property of $\beta$ identifies the pulled-back membership of $\dot h$ with the generic ideal predicate.
For $b\in B$ its relative truth is $q_j(\delta_C(b))$.
The coordinate expansion \eqref{eq:coordinate-expansion} now gives, with $U$ an arbitrary generalized Scott open,
\[
 \theta_{\dot h}(U)
 =\bigvee_b U(\kappa_Cb)\wedge q_j(\delta_C(b))
 =q_j\left(\bigvee_b U(\kappa_Cb)\wedge\delta_C(b)\right)
 =q_j(U).
\]
The joins are joins in the relative frame; coefficients are the corresponding base truth scalars.
\item \textbf{Naturality.} For the oracle, substitution gives
$\theta_{\dot g}(U)=\theta_{\dot h}(F^{-1}U)=q_j(F^{-1}U)$.
Finally evaluate \eqref{eq:rounded-naturality} at $\dot h$ to obtain the last equation.
\end{enumerate}
\end{proof}


The constructed comparisons identify what an imported completion or truth predicate denotes inside the target topos. They need not be isomorphisms. Native sequence reconstruction supplies a particularly concrete test of this difference, because every finite prefix transports while the infinite generic sequence can be genuinely new.

\section{Native sequences from finite information}
\label{sec:native-sequences}


The native sequence is reconstructed from finite prefixes. The Cohen comparison below determines whether it has a local representative in the imported ground sequence object.

The finite-prefix domain gives a direct account of sequences in the
ground effective topos and in its forcing extensions. We first work
internally in either such topos, denoted $\mathcal T$ in this section.
Write $\N_{\mathcal T}$ for its natural numbers object and
$B_{\mathcal T}=2^{<\N_{\mathcal T}}$ for its coded finite binary
strings, with length $|\sigma|$ and prefix order. All constructions
are also made in slices, with their object parameters retained.

Define the following objects inside $\mathcal T$:
\[
 \begin{aligned}
 \operatorname{Seq}_{\mathcal T}&=2^{\N_{\mathcal T}},\\
 \operatorname{CompPref}_{\mathcal T}
   &=\{s:\N_{\mathcal T}\to B_{\mathcal T}:
       |s(n)|=n\ \wedge\ s(n)\preceq s(n+1)\text{ for all }n\},\\
 \operatorname{TotIdl}_{\mathcal T}(B_{\mathcal T})
   &=\{I\in\Idl_{\mathcal T}(B_{\mathcal T}):
       \forall n\;\exists\sigma\in I\ (|\sigma|=n)\}.
 \end{aligned}
\]
Here $\Idl$ is the ordinary inhabited directed lower-ideal completion
already constructed in the paper. The compatible-prefix object is
the internal limit of the finite-prefix diagram: a generalized
element is exactly a family of compatible finite prefixes. Its
construction uses exponentials and subobjects in $\mathcal T$,
without any extra external countable-limit principle.

\begin{theorem}[Sequences, coherent prefixes, and total ideals]
\label{seq:thm:prefixes}
There are canonical isomorphisms
\[
 \operatorname{Seq}_{\mathcal T}
 \cong\operatorname{CompPref}_{\mathcal T}
 \cong\operatorname{TotIdl}_{\mathcal T}(B_{\mathcal T}).
\]
The associated ideal of $g$ is
\begin{equation}\label{eq:sequence-prefix-ideal}
 \operatorname{Pref}_{\mathcal T}(g)
   =\{\sigma:\sigma=g|_{|\sigma|}\}.
\end{equation}
These identifications and their defining membership equations
commute with substitution of object parameters.
\end{theorem}
\begin{proof}
\leavevmode
\begin{enumerate}[label=\textup{(\roman*)},leftmargin=*,itemsep=3pt,topsep=4pt]
\item \textbf{Sequences coherent prefixes and total ideals.} Construct $g|_n$ by primitive recursion: its zero value is the empty
string and $g|_{n+1}=(g|_n)^\frown g(n)$.
This gives a compatible
family.
Conversely, from $s$ define $g(n)=s(n+1)(n)$, using the
last coordinate of the string of length $n+1$.
Induction gives
$s(n)=g|_n$: the empty strings agree, and the successor step uses
prefix compatibility and the displayed last coordinate.
These
constructions are inverse.
\item \textbf{Ideal and directedness calculation.} The strings $s(n)$ form an ideal.
It contains the empty string;
every shorter prefix of $s(n)$ is $s(k)$ for its length $k$;
and $s(\max(n,m))$ is a common extension of $s(n),s(m)$.
It contains a string of every length.
\item \textbf{Converse construction.} Conversely let $I$ be a total
ideal.
Two strings in $I$ of length $n$ have a common extension in
$I$ by directedness, and hence are equal.
The relation
\[
 R(I,n,\sigma)\quad\Longleftrightarrow\quad
       \sigma\in I\ \wedge\ |\sigma|=n
\]
is therefore total and functional.
Its projection onto $(I,n)$ is
regular epic and monic, hence an isomorphism.
Its inverse constructs
$s_I(n)$ uniformly.
The length-$n$ prefix of $s_I(n+1)$ belongs to
$I$ by lowerness, so uniqueness makes it $s_I(n)$.
\item Finally every $\sigma\in I$ equals $s_I(|\sigma|)$, proving that
the two ideal constructions are inverse and giving
\eqref{eq:sequence-prefix-ideal}.
\item \textbf{Existence and uniqueness.} Only parameterized recursion, finite-string operations, and the
inverse of a total functional relation were used.
Their defining
equations are unique in every slice.
This proves parameter
coherence and shows explicitly why no choice of infinitely many
unrelated witnesses occurs.
\end{enumerate}
\end{proof}

\subsection{The sequence comparison and its domain-theoretic equation}
Let $f:\FF\to\E$ be a geometric morphism. By
\cref{lem:geometric-transport}, $f^*\N_\E\cong\N_\FF$ and
$f^*2\cong2$. Finite binary strings may be encoded by natural
numbers with a leading binary delimiter. Their length, truncation,
and extension operations are primitive recursive. Thus the
pullback of their numerical presentation is canonically
$B_\FF$, by the uniqueness of the transported recursion equations.
No infinite function object is used for this identification.

\begin{proposition}[Canonical sequence transport]
\label{seq:prop:transport}
There is a canonical map
\[
 \beta_f^{\rm seq}:f^*(\operatorname{Seq}_\E)
                  \longrightarrow\operatorname{Seq}_\FF
\]
whose evaluation is the pulled-back ground evaluation. For
$h:\mathcal H\to\FF$, it satisfies
\begin{equation}\label{eq:sequence-composition}
 \beta_{\id}^{\rm seq}=\id,\qquad
 \beta_{f\circ h}^{\rm seq}
   =\beta_h^{\rm seq}\,h^*(\beta_f^{\rm seq}).
\end{equation}
Every finite prefix is transported correctly. In particular the
sequence and ideal-completion comparisons satisfy
\begin{equation}\label{eq:sequence-ideal-comparison}
 \beta^{\Idl}_{f,B_\E}\,f^*(\operatorname{Pref}_\E)
   =\operatorname{Pref}_\FF\,\beta_f^{\rm seq},
\end{equation}
as maps $f^*\operatorname{Seq}_\E\to\Idl_\FF(B_\FF)$.
These equations hold with object parameters.
\end{proposition}
\begin{proof}
\leavevmode
\begin{enumerate}[label=\textup{(\roman*)},leftmargin=*,itemsep=3pt,topsep=4pt]
\item \textbf{Canonical sequence transport.} Pull back $\ev:\operatorname{Seq}_\E\times\N_\E\to2$,
identify its numerical coordinate object with $\N_\FF$, and
transpose.
This constructs $\beta_f^{\rm seq}$.
In
\eqref{eq:sequence-composition} both maps have the same evaluation
at every coordinate, namely the corresponding iterated pullback
of $\ev$.
The exponential universal property proves their equality.
\item \textbf{Evaluation and transposition.} Let $p_\E(g,n)=g|_n$ be the computable prefix function in
\cref{seq:thm:prefixes}.
On the pulled-back sequence object the
functions $f^*p_\E$ and $p_\FF(\beta_f^{\rm seq}(-),-)$ have
the same empty-string value and the same successor operation,
because their coordinate evaluations agree.
Parameterized recursion
therefore gives equality of these prefix functions.
Membership in the left side of \eqref{eq:sequence-ideal-comparison}
is the pulled-back equality $\sigma=p_\E(g,|\sigma|)$, by
\cref{lem:ideal-beta} and \eqref{eq:sequence-prefix-ideal}.
\item \textbf{Pullback construction.} The prefix equality just proved identifies it with membership in
the right side.
Extensionality of the ideal object proves the
equation.
All comparisons were defined by evaluation or membership,
so their pullbacks in slices satisfy the same equations and are
equal by the same uniqueness arguments.
\end{enumerate}
\end{proof}

For a computable prefix-closed binary tree $T$, write
\[
 \operatorname{Path}_{\mathcal T}(T)
   =\{g\in\operatorname{Seq}_{\mathcal T}:
                        \forall n\;g|_n\in T\}.
\]
\begin{corollary}[Native paths and their finite-prefix classifier]
\label{seq:cor:paths}
The object $\operatorname{Path}_{\mathcal T}(T)$ is canonically
the object of total prefix ideals contained in $T$. It is also
the object of points of the locale presented by $T$-cylinders,
binary-splitting covers, and empty cylinders outside $T$.
Under base change the path locale becomes the locale of the
pulled-back finite-prefix presentation. There is a canonical map
\[
 f^*\operatorname{Path}_\E(T)
       \longrightarrow\operatorname{Path}_\FF(f^*T)
\]
induced by $\beta_f^{\rm seq}$, compatible with the total-ideal
identification.
\end{corollary}
\begin{proof}
\leavevmode
\begin{enumerate}[label=\textup{(\roman*)},leftmargin=*,itemsep=3pt,topsep=4pt]
\item \textbf{Ideal and directedness calculation.} \Cref{seq:thm:prefixes} identifies containment of the prefix ideal
in $T$ with the displayed coordinate conditions.
An ideal meeting
the splitting covers has a member of every length, by induction;
conversely a total ideal meets the cover at each of its members by
taking its unique next-length prefix.
This holds uniformly at
variable lengths.
The site classifier therefore identifies its
points with these ideals; the equations for empty cylinders say
exactly that all members lie in $T$.
Their locale base change is
\cref{thm:classifier-basechange}.
\item \textbf{Verification of the equations.} In the context $g:\operatorname{Path}_\E(T),n:\N_\E$ the
membership assertion $p_\E(g,n)\in T$ is valid.
Pulling back this
relation and using the prefix equation from
\cref{seq:prop:transport} proves
$p_\FF(\beta_f^{\rm seq}g,n)\in f^*T$.
Thus the comparison
factors through the native path object.
Equation
\eqref{eq:sequence-ideal-comparison} proves compatibility with
ideals.
This argument transports the displayed family of membership
relations and constructs its factorization; it makes no
identification of the two path objects.
\end{enumerate}
\end{proof}

\subsection{The native Cohen sequence and its newness}
\begin{proposition}[Native Cohen realization by all finite prefixes]
\label{seq:prop:cohen}
For $\pi:\FF=\Sh_\Eff(\mathbb C)\to\Eff$, the generic
coordinate relation constructs a sequence
$g:1\to\operatorname{Seq}_\FF$ such that
\begin{equation}\label{eq:native-sequence-cylinder}
 \llbracket g|_{|\sigma|}=\sigma\rrbracket=[\sigma].
\end{equation}
Its prefix ideal is the generic continuous ideal of the binary-string
site. Both equations hold with finite-string parameters.
\end{proposition}
\begin{proof}
\leavevmode
\begin{enumerate}[label=\textup{(\roman*)},leftmargin=*,itemsep=3pt,topsep=4pt]
\item \textbf{Observation preservation.} The coordinate opens $U_{n,0},U_{n,1}$ of
\cref{ex:cohen-logic} are disjoint and have join top.
Their answer
relation is therefore total and functional.
Inverting its projection
as in \cref{sec:realization} constructs $g:\N_\FF\to2$.
For a string $\sigma$ of length $n$, equality of its prefix with
$g|_n$ has truth $\bigwedge_{k<n}U_{k,\sigma(k)}$.
The
finite-string meet and splitting equations identify this with
$[\sigma]$, proving \eqref{eq:native-sequence-cylinder} by finite
induction.
\item \textbf{Ideal and directedness calculation.} The defining membership equation of
\cref{thm:generic-ideal} is the same cylinder truth.
Hence the
prefix ideal and generic ideal coincide.
All operations are
parameterized finite-string constructions.
\end{enumerate}
\end{proof}

\begin{lemma}[Local representatives of imported objects]
\label{seq:lem:local-representatives}
In the Cohen sheaf topos, a local section of $\pi^*X$ can be
represented, after a covering refinement, by a ground family
$a:T\to X$ on a cylinder $[\sigma]$ over that parameter object.
Such parameterized cylinder charts jointly generate local sections.
\end{lemma}
\begin{proof}
\leavevmode
\begin{enumerate}[label=\textup{(\roman*)},leftmargin=*,itemsep=3pt,topsep=4pt]
\item \textbf{Coherence equations.} On the internal binary-string site, $\pi^*X$ is the associated
sheaf of the constant internal presheaf with value $X$.
In the plus
construction every section is locally represented by a section of
its input presheaf; applying this twice gives local representatives
in the original constant presheaf.
At a site stage these are exactly
base families with value in $X$.
The relative representable charts
generate by \cref{prop:stage-yoneda,thm:site-scott}; for this site
their sheafifications are the string cylinders.
\item \textbf{Verification of the equations.} Refining the two
local presentations and composing their covers gives the stated
charts.
This is the associated-sheaf construction
\cite{maclanemoerdijk,johnstone}, applied internally with all base
parameters retained.
\end{enumerate}
\end{proof}

\begin{proposition}[New generic streams are not transported ground streams]
\label{prop:intensional-new-streams}
For the Cohen projection, the comparison
\[
 \beta_\pi^{\rm seq}:\pi^*(2^\N)
                  \longrightarrow(\pi^*2)^{\pi^*\N}
\]
is not epic. More precisely its pullback along the native generic
sequence $g$ has initial domain: $g$ has no local ground lift.
\end{proposition}
\begin{proof}
\leavevmode
\begin{enumerate}[label=\textup{(\roman*)},leftmargin=*,itemsep=3pt,topsep=4pt]
\item \textbf{Observation preservation.} By \cref{seq:lem:local-representatives}, any proposed local lift
can be covered by cylinder charts $[\sigma]$ over ground parameter
objects $T$, with a representative $a:T\to2^\N$.
Put
$n=|\sigma|$.
The Boolean $a(n)$ splits $T$ into its two decidable
fibres $T_b$.
Over $T_b$, refine the cylinder to
$[\sigma(1-b)]$.
Equation \eqref{eq:native-sequence-cylinder}
makes $g(n)=1-b$ there, whereas the proposed equality with the
imported sequence makes $g(n)=b$.
This cylinder is therefore empty.
\item \textbf{Verification of the equations.} But the stream with prefix $\sigma(1-b)$ and then all zeros gives
a section of that cylinder locale over $T_b$.
Pulling back the
contradiction along the section makes $T_b$ initial.
Both fibres
are initial, so $T$ is initial.
\item \textbf{Failure of the local lift.} Consequently every chart covering the pullback of
$\beta_\pi^{\rm seq}$ along $g$ is initial, and so is the pullback
itself.
The forcing topos is nondegenerate: the all-zero path gives
a section of its structure morphism over the nondegenerate base
$\Eff$.
An epimorphism would have epic pullback along $g$, which
the initial map to its terminal object cannot be.
Hence the
comparison is not epic.
\end{enumerate}
\end{proof}

These results establish the native-sequence pattern inside effective
domain theory: total coherent finite information forms the native
sequence object; finite-prefix comparison agrees with the
ideal-completion comparison; and the Cohen generic has no local
representative in the imported sequence object. The path construction
of \cref{closed:thm:low-basis} takes place in the native object of
\cref{seq:cor:paths}. Its individual finite-run and oracle programs
remain the ones constructed there.

The related geometric distinction between imported objects and
objects formed in the target sheaf category is developed in
\cite[Example~4.2.10 and Lemma~4.2.12]{bhatt-scholze-proetale}
and \cite[Lecture~I, Example~1.5]{scholze-condensed}. These
references provide context for the structural pattern. The sequence,
path, and strictness assertions above are proved over $\Eff$ and
its constructed forcing topoi.

In the Artin gluing of \cref{sec:artin-gluing}, the raw sequence
object is retained by $\Gamma_j$, and restriction to the semantic
open gives its sheafification. The semantic open also has its own
native sequence object. The comparison between them is
$\beta_a^{\rm seq}$ from \cref{seq:prop:transport}; in the
binary-string presentation its target contains the generic $g$.
Thus gluing retains the computational presentation while the native
finite-prefix construction specifies the sequence formed in the
forcing extension.

\section{Semidecision truth and partial evaluation}\label{sec:partial-basechange}

Changing the base topos also changes where semidecisions and partial values live. We compare the imported Rosolini truth object and partial evaluator with their native counterparts, keeping track of the exact domain--value equations. These comparisons will be used by both oracle protocols and concrete numerical dialogue compilers.


There are three effective notions in play: a stable marking $S_j$ in the base, imported semidecision truth in the forcing topos, and the Rosolini truth object constructed natively there. Their classifiers live in different categories, and each comparison below specifies its source and target.

\begin{definition}[Imported truth]\label{def:imported-truth}
Factor the truth comparison through its image:
\[
 f^*\Sigma_\E\xrightarrow{e_f}\Theta_f
 \xrightarrow{\iota_f}\Omega_\FF,\qquad
 \iota_fe_f=\tau_f f^*i.
\]
Thus $\Theta_f$ is the base-changed semidecision object, also denoted $\Sigma_f^{\mathrm{bc}}$.
\end{definition}

\begin{proposition}[Finite operations and native comparison]\label{prop:theta}
The object $\Theta_f$ contains top and bottom and is closed under binary conjunction and disjunction. There is a canonical mono
\[
 \Theta_f\incl\Sigma_\FF
\]
over $\Omega_\FF$. It is an isomorphism if the exponential comparison
\[
 \eps_f:f^*(2^{\N_\E})\longrightarrow2^{\N_\FF}
\]
is regular epic.
\end{proposition}
\begin{proof}
\leavevmode
\begin{enumerate}[label=\textup{(\roman*)},leftmargin=*,itemsep=3pt,topsep=4pt]
\item \textbf{Observation preservation.} Send the pulled-back constants through $e_f$.
For conjunction, form
\[
 f^*\Sigma_\E\times f^*\Sigma_\E
 \cong f^*(\Sigma_\E\times\Sigma_\E)
 \xrightarrow{f^*(\wedge)}f^*\Sigma_\E
 \xrightarrow{e_f}\Theta_f.
\]
Its composite with $\iota_f$ is the conjunction of the two represented truth values, by \cref{lem:geometric-transport}.
Therefore it equalizes the kernel pair of $e_f\times e_f$.
This map is regular epic, so the displayed operation descends uniquely to $\Theta_f\times\Theta_f$.
The argument for disjunction is identical, using finite unions.
Monicity of $\iota_f$ proves uniqueness and all the operation laws.
\item \textbf{Coherence equations.} Recall that $\Sigma_\E$ is the image of the finite-hit map
$\Hit_\E:2^{\N_\E}\to\Omega_\E$.
Transpose pulled-back evaluation to obtain $\eps_f$.
The geometric definition of finite hit gives the commutative square
\[
\begin{tikzcd}
 f^*(2^{\N_\E})\ar[r,"\eps_f"]\ar[d,"f^*\Hit_\E"']&
 2^{\N_\FF}\ar[d,"\Hit_\FF"]\\
 f^*\Omega_\E\ar[r,"\tau_f"']&\Omega_\FF .
\end{tikzcd}
\]
The left finite-hit image is unchanged when precomposed with the regular epimorphism $f^*(2^{\N_\E})\to f^*\Sigma_\E$.
Its image after $\tau_f$ is therefore $\Theta_f$.
\item \textbf{Inverse comparison.} The top-right composite factors through the image $\Sigma_\FF$ of $\Hit_\FF$.
Image universality gives the required inclusion.
If $\eps_f$ is regular epic, $\Hit_\FF\eps_f$ and $\Hit_\FF$ have the same image, so that inclusion is an isomorphism.
\end{enumerate}
\end{proof}

\begin{example}[Strict imported truth over the effective topos]\label{ex:theta-strict}
For the Cohen projection $\pi:\Sh_{\Eff}(\mathbb C)\to\Eff$, the inclusion $\Theta_\pi\incl\Sigma_{\Sh_{\Eff}(\mathbb C)}$ is strict. Let $g$ be the generic stream. Its native finite-hit truth is
\[
 U\;:\!\Longleftrightarrow\;\exists n\,[g(n)=1]
      =\bigvee_{m:\N}[0^m1].
\]
This is natively semidecidable by its displayed stream. If its classifier factored through $\Theta_\pi$, the regular epimorphism $\pi^*\Sigma\onto\Theta_\pi$ would locally represent it by a pulled-back ground truth. Restrict a covering presentation to the all-zero section. After the induced parameter cover and a cylinder refinement, this would identify $U$ on $[0^m]$ with a ground truth $p$.

On $[0^m1]$ the truth $U$ holds; this subcylinder has a section over the same ground parameters, so $p$ holds there. Pullback along the all-zero section of $[0^m]$ instead makes $U$ false and hence makes $p$ false. The parameter piece is initial. Such pieces cannot cover the all-zero section. This contradiction proves that $U$ has no imported classifier. The proof treats local representatives of classifiers in arbitrary parameter contexts.
\end{example}


The finite operations of \cref{prop:theta} do not by themselves prove dependent closure, the dominance property. Nor does pulling back an overt basis prove its overtness for all newly available $\Theta_f$-predicates. The following construction supplies a canonical least scope in which the required positive operations exist, while retaining the distinction from native Rosolini semidecidability.

\begin{proposition}[Least overt dominance generated by finite hit]\label{prop:generated-dominance}
In a forcing topos $\FF$ over $\Eff$, let $S\incl\Omega$ be the image of $\Hit:2^\N\to\Omega$. There is a least dominance $\Sigma^\sharp\incl\Omega$ containing $S$ and making $\N$ $\Sigma^\sharp$-overt. If $S$ itself is an $\N$-overt dominance, then $\Sigma^\sharp=S$.
\end{proposition}
\begin{proof}
\leavevmode
\begin{enumerate}[label=\textup{(\roman*)},leftmargin=*,itemsep=3pt,topsep=4pt]
\item \textbf{Partiality and definedness.} In the power object $\Pow(\Omega)$, form the subobject $\mathcal D$ of subobjects $D\incl\Omega$ that contain $S$, contain top, are closed under dependent existential composition of $D$-classified extents, and satisfy
\[
  (\forall n:\N\;D(v(n)))\Longrightarrow D(\exists n:\N\,v(n))
  \qquad(v:\N\to\Omega).
\]
These are internal predicates on subobjects of $\Omega$.
The object $\mathcal D$ is inhabited because $D=\Omega$ qualifies.
Define
\[
 \Sigma^\sharp(p)\iff\forall D\in\mathcal D\;D(p).
\]
Clearly $S\subseteq\Sigma^\sharp$ and top belongs to $\Sigma^\sharp$.
\item \textbf{Observation preservation.} If $p\in\Sigma^\sharp$ and $q:[p]\to\Sigma^\sharp$, then for every $D\in\mathcal D$ both $p$ and all values of $q$ belong to $D$.
Its dependent closure gives $\exists_{u:[p]}q(u)\in D$, hence membership in $\Sigma^\sharp$.
The same pointwise argument for $v:\N\to\Sigma^\sharp$ proves $\exists_n v(n)\in\Sigma^\sharp$.
Monicity of the inclusion constructs unique uniform operations, so this is an $\N$-overt dominance.
Its finite meets are obtained by dependent closure, and its finite joins by $\N$-overtness applied to a family with two nonbottom entries.
\item \textbf{Universal-property argument.} It lies in every qualifying $D$ by definition and is least.
If $S$ qualifies, the defining intersection lies below $S$, while it contains $S$, proving equality.
\end{enumerate}
\end{proof}


\begin{proposition}[Generated effective scope after base change]\label{prop:generated-imported-scope}
For $f:\FF\to\E$ and a displayed family $B\to I$ of bases in $\FF$, there is a least $\N_\FF$-overt dominance $\Delta_{f,B}\incl\Omega_\FF$ containing the imported truth object $\Theta_f$ and making every fibre $B_i$ relatively $\Delta_{f,B}$-overt. In particular every imported positive approximation relation is $\Delta_{f,B}$-classified.
\end{proposition}
\begin{proof}
\leavevmode
\begin{enumerate}[label=\textup{(\roman*)},leftmargin=*,itemsep=3pt,topsep=4pt]
\item \textbf{Existence and uniqueness.} In the power object $\Pow(\Omega_\FF)$ form the subobject of $D\incl\Omega_\FF$ containing $\Theta_f$, satisfying dependent dominance closure, and closed under existential quantification over $\N_\FF$ and over every displayed fibre $B_i$, for all predicates with values in $D$.
The entire truth object qualifies.
Intersect all qualifying subobjects.
Dependent closure and the specified existential closure properties pass to this intersection by the same pointwise calculation as in \cref{prop:generated-dominance}.
\item \textbf{Tracked implementation.} This constructs the least requested $\N_\FF$-overt dominance with the relative overt quantifiers.
Finite conjunction follows from dominance and finite disjunction from $\N_\FF$-overtness using its coded finite subobjects. A pulled-back positive relation factors through $\Theta_f$ by its image definition, and thence through the constructed inclusion $\Theta_f\incl\Delta_{f,B}$.
\end{enumerate}
\end{proof}

\begin{example}[A generated scope can be the entire truth object]
\label{ex:generated-total-scope}
The generated dominance of \cref{prop:generated-imported-scope} need not coincide with native semidecidability. Take the displayed family of bases to be $\true:1\incl\Omega_\FF$. Then $\Delta_{f,B}=\Omega_\FF$.
\end{example}
\begin{proof}
\leavevmode
\begin{enumerate}[label=\textup{(\roman*)},leftmargin=*,itemsep=3pt,topsep=4pt]
\item \textbf{The generic fibre.} In the parameter context $p:\Omega_\FF$, the displayed fibre is the subterminal object $[p]$. Its constant-top predicate is classified by every dominance.
\item \textbf{Its existential image.} Relative overtness of this displayed family requires $\exists_{u:[p]}\top=p$ to belong to $\Delta$, uniformly in $p$. Thus every truth value belongs to $\Delta$ and $\Delta=\Omega_\FF$. The entire truth object satisfies all the required closure properties, proving the claimed equality. Accordingly, the construction of a least effective scope does not itself identify that scope with the native finite-hit object.
\end{enumerate}
\end{proof}

\begin{definition}[Effective scope of presentation base change]\label{def:effective-basechange}
For the displayed family of bases, an effective presentation scope is an $\N_\FF$-overt dominance $\Delta\incl\Omega_\FF$ containing $\Theta_f$, with those bases relatively $\Delta$-overt. \Cref{prop:generated-imported-scope} constructs the least such scope $\Delta_{f,B}$. The original imported scope occurs exactly when $\Theta_f$ itself has these properties. Computability relative to the native Rosolini object $\Sigma_\FF$ is asserted only for a specifically constructed comparison to that object; it does not follow from the mere existence of $\Delta_{f,B}$.
\end{definition}

\begin{proposition}[Effective relation transport]\label{prop:effective-rel-base}
In an effective presentation scope $\Delta$, if $F$ has positive presentation relation $R_F:B\times B'\to\Sigma_\E$, then $F^f$ has the $\Delta$-classified relation
\[
 R_F^{f,\Delta}
 =\bigl(\Theta_f\incl\Delta\bigr)e_ff^*R_F.
\]
When $\Delta=\Theta_f$ this is the original imported classifier $e_ff^*R_F$. Identities, composites, products, and lifts have the previously constructed comparisons, uniformly in all parameters.
\end{proposition}
\begin{proof}
\leavevmode
\begin{enumerate}[label=\textup{(\roman*)},leftmargin=*,itemsep=3pt,topsep=4pt]
\item \textbf{Effective relation transport.} The underlying predicate of $R_F^{f,\Delta}$ is $\tau_f f^*(iR_F)$, hence the pulled-back relation presenting $m^f$ in \cref{prop:base-approximation}.
Its action is
\[
 F^f(I)(c)\iff\exists b\,[I(b)\wedge R_F^{f,\Delta}(b,c)].
\]
The positive matrix calculus of \cref{thm:relations} applies in the $\Delta$-scope because the intermediate bases are relatively $\Delta$-overt and $\Delta$ is a dominance.
\item \textbf{Composition.} For composition, the full matrix is $\exists c\,[R_F(b,c)\wedge R_G(c,a)]$; inverse image preserves this existential formula, and the $\Delta$-classified version is uniquely determined by its underlying predicate.
Identities use the approximation matrix, and products and lifts use the finite formulas in \cref{prop:base-lift}.
All equations hold in their full parameter contexts.
\end{enumerate}
\end{proof}

\subsection{Exact transport of partial maps and description algebras}
\label{int:sec:transport}
The preceding domain--value comparison extends from primitives to
their finite descriptions. Write $L_{\mathcal E}=L_{\Sigma_{\mathcal E}}$.
For a geometric morphism $f:\mathcal F\to\mathcal E$, use the
constructed dominance $\Delta=\Delta_{f,B}$ of
\cref{prop:generated-imported-scope}, and write $L_\Delta$ for its
partial-element monad. All family formulas below are read sortwise.

For every geometric morphism the action
$(X\leftarrow D\to Y)\mapsto(f^*X\leftarrow f^*D\to f^*Y)$
preserves partial-map composition, since it is defined by pullback.
Pulled-back $\Sigma_{\mathcal E}$-domains have their classifiers in
$\Theta_f$. The following comparison represents these spans in the
constructed dominance $\Delta$.

\begin{construction}[The partial-element comparison]
\label{int:con:partialcomparison}
Let $L_{\mathcal E}X\leftarrow D_X\to X$ be the universal
partial-element span, whose domain classifier is the extent map to
$\Sigma_{\mathcal E}$. Its inverse image has classifier through
$\Theta_f\incl\Delta$. Classification in $\mathcal F$ therefore gives
a unique map
\[
 \ell_{f,X}:f^*L_{\mathcal E}X\longrightarrow L_\Delta(f^*X)
\]
whose universal domain--value span pulls back to
$f^*L_{\mathcal E}X\leftarrow f^*D_X\to f^*X$.
\end{construction}

\begin{proposition}[Coherence of the partial-element comparison]
\label{int:prop:partialcomparison}
The maps $\ell_{f,X}$ are natural in $X$ and preserve units,
multiplication, and strength in the following sense:
\begin{align}
 \ell_{f,X}f^*\eta^{L_{\mathcal E}}_X
   &=\eta^{L_\Delta}_{f^*X},\label{int:eq:ellunit}\\
 \ell_{f,X}f^*\mu^{L_{\mathcal E}}_X
   &=\mu^{L_\Delta}_{f^*X}\,
      L_\Delta(\ell_{f,X})\,\ell_{f,L_{\mathcal E}X}.
      \label{int:eq:ellmultiplication}
\end{align}
The finite sequencing comparison follows from strength. Each equation
identifies the domain and value of the represented partial map.
\end{proposition}
\begin{proof}
\leavevmode
\begin{enumerate}[label=\textup{(\roman*)},leftmargin=*,itemsep=3pt,topsep=4pt]
\item \textbf{Naturality.} Naturality is postcomposition of the value map of the universal span.
The unit is its total-domain case.
Multiplication composes the two
partial spans by pulling back the inner domain along the outer value.
Inverse image preserves that pullback, so the two sides of
\eqref{int:eq:ellmultiplication} classify the same composite span.
Strength forms a product with a total parameter; its domain and value
products are likewise preserved.
Finite sequencing iterates strength
and multiplication, giving the assertion for every finite arity.
\item \textbf{Uniqueness.} Uniqueness of partial-map classification proves the equations and
their coherence with parameter pullback.
No comparison of power
objects or of arbitrary exponentials is used.
\end{enumerate}
\end{proof}

\begin{proposition}[Geometric transport of finite syntax]
\label{int:prop:syntaxtransport}
For the supplied finitary signature $\mathcal S$ there are canonical
isomorphisms
\begin{equation}\label{int:eq:theta}
 \theta_{f,V}:f^*\Syn_{\mathcal S}(V)
   \xrightarrow{\cong}\Syn_{f^*\mathcal S}(f^*V).
\end{equation}
They preserve constructors, variable insertion, substitution, and
quotation into the supplied tree-data extension. They are coherent
with composition of geometric inverse images and parameter pullback.
\end{proposition}
\begin{proof}
\leavevmode
\begin{enumerate}[label=\textup{(\roman*)},leftmargin=*,itemsep=3pt,topsep=4pt]
\item \textbf{Finite witness construction.} Let $H_V$ be the many-sorted signature functor whose summands are
variable leaves, supplied literals, and constructor-labelled finite
products of children. Its initial algebra is the finite-tree object
constructed in \cref{int:con:syntax}. For a presented constructor
family, the finite list of input sorts is part of its data; dependent
sums over that family and products along these finite lists give
the same functor. Inverse image preserves the sums and finite
products, including their parameterized versions. It preserves the
NNO and finite-list recursion: the right adjoint lifts to algebras
for $1+(-)$ and $1+X\times(-)$ by adjunction, so their initial
algebras are preserved. Hence, also for the supplied varying finite
arities, there is a natural comparison
$f^*H_V\cong H_{f^*V}f^*$.
\item \textbf{Verification of the equations.} The same adjunction argument applies to $H_V$, using the induced
adjunction on families over the sort object. An algebra
$H_{f^*V}Y\to Y$ induces an $H_V$-algebra on $f_*Y$ by transposing
\[
 f^*H_Vf_*Y\cong H_{f^*V}f^*f_*Y
 \longrightarrow H_{f^*V}Y\longrightarrow Y.
\]
This lifts $f^*\dashv f_*$ to algebra categories. The lifted left
adjoint preserves the initial algebra, giving \eqref{int:eq:theta}.
\item \textbf{Well-founded induction.} Variable insertion and each constructor are preserved by its
definition. Substitution and data quotation are the specified
finite-tree recursions, so their comparisons follow by uniqueness
of those recursions. The comparison for a composite inverse image
and the composite of the two comparisons preserve the same
constructors; initiality makes them equal. The argument in a slice
gives parameter coherence.
\end{enumerate}
\end{proof}

\begin{theorem}[Partial-map transport]\label{thm:partial-transport}
Pullback preserves equality of partial maps as equality of domain--value spans. In particular, any naturality square of partial evaluators remains an equality of pulled-back domains and values. Under the lifted presentation comparisons, a presented partial map $v:D_C\to(D_{C'})_\bot$ satisfies
\[
 \chi_{f,C'}f^*v=v^f\beta_{f,C}.
\]
Its native definedness, restricted along $\beta_{f,C}$, is exactly the pulled-back definedness of $v$.
\end{theorem}
\begin{proof}
\leavevmode
\begin{enumerate}[label=\textup{(\roman*)},leftmargin=*,itemsep=3pt,topsep=4pt]
\item \textbf{Partial-map transport.} For spans, equality means an isomorphism of domains commuting with their monos and value maps.
Inverse image preserves that isomorphism and all commuting triangles.
Composition is preserved because its domain is a pullback of the second domain along the first value map.
\item \textbf{Naturality.} For the lifted form, apply \eqref{eq:rounded-naturality} to the map into the lifted presentation and then use $\lambda_{f,C'}$.
Composing with native definedness and using \eqref{eq:lift-def-base} gives
\[
 \defin_{f^\dagger C'}v^f\beta_{f,C}
 =\tau_f f^*(\defin_{C'}v).
\]
Thus the domains coincide as subobjects, including the truth comparison.
On that subobject, $\chi$ restricts to $\beta_{f,C'}$, so the same naturality equation compares their values.
This proves both components of partial-map equality.
\end{enumerate}
\end{proof}


\begin{theorem}[Exact geometric interpretation of descriptions]
\label{int:thm:transport}
Transport each primitive by
$o^f=\ell_{f,X_t}f^*o$, using the finite-product comparison on inputs,
and interpret the transported literals in the same way. Let
$\alpha_X^f:\Syn_{f^*\mathcal S}(L_\Delta f^*X)\to L_\Delta f^*X$
be the resulting tree algebra. Then
\begin{equation}\label{int:eq:transportalgebra}
\begin{tikzcd}[column sep=large]
 f^*\Syn(L_{\mathcal E}X)\ar[r,"f^*\alpha_X"]
   \ar[d,"{\Syn(\ell_{f,X})\theta_{f,L_{\mathcal E}X}}"']&
 f^*L_{\mathcal E}X\ar[d,"{\ell_{f,X}}"]\\
 \Syn(L_\Delta f^*X)\ar[r,"\alpha_X^f"']&L_\Delta f^*X.
\end{tikzcd}
\end{equation}
The equation preserves the full termination domain and its returned
value, respects substitution, and is natural in parameters.

For the displayed native semantic sorts $X^f$ and their comparison
maps $\beta_X:f^*X\to X^f$, use the exact primitive domain--value
comparisons of \cref{thm:partial-transport}. Put
\[
 \gamma_X=L_\Delta(\beta_X)\ell_{f,X}.
\]
The native tree algebra $\alpha_{X^f}$ then satisfies
\begin{equation}\label{int:eq:nativetree}
 \gamma_Xf^*\alpha_X
 =\alpha_{X^f}\Syn(\gamma_X)\theta_{f,L_{\mathcal E}X}.
\end{equation}
\end{theorem}
\begin{proof}
\leavevmode
\begin{enumerate}[label=\textup{(\roman*)},leftmargin=*,itemsep=3pt,topsep=4pt]
\item \textbf{Partiality and definedness.} Identify the transported tree object by
\cref{int:prop:syntaxtransport}. At a variable,
\eqref{int:eq:transportalgebra} is the map $\ell_{f,X}$.
At a literal it is the definition of its transported partial value.
At a primitive node, \cref{int:prop:partialcomparison} gives
\[
 \ell_{f,X_t}f^*(o^\sharp)
 =(o^f)^\sharp\Bigl(\prod_k\ell_{f,X_{s_k}}\Bigr).
\]
Apply the induction equations to the children. This proves the square.
\item \textbf{Naturality.} For native sorts, equality of the existing primitive domains
and their values under $\beta$ is exactly
\[
 \gamma_{X_t}f^*o=o^{\mathrm{nat}}\Bigl(\prod_k\beta_{X_{s_k}}\Bigr).
\]
Naturality of strength and multiplication extends it to lifted
operations on partial values. The same three tree clauses therefore
prove \eqref{int:eq:nativetree}. This applies to the actual supplied
comparisons, including non-epic ones.
\item \textbf{Evaluation and transposition.} Both $\ell$ and $L_\Delta\beta$ retain the classified domain of
their partial element. Equality of the two evaluations consequently
identifies their definedness subobjects and their values there.
The substitution law follows from \cref{int:prop:substitution} and
preservation of substitution by $\theta$. Pullback of the primitive
domain squares and uniqueness of the recursion clauses give parameter
coherence. Successive transports agree whenever compared in a common
target dominance containing their transported extents, since they
classify the same pulled-back spans and use the same tree recursion.
\end{enumerate}
\end{proof}

\begin{corollary}[Description kernels and oracle interpretation]
\label{int:cor:kerneltransport}\label{cor:kernel-basechange}
For every finite expression in the implemented signature and every
displayed positive output observation, its transported computation
kernel is the pullback of the original kernel. Under the native
comparison this includes definedness. Substitution of the generic
oracle gives the forcing truth of that full kernel along the localic
realization, as in \cref{thm:native-generic-computation}.
\end{corollary}
\begin{proof}
\leavevmode
\begin{enumerate}[label=\textup{(\roman*)},leftmargin=*,itemsep=3pt,topsep=4pt]
\item \textbf{Description kernels and oracle interpretation.} The kernel is the inverse image of the lifted output observation
under the expression evaluator.
Pull back the relevant square of
\cref{int:thm:transport} along that observation and use its displayed
membership comparison.
Taking the total observation gives the
termination subobject.
Substituting the generic parameter and applying
the existing kernel comparison gives the stated forcing truth.
\end{enumerate}
\end{proof}

\begin{proposition}[Transport of temporal self-interpretation]
\label{int:prop:temporaltransport}
Inverse image acts on internal temporal diagrams and gives canonical
comparisons
\[
 \operatorname{Seq}(f^*)\Prior\cong\Prior\operatorname{Seq}(f^*),
 \qquad
 \operatorname{Seq}(f^*)\triangleright
       \cong\triangleright\operatorname{Seq}(f^*).
\]
It preserves the displayed $p$, $\operatorname{next}$, and indexed
pullback diagrams of \cref{int:con:temporal}, and takes a constructed
guarded solution to the constructed solution of the transported map.
For running families, it preserves their finite-completion
domain--value spans, represented in the generated effective scope.
\end{proposition}
\begin{proof}
\leavevmode
\begin{enumerate}[label=\textup{(\roman*)},leftmargin=*,itemsep=3pt,topsep=4pt]
\item \textbf{Pullback construction.} Preservation of the NNO identifies the internal stage indices.
Prior uses successor reindexing; later uses that reindexing and a
terminal zeroth stage.
Inverse image preserves these operations and
the adjacent restrictions.
The indexed diagrams use pullbacks, so
their comparisons follow from left exactness.
Applying inverse image
to the guarded solution gives a solution of the transported equation;
the stage-zero and successor recurrence proves its uniqueness there.
\item \textbf{Partiality and definedness.} The finite-return domains are obtained by finite pullbacks, tagged
sums, and the internally indexed dependent sum of the successful
strata.
These constructions and their images are preserved.
The
pulled-back classifier factors through $\Theta_f\incl\Delta$, which
represents the resulting partial span by \cref{int:con:partialcomparison}.
This transports the displayed indexed structure and its solutions;
no preservation of all slice exponential objects is used.
\end{enumerate}
\end{proof}


Internally, positive logical formulas retain the same effective quantifier proof. For a dominance $\Theta$ in $\FF$, substitution is pullback of $\Theta$-subobjects, finite conjunction is composition of their monos, and an overt disjunction or existential is the supplied positive direct image. Structural induction therefore constructs a uniform $\Theta$-classifier for every such formula. This is internal effective truth. A stage semidecision in $\E$ instead requires a lift into $S_j$, exactly as in \cref{thm:admissibility}; no equality between these two objects has been used.

\subsection{The actual generic-relative evaluation square}
Suppose $\dot g:1\to\pi^*X$ is the reconstructed generic oracle of \cref{thm:object-generic}. Let $R$ be the program domain and $A$ the input domain, and write $j_R:\Pt R\to R$, $j_A:\Pt A\to A$ for the effective-point inclusions. Define the old-object generic-relative evaluation by the composite
\[
 \pi^*(\Pt R\times\Pt A)
 \xrightarrow{(\pi^*j_R,\dot g,\pi^*j_A)}
 \pi^*(R\times X\times A)
 \xrightarrow{\pi^*\operatorname{Eval}}\pi^*(Y_\bot).
\]
Call it $\operatorname{Eval}_{\dot g}$. If $\dot g^\dagger=\beta_{\pi,X}\dot g$, define the native presented evaluation by substituting
$\beta_{\pi,R}\pi^*j_R$, $\dot g^\dagger$, and
$\beta_{\pi,A}\pi^*j_A$ into $\operatorname{Eval}^\pi$.

\begin{proposition}[Generic-relative domains and values]\label{prop:generic-evaluation}
The two evaluations satisfy
\[
 \chi_{\pi,Y}\operatorname{Eval}_{\dot g}
 =\operatorname{Eval}^{\dagger}_{\dot g^\dagger}.
\]
Their definedness predicates agree through $\tau_\pi$, and on the common defined extent their values agree through $\beta_{\pi,Y}$. For every generalized output observation $V$, the relative truth of the old-object computation terminating in $\pi^*V$ is
\[
 \rho^*K_{e,a,V}.
\]
\end{proposition}
\begin{proof}
\leavevmode
\begin{enumerate}[label=\textup{(\roman*)},leftmargin=*,itemsep=3pt,topsep=4pt]
\item \textbf{Generic-relative domains and values.} Use \cref{cor:beta-product} to compare the triple product.
Naturality of the evaluator from \cref{thm:partial-transport} gives the commuting evaluation square.
Precompose it with the three displayed parameters.
Its oracle component becomes $\dot g^\dagger$ by definition and its other components become the compared effective program and input, proving the first equation.
\item \textbf{Partiality and definedness.} Compose with definedness and use \eqref{eq:lift-def-base}.
Restrict to the resulting equal subobjects; on those subobjects $\chi_{\pi,Y}$ is $\beta_{\pi,Y}$ by \cref{prop:base-lift}.
This proves the domain and value assertions.
\item For the final equation, the computation predicate is the pullback of $V^+$ along the evaluator with its program and input substituted.
That is the predicate $\dot g\in\pi^*K_{e,a,V}$.
The relative Scott point of $\dot g$ is the realization by \cref{thm:object-generic}, so its relative truth is precisely $\rho^*K_{e,a,V}$.
The calculation is with generalized parameters throughout, and hence proves the family equation.
\end{enumerate}
\end{proof}

The proposition also covers halting by taking $V=\top$. Unlike the unconditional native theorem \cref{thm:native-generic-computation}, this comparison has the additional conclusion of an equality involving the pulled-back old evaluator. It makes the two routes agree: first reconstruct an object-valued generic and evaluate, or first take the full oracle kernel and interpret it through the localic realization. The native comparison proves the agreement without identifying the two completion objects in cases where the comparison is not invertible.


\part{Oracle interfaces and localization}\label{part:modalities}

The previous parts provide the forcing semantics and the change-of-base comparisons. An oracle interface introduces another distinction: making some answer locally available does not choose the answer, and existence of an answer does not determine its computational use. We construct the induced localization, its well-founded protocols, and the gluing that retains both operational descriptions and sheafified values.


\section{Queries, answer objects, and operational reductions}\label{sec:answers}


A query has an answer type as well as a truth asserting that some answer exists. This distinction is already visible in a total binary oracle: the existence assertion is true for every coordinate, while the chosen bit remains computational data.

\begin{definition}[Answer interface]\label{def:answer-interface}
In a topos $\FF$, an answer interface is a map $d:A\to Q$. Its fibre over $q$ is the answer object for that query. Its support is the image $P_d\incl Q$. A realization of the interface is a section of $d$.
\end{definition}

The definition allows multivalued answers. It also allows a query object depending on local context. The support forgets the difference between two answers to the same query, so it will be used to define the required localization, while the map $d$ remains the operational datum.

A common construction starts in $\E$ with a query object $Q$, a response map $r:R\to Q$, and a family of answer opens $a:R\to\OO(X)$. For $\rho:L\to X$, the interpreted predicates $\rho^*a(t)$ define a subobject
\[
 \Ans_{\rho,a}\incl\pi^*R.
\]
Its composite with $\pi^*r$ is the generic answer projection.

\begin{proposition}[Exhaustiveness and functionality]\label{prop:answer-reconstruction}
The generic answer projection is epic exactly when
\[
 \bigvee_{t:r(t)=q}\rho^*a(t)=1
\]
in each query context. It is monic exactly when two valid responses to the same query are equal. When both hold, its inverse followed by the inclusion in $\pi^*R$ gives the unique answer map
\[
 \operatorname{ans}:\pi^*Q\longrightarrow\pi^*R
\]
over $\pi^*Q$.
\end{proposition}
\begin{proof}
\leavevmode
\begin{enumerate}[label=\textup{(\roman*)},leftmargin=*,itemsep=3pt,topsep=4pt]
\item \textbf{Kernel-pair descent.} The support of the projection is its regular image, hence is classified by the existential join of the answer predicates in that fibre.
It is all of $\pi^*Q$ precisely when the displayed equation holds.
Monicity means that the two projections of the kernel pair agree.
The kernel pair consists of valid answers $t,t'$ with $r(t)=r(t')$, so this is exactly functionality. A monic regular epimorphism in a topos is an isomorphism.
Its inverse supplies the map, and uniqueness follows because the inverse image of a query has only one member.
\end{enumerate}
\end{proof}

The construction applies to the Cohen coordinate interface from \cref{ex:cohen-logic}. It also applies to a coordinate relation presenting a total increasing sequence or a homogeneous subsequence. A partial functional relation gives instead an answer map defined on its support; its computational domain is effective precisely when that support has the indicated dominance classification.

\begin{proposition}[Effective answer observations]\label{prop:answer-effective}
If a selected family of response opens has its interpretation lifted to $S_j$, its validity has a uniform effective trace. Pullback along a response or query substitution gives the corresponding substituted trace.
\end{proposition}
\begin{proof}
\leavevmode
\begin{enumerate}[label=\textup{(\roman*)},leftmargin=*,itemsep=3pt,topsep=4pt]
\item \textbf{Effective answer observations.} Apply \cref{thm:admissibility} to the family of response opens.
Substitution precomposes its lift, and the pullback universal property identifies the result.
The stage and information-state evaluations are the formulas from \cref{sec:markings}.
\end{enumerate}
\end{proof}

\subsection{Affirmative queries, answer support, and operational information}
\label{sec:affirmative-interfaces}
An affirmative oracle is specified by a query object $Q$ and a family of oracle opens $a:Q\to\OO(X)$. Along a realization $\rho:L\to X$ it gives the relative predicate
\[
 P_{\rho,a}(q)=\rho^*a(q)\in\OO(L).
\]
For the generic oracle this is the proposition $\dot G\in a(q)$. Such a query requests positive evidence of that observation. It does not ask for a Boolean decision: forcing $P_{\rho,a}(q)$ makes the affirmative proposition true, whereas answering a decision query requires specifying which one of two disjoint responses occurs.

For a genuinely multivalued interface retain a response map $r:R\to Q$ and an open family $a:R\to\OO(X)$. Relative truth classification constructs the generic answer subobject $A_{\rho,a}\incl\pi^*R$ with membership predicate $\dot G\in a(v)$. Its answer projection $d:A_{\rho,a}\to\pi^*Q$ has support
\begin{equation}\label{eq:oracle-support}
 P_d(q)=\bigvee_{v:r(v)=q}\rho^*a(v)\in\OO(L).
\end{equation}
This is the regular-image formula: the support asserts that \emph{some} valid response exists, not which response was used. Every selected effective family of response opens $a_\Sigma:R\to M_X$ factors through $\Adm(\rho)$ precisely on its maximal subobject of forcing-effective responses; on that subobject the stage or information-state observation maps of \cref{sec:markings} semidecide response validity. These calculations do not use exhaustiveness or functionality.

\begin{proposition}[Answer existence does not determine the operational oracle]\label{prop:answer-support-quotient}
If the answer projection $d:A\to Q$ in a topos is epic, its support predicate is constantly true and its least answer-existence localization is the identity. The fibres of $d$ and the derived response transformations are not determined by that support. In particular, distinct epic answer maps can induce the same identity modality while having different answer objects and different finite protocol objects.
\end{proposition}
\begin{proof}
\leavevmode
\begin{enumerate}[label=\textup{(\roman*)},leftmargin=*,itemsep=3pt,topsep=4pt]
\item \textbf{Unit and identity laws.} The image factorization of an epimorphism has image $Q$, so $P_d(q)=\top$ in every query context.
The identity local operator forces this predicate, and is least.
To see that the fibres are not reconstructed, take $Q=1$ in $\Eff$ and the epic maps $1\to1$ and $2\to1$.
Both supports are true.
Their answer fibres are respectively $1$ and $2$, and the query-continuation factors $Z^1$ and $Z^2$ in the protocol constructor need not be isomorphic.
Hence the support forgets information retained by actual response objects.
\end{enumerate}
\end{proof}

\begin{definition}[Container morphism]\label{def:oracle-reduction}
A reduction from $d:A\to Q$ to $d':A'\to Q'$ consists of $t:Q\to Q'$ and an answer transformation $v:Q\times_{Q'}A'\to A$ over $Q$.
\end{definition}
Thus an answer to the translated query determines an answer to the original query. For two reductions, pull the second answer transformation back along the first query map and compose. The canonical pullback associator identifies the two composites of three reductions, since both send a final answer through the same two answer transformations. This proves associativity on the actual response data.

A stronger separation is constructed in \cref{obs:thm:crown}.
Its two epic answer maps are even locally isomorphic, but
their container-morphism order is strict. The distinction is
measured by a nonzero torsor class in the same effective
domain sheaf setting. Thus the missing global response
transformation has a cohomological obstruction, beyond
the difference in fibre cardinalities above.

Finite adaptive interaction is represented without discarding those data: after an answer $a\in A_q$, the next query may be a function of $a$. For a prescribed finite interaction protocol its iterated dependent sums form the object of transcripts. Evaluation of that protocol is proved by induction on its finite number of steps. At a terminal step it is the specified value; at a query step it is substitution of the received answer into the continuation. The resulting domain is the iterated dependent extent of the answers used. The partial-map and evaluator constructions of \cref{thm:canonicalcoding} implement this operation whenever those extents are computationally classified. Taking support at the end expresses successful termination, while the transcript object retains the answers and their order.


The answer map, not its support alone, determines which continuation programs may depend on which replies. The localization construction will make those answers locally available through a contextual cover while keeping the absence of a canonical selection explicit.

\section{Contextual rules and the cover-answer construction}\label{sec:rules}

A query may become answerable only after passing to a cover of the current context. To encode this without choosing a global answer, we construct the associated contextual coverage and its least localization. The resulting sheafification solves the existence problem while retaining answer dependence in its protocol presentation.


We first construct generated localizations directly as least solutions of equations. This gives existence within the elementary-topos setting without stipulating a saturation operator.

\begin{lemma}[Meets of nuclei and generated rules]\label{lem:generated-nucleus}
For an internal family of nuclei $(k_t)$ on a frame $H$, its pointwise meet is a nucleus. For a family of pairs $(a_r,b_r)$, there is a least nucleus $j_\mathcal R$ satisfying
\[
 j_\mathcal R(a_r)\leq j_\mathcal R(b_r)\quad\text{for all }r.
\]
It is the meet of all validating nuclei.
\end{lemma}
\begin{proof}
\leavevmode
\begin{enumerate}[label=\textup{(\roman*)},leftmargin=*,itemsep=3pt,topsep=4pt]
\item \textbf{Meets of nuclei and generated rules.} Put $k(x)=\bigwedge_tk_t(x)$.
It is monotone, inflationary, and finite-meet preserving.
Since $k(x)\leq k_t(x)$,
$k_t(k(x))\leq k_t^2(x)=k_t(x)$.
Also $k(k(x))\leq k_t(k(x))$ for each $t$.
Taking the meet gives $k^2(x)\leq k(x)$, and inflationarity gives equality.
\item \textbf{Localization and coverage.} For a nucleus $k$, the equation $k(a)\leq k(b)$ is equivalent to $a\leq k(b)$.
This latter form is preserved by pointwise meets of validating nuclei.
The constant-top nucleus validates every rule, so the class of validators is inhabited.
It is an internal subobject of the object of endomorphisms of $H$, cut out by the nucleus laws and the rule equations.
Its meet therefore exists internally and is the least validator.
\end{enumerate}
\end{proof}

\subsection{Elementary localizations}
The least-rule construction includes the indicated localizations with explicit formulas.

\begin{proposition}[Open and closed requirements]\label{prop:elementary-nuclei}
The least nucleus imposing $1\vdash a$ is
$j_a(x)=a\Rightarrow x$. The least nucleus imposing $a\vdash0$ is
$c_a(x)=a\vee x$. If two nuclei commute, their composite is their join in the ordered set of nuclei.
\end{proposition}
\begin{proof}
\leavevmode
\begin{enumerate}[label=\textup{(\roman*)},leftmargin=*,itemsep=3pt,topsep=4pt]
\item \textbf{Adjunction and universal property.} The right adjoint $a\Rightarrow-$ preserves meets.
It is inflationary because $a\wedge x\leq x$, and
$a\Rightarrow(a\Rightarrow x)=(a\wedge a)\Rightarrow x=a\Rightarrow x$.
Thus $j_a$ is a nucleus with $j_a(a)=1$.
If $k(a)=1$, apply $k$ to
$a\wedge(a\Rightarrow x)\leq x$:
\[
 k(a)\wedge k(a\Rightarrow x)\leq k(x).
\]
Since $k(a)=1$ and $a\Rightarrow x\leq k(a\Rightarrow x)$, this gives
$j_a(x)\leq k(x)$, proving minimality.
\item \textbf{Verification of the equations.} The map $c_a$ is inflationary and idempotent; frame distributivity gives finite-meet preservation.
It sends $a$ and $0$ to the same value.
If $k(a)=k0$, then $a\leq k0\leq kx$ and $x\leq kx$, so $a\vee x\leq kx$.
This proves its minimality.
\item \textbf{Universal-property argument.} If $jk=kj$, the composite preserves finite meets and is inflationary.
Its square is $j^2k^2=jk$.
It is above both factors.
If $h$ is above both, then $jk(x)\leq h(hx)=hx$, so it is their least upper bound.
\end{enumerate}
\end{proof}

Thus requiring a truth to hold restricts to its open sublocale, while requiring it to be false gives the complementary closed sublocale. For arbitrary families the least-validator construction remains available without a commuting-composite formula. In particular the proofs of tower coherence use the relative quotient maps, not an unproved identity between a join of nuclei and an ordinary composite.

\begin{proposition}[Nuclei and local operators]\label{prop:nucleus-topology}
For $\FF=\Sh_\E(L)$ and $H=\OO(L)$ there is an order isomorphism
\[
 \Nuc_\E(H)\cong\LT(\FF).
\]
For a nucleus $k$, its closure on a predicate, tested by a local section over an open $u$, replaces its local truth $v\leq u$ by $u\wedge k(v)$. On base-indexed relative subterminals its action is the familywise application of $k$.
\end{proposition}
\begin{proof}
\leavevmode
\begin{enumerate}[label=\textup{(\roman*)},leftmargin=*,itemsep=3pt,topsep=4pt]
\item \textbf{Pullback construction.} Let $M\incl A$ and let $a:u\to A$ be a local section.
Its pullback predicate is an open $v\leq u$.
Declare the section to lie in the closure when $u\leq k(v)$.
This declaration is stable under restriction to $u'\leq u$, because
\[
 u'\wedge k(u'\wedge v)=u'\wedge k(v).
\]
It is local under a cover: the condition says that the entire covered open lies in $k(v)$, and inclusion of opens can be tested on their join.
Thus it defines a subobject of $A$.
\item \textbf{Verification of the equations.} Inflationarity, monotonicity, finite-meet preservation, and idempotence follow from the corresponding laws of $k$ on these local truth values.
Pullback of a predicate along a sheaf map pulls back the local section and its local truth, so the closure is universal.
It therefore gives a Lawvere--Tierney topology.
\item \textbf{Converse construction.} Conversely a universal closure restricts to subterminals in every base context.
Through \cref{prop:relative-truth} these restrictions define a map $k:H\to H$, with the nucleus laws inherited from closure.
Every predicate is tested by local sections, and universality identifies its closure there with the subterminal calculation.
Thus the two constructions are inverse.
Their orders are both inclusion of closures, so the correspondence is an order isomorphism.
\item \textbf{Functoriality and coherence.} This is the explicit local form of the sublocale--subtopos correspondence \cite{johnstone,maclanemoerdijk}.
\end{enumerate}
\end{proof}

\begin{definition}[Oracle localization]\label{def:oracle-localization}
For an answer map $d:A\to Q$, let $o_d$ be the least Lawvere--Tierney topology making its support $P_d\incl Q$ dense. Equivalently, its sheafification makes $d$ epic.
\end{definition}
The equivalence follows because sheafification preserves the regular epi--mono factorization, and a mono becomes an isomorphism precisely when it is dense. Existence is the meet construction on local operators; universal closure may be tested on the support inclusion. For a localic sheaf topos, \cref{prop:nucleus-topology} gives the corresponding frame nucleus.

\begin{theorem}[Contextual rule compilation]\label{thm:rule-compiler}
Let $a,b:R\to H$ be an $\E$-indexed family of antecedents and consequents. In $\FF$, form
\[
 Q_\mathcal R=\sum_{r:\pi^*R}\widetilde a_r,\qquad
 A_\mathcal R=\sum_{r:\pi^*R}(\widetilde a_r\times\widetilde b_r),
 \qquad d_\mathcal R:A_\mathcal R\incl Q_\mathcal R.
\]
The nucleus corresponding to $o_{d_\mathcal R}$ is exactly the least nucleus validating $a_r\vdash b_r$ for all $r$.
\end{theorem}
\begin{proof}
\leavevmode
\begin{enumerate}[label=\textup{(\roman*)},leftmargin=*,itemsep=3pt,topsep=4pt]
\item \textbf{Localization and coverage.} A query is a rule index with evidence of its antecedent.
Its answer predicate is the consequent in that context.
By \cref{prop:nucleus-topology}, a nucleus $k$ makes this predicate dense precisely when
$a_r\leq k(b_r)$ in every index context.
By \cref{lem:generated-nucleus}, this is equivalent to $k(a_r)\leq k(b_r)$.
Thus the two constructions have exactly the same validating localizations, and their least elements correspond.

Naturality is expressed by the commutative diagram
\[
\begin{tikzcd}[column sep=large,row sep=large]
A_{\mathcal R}\ar[r,hook,"d_{\mathcal R}"]\ar[d]&Q_{\mathcal R}\ar[d]\\ \pi^*R\ar[r,"\id"']&\pi^*R
\end{tikzcd}
\]
\end{enumerate}
\end{proof}

The antecedent belongs to the query context. Replacing the query object by just $R$ would require $b_r$ globally, even where its antecedent is unavailable. The dependent sum in the theorem is the construction that retains this context.

\subsection{Exact comparison with Ahman--Bauer oracle modalities}
\label{sec:ahman-bauer-comparison}
The propositions supplied by an oracle and the answers carried by its response object play different roles. Ahman and Bauer \cite{ahmanbauer} associate to a predicate in a topos its least oracle modality, and also study computation trees and equifoliate trees retaining the intensional answer data. We first identify their \emph{propositional} construction exactly with the relative nucleus just constructed, without asserting an identification of those computation trees.

Write $H=\OO(L)$ and $\FF=\Sh_\E(L)$. Let
\[
 \sigma_L:\pi^*H\longrightarrow\Omega_\FF
\]
be the map classifying the relative subterminal indexed by an open of $L$, supplied by \cref{prop:relative-truth}. An $\E$-indexed predicate $P:Q\to H$ determines an internal predicate
\[
 \widetilde P=\sigma_L\pi^*P:\pi^*Q\longrightarrow\Omega_\FF.
\]
Let $o_P$ be the least nucleus on $H$ with $o_P(P(q))=1$ for all generalized $q:Q$, as obtained from the rule compiler by the rules $1\vdash P(q)$. Let $o_{\widetilde P}$ be the Ahman--Bauer oracle modality in $\FF$, characterized as the least Lawvere--Tierney topology forcing $\widetilde P$.

\begin{theorem}[Exact oracle-modality comparison]\label{thm:ab-comparison}
Under the nucleus--Lawvere--Tierney order isomorphism of \cref{prop:nucleus-topology},
\begin{equation}\label{eq:ab-modality-comparison}
 \widehat{o_P}=o_{\widetilde P}.
\end{equation}
For an affirmative query family $a:Q\to\OO(X)$ and a realization $\rho:L\to X$, the same equality holds with $P=P_{\rho,a}=\rho^*a$. All equations are internal in the full query context.
\end{theorem}
\begin{proof}
\leavevmode
\begin{enumerate}[label=\textup{(\roman*)},leftmargin=*,itemsep=3pt,topsep=4pt]
\item \textbf{Unit and identity laws.} Fix a nucleus $k:H\to H$ and write $\widehat k$ for its Lawvere--Tierney topology.
The local-section description in \cref{prop:nucleus-topology} gives the familywise identity
\[
 \widehat k\,\sigma_L\pi^*P
   =\sigma_L\pi^*(kP).
\]
Indeed, on a generalized stage $u$ the closure of the relative subterminal represented by $P(q)$ is represented by $u\wedge k(P(q))$; this is precisely the relative truth classified by $k(P(q))$.
Consequently $\widehat k$ forces $\widetilde P$ if and only if the family $kP:Q\to H$ is constantly top.
\item \textbf{Naturality and reindexing.} This equivalence can be tested after every parameter substitution and so does not replace generalized queries by global ones.
The two predicates therefore have exactly the corresponding validating nuclei and local operators.
Since the correspondence is an order isomorphism, their least validators agree, yielding \eqref{eq:ab-modality-comparison}.
Substituting $\rho^*a$ proves the last assertion.
\end{enumerate}
\end{proof}

\begin{theorem}[Every topology has a canonical propositional oracle]\label{thm:canonical-modal-oracle}
Let $\mathcal T$ denote $\Eff$ or one of its forcing sheaf topoi, and let $j:\Omega_\mathcal T\to\Omega_\mathcal T$ a Lawvere--Tierney topology. Define the query object and its predicate by
\[
 Q_j=\sum_{p:\Omega_\mathcal T}j(p),\qquad
 P_j(p,h)=p.
\]
Then $o_{P_j}=j$. Moreover, for every predicate $P:A\to\Omega_\mathcal T$,
\begin{equation}\label{eq:modal-container-adjunction}
 o_P\leq j\quad\Longleftrightarrow\quad
 \text{there is a propositional-container morphism }P\to P_j.
\end{equation}
Here a container morphism $P\to Q$ consists of a query translation $t:A\to B$ and implications $Q(t(a))\Rightarrow P(a)$. These equivalences exhibit the oracle-modality assignment as a left adjoint retraction onto the poset of modalities from the preorder reflection of propositional containers.
\end{theorem}
\begin{proof}
\leavevmode
\begin{enumerate}[label=\textup{(\roman*)},leftmargin=*,itemsep=3pt,topsep=4pt]
\item \textbf{Oracle interpretation.} The topology $j$ forces $P_j$: a query $(p,h)$ carries an inhabitant $h:j(p)$.
Thus $o_{P_j}\leq j$.
Conversely, under a proof of $j(p)$ the pair $(p,h)$ is a query, and $o_{P_j}$ forces its proposition $p$.
Therefore $j(p)\Rightarrow o_{P_j}(p)$ for every $p$, giving $j\leq o_{P_j}$.
For \eqref{eq:modal-container-adjunction}, $o_P\leq j$ means that $j$ forces each $P(a)$.
The query translation
\[
 a\longmapsto(P(a),\,\text{the supplied proof of }j(P(a)))
\]
together with the identity implication $P(a)\Rightarrow P(a)$ is a container morphism $P\to P_j$.
\item \textbf{Converse construction.} Conversely, a morphism supplies, for every $a$, a pair $(p,h:j(p))$ and an implication $p\Rightarrow P(a)$.
Monotonicity of $j$ makes $j(P(a))$ true.
Thus $j$ forces $P$ and $o_P\leq j$.
These maps are defined uniformly in their contexts.
Taking $P=P_j$ and the first paragraph gives the retraction equation.
\end{enumerate}
\end{proof}


\begin{theorem}[The cover-answer oracle]\label{thm:cover-answer}
Let $\mathcal R$ be the specified generating covers of a condition preorder $P$, with generated topology $J$. In $\mathcal A_P=\Sh_\E(\Alex^\uparrow(P))$, let $\dot I$ be the generic ideal. Form
\[
 Q_\mathcal R=\sum_{(p,S):\mathcal R}\dot I(p),\qquad
 A_\mathcal R=\sum_{(p,S):\mathcal R}\sum_{q\in S}\dot I(q).
\]
Lowerness gives an answer projection $d_\mathcal R:A_\mathcal R\to Q_\mathcal R$. Its oracle localization is the local operator corresponding to $j_J$.
\end{theorem}
\begin{proof}
\leavevmode
\begin{enumerate}[label=\textup{(\roman*)},leftmargin=*,itemsep=3pt,topsep=4pt]
\item \textbf{Localization and coverage.} Since $q\in S$ implies $p\leq q$, generic membership of $q$ gives membership of $p$.
This defines the projection.
The query context has open $\uparrow p$, and the answer support has open
$\bigvee_{q\in S}\uparrow q=S$. A nucleus makes the support dense in this context exactly when
\[
 k(\uparrow p)\leq k(S).
\]
The reverse inequality holds automatically because $S\subseteq\uparrow p$.
Hence these are precisely the cover equations.
\item \textbf{Verification of the equations.} Their least validating nucleus is $j_J$: its topology contains the generators, and any validating nucleus induces a topology containing those generators by the semantic-cover formula of \cref{thm:basis}.
Minimality in either construction gives equality.
\item \textbf{Localization and coverage.} The answer object itself is the dependent sum over refinements.
Its support determines the localization, but its points still record which refinement answered the query.
\end{enumerate}
\end{proof}

\begin{corollary}[The forcing coverage as an oracle modality]\label{cor:cover-ab}
For an internal condition site $(P,J)$ with generating covers $\mathcal R$, the Lawvere--Tierney topology on the Alexandrov sheaf topos $\Sh_\E(\Alex^\uparrow(P))$ determined by $j_J$ is the Ahman--Bauer modality of the contextual cover-answer predicate of \cref{thm:cover-answer}. Its associated sheaf subtopos is $\Sh_\E(P,J)$.
\end{corollary}
\begin{proof}
\leavevmode
\begin{enumerate}[label=\textup{(\roman*)},leftmargin=*,itemsep=3pt,topsep=4pt]
\item \textbf{The forcing coverage as an oracle modality.} Apply \cref{thm:ab-comparison} to the answer support computed in \cref{thm:cover-answer}.
Its least validating nucleus is $j_J$, and \cref{thm:site-scott} identifies the corresponding subtopos with $\Sh_\E(P,J)$.
\end{enumerate}
\end{proof}


\subsection{Oracle sheafification on all objects}\label{sec:full-oracle-sheaves}
The comparison in \cref{thm:ab-comparison} is an equality of
Lawvere--Tierney topologies. We now recover the complete associated
sheaf construction. Work in $\mathcal T=\Eff$ or in one of its forcing sheaf topoi,
let $d:A\to Q$ be an answer map, put
$P_d(q)=\exists_{a:A_q}\top$, and write $j=o_d$.
In this subsection $a_jX$ also denotes the underlying object in
$\mathcal T$ of the associated $j$-sheaf.

\begin{lemma}[Explicit inductive oracle closure]\label{lem:oracle-inductive-closure}
For $s\in\Omega_{\mathcal T}$ put
\[
 \Phi_s(z)=s\vee\exists_{q:Q}(P_d(q)\Rightarrow z),
 \qquad
 J_d(s)=\bigwedge\{z\mid\Phi_s(z)\leq z\}.
\]
Then $J_d=o_d$. It is characterized by the induction principle
\[
 \frac{s\leq z\quad
   \forall q\,[(P_d(q)\Rightarrow z)\leq z]}
  {J_d(s)\leq z}.
\]
\end{lemma}
\begin{proof}
\leavevmode
\begin{enumerate}[label=\textup{(\roman*)},leftmargin=*,itemsep=3pt,topsep=4pt]
\item \textbf{Explicit inductive oracle closure.} The impredicative meet exists in the internal complete Heyting
algebra $\Omega_{\mathcal T}$.
The displayed operator is monotone
in $z$, so the meet of its prefixed points is its least fixed point.
It is monotone in $s$, inflationary and idempotent.
It forces each
$P_d(q)$, because $(P_d(q)\Rightarrow J_d(P_d(q)))\leq
J_d(P_d(q))$, and the antecedent holds at $z=J_d(P_d(q))$.
\item \textbf{Oracle interpretation.} To verify finite meets, let $t=J_d(s\wedge r)$.
The induction
principle, first applied to $J_d(r)$ with target $s\Rightarrow t$,
and then to $J_d(s)$ with target $J_d(r)\Rightarrow t$, gives
$J_d(s)\wedge J_d(r)\leq t$; the reverse inequality follows
by monotonicity.
Explicitly, the query step of the first induction
uses
$P_d(q)\Rightarrow(s\Rightarrow t)$ and the fact that
$s\wedge P_d(q)\Rightarrow t$; the second induction repeats the
same implication calculation with $J_d(r)$ in place of $r$.
Thus $J_d$ is a local operator.
\item \textbf{Forcing construction.} If $k$ is another local operator
forcing every $P_d(q)$, its value $k(s)$ is a prefixed point of
$\Phi_s$: $s\leq k(s)$ and
$(P_d(q)\Rightarrow k(s))\leq k(s)$, the latter by applying
$k$, using $kP_d(q)=\top$, and idempotence.
Consequently
$J_d(s)\leq k(s)$.
This identifies the least forcing operator
with $o_d$.
\end{enumerate}
\end{proof}

\begin{theorem}[Oracle sheaves as uniquely determined algebras]
\label{thm:oracle-sheaf-algebras}
An object $Z$ is a $j$-sheaf if and only if the restrictions
\[
 \operatorname{res}_{q,Z}:Z\longrightarrow Z^{P_d(q)}
\]
are isomorphisms, uniformly in $q:Q$. Their inverses are the
unique oracle-query operations
$\operatorname{ask}_{q,Z}:Z^{P_d(q)}\to Z$. They satisfy
\[
 u:P_d(q)\ \Longrightarrow\
 \operatorname{ask}_{q,Z}(h)=h(u),\qquad
 P_d(q)\Rightarrow z=z'\ \Longrightarrow\ z=z'.
\]
Every morphism of $j$-sheaves preserves these operations.
\end{theorem}
\begin{proof}
\leavevmode
\begin{enumerate}[label=\textup{(\roman*)},leftmargin=*,itemsep=3pt,topsep=4pt]
\item \textbf{Oracle sheaves as uniquely determined algebras.} Since $jP_d(q)=\top$, a $j$-sheaf extends sections uniquely
along $P_d(q)\incl1$, proving the forward implication.
Conversely,
the indicated inverse equations imply the second equality rule.
Apply the induction principle of
\cref{lem:oracle-inductive-closure} to the equality predicate
$z=z'$: $j(z=z')$ implies equality because equality is stable under
each generating query rule.
Thus $Z$ is $j$-separated.
\item \textbf{Forcing construction.} For a general $j$-dense mono $B'\incl B$ and $h:B'\to Z$,
the predicate ``the graph of $h$ has a unique compatible value over
$b$'' contains $B'$.
The uniqueness follows from separatedness.
It is closed under the generating query rule: on $P_d(q)$ the
candidate values give $P_d(q)\to Z$, and
$\operatorname{ask}_{q,Z}$ glues them; its two equations ensure
compatibility with the original graph.
Induction on $J_d$ shows
that the predicate holds at every $b$, since the original subobject
is dense.
\item \textbf{Existence and uniqueness.} The unique-value relation is therefore the graph of a
map $B\to Z$, uniquely extending $h$.
This is the sheaf property.
For a sheaf map $f:Z\to Z'$, its two possible composites with
an $\operatorname{ask}$ operation agree after restriction to
$P_d(q)$.
That subterminal is dense and $Z'$ is separated, so
they agree globally.
This is the oracle-sheaf-algebra calculation
of \cite[Section~4]{ahmanbauer} expressed through the subobject
classifier of the present topos.
\end{enumerate}
\end{proof}

\begin{theorem}[Ordinary sheafification has the quotient-inductive
oracle universal property]\label{thm:oracle-qit-recovery}
The sheafification $a_jX$ carries
\[
 \operatorname{return}=\eta_X:X\to a_jX,\qquad
 \operatorname{ask}_{q}=\operatorname{res}_{q,a_jX}^{-1}.
\]
These operations satisfy both Ahman--Bauer conditional sheaf
equalities and, for every $j$-sheaf $Z$ and $f:X\to Z$, admit a
unique extension $\bar f:a_jX\to Z$ preserving $\operatorname{return}$
and $\operatorname{ask}$. Thus ordinary sheafification is the
conditional quotient-inductive oracle sheafification, by its
actual universal property.
\end{theorem}
\begin{proof}
\leavevmode
\begin{enumerate}[label=\textup{(\roman*)},leftmargin=*,itemsep=3pt,topsep=4pt]
\item \textbf{Oracle interpretation.} \Cref{thm:oracle-sheaf-algebras} constructs both query operations
and their equalities.
By the reflector adjunction there is exactly
one $\bar f$ with $\bar f\eta_X=f$.
The last assertion of that
theorem shows that $\bar f$ preserves the query operations:
under any $u:P_d(q)$ both sides restrict to $\bar f h(u)$, and
density plus separatedness gives global equality.
This is exactly
the stipulated conditional quotient-inductive induction principle
of \cite[Section~4.1]{ahmanbauer}; no additional inductive type
former is used.

The universal property is represented by the commutative diagram
\[
\begin{tikzcd}[column sep=large,row sep=large]
X\ar[r,"\eta_X"]\ar[dr,"f"']&a_jX\ar[d,"\bar f"]\\ &Z
\end{tikzcd}
\]
\end{enumerate}
\end{proof}

\begin{theorem}[Modal-singleton classifier for sheafification]
\label{thm:oracle-singleton-classifier}
Let $\Omega_j\incl\Omega_{\mathcal T}$ be the $j$-fixed truth object,
and put $e_x(y)=j(y=x)$. Then
\[
 \operatorname{MSing}_j(X)=
 \left\{u:X\to\Omega_j\ \middle|\
 j\bigl(\exists x:X\ \forall y:X\ [u(y)=e_x(y)]\bigr)\right\}
\]
is a $j$-sheaf, and the map $x\mapsto e_x$ exhibits it as the
sheafification of $X$. The resulting isomorphism
$\operatorname{MSing}_j(X)\cong a_jX$ is natural in $X$.
\end{theorem}
\begin{proof}
\leavevmode
\begin{enumerate}[label=\textup{(\roman*)},leftmargin=*,itemsep=3pt,topsep=4pt]
\item \textbf{Evaluation and transposition.} The object $\Omega_j$ is a $j$-sheaf.
Exponentiation by an
arbitrary $X$ preserves sheaves, since products with $X$ preserve
$j$-dense monos and the defining extension bijections transpose
across the exponential adjunction.
The displayed predicate is
$j$-closed: applying $j$ to it changes nothing, and its parameters
range over the sheaf $(\Omega_j)^X$. A closed subobject of a sheaf
is a sheaf, hence so is $\operatorname{MSing}_j(X)$.
\item \textbf{Kernel-pair descent.} For $u$ in this object let $R_u=\{x\mid u=e_x\}$.
Its support
is $j$-dense.
If $x,x'\in R_u$, then evaluating the equality
$e_x=e_{x'}$ at $x$ yields $j(x=x')$.
For a map $f:X\to Z$ into a
$j$-sheaf, separatedness gives $f(x)=f(x')$.
Thus the map
$R_u\to Z$, $x\mapsto f(x)$, equalizes the kernel pair of its
image and descends to its support.
The latter is a $j$-dense
subterminal, so the sheaf property extends this map uniquely
to an element of $Z$.
The construction is uniform in $u$ and
therefore transposes to $\bar f:\operatorname{MSing}_j(X)\to Z$.
\item \textbf{Adjunction and universal property.} It satisfies $\bar f(e_x)=f(x)$.
Uniqueness follows by testing
on the same $j$-dense support and using separatedness.
Hence
$x\mapsto e_x$ has the sheafification universal property, which
produces the natural isomorphism.
\end{enumerate}
\end{proof}

\begin{theorem}[The entire forcing sheaf category is a cover-oracle
sheaf category]\label{thm:full-cover-oracle}
For an internal condition site $(P,J)$ generated by covers
$\mathcal R$, let $\mathcal T_0$ be the internal presheaf topos
$\widehat P\simeq\Sh_\E(\Alex^\uparrow(P))$.
In $\mathcal T_0$ let $y(p)$ be the relative representable and
$S\incl y(p)$ its covering sieve subobject. Define the contextual
query and answer objects
\[
 Q_{\mathcal R}=\sum_{(p,S):\mathcal R}y(p),\qquad
 A_{\mathcal R}=\sum_{(p,S):\mathcal R}S,
 \qquad d_{\mathcal R}:A_{\mathcal R}\incl Q_{\mathcal R}.
\]
Then there is an equivalence of reflective subcategories
\[
 \boxed{\Sh_\E(P,J)\simeq
 \Sh_{o_{d_{\mathcal R}}}(\mathcal T_0)}
\]
and an isomorphism between their associated-sheaf reflectors.
Consequently the sheafification of every presheaf $X$ has the
oracle quotient-inductive and modal-singleton universal properties
of \cref{thm:oracle-qit-recovery,thm:oracle-singleton-classifier}.
\end{theorem}
\begin{proof}
\leavevmode
\begin{enumerate}[label=\textup{(\roman*)},leftmargin=*,itemsep=3pt,topsep=4pt]
\item \textbf{Image and quotient construction.} The component $S\incl y(p)$ is the regular image of the refinement-valued answer projection in \cref{thm:cover-answer}: the representable $y(p)$ is the generic-membership subterminal, and the union of the answer subterminals $y(q)$, $q\in S$, is $S$.
Thus the answer object here is the support mono of that earlier interface.
By \cref{prop:answer-support-quotient}, passing to this regular image preserves the generated modality.
\item Consequently \cref{thm:cover-answer,cor:cover-ab} give
\[
 o_{d_{\mathcal R}}=\widehat{j_J}
\]
as Lawvere--Tierney topologies on the entire presheaf topos, in every object context.
The associated
$j_J$-subtopos of $\widehat P$ is $\Sh_\E(P,J)$ by the
site-to-localic-topos correspondence of
\cref{prop:stage-yoneda,thm:site-scott}.
Consequently the two
reflective inclusions into $\mathcal T_0$ have the same full
subcategory of objects and arrows; uniqueness of left adjoints
identifies their associated-sheaf reflectors naturally.
\item \textbf{Oracle interpretation.} The oracle-sheaf condition concerns extension and uniqueness
along the answer-support inclusion in \emph{every generalized query
context}. A bijection
$\Hom(y(p),Z)\to\Hom(S,Z)$ tested for only one generator at
its base stage is not by itself the sheaf condition for all its
refinements.
The equality of full Lawvere--Tierney topologies above
gives all those contextual pullbacks automatically.
\item Finally
\cref{thm:oracle-qit-recovery,thm:oracle-singleton-classifier}
apply to the common reflector, proving both stated universal
properties without any further axiom or choice of local answers.
\end{enumerate}
\end{proof}

\begin{corollary}[Every nucleus has a contextual oracle presentation]\label{cor:every-oracle}
For a nucleus $j$ on $H$, the contextual rules $j(u)\vdash u$, indexed by $u\in H$, generate exactly $j$.
\end{corollary}
\begin{proof}
\leavevmode
\begin{enumerate}[label=\textup{(\roman*)},leftmargin=*,itemsep=3pt,topsep=4pt]
\item \textbf{Every nucleus has a contextual oracle presentation.} The nucleus $j$ validates the rules by idempotence.
If $k$ validates them, then
$j(u)\leq k(u)$ for every $u$, so $j\leq k$.
Apply \cref{thm:rule-compiler}.
\end{enumerate}
\end{proof}

\begin{proposition}[Sums and reductions]\label{prop:modal-reductions}
The oracle localization of a sum of answer interfaces is the join of their oracle localizations. A reduction as in \cref{def:oracle-reduction} gives $o_d\leq o_{d'}$.
\end{proposition}
\begin{proof}
\leavevmode
\begin{enumerate}[label=\textup{(\roman*)},leftmargin=*,itemsep=3pt,topsep=4pt]
\item \textbf{Sums and reductions.} After any sheafification, a coproduct answer map is epic exactly when every component is epic, since pulling back along a summand recovers that component and the summands jointly cover the codomain.
Its validators are therefore the common upper bounds of the component localizations.
\item \textbf{Pullback construction.} For a reduction, a localization making $d'$ epic also makes its pullback $Q\times_{Q'}A'\to Q$ epic.
This map factors through the sheafified $d$, so $d$ is epic there.
Minimality gives $o_d\leq o_{d'}$.
\end{enumerate}
\end{proof}

The least-localization characterization agrees with the propositional oracle modality studied by Ahman and Bauer \cite{ahmanbauer}: both are least among the modalities forcing the support predicate. Their computational account retains answer-dependent interaction. Here the support is used only for the localization, and the original answer map remains part of the interface.


\begin{corollary}[Dense forcing as a cover-answer localization]
\label{cal:cor:dense-oracle}
Transport the extension-density nucleus through the Scott--Alexandrov
isomorphism: $\overline{\Dens}=\Phi_P\Dens\Phi_P^{-1}$. Index the
extension-dense sieves by $(p,S)$ and form the contextual answer map
\begin{equation}\label{cal:eq:answer}
 \delta_{\mathrm{dens}}:
 \coprod_{(p,S)}\coprod_{q\in S}yq
 \longrightarrow\coprod_{(p,S)}yp.
\end{equation}
Then there are canonical equivalences over the base
\begin{equation}\label{cal:eq:identity}
 \Sh_\E(P,\Jdens)
 \simeq\Sh_\E(L_{\overline{\Dens}})
 \simeq\Sh_{o_{\delta_{\mathrm{dens}}}}(\widehat P).
\end{equation}
They preserve the generic membership model and the associated reflectors.
\end{corollary}
\begin{proof}
\leavevmode
\begin{enumerate}[label=\textup{(\roman*)},leftmargin=*,itemsep=3pt,topsep=4pt]
\item \textbf{Dense forcing as a cover-answer localization.} The image of the answer map over the summand $yp$ is the sieve $S$.
The full-cover oracle theorem therefore makes precisely these sieves cover,
giving $\Jdens$.
Apply \cref{thm:site-scott,thm:full-cover-oracle} to
identify the entire sheaf categories and their reflectors.
The generators
$yp$ have the same interpretations, so their generic membership agrees.
\end{enumerate}
\end{proof}
The answer map can have several witnesses. Its sheafification is epic;
no global answer selector or identification of witnesses is involved.

\subsection{Instance reducibility, modal reducibility, and query presentations}
\label{sec:oracle-reducibility}
For predicates $P:A\to\Omega_\mathcal T$ and $Q:B\to\Omega_\mathcal T$ in a topos $\mathcal T$, Bauer's \emph{instance reducibility} \cite{bauerinstance} is the internal assertion
\begin{equation}\label{eq:instance-reducibility}
 P\leq_{\mathrm I}Q\quad:\!\Longleftrightarrow\quad
 \forall a:A\ \exists b:B\;\bigl(Q(b)\Rightarrow P(a)\bigr).
\end{equation}
It is weaker than the data of a chosen one-query translator $t:A\to B$ together with the implications $Q(t(a))\Rightarrow P(a)$: the existential in \eqref{eq:instance-reducibility} need not select such a map. Modal reducibility is the inequality $o_P\leq o_Q$ in the ordered set of Lawvere--Tierney topologies.

\begin{proposition}[Reduction hierarchy]\label{prop:oracle-reduction-hierarchy}
A reduction of answer interfaces in the sense of \cref{def:oracle-reduction} induces an instance reduction between their support predicates. An instance reduction $P\leq_{\mathrm I}Q$ induces $o_P\leq o_Q$. In fact
\[
 o_P\leq o_Q\quad\Longleftrightarrow\quad
 \forall a:A\;o_Q(P(a)).
\]
Neither converse is required for these implications, even for finite query sets.
\end{proposition}
\begin{proof}
\leavevmode
\begin{enumerate}[label=\textup{(\roman*)},leftmargin=*,itemsep=3pt,topsep=4pt]
\item \textbf{Reduction hierarchy.} For an answer-interface reduction with query map $t$, an inhabitant of the translated answer fibre is sent by its given answer transformation to an inhabitant of the original fibre.
Taking regular images gives $P_{d'}(t(q))\Rightarrow P_d(q)$.
This yields the instance reduction with displayed witness $t(q)$.
\item \textbf{Existence and uniqueness.} Suppose $P\leq_{\mathrm I}Q$.
In the $o_Q$-sheaf subtopos every $Q(b)$ is true.
Given $a:A$, eliminate the existential in \eqref{eq:instance-reducibility} into the proposition $o_Q(P(a))$: the predicate inclusion $Q(b)\leq P(a)$ implies $o_Q(Q(b))\leq o_Q(P(a))$ by monotonicity, while $o_Q(Q(b))$ is true.
Thus $o_Q$ forces $P$ and leastness gives $o_P\leq o_Q$.
\item Finally, the displayed biconditional follows directly from the characterization of $o_P$ as the least modality forcing all its queried propositions.
\end{enumerate}
\end{proof}

\begin{example}[Two queries are not a single query]\label{ex:two-query-modality}
Let $S$ be the four-element poset $\{00,10,01,11\}$ ordered by
$00<10<11$ and $00<01<11$, with $10$ and $01$ incomparable. Let
$H=\Up(S)$, the frame of Alexandrov opens, and set
\[
 a=\{10,11\},\qquad b=\{01,11\},\qquad p=a\wedge b=\{11\}.
\]
In $\Sh(H)$ consider the two-query predicate $Q:2\to H$ with
$Q(0)=a$, $Q(1)=b$, and the one-query predicate $P:1\to H$ with
$P(*)=p$. For every nucleus $k$ on $H$,
\[
 k(p)=k(a\wedge b)=k(a)\wedge k(b).
\]
Thus $k$ forces $P$ exactly when it forces both instances of $Q$, and
$o_P=o_Q$. In contrast, the truth value of the internal instance-reduction
assertion $P\leq_{\mathrm I}Q$ is
\[
 (a\Rightarrow_H p)\vee(b\Rightarrow_H p)=b\vee a
      =\{10,01,11\}\ne1_H.
\]
Indeed, the Heyting implication between upper sets consists of those
$s\in S$ for which every $t\geq s$ belonging to the antecedent also
belongs to the consequent; this gives $a\Rightarrow_H p=b$ and
$b\Rightarrow_H p=a$. Since the constant query type $2=1+1$ interprets
the existential over the two queries by their join, the instance reduction
is not valid internally. This separates equality of generated modalities
from instance reduction, without confusing failure of a global query choice
with failure of the internal existential assertion.
\end{example}

\begin{proposition}[Epimorphic query-presentation invariance]\label{prop:query-cover-invariance}
Let $e:C\onto Q$ be a regular epimorphism in $\mathcal T$ and $P:Q\to\Omega_\mathcal T$ a predicate. Then
\begin{equation}\label{eq:query-cover-invariance}
 o_{P\circ e}=o_P.
\end{equation}
The same equality holds for $H$-valued relative predicates in a localic sheaf topos, under the nucleus--local-operator correspondence.
\end{proposition}
\begin{proof}
\leavevmode
\begin{enumerate}[label=\textup{(\roman*)},leftmargin=*,itemsep=3pt,topsep=4pt]
\item \textbf{Epimorphic query-presentation invariance.} For any local operator $k$, it forces $P$ exactly when $kP:Q\to\Omega$ is constantly top.
Composing with the epimorphism $e$ reflects equality of morphisms, so this is equivalent to $kPe$ being constantly top.
The two predicates have identical validating operators and hence identical least validators.
For relative predicates, run this argument in the relevant base topos for the $H$-valued family and then use \cref{thm:ab-comparison}.
\item \textbf{Oracle interpretation.} No section of $e$ is chosen, and the equality concerns the generated propositional modality.

Naturality is expressed by the commutative diagram
\[
\begin{tikzcd}[column sep=large,row sep=large]
C\ar[r,two heads,"e"]\ar[dr,"P e"']&Q\ar[d,"P"]\\ &\Omega
\end{tikzcd}
\]
\end{enumerate}
\end{proof}

\begin{remark}[An open modality is not oracle use]\label{rem:open-vs-oracle}
The open nucleus generated by the single proposition $\forall a\,P(a)$ is generally different from the least nucleus forcing each $P(a)$ separately. For example, on the Cantor locale the clopens $U_n=\{x\mid x(n)=1\}$ have empty frame meet $\bigwedge_nU_n=0$, while the locale generated by imposing all coordinate requirements $U_n$ retains the all-ones point. Treating repeated querying as implication from the open $\bigwedge_nU_n$ would give the indiscrete localization instead. The comparison is the operational reason to keep the individual indexed queries, as emphasized in \cite{ahmanbauer}.
\end{remark}

\subsection{Projective query presentations in the effective topos}
\label{sec:projective-oracle-representations}
The equality \eqref{eq:query-cover-invariance} has a direct computational interpretation in $\Eff$. There a partitioned-projective presentation $e:C\onto Q$ provides names and realizers for generalized queries. A predicate $P:Q\to\Omega$ is represented on the cover by $Pe:C\to\Omega$, and the generated propositional modality is unchanged by that replacement:
\[
 o_{Pe}=o_P.
\]
The calculation uses only regular epimorphy of the presentation and is stable under the forcing inverse-image functors used here. It does not produce a numerical enumeration of an arbitrary query object, nor does it identify the response containers. The semantic oracle and its projective names are related by \cref{sec:pres}.

In particular this presentation invariance is logically weaker than the realizability-theoretic concrete descriptions of oracle modalities in \cite[Section~5]{ahmanbauer}. Those descriptions use specific realizers for adaptive query computation. The present paper proves the abstract modal comparison for any localic forcing topos and transports actual evaluator domains and values through the comparisons of \cref{thm:partial-transport}; their conjunction does not by itself identify a general forcing sheaf topos with a relative realizability topos. For computational reducibility in $\Eff$, compare also Bauer's instance degrees and their extended Weihrauch interpretation \cite{bauerinstance}, and Kihara's description of local operators \cite{kihara}.


The contextual rule construction gives both a least modality and its operational answer interpretation. The next gluing construction retains the raw answer descriptions alongside their localized meanings, allowing subsequent comparisons to distinguish semantics from implementation.

\section{Artin gluing of computational presentations and forcing semantics}
\label{sec:artin-gluing}

Localization identifies semantic values that become equal after sheafification, but the original computational descriptions may still matter. Artin gluing places these layers in one topos. The open component records localized meaning; the closed component retains raw data; and the comparison maps specify how names, observations, and execution witnesses pass between them.


For the condition site already constructed, write
\[
 \mathcal R=\widehat P,\qquad
 \mathcal S=\Sh_\E(P,J),\qquad
 a:\mathcal R\rightleftarrows\mathcal S:i
\]
for sheafification and its fully faithful right adjoint. The unit is
$\eta$. By \cref{thm:full-cover-oracle}, this is also the contextual
cover-oracle adjunction. Artin gluing retains the source, target, and
unit of this actual adjunction in one topos.

The general gluing theorem originates in
\cite[Expos\'e~IV, Th\'eor\`eme~9.5.4]{sga4-gluing}.
We use its elementary-topos form in
\cite[Section~2.2]{faul-manuell-siqueira}; thus its use here does not
require the ground effective topos to be a Grothendieck topos over
$\mathbf{Set}$. The constructions below identify the particular
gluing, its internal frame, and its action on the retained computation
objects.

\begin{theorem}[Gluing the constructed localization]
\label{glue:thm:topos}
The comma category
\[
 \mathcal G_j=(\mathcal R\downarrow i)
\]
is a topos. Its objects are triples $(X,Y,\alpha:X\to iY)$ and its
morphisms $(f,g)$ satisfy $ig\circ\alpha=\alpha'\circ f$.
Write $r(X,Y,\alpha)=X$ and $s(X,Y,\alpha)=Y$. The functor
\[
 \Gamma_jX=(X,aX,\eta_X)
\]
is fully faithful and preserves finite limits and all colimits that
exist. There are adjunctions
\[
 \Gamma_j\dashv r\dashv c_*,\qquad s\dashv o_*,
 \qquad
 c_*X=(X,1,! ),\quad o_*Y=(iY,Y,\id).
\]
The semantic topos $\mathcal S$ is the open subtopos selected by
$o=(0,1,0\to i1)$; its complementary closed subtopos is
$\mathcal R$. In particular
\begin{equation}\label{eq:glue-modal-reconstruction}
 r\Gamma_j=\id,\qquad s\Gamma_j=a,\qquad
 o_*s\Gamma_j\cong\Gamma_jia.
\end{equation}
\end{theorem}
\begin{proof}
\leavevmode
\begin{enumerate}[label=\textup{(\roman*)},leftmargin=*,itemsep=3pt,topsep=4pt]
\item \textbf{Gluing the constructed localization.} The right adjoint $i$ preserves finite limits, so the cited elementary
Artin gluing theorem applies.
Finite limits in the comma category are
computed componentwise.
Its open and closed embeddings have precisely
the right adjoints $o_*$ and $c_*$ displayed above; their inverse
images are $s$ and $r$.
\item \textbf{Unit and identity laws.} A morphism $\Gamma_jX\to(X',Y',\alpha')$ is determined by a map
$f:X\to X'$: its second component is the unique adjunct
$aX\to Y'$ of $\alpha'f:X\to iY'$.
This proves
$\Gamma_j\dashv r$.
Since $r\Gamma_j=\id$ with identity unit,
$\Gamma_j$ is fully faithful.
Equivalently, between two objects
$\Gamma_jX,\Gamma_jX'$ the second component is forced to be $af$.
The formulas for $r\dashv c_*$ and $s\dashv o_*$ follow by the same
description of commuting squares.
\item \textbf{Inverse comparison.} Left exactness of $a$ and the
componentwise limit calculation prove left exactness of $\Gamma_j$;
its adjunction proves preservation of existing colimits.
Finally
\[
 o_*s\Gamma_jX=(iaX,aX,\id),
\]
which is canonically isomorphic to $\Gamma_jiaX$ by the counit
$aiaX\cong aX$ and the triangle identity.
This proves all three
equations in \eqref{eq:glue-modal-reconstruction}.

The defining maps satisfy the following commutative diagram
\[
\begin{tikzcd}[column sep=large,row sep=large]
\mathcal R\ar[r,"\Gamma_j"]\ar[dr,"a"']&\mathcal G_j\ar[d,"s"]\\ &\mathcal S
\end{tikzcd}
\]
\end{enumerate}
\end{proof}

The open modality in \eqref{eq:glue-modal-reconstruction} lives in
the constructed gluing topos. It is obtained from the full
sheafification adjunction and retains its indexed query semantics.
This differs from replacing the query family by its universal
conjunction inside the original topos, the operation distinguished
in \cref{rem:open-vs-oracle}.

\subsection{The glued frame over the information domain}
Put $H=\Up(P)$, $H_j=\Fix(j_J)$, and let $i_H:H_j\incl H$
be the fixed-frame inclusion. Define the internal ordered object
\[
 K_j=\{(u,v)\in H\times H_j:u\leq i_Hv\}.
\]
This is the frame form of Artin gluing
\cite[Definition~2.6]{faul-manuell-frames}. Its application here has
the following explicit structure.

\begin{theorem}[Internal frame and domain presentation of the gluing]
\label{glue:thm:frame}
The object $K_j$ is a frame, with finite meets and arbitrary joins
computed in the two component frames. The map
\[
 \gamma_j:H\longrightarrow K_j,\qquad
 u\longmapsto(u,j_Ju)
\]
is an injective frame homomorphism. With $o=(0,1)$, its open
modality satisfies
\begin{equation}\label{eq:glue-frame-modality}
 o\Rightarrow(u,v)=(i_Hv,v),\qquad
 o\Rightarrow\gamma_j(u)=\gamma_j(j_Ju).
\end{equation}
Over $\E$ there is a canonical equivalence
\[
 \mathcal G_j\simeq\Sh_\E(K_j),
\]
where the right side denotes sheaves on the locale with frame $K_j$.
Under $\Phi:H\cong\Sc(\Idl P)$ and
$\bar j=\Phi j_J\Phi^{-1}$, the componentwise map
$(u,v)\mapsto(\Phi u,\Phi v)$ identifies $K_j$ with the same
gluing constructed from the Scott frame and $\Fix(\bar j)$.
\end{theorem}
\begin{proof}
\leavevmode
\begin{enumerate}[label=\textup{(\roman*)},leftmargin=*,itemsep=3pt,topsep=4pt]
\item \textbf{Internal frame and domain presentation of the gluing.} If $u_\lambda\leq i_Hv_\lambda$, then
\[
 \bigvee_\lambda u_\lambda
 \leq i_H\left(\bigvee_\lambda^{H_j}v_\lambda\right).
\]
Finite meets remain in $K_j$ because $i_H$ preserves them.
These componentwise operations satisfy the frame distributive law.
The bottom is $(0,j_J0)$ and the top is $(1,1)$, so the argument
also covers non-dense localizations.
\item \textbf{Localization and coverage.} The nucleus laws give
\[
 j_J(u\wedge u')=j_Ju\wedge j_Ju',\qquad
 j_J\Bigl(\bigvee_\lambda u_\lambda\Bigr)
   =\bigvee_\lambda^{H_j}j_Ju_\lambda.
\]
Hence $\gamma_j$ is a frame homomorphism; its first projection
recovers $u$.
To compute implication by $o$, observe that
$(x,y)\wedge o\leq(u,v)$ is equivalent to $y\leq v$.
The greatest permitted pair is $(i_Hv,v)$, proving
\eqref{eq:glue-frame-modality} by idempotence.
\item \textbf{Adjunction and universal property.} We verify the relative sheaf assertion in each parameter context.
Let $\rho_R:\mathcal R\to\E$ and
$\rho_S:\mathcal S\to\E$ be the structure morphisms, so
$a\rho_R^*\cong\rho_S^*$.
The composite
$\Gamma_j\rho_R^*$ is left exact and left adjoint to $\rho_{R*}r$,
and therefore defines the structure morphism $\mathcal G_j\to\E$.
A subobject of $\Gamma_j\rho_R^*T$ is a pair
\[
 U\incl\rho_R^*T,\qquad V\incl\rho_S^*T,
 \qquad U\longrightarrow iV
\]
over the unit.
\item \textbf{Unit law.} The unit-pullback correspondence identifies $V$ with
a $j_J$-closed predicate $v$ in the raw context.
The last arrow
exists precisely when $U\leq v$, and is then unique.
Relative truth
classification thus identifies these subobjects with $T$-families
of elements of $K_j$.
\item \textbf{Identification of the resulting maps.} These relative subterminals generate.
For an object $(X,Y,\alpha)$,
take raw local charts $U\incl\rho_R^*T$ mapping to $X$.
Each gives
a map $\Gamma_jU\to(X,Y,\alpha)$ by adjunction.
They cover the raw
component.
Take also semantic local charts
$V\incl\rho_S^*T$ mapping to $Y$; the objects $(0,V,!)$ cover the
semantic component.
Both types are subobjects of
$\Gamma_j\rho_R^*T$.
Their combined family is epic: finite limits
and colimits in this comma category are componentwise, so a map
epic on both components is the coequalizer of its kernel pair.
\item \textbf{Existence and uniqueness.} Raw and semantic charts exist in their internal parameterized
families because both structure morphisms are localic.
The relative
localic criterion used in \cref{prop:stage-yoneda} now identifies the gluing with
$\Sh_\E(K_j)$ and its frame with the one just calculated.
\item Finally \cref{prop:scott-alex} gives $\Phi$ and its inverse.
They carry $j_J$ to $\bar j$ and its fixed points to those of
$\bar j$.
Their componentwise extensions preserve the defining
inequality of $K_j$, all frame operations, and
\eqref{eq:glue-frame-modality}.
They are inverse maps, as asserted.

The defining maps satisfy the following commutative diagram
\[
\begin{tikzcd}[column sep=large,row sep=large]
H\ar[r,"\gamma_j"]\ar[dr,"j_J"']&K_j\ar[d,"\pi_2"]\\ &H_j
\end{tikzcd}
\]
\end{enumerate}
\end{proof}

The raw and semantic projections from $K_j$ are frame homomorphisms.
The semantic component of $\gamma_j(u)$ is its saturation, while
the first component retains $u$ itself. Thus distinct raw predicates
remain distinct even when their semantic truth values agree. The
Scott-frame identification in \cref{glue:thm:frame} places this
construction over the information domain already used by the paper.

\subsection{Parameters, code objects, and evaluation relations}
\begin{proposition}[Gluing in every parameter context]
\label{glue:prop:slices}
For $Z=(X,Y,\alpha)$ define
\[
 i_\alpha:\mathcal S/Y\longrightarrow\mathcal R/X,
 \qquad (Y'\to Y)\longmapsto
       (X\times_{iY}iY'\to X).
\]
There is a canonical equivalence
\[
 \mathcal G_j/Z\simeq(\mathcal R/X\downarrow i_\alpha).
\]
For $Z=\Gamma_jT$ this is the gluing of the slice localization
$a_T:\mathcal R/T\rightleftarrows\mathcal S/aT:i_T$, where
$i_T=i_{\eta_T}$. The induced embeddings $\Gamma_{j,T}$ commute
with parameter pullback up to the canonical coherent isomorphisms.
\end{proposition}
\begin{proof}
\leavevmode
\begin{enumerate}[label=\textup{(\roman*)},leftmargin=*,itemsep=3pt,topsep=4pt]
\item \textbf{Pullback construction.} An object over $Z$ consists of $f:X'\to X$, $g:Y'\to Y$, and
$\alpha':X'\to iY'$ with $ig\alpha'=\alpha f$.
The pullback
universal property identifies this square with a unique arrow
$X'\to X\times_{iY}iY'$ over $X$.
This correspondence is inverse
on objects and morphisms and proves the equivalence.
The functor
$i_\alpha$ preserves finite limits, since it is the slice functor
induced by $i$ followed by pullback along $\alpha$.
\item \textbf{Unit and identity laws.} When $\alpha=\eta_T$, the adjunction $a_T\dashv i_T$ is obtained
by transposing the same square under $a\dashv i$.
Its right adjoint
is fully faithful by the counit.
For $T'\to T$, left exactness of
$\Gamma_j$ identifies $\Gamma_j(T'\times_T X')$ with
$\Gamma_jT'\times_{\Gamma_jT}\Gamma_jX'$.
Identity and iterated
pullback comparisons are the unique maps preserving their
projections, so their coherence equations follow from the same
universal property.
\end{enumerate}
\end{proof}

Ground code diagrams enter $\mathcal R$ by $\rho_R^*$; the next
proposition applies $\Gamma_j$ to those actual diagrams, and also
to the contextual code diagrams already formed in $\mathcal R$.

\begin{proposition}[Retention of computational presentations]
\label{glue:prop:codes}
Apply $\Gamma_j$ to any of the paper's displayed code objects,
name covers, kernel pairs, and partial domain--value spans in
$\mathcal R$. Their defining finite-limit, image, and quotient
equations are preserved; the raw projection recovers the original
diagram, and the semantic projection gives its sheafification.
The construction is coherent in parameter contexts. In particular,
a presented family $R_n\incl C_n\times X$ of raw arithmetic
predicates retains both its code index $C_n$ and its actual relation.
\end{proposition}
\begin{proof}
\leavevmode
\begin{enumerate}[label=\textup{(\roman*)},leftmargin=*,itemsep=3pt,topsep=4pt]
\item \textbf{Kernel-pair descent.} \Cref{glue:thm:topos} gives finite-limit preservation and preservation
of the kernel-pair coequalizers presenting regular epimorphisms.
Consequently it preserves regular images, name-cover equations, and
the monic domain leg of a partial map.
It also preserves the actual
quotients under discussion.
The equations $r\Gamma_j=\id$ and
$s\Gamma_j=a$ recover the two claimed diagrams.
Apply
\cref{glue:prop:slices} to their parameter pullbacks.
\item \textbf{Tracked implementation.} For the last assertion use the monomorphism
$\Gamma_jR_n\incl\Gamma_jC_n\times\Gamma_jX$; no replacement
of its specified code index is involved.
\end{enumerate}
\end{proof}


\subsection{The computational structure sheaf}
\label{sec:computational-structure}
The tracked presentations and their interpreted description algebras
form one internal many-sorted computational structure. Its realization
in the gluing retains the raw intensional maps and their oracle
interpretation. Represented denotations and the native ambient object
have separate sorts, retaining the non-epic comparison of
\cref{prop:beta-not-epi}.

\begin{definition}[A presentation and execution structure]\label{struct:def:core}
A basic presentation and execution structure in a topos consists of objects and maps
\[
 C\xrightarrow{q} V\xrightarrow{v} Z,\qquad
 D\xrightarrow{m}V\times X,\qquad e:D\longrightarrow Y,
\]
where $q$ is regular epic and $v,m$ are monic, together with the pullback
\[
\begin{tikzcd}
 E\ar[r,"p"]\ar[d,hook,"n"']&D\ar[d,hook,"m"]\\
 C\times X\ar[r,"q\times\id"']&V\times X
\end{tikzcd}
\qquad
 W\xrightarrow{w}E
\]
in which $w$ is regular epic. The value of a witness is $epw:W\to Y$.
There is also an answer interface $d:A\to Q$. Its support is retained by its image factorization $A\onto B_d\incl Q$. No epimorphy of $d$ is required in the raw structure.
\end{definition}

Here $C$ is a name or program object, $V$ its represented denotations, $Z$ an ambient value object, $D$ the exact semantic evaluation domain, $E$ its implementation domain, and $W$ a space of successful witnesses. A family of types and operations is handled by repeating this pattern in their parameter contexts, using a fixed small signature. This does not place all objects of a topos into one small internal category.

The kernel pair $C\times_V C$ records denotational equality. It is a relation on names, not their literal equality. For each displayed compiler, add its actual map on names, its semantic operation, and the commuting interpretation diagram. For example,
\[
 \widehat c:C\times C\to C,\qquad c:V\times V\to V,
 \qquad q\widehat c=c(q\times q).
\]
Associativity can hold for $c$ while the two corresponding name expressions are only kernel-pair related. Nothing in this structure equates their raw names.

For the fixed implemented signature, include the following constructed
instances among the sorts and operations of $\Comp$:
\begin{enumerate}[label=\textup{(\roman*)}]
\item exposed assembly presentations, their counits and duplication,
chosen cartesian comparisons, and the exposed maps of selected trackers;
\item description objects, constructors, substitution, tree-data
quotation, and the exact tree algebras of \cref{int:thm:exactsyntax};
\item their domain--value spans and witness covers, with the native
comparisons $\theta,\ell,\gamma$ of \cref{int:thm:transport}.
\end{enumerate}
Include their proved equations as displayed diagrams. Thus the
computational structure records the actual intensional operations;
its interpretation equations are algebra-morphism equations for
descriptions. The transported exposure instances are the diagrams
$a(\RealQ X)$ on the selected assembly presentations.
\begin{proposition}[Exact operational meaning]\label{struct:prop:exact}
\Cref{struct:def:core} presents a partial evaluator on names with exactly the pullback of the semantic domain. It satisfies, internally in every context,
\[
 (c,x)\in E\quad\Longleftrightarrow\quad
 \exists t:W\;nw(t)=(c,x).
\]
All witnesses for the same $(c,x)$ return the same element of $Y$.
\end{proposition}
\begin{proof}
\leavevmode
\begin{enumerate}[label=\textup{(\roman*)},leftmargin=*,itemsep=3pt,topsep=4pt]
\item \textbf{Exact operational meaning.} The pullback square says precisely that $(c,x)$ belongs to $E$ when $(q(c),x)$ belongs to $D$.
Since $w$ is regular epic, its image is all of $E$, proving the displayed equivalence.
If $nw(t)=nw(t')$, monicity of $n$ gives $w(t)=w(t')$, and applying $ep$ gives equality of their values.
This argument preserves the witnesses themselves; it does not identify distinct executions.
\end{enumerate}
\end{proof}

\begin{example}[The existing numerical model]\label{struct:ex:numerical}
Take $C=\N$, $V=\PC$, $q(e)=\varphi_e$, $Z=V$, and $X=Y=\N$. Use the semantic evaluation domain for $D$ and the accepted finite execution records for $W$. The name cover and its kernel pair are \cref{prop:codecover}; \cref{thm:operational-certificates} identifies the witness image with the exact implementation domain. The compiler equations of \cref{prop:compilers} complete the numerical model of \cref{struct:def:core}.
\end{example}

For arbitrary presented carriers, the same construction uses the witness covers and adequacy squares of \cref{thm:partialadequacy,cor:presented-execution-witnesses}. A witness need not be a decidably valid natural-number code for every general carrier. Its actual carrier and evidence maps remain in the structure.

\paragraph{A regular presentation of the structure}
\begin{proposition}[Syntactic presentation]\label{struct:prop:syntax}
The data in \cref{struct:def:core}, together with any fixed small family of the displayed implementation equations, admit a regular syntactic presentation $\Comp$. In a topos $\mathcal T$, its models are equivalently regular functors
\[
 \Ocomp:\mathcal C_{\Comp}\longrightarrow\mathcal T
\]
from its regular syntactic category. Equivalently these are finite-limit-preserving functors carrying its regular covers to regular epimorphisms.
\end{proposition}
\begin{proof}
\leavevmode
\begin{enumerate}[label=\textup{(\roman*)},leftmargin=*,itemsep=3pt,topsep=4pt]
\item \textbf{Coherence equations.} Introduce one sort for each displayed object and one operation for each displayed arrow.
Express monicity by the sequent $m(d)=m(d')\Rightarrow d=d'$, in its indicated variable context.
Express the cover $q$ by local existence, $\top\Rightarrow\exists c\,[q(c)=v]$, and similarly for $w$.
The square commutes by an equation.
Its universal property says that compatible pairs have a unique lift to $E$; existence is regular existential quantification and uniqueness is an equality sequent.
\item \textbf{Image and quotient construction.} An image factorization is expressed by a monic second leg and a regular-epimorphic first leg.
All compiler identities are equations.
These sequents present the diagrams already constructed above.
\item \textbf{Inverse comparison.} Objects of $\mathcal C_{\Comp}$ are regular formulas in finite contexts, modulo provable equivalence.
Arrows are provably functional relations.
Interpret a context as a finite product, conjunction as intersection, and existential quantification as regular image.
This gives a regular functor from any model.
\item \textbf{Converse construction.} Conversely, evaluate a regular functor on the named sorts, operations, and displayed formulas.
Preservation of finite limits and images recovers precisely the diagrams above.
The two procedures are inverse up to their canonical interpretation isomorphisms, also on homomorphisms.
This is the regular syntactic-category construction; see \cite[Sections 5--6]{caramello}.
\end{enumerate}
\end{proof}

This equivalence organizes the structure; by itself it is not a theorem that every abstract model of $\Comp$ is a numerical machine model. Computational content comes from the specified model in \cref{struct:ex:numerical} and its actual general-carrier and oracle constructions. Natural numbers and list objects used there are the actual ambient objects. Likewise the exposure and the initial description algebra are the constructions of \cref{int:thm:exposure,int:prop:syntax}; naming their sorts in an abstract regular model does not establish those universal properties.

Positive domain classifications can be included by their displayed pullback squares against a specified positive truth family. This does not assert that an imported truth family is the native subobject classifier or native dominance. Similarly, actual function objects and protocol objects may be named as sorts with their evaluation maps; their universal properties must be proved separately where they are used.

\paragraph{Morphisms and exact domains}
An ordinary homomorphism is a family of maps on the sorts commuting with the operations. Such a homomorphism preserves displayed membership and successful evaluation, but need not reflect definedness. An implementation comparison intended to preserve \emph{exact} partial domains is required to have a cartesian square on the relevant domain mono. This is a condition on the comparison being described, and is already proved for \cref{thm:partialadequacy,thm:partial-transport}.

\begin{lemma}[Transport of exact implementations]\label{struct:lem:strong}
Geometric inverse image carries models of $\Comp$ to models, and carries implementation comparisons with cartesian domain squares to comparisons with cartesian domain squares. These operations commute coherently with parameter pullback and composition of geometric morphisms.
\end{lemma}
\begin{proof}
\leavevmode
\begin{enumerate}[label=\textup{(\roman*)},leftmargin=*,itemsep=3pt,topsep=4pt]
\item \textbf{Transport of exact implementations.} Inverse image preserves finite limits, regular epimorphisms, and regular images.
It therefore preserves every defining diagram of the model and each specified cartesian square.
The comparison for successive inverse images is their canonical natural isomorphism; for parameter pullbacks it is the unique isomorphism preserving the pullback projections.
Uniqueness gives the composition and identity coherence equations.
\end{enumerate}
\end{proof}

\paragraph{Implementation equations modulo names}
Compiler translations that commute after denotation have a canonical groupoid interpretation.

\begin{proposition}[Coherent implementation comparisons]\label{struct:prop:pseudomaps}
A presentation $q:C\onto V$ determines an internal thin groupoid with objects $C$ and arrows $C\times_VC$. A commuting implementation square determines an internal functor. An implementation equation holding after denotation determines a unique invertible comparison between the corresponding functors on names; all coherence diagrams between such comparisons commute. These structures and comparisons are retained faithfully by $\Gamma_j$.
\end{proposition}
\begin{proof}
\leavevmode
\begin{enumerate}[label=\textup{(\roman*)},leftmargin=*,itemsep=3pt,topsep=4pt]
\item \textbf{Coherent implementation comparisons.} Source and target are the two kernel-pair projections.
Identity, inverse, and composition send $c$ to $(c,c)$, $(c,c')$ to $(c',c)$, and $((c,c'),(c',c''))$ to $(c,c'')$.
The defining equalities of $q$ make these maps well-defined, and their groupoid equations follow from their projections. A commuting square $q'\widehat f=fq$ carries the kernel pair of $q$ into that of $q'$, giving a functor.
\item \textbf{Uniqueness.} Suppose two operations on names have equal denotations in the target.
Their paired map factors through the target kernel pair and gives the comparison arrow in each input context.
Naturality and all associativity or substitution coherence diagrams follow because a thin groupoid has at most one arrow between fixed endpoints.
Invertibility is furnished by exchanging the endpoints.
\Cref{glue:thm:topos} preserves all these finite-limit diagrams, while full faithfulness recovers every actual operation and comparison map from the raw restriction.
\end{enumerate}
\end{proof}

Selected proof-name covers of a kernel pair, and their supplied witness maps, are retained as further sorts. Thinness proves coherence in the denotational relation; equations between chosen proof-name lifts are their displayed equations, not consequences of thinness downstairs.


\begin{proposition}[Descent of computational structures]
\label{struct:prop:descent}
For a regular epimorphism $p:U\onto T$, models of $\Comp$ in
$\mathcal T/T$ are equivalent to models in $\mathcal T/U$ with
isomorphism descent data satisfying the cocycle equation.
Homomorphisms and comparisons with cartesian domain squares descend
under this equivalence.
\end{proposition}
\begin{proof}
\leavevmode
\begin{enumerate}[label=\textup{(\roman*)},leftmargin=*,itemsep=3pt,topsep=4pt]
\item \textbf{Pullback construction.} Apply the effective object-and-arrow descent construction of
\cref{obs:prop:descent} to each sort and operation.
Equations are
local along $p$.
Compare a descended finite-limit diagram with its
actual limit: the comparison is an isomorphism after pullback along
$p$, hence an isomorphism.
Regular epimorphy is local because its
image is the whole target exactly when this holds after that
pullback.
Thus all regular defining diagrams and cartesian domain
squares descend.
The same argument on arrows gives full faithfulness.
\end{enumerate}
\end{proof}
Every sort of a fixed model in $\Sh_\E(K_j)$ is a relative sheaf,
with operations morphisms of sheaves. Allowing models to vary with
isomorphism descent data gives the associated stack of structures.

\paragraph{Models in the gluing.}
\begin{theorem}[Models in the gluing topos]\label{struct:thm:modelgluing}\label{int:thm:structuredgluing}
There is a canonical equivalence
\[
 \Mod_{\Comp}(\mathcal G_j)\simeq
 \bigl(a_*\downarrow\id_{\Mod_{\Comp}(\mathcal S)}\bigr),
 \qquad a_*:\Mod_{\Comp}(\mathcal R)\to\Mod_{\Comp}(\mathcal S).
\]
The right side has objects $(M,N,\theta:a_*M\to N)$ and commuting pairs of model homomorphisms. In particular,
\[
 \Gamma_{j,*}:M\longmapsto(M,a_*M,\id)
\]
is fully faithful, with $r_*\Gamma_{j,*}M=M$ and $s_*\Gamma_{j,*}M=a_*M$. The equivalence retains the displayed exposure and description
operations and respects cartesian-domain comparisons. For the native
interpretations, the comparison component on description evaluation is
\[
\begin{tikzcd}[column sep=large]
 a\Syn(LX)\ar[r,"a\alpha_X"]
 \ar[d,"{\Syn(\gamma_X)\theta_{a,LX}}"']&aLX\ar[d,"\gamma_X"]\\
 \Syn(L_\Delta X^a)\ar[r,"\alpha_{X^a}"']&L_\Delta X^a,
\end{tikzcd}
\]
where $\gamma_X=L_\Delta\beta_X\ell_{a,X}$. It compares exact
termination domains and output values and respects substitution and
common-refinement interpretation.
\end{theorem}
\begin{proof}
We construct the interpretation objectwise, verify its comparison on the retained generating data, and then establish descent and parameter coherence. At each step, the named and semantic layers are compared by their displayed interpretation maps rather than by an identification of underlying codes.

\leavevmode
\begin{enumerate}[label=\textup{(\roman*)},leftmargin=*,itemsep=3pt,topsep=4pt]
\item \textbf{Coherence equations.} A structure in $\mathcal G_j$ assigns a triple to each sort.
Each operation is a commuting pair of arrows, and its commutation condition is exactly naturality of a transformation $a_*M\to N$ after adjunction.
Finite-limit and regular-image formulas are computed componentwise by \Cref{glue:thm:topos}.
Thus the regular defining relations hold precisely when they hold in both $M$ and $N$.
This gives inverse constructions on objects and on arrows.
For structures of the form $\Gamma_{j,*}M$, the semantic component of a homomorphism is forced to be $af$ sortwise.
\item \textbf{Verification of the equations.} Full faithfulness and recovery follow.
Cartesianity is computed componentwise and is reflected by the
fully faithful finite-limit-preserving embedding $\Gamma_j$.
For the native description component, \cref{int:thm:transport} gives
the displayed algebra square and its exact domains; factor the name
maps through their represented images and retain their witness covers.
Substitution is preserved by $\theta$, and
\eqref{int:eq:refinement} compares the interpretations on common
refinements.
Their kernel-pair coherence is
\cref{struct:prop:pseudomaps}.
\item \textbf{Universal-property argument.} This constructs the intensional instances
of the model comparison.
Full faithfulness retains distinct raw
exposed maps even when the open restriction identifies them.
\end{enumerate}
\end{proof}

This theorem is the main algebraic organization result: a comparison between two computational structures is itself \emph{one} computational structure in a sheaf topos. It includes noninvertible comparison maps.

\begin{proposition}[Intrinsic recovery of both sides]\label{struct:prop:open}
Adjoin to the signature one distinguished subterminal sort $O\incl1$, giving the marked presentation $\Comp^\bullet$, and interpret it in $\mathcal G_j$ as
\[
 o=(0,1,0\to i1).
\]
Then the open restriction determined by $o$ is $\mathcal S$, and its closed complement is $\mathcal R$. Restricting $(o,\Gamma_{j,*}M)$ recovers $a_*M$ and $M$, respectively. The comparison map of each sort is recovered from the unit of the open-subtopos adjunction.
\end{proposition}
\begin{proof}
\leavevmode
\begin{enumerate}[label=\textup{(\roman*)},leftmargin=*,itemsep=3pt,topsep=4pt]
\item \textbf{Intrinsic recovery of both sides.} An object over $o$ has raw component initial and is therefore $(0,Y,!)$.
Hence $\mathcal G_j/o\simeq\mathcal S$.
Its inclusion has direct image $o_*Y=(iY,Y,\id)$ and inverse image $s$.
The complementary closed inclusion has direct image $c_*X=(X,1,!)$ and inverse image $r$, as in the gluing construction.
For $Z=(X,Y,\alpha)$, apply $r$ to the unit $Z\to o_*sZ$; the result is exactly $\alpha:X\to iY$.
Apply this sortwise and use \cref{struct:thm:modelgluing}.
Thus the chosen open, included in the one structure, determines the recovery functors and the comparison.
\end{enumerate}
\end{proof}

The distinguished open is necessary for this intrinsic recovery statement. A topos and an unmarked diagram do not by themselves specify which recollement was intended.

\begin{proposition}[Recovery of constructions from their generators]\label{struct:prop:generators}
Retaining the distinguished open, the chosen basis and approximation relations, and the answer interface suffices to recover their raw and native ideal completions and protocol monads, up to their canonical isomorphisms. It is not necessary to put a separate sort for every object of either topos into the signature. Structured equivalences preserve these recovered constructions and their universal properties.
\end{proposition}
\begin{proof}
\leavevmode
\begin{enumerate}[label=\textup{(\roman*)},leftmargin=*,itemsep=3pt,topsep=4pt]
\item \textbf{Ideal and directedness calculation.} \Cref{struct:prop:open} recovers both topoi and the indicated generator objects and relations in them.
Their ideal objects are then constructed by the actual internal ideal predicates in each recovered topos.
Rounded completions use the given approximation relation.
For any result object, the answer interface determines the polynomial whose initial algebra is the protocol object; initiality determines its return, recursion, and substitution operations uniquely.
\item \textbf{Adjunction and universal property.} The fringe functor $r o_*=i$, together with its uniquely determined left adjoint $a$, recovers the localization as well.
Thus coherence and semantic interpretation are constructed in the same recovered environments.
An equivalence of topoi preserves these universal constructions: transport along it and its quasi-inverse identifies their universal properties.
The free term monad is recovered as the finite-tree algebra of its displayed constructors, with interpretation determined by its primitive operations.
The fixed generator signature therefore suffices.
\end{enumerate}
\end{proof}


By \cref{glue:thm:frame}, this is a structure sheaf on
$\Sh_\E(K_j)$, and the Scott--Alexandrov frame isomorphism transports
the whole marked structure to its information-domain presentation.
The extra subterminal sort is interpreted as $o$ after gluing; it is
not claimed to be the image under $\Gamma_j$ of a raw subterminal.
Write $\Ocomp$ for the resulting marked computational structure and
its explicitly displayed extensions below. Each extension adds its
constructed sorts and commuting diagrams to the same regular
presentation.

\paragraph{Native completions and represented images.}

For an approximation presentation $C$ in $\mathcal R$, \cref{prop:base-approximation} constructs the native completion $D_{a^\dagger C}$ and a comparison
\[
 \beta_C:aD_C\longrightarrow D_{a^\dagger C}.
\]
Its adjunct defines an object of $\mathcal G_j$:
\[
 \mathbf D_C=(D_C,D_{a^\dagger C},\widetilde\beta_C).
\]
For a presented Scott-continuous map $F:D_C\to D_{C'}$, the proved equation
\[
 \beta_{C'}aF=F^a\beta_C
\]
makes $(F,F^a)$ a morphism $\mathbf D_C\to\mathbf D_{C'}$. Its identity and composition equations give functoriality. Thus native completions fit the structure directly, without making $\beta_C$ invertible or epic.

\begin{proposition}[A name cover has a represented image]\label{struct:prop:image}
For any map $\mathbf n:\mathbf C\to\mathbf Z$ in $\mathcal G_j$, its factorization
\[
 \mathbf C\onto\mathbf V\incl\mathbf Z
\]
is computed by the image factorizations on the raw and semantic components. In particular, old names for $D_C$ cover their native represented image, which need not exhaust $D_{a^\dagger C}$.
\end{proposition}
\begin{proof}
\leavevmode
\begin{enumerate}[label=\textup{(\roman*)},leftmargin=*,itemsep=3pt,topsep=4pt]
\item \textbf{A name cover has a represented image.} Use componentwise finite limits and kernel-pair coequalizers from \Cref{glue:thm:topos}.
The component comparison between the images is induced by the commutative square for $\mathbf n$ and the image universal property.
The resulting first leg is regular epic and second leg monic on both components, hence in $\mathcal G_j$.
The final assertion specializes this calculation to the old-name map followed by $\beta_C$.
\end{enumerate}
\end{proof}

This explains the separate sorts $V$ and $Z$ in \cref{struct:def:core}. Requiring names to cover every native ambient value would contradict \cref{prop:beta-not-epi}. Factoring through the represented image constructs the correct cover while retaining the entire ambient native object.

The kernel pair of a map into a native ambient object must likewise be computed for that map. It is not automatically the transported old kernel pair: a further comparison can identify additional denotations. Both relations and their induced map can be retained in the same structure.

For a partial evaluator, retain its raw and native domain--value spans and the comparison between them. The exact domain comparisons of \cref{thm:partial-transport} ensure that the two domain monos and the compatible value maps define a partial-operation span in $\mathcal G_j$. Its exactness follows componentwise. No totality of a forcing observation operation outside its constructed admissibility domain is used.


Positive observations enter by the constructed map
$a\Sigma_{\mathcal R}\onto\Theta_a\incl\Sigma_{\mathcal S}$ of
\cref{prop:theta}; this is an additional comparison sort with its actual image and comparison to the native positive truth object. Exact admissibility and quantifier domains enter by
their defining pullbacks. For protocols the required higher-order
comparison and its execution equations are proved in
\cref{struct:thm:Wcomparison,struct:thm:transcript}.

\paragraph{Morphisms of computationally structured topoi.}
For the fixed computational signature, a structured morphism is
\[
 (f,f^\sharp):(\mathcal T,\Ocomp)\longrightarrow
                  (\mathcal T',\Ocomp'),\qquad
 f^\sharp:f^*\Ocomp'\longrightarrow\Ocomp.
\]
The structure map is a model homomorphism. Exact implementation
morphisms have cartesian domain squares; marked morphisms carry the
pulled-back distinguished open isomorphically to the distinguished
open. Composition has structure map
$f^\sharp\circ f^*g^\sharp$, with the canonical inverse-image
coherence. \Cref{struct:lem:strong} proves stability of the domain
conditions, and associativity follows from that coherence.

This is an ordinary relative structured topos. In the specialized
geometry of \cite[Definition~1.2.8]{lurie}, local structure maps are
cartesian on every admissible arrow. Declaring every name cover
admissible would therefore impose cartesian squares on those covers;
a general compiler does not induce the required fibre isomorphisms.
The computational homomorphisms above retain their stated domain
conditions. The construction uses the actual relative elementary
topoi and does not replace them by infinity-topoi.

Gluing's use in organizing computational predicates through an
internal language also appears in synthetic Tait computability
\cite[Sections~2.4 and~4.5]{sterling-tait}. Here the structure retains
the displayed name, execution, and native comparison diagrams.

\subsection{A marked structure for three semantic layers}
\label{sec:flagged-gluing}
For a requirement nucleus $j_0\leq\Dens$ put
$\mathcal R=\widehat P$, $\mathcal S_0=\Sh_\E(P,J_0)$, and
$\mathcal S_d=\Sh_\E(P,\Jdens)$. Write the successive adjunctions as
\[
 a_0:\mathcal R\rightleftarrows\mathcal S_0:i_0,
 \qquad b:\mathcal S_0\rightleftarrows\mathcal S_d:i_1.
\]
The following gluing retains the requirement layer when its dense forcing
localization is introduced.

\begin{theorem}[Flagged gluing]\label{cal:thm:gluing}
The category $\mathcal G$ whose objects are
\[
 (X_0,X_1,X_2;\ \alpha_{01}:X_0\to i_0X_1,
                         \ \alpha_{12}:X_1\to i_1X_2)
\]
is a topos. The functor
\[
 \Gamma X=(X,a_0X,ba_0X;\eta^0_X,\eta^1_{a_0X})
\]
is fully faithful and preserves finite limits and colimits. The nested
subterminals
\[
 o_2=(0,0,1)\ \subseteq\ o_1=(0,1,1)\ \subseteq\ 1
\]
recover, respectively, the dense layer, the two-layer semantic gluing, and the
raw complement as follows:
\[
 \mathcal G/o_2\simeq\mathcal S_d,
 \quad \mathcal G/o_1\simeq(\mathcal S_0\downarrow i_1),
 \quad \mathcal G\setminus o_1\simeq\mathcal R.
\]
Inside $\mathcal G/o_1$, the closed complement of $o_2$ is $\mathcal S_0$.
The two comparison maps are recovered from the corresponding open-inclusion
units.
Writing $H_0=\Fix(j_0)$ and $H_d=\Fix(\Dens)$, the internal frame
\[
 K_{0,d}=\{(u,v,w)\in H\times H_0\times H_d:u\leq v\leq w\}
\]
has componentwise finite meets and joins in the respective fixed frames,
and $\mathcal G\simeq\Sh_\E(K_{0,d})$. Its retained truth embedding is
$u\mapsto(u,j_0u,\Dens u)$.
\end{theorem}
\begin{proof}
\leavevmode
\begin{enumerate}[label=\textup{(\roman*)},leftmargin=*,itemsep=3pt,topsep=4pt]
\item \textbf{Flagged gluing.} Let $\mathcal G_{12}=(\mathcal S_0\downarrow i_1)$.
Its raw projection
$r_{12}$ is right adjoint to the fully faithful finite-limit-preserving functor
$\Gamma_{12}:Y\mapsto(Y,bY,\eta^1_Y)$.
Therefore
$i_0r_{12}$ is right adjoint to the finite-limit-preserving functor
$\Gamma_{12}a_0$.
Artin gluing gives a topos
$(\mathcal R\downarrow i_0r_{12})$, whose objects are exactly the displayed
triples.
\item \textbf{Verification of the equations.} Equivalently, write the arrows as $a_0X_0\to X_1$ and $bX_1\to X_2$.
Finite limits and colimits are then computed componentwise, since $a_0,b$
preserve them. A map between two objects $\Gamma X$ is forced by its raw
component $f$: its other components are $a_0f$ and $ba_0f$.
This proves full
faithfulness and the preservation assertions.
\item \textbf{Universal-property argument.} An object over $o_1$ has $X_0=0$, and an object over $o_2$ has
$X_0=X_1=0$, proving the open identifications.
The closed inclusion of the raw
component sends $X$ to $(X,1,1)$.
The analogous closed inclusion in
$\mathcal G_{12}$ sends $Y$ to $(Y,1)$.
\item Finally, for a gluing object
$(X,Y,\alpha)$, the raw component of its map to the direct image of its open
restriction is exactly $\alpha$.
Apply this observation at the two successive
gluings.
\item \textbf{Identification of the resulting maps.} For the frame statement, the inequalities defining $K_{0,d}$ are preserved
by componentwise finite meets and fixed-frame joins.
For example,
$j_0(\bigvee v_\lambda)\leq\Dens(\bigvee w_\lambda)$ follows from
$v_\lambda\leq w_\lambda$ and $j_0\Dens=\Dens$.
Frame distributivity
is therefore componentwise.
In each base context, subobjects of the
retained terminal family are exactly triples of predicates $u\leq v\leq w$,
with the latter two closed at their respective stages.
\item \textbf{Localization and coverage.} The relative local
charts in \cref{glue:thm:frame}, applied successively to the two gluings,
generate all three components.
This identifies the relative sheaf topos
with $\Sh_\E(K_{0,d})$.
The nucleus laws make the displayed retained truth
map an injective frame homomorphism.
\end{enumerate}
\end{proof}

The regular signature is the presentation-and-execution structure of \cref{struct:def:core,struct:prop:syntax}.

\begin{corollary}[A single algebraic presentation]\label{cal:cor:structure}
Adjoining the flag $o_2\hookrightarrow o_1\hookrightarrow1$ gives a single
regular signature. The retained raw computational structure has a model in
$\mathcal G$ via $\Gamma$. Its three restrictions recover the raw structure
and its two successive reflections. More generally, a model of this signature
in $\mathcal G$ with this distinguished flag is precisely a chain of regular
models and homomorphisms
\[
 a_0M_0\longrightarrow M_1,\qquad bM_1\longrightarrow M_2,
\]
together with the marked flag. Cartesian domain comparisons remain cartesian
under restriction and geometric inverse image.
\end{corollary}
\begin{proof}
\leavevmode
\begin{enumerate}[label=\textup{(\roman*)},leftmargin=*,itemsep=3pt,topsep=4pt]
\item \textbf{A single algebraic presentation.} Interpret each sort as a gluing triple.
Operation commutation is precisely the
homomorphism condition on the two arrows.
Finite limits and regular images are
componentwise, so the defining regular diagrams hold exactly componentwise.
The same calculation applies to each domain pullback square.
Theorem~\ref{cal:thm:gluing}
recovers the marked layers and their arrows.
\end{enumerate}
\end{proof}

The native ideal completions, numerical function objects, and protocol initial
algebras are reconstructed in the recovered layers from the retained generators,
as in \cref{struct:prop:generators}. Their actual comparison maps are retained. Numerical adequacy for the
designated model is established in \cref{cal:thm:dialogues,cal:thm:adequacy}.


Gluing provides a single setting in which raw names, requirement-level predicates, and localized values can be compared without identifying them. The tower calculations that follow compose these comparisons, including their individual effective restrictions.

\section{Relative localizations and coherent towers}\label{sec:towers}

Oracle constructions are frequently iterated, and each stage can admit a different collection of effective observations. We compare the localization maps in a tower and the maximal effective restrictions of their composites. Coherence is expressed by the actual pullback and nucleus comparisons, rather than by treating successive mark selections as interchangeable.


\begin{theorem}[Nuclei above a fixed layer]\label{thm:tower}
For a nucleus $j$ on $H$ there is an order isomorphism
\[
 \Nuc(\Omega_j)\cong\{k\in\Nuc(H)\mid j\leq k\},
 \qquad n\longmapsto i_jnq_j.
\]
Its fixed-frame identifications, locale inclusions, and individually marked stable fragments are coherent for finite towers.
\end{theorem}
\begin{proof}
\leavevmode
\begin{enumerate}[label=\textup{(\roman*)},leftmargin=*,itemsep=3pt,topsep=4pt]
\item \textbf{Nuclei above a fixed layer.} Set $k=i_jnq_j$.
It is inflationary because $x\leq jx\leq i_jnq_jx$, and it preserves finite meets.
The equation $q_ji_j=\id$ gives
$k^2=i_jn^2q_j=k$.
Clearly $k\geq j$.
\item \textbf{Converse construction.} Conversely, if $j\leq k$, monotonicity and idempotence imply
\[
 kj=k=jk.
\]
Indeed $kx\leq kjx\leq k^2x$ and $kx\leq jkx\leq k^2x$.
Therefore $k$ sends fixed points of $j$ to fixed points of $j$, and its restriction defines $n$.
The two assignments are inverse and order preserving.
\item \textbf{Pullback construction.} An $n$-fixed element of $\Omega_j$ is exactly a $k$-fixed element of $H$.
Their inclusions in $H$ agree.
For a marking $M\incl H$, the sequential stable marking is
\[
 \Fix(n)\times_{\Omega_j}(\Omega_j\times_HM)
 \cong\Fix(k)\times_HM.
\]
Pullback pasting supplies this isomorphism.
For three or more layers, every comparison preserves the final truth and its original marked representative.
Uniqueness in the pullback gives equality of all composites.
\item \textbf{Image and quotient construction.} The frame inclusions and quotient maps have the same absorption equations, which gives their coherence and that of the opposite locale diagrams.
\end{enumerate}
\end{proof}

\begin{proposition}[Refined site]\label{prop:refined-site}
Let $d_p=[p]_J$ be the original stage basis and let $n$ be a nucleus on $\Omega_J(P)$. Put $d'_p=q_n(d_p)$. The topology
\[
 S\in J_n(p)\iff d'_p=\bigvee_{q\in S}^{\Fix(n)}d'_q
\]
extends $J$, presents $\Fix(n)$, and has ambient saturation $i_Jnq_J$.
\end{proposition}
\begin{proof}
\leavevmode
\begin{enumerate}[label=\textup{(\roman*)},leftmargin=*,itemsep=3pt,topsep=4pt]
\item \textbf{Image and quotient construction.} The quotient preserves the basis top, meet, and density equations.
Apply \cref{thm:basis}.
Every old cover equation remains valid after $q_n$, so $J\subseteq J_n$.
The fixed points of the new ambient nucleus are exactly the old stable predicates whose frame elements are fixed by $n$.
This is the fixed object of $i_Jnq_J$, and nuclei with the same fixed points agree: each sends $x$ to the meet of fixed points above $x$.
Hence the saturation formula follows.
\end{enumerate}
\end{proof}

\begin{corollary}[Union of requirements]\label{cor:requirements-join}
For two rule families on the same frame,
$j_{\mathcal R\cup\mathcal S}=j_\mathcal R\vee j_\mathcal S$.
The same statement for relative rules is obtained by the tower isomorphism.
\end{corollary}
\begin{proof}
\leavevmode
\begin{enumerate}[label=\textup{(\roman*)},leftmargin=*,itemsep=3pt,topsep=4pt]
\item \textbf{Union of requirements.} A nucleus validates the union precisely when it validates both families.
Validity is upward closed: if $k\leq h$, then $hk=h$, so applying $h$ to $k(a)\leq k(b)$ gives $h(a)\leq h(b)$.
Thus the validators of the union are exactly the common upper bounds of the two least validators.
\end{enumerate}
\end{proof}

The join here is a join in the ordered set of nuclei. The ordinary composite of two nuclei need not be idempotent. The least-validator construction supplies the join with its correct universal property.

For $j\leq k$, the quotient map $\Omega_j\to\Omega_k$ need not carry every member of $S_j$ into $S_k$. Its maximal effective restriction is \cref{thm:lax}. In contrast, an individually effective $k$-fixed truth is also $j$-fixed, and its inclusion $S_k\incl S_j$ is always defined. These are different arrows: one saturates a truth further, while the other regards an already fixed truth in the earlier frame. The pullback comparison in \cref{thm:tower} concerns the second operation and does not assert totality of the first.

\section{Decision interfaces and Booleanization}\label{sec:boolean}

A decision query is a special answer interface with two complementary candidate answers. Making its support locally true does not automatically make the original predicate computationally decidable in the ground category. We compute the generated decision localization and compare the resulting forcing truth with the corresponding semidecision domain.


For a predicate $S\incl Q$, form the Boolean answer relation
\[
 E(q,0)\iff S(q),\qquad E(q,1)\iff\neg S(q).
\]
The two answers are disjoint, and its support is $S(q)\vee\neg S(q)$.

\begin{proposition}[Dense decision reconstruction]\label{prop:decision}
In the current topos let $j_S$ be the least oracle localization of
the displayed interface $d_S:E_S\incl Q$. Then $j_S\leq\neg\neg$,
and $a_{j_S}d_S$ is an isomorphism. Its inverse determines the unique
Boolean decision of the sheafified predicate $a_{j_S}S$.
An indexed family is treated by its coproduct answer interface.
\end{proposition}
\begin{proof}
\leavevmode
\begin{enumerate}[label=\textup{(\roman*)},leftmargin=*,itemsep=3pt,topsep=4pt]
\item \textbf{Localization and coverage.} Disjointness makes $d_S$ monic, with image $S\vee\neg S$.
Double negation sends this image to top: its negation would imply
both $\neg S$ and $\neg\neg S$.
Thus double negation inverts
$d_S$, and minimality of the oracle localization gives
$j_S\leq\neg\neg$.
In particular $j_S0=0$, proving density.
Sheafification preserves the mono and makes its image all of
$a_{j_S}Q$, so it is also epic and hence invertible.
The two answer fibres are disjoint and cover $a_{j_S}Q$; consequently
the transported negative fibre is the actual complement of
$a_{j_S}S$.
\item \textbf{Oracle interpretation.} Their Boolean characteristic map is unique.
For a family, apply the same argument in the joint query context;
\cref{prop:modal-reductions} supplies the coproduct localization.
\end{enumerate}
\end{proof}

\begin{theorem}[Double negation comparison]\label{thm:double-negation}
Under \cref{prop:nucleus-topology}, the internal frame nucleus
$b(u)=\neg\neg u$ corresponds to the double-negation Lawvere--Tierney topology of $\Sh_\E(L)$.
\end{theorem}
\begin{proof}
\leavevmode
\begin{enumerate}[label=\textup{(\roman*)},leftmargin=*,itemsep=3pt,topsep=4pt]
\item \textbf{Double negation comparison.} A nucleus is dense when $k0=0$.
This is equivalent, under the relative truth correspondence, to its closure fixing the empty subterminal.
Thus the nucleus--topology correspondence preserves and reflects density.
\item \textbf{Localization and coverage.} Double negation is a dense nucleus in any Heyting frame.
If $k$ is another dense nucleus, then
\[
 k(u)\wedge\neg u
 \leq k(u)\wedge k(\neg u)
 =k(u\wedge\neg u)=k0=0.
\]
Consequently $k(u)\leq\neg\neg u$.
Thus double negation is the greatest dense nucleus.
The same argument in the subobject classifier gives the greatest dense local operator.
The order isomorphism preserving density identifies them.
\end{enumerate}
\end{proof}

This theorem compares frame and topos double negation in the same forcing
universe. The extension-density nucleus $\Dens$ of
\cref{cal:thm:density} is below it and can be strictly smaller in a
non-Boolean base. Its equality with double negation belongs to the classical
calibration. We next calculate an effective obstruction in the Cohen universe.

\begin{theorem}[Cohen Booleanization is not globally effective]\label{thm:boolean-obstruction}
In $\Eff$, double negation on the structural Cohen frame does not restrict to its full effective fixed fragment.
\end{theorem}
\begin{proof}
\leavevmode
\begin{enumerate}[label=\textup{(\roman*)},leftmargin=*,itemsep=3pt,topsep=4pt]
\item \textbf{Cohen Booleanization is not globally effective.} Use the external enumeration $(W_e)$ of c.e.\ sets to describe a uniform tracked family.
Put $D_n=[0^n1]$.
These pairwise disjoint cylinders have dense union: every string containing a $1$ already lies in the corresponding cylinder, and an all-zero string can be extended by $1$. A finite initial union $\bigvee_{n<k}D_n$ misses $[0^k]$.
\item \textbf{Verification of the equations.} Enumerate $D_n$ into $U_e$ when $W_e$ has acquired $n+1$ distinct elements.
Membership of a finite string in this union is uniformly semidecidable: find its first $1$, say at $n$, and wait for $n+1$ distinct enumerated values.
The predicate is upper.
It is stable under binary splitting, since if the string already has a first $1$ both children have the same first $1$, while an all-zero string has an all-zero child.
Hence $U_e(\sigma)\iff U_e(\sigma0)\wedge U_e(\sigma1)$, which by the saturation formula makes it $j_C$-fixed.
This constructs a morphism $\N\to S_{j_C}$.
\item \textbf{Forcing construction.} If $W_e$ is infinite, $U_e=\bigvee_nD_n$ is dense, so $\neg\neg U_e=1$.
If it has exactly $k$ members, $U_e=\bigvee_{n<k}D_n$, and the nonzero cylinder $[0^k]$ lies in its negation.
Thus
\[
 \neg\neg U_e=1\iff W_e\text{ is infinite}
\]
in the external calibration.
\item \textbf{Evaluation and transposition.} A uniform effective restriction of double negation, followed by evaluation at the empty string, would therefore give a morphism $\N\to\Sigma$ whose external support is the set of indices of infinite c.e.\ sets. A map from the natural numbers object to $\Sigma$ has a uniformly c.e.\ support: pull back the finite-hit epimorphism $2^\N\onto\Sigma$ and use the partitioned-assembly projectivity of $\N$, as in \cref{thm:finitehit}, to obtain a tracked binary stream family; search for a hit.
The infinitude index set is not c.e.
\item \textbf{Identification of the resulting maps.} For completeness, reduce nonhalting to it: enumerate one new integer at each stage until the given machine halts.
The resulting c.e.\ set is infinite exactly when the machine never halts.
If nonhalting were c.e., dovetailing its enumeration with the halting enumeration would decide halting, contradicting diagonalization.
This is the required contradiction.
\end{enumerate}
\end{proof}

The example also shows why an effective family of generating assertions is not itself a construction of an effective saturation procedure. The support determines which localization is required; the effective action of that localization is a separate calculation governed by Part~\ref{part:effective}.


\section{Intensional enhancement: oracle protocols, equifoliate trees, and names}\label{sec:intensional-oracles}

A modality says which propositions become locally true, but not how a computation adaptively obtains the answers it needs. We construct the finite and well-founded interaction objects that supply this missing operational meaning. Their kernel-pair quotient gives the semantic sheafification, while the retained protocols preserve the finite transcripts and the choices on which subsequent queries depend.


\subsection{Finite interaction as a dependent construction}
For completeness, a fixed finite bound on interaction admits an explicit transcript construction in each forcing stage. Fix an answer map $d:A\to Q$ and a result object $Y$. Define, for each finite bound, protocol objects by
\[
 P_0=Y,\qquad
 P_{n+1}=Y+\sum_{q:Q}P_n^{A_q}.
\]
The dependent exponential exists in the slice over $Q$. A protocol is either an immediate result or a query together with a continuation for each possible answer. No decision of the answer support is built into this definition.

Define an object $R_n\to P_n$ of successful transcripts and a value map $R_n\to Y$ by recursion. At level zero, $R_0=Y$ over $P_0=Y$ with identity value. Over a terminal protocol $y$ at the next level the transcript fibre is $1$ with value $y$. Over a query protocol $(q,k)$ its fibre is
\[
 R_{n+1}(q,k)=\sum_{a:A_q}R_n(k(a)),
\]
and its value is the value of the selected continuation transcript. This formula constructs the object by dependent sum and pullback.

\begin{proposition}[Finite protocol adequacy]\label{prop:finite-protocol}
Membership in the support of $R_n\to P_n$ is exactly successful execution of the protocol with some valid finite answer transcript. If $d$ is epic, every protocol has a transcript locally, so $R_n\to P_n$ is epic. These statements are preserved by any sheafification.
\end{proposition}
\begin{proof}
\leavevmode
\begin{enumerate}[label=\textup{(\roman*)},leftmargin=*,itemsep=3pt,topsep=4pt]
\item \textbf{Finite protocol adequacy.} Induct on $n$.
At a terminal protocol the claim is witnessed by the unique empty transcript.
At a query protocol, an element is exactly a valid answer $a:A_q$ and a successful transcript for $k(a)$, by the displayed dependent-sum formula.
This is the execution clause defined by recursion, proving the first assertion.
\item \textbf{Pullback construction.} If $d$ is epic, every $q$ locally has an answer.
By the induction hypothesis, the selected continuation locally has a transcript.
The composite local witnesses give an element of $R_{n+1}(q,k)$; hence its support is top.
Equivalently, the transcript projection is the composite of the pullback of a regular epimorphism for the continuation with the answer projection, so it is epic. A sheafification preserves these finite limits, dependent sums, and regular images.
\item \textbf{Evaluation and transposition.} The pulled-back transcript construction for the given protocol and continuations therefore has the same support and value equations.
No assertion that sheafification preserves the entire protocol exponential is required.
\end{enumerate}
\end{proof}

When answers are functional and their partial supports are $\Sigma$-classified, the execution relation is a partial map with the iterated dependent extent of the transcript. Dominance composes these extents. When answers are genuinely multivalued, $R_n$ retains that multivaluedness, and one does not replace it by a function without constructing a selection. This is the finite operational content of the answer interface that its support alone would forget.

\subsection{Bounded protocols and modal computation}
\label{sec:finite-modal-protocols}
The objects $P_n$ above have bounded query depth, but their branching fibres $A_q$ need not be finite. The distinction between bounded depth and a finite \emph{number} of possible oracle answers matters when comparing these objects with the computation trees of \cite{ahmanbauer}. The next results provide a finite-depth semantic comparison in the original topos, without identifying $P_n$ with a quotient-inductive or equifoliate-tree object.

\begin{lemma}[Finite-depth fold and coherence]\label{lem:protocol-fold}
For any object $Z$ with an operation $\alpha:Y+\sum_{q:Q}Z^{A_q}\to Z$, recursion defines maps $h_n:P_n\to Z$ satisfying
\[
 h_0(y)=\alpha(\mathrm{inl}(y)),\qquad
 h_{n+1}(\mathrm{inl}(y))=\alpha(\mathrm{inl}(y)),\quad
 h_{n+1}(q,k)=\alpha\bigl(q,h_n\circ k\bigr).
\]
There are canonical inclusions $\iota_n:P_n\to P_{n+1}$ obtained by retaining a result and applying $\iota_{n-1}$ to each query continuation, and $h_{n+1}\iota_n=h_n$ for every $n$.
\end{lemma}
\begin{proof}
\leavevmode
\begin{enumerate}[label=\textup{(\roman*)},leftmargin=*,itemsep=3pt,topsep=4pt]
\item \textbf{Unit and identity laws.} The dependent polynomial sends a map $h:P_n\to Z$ to the family of postcomposition maps $P_n^{A_q}\to Z^{A_q}$, natural in $q$.
Its coproduct with the identity on $Y$, followed by $\alpha$, defines $h_{n+1}$ from $h_n$.
The same polynomial construction with $\iota_n$ defines $\iota_{n+1}$.
The base equality is $h_1(\mathrm{inl}(y))=h_0(y)$.
At a query node the induction hypothesis gives $(h_{n+1}\iota_n)\circ k=h_n\circ k$, hence the defining equations for $\alpha$ give the successor equality.
Coproduct elimination proves the equation for all protocols.
No recursion over arbitrary unbounded well-founded trees is used.
\end{enumerate}
\end{proof}

Fix the least oracle modality $j=o_d$ of the answer-support predicate of $d:A\to Q$, and let $a_j:\mathcal T\to\Sh_j(\mathcal T)$ be the associated sheafification. The effect of the modality is that $a_jd$ is epic. Write $v_n:R_n\to Y$ for the value of a successful transcript and put
\begin{equation}\label{eq:finite-modal-coherence}
 E_n=\bigl\{p:P_n\mid
   \forall r,r'\in R_n(p),\ j\bigl(v_n(r)=_Yv_n(r')\bigr)\bigr\}.
\end{equation}
This defines a subobject by universal quantification in $\mathcal T$. It says that a protocol returns at most one result \emph{up to $j$-equality}; it does not assert that the response fibre has a chosen element.


\subsection{Well-founded protocols and their coherent modality}
\label{sec:well-founded-protocols}
The finite-depth objects are useful for numerical bounds, but the general
adaptive interface does not require one. In $\Eff$ and its forcing sheaf topoi,
polynomial W-types exist without an additional inductive-type axiom:
every elementary topos with a natural numbers object has W-types
\cite{moerdijkpalmgren}. Apply this existence theorem separately in each ambient topos to the polynomial associated with $d:A\to Q$; its polynomial formulation is developed in \cite{gambinokock}. Thus the native W-type is constructed in the target topos itself. The concrete site description of \cite{vandenbergmoerdijk} corrects the earlier presheaf-tree description and explicitly works in ZFC with choice. We do not use that concrete description as the existence proof here: the general topos theorem supplies the objects, and evaluation and initiality supply the comparison below.

\begin{definition}[Well-founded oracle protocol]\label{def:well-founded-protocol}
For an object $Y$, let $W_d(Y)$ be the initial algebra for
\[
  \mathcal P_Y(Z)=Y+\sum_{q:Q}Z^{A_q}.
\]
Its constructors are $\operatorname{return}(y)$ and
$\operatorname{ask}(q,k)$ with $k:A_q\to W_d(Y)$. Initiality gives
recursion and induction, uniformly in $Y$, $Q$, and $d$. The finite-depth
protocols of \cref{sec:finite-modal-protocols} have canonical maps
$\iota_n^W:P_n\to W_d(Y)$, defined by recursion on $n$ and the
constructors.
\end{definition}


\paragraph{Comparison with native protocols.}
Write $a:\mathcal R\rightleftarrows\mathcal S:i$ for a left exact
localization, in particular the forcing or oracle localization above.
For the polynomial of \cref{def:well-founded-protocol}, put
\[
 P_Y(Z)=Y+\sum_{q:Q}Z^{A_q},\qquad
 P'_{aY}(Z')=aY+\sum_{q':aQ}(Z')^{(aA)_{q'}}.
\]
Write $W=W_d(Y)$ and $W'=W_{ad}(aY)$ for their actual initial
algebras, constructed separately by the elementary-topos existence theorem \cite{moerdijkpalmgren}. The following comparison is built from evaluation, adjunction, and initiality, and does not require preservation of exponentials.
For $f:Y\to W_d(Z)$, the Kleisli extension $f^\star:W_d(Y)\to W_d(Z)$
is substitution, defined by the fold equations
\[
 f^\star(\operatorname{return}(y))=f(y),\qquad
 f^\star(\operatorname{ask}(q,k))
   =\operatorname{ask}(q,f^\star\circ k).
\]

\begin{lemma}[Polynomial comparison]\label{struct:lem:polynomial}
There is a natural map
\[
 \lambda_Z:aP_Y(Z)\longrightarrow P'_{aY}(aZ)
\]
preserving the result summand and sending an imported continuation to the continuation obtained by applying $a$ to its evaluation map. These comparisons commute with successive inverse images and parameter pullback.
\end{lemma}
\begin{proof}
\leavevmode
\begin{enumerate}[label=\textup{(\roman*)},leftmargin=*,itemsep=3pt,topsep=4pt]
\item \textbf{Unit and identity laws.} Pull back the evaluation map over $Q$ to obtain
\[
 a\left(\sum_{q:Q}Z^{A_q}\right)\times_{aQ}aA\longrightarrow aZ.
\]
Transpose in the slice over $aQ$ to obtain the query-summand comparison.
Combine it with the identity of $aY$ using preservation of coproducts.
Naturality follows by postcomposition of evaluation.
Two iterated comparison maps have the same evaluation, so exponential adjunction identifies them.
The same argument in slices proves parameter coherence.
The comparison retains its specified direction.
\end{enumerate}
\end{proof}

\begin{theorem}[Canonical native protocol comparison]\label{struct:thm:Wcomparison}
There is a unique map $\kappa_Y:aW_d(Y)\to W_{ad}(aY)$ satisfying
\[
 \kappa_Y\,a(\operatorname{sup})
 =\operatorname{sup}'\,P'_{aY}(\kappa_Y)\lambda_W.
\]
It preserves returns and oracle queries with their imported continuations. It is natural under relabelling results and, for $f:Y\to W_d(Z)$, satisfies
\[
 \kappa_Z\,a(f^\star)
   =(\kappa_Z\,af)^\star\kappa_Y.
\]
\end{theorem}
\begin{proof}
\leavevmode
\begin{enumerate}[label=\textup{(\roman*)},leftmargin=*,itemsep=3pt,topsep=4pt]
\item \textbf{Native polynomial algebra.} Give $iW'$ a $P_Y$-algebra structure by taking the adjunct of the composite
\[
 aP_Y(iW')\xrightarrow{\lambda}P'_{aY}(aiW')
 \xrightarrow{P'(\epsilon)}P'_{aY}(W')
 \xrightarrow{\operatorname{sup}'}W',
\]
where $\epsilon:ai\cong1$.
Initiality of $W$ yields a unique algebra map $W\to iW'$.
Its adjunct is $\kappa_Y$.
Transposing the algebra equation proves the displayed identity; the same transposition proves uniqueness. Explicitly, the constructor equation is the commutative diagram
\[
\begin{tikzcd}[column sep=large,row sep=large]
 aP_Y(W)\ar[r,"\lambda_W"]\ar[d,"{a(\operatorname{sup})}"']
  &P'_{aY}(aW)\ar[r,"{P'_{aY}(\kappa_Y)}"]&P'_{aY}(W')\ar[d,"\operatorname{sup}'"]\\
 aW\ar[rr,"\kappa_Y"']&&W'.
\end{tikzcd}
\]
\item \textbf{Initiality and comparison map.} For result relabelling, the two proposed composites satisfy the same recursion equations, so initiality identifies them.
For substitution, the compared map on return labels is $\kappa_Zaf:aY\to W_{ad}(aZ)$.
The two maps obtained by substituting before or after comparison agree on returns.
At an oracle node their continuations agree by the induction hypothesis and \cref{struct:lem:polynomial}. W-induction in $\mathcal R$, followed by adjunction, proves the substitution equation.
\item \textbf{Constructor coherence.} This concerns a supplied continuation map, and does not identify the whole imported function object with its native counterpart.
\end{enumerate}
\end{proof}

The comparison is also coherent for successive localizations: the direct comparison and the composite of the two comparisons satisfy the same constructor equation, by the coherence of $\lambda$, so uniqueness identifies them. In parameter slices the same proof applies to the slice adjunction. Reindexing between slices preserves the relevant polynomial constructions and their initial algebras, since slice pullback has a right adjoint and satisfies dependent-product base change. The comparison is consequently stable under parameter substitution.


\begin{lemma}[Successful-list execution]\label{lem:well-founded-transcripts}
Let $\operatorname{List}(A)=\sum_{n:\N}A^{\Fin(n)}$. There is a
subobject $\operatorname{Exec}_d\incl W_d(Y)\times\operatorname{List}(A)
\times Y$ determined by
\begin{align*}
 \operatorname{Exec}_d(w,[],y)
   &\iff w=\operatorname{return}(y),\\
 \operatorname{Exec}_d(w,a::\ell,y)
   &\iff\exists q,k\,[w=\operatorname{ask}(q,k)\wedge d(a)=q
                  \wedge\operatorname{Exec}_d(k(a),\ell,y)].
\end{align*}
Write $R_d(Y)$ for this subobject and $r:R_d(Y)\to W_d(Y)$,
$v:R_d(Y)\to Y$ for its projections. Their fibres satisfy natural
isomorphisms
\[
 R_d(Y)_{\operatorname{return}(y)}\cong1,
 \qquad
 R_d(Y)_{\operatorname{ask}(q,k)}
       \cong\sum_{a:A_q}R_d(Y)_{k(a)},
\]
and the isomorphisms preserve terminal values.
\end{lemma}
\begin{proof}
\leavevmode
\begin{enumerate}[label=\textup{(\roman*)},leftmargin=*,itemsep=3pt,topsep=4pt]
\item \textbf{Finite witness construction.} Finite lists are constructed from the NNO and finite dependent powers.
Recursion on list length in the power object of $W_d(Y)\times Y$
constructs the displayed relation; it does not require deciding equality
or selecting a query answer.
At a return node only the empty list
qualifies.
At a query node a successful nonempty list has a uniquely
specified head $a$, satisfying $d(a)=q$, and its tail is a successful
list for $k(a)$.
Removing or prepending that head gives the two inverse
maps of the second fibre isomorphism.
The last output coordinate is
unchanged.
The construction is internal in every parameter context.
\end{enumerate}
\end{proof}


\begin{theorem}[Finite transcript comparison]
\label{struct:thm:transcript}
The map $\kappa_Y:aW_d(Y)\to W_{ad}(aY)$ induces an isomorphism
\[
 aR_d(Y)\ \cong\
 aW_d(Y)\times_{W_{ad}(aY)}R_{ad}(aY),
\]
preserving the answer list and returned value. Thus finite successful
execution of an imported tree is preserved and reflected by its
native interpretation. Its support is the pullback of the native
successful-execution support along $\kappa_Y$.
\end{theorem}
\begin{proof}
\leavevmode
\begin{enumerate}[label=\textup{(\roman*)},leftmargin=*,itemsep=3pt,topsep=4pt]
\item \textbf{Finite transcript comparison.} Finite-list objects and their constructors are preserved by $a$: their polynomial $1+A\times-$ commutes with $a$, and initiality transfers by the same adjoint argument as above with an invertible polynomial comparison.
Thus $a\operatorname{List}(A)\cong\operatorname{List}(aA)$.
\item \textbf{Verification of the equations.} The same argument preserves the NNO.
We compare the execution subobjects at each list length by induction on that NNO.
Lambek's lemma gives the disjoint decompositions
\[
 aW\cong aY+a\Bigl(\sum_qW^{A_q}\Bigr),\qquad
 W'\cong aY+\sum_{q'}(W')^{(aA)_{q'}}.
\]
The defining equation for $\kappa_Y$ preserves the two tags.
At length zero, an imported return labeled $y_0$ maps to the native return labeled $y_0$, and a query cannot map to a return.
\item Therefore native execution with empty list and output $y$ holds exactly when the imported tree is $a\operatorname{return}(y)$.
This is the image under $a$ of the raw length-zero relation.
\item \textbf{Evaluation and transposition.} For the successor step, work on the query summand of the first decomposition, with imported continuation data $k$.
The native root has the same query label.
An answer $b$ is valid on either side exactly when $(ad)(b)$ equals that label.
The comparison of continuation evaluation gives
\[
  k_{\mathsf{native}}(b)=\kappa_Y(k_{\mathsf{raw}}(b)),
\]
where $k_{\mathsf{raw}}(b)\in aW$ is obtained from the pulled-back evaluation map, not from a claimed preservation of exponentials.
\item \textbf{Induction.} By induction, native success of the tail from $k_{\mathsf{native}}(b)$ is equivalent to the transported raw success relation from $k_{\mathsf{raw}}(b)$.
The return summand has no nonempty successful transcript on either side.
These are exactly the successor execution clauses, since $a$ preserves their finite limits and existential images.
They identify the finite-list and output coordinates together with their classified supports.
\item \textbf{Pullback construction.} The length cases give equality of the subobjects of $aW\times\operatorname{List}(aA)\times aY$ in question, yielding the displayed isomorphism.
Finally, regular images commute with pullback in a topos, proving the support statement.
\end{enumerate}
\end{proof}

The theorem permits arbitrarily branching, unbounded well-founded protocols. Only each successful transcript is finite. The induction ranges over individually finite successful transcripts.

\begin{corollary}[A native result object]\label{struct:cor:nativeresult}
For a constructed map $b:aY\to Z$ in $\mathcal S$, the composite
\[
 aW_d(Y)\xrightarrow{\kappa_Y}W_{ad}(aY)
       \xrightarrow{W_{ad}(b)}W_{ad}(Z)
\]
induces an isomorphism
\[
 aR_d(Y)\cong aW_d(Y)\times_{W_{ad}(Z)}R_{ad}(Z),
\]
preserving answer lists and transforming returned values by $b$.
In particular this applies when $b$ is a non-epic completion comparison.
\end{corollary}
\begin{proof}
\leavevmode
\begin{enumerate}[label=\textup{(\roman*)},leftmargin=*,itemsep=3pt,topsep=4pt]
\item \textbf{Oracle interpretation.} Relabeling return leaves changes neither query nodes nor their answer lists.
Induction on list length identifies the transcript object of a tree with that of its relabelling, preserving the list and applying $b$ to the final value.
At length zero the value is the given leaf label; at a successor length the first answer and continuation are unchanged.
Combine this identification with \cref{struct:thm:transcript}.
Neither injectivity nor surjectivity of $b$ is used.
\end{enumerate}
\end{proof}


For any object $Z$ write $\eta_Z:Z\to a_jZ$ for the associated-sheaf
unit. The equality predicate of the image of $\eta_Z$ is the $j$-closure
of the diagonal:
\begin{equation}\label{eq:sheaf-unit-equality}
 \eta_Z(z)=\eta_Z(z')\quad\Longleftrightarrow\quad j(z=z').
\end{equation}
This is the separated-reflection calculation for a
Lawvere--Tierney topology \cite{johnstone,maclanemoerdijk}; it applies
in every object context.

\begin{theorem}[Well-founded modal transcript and output theorem]
\label{thm:well-founded-modal-adequacy}
Let $j=o_d$ be the least modality making $d:A\to Q$ epic, and define
\[
 E_d(Y)=\left\{w:W_d(Y)\;\middle|\;
   \forall r,r'\in R_d(Y)_w,\quad
   j\bigl(v(r)=v(r')\bigr)\right\}.
\]
Then the following maps are constructed without any branching-finiteness
or depth bound.
\begin{enumerate}
\item The transcript projection $a_jR_d(Y)\to a_jW_d(Y)$ is a regular
 epimorphism; so is its pullback over $a_jE_d(Y)$.
\item The transcript values descend uniquely to
\[
  \operatorname{value}_d^Y:a_jE_d(Y)\longrightarrow a_jY,
\]
with $\operatorname{value}_d^Y\,a_jr=a_jv$ after restriction of the
transcript object to $E_d(Y)$.
\item The assignments $Y\mapsto E_d(Y)$ form a submonad of the
well-founded oracle-tree monad $W_d$. The modal output maps are
compatible with its return and substitution operations. Pullback along the finite-depth constructor maps gives the bounded transcript and output constructions.
\end{enumerate}
\end{theorem}
\begin{proof}
\leavevmode
\begin{enumerate}[label=\textup{(\roman*)},leftmargin=*,itemsep=3pt,topsep=4pt]
\item \textbf{Well-founded induction.} Let $s(w)=\exists r\in R_d(Y)_w\,\top$.
We prove $j(s(w))=\top$ by
W-induction. A return node has the empty transcript.
At a query node
$w=\operatorname{ask}(q,k)$, the induction hypothesis says
$j(s(k(a)))=\top$ for every $a:A_q$.
By the second fibre
isomorphism of \cref{lem:well-founded-transcripts},
$s(k(a))\leq s(w)$ in the context $a:A_q$.
Thus
$P_d(q)\leq j(s(w))$, by existential elimination, where
$P_d(q)=\exists a:A_q\,\top$.
Since $j(P_d(q))=\top$ and $j$ is
idempotent, $j(s(w))=\top$.
This completes the induction.
\item \textbf{Pullback construction.} Factoring
$r$ through its support mono and applying sheafification proves the
first regular-epimorphism assertion.
Pullback along $a_jE_d(Y)$
preserves that regular epimorphism, and left exactness identifies the
pulled-back object with the sheafification of the restricted
transcript relation.
\item \textbf{Kernel-pair descent.} For two transcripts with the same protocol in $E_d(Y)$, the definition
of that subobject and \eqref{eq:sheaf-unit-equality} make their output
values equal after sheafification.
Left exactness transfers this
equality to the kernel pair of the sheafified regular epimorphism.
The latter is the coequalizer of its kernel pair, so the sheafified
value descends to one and only one $\operatorname{value}_d^Y$.
\item For the final assertion, the substitution fold preceding
\cref{struct:thm:Wcomparison} gives the multiplication of $W_d$;
return is its unit. W-induction proves the unit and associativity
equations, also obtained from the polynomial free-monad theorem
\cite{gambinokock}.
Relabelling returns by $f:Y\to Z$
preserves coherence because $j(y=y')$ implies $j(fy=f y')$.
Returns themselves have one transcript and are coherent.
\item \textbf{Identification of the resulting maps.} For substitution let $f:Y\to E_d(Z)$ and $w\in E_d(Y)$.
W-recursion on $w$ proves the transcript decomposition
\[
 R_{\operatorname{bind}(w,f)}
       \cong\sum_{r:R_w}R_{f(v(r))},
\]
with the terminal value given by the inner transcript.
Define
$\delta_Z=\operatorname{value}_d^Z\eta_{E_d(Z)}$ and
$g=\delta_Zf:Y\to a_jZ$, using the already established output
map for result type $Z$.
Consider two completed transcripts with
outer terminal values $y,y'$ and inner terminal values $z,z'$.
Coherence of $w$ gives $j(y=y')$.
\item \textbf{Monad compatibility.} Since $g$ takes values in the
$j$-sheaf $a_jZ$, its equality is $j$-closed and therefore
$g(y)=g(y')$.
The two inner descent equations give
$g(y)=\eta_Z(z)$ and $g(y')=\eta_Z(z')$.
Hence
$\eta_Z(z)=\eta_Z(z')$, or equivalently $j(z=z')$ by
\eqref{eq:sheaf-unit-equality}.
Thus the substituted tree is
coherent.
Return is coherent; the unit and associativity laws
restrict from the W-type free monad because $E_d$ is a subobject.
For compatibility, let $g^\star:a_jY\to a_jZ$ be the unique sheaf extension of $g=\delta_Zf$.
\item \textbf{Forcing construction.} On an outer transcript returning $y$, an inner transcript returning $z$ gives $\eta_Z(z)=g(y)$.
The inner transcript projection is $j$-dense by the first assertion, and equality in $a_jZ$ is $j$-closed.
Thus the value of the substituted tree equals $g(y)$ already in the outer transcript context.
There $g^\star\delta_Y(w)=g^\star\eta_Y(y)=g(y)$ as well.
\item \textbf{Naturality and reindexing.} Density of the outer transcript projection and separatedness now give
\[
 \delta_Z(\operatorname{bind}(w,f))=g^\star\delta_Y(w).
\]
The return equation is $\delta_Y(\operatorname{return}(y))=\eta_Y(y)$, and relabelling values gives naturality by the same descent uniqueness.
These prove compatibility with the monad operations.
Restriction along $\iota_n^W$ has the finite-depth transcript fibres by induction on $n$, giving the bounded construction.
\end{enumerate}
\end{proof}

\begin{theorem}[Finite-depth modal adequacy]\label{thm:finite-modal-adequacy}
For every $n$ the sheafified successful-transcript projection
\[
 a_j(R_n|_{E_n})\longrightarrow a_jE_n
\]
is a regular epimorphism. Moreover, there is a unique map of $j$-sheaves
\[
 \operatorname{value}_n:a_jE_n\longrightarrow a_jY
\]
whose composite with that epimorphism is the sheafification of the transcript-value map. These value maps are coherent with the inclusions $\iota_n$ and with the evident inclusions $E_n\to P_n$. In particular the resulting modal output is independent of the chosen finite successful oracle transcript.
\end{theorem}
\begin{proof}
\leavevmode
\begin{enumerate}[label=\textup{(\roman*)},leftmargin=*,itemsep=3pt,topsep=4pt]
\item \textbf{Well-founded induction.} Use the constructor map $\iota_n^W:P_n\to W_d(Y)$ of \cref{def:well-founded-protocol}.
Induction on $n$, using the return and query fibre isomorphisms of \cref{lem:well-founded-transcripts}, identifies its pulled-back transcript object with $R_n$, preserving the value map.
Hence its inverse image of $E_d(Y)$ is exactly $E_n$.
Pull back the regular epimorphism and its descended value from \cref{thm:well-founded-modal-adequacy}; left exactness of $a_j$ gives the asserted maps.
\item \textbf{Identity law.} The identity $\iota_{n+1}^W\iota_n=\iota_n^W$ follows from the fold equations of \cref{lem:protocol-fold}.
Both induced values agree on the transcript cover, so uniqueness of descent proves their coherence.
\end{enumerate}
\end{proof}

\begin{remark}[Exact recovery of oracle sheafification]\label{rem:equifoliate-boundary}
The subsequent comparison identifies the coherent W-type submonad
with Ahman--Bauer's equifoliate trees, and constructs its natural
map and modal kernel-pair quotient into ordinary sheafification.
The latter is identified by its conditional quotient-inductive
universal property in \cref{thm:oracle-qit-recovery}.
Neither construction erases the intensional tree before its
semantic quotient is taken.
\end{remark}

\subsection{Equifoliate trees as the coherent-protocol subobject}
\label{sec:intensional-equifoliate}
Retain the actual answer fibres of $d:A\to Q$. Their support
predicate determines $j=o_d$, but it does not determine the
continuations or the transcript object. For $x:Y$ and
$w:W_d(Y)$ define the full modal membership predicate
\[
 x\in_jw\quad:\!\Longleftrightarrow\quad
 \forall r:R_d(Y)_w\ j(x=v(r)).
\]
The transcript decompositions of \cref{lem:well-founded-transcripts}
give, uniformly in all parameters,
\begin{equation}\label{eq:intensional-membership-recursion}
 x\in_j\operatorname{return}(y)\iff j(x=y),\qquad
 x\in_j\operatorname{ask}(q,k)
 \iff\forall a:A_q\,x\in_jk(a).
\end{equation}

\begin{lemma}[Modal transcript coherence and branch transfer]
\label{lem:equifoliate-transcript-bridge}
For every $w:W_d(Y)$ and $x:Y$ let
\[
 \begin{aligned}
 s(w)&=\exists r:R_w\,\top,\\
 C(w)&=\forall r,r':R_w\ j(v(r)=v(r')),\\
 M(x,w)&=\forall r:R_w\ j(x=v(r)).
 \end{aligned}
\]
Every tree has $j(s(w))=\top$. Moreover $M$ has the recursion
\[
 M(x,\operatorname{return}(y))\iff j(x=y),\qquad
 M(x,\operatorname{ask}(q,k))
    \iff\forall a:A_q\ M(x,k(a)),
\]
and coherence satisfies the precise equivalence
\begin{align*}
 C(\operatorname{return}(y))&\iff\top,\\
 C(\operatorname{ask}(q,k))&\iff
   \bigl(\forall a:A_q\ C(k(a))\bigr)
   \ \wedge\\[-2pt]
 &\hspace{1.2em}
   \bigl(\forall a,b:A_q\ \forall x:Y\,
           [M(x,k(a))\Rightarrow M(x,k(b))]\bigr).
\end{align*}
No answer, successful transcript, or global branch is chosen.
\end{lemma}
\begin{proof}
\leavevmode
\begin{enumerate}[label=\textup{(\roman*)},leftmargin=*,itemsep=3pt,topsep=4pt]
\item \textbf{Modal transcript coherence and branch transfer.} The transcript-support assertion is the W-induction proved in
\cref{thm:well-founded-modal-adequacy}.
The two equations for
$M$ follow immediately from the return and query transcript-fibre
isomorphisms of \cref{lem:well-founded-transcripts}.
\item \textbf{Oracle interpretation.} Suppose the parent query tree is coherent.
Two transcripts of one
child, preceded by the same answer, are transcripts of the parent,
so each child is coherent.
To verify branch transfer, fix $a,b,x$
and suppose $M(x,k(a))$.
For every $r_b:R_{k(b)}$ and every
$r_a:R_{k(a)}$, transitivity of $j$-equality gives
\[
 j(x=v(r_a))\wedge j(v(r_a)=v(r_b))
       \Longrightarrow j(x=v(r_b)).
\]
The conclusion is $j$-closed.
Since
$j(\exists r_a:R_{k(a)}\,\top)=\top$, existential elimination
followed by $j$-closure establishes the conclusion without
selecting $r_a$.
Quantifying over $r_b$ proves
$M(x,k(b))$.
\item \textbf{Converse construction.} Conversely suppose the displayed child-coherence and transfer
predicates.
Given parent transcripts $(a,r_a)$ and $(b,r_b)$,
child coherence implies $M(v(r_a),k(a))$ by taking every other
transcript of that child.
Transfer gives $M(v(r_a),k(b))$;
evaluation at $r_b$ yields $j(v(r_a)=v(r_b))$.
This proves parent
coherence.
The return case has its unique empty transcript, so it
is coherent.
Every equation holds in its full parameter context.
\end{enumerate}
\end{proof}

\begin{theorem}[Exact equifoliate identification]
\label{thm:intensional-equifoliate}
The object $E_d(Y)$ of \cref{thm:well-founded-modal-adequacy} is
canonically the Ahman--Bauer equifoliate-tree subobject of $W_d(Y)$.
Equivalently, its coherence predicate is generated inductively by
\begin{enumerate}
\item every return is equifoliate;
\item $\operatorname{ask}(q,k)$ is equifoliate if every $k(a)$ is,
and for all $a,b:A_q$ and $x:Y$ one has
$(x\in_j k(a))\Rightarrow(x\in_j k(b))$.
\end{enumerate}
The identification is the identity on underlying trees, not a
quotient of their query histories.
\end{theorem}
\begin{proof}
\leavevmode
\begin{enumerate}[label=\textup{(\roman*)},leftmargin=*,itemsep=3pt,topsep=4pt]
\item \textbf{Oracle interpretation.} The membership $x\in_jw$ is $M(x,w)$ from
\cref{lem:equifoliate-transcript-bridge}, whose leaf and query
recursion is exactly the membership recursion of
\cite[Section~4.2, Definition~33]{ahmanbauer}.
The same bridge proves that the
transcript-coherence predicate $C(w)$ has the return clause and
precisely the child-coherence and branch-transfer query clause of
\cite[Definition~34]{ahmanbauer}. W-induction therefore gives
$C(w)\iff\operatorname{equi}(w)$ for every tree, uniformly in its
parameters.
\item \textbf{Image and quotient construction.} These are two inclusions of subobjects of the same
unquotiented W-type, hence they identify $E_d(Y)$ with the
Ahman--Bauer equifoliate subobject.
No quotient or excluded-middle
decision of a query support is used.
\end{enumerate}
\end{proof}

\subsection{Exact modal quotient and substitution}
\label{sec:intensional-modal-quotient}
Let $\eta_Y:Y\to a_jY$ be the sheafification unit. The map
$\operatorname{value}_d^Y$ of
\cref{thm:well-founded-modal-adequacy} determines
\[
 \delta_Y:E_d(Y)\longrightarrow a_jY,\qquad
 \delta_Y=\operatorname{value}_d^Y\eta_{E_d(Y)}.
\]
Here and below $a_jY$ is also viewed as an object of the original
topos via the full inclusion of the sheaf subtopos.

\begin{theorem}[Transcript interpretation and exact modal equality]
\label{thm:intensional-kernel-pair}
For every coherent tree $w$ and every successful transcript $r:R_w$,
\[
 \delta_Y(w)=\eta_Y(v(r)),\qquad
 x\in_jw\iff\eta_Y(x)=\delta_Y(w).
\]
The kernel pair of $\delta_Y$ is the subobject
\begin{equation}\label{eq:intensional-modal-equivalence}
 w\sim_jw'\iff
 \forall r:R_w\ \forall r':R_{w'}\ j(v(r)=v(r')).
\end{equation}
After sheafification the induced map is a regular epimorphism and
presents $a_jY$ as an effective quotient:
\[
 a_jY\cong\coeq_{\Sh_j(\mathcal T)}
 \left(a_j(\sim_j)\rightrightarrows a_jE_d(Y)\right).
\]
In particular every element of $a_jY$ is locally the interpretation
of an actual equifoliate protocol.
\end{theorem}
\begin{proof}
\leavevmode
\begin{enumerate}[label=\textup{(\roman*)},leftmargin=*,itemsep=3pt,topsep=4pt]
\item \textbf{Transcript interpretation and exact modal equality.} The first equation is the descent equation in
\cref{thm:well-founded-modal-adequacy}.
If $x\in_jw$,
the sheafified transcript values are constantly $\eta_Y(x)$.
They descend along the epic sheafified transcript projection;
uniqueness gives $\delta_Y(w)=\eta_Y(x)$.
The converse follows
by evaluating the first equation for each transcript and using
\eqref{eq:sheaf-unit-equality}.
\item \textbf{Inverse comparison.} If $\delta_Y(w)=\delta_Y(w')$, the first equation and
\eqref{eq:sheaf-unit-equality} give every equation in
\eqref{eq:intensional-modal-equivalence}.
Conversely, under that
relation, the two sheafified terminal values agree on the product
of the two successful-transcript objects.
Both transcript supports
are $j$-dense by W-induction; the product support is also dense
by meet preservation of $j$.
Thus
$j(\delta_Y(w)=\delta_Y(w'))=\top$.
Equality in the sheaf
$a_jY$ is $j$-closed, so the interpretations agree.
This is the
exact kernel pair, including empty or locally inhabited answer fibres.
\item \textbf{Forcing construction.} The image of $\delta_Y$ contains the image of $\eta_Y$, since
$\delta_Y\operatorname{return}=\eta_Y$.
Sheafification of $\eta_Y$
is invertible, hence its image is $j$-dense; consequently the image
of $\delta_Y$ is dense and $a_j\delta_Y$ is a regular epimorphism.
Left exactness of $a_j$ carries the calculated kernel pair to
the kernel pair of this regular epimorphism.
Regular epimorphisms are
coequalizers of their kernel pairs in a topos, giving the displayed
quotient and its universal property.
\end{enumerate}
\end{proof}

Under \cref{thm:intensional-equifoliate}, this is the comparison of
\cite[Theorems~38--39]{ahmanbauer}. Its surjectivity is modal:
the dense image does not select ground-level tree representatives.
The monad structure compared next is \cite[Theorem~37]{ahmanbauer}.

\begin{theorem}[The equifoliate monad maps to the sheafification monad]
\label{thm:intensional-monad-morphism}
The W-type return and substitution operations restrict to $E_d$,
and $\delta:E_d\Rightarrow a_j$ respects unit and Kleisli
substitution. In particular, if $f:Y\to E_d(Z)$ and
$w:E_d(Y)$, then
\begin{equation}\label{eq:intensional-bind-equation}
 \delta_Z(\operatorname{bind}_{E}(w,f))
 = (\delta_Z f)^\star(\delta_Yw),
\end{equation}
where $(-)^\star$ is Kleisli extension for the sheafification monad, equivalently unique extension into the $j$-sheaf $a_jZ$.
\end{theorem}
\begin{proof}
\leavevmode
\begin{enumerate}[label=\textup{(\roman*)},leftmargin=*,itemsep=3pt,topsep=4pt]
\item \textbf{Naturality and reindexing.} The return and substitution operations, their closure in $E_d$, and the displayed equation were proved in \cref{thm:well-founded-modal-adequacy} from transcript decomposition and descent into the sheaf $a_jZ$.
Its return equation and naturality give a natural monad morphism.
The identification with the equifoliate construction is \cref{thm:intensional-equifoliate}, so these are its unit and Kleisli-substitution equations.
\end{enumerate}
\end{proof}


\begin{corollary}[Compatibility with native coherent execution]
\label{struct:cor:coherent}
For $a=a_j$ and $j=o_d$, the restriction of $\kappa_Y$ to
$a_jE_d(Y)$ lands in the native subobject of protocols whose successful
transcript values agree. Native evaluation on this image equals
$a_j\delta_Y:a_jE_d(Y)\onto a_jY$. Return and substitution retain
their semantic equations under this comparison.
\end{corollary}
\begin{proof}
\leavevmode
\begin{enumerate}[label=\textup{(\roman*)},leftmargin=*,itemsep=3pt,topsep=4pt]
\item \textbf{Kernel-pair descent.} After sheafification, the two value maps on the kernel pair of the
transcript projection over $a_jE_d(Y)$ are equal by modal coherence.
\Cref{struct:thm:transcript} identifies this with the kernel pair of
the pulled-back native transcript projection.
Hence the compared
protocol is coherent.
Since $a_jd$ is epic, W-induction gives a
successful transcript locally for every native tree: returns have the
empty transcript, and a query locally has an answer followed by a
transcript supplied by induction.
Coherent values therefore descend
along this epimorphism.
\item \textbf{Monad compatibility.} They agree with $a_j\delta_Y$ on the compared
transcript cover, so the descended maps coincide.
Return and
substitution follow from
\cref{struct:thm:Wcomparison,thm:intensional-monad-morphism}.
\end{enumerate}
\end{proof}

The triples $(W_d(Y),W_{ad}(aY),\widetilde\kappa_Y)$, continuation
sorts, transcripts, and constructors form compatible structure data
in $\mathcal G_j$. Their higher-order universal properties are
recovered on the raw and native restrictions by
\cref{struct:prop:generators}. This does not assert that those triples
are the W-types or exponentials formed internally in $\mathcal G_j$,
or that $a_jE_d(Y)$ exhausts the native coherent subobject.

\subsection{Denotationally coherent oracle names}
\label{sec:intensional-oracle-codes}
Let $q:C\onto Y$ be any regular-epimorphic name presentation as in
\cref{sec:pres}. Its names need not have decidable equality, and distinct
names can denote one semantic element. Define
\[
 \operatorname{OCode}_{d}(q)
 =\left\{w:W_d(C)\ \middle|\
   \forall r,r'\in R_w\ j(qv(r)=qv(r'))\right\}.
\]
Equivalently this is the inverse image of $E_d(Y)$ under the
return-relabelling map $W_d(q):W_d(C)\to W_d(Y)$.

\begin{theorem}[Oracle-relative name presentation]
\label{thm:intensional-oracle-coding}
There is a canonical regular-epimorphic interpretation of sheaves
\[
 q_d:a_j\operatorname{OCode}_{d}(q)\onto a_jY
\]
whose value at a successful transcript $r$ is $\eta_Y(qv(r))$.
Let $R_{d,q}\incl\operatorname{OCode}_{d}(q)^2$ denote the relation
\[
 w\equiv_{d,q}w'\iff
 \forall r:R_w\ \forall r':R_{w'}\ j(qv(r)=qv(r')).
\]
Then its sheafification is the kernel pair of $q_d$:
$ a_jR_{d,q}\cong a_j\operatorname{OCode}_{d}(q)
\times_{a_jY}a_j\operatorname{OCode}_{d}(q)$.
A commuting implementation square $q'\widehat f=fq$ induces a
commuting square of these oracle name presentations, by relabelling
returned names with $\widehat f$. The construction retains the
unquotiented query and answer transcripts on its name side.
\end{theorem}
\begin{proof}
\leavevmode
\begin{enumerate}[label=\textup{(\roman*)},leftmargin=*,itemsep=3pt,topsep=4pt]
\item \textbf{Oracle interpretation.} W-recursion relabelling returns by $q$ preserves query nodes and
transcripts.
It carries the displayed coherence condition to that
for $E_d(Y)$, hence restricts to a map
$\operatorname{OCode}_{d}(q)\to E_d(Y)$.
Compose its sheafification
with $a_j\delta_Y$ to obtain $q_d$.
The ordinary returned-name
inclusion $C\to\operatorname{OCode}_{d}(q)$ satisfies
$q_d a_j\operatorname{return}=a_jq$.
Since $q$ is regular epic and
sheafification preserves such epis, $a_jq$ is epic.
Consequently
$q_d$ is epic.
\item \textbf{Kernel-pair descent.} The pre-sheafification equality relation is the pullback of the
kernel pair in \cref{thm:intensional-kernel-pair} under return
relabelling by $q$.
Since $a_j$ is left exact, sheafifying this
pullback relation gives the kernel pair of $q_d$.
It records equality
of actual sheafified denotations, not equality of returned names.
\item Finally $q'\widehat f=fq$ transfers every
modal equality after denotation along $f$.
The induced return
relabelling therefore preserves the code object, and the
transcript-value equation proves commutativity of its interpretation
square.
No functorial choice of a projective cover is required.
\end{enumerate}
\end{proof}

The quotient in \cref{thm:intensional-oracle-coding} presents
$a_jY$, the sheafification of the \emph{supplied} denotation object.
It does not assert that every native forced completion object is a
pullback of $Y$. The precise old-to-native representability comparison
is constructed in \cref{thm:intensional-old-image}.


\subsection{Native finite execution}
\begin{proposition}[Native oracle execution checker]\label{prop:intensional-native-checker}
For a constructed total Boolean oracle operation $g:Q\to2$ in $\FF$, there is a total Boolean checker for finite numerical execution records with supplied $Q$-labels at oracle-read nodes. It checks the ordinary transitions and, at a node labelled $(q,b)$, the equation $g(q)=b$. Its accepted-run object is compatible with renaming, parameter substitution, and the checked wrapper constructions of \cref{exec:thm:codefeedback}. For numerically coded queries, including the Cohen coordinate oracle, the termination support is classified by the native finite-hit object $\Sigma_\FF$.
\end{proposition}
\begin{proof}
\leavevmode
\begin{enumerate}[label=\textup{(\roman*)},leftmargin=*,itemsep=3pt,topsep=4pt]
\item \textbf{Finite witness construction.} Represent a record by its finite numerical shape and the finite list of query labels used at its read nodes. A shape specifies each lookup position, the recorded Boolean, and the ensuing numerical state.
Structural recursion checks the shape and each ordinary transition; at a read node it applies $g$ to the supplied label and uses decidable equality on $2$.
Only finitely many such operations are performed, so the checker is total.
It does not test equality of arbitrary $Q$-labels or the validity of arbitrary program-name presentations.
\item \textbf{Coherence equations.} Reindexing the supplied labels commutes with lookup.
The read case of each wrapper equation evaluates the same supplied query with the same $g$; the ordinary cases are the finite-run equations of \cref{thm:operational-certificates,exec:thm:codefeedback}.
Induction on the finite record proves the equations.
\item \textbf{Oracle interpretation.} When the query labels are numerical, the complete records have numerical encodings.
Enumerate the codes, reject malformed ones, and use the checker result as the next bit of a binary stream.
This is a total internal stream, and its finite-hit truth is exactly existence of an accepted terminating record.
Hence its support lies in $\Sigma_\FF$ directly by the definition of that object.
For general $Q$ the labels remain supplied data; existential search over them is treated by the exact positive quantifier domains of \cref{thm:maximal-positive-quantifiers}.
\item \textbf{Finite witness construction.} For a general answer relation $A\to Q$, its finite transcript object and answer evidence are retained as in \cref{sec:intensional-oracles}; the Boolean checker for a total Boolean operation does not replace those answer covers.
\end{enumerate}
\end{proof}


\begin{theorem}[An internal oracle partial combinatory algebra]\label{struct:thm:pca}
For the constructed total Boolean oracle $g:\N\to2$, define
\[
 u\mathbin{\cdot_g}x\simeq U^g(u,x)
\]
by the numerical oracle interpreter of \cref{prop:intensional-native-checker}. This makes the native numerical object an internal partial combinatory algebra, with native semidecidable application domain. Its combinators and specialization operations are obtained by fixed numerical compiler constructions.
\end{theorem}
\begin{proof}
\leavevmode
\begin{enumerate}[label=\textup{(\roman*)},leftmargin=*,itemsep=3pt,topsep=4pt]
\item \textbf{An internal oracle partial combinatory algebra.} Determinism of the finite execution rules makes the application relation functional.
Its domain is the finite-hit predicate of \cref{prop:intensional-native-checker}.
The code generator taking $a$ to a program that returns the literal $a$ is total and numerical.
Let $k$ be a program implementing that generator.
Then $k\cdot_g a$ is defined and $(k\cdot_g a)\cdot_g b=a$ for every $a,b$.
\item \textbf{Tracked implementation.} For $s$, compile the following delayed operation.
On input $a$, return the code of a program that, on input $b$, returns the code of a program waiting for $c$.
Only when $c$ is supplied does it execute $a$ on $c$, execute $b$ on $c$, and execute the first returned program on the second returned value.
The first two stages write finite syntax and do not run the supplied programs.
Consequently $s\cdot_g a$ and $(s\cdot_g a)\cdot_g b$ are always defined.
\item \textbf{Partiality and definedness.} The last stage gives
\[
 ((s\cdot_g a)\cdot_g b)\cdot_g c
 \simeq (a\cdot_g c)\cdot_g(b\cdot_g c),
\]
including equality of definedness: the left run terminates exactly when the two preliminary runs and their final application all terminate.
In either direction their finite records are combined or extracted by the checked wrappers of \cref{thm:operational-certificates,exec:thm:codefeedback,prop:intensional-native-checker}.
All code-generation operations are primitive recursive on syntax and are independent of $g$; each oracle read is delegated to the supplied $g$.
\item \textbf{Functoriality and coherence.} This proves the internal combinatory equations uniformly.
The literal-insertion wrapper gives specialization with the same domain-and-value correctness.
\end{enumerate}
\end{proof}

This is the native numerical component of \cref{struct:def:core}; its accepted-run projection supplies the witness cover. The full structure also retains names, denotational relations, and execution records. General multivalued answer interfaces remain in the larger structure; they are not turned into deterministic application by an unconstructed answer selection. The following comparison identifies its numerical application with the oracle-dialogue construction.


\subsection{Exact comparison with numerical oracle dialogues}
Let $g:\N\to2$ be a designated total oracle in a semantic layer. Write
$U^g(e,x)$ for the interpreter of \cref{struct:thm:pca}. Its accepted runs
contain ordinary finite machine transitions and query-answer steps checked
against $g$. Write $a\mathbin{\star_g}x$ for the dialogue application on
$K_1[g]$: a program is repeatedly run on the input and the list of answers
obtained so far, issuing either a new query or a final result. This is the
construction in \cite[Theorem 1.2]{relative-recursion}.

\begin{theorem}[Two actual compiler translations]\label{cal:thm:dialogues}
There are uniform numerical compilers $c$ and $r$, independent of $g$, such
that
\[
 U^g(e,x)\simeq c(e)\mathbin{\star_g}x,
 \qquad
 a\mathbin{\star_g}x\simeq U^g(r(a),x),
\]
where $\simeq$ means equality of partial values and equality of domains.
They induce translations of finite successful witnesses. Hence the two
interpreters represent exactly the same partial numerical functions and give
equivalent PCAs under applicative morphisms.
\end{theorem}
\begin{proof}
\leavevmode
\begin{enumerate}[label=\textup{(\roman*)},leftmargin=*,itemsep=3pt,topsep=4pt]
\item \textbf{Oracle interpretation.} On input $(x,u)$, where $u$ is the finite answer list, $c(e)$ replays the
machine $e$ on $x$.
At its successive oracle calls it consumes the successive
entries of $u$.
At the first unanswered call it returns the query tag and
query number.
If the machine halts it returns the result tag and value.
An ordinary divergence remains divergence.
The program never inspects $g$;
the dialogue mechanism supplies its answers.
Compilation is effective because
replay, list access, and machine simulation are ordinary computable operations.
\item \textbf{Converse construction.} Conversely $r(a)$ starts with the empty answer list, simulates the ordinary
program $a$ on $(x,u)$, and examines its result. A query tag triggers the actual
read $g(q)$ and appends the answer; a result tag halts with its value.
An
invalid tag causes divergence.
This is an oracle-machine program uniformly
compiled from $a$.
\item \textbf{Induction.} Induction on a finite run proves the first correspondence: after $n$ replies
the replay reaches exactly the same machine state as the original execution
after its $n$th read.
Induction on a dialogue proves the second: its successive
ordinary computations are precisely the simulations performed by $r(a)$.
The witness translations replay or concatenate these finite computations.
Both directions preserve the final value, and either successful witness
constructs a successful witness on the other side.
This proves exact domains.
\item \textbf{Unit and identity laws.} For the PCA assertion, use the identity relation on the underlying numerical
carriers as an applicative morphism in each direction.
The translated universal
simulator, with currying compiled by $s$-$m$-$n$, tracks the other application:
it takes its program number and input as \emph{data}.
Thus each application
is representable in the other PCA.
The composites of these identity relations
are identity relations.
This is applicative equivalence, not literal equality
of the two application tables or of compiled program names.
\end{enumerate}
\end{proof}

\begin{proposition}[Exact transport of the designated execution relation]
\label{cal:prop:transport}
Under geometric inverse image, the finite execution relation for $U^g$
transports to that for the pulled-back oracle. Its image, the exact application
domain, transports with it. The same holds for the two compiler comparisons.
\end{proposition}
\begin{proof}
\leavevmode
\begin{enumerate}[label=\textup{(\roman*)},leftmargin=*,itemsep=3pt,topsep=4pt]
\item \textbf{Existence and uniqueness.} Finite records use the natural-number and finite-list objects, the fixed
ordinary transition relations, and finitely many equations $g(q)=b$.
Acceptance is given by finite conjunction and existentially supplied finite
records.
Geometric inverse image preserves these constructions and regular
images.
Therefore it preserves both the record relation and its image in
program--input--output space.
The same finite-record proof describes the
compiled simulations, giving their comparisons and coherence after pullback.
\end{enumerate}
\end{proof}

\Cref{cal:prop:transport} is an exact statement for the designated numerical interpreter. It does
not assert that inverse image preserves every function object, or that the
entire forcing topos is the realizability topos of $K_1[g]$. The universal
property of relative recursion in \cite{relative-recursion} concerns the PCA
construction; \cref{cal:cor:dense-oracle} concerns the forcing localization.
The marked structure connects these different categorical statements.

\section{Arithmetic oracles and their sheaf interpretation}\label{sec:arithmetic-oracles}

The categorical oracle operations can now be tested against numerical computation. The presented arithmetic hierarchy and decision interfaces provide successive oracle stages; the numerical interpreter specifies what can actually be decided at each stage. We distinguish these constructive execution statements from the classical interpretations of arithmetic and the hierarchy--jump comparison.


Use the presented hierarchy of \cref{sec:presented-arithmetic} and
the dense decision construction of \cref{prop:decision}. Computability
continues to mean execution by the retained numerical interpreter or
its specified oracle extension. New sequences belong to its native function objects, such as
$(a_j2)^{a_j\N}$. \Cref{seq:thm:prefixes,seq:prop:transport}
constructs these objects from coherent prefixes and gives the map
from $a_j(2^\N)$ with its finite-prefix equations. Geometric path
and oracle presentations are base changed as locales, as in
\cref{thm:classifier-basechange,seq:cor:paths}.

The selection operations in the constructions of this section and
the bushy, Ramsey, and closed-class applications are either the
unique answer of \cref{prop:decision}, a finite
decidable selection, or the least witness of an inhabited decidable
predicate on $\N$. In the last case existence of a least witness follows
by induction on a supplied witness; the resulting least-witness relation
is total and functional. Unique-value comprehension, not a choice axiom,
produces the map. Composite dense inclusions are dense, because the
composite inverse image still reflects the empty subobject.

\subsection{Halting decisions and iterated oracle jumps}
For the ground finite-run interpreter of
\cref{thm:operational-certificates} put
\[
 H(e)\ \Longleftrightarrow\ \exists t\,
       \operatorname{Accept}(e,e,t).
\]
Here $\operatorname{Accept}(e,x,t)$ tests a numerical code $t$ for
an output value together with an accepted finite run of program $e$
on input $x$; it is the Boolean checker of \cref{thm:operational-certificates}
with the output and run packed into one number.
Apply \cref{prop:decision}. After exchanging the two output labels if
necessary, write $h(e)=1$ for a positive halting answer. For a constructed
Boolean oracle $A:\N\to2$, the native finite-dialogue checker of \cref{prop:intensional-native-checker} similarly defines
\[
 H^A(e)\ \Longleftrightarrow\ \exists t\,
       \operatorname{Accept}^A(e,e,t).
\]
This uses the explicit numerical oracle interpreter, not all morphisms
of the surrounding topos. Repeating \cref{prop:decision} constructs
$A',A'',\ldots$ as successive oracle extensions.

\begin{proposition}[Strictness at the program level]\label{arith:prop:jump-strict}
No program using only $A$ decides $H^A$ on all inputs. Nevertheless the
next decision localization constructs such a decision as a new oracle.
\end{proposition}
\begin{proof}
\leavevmode
\begin{enumerate}[label=\textup{(\roman*)},leftmargin=*,itemsep=3pt,topsep=4pt]
\item \textbf{Evaluation and transposition.} If $d^A$ were a total decision program with value $1$ exactly on $H^A$,
compile a program $q$ which, on input $n$, runs $d^A(n)$, halts when its
answer is $0$, and otherwise loops.
The compiler produces its numerical
index $q$.
Evaluation at $q$ gives
\[
 H^A(q)\quad\Longleftrightarrow\quad d^A(q)=0
 \quad\Longleftrightarrow\quad\neg H^A(q),
\]
a contradiction by the two possible Boolean answers.
The compiler and finite-run correctness are \cref{sec:numericalcodes,thm:operational-certificates,prop:intensional-native-checker}.
\end{enumerate}
\end{proof}

The external numerical calibration uses van Oosten's
oracle-adjunction theorem, which constructs $A[f]$ with
its universal property, identifies $K_1[f]$ with ordinary oracle
application, and constructs the associated geometric inclusion
\cite[Theorem 1.2, Corollary 1.3, Section 2]{relative-recursion}. The least-local-operator
characterization and its realizability interpretation are recorded in
\cite[Section 6.1]{faber-oosten}. This supplies the categorical connection
with relative recursion. It does not make a program presentation
recoverable after forgetting it.

\subsection{The exact arithmetic translation under localization}
\begin{proposition}[Closed-subobject interpretation]\label{arith:prop:closed-interpretation}
For a local operator $j$ and ground objects $X$, subobjects of $a_jX$
are identified, by pullback along the unit, with $j$-closed subobjects of
$X$. Under this identification the logical operations are
\[
 \begin{array}{c|c}
 \text{operation}&\text{operation on closed predicates}\\\hline
 \top,\bot&\top,j\bot\\
 A\wedge B&A\wedge B\\
 A\vee B&j(A\vee B)\\
 A\Rightarrow B&A\Rightarrow B\\
 \exists_\pi A&j(\exists_\pi A)\\
 \forall_\pi A&\forall_\pi A.
 \end{array}
\]
Atoms are interpreted by closing their actual subobjects. These clauses,
applied by recursion to the retained formulas, give exactly their
interpretation in the sheaf topos.
\end{proposition}
\begin{proof}
\leavevmode
\begin{enumerate}[label=\textup{(\roman*)},leftmargin=*,itemsep=3pt,topsep=4pt]
\item \textbf{Pullback construction.} The closed-subobject correspondence is the unit pullback
correspondence for sheafification.
Finite intersections of closed
subobjects are closed; joins and images are followed by closure.
If $B$ is closed, then
$j(A\Rightarrow B)\wedge A\leq
j((A\Rightarrow B)\wedge A)\leq jB=B$,
so implication into $B$ is closed.
For a closed $A$ on the projection
domain, pullback stability of closure gives
\[
 \pi^*j(\forall_\pi A)
 =j(\pi^*\forall_\pi A)\leq jA=A.
\]
Adjunction gives $j(\forall_\pi A)\leq\forall_\pi A$, so the universal
quantifier is closed as well.
\item \textbf{Verification of the equations.} Their adjunction universal properties on
closed predicates identify them with the sheaf operations.
Induction on
formulas finishes the proof.
\end{enumerate}
\end{proof}

For $j=\neg\neg$ this is a construction of a Boolean interpretation;
excluded middle is not appended as an axiom to $\Eff$. The level-zero
matrices remain the originally named programs. Replacing them by
\emph{all} decidable subobjects in the Boolean topos would erase the
computational hierarchy, since every subobject there is decidable.
Likewise, $a_j(\forall x\,R)$ must not be silently replaced by
$\forall x\,a_jR$. The closed interpretation defined by recursion on formulas supplies the arithmetic semantics.

\subsection{The degree-free Post calibration}
In the classical numerical interpretation, let $A:\N\to2$ be a Boolean oracle,
let $A^{(0)}=A$, and let $A^{(n+1)}$ be the characteristic function of
$H^{A^{(n)}}$. Write $\SigA_n^A,\PiA_n^A$ for the presented hierarchy
with $A$-computable matrices. The same calculation can be interpreted in
the constructed Boolean localization using the retained interpreters and
\cref{arith:prop:closed-interpretation}.

\begin{theorem}[Arithmetic complexity as relative computation]\label{arith:thm:post}
For $n\geq0$, a presented numerical predicate is in $\SigA_{n+1}^A$
exactly when it is the domain of a partial program using $A^{(n)}$.
It belongs to both $\SigA_{n+1}^A$ and $\PiA_{n+1}^A$ exactly when
it is decided by a total program using $A^{(n)}$. The translations
are effective on the indicated formula and program presentations.
For $n\geq1$, $\SigA_n^A\ne\PiA_n^A$.
\end{theorem}
\begin{proof}
\leavevmode
\begin{enumerate}[label=\textup{(\roman*)},leftmargin=*,itemsep=3pt,topsep=4pt]
\item \textbf{Initial construction.} First record the finite-transcript calculation.
If a Boolean oracle $B$
has a $\SigA_m^A$ definition with $m\geq1$, the correctness of a finite
$B$-oracle transcript is a finite conjunction of $\SigA_m^A$ predicates
for positive answers and $\PiA_m^A$ predicates for negative answers.
Each can be written with a $\SigA_{m+1}^A$ prefix by adding dummy
quantifiers. A finite conjunction of these prefixes is again of that
form: code the finite list of witnesses at each existential block, and
the finite list of challenges at each universal block.
\item \textbf{Tracked implementation.} The decidable
matrix checks the components and the ordinary machine transitions.
The same finite-list coding works when the number of queries is given
by the transcript.
Existentially quantifying the transcript merges with
the first existential block.
Thus a program domain relative to such a
$B$ is $\SigA_{m+1}^A$, uniformly.
For $B=A$ the finite transcript is
decidable using $A$, so its existence is $\SigA_1^A$ directly.
Classical negation of a displayed prefix is used here; this paragraph
is in the stated Boolean interpretation.
\item It follows inductively that $A^{(m)}$ has a $\SigA_m^A$ definition
for every $m\geq1$: the first jump is the existence of an $A$-oracle
halting transcript, and the transcript calculation gives the induction
step.
Applying the same calculation to an arbitrary program using
$A^{(n)}$ proves the domain-to-formula direction of the theorem.
\item \textbf{Converse construction.} For the converse, a $\SigA_1^A$ formula is semidecided by searching its
$A$-computable matrix.
Inductively, a $\PiA_n^A$ predicate is the complement
of a $\SigA_n^A$ predicate, whose semidecision program uses
$A^{(n-1)}$ by induction.
Compile the input into a program that ignores
its own input and performs this semidecision.
One query to $A^{(n)}$
decides whether that program halts, and hence decides the original
$\PiA_n^A$ predicate.
Searching the first quantifier of
$\exists z\,R(x,z)$ now semidecides every $\SigA_{n+1}^A$ predicate
using $A^{(n)}$.
\item \textbf{Identification of the resulting maps.} All compilations are finite syntactic transformations
and specialization operations already supplied by the interpreter.
\item \textbf{Oracle interpretation.} If both a predicate and its complement have these semidecision
programs, dovetail them; exactly one eventually returns, giving a
total decision.
Conversely, from a decision program construct the two
semidecision programs by returning on the chosen answer and looping on
the other.
This proves the intersection assertion.
\item Finally
$A^{(n)}\in\SigA_n^A$.
If it also belonged to $\PiA_n^A$, the
intersection assertion would give an $A^{(n-1)}$-program deciding its
own jump, contrary to \cref{arith:prop:jump-strict}.
This proves strictness.
\end{enumerate}
\end{proof}

This obtains the hierarchy--jump theorem without forming degrees.
The formula/program presentations, and the explicit location of the
classical prefix dualities, are essential to the statement. The
constructive analysis of those dualities is discussed in \cite{forster-kirst-mueck};
they are not claimed as additional equations in the ground doctrine.

\part{Regions, constraints, and forcing questions}\label{part:regions}

A forcing condition can constrain a region of possible oracles rather than identify a single value. We therefore study containment and positive possibility as properties of a locale map over the oracle space. Their universal saturation explains how stage questions represent geometric observations; finite witnesses and concrete forcing orders then show how these abstract operations are implemented.


\section{Relative regions and saturated universal observation}\label{sec:regions}

Not every forcing datum is a point or an answer map. A region records a collection of allowed oracle outcomes with its own possible failure or nonemptiness. We express stage containment as factorization of locale maps, then calculate the universal saturation that transports those constraints into forcing observations.


A region is a locale map $k:K\to X$, regarded as an object of $\Loc(\E)/X$. It need not be monic. For an open $U$ of $X$, write $K\leq_XU$ when $k$ factors through the open sublocale $X_U\incl X$. Equivalently $k^*U=1_K$. Write $K[U]=K\times_XX_U$ for its restriction.

\begin{lemma}[Containment calculus]\label{lem:containment}
Containment preserves top and finite conjunction, is monotone in $U$, and is preserved by a map of regions $K'\to K$ over $X$. It can be checked on a jointly surjective family of regions over $K$. Moreover $K\leq_X0$ precisely when $K$ is the initial locale.
\end{lemma}
\begin{proof}
\leavevmode
\begin{enumerate}[label=\textup{(\roman*)},leftmargin=*,itemsep=3pt,topsep=4pt]
\item \textbf{Containment calculus.} All but the covering assertion follow from preservation of finite meets, top, and bottom by $k^*$.
For a jointly surjective family $e_a:K_a\to K$, the map $\coprod_aK_a\to K$ has injective inverse image.
If all $e_a^*k^*U=1$, that inverse image identifies $k^*U$ with $1$, so they are equal.
Finally $k^*0=0_K$, and $0_K=1_K$ is exactly degeneracy of the frame.
\end{enumerate}
\end{proof}

\begin{definition}[Condition regions]\label{def:region-system}
A region system consists of $k_p:K_p\to X$ and maps $r_{p,q}:K_q\to K_p$ over $X$ for $p\leq q$, satisfying identity and composition. It is indexed in the base: for $p:T\to P$ the region lives internally over $T$, and the maps commute with parameter substitution.

Put $m_U(p)\iff K_p\leq_XU$. Local join splitting means that, whenever $K_p\leq_X\bigvee_iU_i$, a cover $S\in J(p)$ exists with
$\forall q\in S\,\exists i\,K_q\leq_XU_i$.
Its nullary case is local consistency: an empty region has an empty covering sieve. Cover exhaustivity means that $\coprod_{q\in S}K_q\to K_p$ is surjective for every cover. A system with local join splitting and cover exhaustivity is called exact.
\end{definition}

These clauses describe the constructions being combined: splitting is a local way to realize a join, and exhaustivity lets the displayed regions detect the structural covers. They do not assert absolute compactness of a region. The canonical regions below satisfy all clauses by construction.

\begin{theorem}[Region realization]\label{thm:regions}
Local join splitting constructs a frame homomorphism
\[
 \rho_K^*:\OO(X)\to\Omega_J(P),\qquad
 \rho_K^*U=q_J(m_U).
\]
For an exact system the predicates $m_U$ are already stable and
\[
 [p]_J\leq\rho_K^*U\iff K_p\leq_XU.
\]
The resulting realization is uniquely determined by this exact stage equation.
\end{theorem}
\begin{proof}
\leavevmode
\begin{enumerate}[label=\textup{(\roman*)},leftmargin=*,itemsep=3pt,topsep=4pt]
\item \textbf{Region realization.} Refinement maps and \cref{lem:containment} make $m_U$ upper.
The same lemma gives $m_1=1$ and $m_{U\wedge V}=m_U\wedge m_V$.
The nullary splitting clause gives $m_0\leq j_J0$, and hence $q_Jm_0=0_{\Omega_J}$.
\item \textbf{Localization and coverage.} For a family $(U_i)$, monotonicity gives
$j_J(\bigvee_i m_{U_i})\leq j_Jm_{\bigvee_iU_i}$.
If $m_{\bigvee_iU_i}(p)$, splitting gives a cover contained in $\bigvee_i m_{U_i}$; therefore $m_{\bigvee_iU_i}\leq j_J(\bigvee_i m_{U_i})$.
Applying the nucleus proves the reverse inequality.
Thus $\rho_K^*$ preserves all joins and finite meets.
\item \textbf{Forcing construction.} If a cover of $p$ satisfies $m_U$ throughout, cover exhaustivity and \cref{lem:containment} give $m_U(p)$.
Hence $m_U$ is stable.
The exact stage equation follows from \cref{cor:stage-truth}.
It determines $\rho_K^*U$ on every stage, and stages are join-dense, proving uniqueness.
\end{enumerate}
\end{proof}

There is an adjoint description of raw containment. For the terminal map $!_p:K_p\to1$, inverse image has a right adjoint on frames. Then
\[
 m_U(p)=(!_p)_*(k_p^*U)\in\Omega.
\]
Indeed a truth value $v$ is below this right adjoint precisely when $!_p^*v\leq k_p^*U$, which says that the region, restricted to the context $v$, lies in $U$. Thus the raw must operation is a universal quantifier; saturation turns it into the forcing inverse image.

\begin{proposition}[Geometric regional interpretation]\label{prop:region-logic}
For an exact region system and a geometric oracle assertion with open interpretation $\llbracket\varphi\rrbracket\in\OO(X)$,
\[
 p\Vdash\varphi(\dot G)\iff K_p\leq_X\llbracket\varphi\rrbracket.
\]
Disjunction is witnessed locally by the splitting cover. Assertions in additional sheaf contexts are interpreted by \cref{thm:kj}.
\end{proposition}
\begin{proof}
\leavevmode
\begin{enumerate}[label=\textup{(\roman*)},leftmargin=*,itemsep=3pt,topsep=4pt]
\item \textbf{Geometric regional interpretation.} Combine \cref{thm:generic-observation,thm:regions}.
The local disjunction clause is exactly local join splitting in one direction and cover exhaustivity in the other.
Additional context variables change the subobject being interpreted, so their semantics is the Kripke--Joyal forcing already established.
\end{enumerate}
\end{proof}

\begin{theorem}[Selected effective regional truth]\label{thm:effective-regions}
Let $U:T\to M_X$ be a selected family. A jointly positive predicate $C:P\times T\to\Sigma$ satisfying
\[
 iC(p,t)\iff j_J(m_{m_XU(t)})(p)
\]
constructs its unique stable effective lift along $\rho_K$. For an exact region system it suffices that $C$ classify raw containment. If raw containment is positive and $j_J$ has a global effective restriction, applying that restriction constructs the required $C$.
\end{theorem}
\begin{proof}
\leavevmode
\begin{enumerate}[label=\textup{(\roman*)},leftmargin=*,itemsep=3pt,topsep=4pt]
\item \textbf{Selected effective regional truth.} The displayed full predicate is upper and stable because it is saturated.
Its positive classifier therefore transposes to a map $T\to\Up_\Sigma(P)$ representing that same stable family.
Pair it with $\rho_K^*m_XU$ in the pullback defining $S_J$.
This constructs the lift and proves uniqueness.
\item \textbf{Evaluation and transposition.} For exact regions, \cref{thm:regions} identifies raw containment with its saturation.
In the final case, let $A_\Sigma:T\to\Up_\Sigma(P)$ be the positive raw trace.
The composite $j_{J,\Sigma}A_\Sigma$ has underlying predicate $j_Jm_{m_XU}$ by the defining square for global restriction, so its evaluation is the required $C$.
The site-to-Scott comparison transports the lift when the condition presentation is effective.
\end{enumerate}
\end{proof}


Universal region saturation expresses what every allowed oracle must satisfy; it does not by itself supply a positive witness that some oracle satisfies a test. The next section introduces that existential possibility and specifies when it admits finite effective evidence.

\section{Positive possibility and finite witness questions}\label{sec:possibility}

Containment expresses that every possible oracle in a region satisfies an observation. Existential possibility asks a different question, and it becomes positive only when the relevant locale is overt. We separate these tests and compare their finite witnesses in concrete forcing presentations.


For an overt locale $K$, the terminal inverse image has a left adjoint
$\LPos_K:\OO(K)\to\Omega$. Define
\[
 \May_K(p,U)=\LPos_{K_p}(k_p^*U),\qquad
 \operatorname{Impossible}_K(p,U)\iff k_p^*U=0.
\]

\begin{proposition}[May and impossible calculus]\label{prop:may}
Positive may preserves arbitrary joins and bottom. For overt regions,
\[
 \operatorname{Impossible}_K(p,U)\iff\neg\May_K(p,U).
\]
If $K_p$ is positive, raw must implies positive may. Impossible is persistent under refinement, takes joins to conjunctions, and is stable under structural covers in a cover-exhaustive system.
\end{proposition}
\begin{proof}
\leavevmode
\begin{enumerate}[label=\textup{(\roman*)},leftmargin=*,itemsep=3pt,topsep=4pt]
\item \textbf{Pullback construction.} The composite $\LPos_{K_p}k_p^*$ preserves joins because both factors do.
By the adjunction,
\[
 \neg\LPos_{K_p}(a)
 \iff\LPos_{K_p}(a)\leq0
 \iff a\leq!_p^*0
 \iff a=0.
\]
For a positive region, $\LPos_{K_p}(1)=1$, so containment of the whole region gives may.
Pullback along a refinement preserves zero, proving persistence of impossibility.
Preservation of joins by $k_p^*$ gives its join-to-conjunction law.
Under a covering family, injectivity of the joint inverse image reflects equality with zero, exactly as in \cref{lem:containment}, proving stability.
\end{enumerate}
\end{proof}

\begin{corollary}[Complemented region observations]\label{cor:complemented-regions}
Let $U\in\OO(X)$ have a complement $U^c$, so
$U\wedge U^c=0$ and $U\vee U^c=1$. For any condition region,
\[
 K_p\leq_XU
   \quad\Longleftrightarrow\quad
 \operatorname{Impossible}_K(p,U^c).
\]
If the region $K_p$ is overt, then
\[
 \May_K(p,U)\vee\May_K(p,U^c)
       =\LPos_{K_p}(1_{K_p}).
\]
For a positive overt region this last truth is top. These are
complemented-observation calculations, not a general decision
principle for arbitrary oracle opens.
\end{corollary}
\begin{proof}
\leavevmode
\begin{enumerate}[label=\textup{(\roman*)},leftmargin=*,itemsep=3pt,topsep=4pt]
\item \textbf{Complemented region observations.} Inverse image preserves the complementary equations.
Thus
$k_p^*U$ is complementary to $k_p^*U^c$, and one is top exactly when
the other is bottom, giving the first equivalence.
For overt $K_p$,
$\LPos_{K_p}$ preserves joins by its left adjoint property, hence
\[
 \LPos_{K_p}(k_p^*U)\vee\LPos_{K_p}(k_p^*U^c)
 =\LPos_{K_p}(k_p^*(U\vee U^c))
 =\LPos_{K_p}(1_{K_p}).
\]
Positivity identifies the final truth with top.
\end{enumerate}
\end{proof}

May need not persist under strengthening: a refinement can remove the part where $U$ is positive. Nor does $\May\vee\neg\May$ follow constructively. Must can imply may on a positive region, so the three expressions do not form three disjoint truth values.

A condition-theoretic question is
\[
 \May_{\mathrm{ref}}(p,U)\iff
 \exists q\geq p\,K_q\leq_XU.
\]
A degenerate refinement witnesses this for every $U$. The positive version includes $\LPos_{K_q}(1)$ in the existential predicate.

\begin{proposition}[Positive refinement witnesses]\label{prop:positive-extension}
A positive refinement witnessing must also witnesses positive may at the original region. For a family on which the extension construction produces such a refinement from every positive compatible region, positive may and positive refinement may agree.
\end{proposition}
\begin{proof}
\leavevmode
\begin{enumerate}[label=\textup{(\roman*)},leftmargin=*,itemsep=3pt,topsep=4pt]
\item \textbf{Monad compatibility.} For a region map $r:K_q\to K_p$, the unit of $\LPos_{K_p}\dashv!_p^*$ gives
$a\leq!_p^*\LPos_{K_p}(a)$.
Pulling back along $r$ and applying $\LPos_{K_q}$ gives
$\LPos_{K_q}(r^*a)\leq\LPos_{K_p}(a)$.
Take $a=k_p^*U$. A positive refining region contained in $U$ makes the left side top.
This proves the first assertion.
The stipulated extension construction supplies the reverse direction for its displayed family, so the predicates coincide there.
\end{enumerate}
\end{proof}

In an effective presentation the classifier for raw refinement may is explicitly
\[
 \exists_P^\Sigma q\,[\,p\leq_\Sigma q\wedge C(q,t)\,]
\]
when the refinement classifier, overt search, and selected containment trace $C$ are the ones constructed for the presentation. The positive version inserts the selected positive-region classifier. This formula explains why the existence of an extensional refining condition by itself is not a search procedure.

\subsection{A finite bushy-witness calculation}
Work in $\Eff$ for this finite-witness calculation, so that the coded searches use \cref{thm:relativefoundation,prop:enumerated-overt}. For a stem $\sigma\in\N^{<\N}$ and a branching grade $b:\N\to\N$ with $b(n)\geq2$, a finite tree above $\sigma$ is $b$-bushy when every non-leaf node $\tau$ has at least $b(|\tau|)$ immediate successors in the tree. A predicate $U$ on finite strings is $b$-big above $\sigma$ when such a finite tree has all leaves in $U$.

\begin{proposition}[Positive bushy witnesses]\label{prop:bushy-positive}
For a $\Sigma$-classified $U$, the bigness predicate is $\Sigma$-classified uniformly in a displayed effective grade. Finite grafting composes witnesses. If $U\cup V$ is $(b+c-1)$-big above $\sigma$, then $U$ is $b$-big or $V$ is $c$-big above $\sigma$, in the corresponding finite decidable leaf-labelled formulation.
\end{proposition}
\begin{proof}
\leavevmode
\begin{enumerate}[label=\textup{(\roman*)},leftmargin=*,itemsep=3pt,topsep=4pt]
\item \textbf{Positive bushy witnesses.} A finite coded tree, its branching counts, and a finite family of leaf labels give an overt search over codes.
Validity of the tree and its bushiness are finite decidable checks.
The conjunction of the leaf predicates is positive by \cref{lem:finite-compact}.
Existential search therefore gives a positive classifier.
\item \textbf{Finite witness construction.} For grafting, attach at each leaf a displayed finite witness rooted there, identifying that leaf with the root. A finite union of finite coded trees is finite, and every non-leaf belongs either to the original tree or to one of the attached trees with its required successors.
Leaves are the leaves of the attached witnesses.
\item \textbf{Universal-property argument.} For the partition assertion, label a leaf red or blue according to its supplied membership in $U$ or $V$.
Proceed upwards through the finite tree.
Label a node red if at least $b(n)$ children have been labelled red, where $n$ is its length; otherwise fewer than $b(n)$ are red, so among at least $b(n)+c(n)-1$ children at least $c(n)$ are blue, and label it blue.
This finite counting statement is decidable.
The nodes reached by retaining the indicated children form a $b$-bushy red witness or a $c$-bushy blue witness according to the root label.
\end{enumerate}
\end{proof}

The supplied leaf labels are legitimate local data when $U\cup V$ is used as a positive union: a finite family of existential choices can be eliminated into the finite conclusion by induction. No uniform decision of membership in arbitrary $\Sigma$-predicates is inferred. This finite witness calculation does not decide smallness. The next construction proves the negative-invariant operations and obtains their total successor in the explicit decision localization; the retained forbidden set and grade record the resulting promise history.

\subsection{Negative bushy constraints and constructed successors}\label{sec:bushy-negative}
Specialize the preceding grade to a constant $k\geq2$.
Finite trees are coded, prefix closed above their root, and contain the
root. A tree is $k$-bushy if every non-leaf has at least $k$ immediate
successors. A singleton-root tree is allowed. Put
\[
 C_kB(\sigma)\ \Longleftrightarrow\ \BigWitness_k(B,\sigma)
 \ \Longleftrightarrow\
 \exists T\,[T\text{ finite $k$-bushy above $\sigma$ and }
                       \operatorname{Leaf}(T)\subseteq B].
\]
Smallness is $\neg C_kB(\sigma)$.

\begin{lemma}[Closure and smallness arithmetic]\label{bush:lem:closure}
For every predicate $B$ on finite strings, $C_k$ is monotone, extensive,
and idempotent. Moreover
\[
 \Small_m(B,\sigma)\wedge\Small_n(C,\sigma)
 \ \Longrightarrow\ \Small_{m+n-1}(B\cup C,\sigma).
\]
For a semidecidable $B$, $C_kB$ is semidecidable, uniformly in its code.
\end{lemma}
\begin{proof}
\leavevmode
\begin{enumerate}[label=\textup{(\roman*)},leftmargin=*,itemsep=3pt,topsep=4pt]
\item \textbf{Closure and smallness arithmetic.} Extensivity uses the singleton tree; monotonicity uses the same witness.
For idempotence, take a finite witness with leaves in $C_kB$ and graft
a witnessing $k$-bushy $B$-tree at each leaf.
Finite existential
elimination assembles the witnesses.
The result is a finite $k$-bushy
$B$-tree, proving $C_kC_kB\subseteq C_kB$.
\item \textbf{Verification of the equations.} For the smallness formula, a counterexample would give a finite
$(m+n-1)$-bushy witness for $B\cup C$.
The leaf-labelled partition
calculation of \cref{prop:bushy-positive}, with constant grades,
produces either an $m$-bushy $B$-witness or an $n$-bushy $C$-witness.
Both contradict the displayed smallness predicates.
\item Finally enumerate finite tree
codes and finite halting certificates for leaf membership.
This is
exactly a semidecision of bigness.
\end{enumerate}
\end{proof}

\begin{remark}[A closure that must not be confused with upward closure]
Even for upward-closed $B$, $C_kB$ need not be upward closed.
For example, let $k=2$ and let $B$ consist of strings extending $00$ or
$01$. Then $0\in C_2B$ but $02\notin C_2B$. Thus calling $C_kB$ an
upward-closed forbidden region would be incorrect. The following
normalization supplies both properties on the working cone.
\end{remark}

For a root $\sigma$ define, on its extension cone,
\[
 N_{\sigma,k}B(\tau)\ \Longleftrightarrow\
 \exists\rho\,[\sigma\preceq\rho\preceq\tau\ \wedge\ C_kB(\rho)].
\]
\begin{proposition}[Constructed forbidden-region normalization]\label{bush:prop:normalization}
$N_{\sigma,k}B$ is upward closed and $k$-closed on the cone of
$\sigma$, contains $B$ on that cone, and is $k$-small at $\sigma$
whenever $B$ is. This normalization is idempotent and preserves
semidecidability.
\end{proposition}
\begin{proof}
\leavevmode
\begin{enumerate}[label=\textup{(\roman*)},leftmargin=*,itemsep=3pt,topsep=4pt]
\item \textbf{Observation preservation.} Upward closure and inclusion are immediate from the formula and the
singleton witness.
Suppose a finite $k$-bushy tree above $\tau\geq\sigma$
has leaves in the normalization.
For each leaf select its witnessing
prefix $\rho$ and its finite $B$-tree.
These are finitely many witnesses.
If one selected $\rho$ is a prefix of $\tau$, it directly witnesses
$N_{\sigma,k}B(\tau)$.
Otherwise all selected $\rho$ extend $\tau$.
Prune the original tree at the prefix-minimal selected nodes.
Every
branch meets one of them; every unpruned non-leaf retains all its
children.
\item \textbf{Verification of the equations.} The pruned tree is still $k$-bushy.
Grafting the selected
$B$-trees at its leaves gives $C_kB(\tau)$, again proving membership
in the normalization.
The case distinction uses only the decidable
prefix relation on finitely many already selected strings.
\item \textbf{Observation preservation.} Taking $\tau=\sigma$, bigness of the normalized set implies
$C_kB(\sigma)$: in the first case the selected prefix can only be
$\sigma$, and the second case gives the same conclusion by grafting.
This proves preservation of smallness. A $k$-closed upward set on the
cone is fixed by the normalization, proving idempotence.
Its defining
finite-prefix existential and \cref{bush:lem:closure} give the uniform
semidecision claim.
\end{enumerate}
\end{proof}

\begin{theorem}[An explicit safe extension]\label{bush:thm:successor}
In the decision localization for the family $C_kB(\tau)$, a root
$\sigma$ at which $B$ is $k$-small has a canonically selected immediate
successor at which $B$ is $k$-small. Iteration produces a total branch
whose prefixes extending that root avoid $B$. For semidecidable $B$ and a presented finite grade $k$,
all these decisions reduce uniformly to the halting oracle.
\end{theorem}
\begin{proof}
\leavevmode
\begin{enumerate}[label=\textup{(\roman*)},leftmargin=*,itemsep=3pt,topsep=4pt]
\item \textbf{Finite witness construction.} Consider $\sigma0,\ldots,\sigma(k-1)$.
If bigness held at all of
them, graft their finite witnesses below the root, making $B$ big at
$\sigma$.
Thus they cannot all be big.
In the decision extension their
bigness predicates are decidable, so a finite scan finds the first
small child.
This is a total function on the small-root states.
Recursion constructs extensions of every larger length, each still small.
The
singleton-tree implication $B\subseteq C_kB$ shows that none belongs
to $B$.
\item \textbf{Oracle interpretation.} For an upward-closed $B$, the earlier prefixes of the starting
root avoid $B$ as well.
For semidecidable $B$, compile the tree-witness search of
\cref{bush:lem:closure} to a halting query.
This constructs the reduction.
\end{enumerate}
\end{proof}

An update adding an $n$-small forbidden set to an $m$-small one uses
grade $m+n-1$ and then \cref{bush:prop:normalization}. Consequently
smallness preservation, closure, and safe extension are proved operations;
they need not be hidden in an admissibility package. The underlying
finite combinatorics are the lemmas of
\cite[Lemmas 2.6--2.8]{khan-miller}; the normalization above also resolves the
upward-closure issue explicitly.

A total function $g:\N\to\N$ is \textbf{diagonally noncomputable relative to $A$} when $\varphi_e^A(e)\converges$ implies $g(e)\ne\varphi_e^A(e)$ for every $e$. Write $\DNC^A$ for this class; without an oracle, write $\DNC$.

\begin{corollary}[A complete diagonal-avoidance instance]\label{bush:cor:diagonal}
Let
\[
 B_{\mathrm{diag}}(\sigma)\ \Longleftrightarrow\
 \exists e<|\sigma|\ \exists t\,
          [\varphi_e(e)\text{ halts by }t\text{ with value }\sigma(e)].
\]
This is upward closed, semidecidable, and $2$-small at the empty string.
The preceding successor therefore constructs, in the halting decision
extension, a total $g$ with
$\varphi_e(e)\converges\Rightarrow g(e)\ne\varphi_e(e)$.
\end{corollary}
\begin{proof}
\leavevmode
\begin{enumerate}[label=\textup{(\roman*)},leftmargin=*,itemsep=3pt,topsep=4pt]
\item \textbf{A complete diagonal-avoidance instance.} Suppose a finite $2$-bushy tree had all leaves in $B_{\mathrm{diag}}$.
Extract a diagonal convergence certificate at each leaf and take a
common bound on their times.
At each internal node of length $e$, at
most one child has the diagonal value seen by that bound.
Follow a
different child, chosen from the at least two available.
This terminates
at a leaf whose chosen convergence certificate was included in the
bound; the path avoided that value at its corresponding coordinate, a
contradiction.
\item \textbf{Verification of the equations.} The successor theorem now applies and every finite
diagonal violation would appear in one of its forbidden prefixes.
\end{enumerate}
\end{proof}

\paragraph{The resulting bushy condition presentation.}
Use coded states $(\sigma,B,k)$, where $B$ is an enumerated
upward-closed forbidden set and is $k$-small above $\sigma$.
Upward closure of a raw enumeration is obtained by enumerating all
extensions of each enumerated string. A strengthening lengthens the
stem, enlarges the forbidden set, and has its own verified smallness
grade. The preceding normalization, union calculation, and successor
are actual operations on these states: their proofs establish the
output smallness invariant. For global forbidden sets the normalized
update is $B\cup N_{\sigma,k}B$; its restriction to the working cone
is exactly the normalization proved above. No decidable validity test
for arbitrary codes is presumed. The diagonal state of
\cref{bush:cor:diagonal} supplies a constructed starting condition.

In the indicated halting decision extension, the sets of conditions
whose stems have length at least $n$ are dense, by
\cref{bush:thm:successor}. Give this order the coverage generated by
those length refinements. Its generic directed filter has a total
union of stems. For any condition in that filter, the union avoids
its forbidden set: a forbidden prefix and a sufficiently long later
stem would, by upward closure and a common strengthening, put a valid
stem in its own forbidden set. The singleton-tree test excludes that.
This constructs the condition site, its generic object, and the
negative-preservation statement used by \cref{thm:forcing-domain-oracle-identity}.
\Cref{cal:thm:bushy} below constructs a further requirement refinement by these operations.

\subsection{Full bushy carriers and a uniform requirement compiler}
The basic order in \cite[Section 3]{khan-miller} permits arbitrary forbidden sets
$B$, whereas the preceding implemented fragment uses c.e. ones. Retaining the full forbidden-set carrier, with grades $k\geq2$, gives the
condition subobject of
$\N^{<\N}\times\mathcal P(\N^{<\N})\times\N$ given by upwardness and
smallness. This retains the existing smallness operations, with answers to
the actual bigness queries. It does not grant arbitrary $B$ a c.e. code.

\begin{proposition}[Forgetting a displayed bushiness grade]
\label{cal:prop:bushy-codes}
On the full carrier, the projection $(\sigma,B,k)\mapsto(\sigma,B)$,
whose target records existence of a finite smallness grade, identifies the
dense sheaf topoi of the two presentations.
\end{proposition}
\begin{proof}
\leavevmode
\begin{enumerate}[label=\textup{(\roman*)},leftmargin=*,itemsep=3pt,topsep=4pt]
\item \textbf{Forgetting a displayed bushiness grade.} A stronger pair lifts by its locally supplied grade.
Two grades over one
pair have a common strengthening with that same pair and either valid grade:
grade is not an additional restraint in the order.
The projection is regular
epic by the existential definition of its target.
Apply
\cref{cal:lem:codes}.
\end{enumerate}
\end{proof}

\begin{proposition}[The enumerated fragment is not dense]\label{cal:prop:bush-notdense}
The fragment with c.e. forbidden sets is not dense in the full basic bushy order.
\end{proposition}
\begin{proof}
\leavevmode
\begin{enumerate}[label=\textup{(\roman*)},leftmargin=*,itemsep=3pt,topsep=4pt]
\item \textbf{Observation preservation.} In the classical interpretation take
\[
 B_{0'}=\{\sigma:\exists e<|\sigma|\,
                         \varphi_e^{0'}(e)\downarrow=\sigma(e)\}.
\]
The finite diagonal-avoidance argument makes this upward set $2$-small at
the empty string.
If a valid c.e. condition $(\tau,C,k)$ strengthened
$(\varnothing,B_{0'},2)$, then $C\supseteq B_{0'}$.
\Cref{bush:thm:successor}
would construct a branch $g\leq_T0'$ above $\tau$ avoiding $C$.
Upwardness
also excludes forbidden prefixes shorter than $\tau$.
\item Thus $g\in\DNC^{0'}$, which is impossible for a function computable in $0'$: apply its own
$0'$-program index to itself.
\end{enumerate}
\end{proof}

We can nevertheless implement a substantial requirement uniformly in the
enumerated fragment. The following implements the diagonalization in \cite[Claim 3.6]{khan-miller}, using \cref{bush:prop:normalization,bush:thm:successor}. The existence theorem is due to Ambos-Spies,
Kjos-Hanssen, Lempp, and Slaman \cite[Theorem 1.8]{ambos-spies-etal}; the point here is
its realization by the designated operations, including uniform treatment of
partial bound programs.

\begin{theorem}[A uniform requirement compiler]\label{cal:thm:bushy}
From a valid c.e. condition $p=(\sigma,B,k)$ and numerical program indices
$(\Gamma,h)$ one can compute, using $0'$, a valid stronger c.e. condition.
For every total positive-valued $h$, every branch compatible with the output
fails to make $\Gamma$ a total $h$-bounded DNC function. The construction is
total on all pairs of indices, including partial $h$.
\end{theorem}
\begin{proof}
\leavevmode
\begin{enumerate}[label=\textup{(\roman*)},leftmargin=*,itemsep=3pt,topsep=4pt]
\item \textbf{Observation preservation.} Normalize $B$.
Let $\Gamma_h$ be the clipped functional: compute $h(x)$ and
$\Gamma(x)$, returning the latter only if it is smaller than the former.
Compile an ordinary program $\theta$ which, on input $m$, searches for a
finite $k$-bushy tree above $\sigma$ with all leaf computations
$\Gamma_h^\tau(m)$ converging to one common value.
Let $e$ be its index.
No self-reference operation is needed: $e$ is obtained from these finite
parameters by the ordinary compiler.
\item \textbf{Verification of the equations.} Execute the following finite case
procedure:
\[
\begin{array}{ll}
 h(e)\uparrow:&\text{return the normalized input};\\
 h(e)=0:&\text{return the normalized input};\\
 h(e)=b>0:& A=\{\tau:\Gamma_h^\tau(e)\downarrow\},\quad
                  \text{decide }C_{bk}A(\sigma);\\
 \neg C_{bk}A(\sigma):&
      \text{return the normalization of }(\sigma,B\cup A,k+bk-1);\\
 C_{bk}A(\sigma):&
      \text{run }\theta(e),\text{ choose a leaf }\tau\notin B
        \text{ of its witness, return }(\tau,B,k).
\end{array}
\]
All decisions here are halting queries on ordinary c.e. searches, including
membership of the normalized $B$.
\item \textbf{Universal-property argument.} In the last case, finite leaf-colouring
turns a $bk$-bushy $A$-witness into a $k$-bushy constant-output witness, so
$\theta(e)$ terminates.
Smallness supplies a leaf outside $B$; $k$-closure
makes the new root small.
Its value equals $\varphi_e(e)$.
In the preceding
case, smallness arithmetic validates the update and avoidance of $A$ prevents
convergence.
The first case is irrelevant for total $h$; the second precludes
any total $h$-bounded function.
These alternatives prove correctness and
totality of the compiled operation.
\end{enumerate}
\end{proof}

\begin{corollary}\label{cal:cor:bushy}
The designated halting decision extension constructs a DNC function $g$
which computes no computably bounded DNC function. Externally the stage
construction gives $g\leq_T0'$.
\end{corollary}
\begin{proof}
\leavevmode
\begin{enumerate}[label=\textup{(\roman*)},leftmargin=*,itemsep=3pt,topsep=4pt]
\item \textbf{Construction of the comparison.} Start from the c.e. diagonal forbidden set and grade $2$.
Enumerate all pairs
of numerical indices $(\Gamma,h)$, apply Theorem~\ref{cal:thm:bushy} to each, and
interleave the safe length-extension operation.
Each resulting forbidden set
still has an ordinary c.e. index.
The chosen sequence of such indices is
$0'$-computable, which is sufficient to compute every stage and each eventual
coordinate of the union $g$.
\item \textbf{Forcing construction.} Diagonal avoidance persists because forbidden sets only increase. A negative
requirement outcome puts its convergence set into the permanent forbidden
set; a positive outcome puts a finite diagonalizing computation into the
permanent stem.
Thus every requirement persists.
If a total
$h$-bounded DNC function were computed from $g$, its functional and total
computable bound would appear among the enumerated pairs, contradicting the
corresponding stage.
Bounds of the form $\leq h$ are replaced by $<h+1$.
No enumeration or decision of the set of total computable functions is used.
\end{enumerate}
\end{proof}

This also gives new generating requirement covers on the c.e. condition site:
for each total bound and functional, every condition has the displayed
requirement refinement. On arbitrary indices one can retain the whole output
procedure, with its recorded inactive or successful outcome. The full basic
order uses the bigness decisions carried by the condition.
The richer tree-and-growth-rate conditions used for minimal-degree arguments
in \cite[Section 5]{khan-miller} require additional operations; they are not identified
with this basic order.

\section{Canonical stages and explicit region comparisons}\label{sec:stage-regions}

The abstract region calculus needs a direct connection with the condition sites used for forcing. Every realization has canonical stage regions obtained by restriction to its basic opens. We establish their exactness and then use explicit examples to show when a stage observation is or is not effectively representable.


For any realization $\rho:L\to X$, let $V_p=L_{[p]_J}$ be its open stage and let $v_p:V_p\incl L$.

\begin{theorem}[Canonical exact regions]\label{thm:canonical-regions}
The maps $\rho v_p:V_p\to X$ form an exact region system, and its region realization is $\rho$.
\end{theorem}
\begin{proof}
\leavevmode
\begin{enumerate}[label=\textup{(\roman*)},leftmargin=*,itemsep=3pt,topsep=4pt]
\item \textbf{Refinement comparison.} Refinement is the open inclusion $V_q\incl V_p$.
Covers are jointly surjective because their stage opens join to $[p]_J$.
For an open $U$,
\[
 V_p\leq_XU
 \iff[p]_J\wedge\rho^*U=[p]_J
 \iff[p]_J\leq\rho^*U.
\]
For a family $(U_i)$, this inequality at its join says that $p$ belongs to the saturation of the union of the stable predicates for $\rho^*U_i$.
The saturation formula supplies a cover on each of whose members one of these predicates holds.
This is local join splitting, including its empty case.
The displayed equation identifies the resulting inverse image with $\rho^*$.
\end{enumerate}
\end{proof}

\begin{corollary}[Exactness is derived for every canonical realization]
\label{cor:canonical-exactness-no-axiom}
Every constructed locale realization has an exact region system,
namely its stage restrictions from \cref{thm:canonical-regions}.
Thus exactness is not an additional existence axiom. For a
promise-bearing displayed region system, its own exactness is
established by its comparison maps and their surjectivity as in
\cref{prop:stage-comparison}; equality of raw must theories does
not itself construct such maps.
\end{corollary}
\begin{proof}
\leavevmode
\begin{enumerate}[label=\textup{(\roman*)},leftmargin=*,itemsep=3pt,topsep=4pt]
\item \textbf{Exactness is derived for every canonical realization.} The first assertion is \cref{thm:canonical-regions} applied to
the given realization.
The displayed comparison assertion follows
from \cref{prop:stage-comparison}, and its final qualification
follows from the nonidentical region examples immediately after
that proposition.
\end{enumerate}
\end{proof}

\begin{corollary}[Exact regions and canonical stages have the same must theory]\label{cor:regions-stage-must}
For an exact region realization $\rho_K:L\to X$ and its canonical
open stage $V_p=L_{[p]_J}$, every oracle open $U\in\OO(X)$ satisfies
\[
 K_p\leq_XU\quad\Longleftrightarrow\quad
 [p]_J\leq\rho_K^*U\quad\Longleftrightarrow\quad
 V_p\leq_XU.
\]
These equivalences are uniform in the condition and oracle-open
parameters. They identify the raw must theories, without identifying
the regions or supplying comparison maps between them.
\end{corollary}
\begin{proof}
\leavevmode
\begin{enumerate}[label=\textup{(\roman*)},leftmargin=*,itemsep=3pt,topsep=4pt]
\item \textbf{Exact regions and canonical stages have the same must.} The first equivalence is the exact stage equation of
\cref{thm:regions}.
The second is the canonical-stage containment
equation in \cref{thm:canonical-regions}, applied to $\rho_K$.
Both are equalities of predicates in the full parameter context, so
their composite is uniform.
\end{enumerate}
\end{proof}

\begin{proposition}[Universal property of a canonical stage]\label{prop:stage-universal}
A locale map $h:K\to L$ factors uniquely through $V_p$ precisely when $h^*[p]_J=1$. If $h$ also lies over $X$ through $\rho$, this is the universal factorization in the category of regions over $X$.
\end{proposition}
\begin{proof}
\leavevmode
\begin{enumerate}[label=\textup{(\roman*)},leftmargin=*,itemsep=3pt,topsep=4pt]
\item \textbf{Pullback construction.} The open sublocale $V_p$ is the pullback of truth along the point of the Sierpi\'nski locale classifying $[p]_J$.
Pulling that square back along $h$ gives the open of $K$ classified by $h^*[p]_J$. A factorization exists precisely when this open is all of $K$; it is unique because the open inclusion is monic.
The equation over $X$ follows by composition with $\rho$, so the same factorization has the relative universal property.
\end{enumerate}
\end{proof}

\begin{proposition}[Displayed stage comparison]\label{prop:stage-comparison}
Suppose maps $\lambda_p:V_p\to K_p$ over $X$ commute with refinement, and each $\lambda_p$ is surjective. Then $K_p\leq_XU$ is equivalent to $[p]_J\leq\rho^*U$. The displayed regions are exact and give the same realization. A section $\kappa_p$ of $\lambda_p$ is a concrete proof of the required surjectivity.
\end{proposition}
\begin{proof}
\leavevmode
\begin{enumerate}[label=\textup{(\roman*)},leftmargin=*,itemsep=3pt,topsep=4pt]
\item \textbf{Inverse comparison.} If the displayed region is contained in $U$, so is the canonical stage.
Conversely the equality of inverse images with top can be reflected through the injective map $\lambda_p^*$.
This proves the stage equivalence.
Local splitting follows by applying the canonical stage splitting and the same equivalence at each refinement.
For a covering sieve, the commuting comparisons show that the composites $V_q\to K_q\to K_p$ jointly cover $K_p$, because the $V_q$ cover $V_p$ and $\lambda_p$ is surjective.
Hence the $K_q$ themselves jointly cover $K_p$.
Finally a split epimorphism has injective inverse image.
\end{enumerate}
\end{proof}

Equal containment theories alone do not produce maps $\lambda_p$. For example $X\to X$ and $X+X\to X$ have the same must observations but can have different fibres. A displayed comparison is therefore retained as a map, with its own image and kernel pair.

For a further localization $i:L'\incl L$, the canonical stage at $p$ is $V'_p=V_p\times_LL'$, and its realization is $\rho i$. These are exact by \cref{thm:canonical-regions}. For an original displayed region system, the saturated interpretation becomes the further quotient of $\rho_K^*$. Whether raw containment already expresses this new truth is exactly the cover-exhaustivity calculation for the new covers; it is not encoded by relabelling the old regions.

\subsection{Mathias regions, bounded observations, and an obstruction}\label{sec:mathias-calculation}
We give the promised small construction in $\Eff$, using the reservoir-and-stem form of Mathias forcing \cite{mathias}. A reservoir is a strictly increasing map $r:\N\to\N$. Its range is decidable: to test $n\in\im(r)$ it suffices to inspect $r(0),\ldots,r(n)$, since $r(k)\geq k$. Its tail above $a$ has a canonical increasing enumeration, found by the bounded search for the first index with $r(k)>a$. A condition is $(s,r)$, where $s$ is finite increasing and all reservoir values exceed its last member. Refinement extends $s$ using reservoir values and replaces the reservoir by a subreservoir above the new stem. The latter inclusion is a full internal predicate; only the finite membership tests used below are decidable.

The next-element children are
\[
 (s,r)[k]=(s^\frown r(k),r_{>r(k)}),
\]
and they generate the structural coverage. Pullback stability is explicit. A strengthening with a longer stem already refines the child named by its first new member. A strengthening with the same stem has every one of its children refining a child of the original reservoir. This verifies the required coverage stability.

Let $\Inc(\N)$ be the locale of total strictly increasing sequences, presented by finite increasing-prefix cylinders. A continuous condition ideal has compatible stems, and next-element continuity makes their lengths unbounded. Hence its coordinate relations define a total increasing sequence. For the universal ideal this constructs a realization $\rho_M:L_M\to\Inc(\N)$ by its prefix opens, using \cref{prop:site-classifier}. If a condition belongs to the ideal, every later value of the sequence belongs to its reservoir: choose a condition in the ideal whose stem reaches that value, take a common strengthening with the original condition, and use the refinement clause.

Define $K_{s,r}$ as the sublocale of increasing sequences extending $s$ and having all later values in $\im(r)$. These are geometric coordinate conditions indexed by $\N$. The preceding argument gives a stage factorization
$\lambda_{s,r}:V_{s,r}\to K_{s,r}$.

\begin{lemma}[Uniform frame-valued Mathias stage membership]\label{lem:mathias-open-membership}
Fix a displayed Mathias region $K_{s,r}$ as above, and let $p$ be an arbitrary Mathias condition. For a finite increasing prefix $\tau$ of length $m\geq|s|$ compatible with $(s,r)$, define
\[
 d_m(\tau)=
 \begin{cases}
 (s,r),&m=|s|,\\
 (\tau,r_{>\tau(m-1)}),&m>|s|.
 \end{cases}
\]
Define the finite compatibility predicate
\[
 \operatorname{Compat}_{s,r}(\tau)\iff
 s\preceq\tau\ \wedge\
 \forall k<|\tau|\;\bigl(|s|\leq k
              \Rightarrow\tau(k)\in\im(r)\bigr),
\]
where $\tau$ is a finite increasing sequence. Membership in $\im(r)$
is decidable by bounded search because $r$ is strictly increasing.
Consequently compatibility is decidable, and $d_m(\tau)$ is a valid
Mathias condition for every prefix occurring in the following index.
The family is an open predicate on $K_{s,r}$, uniformly in $p$:
\begin{equation}\label{eq:mathias-stage-open-membership}
 W_p=\bigvee_{m\geq|s|}\ \bigvee_{\substack{
      \tau\in\operatorname{IncFin}(\N),\;|\tau|=m\\
      \operatorname{Compat}_{s,r}(\tau)}}
 \bigl([\tau]|_{K_{s,r}}\wedge!^*[\,p\leq d_m(\tau)\,]\bigr).
\end{equation}
Here $\operatorname{IncFin}(\N)$ is the object of coded finite strictly
increasing sequences. The index contains only compatible prefixes, so
both factors of every displayed meet are defined. No decision of the
non-emptiness of $K_{s,r}$ is made.
The predicates $p\mapsto W_p$ form a frame-valued continuous Mathias ideal satisfying the next-element cover equations, and its induced locale map is a section of $\lambda_{s,r}$.
\end{lemma}
\begin{proof}
\leavevmode
\begin{enumerate}[label=\textup{(\roman*)},leftmargin=*,itemsep=3pt,topsep=4pt]
\item \textbf{Pullback construction.} The scalar truth $[p\leq d_m(\tau)]$ lies in the base topos; its pullback along the terminal locale map is an open of $K_{s,r}$.
Each cylinder $[\tau]|_{K_{s,r}}$ is open, so \eqref{eq:mathias-stage-open-membership} is a join of opens and hence an open.
On a generalized sequence $g$, the defining predicate is
\[
 W_p(g)\quad\Longleftrightarrow\quad
 \exists m\geq|s|\;[p\leq d_m(g|_m)].
\]
The chain $d_m(g|_m)$ is increasing: the next value belongs to the retained reservoir and the new reservoir is the appropriate tail.
\item Thus its generated lower set is inhabited and lower.
Given two witnesses for $W_p$ and $W_q$, their maximum stage witnesses a common strengthening, giving the frame-valued directedness equation.
\item \textbf{Forcing construction.} For a next-element cover of $p$ and a witness $p\leq d_m(g|_m)$, choose a later stage $d_l$ whose stem strictly extends the stem of $p$.
The first newly appended element of that stem names a next-element child $p[k]$ of $p$, and $p[k]\leq d_l$.
This proves that the join of the child membership opens contains $W_p$; the reverse inclusion follows from lowerness.
The equation is uniform in $p$ and in the parameters of the generic sequence.
\item Consequently the frame-valued family satisfies every generating Mathias cover, and the topology-detected-by-an-ideal lemma supplies all generated covers.
The site classifier \cref{prop:site-classifier} constructs the locale map $K_{s,r}\to V_{s,r}$.
Each generating chain contains $(s,r)$ and has the given sequence of stems, so it factors through that stage and its image under $\lambda_{s,r}$ is the identity on the coordinate generators of $K_{s,r}$.
\end{enumerate}
\end{proof}

\begin{lemma}[A section of the Mathias stage comparison]\label{lem:mathias-section}
The map $\lambda_{s,r}$ has a canonical section.
\end{lemma}
\begin{proof}
\leavevmode
\begin{enumerate}[label=\textup{(\roman*)},leftmargin=*,itemsep=3pt,topsep=4pt]
\item \textbf{Ideal and directedness calculation.} The uniform frame-valued construction of \cref{lem:mathias-open-membership} defines, in the relative frame of $K_{s,r}$, an ideal satisfying all the displayed Mathias cover equations.
Applying \cref{prop:site-classifier} gives its unique classifying locale map $\kappa_{s,r}:K_{s,r}\to V_{s,r}$.
Its generic chain starts with $(s,r)$ and has precisely the prefixes of the generic increasing sequence of $K_{s,r}$.
Therefore the sequence reconstructed after applying $\lambda_{s,r}$ is the original one on every coordinate.
Coordinate generation of the increasing-sequence locale gives $\lambda_{s,r}\kappa_{s,r}=\id$.
\end{enumerate}
\end{proof}

\Cref{prop:stage-comparison} now proves exact regional semantics for these Mathias regions. This constructs the comparison at the level of frame homomorphisms and locales.


\subsection{Mathias regions as overt Baire locales}
\label{sec:mathias-overt}
The positive possibility of an arbitrary open can be calculated more
generally than the bounded universal-must predicate. Write $\mathbb B_{\mathrm{loc}}$ for the locale presented by finite
numerical-string cylinders and their immediate-extension covers. This
is a geometric frame presentation, not an identification with a
particular old-object exponential $\N^\N$.

\begin{theorem}[Canonical Baire presentation of a Mathias region]
\label{thm:mathias-baire-region}
For every Mathias condition $p=(s,r)$ there is a canonical locale
isomorphism $K_p\cong\mathbb B_{\mathrm{loc}}$, natural in the condition's numerical
parameters. In particular $K_p$ is overt, with positivity computed on
its Baire cylinder expansion.
\end{theorem}
\begin{proof}
\leavevmode
\begin{enumerate}[label=\textup{(\roman*)},leftmargin=*,itemsep=3pt,topsep=4pt]
\item \textbf{Existence and uniqueness.} For $g\in K_p$, each tail value $g(|s|+n)$ belongs to the range
of the strictly increasing $r$.
Its unique reservoir index $i_n$
is defined by finite search below $g(|s|+n)$, since $r(k)\geq k$.
The $i_n$ are strictly increasing.
Encode them by numerical gaps
\[
 a_0=i_0,\qquad a_{n+1}=i_{n+1}-i_n-1.
\]
Conversely define $i_0=a_0$ and
$i_{n+1}=i_n+1+a_{n+1}$, and then
$g=s^\frown(r(i_0),r(i_1),\ldots)$.
We verify these formulas on the geometric cylinder generators.
\item \textbf{Verification of the equations.} For a finite numerical gap string $a$ of length $k$, form by recursion
$i_0,\ldots,i_{k-1}$ and define the compatible increasing prefix
$\tau_{p,a}=s^\frown(r(i_0),\ldots,r(i_{k-1}))$.
The inverse image
of the Baire cylinder $[a]$ is
\[
   \Phi_p^*[a]=[\tau_{p,a}]\big|_{K_p}.
\]
Conversely, if $\tau$ is an increasing prefix of length $m$, put
$k=\max(0,m-|s|)$.
\item \textbf{Forcing construction.} Its inverse image under the gap-to-sequence
assignment is the open
\[
 \Psi_p^*([\tau]|_{K_p})
   =\bigvee_{a\in\N^k:\,\tau\preceq\tau_{p,a}}[a].
\]
For $k=0$ the only gap prefix is empty, so this is either top or
bottom according to the finite stem compatibility test.
At larger
$k$, equality of the first $k$ reservoir indices is equivalent to
equality of the corresponding numerical gaps, by the mutually inverse
recursions.
\item \textbf{Functoriality and coherence.} These generator equations preserve finite intersections
of compatible cylinders and every next-coordinate cover: all possible
next numerical gaps correspond bijectively to all allowed next
reservoir indices.
They therefore extend to inverse frame maps.
Their composites fix every cylinder by the same finite index
calculation, hence are identities by basis generation.
This proves
the locale isomorphism without appealing to an external point-set
identification.
\item \textbf{Pullback construction.} For overtness, let $J_{\mathbb B_{\mathrm{loc}}}$ be the generated
immediate-extension topology.
The sieves $S$ at $\sigma$ satisfying
\[
 \forall\tau\geq\sigma\;\exists\xi\geq\tau\;S(\xi)
\]
form a topology: maximality and pullback stability are direct, and
local character follows by successive existential refinement.
The
immediate-extension sieve belongs to this topology, so every
$J_{\mathbb B_{\mathrm{loc}}}$-cover is inhabited.
Hence every constant base truth
predicate is stable under these covers.
\item \textbf{Observation preservation.} For a stable Baire open $U$,
put $\operatorname{Pos}(U)=\exists\sigma\,U(\sigma)$.
For any base
truth $v$, the stage expansion and constancy give
\[
 \operatorname{Pos}(U)\leq v
 \quad\Longleftrightarrow\quad U\leq!^*v.
\]
Thus $\operatorname{Pos}\dashv!^*$, exhibiting overtness.
Transfer
this operation along the established frame isomorphism.
\end{enumerate}
\end{proof}

For a finite increasing prefix $\tau$ put
\[
 \operatorname{Compat}(p,\tau)\iff
 \tau\preceq s\ \vee\
 \left(s\preceq\tau\ \wedge\
  \forall |s|\leq k<|\tau|\;\tau(k)\in\operatorname{im}(r)\right).
\]
This predicate is decidable uniformly in $p,\tau$.
On its extent define the refinement constructor
\[
 q(p,\tau)=
 \begin{cases}
   p,&\tau\preceq s,\\
   (\tau,r_{>\tau(|\tau|-1)}),&s\prec\tau.
 \end{cases}
\]
The tail on the second line exists because $\tau$ is nonempty there;
its increasing enumeration is the bounded-search construction already
used in the Mathias section.

\begin{theorem}[Positive Mathias possibility is positive refinement]
\label{thm:mathias-positive-refinement}
For every open $U$ of $\Inc(\N)$, with canonical cylinder expansion
$U=\bigvee_{\tau:E_U(\tau)}[\tau]$, one has the uniform full-truth
identity
\begin{equation}\label{eq:mathias-positive-extension}
 \begin{aligned}
  \May_K(p,U)
   &\iff\exists\tau\,[E_U(\tau)\wedge\operatorname{Compat}(p,\tau)]\\
   &\iff\exists q\geq p\,
          [K_q\leq_XU\ \wedge\ \operatorname{Pos}_{K_q}(1)].
 \end{aligned}
\end{equation}
If the cylinder family $E_U$ is jointly $\Sigma$-classified, the
first expression has a joint $\Sigma$-classifier in $\Eff$.
The witness $q$ in the second line is constructed by the displayed
refinement constructor, not chosen from an arbitrary existential
extension assertion.
\end{theorem}
\begin{proof}
\leavevmode
\begin{enumerate}[label=\textup{(\roman*)},leftmargin=*,itemsep=3pt,topsep=4pt]
\item \textbf{Pullback construction.} A prefix $\tau$ incompatible with $p$ has empty pullback cylinder
in $K_p$: it either contradicts the fixed stem or requests an element
outside the reservoir.
If $\operatorname{Compat}(p,\tau)$, the
constructor $q(p,\tau)$ is a valid strengthening and
$K_{q(p,\tau)}\subseteq[\tau]$.
It has a point given by the displayed
stem followed by its enumerated reservoir tail, so its top open is
positive.
Thus $K_p\cap[\tau]$ is positive precisely when the finite
compatibility predicate holds.
Positivity of $K_p$ preserves the
cylinder join of $U$, proving the first equivalence.
\item \textbf{Inverse comparison.} A compatible cylinder witness $\tau$ yields the indicated positive
strengthening $q(p,\tau)$, whose region is contained in
$[\tau]\subseteq U$, proving the forward implication of the second
equivalence.
\item \textbf{Converse construction.} Conversely a positive refinement contained in $U$
witnesses positive possibility at $p$ by
\cref{prop:positive-extension}.
For effectivity,
$\operatorname{Compat}$ is decidable and the coded finite-prefix
object is $\Sigma$-overt in $\Eff$, by
\cref{prop:enumerated-overt}; the displayed existential is therefore
a uniform $\Sigma$-map.
No finite-use bound on $U$ is required.
\end{enumerate}
\end{proof}


For a finite Boolean assignment $b:2^u$, let $E_b$ say that membership in the range of the sequence below $u$ has pattern $b$. These are clopens: membership of $n$ is decided by the first $n+1$ entries of a strictly increasing sequence. Write $p=(s,r)$ and define the decidable guard
\[
 G_s(n)\iff\exists k<|s|\,[k+1=|s|\ \wedge\ n\leq s(k)].
\]
Thus $G_s$ means ``at or below the last stem member,'' and is false for the empty stem. Define the following explicitly typed finite predicate:
\begin{equation}\label{eq:mathias-compatibility}
\begin{aligned}
 \operatorname{Compat}_{\mathrm{pat}}(p,b)\iff{}&
 \forall n<u\,[G_s(n)\Rightarrow(b(n)=1\iff n\in s)]\\[-2pt]
 &\wedge\ \forall n<u\,[\neg G_s(n)\wedge b(n)=1\Rightarrow n\in\im(r)].
\end{aligned}
\end{equation}
Every test is decidable, including $n\in\im(r)\iff\exists k\leq n\,[r(k)=n]$. Given compatibility, append to $s$ the finitely many requested new members $n<u$ with $\neg G_s(n)$ and $b(n)=1$, in increasing order. Then append the reservoir tail beyond both the stem and all coordinates below $u$. This constructs, uniformly in $p,u,b$ with its compatibility proof, an increasing sequence in $K_p$ with precisely the required membership pattern. Conversely every point of $K_p$ satisfies \eqref{eq:mathias-compatibility}.

\begin{theorem}[Bounded Mathias observation]\label{thm:mathias-bounded}
For $D:2^u\to\Sigma$, put $U_D=\bigvee_b(iD(b)\wedge E_b)$. Then
\[
 [p]\leq\rho_M^*U_D
 \iff
 \forall b:2^u\,[\operatorname{Compat}_{\mathrm{pat}}(p,b)\Rightarrow iD(b)].
\]
The right side has a uniform $\Sigma$-classifier. Consequently all uniformly bounded-use positive computation observations have effective stable traces.
\end{theorem}
\begin{proof}
\leavevmode
\begin{enumerate}[label=\textup{(\roman*)},leftmargin=*,itemsep=3pt,topsep=4pt]
\item \textbf{Bounded Mathias observation.} The clopens $E_b$ partition $K_p$, with the incompatible pieces empty.
If every compatible $b$ satisfies $D(b)$, this partition shows $K_p\leq U_D$.
Conversely, for a compatible $b$, pull containment back along the constructed point with pattern $b$.
It gives $D(b)$, because all other pattern opens are false at that point.
Exact stage semantics gives the equivalence.
\item \textbf{Existence and uniqueness.} By \eqref{eq:mathias-compatibility}, compatibility is a finite decidable predicate: the only reservoir test is the bounded search $\exists k\leq n\,[r(k)=n]$.
Thus the displayed implication into $D(b)$ has the positive classifier $\neg\operatorname{Compat}_{\mathrm{pat}}(p,b)\vee D(b)$.
Its finite conjunction over $2^u$ is positive by \cref{lem:finite-compact}.
\item \textbf{Universal-property argument.} The explicit compatible-point constructor proves the reverse containment implication, and the classifier is uniform in $p,u,D$; it therefore gives a stable effective family. A computation with use bounded by $u$ has a positive termination-and-output predicate on each complete finite assignment $b$, by finite-oracle simulation and the evaluator.
Substitution yields the last assertion.
\end{enumerate}
\end{proof}

\begin{theorem}[An unbounded-use Mathias kernel]\label{thm:mathias-unbounded}
Here use is measured by the integer coordinates queried in the membership oracle. The effective oracle open $U_{\mathrm{even}}$ saying that the least member of the selected sequence is even has no uniformly effective stable trace along $\rho_M$.
\end{theorem}
\begin{proof}
\leavevmode
\begin{enumerate}[label=\textup{(\roman*)},leftmargin=*,itemsep=3pt,topsep=4pt]
\item \textbf{An unbounded-use Mathias kernel.} The open is the halting kernel of the program which queries membership successively, stopping at the first positive answer, and halting exactly when that integer is even.
Each terminating run is finite, so it is an effective positive kernel.
\item \textbf{Verification of the equations.} Let $H_e(n)$ be the decidable test that machine $e$ has halted by stage $n$, with values in $\{0,1\}$, and define
$r_e(n)=2n+H_e(n)$.
This is a uniformly computable strictly increasing reservoir.
Set $p_e=(\varnothing,r_e)$.
If $e$ never halts, every reservoir member is even and $K_{p_e}\leq U_{\mathrm{even}}$.
\item \textbf{Converse construction.} Conversely, for each $n$ the sequence obtained from the tail of $r_e$ starting at $n$ belongs to $K_{p_e}$.
Containment in $U_{\mathrm{even}}$ forces $r_e(n)$ even, hence $H_e(n)=0$.
Thus containment is equivalent to nonhalting.
\item \textbf{Observation preservation.} By exact stage semantics, a uniform effective trace would give a morphism $\N\to\Sigma$ classifying nonhalting.
The finite-hit extraction argument used in \cref{thm:boolean-obstruction} would make nonhalting c.e., a contradiction.
Hence this effective oracle kernel does not belong to the admissible family of the Mathias realization.
\end{enumerate}
\end{proof}

This calculation retains the positive finite transcript semantics and the full unbounded kernel simultaneously. It also provides a concrete reason to keep selected admissibility in the general theory.


\section{Effectively closed classes and a controlled-jump path}\label{sec:closed-class-forcing}

Closed classes provide a setting where finite information can be maintained while one decides increasingly complex oracle questions. We use computable trees and a constructive stage extension to obtain a path with a controlled jump. The proof records which decisions are built into the path rather than postulating a generic choice of branches.


Let $T\subseteq2^{<\N}$ be a computable prefix-closed tree with a node at
every length. A nonempty $\Pi^0_1$ class has such a presentation.
For a computable tree and a string define the semidecidable exhaustion
predicate
\[
 E(T,\sigma)\ \Longleftrightarrow\
 \exists n\geq|\sigma|\ \forall\tau\in2^n\,
          [\sigma\preceq\tau\Rightarrow\tau\notin T].
\]
The inner test is finite. Its search compiles uniformly to a halting
query; the oracle $h$ therefore decides $E$ for all tree codes actually
constructed below.

\begin{lemma}[Extendible-child selection]\label{closed:lem:child}
When $\neg E(T,\sigma)$ holds, the halting decision extension computes
a child $\sigma b$ with $\neg E(T,\sigma b)$.
\end{lemma}
\begin{proof}
\leavevmode
\begin{enumerate}[label=\textup{(\roman*)},leftmargin=*,itemsep=3pt,topsep=4pt]
\item \textbf{Extendible-child selection.} If both children were exhausted, choose their two witnessing levels and
take the larger.
Prefix closure makes both subtrees empty at that level,
which witnesses $E(T,\sigma)$.
Decide their exhaustion predicates and
choose the first non-exhausted child.
The finite empty-level tests and
the selected witnesses are those of the displayed formula.
\end{enumerate}
\end{proof}

For the oracle interpreter, let
$\operatorname{Halt}_e(\sigma)$ mean that the computation on input $e$
halts within $|\sigma|$ steps, asks only coordinates below $|\sigma|$,
and obtains the answers specified by $\sigma$. This is decidable and
persistent under extension. Every genuine halting computation satisfies
this test at a sufficiently long prefix. Define the computable subtree
\[
 S(T,e)=\{\sigma\in T:\neg\operatorname{Halt}_e(\sigma)\}.
\]

\begin{theorem}[An explicit path and its jump program]\label{closed:thm:low-basis}
From the computable tree $T$ above, programs using $h$ construct a path
$G\in[T]$ and a total Boolean function $b$ such that
\[
 b(e)=1\quad\Longleftrightarrow\quad\Phi_e^G(e)\converges.
\]
Both program indices are uniform in the index of $T$. This is the
low-basis construction stated before forming Turing degrees.
\end{theorem}
\begin{proof}
\leavevmode
\begin{enumerate}[label=\textup{(\roman*)},leftmargin=*,itemsep=3pt,topsep=4pt]
\item \textbf{An explicit path and its jump program.} Maintain a computable infinite tree $T_e$ with a stem $\sigma_e$ of
length $e$, beginning with $T_0=T$ and the empty stem.
At stage $e$
construct an index for the computable tree $S_e=S(T_e,e)$ and ask whether
$E(S_e,\sigma_e)$ holds.
\item \textbf{Observation preservation.} If it does not, set $U_e=S_e$ and $b(e)=0$.
If it does, set
$U_e=T_e$ and $b(e)=1$.
In either case $U_e$ is non-exhausted above
$\sigma_e$.
Apply \cref{closed:lem:child} to obtain the first non-exhausted
child $\sigma_{e+1}$, and put
\[
 T_{e+1}=\{\tau\in U_e:
                 \tau\preceq\sigma_{e+1}\text{ or }
                 \sigma_{e+1}\preceq\tau\}.
\]
This is a computable infinite tree with the required stem.
Its index is
obtained by an ordinary compiler from the earlier index and the finite
selected string.
\item \textbf{Oracle interpretation.} Although the sequence of indices uses $h$, each
individual index denotes a computable tree; no extra oracle is hidden in
the exhaustion query.
This gives a total computable state-update
operation relative to $h$.
\item \textbf{Observation preservation.} Let $G=\bigcup_e\sigma_e$.
For any fixed $e$, all sufficiently late
stems lie in $T_e$.
Prefix closure therefore gives every prefix of $G$
in $T_e$, and in particular $G\in[T]$.
If $b(e)=0$, every prefix of $G$ lies in $S_e$. A finite halting run
would violate this at a sufficiently long prefix, so
$\Phi_e^G(e)\diverges$.
If $b(e)=1$, the positive exhaustion answer supplies a level $n$ with
no node of $S_e$ above $\sigma_e$.
\item \textbf{Forcing construction.} The prefix $G\mathbin{\upharpoonright}n$
is in $T_e$ and extends $\sigma_e$, so it satisfies
$\operatorname{Halt}_e$ and the computation halts.
Running the construction through stage $e$ computes $b(e)$ using $h$;
running it through the required stem length computes $G(n)$.
These are
the two uniform programs asserted in the theorem.
\end{enumerate}
\end{proof}

All assertions about the new path use its actual finite-run checker.
The proof is thus valid in the constructed halting decision extension as
well as in the classical oracle interpretation. In particular, it does
not appeal to a point of the original locale, countable choice, or an
unproved preservation of universal arithmetic truth. This is the
classical low-basis phenomenon of Jockusch--Soare \cite{jockusch-soare}, with its
path and jump reductions exposed as programs. The proof above supplies its internal realization by the decision
oracle and the retained finite-run interpreter.

Here $[T]$ in the extension is the native path object constructed
in \cref{seq:cor:paths}: its elements are total prefix ideals
contained in the tree, equivalently points of the base-changed
finite-cylinder locale. The algorithm above constructs its
compatible prefixes and hence its native sequence by
\cref{seq:thm:prefixes}. The comparison from the ground path object
is induced by \cref{seq:prop:transport}; it is not used to supply
this path. For a computable tree with no computable path, the ground
path object is empty and its sheafification remains empty, whereas
the displayed oracle program constructs a path in the native object.

\paragraph{The closed-class condition site.}
In the halting decision extension, take conditions $(T,\sigma)$ with
$T$ a computable-tree index, every node of $T$ comparable
with $\sigma$, and $\neg E(T,\sigma)$. Order them by subtree
inclusion and stem extension. This is the tree presentation
of forcing by nonempty effectively closed subclasses of Cantor space,
with a displayed stem; \cref{cal:prop:closed-codes} gives the exact
presentation comparison. For each $e$ define the two upper sets
\[
 \begin{split}
 D_e^0&=\{(T,\sigma):\forall\tau\in T\,
                              \neg\operatorname{Halt}_e(\tau)\},\\
 D_e^1&=\{(T,\sigma):\exists n\geq|\sigma|\ 
             \forall\tau\in T\cap2^n\,
                              \operatorname{Halt}_e(\tau)\}.
 \end{split}
\]
The stage construction proves $D_e^0\cup D_e^1$ dense below every
condition. The positive case uses the exhaustion certificate; the
negative case replaces $T$ by $S(T,e)$. Both sets persist under
strengthening: a smaller tree preserves the negative clause, and a
larger level beyond the new stem preserves the positive clause by
persistence of the finite halting test. The length sets are dense by
\cref{closed:lem:child}. Generate the coverage from these dense refinements using \cref{prop:site-nucleus}.
The directed union of its generic stems is a total path, and the
$D_e$ covers give its halting decisions exactly as in the theorem.
The constructed chain meets every one of these generators, so it is
a point of this explicitly presented relative locale. Cover continuity is with respect to this explicitly generated coverage.

To verify cover continuity without an additional genericity axiom,
observe that each generating dense set is upper. For a generating
cover at $p$ in the filter, directedness combines $p$ with a filter
member meeting that dense set, giving a member of the restricted
cover. For intersections take common stronger members;
for transitive composites first meet the outer cover and then the
cover at that member. These operations preserve the property of being
met by the filter, proving it for the generated Grothendieck coverage.

\begin{proposition}[Exact closed-class presentation comparison]
\label{cal:prop:closed-codes}
In the classical calibration, the literal computable-tree condition order and
the order of represented nonempty $\Pi^0_1$ subclasses have canonically
equivalent dense sheaf topoi, preserving the generic path.
\end{proposition}
\begin{proof}
\leavevmode
\begin{enumerate}[label=\textup{(\roman*)},leftmargin=*,itemsep=3pt,topsep=4pt]
\item \textbf{Observation preservation.} Send $(T,\sigma)$ to $[T]\cap[\sigma]$. A stronger represented class with
representative $(S,\tau)$ lifts to $(T\cap S,\rho)$, where $\rho$ is the
longer stem.
Nonemptiness makes the stems comparable and puts $\rho$ in both
trees.
The intersection represents the smaller class and is a literal
strengthening.
The same operation contracts two representatives of the same
class.
The target consists of represented classes, so the projection is
surjective.
Apply \cref{cal:lem:codes}; the common finite prefixes identify
the generic path.
\item \textbf{Localization and coverage.} The length and jump-decision covers above are dense in the
halting decision base by their explicit constructors, so
\cref{cal:thm:localize} also applies to that requirement site.
\end{enumerate}
\end{proof}

\section{Restrained functional forcing}\label{sec:restrained-forcing}

A partial functional can be built by finite axioms only when later axioms respect previously committed uses and outputs. We make those restraints explicit in a finite condition order. The resulting extension and use-growth operations are what connect the functional forcing to the general coverage and evaluation calculus.


A finite computation entry, traditionally called a functional axiom,
is a triple $(n,b,\sigma)$ with $b<2$ and $\sigma\in2^{<\N}$.
A finite use-monotone functional $\Phi$ is a finite set of such triples
satisfying the following finite tests:
\begin{enumerate}
\item Compatible oracle strings for the same input must be the identical
axiom (same string and output).
\item If the string of one axiom properly extends that of another, its
input is larger.
\item An axiom for input $n$ has, along its string, an axiom for every
smaller input.
\end{enumerate}
These are the finite version of \cite[Definition 1.1.3]{slaman-woodin}.
A condition is $(\Phi,F)$, where $F$ is a finite list of protected
binary sequences. Strengthening extends $\Phi$, gives every new axiom
use greater than all old uses, extends the protected collection, and
leaves $\Phi(Y)$ unchanged for every previously protected $Y$.
One may append to the list and quotient by its represented finite set;
checking that new axioms respect a displayed list needs only finite
prefix queries. It never requires deciding equality of arbitrary reals.

\begin{theorem}[Constructed extension away from restraints]\label{ks:thm:extension}
For a condition $(\Phi,F)$ and a sequence $X$ with supplied finite
apartness witnesses from the members of $F$, one can add an axiom
giving either chosen output at the next undefined input along $X$.
The result is a valid strengthening. Repetition realizes any prescribed
finite continuation along $X$ while retaining all the restraints.
\end{theorem}
\begin{proof}
\leavevmode
\begin{enumerate}[label=\textup{(\roman*)},leftmargin=*,itemsep=3pt,topsep=4pt]
\item \textbf{Constructed extension away from restraints.} For each $Y\in F$ let $a_Y$ be a coordinate at which $X$ and $Y$
differ.
Select $L$ greater than all old uses and all $a_Y$; put
$\sigma=X\mathbin{\upharpoonright}L$.
Querying these finitely many
bits computes $\sigma$.
The old axioms applicable to $X$ are a finite
initial segment of inputs, by the third use-monotonicity clause.
Let
$n$ be its length and let $b$ be the desired next output.
Add
$(n,b,\sigma)$.
\item \textbf{Observation preservation.} Every old axiom compatible with $\sigma$ has shorter use and is
applicable to $X$, so its input is less than $n$.
Thus no functional
conflict is created and the proper-prefix rule holds.
The old initial
segment supplies all smaller inputs required by the third clause.
For every protected $Y$, the coordinate $a_Y<L$ makes $\sigma$
incompatible with $Y$, so no new computation is added there.
The new
use is larger than every old use.
These verify every strengthening
clause.
Repeating the same finite construction proves the final claim.
\end{enumerate}
\end{proof}

The apartness witnesses in this statement are the exact domain of the
extension operation, not a new global separation axiom. They can also be
found by search from a supplied positive apartness assertion
$\exists n\,[X(n)\ne Y(n)]$. No assertion that distinct binary
sequences are constructively apart is used.

There is also a successor without a supplied unprotected sequence.
Given $(\Phi,F)$ and a desired use bound $m$, choose
$L>m$ larger than all old uses with $2^L>|F|$. Query the first $L$
bits of the finitely many members of $F$. Their prefixes occupy at
most $|F|$ strings in $2^L$, so choose the first unused string
$\sigma$. The finite old computations along $\sigma$ form an initial
segment of inputs. Add output $0$ at its next input using $\sigma$.
The checks in \cref{ks:thm:extension} apply verbatim and prove a valid extension
whose new use exceeds $m$. Thus both ``protect this supplied real''
and ``increase the finite functional beyond this bound'' are genuine
dense refinement operations. Their generated coverage supplies a
functional-forcing site, and the union calculation below gives its
generic semantics. Stronger classical genericity requirements can use
the dense coverage in the classical calibration; they are not needed
for these two constructive operations.

\begin{proposition}[The generic functional and its domain]\label{ks:prop:generic}
The directed union of the finite functional components is a functional.
In the site generated by the use-growth and protection covers above,
membership of a finite axiom in that union is decidable. Its evaluation
is the native finite-hit observation
\[
 \Phi_G(X)(n)=b\quad\Longleftrightarrow\quad
 \exists\sigma\preceq X\ [(n,b,\sigma)\in\Phi_G].
\]
The order with restraints maps to the ideal domain of finite
use-monotone functionals by forgetting the protected list. This map
preserves the displayed evaluation. The restraints remain part of the
condition presentation.
\end{proposition}
\begin{proof}
\leavevmode
\begin{enumerate}[label=\textup{(\roman*)},leftmargin=*,itemsep=3pt,topsep=4pt]
\item \textbf{The generic functional and its domain.} Any pair of potentially conflicting axioms is already contained in one
common stronger finite condition, where the finite functional tests
exclude the conflict.
An axiom and its finitely many smaller-input
witnesses occur in a finite condition, so use monotonicity is also
preserved.
Evaluation is an existential union of finite-prefix tests;
the finite axiom witnessing it belongs to some condition.
\item \textbf{Converse construction.} Conversely
such a finite witness persists in every stronger condition.
For an axiom $a=(n,b,\sigma)$, cover by conditions whose maximum use
exceeds $|\sigma|$.
In each such condition, finite membership decides
whether $a$ has already appeared.
If it has not, the end-extension
rule prohibits its appearance in any later condition.
These locally
compatible Boolean decisions therefore glue to a decision of
$a\in\Phi_G$.
Invalid triples are uniformly rejected.
\item \textbf{Evaluation and transposition.} In the
displayed evaluation equation, enumeration of numerical string codes
now performs a decidable test of this axiom-membership map and of the
finitely many prefix bits of $X$.
Its hit stream proves precisely the
native $\Sigma$-classification asserted.
Finite valid functionals under inclusion form a basis; an ideal gives
their union.
The finite consequences of any finitely many axioms can
be included from the original finite witnesses, so the union has the
same directed finite representation.
\item \textbf{Ideal and directedness calculation.} Forgetting restraints is monotone
and hence extends to ideal completion.
The finite-witness evaluation
equation proves the observation assertion.
\end{enumerate}
\end{proof}

There is a necessary domain restriction here. For each ground sequence
$Y$, appending $Y$ to the protected list is always a strengthening.
The set of conditions protecting $Y$ is therefore dense. Once met,
the final functional on $Y$ is frozen at the finite value it had then.
Thus a generic filter meeting these sets does \emph{not} make
$\Phi_G(Y)$ total for every ground $Y$. The unprotected-sequence
extension in \cref{ks:thm:extension} and this freezing property are different
operations. Their coexistence is part of Kumabe--Slaman forcing, not a
defect that can be removed by a totality assertion. The full classical
requirement argument is in \cite[Section 8.2]{slaman-woodin}; the extension and
domain calculations above are the pieces needed to implement its
condition presentation in the effective-domain framework.

\begin{proposition}[Protected-list and protected-set presentations]
\label{cal:prop:ks-codes}
The finite-list implementation of protected reals and the finite-set order of
\cite[Definition 8.2.2]{slaman-woodin} have canonically equivalent dense
sheaf topoi. Their generic finite-axiom relations agree.
\end{proposition}
\begin{proof}
\leavevmode
\begin{enumerate}[label=\textup{(\roman*)},leftmargin=*,itemsep=3pt,topsep=4pt]
\item \textbf{Protected-list and protected-set presentations.} Send $(\Phi,\ell)$ to $(\Phi,\im\ell)$.
To lift a stronger protected set,
take any local finite enumeration of it and append the old list.
To contract
two representatives of the same set and functional, concatenate their lists.
The functional is unchanged in a fibre, so its restraint conditions remain
valid.
These are the lifting and contraction operations of
\cref{cal:lem:codes}; finite lists cover the finite-set object.
Forgetting their order and repetitions does not change the finite functional,
so its union and evaluation relation are preserved.
\item \textbf{Localization and coverage.} Protection and use-growth
are dense by the explicit extension constructions above.
Thus
\cref{cal:thm:localize} connects their generated coverage to the dense one.
\end{enumerate}
\end{proof}

\section{Intensional enhancement: finite regional witnesses}
\label{sec:intensional-regions}


The regional must and may predicates are full frame calculations.
A positive witness description is additional data: it identifies
which finite search certified a particular positive observation.
The examples above supply three such displayed description objects:
finite binary covers, finite bushy trees with labelled leaves, and
Mathias-compatible finite assignment patterns. They need not be
identified with each other or with their semantic support.

\begin{proposition}[Sound finite witnesses and their positive support]
\label{prop:intensional-region-witnesses}
Let $C\to P\times T$ be any of these displayed finite-certificate
families. Suppose its verifier is the conjunction of a decidable
structural check with finitely many $\Sigma$-classified leaf tests,
and that its code fibres have the overt search constructed in the
corresponding example. The accepted-certificate support
\[
 W(p,t)\iff\exists c:C_{p,t}\,\operatorname{Accept}(p,t,c)
\]
is jointly $\Sigma$-classified. If every accepted certificate
covers its root by regions contained in the target open $U_t$,
then $W(p,t)$ implies $[p]_J\leq\rho^*U_t$. If the example's
completeness calculation converts this latter truth into a
certificate, the equivalence holds and $W$ is the effective
stable trace on that family.
\end{proposition}
\begin{proof}
\leavevmode
\begin{enumerate}[label=\textup{(\roman*)},leftmargin=*,itemsep=3pt,topsep=4pt]
\item \textbf{Observation preservation.} Finite conjunction preserves $\Sigma$ by
\cref{lem:finite-compact}; overt existential search over the
specified certificate codes preserves it by
\cref{prop:enumerated-overt,prop:positive-logic}.
Soundness turns the leaf membership checks into covering opens
lying below $\rho^*U_t$; their join contains $[p]_J$.
The
converse uses the explicitly exhibited normalization or compatible
finite assignment construction for the selected example, not a
choice from an arbitrary cover.
\item \textbf{Pullback construction.} The two directions identify the
$\Sigma$-classifier with the stable forcing truth, and the
pullback of \cref{def:stable-mark} supplies its unique lift.
\end{enumerate}
\end{proof}

For Cohen, the certificate is a finite binary level and its leaf
checks; completeness is \cref{thm:cohen-saturation}. For bushy
search it is an actual finite tree and its leaf labels, with
grafting justified in \cref{prop:bushy-positive}. For bounded
Mathias computation it is a complete finite membership pattern
and its compatibility verification, with the effective universal
check in \cref{thm:mathias-bounded}. These witnesses remain
available as finite execution and witness certificates, as constructed in \cref{thm:operational-certificates}. Their existence does not decide negative forcing truths:
\cref{thm:mathias-unbounded} provides a concrete effective kernel
for which no such globally positive stage trace exists.

\part{Histories, promises, and least reconstruction}\label{part:histories}

The image of a realization records which oracle is obtained, but not necessarily the conditions or promises that produced it. Its relative fibre retains precisely those histories. We build the canonical history site, ask when a least history can be selected, and compare different kinds of committed and residual information in partition, Mathias, and perfect-tree examples.


\section{Generalized fibres and canonical history sites}\label{sec:histories}


For a realization $\rho:L\to X$, the oracle does not generally determine its forcing history. The appropriate fibre is a locale, retaining generalized histories and their open observations.

\begin{definition}[History locale]\label{def:history}
For a locale map $x:T\to X$, define $\Hist_\rho(x)$ by
\[
\begin{tikzcd}
 \Hist_\rho(x)\ar[r,"h_x"]\ar[d,"p_x"']&L\ar[d,"\rho"]\\
 T\ar[r,"x"']&X.
\end{tikzcd}
\]
It is regarded as a locale over $T$.
\end{definition}
Here $T$ can be any internal locale. When a parameter is an object of $\E$, its discrete locale gives the corresponding special case. Interpreting the diagram internally in $\Sh_\E(T)$ turns $x$ into a point and the history locale into its ordinary relative fibre.

\begin{proposition}[Reindexing and fibre interpretation]\label{prop:history-reindex}
For $u:T'\to T$ there is a canonical isomorphism
\[
 T'\times_T\Hist_\rho(x)\cong\Hist_\rho(xu).
\]
These isomorphisms satisfy identity and composition coherence. A map $S\to\Hist_\rho(x)$ over $T$ is exactly a map $h:S\to L$ with $\rho h=xv$, where $v:S\to T$ is its structure map.
\end{proposition}
\begin{proof}
\leavevmode
\begin{enumerate}[label=\textup{(\roman*)},leftmargin=*,itemsep=3pt,topsep=4pt]
\item \textbf{Reindexing and fibre interpretation.} Both sides of the first isomorphism are the pullback of $\rho$ along $xu$.
Their universal maps preserve the projections to $T'$ and $L$, so they are inverse.
The same uniqueness proves coherence for two successive changes of parameter.
The final statement is precisely the pullback universal property.
\end{enumerate}
\end{proof}

\begin{lemma}[Relative fibre inclusion]\label{lem:relative-fibre-basis}
Internally over $T$, the fibre is a sublocale of the pulled-back $L$. If $(d_p)$ is a basis of $L$, the inverse images of its pulled-back generators form a relative join basis for the fibre.
\end{lemma}
\begin{proof}
\leavevmode
\begin{enumerate}[label=\textup{(\roman*)},leftmargin=*,itemsep=3pt,topsep=4pt]
\item \textbf{Relative fibre inclusion.} In $\Sh_\E(T)$ the map $x$ is a point of the pulled-back $X$.
Its inverse image on frames has the terminal inverse image as a section, since the composite of the point with the terminal map is the identity.
Thus it is surjective and the point is a sublocale inclusion.
Pulling this inclusion back along the pulled-back $\rho$ gives the fibre inclusion.
\item \textbf{Verification of the equations.} A sublocale inverse image is a surjective frame map.
Apply it to the join expansion of every open in the ambient pulled-back locale.
The images of the basis generators therefore generate the fibre frame.
This is a relative statement: coefficients are truth values in the base $\Sh_\E(T)$.
\end{enumerate}
\end{proof}

\begin{theorem}[Canonical history site]\label{thm:history-site}
Let $a_p$ be the relative image in the fibre frame of the stage $d_p$. The preorder
\[
 p\preceq_x q\iff a_q\leq a_p
\]
with its semantic basis coverage presents $\Hist_\rho(x)$ internally over $T$. Every displayed history basis satisfying the basis and cover comparison equations gives the same relative locale, with its indicated maps.
\end{theorem}
\begin{proof}
\leavevmode
\begin{enumerate}[label=\textup{(\roman*)},leftmargin=*,itemsep=3pt,topsep=4pt]
\item \textbf{Canonical history site.} \Cref{lem:relative-fibre-basis} supplies join density.
Apply \cref{cor:canonical-basis} inside $\Sh_\E(T)$.
For a displayed history basis, apply \cref{prop:displayed-basis} in that same base.
Its explicit inverse maps identify the selected history stages and therefore the maps represented by their observations.
\end{enumerate}
\end{proof}

The relative qualification prevents a loss of parameter information. A basis pulled back from $L$ generates the relative frame over $T$; its coefficients include opens of $T$. It is not asserted to be an absolute basis over $\E$ without those coefficients.

\begin{proposition}[History as a geometric model]\label{prop:history-model}
If $L$ is presented by continuous ideals and $X$ by oracle coordinates, the history locale over a point $x$ is presented by the ideal equations and the equations identifying each interpreted oracle coordinate with its value at $x$.
\end{proposition}
\begin{proof}
\leavevmode
\begin{enumerate}[label=\textup{(\roman*)},leftmargin=*,itemsep=3pt,topsep=4pt]
\item \textbf{Ideal and directedness calculation.} In any test locale over $T$, a model of the ideal equations determines a map to the pulled-back $L$ by \cref{prop:site-classifier}.
The additional coordinate equations say that its composite with $\rho$ equals $x$: equality of the coordinate inverse images implies equality on all opens by their generating property.
The pullback universal property of \cref{prop:history-reindex} therefore gives a unique map to the history locale.
\item \textbf{Converse construction.} Conversely such a map gives these relations by pulling back the universal model.
The two constructions fix the membership relations and are inverse.
Since this is proved for every test locale, it identifies the locale presentations, without any appeal to enough points.
\end{enumerate}
\end{proof}


The relative fibre gives a locale of histories, not a selected history. We can therefore ask a separate universal question: when does a section choose the least history over every realized oracle? The order-enriched adjunction provides the answer.

\section{Sections and least histories}\label{sec:least-histories}

A realization may have many histories above one oracle. Selecting a section need not select a least one, and a least section is an order-theoretic property rather than mere set-theoretic choice. We characterize it by an adjunction and construct its consequences for observation and reconstruction.


Order locale maps by $f\leq g$ when $f^*\leq g^*$ pointwise. This convention agrees with inclusion of ideal models for the locales presented above. Composition is monotone in both variables.

\begin{theorem}[Least-history adjunction]\label{thm:least-history}
Let $\kappa:X\to L$ be a section of $\rho$. The following are equivalent:
\begin{enumerate}
\item for every generalized oracle $x:T\to X$ and every $h:T\to L$ with $\rho h=x$, one has $\kappa x\leq h$;
\item $\kappa\rho\leq\id_L$;
\item $\kappa\dashv\rho$ in this locally ordered category, with identity unit $\rho\kappa=\id_X$.
\end{enumerate}
The statement is natural in the generalized oracle.
\end{theorem}
\begin{proof}
\leavevmode
\begin{enumerate}[label=\textup{(\roman*)},leftmargin=*,itemsep=3pt,topsep=4pt]
\item \textbf{Unit and identity laws.} Apply the first clause to $T=L$, $x=\rho$, and $h=\id_L$ to obtain the second.
Conversely compose the second inequality with any $h$ to get
$\kappa x=\kappa\rho h\leq h$.
The two inequalities of an adjunction are $\id_X\leq\rho\kappa$ and $\kappa\rho\leq\id_L$; the first is the given identity, proving the equivalence with the third clause.
Precomposition with a map of parameter locales preserves the equations and inequalities, proving naturality.

The universal property is represented by the commutative diagram
\[
\begin{tikzcd}[column sep=large,row sep=large]
X\ar[r,"\kappa"]\ar[dr,"\id_X"']&L\ar[d,"\rho"]\\ &X
\end{tikzcd}
\]
\end{enumerate}
\end{proof}

This is a characterization of least reconstruction, not a selection from an arbitrary inhabited fibre. When a least section is constructed, the adjunction packages exactly its fibrewise universal property. Its effective behavior is proved through the marking maps for $\kappa^*$; the adjunction alone has no effective classification content.


\begin{theorem}[Universal oracle-derived interior and criterion for
least histories]\label{thm:least-history-interior}
For any locale map $\rho:L\to X$, its inverse-image frame map has a
right adjoint
\[
 \rho_*(U)=\bigvee\{V\in\OO(X)\mid\rho^*V\leq U\}.
\]
The map $\rho^*\rho_*:\OO(L)\to\OO(L)$ is a deflationary,
idempotent, monotone interior operator. It sends $U$ to the
largest oracle-derived open beneath $U$. Moreover the following
are equivalent:
\begin{enumerate}
\item $\rho$ admits a least-history section $\kappa:X\to L$;
\item $\rho$ is a surjective locale map and $\rho_*$ preserves
arbitrary joins.
\end{enumerate}
When they hold, $\rho_*=\kappa^*$ and the reconstruction is
unique. Thus the existence of a least section is characterized by
the \emph{actual frame right adjoint}, not inferred merely from
inhabited fibres.
\end{theorem}
\begin{proof}
\leavevmode
\begin{enumerate}[label=\textup{(\roman*)},leftmargin=*,itemsep=3pt,topsep=4pt]
\item \textbf{Universal oracle-derived interior and criterion for.} The formula for $\rho_*$ is the right adjoint to $\rho^*$ in
internal complete lattices.
Its counit says
$\rho^*\rho_*U\leq U$.
Every $\rho^*V\leq U$ lies below
$\rho^*\rho_*U$ by the defining adjunction.
The triangle
identities imply $(\rho^*\rho_*)^2=\rho^*\rho_*$, and the
construction is monotone.
\item \textbf{Adjunction and universal property.} If $\kappa$ is a least section, then
$\kappa^*\rho^*=\id$ and $\rho^*\kappa^*\leq\id$ by
\cref{thm:least-history}.
For $V\in\OO(X)$ and $U\in\OO(L)$,
\[
 \rho^*V\leq U
 \iff V=\kappa^*\rho^*V\leq\kappa^*U
 \iff V\leq\kappa^*U.
\]
Hence $\kappa^*$ is the right adjoint, so $\rho_*=\kappa^*$
and preserves joins.
Moreover $\rho^*$ has a left inverse,
so it is order-reflecting and $\rho$ is surjective.
\item \textbf{Converse construction.} Conversely assume $\rho^*$ is injective and $\rho_*$ preserves
all joins.
As a right adjoint $\rho_*$ preserves all meets; thus
it is a frame homomorphism and defines a locale map $\kappa$.
Since $\rho^*$ is an order embedding, the unit and counit of the
adjunction imply $\rho_*\rho^*=\id$.
Consequently
$\rho\kappa=\id_X$ and the counit is
$\kappa\rho\leq\id_L$.
The latter is exactly the least-history
condition.
The right adjoint to $\rho^*$ is unique, hence so is the corresponding locale map.
\end{enumerate}
\end{proof}

\begin{example}[Surjectivity does not construct a least history]
\label{ex:surjection-no-least-history}
For the discrete locale map $\rho:2\to1$, the right adjoint is
universal quantification over the two points. It satisfies
\[
 \rho_*\{0\}=\rho_*\{1\}=0,
 \qquad \rho_*(\{0\}\vee\{1\})=1.
\]
It does not preserve joins, although $\rho$ is surjective and has
sections. By \cref{thm:least-history-interior} it has no least-history
section. The following examples instead construct their least
histories and verify the requisite frame equations.
\end{example}


\subsection{Canonical positive domains for history observation}
\label{sec:history-effective}
Work in the relative topos $\mathcal T=\Sh_\E(T)$ and let
$F_x=\Hist_\rho(x)$ be the relative history locale. Let $H$ be its
canonical semantic stage preorder, with join-dense opens
$b_h\in\OO(F_x)$ and semantic coverage as in
\cref{thm:history-site}. Let $\iota_\Theta:\Theta\incl\Omega_\mathcal T$
be a selected positive truth object with finite conjunction and the
restricted positive quantifiers of \cref{thm:maximal-positive-quantifiers}.
No global overtness of $H$ or effective membership relation is assumed.

Write $\ContIdl(H)$ for the object of inhabited, lower, directed,
cover-continuous ideals of $H$. Its universal membership is
$m(I,h)=I(h)$. The full frame-valued history observation is
\begin{equation}\label{eq:history-full-evaluation}
 \ev_x(I,U)=\exists h:H\,[I(h)\wedge b_h\leq U],
 \qquad I:\ContIdl(H),\ U:\OO(F_x).
\end{equation}
The basis expansion proves that this is the truth of $U$ at the
locale point classified by $I$: its inverse-image frame homomorphism
takes each stage $b_h$ to $I(h)$, and the displayed join is the
result of evaluating the expansion of $U$.

\begin{theorem}[Maximal positive history-observation domain]
\label{thm:maximal-history-observation}
There is a canonical pullback
\[
\begin{tikzcd}
 D_x^{\mathrm{out}}\ar[r,"\Obs_x^{\mathrm{out}}"]\ar[d,hook]&
 \Theta\ar[d,hook,"\iota_\Theta"]\\
 \ContIdl(H)\times\OO(F_x)\ar[r,"\ev_x"']&\Omega_\mathcal T.
\end{tikzcd}
\]
It is the largest parameter-stable subobject of history/open pairs
whose actual observation has a positive classifier. Its classifier
$\Obs_x^{\mathrm{out}}$ is unique and defined jointly in both
parameters.
\end{theorem}
\begin{proof}
\leavevmode
\begin{enumerate}[label=\textup{(\roman*)},leftmargin=*,itemsep=3pt,topsep=4pt]
\item \textbf{Maximal positive history-observation domain.} Apply \cref{thm:admissibility-coreflection} in $\mathcal T$ to $\iota_\Theta$ and the family $\ev_x$.
The counit is the displayed domain mono; its universal property gives the unique positive classifier jointly in both parameters, maximality and substitution coherence.
\end{enumerate}
\end{proof}

There is also a canonical smaller domain on which the particular
\emph{positive overt-search formula} can be evaluated. Define its
positive membership and positive stage-trace objects by the pullbacks
\[
 I_\Theta=\ContIdl(H)\times_{\Omega^H}\Theta^H,
 \qquad
 O_\Theta=\OO(F_x)\times_{\Omega^H}\Theta^H,
\]
using respectively $I\mapsto(h\mapsto I(h))$ and
$U\mapsto(h\mapsto b_h\leq U)$. Both right-hand arrows to
$\Omega^H$ are the exponentiated inclusion $\iota_\Theta^H$.
On $I_\Theta\times O_\Theta$ their positive classifiers give
\[
 w(I,U)(h)=I_\Theta(I,h)\wedge_\Theta O_\Theta(U,h)
 \quad\in\Theta^H.
\]
Let $Q_{!_H}^{\exists}\incl\Theta^H$ be the maximal positive
effectivity domain of existential quantification along $!_H:H\to1$,
constructed in \cref{thm:maximal-positive-quantifiers}. Define
\begin{equation}\label{eq:history-search-domain}
 D_x^{\mathrm{search}}=
 (I_\Theta\times O_\Theta)\times_{\Theta^H}Q_{!_H}^{\exists}.
\end{equation}

\begin{theorem}[Computed positive-search history domain]
\label{thm:history-search-domain}
There is a canonical mono
$D_x^{\mathrm{search}}\incl D_x^{\mathrm{out}}$. On it the
actual history observation has the positive classifier
\begin{equation}\label{eq:history-observation}
 \Obs_{\mathrm{Hist}}(I,U)
   =\exists_{h:H}^{\Theta}
        \bigl(I_\Theta(I,h)\wedge_\Theta O_\Theta(U,h)\bigr).
\end{equation}
Every family on which positive membership, positive stage traces,
and then the displayed positive existential operation are all defined
factors uniquely through $D_x^{\mathrm{search}}$. This construction
is stable under parameter substitution and does not postulate that any
of these three partial effective operations is total.
\end{theorem}
\begin{proof}
\leavevmode
\begin{enumerate}[label=\textup{(\roman*)},leftmargin=*,itemsep=3pt,topsep=4pt]
\item \textbf{Computed positive-search history domain.} On $I_\Theta\times O_\Theta$, positive conjunction constructs
$w$ uniformly.
The last pullback in
\eqref{eq:history-search-domain} supplies exactly the classifier of
its existential truth.
After composing with $\iota_\Theta$,
\eqref{eq:history-observation} equals the underlying predicate
\eqref{eq:history-full-evaluation}.
The universal property of
$D_x^{\mathrm{out}}$ therefore induces the required mono.
\item \textbf{Pullback construction.} Each of the first two positive-factorization conditions is the
universal property of its defining pullback; the final one is the
universal property of $Q_{!_H}^{\exists}$.
Pullback pasting proves
maximality for this execution order.
Beck--Chevalley for the full
existential, together with uniqueness of positive classifiers, proves
parameter stability.
\end{enumerate}
\end{proof}

\begin{corollary}[Uniform effective history observation]
\label{thm:history-effective-observation}
If the relative history basis is $\Theta$-overt and a selected family
of continuous history ideals has jointly $\Theta$-classified
membership, while a selected family of fibre opens has jointly
$\Theta$-classified traces
\[
 \iota_\Theta a_A(h,k)\iff b_h\leq A(k),
\]
then their entire product factors through
$D_x^{\mathrm{search}}$, and its effective observation is
\[
 \Obs_{\mathrm{Hist}}(I,k)=
 \exists_{h:H}^{\Theta}\bigl(I(h)\wedge_\Theta a_A(h,k)\bigr).
\]
The underlying truth is the value of $A(k)$ at the history
classified by $I$.
\end{corollary}
\begin{proof}
\leavevmode
\begin{enumerate}[label=\textup{(\roman*)},leftmargin=*,itemsep=3pt,topsep=4pt]
\item \textbf{Uniform effective history observation.} The two jointly positive classifiers supply maps into $I_\Theta$
and $O_\Theta$.
Overtness of $H$ makes
$Q_{!_H}^{\exists}\to\Theta^H$ an isomorphism, so the product map
factors through $D_x^{\mathrm{search}}$.
Apply
\cref{thm:history-search-domain}.
The assertion about full truth is
\eqref{eq:history-full-evaluation}.
\end{enumerate}
\end{proof}

\begin{proposition}[Effective least-history comparison]
\label{prop:effective-least-history}
Let $\kappa:X\to L$ be a least-history section of
\cref{thm:least-history}. Its relative history membership and an open
family have a uniform positive observation on the pullback of
$D_x^{\mathrm{out}}$ along the section. On the further
$D_x^{\mathrm{search}}$ domain the value is given by
\eqref{eq:history-observation}. If its membership and the open traces
have the displayed positive classifiers and $H$ is $\Theta$-overt,
this observation is total on the selected families. The leastness
inequality is preserved by geometric base change of the section and
history comparison maps.
\end{proposition}
\begin{proof}
\leavevmode
\begin{enumerate}[label=\textup{(\roman*)},leftmargin=*,itemsep=3pt,topsep=4pt]
\item \textbf{Pullback construction.} Pullback of the universal output and search domains selects exactly
the pairs involving the history of $\kappa$.
Their positive outputs
are the pullbacks of the classifier maps by
\cref{thm:maximal-history-observation,thm:history-search-domain}.
The totality assertion is the preceding corollary.
Leastness is the
frame inequality $\kappa\rho\leq\id_L$, preserved by the functor on
internal locales under geometric base change.
Its pullback in the
relative history fibre compares the reindexed least ideal with each
reindexed history, as in \cref{thm:least-history}.
\end{enumerate}
\end{proof}

The unrestricted membership tests for a Mathias promise are still
full predicates, not automatically semidecisions. The two constructed
domains specify exactly where those tests or their resulting
observations become positive; no separate effective name for a
geometric least section is assumed.

\subsection{A displayed history basis, with all equations}
The canonical semantic order is always available. A combinatorial history presentation is useful because its order and covers may have simpler formulas. Here is the exact comparison that such a presentation must establish, internally in the current parameter topos.

Let $A=\OO(\Hist_\rho(x))$, and let $a:B_x\to A$ be the proposed basic history opens. For a proposed refinement order on $B_x$, verify
\[
 a_b\geq a_c\quad(b\leq c),\qquad
 \bigvee_ba_b=1,\qquad
 a_b\wedge a_c=\bigvee_{d\geq b,c}a_d,\qquad
 u=\bigvee_{a_b\leq u}a_b.
\]
For its proposed coverage $J_x$, verify soundness and detection:
\[
 S\in J_x(b)\quad\Longleftrightarrow\quad
 a_b=\bigvee_{c\in S}a_c.
\]
It is enough to prove the right-to-left implication for the generated topology; the left-to-right implication can be checked on the displayed generators and propagated by the topology axioms.

\begin{proposition}[Exact displayed history comparison]\label{prop:displayed-history}
These equations construct inverse frame maps
\[
 \Omega_{J_x}(B_x)\longrightarrow A,\quad
 U\longmapsto\bigvee_{b:U(b)}a_b,
 \qquad
 A\longrightarrow\Omega_{J_x}(B_x),\quad
 u\longmapsto(b\mapsto a_b\leq u).
\]
Conversely, a presentation isomorphism carrying its basic stage $[b]$ to $a_b$ implies all the displayed equations and the cover-detection equivalence.
\end{proposition}
\begin{proof}
\leavevmode
\begin{enumerate}[label=\textup{(\roman*)},leftmargin=*,itemsep=3pt,topsep=4pt]
\item \textbf{Inverse comparison.} The forward construction is \cref{thm:basis,prop:displayed-basis}, applied in the relative base.
For necessity, stage generators satisfy top, refinement, and common-refinement meet equations, and they are join-dense. A frame isomorphism transports each equation. A sieve covers $b$ precisely when $[b]$ is the join of its member stages, by the saturation formula and \cref{cor:stage-truth}; transport this equality through the isomorphism to obtain cover detection.
Thus the equations are both sufficient and necessary for the specified generator-preserving comparison.
\end{enumerate}
\end{proof}

\begin{corollary}[Direct comparison of history sites]\label{cor:history-direct-comparison}
Work internally over the oracle parameter $T$, and put
$A=\OO(\Hist_\rho(x))$. Let $c_p\in A$ be the canonical fibre-stage
opens of \cref{thm:history-site}, and let $b_h\in A$ be a displayed
history basis satisfying \cref{prop:displayed-history}, with semantic
topologies $J_{\mathrm{can}}$ and $J_x$. The canonical isomorphism from
the displayed to the canonical history forcing frame is
\begin{equation}\label{eq:history-direct-comparison}
 \begin{aligned}
 \Theta_{H,x}:\Omega_{J_x}(B_x)&\longrightarrow
                       \Omega_{J_{\mathrm{can}}}(P),\\
 \Theta_{H,x}(U)(p)&\iff
                 c_p\leq\bigvee_{h:U(h)}b_h.
 \end{aligned}
\end{equation}
Its inverse on a canonical stable predicate $V$ is
\[
 \Theta_{H,x}^{-1}(V)(h)
    \iff b_h\leq\bigvee_{p:V(p)}c_p.
\]
In particular the inverse image of the canonical stage $[p]_x$ is the
displayed stable predicate
\[
 h\longmapsto(b_h\leq c_p).
\]
These maps preserve the relative history observations and commute
with substitutions of the generalized oracle parameter.
\end{corollary}
\begin{proof}
\leavevmode
\begin{enumerate}[label=\textup{(\roman*)},leftmargin=*,itemsep=3pt,topsep=4pt]
\item \textbf{Inverse comparison.} The two applications of basis reflection give explicit frame
isomorphisms to $A$:
\[
 \Phi_H(U)=\bigvee_{h:U(h)}b_h,\qquad
 \Phi_{\mathrm{can}}(V)=\bigvee_{p:V(p)}c_p.
\]
Their inverse maps test the frame inequalities $b_h\leq-$ and
$c_p\leq-$, respectively.
Thus
$\Theta_{H,x}=\Phi_{\mathrm{can}}^{-1}\Phi_H$ and
$\Theta_{H,x}^{-1}=\Phi_H^{-1}\Phi_{\mathrm{can}}$, which are exactly
the formulas displayed.
Since $\Phi_{\mathrm{can}}([p]_x)=c_p$, the
stage formula follows.
\item \textbf{Naturality and reindexing.} Each map preserves its distinguished opens, and
the basis-reflection formulas are internal in the parameter base; this
proves the observation and reindexing assertions.
\end{enumerate}
\end{proof}

The final assertion prevents an arbitrary locale isomorphism from being read as a comparison of two chosen history bases. The generator equation is part of the assertion being proved. Once it holds, all history observations and their relative parameter maps are identified by the explicit inverse formulas.

\subsection{Displayed promises}
Frequently a condition is represented by a committed component $s$ and a promise component $r$, with a map $a_x:B_x\to P$ assigning candidate histories over $x$. The geometric presentation method is to put the original refinement relation on these candidates and pull the structural requirements back to them. A candidate may then acquire an empty cover when it conflicts with a required coordinate. This constructs admissibility through coverage. The Mathias calculation below makes this procedure explicit.

\section{Partition trees and constructed promises}\label{sec:partition-promises}


Here the purpose is to construct the bookkeeping shared by partition
forcing, not to replace the distinct combinatorial arguments for each
principle. Use the coded finite-partition trees and refinement maps of
\cite[Definitions 3.7--3.9]{patey-jumps}. A morphism $S\to T$ includes a
finite map $f:\operatorname{parts}(S)\to\operatorname{parts}(T)$ and
the nodewise assertions that every new part is included in its indicated
old part. Composites use $f\circ g$; set inclusion proves their nodewise
assertions. Identities use identity part maps. Code-padding restrictions
in that presentation are preserved by composition.

Point this category by objects $(\nu,T)$ and arrows whose part map sends
the distinguished new part to the distinguished old part. Let $U[T]$
be the class of pointed refinements below $T$. A promise $C$ is the
upward-closed class of such pointed refinements described in the cited
definition, meeting each infinite refinement in a nonempty part.
This property must be proved for a constructed $C$; naming $C$ is not
a proof of it.

\begin{proposition}[Initial promise and exact restriction]\label{partition:prop:promises}
The maximal class $U[T]$ is a promise. For every actual promise $C$
and refinement $S\to T$, restriction is the constructed class
\[
 C[S]=C\cap U[S].
\]
It is a promise on $S$, and for $R\to S\to T$ one has
$C[S][R]=C[R]$ on the nose. Pointed refinement and this restriction
therefore define a split reindexing construction.
\end{proposition}
\begin{proof}
\leavevmode
\begin{enumerate}[label=\textup{(\roman*)},leftmargin=*,itemsep=3pt,topsep=4pt]
\item \textbf{Initial promise and exact restriction.} An infinite finite-partition tree has a node past the start of its tail.
The colour of any one position in that node exhibits a nonempty part.
It belongs to $U[T]$, proving the maximal promise property without
selecting an infinite path or an infinitely recurring part.
If $R\to S$, then $R\to T$ by composition.
The promise property of
$C$ supplies a nonempty part at $R$; that pointed object also lies in
$U[S]$, hence in $C[S]$.
Upward closure is inherited on the restricted
cone.
\item Finally $U[R]\subseteq U[S]$, so
$C\cap U[S]\cap U[R]=C\cap U[R]$.
The identity and associativity
equations for the part maps and these equalities are the claimed
coherence equations.
\end{enumerate}
\end{proof}

For the primitive recursive tree presentation, this initial promise
also has the required oracle code. With a fixed finite part map,
refinement is a universal numerical test on finite nodes, with a
primitive recursive matrix. Whether a tree has a node at every level
has the same form, because each level is finite. The halting oracle
decides the counterexample searches. If the part map is not already
displayed, test the finitely many maps between the two finite part
sets. These bounded operations yield a primitive recursive procedure
relative to the halting oracle for the maximal promise and for its
restrictions. Thus its coding requirement is constructed as well.

\paragraph{The full partition-forcing carrier.}
For the specific stable computable tournament of
\cite[Section 3]{patey-jumps}, a condition also retains the finite-stem tuple
$\vec F$, the infinite primitive recursive partition tree $T$, and the
promise $C$ primitive recursive relative to $0'$. Each $(F_\nu,\dom(T))$ is the
source's Erd\H{o}s--Moser condition. A refinement retains its finite part map,
the corresponding stem extensions, refinement of the restricted tree, and
promise inclusion, exactly as in Definition 3.13 of that source.
The split restriction proved above acts on this full carrier; it does not
forget the stem or refinement fields. This is the coded order used in
\cref{forcing:thm:raw-saturation}, where the generic numerical name is also
identified explicitly.

The displayed construction proves reindexing and its coherence. It does not
assert that an arbitrarily pruned class is a promise. For the
requirement-dependent pruning in Patey's actual forcing, its proofs are
already available: Lemmas 3.11--3.12 establish the requisite part
normalizations, Lemmas 3.17 and 3.21 give requirement extensions, and
Lemmas 3.22--3.23 prove the truth lemmas. These results supply the requirement extensions for that presentation; \cref{forcing:thm:raw-saturation} identifies its higher-level forcing predicates after saturation. In particular, a nonempty part
must not be confused with a part that is infinite on some path. The
latter contains an infinite-path assertion and is not supplied by the
finite argument above.

\section{Committed information and residual history}\label{sec:promises}


Suppose a realization factors as $L\xrightarrow{\tau}Z\xrightarrow{\sigma}X$. The first map records committed information and the second its extensional oracle interpretation. Pulling this factorization back along $x:T\to X$ gives
\[
 \Hist_\rho(x)\longrightarrow T\times_XZ\longrightarrow T.
\]
The image of the first map records the committed values realized by histories over $x$. The fibres of that map retain the residual promises over a committed value.

\begin{proposition}[Kernel pairs and residual fibres]\label{prop:history-kernel}
The kernel pair $R=L\times_XL$ classifies pairs of histories with the same realized oracle. For the factorization above, equality of committed information implies equality of oracle values, giving a canonical map
\[
 L\times_ZL\longrightarrow L\times_XL.
\]
The image factorization of a realization and the fibres of its surjective leg retain its residual history.
\end{proposition}
\begin{proof}
\leavevmode
\begin{enumerate}[label=\textup{(\roman*)},leftmargin=*,itemsep=3pt,topsep=4pt]
\item \textbf{Pullback construction.} A map $S\to R$ is a pair $h_0,h_1:S\to L$ with $\rho h_0=\rho h_1$, by the pullback property.
If $\tau h_0=\tau h_1$, applying $\sigma$ gives this equation, which defines the displayed map.
For the image factorization $L\onto\im(\rho)\incl X$, the mono makes equality in the image equivalent to equality after inclusion in $X$, so its surjective leg has the same kernel pair as $\rho$.
Pulling it back at an image value gives exactly the histories realizing that value.
\end{enumerate}
\end{proof}

An oracle observation $U$ pulls back to equal opens along the two kernel-pair projections. The converse, that every coherent open is an inverse image from the oracle, is a descent question. Part~\ref{part:descent} constructs the exact descended fragment and separates it from bare coherence.

Effective history observations are obtained by the same marking pullbacks as before. For a displayed relative history presentation, choose its positive trace object from the actual presentation. The pullback of an oracle or condition observation is effective precisely on the maximal admissibility object for that inverse-image map. This applies in parameter contexts and does not require every history order to be positive.

\begin{proposition}[Localizing histories]\label{prop:history-localization}
For $i:L'\incl L$, the history locale of $\rho i$ is
\[
 \Hist_{\rho i}(x)\cong\Hist_\rho(x)\times_LL'.
\]
Its canonical history generators are the restrictions of those of $\Hist_\rho(x)$, and localization towers give coherent such comparisons.
\end{proposition}
\begin{proof}
\leavevmode
\begin{enumerate}[label=\textup{(\roman*)},leftmargin=*,itemsep=3pt,topsep=4pt]
\item \textbf{Localizing histories.} Both expressions are the pullback of $\rho i$ along $x$.
The generator statement follows by applying the pullback to the canonical stage maps, and the relative basis comparison follows from \cref{thm:history-site}.
For a tower, all isomorphisms preserve the projections into the final source locale and the parameter base, so their composites agree by uniqueness.
\end{enumerate}
\end{proof}

\subsection{Residual histories at a displayed region}
Let $\lambda_p:V_p\to K_p$ be a stage comparison. For a generalized region value $v:T\to K_p$, define its residual stage history by
\[
 R_p(v)=T\times_{K_p}V_p.
\]
A map $S\to R_p(v)$ over $T$ is a stage history $h:S\to V_p$ with $\lambda_ph=v\,s$. This is a direct pullback construction, and it retains histories that agree on the region value but differ in their promises.

\begin{proposition}[Residual stage calculus]\label{prop:residual-stage}
The kernel pair of $\lambda_p$ classifies pairs of stage histories with the same displayed region value. Its image factorization has the same kernel pair. A commuting refinement square
\[
\begin{tikzcd}
 V_q\ar[r]\ar[d,"\lambda_q"']&V_p\ar[d,"\lambda_p"]\\
 K_q\ar[r,"r_{p,q}"']&K_p
\end{tikzcd}
\]
induces $R_q(v)\to R_p(r_{p,q}v)$. A section of $\lambda_p$ gives a distinguished point of every such relative residual history.
\end{proposition}
\begin{proof}
\leavevmode
\begin{enumerate}[label=\textup{(\roman*)},leftmargin=*,itemsep=3pt,topsep=4pt]
\item \textbf{Pullback construction.} The first assertion is the universal property of $V_p\times_{K_p}V_p$.
Passing through the image mono does not change equality of the two values, so it does not change that pullback.
For refinement, compose a lift to $V_q$ with $V_q\to V_p$.
The commuting square verifies its required value in $K_p$, hence the residual pullback gives the induced map.
The maps preserve identity and composition because their composites into $V_p$ do.
If $\kappa_p$ is a section, $(\id_T,\kappa_pv)$ is a section of $R_p(v)\to T$, by $\lambda_p\kappa_pv=v$.
\end{enumerate}
\end{proof}

This proves the relationship between a region section and a pointed history. It does not identify the entire residual locale with its distinguished point. In particular, the Mathias stage section constructs exact must semantics while leaving all additional compatible reservoir histories present.

\begin{theorem}[Monicity and history reconstruction]\label{thm:history-monic}
For a realization $\rho:L\to X$, the following are equivalent:
\begin{enumerate}
\item $\rho$ is monic;
\item the diagonal $L\to L\times_XL$ is an isomorphism;
\item every generalized oracle has at most one generalized history, in every further parameter locale.
\end{enumerate}
If $\rho$ also has a section, these are equivalent to $\rho$ being an isomorphism. A least-history section, when it exists, is unique.
\end{theorem}
\begin{proof}
\leavevmode
\begin{enumerate}[label=\textup{(\roman*)},leftmargin=*,itemsep=3pt,topsep=4pt]
\item \textbf{Monicity and history reconstruction.} If $\rho$ is monic, the two projections from its kernel pair agree.
That common projection is inverse to the diagonal, by the two pullback projection equations.
Conversely an invertible diagonal makes any pair $h_0,h_1$ with $\rho h_0=\rho h_1$ equal, proving monicity.
The same equality is exactly uniqueness of lifts in the generalized-fibre description, proving equivalence with (3).
\item \textbf{Forward implication.} If $\rho\kappa=\id$ and $\rho$ is monic, then
$\rho\kappa\rho=\rho=\rho\id$ implies $\kappa\rho=\id$.
Thus $\rho$ is invertible.
Conversely an isomorphism is monic and sectioned.
Finally, for two least sections $\kappa,\kappa'$, each lies below the other on every generalized oracle; in particular on the identity of $X$.
The order on locale maps is antisymmetric because it is pointwise order on inverse-image maps.
Hence $\kappa=\kappa'$.
\end{enumerate}
\end{proof}

There is a useful fibrewise description of the kernel pair. Pull $\rho:L\to X$ back along $\rho$ itself. The result is $L\times_XL\to L$, whose fibre over a history $h$ consists of all histories with the same realized oracle as $h$. Consequently the kernel pair is the universal family of observationally equivalent histories, while the fibre over one $x$ is its specialization at that oracle.

\subsection{A full fibre construction for Mathias promises}
Let $g$ be a total increasing sequence in a parameter topos. All condition objects and relations in what follows are pulled to that base. A candidate future is a pair $(m,r)$ for which $(g|_m,r)$ is a valid Mathias condition. Put $a_g(m,r)=(g|_m,r)$ and define its refinement by the original Mathias relation. The next candidate is
\[
 \sigma_g(m,r)=(m+1,r_{>g(m)}),
\]
defined when $g(m)\in\im(r)$. At $(m,r)$ impose the singleton cover by this candidate when it is defined, and the empty cover otherwise.

\begin{lemma}[Candidate cover stability]\label{lem:future-covers}
These candidate covers are stable under refinement in the generated-coverage sense.
\end{lemma}
\begin{proof}
\leavevmode
\begin{enumerate}[label=\textup{(\roman*)},leftmargin=*,itemsep=3pt,topsep=4pt]
\item \textbf{Observation preservation.} Suppose $(m,r)\leq_g(n,u)$.
If $m<n$, the first new stem value $g(m)$ belongs to $\im(r)$, and $(n,u)$ already refines $\sigma_g(m,r)$.
Thus the pulled-back sieve is maximal.
If $m=n$ and $g(n)\in\im(u)$, reservoir inclusion gives $g(m)\in\im(r)$, and taking tails preserves that inclusion; the child of $(n,u)$ refines the child of $(m,r)$.
If $m=n$ and the membership fails, the generator at $(n,u)$ is empty and is contained in every pulled-back sieve.
These cases exhaust the possibilities because stem length is monotone under refinement.
\end{enumerate}
\end{proof}

\begin{theorem}[Reconstruction of the Mathias history fibre]\label{thm:mathias-fibre}
The candidate-future site presents the fibre of $\rho_M$ over $g$. Its continuous ideals correspond to continuous Mathias ideals realizing $g$ by
\[
 A_I(m,r)\iff I(a_g(m,r)),\qquad
 I_A(p)\iff\exists(m,r)\,[A(m,r)\wedge p\leq a_g(m,r)].
\]
The correspondence is geometric in the fixed pulled-back relations and coherent under further change of parameter.
\end{theorem}
\begin{proof}
\leavevmode
\begin{enumerate}[label=\textup{(\roman*)},leftmargin=*,itemsep=3pt,topsep=4pt]
\item \textbf{Reconstruction of the Mathias history fibre.} Let $I$ realize $g$.
Every stem in $I$ is a prefix of $g$, and every subsequent value of $g$ lies in that condition's reservoir, by the directedness and next-element argument in \cref{sec:mathias-calculation}.
Thus each member is of the form $a_g(m,r)$ and $A_I$ is inhabited, lower, and directed.
\item \textbf{Ideal and directedness calculation.} If $A_I(m,r)$, continuity of $I$ meets a next-element cover. A child in the ideal must have next value $g(m)$, because its stem is a prefix of $g$.
Lowerness therefore puts $a_g(\sigma_g(m,r))$ in $I$.
The membership $g(m)\in\im(r)$ is also obtained.
Hence $A_I$ meets the candidate generator.
\item \textbf{Converse construction.} Conversely let $A$ be a continuous candidate ideal.
The formula for $I_A$ makes it inhabited and lower; directedness follows by taking a common upper bound in $A$ for two witnesses.
Continuity for the candidate cover implies
\[
 A(m,r)\Rightarrow
 g(m)\in\im(r)\ \wedge\ A(\sigma_g(m,r)).
\]
Indeed an empty generator cannot be met, and a member in the sieve of a singleton puts its generator in the lower ideal.
\item \textbf{Observation preservation.} Now take $p=(s,v)\in I_A$, witnessed by $(m,r)$.
If $|s|<m$, the first new value $g(|s|)$ supplies a child of $p$ below $a_g(m,r)$.
If $|s|=m$, the displayed closure supplies a child below $a_g(\sigma_g(m,r))$.
In both cases that child belongs to $I_A$.
Thus the original next-element generators are met.
\item \textbf{Ideal and directedness calculation.} To justify passage from generators to all covers, for any ideal $I$ let $J_I(p)$ be the sieves $S$ such that $I(p)\Rightarrow\exists q\in I\cap S$.
Maximality uses $p$; stability uses a common upper bound in $I$; transitivity uses two successive existential witnesses.
Hence $J_I$ is a topology.
Meeting the generators gives continuity for the generated topology.
This applies to both ideals above.
\item \textbf{Forcing construction.} Iterating the deterministic candidate child produces prefixes of $g$ of every larger length.
All stems in $I_A$ are prefixes of $g$, so its realized sequence is exactly $g$.
\item \textbf{Existence and uniqueness.} For the inverse equations, every member $p$ of $I$ is some $a_g(m,r)$, and lowerness gives the converse containment, so $I_{A_I}=I$.
Likewise
$A_{I_A}(m,r)$ means that some member of $A$ refines $(m,r)$, and lowerness makes this equivalent to $A(m,r)$.
These proofs are finite, existential, and use the fixed relations, so hold in every further parameter context.
Applying them to universal models in arbitrary test locales gives an isomorphism of the classifying locales by \cref{prop:history-model}, not just a bijection of external points.
\end{enumerate}
\end{proof}

\begin{proposition}[Frame-valued form of the Mathias fibre equivalence]\label{prop:mathias-fibre-framevalued}
The equivalence in \cref{thm:mathias-fibre} is natural in every relative test locale $Z$. In particular, interpreting the candidate predicate $A$ as a family of opens of $Z$, the formulas
\[
 I_A(p)=\bigvee_{(m,r):\,p\leq a_g(m,r)}A(m,r),
 \qquad A_I(m,r)=I(a_g(m,r))
\]
are inverse transformations of frame-valued models of the two respective geometric theories. The resulting isomorphism of their classifying locales is therefore an isomorphism of locales over the generalized oracle parameter.
\end{proposition}
\begin{proof}
\leavevmode
\begin{enumerate}[label=\textup{(\roman*)},leftmargin=*,itemsep=3pt,topsep=4pt]
\item \textbf{Ideal and directedness calculation.} All displayed joins take place in $\OO(Z)$ and are indexed by predicates in the fixed relative parameter context.
For a candidate model $A$, lowerness and directedness of $I_A$ follow from the corresponding frame equations for $A$, by distributing finite meets over the displayed joins and applying the common-refinement equation.
The candidate singleton and empty covers imply the original next-element cover equations, by the two cases $|s_p|<m$ and $|s_p|=m$ used in the proof of \cref{thm:mathias-fibre}.
\item \textbf{Converse construction.} Conversely, evaluation of the original model $I$ at a candidate is a model of the candidate coverage, because an original next-element cover compatible with $g$ has exactly the next value $g(m)$.
\item \textbf{Unit and identity laws.} The composite $I\mapsto A_I\mapsto I_{A_I}$ is the identity by lowerness and the fact that every nonzero original membership stage has a stem prefix of $g$.
The other composite is the identity because $A$ is lower and every witness in the displayed join defining $A_{I_A}(m,r)$ refines $(m,r)$.
These proofs concern frame inequalities and use only finite meets and existential joins, so commute with arbitrary frame homomorphisms $\OO(Z)\to\OO(Z')$.
Naturality in $Z$ and the universal properties of the two geometric classifiers give the claimed locale isomorphism.
\end{enumerate}
\end{proof}

The construction explains the role of the empty candidate covers. A promise contradicting the next required value is eliminated by the topology. Repeated singleton continuation then enforces all later membership requirements. The infinitary compatibility of a reservoir with the whole sequence has been derived from continuity.

\section{Least Mathias and perfect-tree histories}\label{sec:least-history-examples}

The abstract least-history criterion can be tested on familiar forcing orders. Mathias conditions admit a canonical full-tail history under the stated realization, while perfect-tree histories retain branching information that a path alone does not determine. These comparisons locate exactly which reconstruction maps exist.


\begin{proposition}[An explicit global least Mathias history]
\label{prop:mathias-least-frame}
For a total increasing generic sequence $g$ let
$c_{g,0}=(\varnothing,\operatorname{id}_{\N})$ and, for $m>0$,
let $c_{g,m}$ have stem $g|_m$ and the full natural-number
reservoir above its last member. For every Mathias condition $p$ put
\[
 W_p(g)\iff\exists m:\N\,[p\leq c_{g,m}].
\]
These opens form a continuous frame-valued condition ideal and
hence define the least section of \cref{thm:mathias-least}.
\end{proposition}
\begin{proof}
\leavevmode
\begin{enumerate}[label=\textup{(\roman*)},leftmargin=*,itemsep=3pt,topsep=4pt]
\item \textbf{Ideal and directedness calculation.} The conditions $c_{g,m}$ form an increasing sequence.
Their
lower closure is inhabited, lower and directed, with maximum of
two indices providing directedness.
Each $c_{g,m}$ has a
next-element child below $c_{g,m+1}$, namely the child selecting
$g(m)$, so the generated ideal meets the next-element cover.
\Cref{lem:ideal-topology,prop:site-classifier} then construct a
locale map from these frame-valued opens.
The sequence of its
stems has exactly the prefixes of $g$, hence its realization is
$g$.
\item \textbf{Scott continuity.} Every continuous Mathias ideal realizing $g$ contains a
condition with stem at least $g|_m$; by lowerness it contains
$c_{g,m}$.
The constructed ideal is therefore the least one,
proving the claim without a global selection of arbitrary histories.
\end{enumerate}
\end{proof}

\begin{theorem}[The least full-tail Mathias history]\label{thm:mathias-least}
The Mathias realization has a canonical least section. Its history at $g$ is the lower ideal generated by the conditions with stem $g|_m$ and full reservoir tail above the last stem value. The realization is not monic.
\end{theorem}
\begin{proof}
\leavevmode
\begin{enumerate}[label=\textup{(\roman*)},leftmargin=*,itemsep=3pt,topsep=4pt]
\item \textbf{The least full-tail Mathias history.} The section and its fibrewise leastness are constructed in \cref{prop:mathias-least-frame}; \cref{thm:least-history} identifies their universal property with the adjunction.
\item \textbf{Scott continuity.} For nonmonicity, use the computable sequence $g(n)=2n$.
Besides its full-tail history, take the history generated by prefixes of $g$ with the even reservoir tail.
Both are continuous and realize $g$.
The latter contains the empty-stem even-reservoir condition.
The former does not, because no full tail is contained in the even numbers.
The corresponding points of $L_M$ are distinct and have the same image under $\rho_M$.
\end{enumerate}
\end{proof}

\subsection{Mathias histories, Ellentuck covers, and dense forcing}
In the classical interpretation, an increasing reservoir enumeration and its
infinite range determine one another. Thus the preceding conditions
$(s,r)$ give the Mathias order $(s,A)$: strengthen the stem using $A$,
and shrink the remaining infinite reservoir. The corresponding regions
\[
 [s,A]=\{x\in[\N]^\omega:s\text{ is an initial segment of }x,
                                  \ x\setminus s\subseteq A\}
\]
are the Ellentuck basis; see \cite{dobrinen-ramsey}. Let $J_{\rm next}$ be the
next-element coverage of \cref{sec:mathias-calculation}, and $J_{\rm Ell}$ the coverage by unions
of these basis regions.

\begin{proposition}\label{cal:prop:mathias}
Classically,
\[
 J_{\rm next}\ \subsetneq\ J_{\rm Ell}\ \subseteq\ \Jdens.
\]
The two structural sites have the same dense completion, namely
$\Sh(\RegOpen(\mathbb M))$, but they are not identified as sites over their
common condition order.
\end{proposition}
\begin{proof}
\leavevmode
\begin{enumerate}[label=\textup{(\roman*)},leftmargin=*,itemsep=3pt,topsep=4pt]
\item \textbf{Forcing construction.} Every increasing infinite sequence has a first new element after its stem;
the next-element regions therefore cover the original region.
If a family of basis regions covers $[s,A]$, take a stronger $(t,B)$ and a
point $x\in[t,B]$, for example its stem followed by its enumerated reservoir.
Some member $[u,C]$ contains $x$.
The stems $t,u$ are comparable; take the
longer one and the infinite common tail of $x$ in $B\cap C$.
This is a common
refinement.
Hence the family is extension-dense.
\item \textbf{Scott continuity.} For strictness, start at $(\varnothing,\N)$ and put $g(n)=2n$.
The downward
closure $I$ of the chain
\[
 c_m=(g\mathbin{\upharpoonright} m,\{n:n>g(m-1)\})
\]
(with full reservoir at $m=0$) is a $J_{\rm next}$-continuous filter.
Indeed, if $p\in I$, some $c_m$ strengthens $p$; lengthen this chain until it
passes the next stem coordinate of $p$.
The resulting condition strengthens
the next-element child selected by $g$.
\item \textbf{Localization and coverage.} Let $S$ be the sieve generated by $(\varnothing,2\N)$ and all full-tail
conditions whose stem contains an odd number.
Their regions cover all
infinite increasing sequences: a sequence is entirely even or has an odd
entry.
However, no $c_m$ refines any generator.
Its reservoir contains odd
numbers, and its stem contains none.
Thus $I$ misses $S$, so $S$ is not a
$J_{\rm next}$-cover.
The final assertion follows from
\cref{cal:thm:localize} applied to the two intermediate nuclei.
\end{enumerate}
\end{proof}

Thus \cref{thm:mathias-least,thm:mathias-baire-region} describe the
next-element history locale, while the constructed additional coverages give
the Ellentuck and classical forcing semantics.

\subsection{Perfect trees and the same least-history pattern}
The calculation is internal to $\Eff$. A tree is a prefix-closed
subobject of $2^{<\N}$; pruning and perfectness mean that every
member has arbitrarily long extensions and two incompatible
extensions in the tree. No decidability of tree membership is used.
Let $P_S$ be the object of nonempty pruned perfect binary trees,
ordered by reverse inclusion. For a finite string $\sigma$, put
\[
 C_\sigma=\{\tau:\tau\preceq\sigma\text{ or }\sigma\preceq\tau\}.
\]
For $u\in T$, the restriction $T|u=T\cap C_u$ is again a condition:
below $u$ one first extends to $u$ and then uses its extensions in
$T$; above $u$ one uses pruning and perfectness in $T$ directly.
Cover $T$ by the restrictions $T|u$ for $u\in T\cap2^n$.
Under a refinement $S\subseteq T$, its level restrictions $S|u$
refine the corresponding restrictions $T|u$, so these families
generate the stated coverage.

Write $b_T\in\OO(L_S)$ for its stage generators and $[\sigma]$
for the Cantor cylinders. Their defining relations include
\[
 b_{2^{<\N}}=1,\qquad
 b_T\wedge b_S=\bigvee_{R\in P_S,\ R\subseteq T\cap S}b_R,
 \qquad b_T=\bigvee_{u\in T\cap2^n}b_{T|u}.
\]
In particular $b_{C_\epsilon}=1$ and
\[
 b_{C_\sigma}=b_{C_{\sigma0}}\vee b_{C_{\sigma1}},\qquad
 b_{C_{\sigma0}}\wedge b_{C_{\sigma1}}=0.
\]
There is no nonempty pruned tree inside the intersection of the two
successor trees. These are the Cantor presentation relations, so
they construct the realization
\[
 \rho_S:L_S\to\mathbb C,\qquad
 \rho_S^*[\sigma]=b_{C_\sigma}.
\]
The level cover of $T$ gives
$b_T\leq\bigvee_{u\in T\cap2^n}\rho_S^*[u]$ for every $n$;
thus this map realizes the branch determined by the generic history.

For precision, the path sublocale of an arbitrary internal tree is
\[
 \operatorname{Path}(T)=\bigwedge_{n:\N}\mathsf o(O_{T,n}),
 \qquad
 O_{T,n}=\bigvee_{\sigma\in2^n}!^*T(\sigma)\wedge[\sigma],
\]
where $\mathsf o(U)$ is the open sublocale corresponding to $U$.
It classifies branches all of whose prefixes belong to $T$; it need
not be a closed sublocale.

\begin{proposition}[Least cylinder history and nontrivial fibre]
\label{prop:perfect-history}
There is a least section $\kappa_S:\mathbb C\to L_S$, with
\[
 \kappa_S^*(b_T)=\operatorname{int}\operatorname{Path}(T)
 =\bigvee_{C_\sigma\subseteq T}[\sigma].
\]
It satisfies $\rho_S\kappa_S=\id$,
$\kappa_S\rho_S\leq\id$, and $\kappa_S^*=(\rho_S)_*$.
The realization $\rho_S$ is not monic.
\end{proposition}
\begin{proof}
\leavevmode
\begin{enumerate}[label=\textup{(\roman*)},leftmargin=*,itemsep=3pt,topsep=4pt]
\item \textbf{Ideal and directedness calculation.} Put $W_T=\bigvee_{C_\sigma\subseteq T}[\sigma]$.
We verify the
frame-valued ideal relations directly.
The full tree gives top,
and $R\subseteq T$ gives $W_R\leq W_T$.
Expand
$W_T\wedge W_S$ as a join of intersections of cylinders.
Incompatible cylinders have intersection zero; compatible ones
intersect in the longer cylinder $[\nu]$, with
$C_\nu\subseteq T\cap S$.
Taking $R=C_\nu$ proves
\[
 W_T\wedge W_S=\bigvee_{R\in P_S,\ R\subseteq T\cap S}W_R.
\]
For the reverse inequality each such $W_R$ is below both factors.
\item \textbf{Observation preservation.} For a level cover, take $C_\sigma\subseteq T$ and
$m=\max(n,|\sigma|)$.
The finite cylinder decomposition
\[
 [\sigma]=\bigvee_{\nu\in2^m,\ \sigma\preceq\nu}[\nu]
\]
has each $C_\nu\subseteq T|(\nu|_n)$.
Hence
$W_T\leq\bigvee_{u\in T\cap2^n}W_{T|u}$; the reverse
inequality is monotonicity.
The classifier universal property of
\cref{prop:site-classifier} therefore constructs the unique map
$\kappa_S$ with inverse images $W_T$.
\item \textbf{Universal-property argument.} To identify these opens intrinsically, we prove
\[
 [\sigma]\subseteq\operatorname{Path}(T)
 \quad\Longleftrightarrow\quad C_\sigma\subseteq T.
\]
For the forward implication, take any finite $\tau$ comparable
with $\sigma$ and extend the longer of them by zeros.
This explicit
branch lies in $[\sigma]$; its path condition gives $\tau\in T$.
This construction works in the parameter context of $T,\sigma,\tau$,
so it does not appeal to enough points.
\item \textbf{Converse construction.} Conversely, every prefix of
a branch in $[\sigma]$ belongs to $C_\sigma$, which proves the
factorization through each open sublocale in the displayed meet.
The cylinder basis now gives
$W_T=\operatorname{int}\operatorname{Path}(T)$.
\item \textbf{Unit and identity laws.} Since $W_{C_\sigma}=[\sigma]$, the composite $\rho_S\kappa_S$
is the identity on cylinder generators.
Moreover
\[
 (\kappa_S\rho_S)^*(b_T)
 =\bigvee_{C_\sigma\subseteq T}b_{C_\sigma}\leq b_T.
\]
The stages generate the frame, so $\kappa_S\rho_S\leq\id$.
For every generalized history $h$ realizing $g$ this gives
$\kappa_Sg=\kappa_S\rho_Sh\leq h$, proving leastness.
The same equations give $\rho_S^*\dashv\kappa_S^*$ and thus
$\kappa_S^*=(\rho_S)_*$, as in \cref{thm:least-history-interior}.
\item \textbf{Ideal and directedness calculation.} For nonmonicity, let $g=000\cdots$ and let $Q$ consist of strings
whose even-indexed coordinates are zero.
It is a decidable perfect
tree.
The two lower ideals
\[
 I_{\min}=\{T:\exists n\ C_{0^n}\subseteq T\},\qquad
 I_Q=\{T:\exists n\ Q\cap C_{0^n}\subseteq T\}
\]
are inhabited and directed, using the maximum of two indices.
They meet every level cover: increase the index to the requested
level and use the all-zero restriction.
Hence they give points of
$L_S$ by \cref{thm:points-ideals}, and both realize $g$.
\item \textbf{Forcing construction.} But $Q\in I_Q$, whereas $Q\notin I_{\min}$, since every full
cylinder has an extension with a $1$ in a later even coordinate.
Thus the two histories are distinct and $\rho_S$ is not monic.
This internal structural calculation is distinct from the
regular-open forcing completion in perfect-set forcing \cite{sacks}.
\end{enumerate}
\end{proof}

\begin{remark}[Why closedness is not used]
\label{rem:perfect-path-closedness}
For an arbitrary truth value $p:\Omega$, the tree
\[
 T_p=\{\epsilon\}\cup\{0\sigma:\sigma\in2^{<\N}\}
       \cup\{1\sigma:\sigma\in2^{<\N},\ p\}
\]
is nonempty, pruned and perfect, but its path sublocale is the open
$[0]\vee(!^*p\wedge[1])$. Pulling it back along $111\cdots$
gives the open sublocale $p$ of the terminal locale. If this were
closed, its nucleus $x\mapsto(p\Rightarrow x)$ would equal
$x\mapsto q\vee x$ for some $q$. At $x=0$ this gives $q=\neg p$;
at $x=p$ it gives $p\vee\neg p$. Thus a uniform assertion that all
these path sublocales are closed would imply excluded middle.
The preceding cylinder proof establishes the full least-history
statement and interior formula without such an assertion.
\end{remark}

The two examples realize the same categorical pattern with different
promises: reservoir restrictions in one case, tree restrictions in
the other. The least history gives a canonical reconstruction, and
the full fibre records the remaining construction information.


\subsection{The full Sacks carrier and its computable fragment}
The full classical interpretation of the preceding nonempty perfect
binary-tree order is Sacks forcing, as defined in
\cite[Definition 21]{geschke-quickert}. The maps
\[
 T\longmapsto[T],\qquad
 K\longmapsto\{\sigma:K\cap[\sigma]\ne\varnothing\}
\]
identify pruned perfect trees and nonempty perfect closed sets. A node of a
pruned binary tree lies on a branch by successive selection of its least
available child, so these maps are inverse. Consequently $[S]\subseteq[T]$
iff $S\subseteq T$. Finite-level restriction covers are dense by restriction
inside any stronger perfect subtree. \Cref{cal:thm:localize} therefore
constructs the exact Sacks regular-open sheaf topos from the preceding
full perfect-tree presentation.

\begin{proposition}[The computable fragment is not dense]\label{cal:prop:sacks}
Computable pruned perfect binary trees are not dense in the full classical
Sacks order.
\end{proposition}
\begin{proof}
\leavevmode
\begin{enumerate}[label=\textup{(\roman*)},leftmargin=*,itemsep=3pt,topsep=4pt]
\item \textbf{Observation preservation.} Let $T$ consist of the binary strings $\sigma$ such that for every
$e<|\sigma|$, if $\varphi_e(e)$ converges to $0$ or $1$, then
$\sigma(e)\ne\varphi_e(e)$.
Membership is decidable in $0'$.
Each coordinate
has at least one allowed bit. A computable infinite set of indices for
everywhere-divergent programs supplies arbitrarily late unrestricted
coordinates, by padding.
Thus $T$ is nonempty, pruned, and perfect.
Every branch is DNC, hence noncomputable. A computable pruned subtree would
have a computable branch by the least-child recursion.
It therefore cannot
be a nonempty perfect subtree of $T$.
\end{enumerate}
\end{proof}

\Cref{fusion:thm:diagonal} later realizes the same obstruction by protected-front fusion. The full tree object and its computable presentation are therefore retained
as distinct structured data. Their inclusion does not provide a
dense-site equivalence. Classical continuous reading and minimality results
are available in \cite[Section 2]{geschke-quickert}; they require their fusion
construction in addition to the present site identification.

\section{Intensional enhancement: history and promise names}
\label{sec:intensional-histories}


Let $x:T\to X$ be a generalized oracle and
$H_x=\Hist_\rho(x)$ its relative locale. Its canonical history
site has stage object $P_x$ and generic membership opens $a_p$,
as in \cref{thm:history-site}. A history name is a stage $p$ together
with a witness in the universal continuous ideal; this retains the
choice of condition and its promise parameters, while its semantic
interpretation is the corresponding relative history model.

\begin{proposition}[Universal common refinement of named history stages]
\label{prop:intensional-history-names}
For a displayed history presentation $(B_x,J_x)$ satisfying
\cref{prop:displayed-history}, let $b_h\in\OO(H_x)$ be its basic
opens and let $c_p\in\OO(H_x)$ be the canonical history opens.
The incidence family
\[
 e:B_x\times P_x\to\OO(H_x),\qquad e_{h,p}=b_h\wedge c_p
\]
is join-dense. Give its carrier the semantic refinement order
$(h,p)\preceq(h',p')$ when $e_{h',p'}\leq e_{h,p}$ and the semantic
cover topology. This \emph{common refinement site} classifies
exactly the history locale $H_x$. Its two projections compare the
original basic descriptions with their canonical stages, and the
resulting frame comparisons are natural in generalized oracle
parameters. Equal realized oracles do not force equality of these
retained history descriptions.
\end{proposition}
\begin{proof}
\leavevmode
\begin{enumerate}[label=\textup{(\roman*)},leftmargin=*,itemsep=3pt,topsep=4pt]
\item \textbf{Forcing construction.} Both the displayed and canonical basic families are join-dense.
For every $u\in\OO(H_x)$, distributivity and the two generator
expansions yield
\[
 u=\bigvee_{b_h\leq u}b_h
   =\bigvee_{b_h\leq u}\ \bigvee_p(b_h\wedge c_p),
\]
where the last equality uses $\bigvee_pc_p=1$.
Thus the meet
family is join-dense, and its semantic order and cover equations
satisfy \cref{cor:canonical-basis}, which gives the asserted
classifier.
\item \textbf{Existence and uniqueness.} In the relative frame the opens $b_h$ and $c_p$ are
respectively joins of incidence opens lying below them; hence
the canonical maps out of the common classifier preserve each
specified generator.
They coincide with the explicit maps of
\cref{cor:history-direct-comparison} by uniqueness of basis
reflection.
Every construction is given by finite meets and
relative joins of the same basic opens, so it commutes with
parameter substitution as in \cref{prop:history-reindex}.
\item \textbf{Kernel-pair descent.} The kernel pair $L\times_XL$ can contain distinct histories over
one oracle, by \cref{thm:mathias-least}.
Since the comparison
retains their separate membership opens, it does not equate their
descriptions merely from equality of realized values.
\end{enumerate}
\end{proof}

In Mathias forcing a name contains a finite stem and a reservoir,
even when two names reconstruct the same increasing sequence. The
candidate-future site of \cref{thm:mathias-fibre} gives the explicit
relative name presentation. Its empty and singleton covers
respectively eliminate incompatible promises and force the next
committed value. The least full-tail section selects one canonical
history, but it is not an identification of the entire fibre with
that history. The positive observation of any selected history name
is governed by the output and search domains of
\cref{thm:maximal-history-observation,thm:history-search-domain}.

\part{Local candidates and descent of observations}\label{part:descent}

A computation may become available only after a local choice of oracle or answer. The issue is then whether its observations, and separately its descriptions, agree on overlaps sufficiently to descend. We construct the universal observational fragment of a candidate span and its effective marking, keeping semantic descent distinct from the existence of a global name.


\section{Local realizations and candidate relations}\label{sec:candidates}


Here local choices are represented by locale maps. When
a naming fibre is instead an epic sheaf with a simply
transitive group action, \cref{sec:naming-obstructions}
supplies its torsor obstruction. The candidate construction
below applies directly to the retained localic realization
span and requires no conversion of that span into a torsor.

An oracle can be available after a local choice without being globally represented by a function on the original forcing space. The geometric datum is
\[
 \begin{tikzcd}
 &\widetilde L\ar[dl,two heads,"e"']\ar[dr,"t"]&\\
 L&&X .
 \end{tikzcd}
\]
Here $e$ is a localic surjection, meaning that $e^*$ is injective. The right leg is the realized oracle after the choice. The span is retained even when some observations descend through $e$.

\begin{definition}[Candidate image]\label{def:candidate}
Factor $(e,t):\widetilde L\to L\times X$ through its localic image:
\[
 \widetilde L\xrightarrow{q}\Cand(e,t)\xrightarrow{m}L\times X.
\]
Write $c:\Cand(e,t)\to L$ and $u:\Cand(e,t)\to X$ for the projections.
\end{definition}

 the image factorization can be constructed on frames. If $f:Y\to Z$, the nucleus $f_*f^*$ on $\OO(Z)$ gives its image sublocale. The map from $Y$ into that sublocale has injective inverse image, because on fixed elements equality after $f^*$ is reflected by $f_*$. Its composite with the sublocale inclusion is $f$. This constructs the image used here.

\begin{proposition}[Candidate properties]\label{prop:candidate}
The projection $c$ is surjective. The kernel pair of $q$ is the sublocale of $\widetilde L\times_L\widetilde L$ on which the two $t$-values agree. The span descends to $\rho:L\to X$ precisely when its candidate relation is the graph of $\rho$, equivalently when $c$ is an isomorphism.
\end{proposition}
\begin{proof}
\leavevmode
\begin{enumerate}[label=\textup{(\roman*)},leftmargin=*,itemsep=3pt,topsep=4pt]
\item \textbf{Candidate properties.} Since $e=cq$, injectivity of $e^*=q^*c^*$ implies injectivity of $c^*$.
Equality after $q$ is equivalent to equality after $(e,t)$ because $m$ is monic.
Thus a pair in its kernel pair has equal $e$-values and equal $t$-values, proving the second statement by the pullback property.
\item \textbf{Verification of the equations.} If $t=\rho e$, then $(e,t)$ factors as $e$ followed by the graph of $\rho$.
The image of a surjection followed by a sublocale inclusion is that sublocale, so the candidate image is the graph.
\item \textbf{Converse construction.} Conversely a graph projection is an isomorphism and its second component gives $\rho$.
If merely $c$ is known to be an isomorphism, set $\rho=uc^{-1}$; then $t=uq=uc^{-1}e$.
All constructions are unique because $e$ is epic.
\end{enumerate}
\end{proof}

The candidate image records the possible pairs of condition history and oracle value. The surjection $q$ can still have nontrivial fibres. Those fibres are histories of the choice procedure and remain part of the original local realization span.

For a generalized base history $x:T\to L$, the relative candidate locale is $T\times_L\Cand(e,t)\to T$. Pullback of the original span gives the corresponding local histories. These constructions are coherent under parameter substitution by the same pullback argument as \cref{prop:history-reindex}. A point of the relative candidate locale is a candidate; localic surjectivity of its projection is not a supplied section.

\subsection{Choosing a candidate and descending the local oracle}
A section $s:L\to C$ of the candidate projection $c:C\to L$ chooses a candidate in each base context. Its selected realization is $\rho_s=us$. This selection need not recover the original map $t$ on every local choice.

\begin{proposition}[When a selected candidate descends]\label{prop:section-descent}
For a section $s$ of $c$, the following are equivalent:
\[
 \rho_se=t,\qquad se=q,\qquad c\text{ is an isomorphism and }s=c^{-1}.
\]
After pulling back along $e$, there is always the canonical candidate section
\[
 \widetilde L\longrightarrow\widetilde L\times_LC,\qquad
 y\longmapsto(y,qy),
\]
and its candidate value is $t$.
\end{proposition}
\begin{proof}
\leavevmode
\begin{enumerate}[label=\textup{(\roman*)},leftmargin=*,itemsep=3pt,topsep=4pt]
\item \textbf{When a selected candidate descends.} If $\rho_se=t$, the maps $se,q:\widetilde L\to C$ have equal composites with both $c$ and $u$: the first composites are $e$, and the second are $t$.
Since $(c,u)$ is the image inclusion in $L\times X$, it is monic, hence $se=q$.
Conversely this equality gives $use=uq=t$.
\item \textbf{Forward implication.} If $se=q$, then $scq=se=q$.
Surjectivity of $q$ implies $sc=\id_C$, while $cs=\id_L$ is the section equation.
Thus $c$ is an isomorphism with inverse $s$.
The converse follows from $cq=e$.
\item Finally $(\id,q)$ lands in the indicated pullback by $cq=e$, projects identically to $\widetilde L$, and has value $uq=t$, proving the pointed assertion.
\end{enumerate}
\end{proof}

A span is called stable when its left leg remains a kernel-pair coequalizer after every locale pullback. Such a map is in particular stably surjective. The open surjections constructed below form one class. This terminology records a property of the map already present in the span.

\begin{proposition}[Stable local candidates]\label{prop:stable-candidates}
If $e$ is stably surjective, then $c$ is stably surjective. Thus the relative candidate locale is internally locally inhabited after every change of localic base. This does not supply a section of $c$.
\end{proposition}
\begin{proof}
\leavevmode
\begin{enumerate}[label=\textup{(\roman*)},leftmargin=*,itemsep=3pt,topsep=4pt]
\item \textbf{Stable local candidates.} Pull back $e=cq$ along any $h:T\to L$.
The resulting map $e_T$ is surjective by the stated property and factors through $c_T:T\times_LC\to T$.
If $c_T^*$ identifies two opens, applying the remaining inverse image identifies them under $e_T^*$, which is injective.
Hence $c_T$ is surjective. A relative locale's structure map being surjective is precisely its localic inhabitation assertion. A section would additionally be a map splitting that projection, which has not been constructed by this argument.
\end{enumerate}
\end{proof}

\subsection{Discrete answers from local candidates}
Let $a:X\to\disc(A)$ be a total discrete-valued observable. Form the candidate image of $(e,at)$:
\[
 R_a\incl L\times\disc(A).
\]
Its projection to $L$ is surjective, and it records the locally possible answers to this particular observation.

\begin{proposition}[Descent of an answer]\label{prop:answer-descent}
The local answer $at$ descends to a unique map $L\to\disc(A)$ precisely when $R_a\to L$ is an isomorphism. If this relative candidate locale is represented by the discrete locale of an object $B$ in $\Sh_\E(L)$, this criterion is $B\cong1$. Under that representation its value map is an actual object map $B\to\pi^*A$.
\end{proposition}
\begin{proof}
\leavevmode
\begin{enumerate}[label=\textup{(\roman*)},leftmargin=*,itemsep=3pt,topsep=4pt]
\item \textbf{Product calculation.} Apply \cref{prop:candidate} to the new span with right leg $at$.
Under the relative locale correspondence a graph projection is the terminal relative locale.
The discrete-locale construction is fully faithful on objects: a locale map $\disc(B)\to\disc(A)$ is a partition of $B$ by the singleton opens of $A$, hence a total functional relation and thus an object map.
It therefore reflects isomorphisms and sends the terminal object to the terminal locale.
This proves the criterion and the statement about the value map.
\end{enumerate}
\end{proof}

Different observations can have different answer-candidate locales. Some may be discrete and single-valued even when the full candidate locale is neither. The whole oracle and its selected answer interfaces are therefore retained as separate data.

\subsection{Pointing}
After passing to $\Sh_\E(\widetilde L)$, the generic point of $\widetilde L$ gives the oracle point determined by $t$. This is the pointed generic realization. A coordinate interface which is exhaustive and functional there yields an actual answer map by \cref{prop:answer-reconstruction}. Its descent to $\Sh_\E(L)$ is controlled by the following construction.


A candidate span now records the local sources from which an oracle value may be realized. The next construction tests which observations agree under these choices and can therefore be defined on the original forcing space, even when the candidate itself does not descend.

\section{The universal descended observational fragment}\label{sec:descent}


The question here is whether the specified local
observations agree under descent. This differs from
selecting a name in a locally inhabited naming fibre:
\cref{obs:subsec:scope} relates the two questions, and
\cref{ex:selected-descent} separates them explicitly.

\begin{theorem}[Canonical supported descent for an arbitrary candidate span]
\label{thm:supported-descent}
Let $L\xleftarrow{e}\widetilde L\xrightarrow{t}X$ be any
internal locale span, without a surjectivity assumption. Factor
$e$ through its localic image
\[
 \widetilde L\xrightarrow{q}C\xhookrightarrow{c}L,
\]
so that $q$ is surjective. Construct the locale pushout
$D_C=X\amalg_{\widetilde L}C$. Then
\[
 \OO(D_C)=\{(U,V)\in\OO(X)\times\OO(C):
                                      t^*U=q^*V\}.
\]
Its first frame projection is monic, so its opens classify exactly
those oracle observations whose local interpretations descend to
the \emph{actual image} $C$ of the candidate procedure. Each such
observation has a unique descended value, and $D_C$ has the
universal property of a pushout. If $e$ is surjective, $C\cong L$
and this is the whole-base descent construction below.
For marked source observations $S\incl\OO(L)$, their canonical
supported marking is the regular image
$S_C=\im(S\xrightarrow{c^*}\OO(C))\incl\OO(C)$; the maximal
compatible marking on $D_C$ is then
$M_X\times_{\OO(\widetilde L)}S_C$. This last image need not
select a unique \emph{old} positive name for a marked open of $C$.
\end{theorem}
\begin{proof}
\leavevmode
\begin{enumerate}[label=\textup{(\roman*)},leftmargin=*,itemsep=3pt,topsep=4pt]
\item \textbf{Pullback construction.} The localic image exists by the frame nucleus $e_*e^*$:
its fixed frame defines $C$, and the first factor $q$ has
injective inverse image.
The frame pullback displayed above is a
frame under componentwise joins and finite meets, since $q^*$
and $t^*$ preserve these operations.
It represents the pushout
because a compatible cocone into any locale $Y$ is exactly a
frame homomorphism $\OO(Y)$ into this pullback.
If two compatible
pairs have the same first component, injectivity of $q^*$ forces
their second components to agree.
\item Therefore the first projection
is injective and defines a surjective map $X\to D_C$; the second
component of each pair is its unique descended value.
When $e$ is
surjective, its image mono is invertible by image factorization.
\item \textbf{Image and quotient construction.} For the marked assertion, regular-image factorization constructs
$S_C\incl\OO(C)$ without choosing a representative from $S$.
Pulling this marking and $M_X\incl\OO(X)$ back along the two
inverse-image maps into $\OO(\widetilde L)$ constructs exactly
the compatible marked pairs.
Its universal property proves
maximality among markings on $D_C$ preserving these specified
supported observations.
No preservation of arbitrary positive
joins or existence of an old name selector is inferred.
\end{enumerate}
\end{proof}

For the remainder of this section assume the original left leg
$e:\widetilde L\to L$ is surjective, so the supported construction
specializes to the original forcing base. Form the locale pushout
\[
\begin{tikzcd}
 \widetilde L\ar[r,"t"]\ar[d,"e"']&
 X\ar[d,"a"]\\
 L\ar[r,"b"']&D .
\end{tikzcd}
\]
It exists internally because its frame is the pullback of the two frame maps:
\begin{equation}\label{eq:descent-frame}
 \OO(D)=
 \{(U,V)\in\OO(X)\times\OO(L)\mid t^*U=e^*V\}.
\end{equation}
Joins and finite meets are componentwise; the defining equation is preserved because the two maps are frame homomorphisms.

\begin{theorem}[Universal descent of observations]\label{thm:descent-pushout}
The map $a:X\to D$ is surjective. Its inverse image identifies $\OO(D)$ with the subframe
\[
 \Desc(e,t)=\{U\in\OO(X)\mid t^*U\in\im(e^*)\}.
\]
For $U$ in this subframe there is a unique descended truth $\operatorname{desc}(U)\in\OO(L)$ satisfying $e^*\operatorname{desc}(U)=t^*U$. The map $\operatorname{desc}$ is a frame homomorphism. For compatible maps $X\to Y$ and $L\to Y$, there is a unique induced map $D\to Y$.
\end{theorem}
\begin{proof}
\leavevmode
\begin{enumerate}[label=\textup{(\roman*)},leftmargin=*,itemsep=3pt,topsep=4pt]
\item \textbf{Universal descent of observations.} The first projection of \eqref{eq:descent-frame} is injective: if $U$ agrees, then $e^*V$ agrees, and injectivity of $e^*$ gives equality of $V$.
This projection is $a^*$, proving surjectivity and the identification with its displayed image.
\item \textbf{Pullback construction.} The second projection gives the descended truth.
The two components of the pullback preserve all joins and finite meets, so this is a frame homomorphism on $\Desc(e,t)$.
\item For the final assertion, inverse image of a compatible pair of locale maps gives a pair of frame maps $\OO(Y)\to\OO(X),\OO(L)$ whose composites into $\OO(\widetilde L)$ agree.
Their unique factorization through the pullback frame is the inverse image of the required map.
This proves the pushout universal property explicitly.

The universal property is represented by the commutative diagram
\[
\begin{tikzcd}[column sep=large,row sep=large]
E\ar[r,"e"]\ar[d,"t"']&L\ar[d,"b"]\\ X\ar[r,"a"']&D
\end{tikzcd}
\]
\end{enumerate}
\end{proof}

\begin{corollary}[Finest descended quotient]\label{cor:finest-descended-quotient}
The map $a:X\to D$ is the finest localic quotient of $X$ through which
the pointed realization becomes a map on $L$. Explicitly, for any
locale $Y$, map $q:X\to Y$, and map $\sigma:L\to Y$ satisfying
\[
 qt=\sigma e,
\]
there is a unique $\bar q:D\to Y$ with
$\bar q a=q$ and $\bar q b=\sigma$. If $q$ is surjective, then
$\bar q$ is surjective. Thus every quotient of $X$ to which the local
oracle descends is a further quotient of $D$.
\end{corollary}
\begin{proof}
\leavevmode
\begin{enumerate}[label=\textup{(\roman*)},leftmargin=*,itemsep=3pt,topsep=4pt]
\item \textbf{Finest descended quotient.} The pair $(q,\sigma)$ is a compatible cocone on the pushout defining
$D$, so the pushout universal property produces $\bar q$ and its two
equations, uniquely.
If $q$ is surjective and
$\bar q^*u=\bar q^*v$, then
\[
 q^*u=a^*\bar q^*u=a^*\bar q^*v=q^*v.
\]
Injectivity of $q^*$ gives $u=v$, proving surjectivity of $\bar q$.
\end{enumerate}
\end{proof}

One may calculate the descended value using the frame right adjoint:
\[
 \operatorname{desc}(U)=e_*t^*U
 \quad\text{on }\Desc(e,t),
 \qquad
 U\in\Desc(e,t)\iff e^*e_*t^*U=t^*U.
\]
Indeed $e_*e^*=\id$ for a surjective locale map: the unit and order reflection by $e^*$ give equality. The right adjoint need not preserve arbitrary joins on all opens; its restriction here does, by \cref{thm:descent-pushout}.

\begin{corollary}[Full descent criterion]\label{cor:full-descent}
The local realization descends to an oracle map on $L$ precisely when $a:X\to D$ is an isomorphism, equivalently when every oracle open belongs to $\Desc(e,t)$.
\end{corollary}
\begin{proof}
\leavevmode
\begin{enumerate}[label=\textup{(\roman*)},leftmargin=*,itemsep=3pt,topsep=4pt]
\item \textbf{Full descent criterion.} If every open descends, the frame map $\operatorname{desc}:\OO(X)\to\OO(L)$ defines $\rho$ with $t=\rho e$.
Conversely such a $\rho$ supplies $V=\rho^*U$ for every $U$.
By the pushout description the first projection is then an isomorphism.
\end{enumerate}
\end{proof}

\subsection{Coherence and actual descent}
Let $R=\widetilde L\times_L\widetilde L$ with projections $d_0,d_1$. Define
\[
 \Coh(e,t)=
 \{U\in\OO(X)\mid d_0^*t^*U=d_1^*t^*U\}.
\]
It is a subframe, since it is an equalizer of frame homomorphisms.

\begin{proposition}[Descent and coherence]\label{prop:coherence}
There is a canonical inclusion $\Desc(e,t)\incl\Coh(e,t)$. The span itself is coherent when $td_0=td_1$. If $e$ has a section and the span is coherent, it descends. More generally, if $e$ is the coequalizer of its kernel pair, coherent maps out of $\widetilde L$ descend, and $\Desc(e,t)=\Coh(e,t)$.
\end{proposition}
\begin{proof}
\leavevmode
\begin{enumerate}[label=\textup{(\roman*)},leftmargin=*,itemsep=3pt,topsep=4pt]
\item \textbf{Descent and coherence.} For a descended open, its two pullbacks are equal because $ed_0=ed_1$.
This proves the inclusion.
If $s:L\to\widetilde L$ is a section, the pair $(se,\id)$ factors through $R$.
Coherence then gives $tse=t$, so $\rho=ts$ descends the span.
\item \textbf{Kernel-pair descent.} If $e$ is the kernel-pair coequalizer, the universal property gives descent of coherent maps.
Taking opposites, $e^*$ is the equalizer of $d_0^*,d_1^*$ in frames.
Equalizers of frame homomorphisms are computed by the equalizing subframe.
Therefore an open of $\widetilde L$ descends precisely when its two kernel-pair pullbacks agree.
Apply this to $t^*U$.
\end{enumerate}
\end{proof}

This also proves the equivalence of locale-map descent and frame descent for a surjection: both say that the indicated coequalizer, respectively equalizer, is the original map. We now give a constructive class of maps where this property is available.

\begin{lemma}[Open-map base change]\label{lem:open-bc}
An open locale map $e:Y\to Z$ has a left adjoint $e_!\dashv e^*$ satisfying Frobenius. Its pullbacks are open and their left adjoints satisfy Beck--Chevalley. In its kernel-pair square,
\[
 e^*e_!=d_{1!}d_0^*.
\]
\end{lemma}
\begin{proof}
\leavevmode
\begin{enumerate}[label=\textup{(\roman*)},leftmargin=*,itemsep=3pt,topsep=4pt]
\item \textbf{Open-map base change.} This is the open-map pullback and Beck--Chevalley theorem for
locales, including its constructive Frobenius form
\cite{joyaltierney,johnstone}.
Applying Beck--Chevalley to the
pullback of $e$ along itself yields
$e^*e_!=d_{1!}d_0^*$ with the stated projection convention.
No choice principle is part of the open-map theorem.
\end{enumerate}
\end{proof}

\begin{theorem}[Stable descent for open surjections]\label{thm:open-descent}
An open surjection is a kernel-pair coequalizer, and the same is true after every locale pullback.
\end{theorem}
\begin{proof}
\leavevmode
\begin{enumerate}[label=\textup{(\roman*)},leftmargin=*,itemsep=3pt,topsep=4pt]
\item \textbf{Stable descent for open surjections.} First $e_!1=1$.
The adjunction unit gives $1\leq e^*e_!1$, and injectivity of $e^*$ reflects equality with top.
The pullback remains open by \cref{lem:open-bc}, and its left adjoint also sends top to top by Beck--Chevalley; Frobenius then gives $e'_!e'^*v=v$, proving that it is surjective.
\item \textbf{Kernel-pair descent.} If $u\in\OO(Y)$ has equal kernel-pair pullbacks, then
\[
 e^*e_!u=d_{1!}d_0^*u=d_{1!}d_1^*u
 =u\wedge d_{1!}1=u.
\]
Thus every equalizing open lies in the image of $e^*$.
The converse is immediate from the pullback square.
Hence $e^*$ is the equalizer and $e$ is its opposite coequalizer.
Apply the same argument to each open surjective pullback.

The defining maps satisfy the following commutative diagram
\[
\begin{tikzcd}[column sep=large,row sep=large]
E\times_XE\ar[r,shift left=.6ex,"\pi_1"]\ar[r,shift right=.6ex,"\pi_2"']&E\ar[r,"e"]&X
\end{tikzcd}
\]
\end{enumerate}
\end{proof}

For general localic surjections, \cref{thm:descent-pushout} constructs maximal descended observations directly; the image-corrected comparison of \cref{thm:descended-basechange} also exists after arbitrary localic base change. Kernel-pair coequalizer descent is a separate, stronger property proved for open surjections.

\begin{theorem}[Change of localic base through the covered image]\label{thm:descended-basechange}
Let $r=(L\xleftarrow{e}\widetilde L\xrightarrow{t}X)$ be a local realization span with surjective left leg, and let $h:L'\to L$ be any locale map. Form the pulled-back span
\[
 \widetilde L'=L'\times_L\widetilde L,
 \qquad e':\widetilde L'\to L',\qquad t':\widetilde L'\to X.
\]
Factor $e'$ through its localic image
$\widetilde L'\xrightarrow{q}C_h\xrightarrow{c}L'$, with $q$ surjective. Let $D_r$ be the original descended-observation locale, and $D_h^{\mathrm{im}}$ the pushout of
$C_h\xleftarrow{q}\widetilde L'\xrightarrow{t'}X$.
There is a canonical surjective locale map
\[
 \beta_h^{\mathrm{im}}:D_h^{\mathrm{im}}\longrightarrow D_r
\]
with inverse-image formula
\begin{equation}\label{eq:descended-basechange-formula}
 (\beta_h^{\mathrm{im}})^*(U,V)=(U,(hc)^*V).
\end{equation}
If $e'$ is surjective (in particular when $e$ is stably surjective), then $C_h=L'$ and this is the original comparison $(U,V)\mapsto(U,h^*V)$. In the stably surjective case these comparisons are identities for $h=\id$ and compose coherently under further localic base change. No stability property of $e$ is needed for the image-corrected construction.
\end{theorem}
\begin{proof}
\leavevmode
\begin{enumerate}[label=\textup{(\roman*)},leftmargin=*,itemsep=3pt,topsep=4pt]
\item \textbf{Pullback construction.} An open $(U,V)$ of $D_r$ satisfies $t^*U=e^*V$.
Pull this equality back along $\widetilde L'\to\widetilde L$ and use $e'=cq$ to obtain
\[
 t'^*U=e'^*h^*V=q^*(hc)^*V.
\]
Therefore $(U,(hc)^*V)$ belongs to the pullback frame for $D_h^{\mathrm{im}}$.
The assignment preserves finite meets and arbitrary joins componentwise and hence defines $\beta_h^{\mathrm{im}}$.
It is injective on frames because equality of outputs implies equality of the first components $U$, and the first projection of $\OO(D_r)$ is monic by surjectivity of the original $e$.
\item Thus the constructed locale map is surjective.
It satisfies the cocone equations on both sides by applying inverse image to the displayed pair.
If $e'$ is surjective, its image inclusion $c$ is an isomorphism, giving the stated formula over $L'$.
For two successive base changes in the stable case, the direct and sequential maps agree on pairs by
$(U,V)\mapsto(U,k^*h^*V)$; this establishes coherence.
All claims follow from the constructed image factorization and the frame pullback, without inferring arbitrary pullback stability of localic surjections.
\end{enumerate}
\end{proof}

\section{Effective descent and a branching candidate example}\label{sec:effective-descent}


The effective marking below is obtained from the
actual observation maps. A group action is needed before
a local choice can be treated as a torsor; for arbitrary
epic sheaf name covers the canonical augmentation
construction of \cref{obs:prop:augmentation} provides
the corresponding abelian obstruction.

Let $L=L_j$ carry its stable marking $S_j$, and let $X$ carry $M_X$. The pushout has a canonical marking
\[
 M_D=M_X\times_{\OO(\widetilde L)}S_j,
\]
where the maps are $t^*m_X$ and $e^*s_j$. It is a subobject of $\OO(D)$ through \eqref{eq:descent-frame}.

\begin{theorem}[Maximal effective descended marking]\label{thm:marked-descent}
The marking $M_D$ is the largest marking on $D$ for which both maps $X\to D$ and $L\to D$ preserve the specified observations. A selected effective oracle family descends effectively exactly when it factors through the projection $M_D\to M_X$. The pushout is universal also for compatible observation-preserving maps.
\end{theorem}
\begin{proof}
\leavevmode
\begin{enumerate}[label=\textup{(\roman*)},leftmargin=*,itemsep=3pt,topsep=4pt]
\item \textbf{Maximal effective descended marking.} A full open of $D$ is a compatible pair $(U,V)$ with $t^*U=e^*V$.
Pull back the product marking $M_X\times S_j\incl\OO(X)\times\OO(L)$ along this compatible-pair inclusion.
The result is exactly $M_D$.
Thus \cref{thm:admissibility-coreflection} gives maximality, the familywise factorization criterion and parameter coherence.
\item \textbf{Pullback construction.} For compatible observation-preserving maps $X,L\to Y$, the ordinary pushout gives $D\to Y$. A marked open of $Y$ pulls back to marked opens $U,V$ on the two sides, and compatibility says that they have the same pullback to $\widetilde L$.
Thus the induced open of $D$ belongs to $M_D$.
This proves the effective universal property.
Equality and coherence follow from the ordinary universal property and the monic markings.
\end{enumerate}
\end{proof}


\begin{theorem}[Positive operations on the descended marking]
\label{thm:positive-descended-joins}
The maximal descended marking $M_D$ always inherits top and finite
meets when the two markings $M_X$ and $S_j$ carry those operations.
If $M_X$ has its supplied overt-indexed joins and $j$ has a global
positive saturation restriction $j_M$ on the source marking, then
$M_D$ has the same overt-indexed joins, with both inverse-image
marking projections preserving them.

Without a global restricted saturation, for each indexing object
$I$ there is nevertheless a canonical maximal subobject
\[
 Q_I^{\bigvee}= (M_D)^I
    \times_{\OO(D)}M_D,
\]
where the map $(M_D)^I\to\OO(D)$ takes the join of the underlying
full-frame family. On this subobject the full join has a unique
marked representative. For an overt $I$, this is precisely the
maximal domain of the requested positive indexed-join operation.
\end{theorem}
\begin{proof}
\leavevmode
\begin{enumerate}[label=\textup{(\roman*)},leftmargin=*,itemsep=3pt,topsep=4pt]
\item \textbf{Positive operations on the descended marking.} Since $e$ is surjective, $e^*:\OO(L)\to\OO(\widetilde L)$
is an order embedding.
Therefore its composite with the mono
$s_j:S_j\incl\OO(L)$ is monic and order-reflecting.
The
pullback projection $\pr_X:M_D\to M_X$ is consequently a
monomorphism and reflects the inherited frame order.
The other
projection need not be monic; none of the arguments uses that
claim.
\item \textbf{Localization and coverage.} Top and finite meets on the two markings give compatible pairs
because $t^*$ and $e^*$ preserve these operations; their equations
are reflected by $\pr_X$.
For a relatively overt-indexed family
$(U_i,V_i)$ in $M_D$, form the supplied marked join
$U=\bigvee_i^{M_X}U_i$.
The restricted source nucleus provides
a marked fixed-frame join
\[
 V=\bigvee_i^{S_j}V_i,
 \qquad
 e_j(V)=j_M\!\left(\bigvee_i^{M}e_j(V_i)\right),
\]
where the join in $M$ uses the positive existential over the
\emph{index object} $I$.
Its ambient fixed-frame value is
$\bigvee_i^{\Omega_j}s_j(V_i)$.
\item \textbf{Universal-property argument.} Since every original pair was
compatible,
\[
 \begin{aligned}
 t^*m_XU
 &=\bigvee_i t^*m_XU_i
  =\bigvee_i e^*s_jV_i\\
 &=e^*\!\left(\bigvee_i^{\Omega_j}s_jV_i\right)
  =e^*s_jV.
 \end{aligned}
\]
Thus $(U,V)$ is a member of $M_D$ and its full open is the join
in $\OO(D)$.
If another pair $(U',V')$ bounds all inputs, then
$U\leq U'$ by the target join property, and the order-reflecting
mono $\pr_X$ gives $(U,V)\leq(U',V')$.
This proves the
universal join property and its equations, as well as preservation
by both marking projections.
No overtness of the source basis is
used in forming this overt-\emph{indexed} join.
\item \textbf{Pullback construction.} Without global source saturation, take the pullback of
$M_D\incl\OO(D)$ along the ordinary full-frame join
$(M_D)^I\to\OO(D)$.
Its first projection is monic; the pullback
classifies precisely those indexed families whose full join has a
marked representative.
This is the maximal domain, with unique
representative and stability under parameter substitution.
It
does not assert that this partial operation is total.
\end{enumerate}
\end{proof}


The first projection is injective because the descended truth is unique. It consequently identifies $M_D$ with the exact subobject of oracle observations that both descend and have an effective stable truth on the original forcing presentation. This construction retains the full support space, the local answer data, and the selected effective observations separately.

\begin{proposition}[Descended predicates and their oracle modalities]\label{prop:modal-descent}
For a local realization span $L\xleftarrow{e}\widetilde L\xrightarrow{t}X$ with $e$ surjective, let $D$ be the descended-observation pushout and let $P:Q\to M_D$ be a selected indexed family in its effective marking. Write
\[
 P_X=m_X\pr_XP:Q\to\OO(X),\qquad
 P_L=s_j\pr_LP:Q\to\OO(L).
\]
Then $t^*P_X=e^*P_L$ as maps $Q\to\OO(\widetilde L)$. These are the same internal propositional oracle after interpretation in $\Sh_\E(\widetilde L)$, so they generate the same Ahman--Bauer oracle modality there. The forcing-side family $P_L$ has its displayed effective representatives in $S_j$ and hence a uniform stage semidecision. This does not imply that the \emph{modality generated by} $P_L$ preserves the whole effective marking.
\end{proposition}
\begin{proof}
\leavevmode
\begin{enumerate}[label=\textup{(\roman*)},leftmargin=*,itemsep=3pt,topsep=4pt]
\item \textbf{Pullback construction.} Membership of $P$ in the pullback $M_D=M_X\times_{\OO(\widetilde L)}S_j$ is precisely the indexed equality $t^*m_X\pr_XP=e^*s_j\pr_LP$.
Under relative truth classification these two families classify the same subobjects of the pulled-back query object in $\Sh_\E(\widetilde L)$.
The characterization of oracle modalities as least local operators forcing an internal predicate, proved in \cref{thm:ab-comparison}, therefore identifies their generated modalities.
\item \textbf{Evaluation and transposition.} The second projection of $P$ is already a lift through $S_j$, and its first-level semidecision is \cref{prop:stable-universal} and the stage evaluation in \cref{sec:markings}.
No saturation map for the newly generated nucleus has been constructed, so no further effective-preservation conclusion follows.
\end{enumerate}
\end{proof}

\begin{proposition}[Candidate-image invariance of descended observations]\label{prop:candidate-descent}
Replacing $(e,t)$ by the two projections $(c,u)$ of its candidate image leaves $\Desc(e,t)$ and $M_D$ canonically unchanged.
\end{proposition}
\begin{proof}
\leavevmode
\begin{enumerate}[label=\textup{(\roman*)},leftmargin=*,itemsep=3pt,topsep=4pt]
\item \textbf{Pullback construction.} The equation $t^*U=e^*V$ is $q^*u^*U=q^*c^*V$.
Since $q$ is surjective, its inverse image is injective and this is equivalent to $u^*U=c^*V$.
Thus the two pullback frames are the same subobject of $\OO(X)\times\OO(L)$.
Restricting the factors to $M_X$ and $S_j$ gives the same statement for their markings.
This identifies descended observations, while the kernel pair of $q$ still records the histories forgotten in passing to the candidate image.
\end{enumerate}
\end{proof}

\subsection{Comparison with a pointed effective lift}
Suppose $\widetilde L$ has its own chosen forcing presentation and stable marking $\widetilde S\incl\OO(\widetilde L)$. A selected oracle family $K:T\to M_X$ may first have a pointed effective lift
$\widetilde K:T\to\widetilde S$ with underlying full truth $t^*m_XK$. Its effective descent is a map $\bar K:T\to S_j$ satisfying
\[
 e^*s_j\bar K=\widetilde s\,\widetilde K.
\]
No effective restriction of $e^*$ on all of $S_j$ is needed to state this equation.

\begin{proposition}[Effective descent of a pointed family]\label{prop:pointed-effective-descent}
The displayed equation has at most one solution. It is equivalent to the factorization of $K$ through $M_D\to M_X$. If an effective inverse-image map $e_\Sigma^*:S_j\to\widetilde S$ has been constructed, it is also equivalent to
$e_\Sigma^*\bar K=\widetilde K$.
\end{proposition}
\begin{proof}
\leavevmode
\begin{enumerate}[label=\textup{(\roman*)},leftmargin=*,itemsep=3pt,topsep=4pt]
\item \textbf{Effective descent of a pointed family.} The composite $e^*s_j$ is monic, so it reflects equality of solutions.
Substituting the pointed lifting equation makes the displayed condition
$e^*s_j\bar K=t^*m_XK$, which is exactly the pullback condition defining $M_D$.
If $e_\Sigma^*$ exists, both sides of its proposed equality have the same composite with the monic $\widetilde s$ precisely when the full equation holds.
Monicity proves equivalence.
\end{enumerate}
\end{proof}

For a site presentation of $L$, the descended family gives the semidecision
$(p,t)\mapsto e_J\bar K(t)(p)$. For a searchable information presentation, it gives
$(c,t)\mapsto\Obs_C(c,e_j\bar K(t))$.
Both follow by evaluation of the constructed stable lift. They make no assertion about the effective complexity of other descended opens.

\begin{example}[Effective observation without oracle descent]\label{ex:selected-descent}
In $\Eff$ take $L=1$, $\widetilde L=\disc(2)$, $X=\disc(2)$, with $e$ terminal and $t$ the identity. This is an open surjective local realization. There is no $\rho:1\to\disc(2)$ with $\rho e=t$, since either composite is constant. The observation $1_X$, however, pulls back to $1_{\widetilde L}$ and descends to $1_L$. The top observation belongs to the canonical positive marking, so the selected descent is effective. Thus effective descent of a selected family is strictly weaker than descent of the local oracle.
\end{example}

The coherent observations form a subframe because the two kernel-pair inverse images preserve joins and finite meets. The descended observations form a subframe for the same reason applied to the pullback \eqref{eq:descent-frame}. If one starts with a selected family and closes it under geometric operations, its descended fragment is the intersection of that generated subframe with $\Desc(e,t)$. Its effective part is then selected by the marking pullback. This construction avoids calling a positive fragment a frame when it does not have arbitrary saturated joins.

\subsection{Two recurrent colour candidates}
We first construct the recurrent-colour realization span. \Cref{sec:ramsey-construction} then constructs the prehomogeneous input and a recurrent colour in explicitly specified dense oracle localizations. Fix a total colouring $c:[\N]^2\to2$. A finite prehomogeneous stem is $(s,\lambda)$, where $s$ is a finite increasing sequence, $\lambda:|s|\to2$, and
\[
 c(s(r),s(t))=\lambda(r)\qquad(r<t<|s|).
\]
Present the locale $B$ of total prehomogeneous sequences by these finite-prefix generators and their coordinate-extension relations. This is a geometric presentation: the parameters and finite relations are fixed, and totality is the family of cover equations at each coordinate. Its generic coordinates are increasing $g:\N\to\N$ and $\lambda:\N\to2$ with the displayed pair relation.

For each $i\in2$, let $B_i\incl B$ be the localization imposing
\[
 \bigvee_{r\geq n}[\lambda(r)=i]=1\qquad(n\in\N).
\]
Thus the colour $i$ recurs without bound. Define $L=B_0\vee B_1$ as the join of these sublocales in $B$, and put $\widetilde L=B_0+B_1$. The canonical map $e:\widetilde L\to L$ is surjective.

\begin{lemma}[Surjectivity of pointing]\label{lem:ramsey-pointing}
The last assertion follows from the frame formula for the join of two sublocales.
\end{lemma}
\begin{proof}
\leavevmode
\begin{enumerate}[label=\textup{(\roman*)},leftmargin=*,itemsep=3pt,topsep=4pt]
\item \textbf{Surjectivity of pointing.} Let $k_0,k_1$ be their nuclei on $\OO(B)$.
Their join as sublocales is the fixed frame of $k=k_0\wedge k_1$.
The inverse image of the coproduct map sends a $k$-fixed $u$ to $(k_0u,k_1u)$.
If these pairs agree for $u,v$, then
$u=k_0u\wedge k_1u=k_0v\wedge k_1v=v$.
The inverse image is injective.
\end{enumerate}
\end{proof}

Inside $\Sh(B_i)$, recurrence constructs the increasing enumeration $r_0<r_1<\cdots$ of the indices with $\lambda(r)=i$. For each lower bound, the next index exists by recurrence and has a least value because the predicate is decidable on $\N$. The least value is defined by a total functional relation; recursion gives the entire enumeration. Put $h_i(n)=g(r_n)$. For $m<n$,
\[
 c(h_i(m),h_i(n))=\lambda(r_m)=i.
\]
Hence $h_i$ is a total increasing homogeneous sequence.

\paragraph{The homogeneous-sequence classifier.}
Write $\operatorname{IncFin}(\N)$ for the finite strictly increasing
sequences and $\Inc(\N)$ for their total-increasing-sequence locale,
presented by prefix cylinders and the coordinate-extension covers used
above. The locale $\operatorname{Homog}(c)$ is the sublocale of
$\Inc(\N)\times\disc(2)$ classified by pairs $(h,i)$ satisfying
\[
 \forall m<n:\N\qquad c(h(m),h(n))=i.
\]
Explicitly, start with the prefix cylinders $[\tau]$ and colour opens
$[i]$, impose the increasing-prefix intersection and totality equations,
the two-colour partition $[0]\vee[1]=1$, $[0]\wedge[1]=0$, and impose
\[
 [\tau]\wedge[i]=0
 \qquad\text{if some }m<n<|\tau|\text{ has }
                  c(\tau(m),\tau(n))\ne i.
\]
These relations give a geometric frame presentation: each forbidden
finite pattern is eliminated by an equation between opens. In every
test locale a model consists of an increasing coordinate sequence and
a colour satisfying the displayed pair equations. The universal
property of this presentation therefore constructs maps into
$\operatorname{Homog}(c)$ from exactly such coordinate interpretations,
without a choice of a global recurrent colour. The colour coordinate is
uniquely recovered from a homogeneous sequence as
$i=c(h(0),h(1))$, so this presentation also identifies canonically
with the uncoloured homogeneous-sequence classifier.

\begin{proposition}[The local homogeneous realization]\label{prop:ramsey-realization}
The coordinate interpretations of the $h_i$ give a map
$t:B_0+B_1\to\operatorname{Homog}(c)$, and hence a local realization span
\[
 L\ \xleftarrow{e}\ B_0+B_1\ \xrightarrow{t}\ \operatorname{Homog}(c).
\]
The universal descended frame and its effective selected fragment are those of \cref{thm:descent-pushout,thm:marked-descent}.
\end{proposition}
\begin{proof}
\leavevmode
\begin{enumerate}[label=\textup{(\roman*)},leftmargin=*,itemsep=3pt,topsep=4pt]
\item \textbf{The local homogeneous realization.} The preceding computation verifies totality, increasingness, and the homogeneous coordinate equations in each summand.
Their geometric presentation therefore gives a locale map on each summand, and coproduct universality gives $t$.
The left leg is surjective by \cref{lem:ramsey-pointing}, so the general descent constructions apply.
\end{enumerate}
\end{proof}

This construction does not require identifying $L$ with all of $B$ or choosing a recurrent colour globally. Pointing supplies the colour context in which the coordinate reconstruction is valid. The existence of a homogeneous candidate before that choice is represented by a relative locale. Finite recurrence witnesses give finite computational transcripts for the pointed extraction; their effective descent is tested by the actual marking $M_D$.


\subsection{Constructed prehomogeneous and homogeneous oracles}\label{sec:ramsey-construction}
Fix a total colouring $c:[\N]^2\to k$, where $1\leq k<\infty$.
It is a parameter, not a uniform computability assertion about all
colourings. All computations here are relative to its coordinate map.
A state $(s,\lambda)$ consists of a finite increasing sequence and a
same-length colour sequence with
$c(s(a),s(b))=\lambda(a)$ for $a<b<|s|$. Put
\[
 \begin{split}
 R_{s,\lambda}(x)&\ \Longleftrightarrow\
       (\forall a<|s|\ s(a)<x)\ \wedge\
       (\forall a<|s|\ c(s(a),x)=\lambda(a)),\\
 T(s,\lambda,n)&\ \Longleftrightarrow\ \exists x\geq n\,R_{s,\lambda}(x),\\
 F(s,\lambda)&\ \Longleftrightarrow\ \exists n\,\neg T(s,\lambda,n).
 \end{split}
\]
The empty state has $R=\N$. The predicate $R$ is decidable by finitely
many colour queries. The predicates $T$ and $F$ are explicit arithmetic
queries, the latter recording boundedness.

\begin{lemma}[Correct interpretation of the boundedness questions]\label{ramsey:lem:boundedness}
After jointly adjoining decisions for all the displayed $T$ and $F$,
their sheafified interpretations still satisfy
$F(s,\lambda)=\exists n\,\neg T(s,\lambda,n)$, where $T$ is the
actual tail-existence predicate for the transported decidable $R$.
\end{lemma}
\begin{proof}
\leavevmode
\begin{enumerate}[label=\textup{(\roman*)},leftmargin=*,itemsep=3pt,topsep=4pt]
\item \textbf{Correct interpretation of the boundedness questions.} The definition of $T$ is an existential image of a decidable finite
predicate, so geometric inverse image preserves it.
The decision of
$T$ makes the transported $\neg T$ its actual complement, by
\cref{prop:decision}.
Existential images are again preserved, proving
the formula for $F$.
This is why both levels of questions were included;
preservation of a fresh negative quantifier is not being presumed.
\end{enumerate}
\end{proof}

\begin{theorem}[Uniform prehomogeneous construction]\label{ramsey:thm:prehomogeneous}
In this first dense decision extension recursion constructs a total increasing
$g:\N\to\N$ and a map $\lambda:\N\to k$ satisfying
\[
 c(g(a),g(b))=\lambda(a)\qquad(a<b).
\]
\end{theorem}
\begin{proof}
\leavevmode
\begin{enumerate}[label=\textup{(\roman*)},leftmargin=*,itemsep=3pt,topsep=4pt]
\item \textbf{Uniform prehomogeneous construction.} The invariant is a valid finite state with $\neg F(s,\lambda)$.
For every $n$, $\neg T(s,\lambda,n)$ would witness $F$, so
$\neg\neg T(s,\lambda,n)$ holds.
The decision of $T$ gives $T$.
The decidable predicate $R$ thus has a least element $x$; find it by
search.
For each $i<k$ form the child $(sx,\lambda i)$.
\item \textbf{Verification of the equations.} If all children satisfied $F$, take their finitely many bounds and a
common upper bound $M>x$.
Any parent-residual $y\geq M$ has one
colour $c(x,y)=i$, so it would lie in that child's residual beyond its
bound.
Thus there is no such $y$, giving $F(s,\lambda)$, a
contradiction.
The finitely many child predicates $F$ are decidable;
choose the first false one.
The new state is valid and has the same
invariant.
This is a total successor function.
The empty state satisfies
the invariant, since $R=\N$ there.
Natural-number recursion constructs
the required sequences.
The residual equations retained at each step
give the pair-colour equation for every $a<b$.
\end{enumerate}
\end{proof}

In this extension form the second family
\[
 T_\lambda(i,n)=\exists r\geq n\,[\lambda(r)=i],\qquad
 F_\lambda(i)=\exists n\,\neg T_\lambda(i,n),
\]
and adjoin their decisions. They have the correct fresh interpretation
by the proof of \cref{ramsey:lem:boundedness}, now in the first extension.

\begin{theorem}[Constructed homogeneous sequence]\label{ramsey:thm:homogeneous}
The composite dense oracle extension contains a canonically constructed
colour $i<k$ and a total increasing sequence $h$ with
$c(h(a),h(b))=i$ whenever $a<b$.
\end{theorem}
\begin{proof}
\leavevmode
\begin{enumerate}[label=\textup{(\roman*)},leftmargin=*,itemsep=3pt,topsep=4pt]
\item \textbf{Constructed homogeneous sequence.} All colours cannot satisfy $F_\lambda$: a common bound would leave
$\lambda$ with no value at that bound.
Decide the finitely many
$F_\lambda(i)$ and choose the first false one.
As before, the decisions
of $T_\lambda$ turn $\neg F_\lambda(i)$ into
$\forall n\,T_\lambda(i,n)$.
Set $r_0$ equal to the first occurrence
of $i$, and $r_{a+1}$ equal to the first occurrence strictly after
$r_a$.
Each search terminates and is unique, so recursion defines the
sequence.
Set $h(a)=g(r_a)$.
Then
$c(h(a),h(b))=\lambda(r_a)=i$ for $a<b$.
Density of both inclusions
was proved in \cref{prop:decision}.
\end{enumerate}
\end{proof}

Applied to the generic sequence of the preceding prehomogeneous
locale $B$, the second construction supplies an actual recurrent colour
after its specified localization. Thus its two recurrent-colour
sublocales jointly cover that localized universe, and its pointing map
there has a section. Applied to the original colouring, the first
construction supplies a prehomogeneous sequence before the recurrent-colour extraction. Neither assertion identifies the unlocalized sublocale
$B_0\vee B_1$ with $B$. No optimal bound on the number of jumps is
claimed: the exact queries used above are the resource specification.

\subsection{A geometric presentation of Ramsey histories}
The preceding algorithm has a geometric history-to-solution map. For its
given total colouring $c:[\N]^2\to k$, define $\mathrm{Sol}_c$ as the locale of pairs $(i,h)$
where $i<k$, $h:\N\to\N$ is strictly increasing, and
\[
 a<b\ \Longrightarrow\ c(h(a),h(b))=i.
\]
Define $\mathrm{Hist}_c$ by tuples $(g,\lambda,i,r)$ with $g,r$ strictly
increasing, $\lambda:\N\to k$, and
\[
 c(g(a),g(b))=\lambda(a)\quad(a<b),\qquad
 \lambda(r(n))=i\quad(n\in\N).
\]
These locales have geometric presentations by finite-coordinate generators,
functionality and totality covers, and one finite-coordinate relation for
each displayed tuple of numerical indices. They are formed natively in the
chosen base, with explicit comparison to the pulled-back exponential.

\begin{proposition}[A split solution map]\label{cal:prop:ramsey}
There are locale maps
\[
 \mathrm{Hist}_c\xrightarrow{\pi}\mathrm{Sol}_c
 \xrightarrow{s}\mathrm{Hist}_c,
 \quad
 \pi(g,\lambda,i,r)=(i,g\circ r),\quad
 s(i,h)=(h,\mathrm{const}_i,i,\id),
\]
with $\pi s=\id$. In particular $\pi$ is a split surjection of locales.
The construction of \cref{ramsey:thm:prehomogeneous,ramsey:thm:homogeneous} supplies a point of
$\mathrm{Hist}_c$ in its constructed extension and hence a point of
$\mathrm{Sol}_c$ there.
\end{proposition}
\begin{proof}
\leavevmode
\begin{enumerate}[label=\textup{(\roman*)},leftmargin=*,itemsep=3pt,topsep=4pt]
\item \textbf{A split solution map.} If $a<b$, then $r(a)<r(b)$ and
$c(g(r(a)),g(r(b)))=\lambda(r(a))=i$.
Thus $\pi$ respects the solution
relations.
Its inverse image on a coordinate generator is explicitly
\[
 \pi^*[h(n)=m]=\bigvee_j([r(n)=j]\wedge[g(j)=m]),
\]
so this is a geometric map on the presented locales, not just a map on their
set-valued points.
The equations for $s$ verify its defining relations and
give $\pi s=\id$ coordinatewise.
The composite of inverse images is the
identity, so $\pi^*$ is injective, which is locale surjectivity.
\item \textbf{Verification of the equations.} \Cref{ramsey:thm:prehomogeneous} produces the prehomogeneous pair $(g,\lambda)$.
\Cref{ramsey:thm:homogeneous} selects an infinitely recurring colour $i$ and the increasing
enumeration $r$ of selected occurrences.
Those proved equations are exactly
the defining relations of $\mathrm{Hist}_c$, establishing the final assertion.
\end{enumerate}
\end{proof}

This construction gives an exact structured solution presentation. It does
not discard the additional reservoir or partition history of a specific
forcing, and it does not improve the jump bound of the preceding two
decision stages. Those are separate computational questions.

\section{Intensional enhancement: descent of answers and descriptions}
\label{sec:intensional-descent}


Let $L\xleftarrow{e}\widetilde L\xrightarrow{t}X$ be a local
realization span. The descended-observation frame records when a
full semantic truth factors across $e^*$. It does not decide whether
a \emph{description of that truth}, or an answer-dependent
continuation, factors across $e$. The distinction can be expressed
by applying the same kernel-pair test to the actual code carrier.

\begin{theorem}[Separate semantic and code descent]
\label{thm:intensional-code-descent}
Suppose $e$ is an open surjection and $C$ is a displayed code
object in the base. Let $z:\widetilde L\to\disc(C)$ be a locally
supplied code, with denotation $q:C\to Y$ into a displayed
observable type. Write
$R=\widetilde L\times_L\widetilde L$ with projections $d_0,d_1$.
Then:
\begin{enumerate}
\item $z$ descends to a unique code map $\bar z:L\to\disc(C)$
if and only if $zd_0=zd_1$;
\item its denotation $\disc(q)z$ descends if and only if
$\disc(q)zd_0=\disc(q)zd_1$;
\item descent of $z$ implies descent of its denotation, but the
converse need not hold when $q$ identifies distinct codes.
\end{enumerate}
For arbitrary surjective $e$, the exact code and denotation
fragments can instead be constructed as the respective pullback
frames of \cref{thm:descent-pushout}; no unproved effective-descent
condition for $e$ is used. When a semantic oracle observation also
has a marked representative, the pullback $M_D$ adds its exact
positive-effectivity condition.
\end{theorem}
\begin{proof}
\leavevmode
\begin{enumerate}[label=\textup{(\roman*)},leftmargin=*,itemsep=3pt,topsep=4pt]
\item \textbf{Kernel-pair descent.} By \cref{thm:open-descent}, $e$ is the coequalizer of $d_0,d_1$.
Its universal property gives the first statement, with existence
and uniqueness of $\bar z$.
Applying the same property to
$\disc(q)z$ gives the second.
Applying $\disc(q)$ to the
code-coherence equation proves the implication in (3).
For strictness, in $\Eff$ take $L=1$, $\widetilde L=\disc(2)$,
$e$ terminal, $C=2$, $z=\id$, and $Y=1$ with $q$ terminal.
The semantic denotation descends, whereas $z$ does not.
\item \textbf{Pullback construction.} For arbitrary surjections, the respective locale pushouts have
frames of compatible pairs of opens by \cref{eq:descent-frame};
the construction does not presume kernel-pair descent.
Pullback
of the original and stable markings produces $M_D$ by
\cref{thm:marked-descent}.
\end{enumerate}
\end{proof}

The same proof works for an oracle response code that is an actual
discrete-valued observable on the local candidate locale. A
successful answer transcript can be substituted for $C$ once
its interpretation as a displayed object and the corresponding
locale map are supplied. For the Ramsey construction the branch
index $i\in2$ is retained in $B_0+B_1$, while the homogeneous
sequence is interpreted through the map to $\operatorname{Homog}(c)$.
Their kernel-pair equations are separate checks; the existence of
the homogeneous output on the pointed cover supplies no global
colour selector. This is the precise intensional counterpart of
the semantic descent theorem.

\part{Naming obstructions and finite deformations}\label{part:naming-obstructions}

The previous descent constructions identify when local observations agree. Even then, locally available names need not glue to a global selector. We measure this failure by effective descent, torsors, and abelian localization defects, then use finite sites to calculate concrete naming and deformation obstructions.


\section{Abelian recollement and localization defects}
\label{obs:subsec:recollement}

The glued computational structure has an open semantic component and a closed raw component. Passing to abelian-group objects turns this separation into exact sequences that measure what localization forgets. We calculate the associated adjunctions and tensor comparisons before using them to study finite naming defects.


Write $\operatorname{Ab}(\mathcal T)$ for the category of internal
abelian groups. The functor $i$ induces a left exact functor on
these categories, again denoted $i$, and
\[
 \operatorname{Ab}(\mathcal G_j)
 \simeq
 \bigl(\operatorname{Ab}(\mathcal R)
       \downarrow i:\operatorname{Ab}(\mathcal S)
                      \to\operatorname{Ab}(\mathcal R)\bigr).
\]
Indeed, a group law in the comma topos is exactly a group law on
each component for which the comparison arrow is a homomorphism.
The following is the left-exact comma-category recollement
\cite[Section~3.2]{franjou-pirashvili}, with its functors and
exact sequences calculated for the localization in this paper.

\begin{proposition}[Six adjoints for the glued abelian category]
\label{obs:prop:recollement}
For $F=(A,B,\alpha:A\to iB)$ the formulas
\[
 \begin{array}{lll}
 o_!B=(0,B,0),&o^*F=B,&o_*B=(iB,B,\id),\\
 c^*F=A,&c_*A=(A,0,0),&c^!F=\ker\alpha
 \end{array}
\]
give adjunctions
\[
 o_!\dashv o^*\dashv o_*,
 \qquad c^*\dashv c_*\dashv c^!.
\]
The three functors $o_!,o_*,c_*$ are fully faithful, and the
essential image of $c_*$ is the kernel of $o^*$.
There are natural exact sequences
\begin{align}
 0&\longrightarrow o_!B\longrightarrow F
       \longrightarrow c_*A\longrightarrow0,
       \label{obs:eq:open-closed}\\
 0&\longrightarrow c_*\ker\alpha\longrightarrow F
       \longrightarrow o_*B
       \longrightarrow c_*\operatorname{coker}\alpha
       \longrightarrow0.
       \label{obs:eq:defects}
\end{align}
In particular, for the retained abelian presentation
$\Gamma_jA=(A,aA,\eta_A)$,
\begin{equation}\label{obs:eq:unit-defects}
 0\longrightarrow c_*\ker\eta_A\longrightarrow\Gamma_jA
 \longrightarrow o_*aA
 \longrightarrow c_*\operatorname{coker}\eta_A
 \longrightarrow0,
\end{equation}
and both defect groups are annihilated by $a$.
\end{proposition}
\begin{proof}
\leavevmode
\begin{enumerate}[label=\textup{(\roman*)},leftmargin=*,itemsep=3pt,topsep=4pt]
\item \textbf{Six adjoints for the glued abelian category.} A map $(A,B,\alpha)\to(A',B',\alpha')$ is a pair $(u,v)$
with $iv\,\alpha=\alpha'u$.
From $o_!B$ the first component is
zero and the second is unrestricted.
Into $o_*B'$ the first is
forced by the second.
Into $c_*A'$ only the first component
remains.
From $c_*A'$ it must factor through $\ker\alpha$.
These four descriptions prove the adjunctions; their units or
counits on the indicated components are identities and prove
full faithfulness.
An object with $B=0$ has its unique map
$A\to i0=0$, giving the kernel statement.
\item \textbf{Adjunction and universal property.} Kernels are componentwise since $i$ is left exact.
Cokernels
are also componentwise: for a cokernel $B'\to D$, the composite
$A'\to iB'\to iD$ annihilates the first component of the given
map and factors uniquely through its cokernel.
This supplies
the comparison arrow and its universal property.
Thus exactness is detected on the two components.
The maps in \eqref{obs:eq:open-closed} are $(0,\id)$ and
$(\id,0)$; those around $F\to o_*B$ in
\eqref{obs:eq:defects} use $(\alpha,\id)$.
\item \textbf{Universal-property argument.} Their component
kernels and cokernels give the displayed sequences, which are
the corresponding recollement sequences of
\cite[Section~4.1]{franjou-pirashvili}.
\item \textbf{Inverse comparison.} Sheafification on abelian groups is exact: it preserves finite
limits and cokernels.
Applying it to $\eta_A$ gives an
isomorphism by the adjunction identities.
Consequently
$a(\ker\eta_A)=0=a(\operatorname{coker}\eta_A)$, proving the
last assertion.
This identifies the precise abelian defects of
the raw-to-semantic comparison without discarding the raw object.

The defining maps satisfy the following commutative diagram
\[
\begin{tikzcd}[column sep=large,row sep=large]
\mathcal A(\mathcal S)\ar[r,"o_!"]\ar[dr,"\id"']&\mathcal A(\mathcal G_j)\ar[d,"o^*"]\\ &\mathcal A(\mathcal S)
\end{tikzcd}
\]
\end{enumerate}
\end{proof}

For completeness the tensor operations also have an explicit
description. Form the tensor product of internal abelian groups
as the free abelian object on pairs, quotiented by the two
bilinearity relations. This uses the NNO and exact quotients.
The internal group of additive maps is the corresponding
subobject of the exponential, giving the tensor--internal-Hom
adjunction. The inverse images $r,s$ preserve the constructions,
so the tensor product in the gluing is componentwise; its
comparison map uses the canonical lax monoidal map
$iB\otimes iB'\to i(B\otimes B')$.

\begin{corollary}[Projection formulas and relative calculus]
\label{obs:cor:projection}
There are natural isomorphisms
\[
 o_!(B\otimes o^*F)\cong o_!B\otimes F,\qquad
 c_*(A\otimes c^*F)\cong c_*A\otimes F.
\]
The same formulas and exact sequences hold in the parameter
slices of \cref{glue:prop:slices}, with the corresponding
left exact functor $i_\alpha$.
\end{corollary}
\begin{proof}
\leavevmode
\begin{enumerate}[label=\textup{(\roman*)},leftmargin=*,itemsep=3pt,topsep=4pt]
\item \textbf{Projection formulas and relative calculus.} The two sides of the first formula are $(0,B\otimes sF,0)$;
the two sides of the second are $(A\otimes rF,0,0)$.
The identifications are natural componentwise.
In a slice,
\cref{glue:prop:slices} supplies exactly the same comma
description, and its $i_\alpha$ is left exact.
Repeating the
kernel, cokernel, and tensor calculations proves the formulas
there.
The canonical parameter comparisons satisfy their
coherence equations by their component projections.
\end{enumerate}
\end{proof}

This is an underived open--closed six-operation calculus with
abelian coefficients: on the open part $o^!=o^*$, on the closed
part $c_!=c_*$, and $c^!$ is the calculated kernel functor.
Its content is the adjunctions, exact sequences, tensor
operations, and projection formulas just proved. A derived
six-functor formalism for arbitrary maps of forcing topoi is a
different construction and is not used in these results.


The exact sequences identify the algebraic defect between closed raw data and open localized data. Finite cohomological examples will turn that general defect into a calculated obstruction to selecting names or extending an infinitesimal implementation.

\section{Cohomological naming obstructions and finite deformations}
\label{sec:naming-obstructions}


The presentation fibration supports a local-to-global obstruction
theory within the effective-domain framework. We first attach a
canonical abelian torsor to every name cover. We then construct
finite effective sites on which cohomology is computed by explicit
projective resolutions, and exhibit both a nontrivial naming torsor
and an obstructed infinitesimal extension. The coefficient objects
and their operations are constructed in the same forcing topoi.

\subsection{A canonical class attached to a name cover}
In a topos $\mathcal T$ with NNO, a torsor for a group object $G$
is an epimorphism $T\onto1$ with a $G$-action whose map
$G\times T\to T\times T$, $(g,t)\mapsto(t,gt)$, is an
isomorphism. In a slice, replace $1$ by the base object.
A section trivializes a torsor by translation; conversely, a
trivialization supplies a section. We write $H^1(\mathcal T,G)$
for its isomorphism classes, pointed by the trivial torsor.
For abelian $G$ the contracted product gives an abelian group.
These are the torsor classes
\cite[Tags 03AG and 03AJ]{stacks-project}. The definition and
the constructions below do not require an injective resolution.

\begin{proposition}[The augmentation obstruction to choosing names]
\label{obs:prop:augmentation}
For every presentation $q:C\onto X$ in $\mathcal T$, construct
in $\mathcal T/X$
\[
 \epsilon_q:\mathbb Z_X[C]\onto\mathbb Z_X,\qquad
 [c]\longmapsto1,\qquad
 K_q=\ker\epsilon_q,\qquad P_q=\epsilon_q^{-1}(1).
\]
Then $P_q$ is a $K_q$-torsor, giving a canonical class
\[
 o(q)=[P_q]\in H^1(\mathcal T/X,K_q).
\]
A section of $q$ implies $o(q)=0$. More precisely, $o(q)=0$
if and only if $\epsilon_q$ has an additive splitting, equivalently
if there is a global section of its augmentation-one fibre.
The construction is natural under geometric inverse image,
including parameter pullback.
\end{proposition}
\begin{proof}
\leavevmode
\begin{enumerate}[label=\textup{(\roman*)},leftmargin=*,itemsep=3pt,topsep=4pt]
\item \textbf{Coherence equations.} Construct $\mathbb Z_X[C]$ from finite signed words in $C$ by
the congruence imposing commutativity and cancellation; the
finite-list construction and exact quotient give the free
abelian object.
The sum of coefficients defines $\epsilon_q$.
After the epimorphism $C\onto X$, a chosen local name gives
the splitting $n\mapsto n[c]$, so $\epsilon_q$ is epic.
The difference of any two elements of $P_q$ belongs to $K_q$
and is their unique translation.
Moreover $c\mapsto[c]$ maps
$C$ into $P_q$, making the latter locally inhabited.
These
facts prove the torsor equation and its epic support.
\item \textbf{Torsor comparison.} A section of $q$ gives a section $[c]$ of $P_q$. A section of
$P_q$ trivializes the torsor, and conversely.
Since
$\mathbb Z_X$ is free on the terminal object of the slice,
such a section extends uniquely to an additive splitting of
$\epsilon_q$.
This proves the exact vanishing criterion.
A global section of this quotient object is a compatible
formal integral combination of names of total coefficient one;
the argument does not extract a single name from it.
\item Thus nonvanishing obstructs a selector, whereas vanishing
is not being identified with the existence of a selector.
\item \textbf{Adjunction and universal property.} For a geometric morphism $f:\mathcal F\to\mathcal T$, the
adjunction on abelian objects and the free universal property
give, naturally in $B$,
\[
 \Hom_{\operatorname{Ab}(\mathcal F)}
       (f^*\mathbb Z[C],B)
 \cong\Hom_{\mathcal T}(C,f_*UB)
 \cong\Hom_{\mathcal F}(f^*C,UB).
\]
Here $U$ forgets the abelian group structure; $f_*$ commutes
with it because it preserves finite limits.
Therefore
$f^*\mathbb Z[C]\cong\mathbb Z[f^*C]$, compatibly with
augmentation and its unit.
\item \textbf{Naturality and reindexing.} Left exactness and preservation of
epimorphisms identify the pulled-back kernel and
augmentation-one fibre with $K_{f^*q}$ and $P_{f^*q}$.
Their actions and torsor equations agree.
The same proof in
slices gives parameter naturality, with composition coherence
forced by the free universal property.

The defining maps satisfy the following commutative diagram
\[
\begin{tikzcd}[column sep=large,row sep=large]
C\times_XC\ar[r,shift left=.6ex,"\pi_1"]\ar[r,shift right=.6ex,"\pi_2"']&C\ar[r,"q"]&X
\end{tikzcd}
\]
\end{enumerate}
\end{proof}

For later calculations, this torsor definition agrees with
degree-one Yoneda cohomology. An extension
$0\to A\to E\to\mathbb Z\to0$ gives the $A$-torsor over $1$.
Conversely, for an $A$-torsor $T$, the augmentation kernel of
$\mathbb Z[T]$ maps to $A$ by
\[
 \sum_i n_i[t_i]\longmapsto
       \sum_i n_i(t_i-t_0),\qquad \sum_i n_i=0.
\]
Choose $t_0$ only locally. The coefficient equation makes the
result independent of that local choice, so it descends.
Push out the augmentation extension along this map.
Its fibre over $1$ is $T$, as is checked on a trivializing
cover. Conversely, that same local calculation recovers every
extension from its fibre. These natural inverse constructions
identify contracted product with addition of extensions.
This proves $H^1(\mathcal T,A)=
\operatorname{Ext}^1_{\operatorname{Ab}(\mathcal T)}
(\mathbb Z,A)$ without an enough-injectives premise.

\subsection{Finite sites with explicit cohomology resolutions}
\label{obs:subsec:finite-resolution}
The following uses only the projectivity of $1$ in the
\emph{ground} effective topos. It does not assert projectivity
of the terminal object in a forcing topos or in every slice.

\begin{lemma}[Ground points and principal charts]
\label{obs:lem:principal}
The object $1$ is projective in $\Eff$. For a finite decidable
poset $P$, put
\[
 \mathcal F_P=\Eff^P\simeq\Sh_\Eff(\Up(P)),
 \qquad U_p=\mathord\uparrow p.
\]
The convention is that $p\leq q$ gives a restriction
$A(p)\to A(q)$. Each $\mathbb Z[U_p]$ is projective in
$\operatorname{Ab}(\mathcal F_P)$, and the terminal object of
$\mathcal F_P/U_p$ is projective.
\end{lemma}
\begin{proof}
\leavevmode
\begin{enumerate}[label=\textup{(\roman*)},leftmargin=*,itemsep=3pt,topsep=4pt]
\item \textbf{Ground points and principal charts.} In the realizability presentation of an object $X$, the image
of $X\to1$ is the truth value
$\exists x\,E_X(x)$, where $E_X$ is its support predicate.
If the map is epic, a realizer of this truth yields one
supported representative $x$. A constant realizer for that
representative gives a point $1\to X$.
No family of choices
is involved.
Pulling an arbitrary epimorphism back along a
point proves projectivity of $1$.
\item \textbf{Tracked implementation.} This is the single-point
instance of the realizability description
\cite{vanoosten}; it is independent of the arbitrary-carrier
choice issue in \cref{rem:external-projectivity}.
\item \textbf{Adjunction and universal property.} Yoneda and the free abelian universal property give
\[
 \Hom_{\operatorname{Ab}(\mathcal F_P)}(\mathbb Z[U_p],A)
       \cong\Hom_\Eff(1,UA(p)).
\]
Evaluation is exact, and an epimorphism of internal abelian
groups is epic on its underlying objects.
The projectivity
just proved makes this Hom functor exact.
\item Finally, a diagram
on $\mathord\uparrow p$ has a global section exactly when
its value at the initial stage $p$ has a ground section;
the restriction maps supply all other components.
An
epimorphism onto the terminal diagram therefore splits,
proving the slice assertion.
\end{enumerate}
\end{proof}

\begin{proposition}[The face-site projective resolution]
\label{obs:prop:resolution}
Let $K$ be a finite simplicial complex with its vertices
decidably ordered, and let $P_K$ be its poset of nonempty
faces ordered by inclusion. The vertex opens $U_i$ cover
$P_K$, and their intersection $U_\sigma$ is the principal
open $\mathord\uparrow\sigma$ when $\sigma$ is a face and
is empty otherwise. The augmented complex
\[
 \cdots\longrightarrow
 \bigoplus_{\substack{\sigma\in K\\|\sigma|=3}}\mathbb Z[U_\sigma]
 \longrightarrow
 \bigoplus_{\substack{\sigma\in K\\|\sigma|=2}}\mathbb Z[U_\sigma]
 \longrightarrow
 \bigoplus_{\substack{\sigma\in K\\|\sigma|=1}}\mathbb Z[U_\sigma]
 \longrightarrow\mathbb Z\longrightarrow0
\]
with alternating face maps is a finite projective resolution.
For an abelian diagram $A$, applying $\Hom(-,A)$ computes
\[
 H^n(\mathcal F_{P_K},A)
 :=\operatorname{Ext}^n_{\operatorname{Ab}(\mathcal F_{P_K})}
                   (\mathbb Z,A).
\]
In particular, for a constant finite abelian group $G$ the
cochain group in degree $n$ is
$G^{\{\sigma\in K:|\sigma|=n+1\}}$ with the alternating
coboundary.
\end{proposition}
\begin{proof}
\leavevmode
\begin{enumerate}[label=\textup{(\roman*)},leftmargin=*,itemsep=3pt,topsep=4pt]
\item \textbf{The face-site projective resolution.} At a stage $\tau\in P_K$, the displayed chain complex has one
generator for each nonempty subset of the vertices of $\tau$.
It is the augmented simplex complex.
Inserting the least
vertex of $\tau$, with the alternating sign and with a
repeated vertex interpreted as zero, gives a contracting
homotopy.
Thus it is exact at each stage, hence exact as a
diagram.
This stagewise homotopy proves exactness; it need
not be a natural splitting of the whole resolution.
Each nonzero summand is projective by
\cref{obs:lem:principal}, and there are finitely many.
\item \textbf{Evaluation and transposition.} The projective-resolution construction of Yoneda Ext
therefore applies to this actual resolution.
For constant
$G$, evaluation at a principal stage gives precisely $G$,
and the face maps give the stated coboundary.
\end{enumerate}
\end{proof}

The same proof applies to the two-open cover of the four-stage
site below: its overlap is a disjoint union of two principal
opens, so its degree-one term is the sum of their free
abelian objects. All these sites are finite effective bases.
The existing construction $D=\Idl(P)$ and
$\Up(P)\cong\Sc(D)$ places their sheaf topoi over effective
information domains. We do not identify all generalized
ideals of a finite internal basis with principal ideals.

\subsection{A naming torsor separating operational oracles}
\begin{theorem}[Equal existence modality and strict answer reducibility]
\label{obs:thm:crown}
There is a four-stage effective information site carrying an
epic name presentation $T\onto1$ that is locally isomorphic
to the constant two-name presentation $G\onto1$, but has no
global name. Their existence modalities are both the identity,
while their container-morphism order satisfies
\[
 (G\to1)\ \leq\ (T\to1),
 \qquad (T\to1)\ \not\leq\ (G\to1).
\]
The torsor class of $T$ is nonzero in
$H^1(\mathcal F_P,\mathbb F_2)\cong\mathbb F_2$,
and its canonical augmentation obstruction is also nonzero.
\end{theorem}
\begin{proof}
\leavevmode
\begin{enumerate}[label=\textup{(\roman*)},leftmargin=*,itemsep=3pt,topsep=4pt]
\item \textbf{Unit and identity laws.} Take
\[
 P=\{a_0,a_1,b_0,b_1\},\qquad a_i<b_j\quad(0\leq i,j<2),
\]
with no other strict inequalities and with the identity
coverage.
Let $G$ be the constant additive group $\mathbb F_2$.
Set $T(p)=\mathbb F_2$ at every stage.
All restriction maps
are identities except $T(a_1\leq b_1)(x)=x+1$.
There are no nonidentity composable pairs, so this is a
diagram.
Translation by $G$ commutes with all restrictions,
and $G\times T\to T\times T$ is componentwise bijective.
The projection $T\to1$ is componentwise epic.
It is therefore
a torsor and, in particular, a name presentation.
\item \textbf{Torsor comparison.} A global section would give four ground bits with
\[
 x_{b_0}=x_{a_0}=x_{a_1},\qquad
 x_{b_1}=x_{a_0}=x_{a_1}+1,
\]
which is impossible.
The opens $U_0=\uparrow a_0$ and
$U_1=\uparrow a_1$ cover the site.
Choosing $0$ at $a_i$
trivializes the torsor there; on their overlap, the difference
of the two trivializations is $(0,1)$ at $(b_0,b_1)$.
The projective resolution of this cover gives the cochain
complex
\[
 G^2\xrightarrow{\delta}G^2,\qquad
 \delta(u,v)=(v-u,v-u).
\]
Its cokernel is $G$, and $(0,1)$ represents its nonzero
class.
\item \textbf{Homological calculation.} This is sheaf cohomology computed by the displayed
resolution, not an unverified identification of a nerve
calculation with cohomology.
\item \textbf{Unit and identity laws.} For the augmentation class, a section of $P_q$ would be a
compatible pair of integral coefficients at every stage.
The three identity restrictions identify the four coefficient
pairs, while the remaining restriction interchanges them.
Thus a common pair must be $(m,m)$, whose augmentation is
$2m$.
The integer equation $2m=1$ has no solution, so
$P_q$ has no section and $o(q)\ne0$.
\item \textbf{Image and quotient construction.} Both answer projections are epic, so
\cref{prop:answer-support-quotient} gives the identity modality.
The constant map $T\to G$ with value $0$ is an answer
transformation proving the first reduction. A transformation
for the reverse reduction would be a map $G\to T$; its
composite with $0:1\to G$ would be a forbidden global
section.
This proves strictness using the actual definition
of a container morphism.
\end{enumerate}
\end{proof}

The identity of the denotation object $1$ consequently has no
implementation from the identity presentation $1\to1$ to
$T\to1$ in this forcing topos. The canonical refinement is
nevertheless exactly $T\to1$, and on that refinement the
implementation is the identity of $T$, as prescribed by
\cref{prop:refined-implementation}. Thus the cohomological
obstruction locates a failure of unrefined global naming,
while preserving the constructed local implementation.
It does not contradict projectivity of the ground terminal
object, which has not been asserted in $\mathcal F_P$.

\subsection{An obstructed infinitesimal extension of finite names}
\label{obs:subsec:deformation}
We now construct ring objects and a finite naming system
directly on an effective site. This supplies infinitesimal
structure in the coefficient objects; the underlying locale
and its forcing interpretation remain the ones already
constructed from the information basis.

Let $K$ have vertices $1,\ldots,6$, all fifteen edges, and
the following ten maximal faces:
\begin{equation}\label{obs:eq:facets}
 123,\ 124,\ 135,\ 146,\ 156,\ 236,\ 245,\ 256,\ 345,\ 346.
\end{equation}
Use the face site $P_K$ of \cref{obs:prop:resolution};
it has $6+15+10=31$ stages. Orient edges and triangles by
increasing vertex order. Define the integer edge cochain
$z$ by
\[
 z_{25}=z_{26}=z_{34}=z_{36}=z_{45}=1,
 \qquad z_{ij}=0\quad\text{on the other ten edges}.
\]
For an oriented triangle $ijk$,
$(dz)_{ijk}=z_{jk}-z_{ik}+z_{ij}$. Direct substitution in
\eqref{obs:eq:facets} gives
\begin{equation}\label{obs:eq:dz}
 (dz)_{345}=2,\qquad (dz)_\sigma=0
           \quad\text{for every other triangle }\sigma.
\end{equation}
In particular $\bar z=z\bmod2$ is a cocycle. It is not a
coboundary: its sum on the edge cycle $1\,2\,5\,1$ is $1$,
whereas every vertex coboundary has sum zero on that cycle.

Construct the constant finite rings
\[
 R_2=\mathbb F_2[\epsilon]/(\epsilon^2),\qquad
 R_3=\mathbb F_2[\epsilon]/(\epsilon^3)
\]
as bit vectors of lengths two and three, with componentwise
addition and truncated polynomial multiplication. All
operations are finite computable tables, transported from
$\Eff$. Reduction $\rho:R_3\onto R_2$ has square-zero
kernel $I=\epsilon^2\mathbb F_2$. Their units fit into
\begin{equation}\label{obs:eq:units}
 0\longrightarrow\mathbb F_2
  \xrightarrow{\,a\mapsto1+\epsilon^2a\,}
 R_3^\times\longrightarrow R_2^\times\longrightarrow1.
\end{equation}
Here the first term is additive, and the middle terms are
multiplicative abelian groups. Indeed $R_2^\times$ has order
two, while $R_3^\times$ is cyclic of order four, generated
by $1+\epsilon$, whose square is $1+\epsilon^2$.
These are checked identities in the displayed finite rings.

\begin{theorem}[Finite names with a nonzero second obstruction]
\label{obs:thm:deformation}
On $\mathcal F_{P_K}$ there is a locally free rank-one
$R_2$-module $L_2$ with a specified reduction
$L_2/\epsilon L_2\cong\mathbb F_2$. Its reduction map
\[
 q:L_2\onto\mathbb F_2
\]
is a name presentation with locally four names and with
descended addition and scalar operations. There is no global
name for the semantic bit $1$. Moreover no locally free
rank-one $R_3$-module reduces to $L_2$. The obstruction is
the nonzero class
\[
 [b]\in H^2(\mathcal F_{P_K},\mathbb F_2),
 \qquad
 b_\sigma=\begin{cases}1&\sigma=345,\\0&\text{otherwise.}\end{cases}
\]
\end{theorem}
\begin{proof}
The obstruction is calculated from the explicit finite resolution. We identify the deformation cocycle and test its cohomology class before interpreting the failure of an infinitesimal lift in the glued naming system.

\leavevmode
\begin{enumerate}[label=\textup{(\roman*)},leftmargin=*,itemsep=3pt,topsep=4pt]
\item \textbf{The finite gluing module.} On the vertex open $U_i$ take a copy of the free module $R_2$.
For $i<j$ glue by multiplication with
$u_{ij}=(1+\epsilon)^{z_{ij}}$, using inverses for the
opposite order.
Since $(1+\epsilon)^2=1$ in $R_2$,
\eqref{obs:eq:dz} gives $u_{ij}u_{jk}=u_{ik}$ on every
triple overlap.
Effective descent
(\cref{obs:prop:descent}) constructs $L_2$ with these
trivializations.
All transitions reduce to $1$, giving the
specified reduction to the constant bit object.
\item \textbf{The cochain and its boundary.} Locally $q$ is the four-element to two-element map
$a+b\epsilon\mapsto a$.
Multiplication by $u_{ij}$ is
additive, $R_2$-linear, and preserves this map.
Hence its
addition, zero, and scalar operations descend, and
$q(x+y)=q(x)+q(y)$ and $q(rx)=\bar r\,q(x)$ hold globally.
These are actual total implemented operations on the
presentation, obtained from finite tables and descent.
A global name of $1$ would be a section of $L_2$ whose
coefficient in each local chart has constant term $1$.
Such a coefficient is a unit, so the section would be a
global frame.
\item \textbf{Nontriviality of the obstruction.} The transition class of a global frame is a
coboundary in $R_2^\times\cong\mathbb F_2$, contradicting
the cycle calculation for $\bar z$.
Thus the cover does
not split, even though it is locally the displayed finite
name cover.
\item \textbf{Persistence under base change.} Lift the transitions to
$\widetilde u_{ij}=(1+\epsilon)^{z_{ij}}$ in $R_3^\times$.
Their triangle defect is
\[
 \widetilde u_{ij}\widetilde u_{jk}
       \widetilde u_{ik}^{-1}
   =(1+\epsilon)^{(dz)_{ijk}}
   =1+\epsilon^2 b_{ijk},
 \qquad b=(dz)/2\bmod2.
\]
Equation \eqref{obs:eq:dz} gives the asserted support of $b$.
Every other lift of the same transitions has the form
$(1+\epsilon^2v_{ij})\widetilde u_{ij}$; its defect is
$b+\delta v$.
Thus a compatible lifting of transitions
exists exactly when $[b]=0$.
\item \textbf{Homological calculation.} To compute this class, use the actual projective resolution
of \cref{obs:prop:resolution}.
With $\mathbb F_2$
coefficients its cochain complex is
\[
 \mathbb F_2^6\xrightarrow{\delta^0}\mathbb F_2^{15}
                    \xrightarrow{\delta^1}\mathbb F_2^{10}.
\]
Every edge in \eqref{obs:eq:facets} occurs in exactly two
triangles.
Therefore the sum $S$ of all ten triangles is
a cycle over $\mathbb F_2$.
The pairing satisfies
$\langle b,S\rangle=1$, whereas
$\langle\delta v,S\rangle=\langle v,\partial S\rangle=0$.
It follows that $[b]\ne0$ in sheaf cohomology.
\item \textbf{Construction of the comparison.} More explicitly, the edge graph is connected, so
$\operatorname{rank}\delta^0=5$. A dependence among the
triangle boundaries assigns equal coefficients to the two
triangles incident with any edge.
Their adjacency graph
is connected: in the order
\[
 123,\ 124,\ 135,\ 236,\ 146,\ 245,\ 156,\ 345,\ 346,\ 256
\]
each triangle after the first shares an edge with an
earlier one.
Thus the sole dependence is the sum of all
ten boundaries and $\operatorname{rank}\delta^1=9$.
\item \textbf{Verification of the equations.} In particular $H^1\cong\mathbb F_2$ and
$H^2\cong\mathbb F_2$, and the displayed classes are their
nonzero elements.
\item It remains to rule out a lift that is not initially given
in these charts.
Suppose $L_3$ is a locally free rank-one
$R_3$-module with an isomorphism
$L_3\otimes_{R_3}R_2\cong L_2$.
On each principal open
$U_i$, its frame object is a torsor and hence an
epimorphism onto the terminal object of that slice.
By \cref{obs:lem:principal} it has a global frame there.
The reduction of this frame differs from the specified
frame of $L_2$ by an $R_2$-unit on $U_i$.
The map of
units in \eqref{obs:eq:units} is epic; projectivity in
the same slice lifts that unit.
\item \textbf{Functoriality and coherence.} Adjusting the frame by
this lift gives an $R_3$-frame reducing to the specified
one.
Its transition maps are therefore lifts of the
$u_{ij}$, and would make $b+\delta v=0$, a contradiction.
This proves nonexistence of every such $L_3$.
\end{enumerate}
\end{proof}

The reduction $L_2/\epsilon L_2\cong\mathbb F_2$ exhibits
$L_2$ as a first-order deformation of the trivial rank-one
bit module: it is locally free, and hence flat, over the
explicit dual-number ring $R_2$. The theorem proves that
this particular deformation does not extend to $R_3$.
The general geometric pattern is the classification of
line objects by units-torsors and the obstruction from
a square-zero units extension
\cite[Tags 040D and 0C6R]{stacks-project}. Here the name
cover, ring operations, transition maps, resolution, and
nonzero obstruction are constructed on the effective
information site itself. The explicit projective
resolution supplies the cohomology calculation in this
relative realizability setting.

\subsection{Relation to descent and new native objects}
\label{obs:subsec:scope}
The torsor obstruction concerns a locally inhabited fibre
with a specified simply transitive symmetry action.
\Cref{obs:prop:forms} constructs such fibres for locally
isomorphic presentations, and \cref{obs:prop:augmentation}
constructs an abelian torsor from every epic name cover.
For a fixed local value $t:\widetilde L\to X$, descending
$t$ through $e:\widetilde L\to L$ instead asks that this
particular value satisfy the descent equation. The later
\cref{ex:selected-descent} has a split $e$ but a $t$ that
does not descend. A section obstruction for $e$ therefore
cannot describe that example.

Likewise, the Cohen non-lifting theorem
\cref{prop:intensional-new-streams} proves that the fibre
of old names over the native generic sequence is initial.
It has no local names at all and hence is not a torsor.
The candidate locales of \cref{sec:candidates,sec:descent}
retain their localic image and pushout description; a
candidate projection is not assigned a group action or
the status of a sheaf epimorphism merely by terminology.
These distinctions specify which proved construction
applies to each existing computational interface.


\subsection{Transport and persistence of naming obstructions}
\label{obs:subsec:transport}
The finite-limit and image transport of \cref{lem:geometric-transport} applies to the displayed naming diagrams.
For a geometric morphism $f:\mathcal F\to\mathcal T$,
finite-limit and epimorphism preservation carry a $G$-torsor
to an $f^*G$-torsor. The free abelian comparison of
\cref{obs:prop:augmentation} identifies
\[
 f^*K_q\cong K_{f^*q},\qquad f^*P_q\cong P_{f^*q},
 \qquad f^*o(q)=o(f^*q).
\]
Here the last map is the map on torsor classes with its
displayed transported coefficient group. For forms of a
presentation, one transports the action of its actual
automorphism group; the canonical map
$f^*\operatorname{Aut}(C)\to\operatorname{Aut}(f^*C)$
is not presumed invertible. These claims use the
constructed group actions and quotients, without requiring
preservation of freshly formed exponential objects.

\begin{proposition}[Persistence with an independent Cohen oracle]
\label{obs:prop:cohen-product}
Let $L$ be either finite information locale constructed in
\cref{obs:thm:crown,obs:thm:deformation}, and let $C_{\rm Coh}$
be the structural Cohen locale of \cref{ex:cohen-logic}.
Along the projection
\[
 p:L\times C_{\rm Coh}\longrightarrow L
\]
the crown torsor and its augmentation torsor remain
nontrivial, and the finite module $p^*L_2$ still admits
no rank-one lift to $p^*R_3$. Thus the obstructions occur
also in an effective forcing setting containing an
independent native Cohen sequence.
\end{proposition}
\begin{proof}
\leavevmode
\begin{enumerate}[label=\textup{(\roman*)},leftmargin=*,itemsep=3pt,topsep=4pt]
\item \textbf{Persistence with an independent Cohen oracle.} The all-zero computable sequence gives a point
$z:1\to C_{\rm Coh}$: its prefix ideal contains exactly
the all-zero strings and meets the structural prefix
covers.
The frame-map construction of
\cref{thm:points-ideals} verifies this point.
Consequently
$s=(\id,z):L\to L\times C_{\rm Coh}$ is a section of $p$,
and the induced geometric inverse images satisfy
$s^*p^*\cong\id$.
\item \textbf{Torsor comparison.} A section of either pulled-back torsor would pull back
along $s$ to a section of the original torsor, contradicting
\cref{obs:thm:crown}.
Likewise, a locally free rank-one
$p^*R_3$-module lifting $p^*L_2$ would pull back to a
rank-one $R_3$-module lifting $L_2$: inverse image preserves
the local trivializations, their epimorphic covers, and
the tensor quotient defining reduction.
This contradicts
\cref{obs:thm:deformation}.
The structural Cohen factor
supplies its generic sequence by projection.
\item \textbf{Homological calculation.} The argument
uses the product and its specified retraction; it makes
no assertion that the Cohen factor alone has these
cohomology classes.
\end{enumerate}
\end{proof}


\part{Transport of computational presentations}\label{part:presentation-transport}

We have distinguished semantic objects from their names, imported objects from native ones, and full truths from positive observations. The remaining comparison problem is to transport these layers together. This part states the actual domain--value and code-transport maps, specifying where inverse image preserves interpretation and where only a comparison into a native object is available.


\section{Comparisons, localization, and relative replacements}\label{sec:comparison}

The preceding obstructions identify information that cannot always be recovered after localization. For maps that do transport, we now require an exact description of what is preserved. We formulate comparisons for actual name covers, positive classifiers, and partial domain--value spans, with separate attention to their imported and native targets.


\paragraph{External specialization.}
Let $x:\mathcal T\to\FF_j$ be a geometric morphism selecting a generalized forcing model. Applying $x^*$ to an actual partial span gives its externalized domain and value. If $\mathcal T=\mathbf{Set}$ this is an ordinary partial function. If its construction comes from a realizability evaluator, its computational interpretation is the one supplied by its domain-and-value adequacy equation \cite{bauer,hyland,vanoosten}. A map in a topos is not assigned an algorithm merely because it is a morphism; the positive presentation, effective point, and evaluator are the data that supply one.


A comparison of forcing packages specifies which data it transports. The most common realization comparison is a commuting square
\[
\begin{tikzcd}
 L\ar[r,"\phi"]\ar[d,"\rho"']&L'\ar[d,"\rho'"]\\
 X\ar[r,"\psi"']&X'.
\end{tikzcd}
\]
It gives $\rho^*\psi^*=\phi^*(\rho')^*$ on all full observations. If markings are also preserved, restriction gives the corresponding effective transformation on its pullback-defined domain. This single square is a useful organizing language, provided that the conditions, markings, and spans named in it remain part of the data.

\begin{proposition}[Comparison consequences]\label{prop:comparison-levels}
The preceding square induces comparison of the selected observations for which the inverse-image marking maps are supplied. If $\phi,\psi$ are isomorphisms, it induces, for every $x:T\to X$, an isomorphism
\[
 T\times_XL\cong T\times_{X'}L'.
\]
An isomorphism of local realization spans induces an isomorphism of their candidate images and descended observation frames. An isomorphism of marked information presentations induces all these comparisons for the structures built from the named maps.
\end{proposition}
\begin{proof}
\leavevmode
\begin{enumerate}[label=\textup{(\roman*)},leftmargin=*,itemsep=3pt,topsep=4pt]
\item \textbf{Comparison consequences.} The observation equation is obtained by applying inverse image to the square.
Restricting to markings is the universal property of \cref{thm:admissibility,thm:lax}.
For the fibre comparison, the map is $(\id_T,\phi)$, and its inverse is $(\id_T,\phi^{-1})$ after using $\psi^{-1}$ in the commuting square.
The pullback equations verify that both are defined.
\item \textbf{Pullback construction.} An isomorphism of spans identifies the paired maps into the corresponding products.
The localic image factorization is unique up to its unique commuting isomorphism, so the candidate images agree.
The pullback frame \eqref{eq:descent-frame} is preserved under these isomorphisms, as is its marked pullback.
\item Finally a marked presentation isomorphism transports its matrices, nuclei, and inverse-image maps by conjugation; \cref{prop:marked-isomorphism} identifies the effective pullbacks.
Each construction is natural because its defining projections determine it uniquely.
\end{enumerate}
\end{proof}

A raw map of condition preorders gives a forcing map only after its inverse-image transformation respects the cover equations. A forcing-locale isomorphism identifies the full forcing semantics, but does not by itself identify unrelated basic stages or raw information domains. The Cohen comparison in \cref{prop:cohen-comparison,prop:cohen-domains} exhibits this distinction concretely. A realization-compatible isomorphism additionally transports histories. Candidate spans need their own commuting comparison, since a choice procedure can carry information beyond its realized value.

\subsection{Constructing a forcing map from a condition map}
Let $F:P\to Q$ be monotone for the refinement orders. Inverse image of upper predicates gives a frame map $F^{-1}:\Up(Q)\to\Up(P)$. To descend it through the two forcing quotients, the exact equation is
\begin{equation}\label{eq:site-map-condition}
 q_PF^{-1}j_Q=q_PF^{-1}.
\end{equation}

\begin{proposition}[Site-induced forcing comparison]\label{prop:site-map}
Equation \eqref{eq:site-map-condition} constructs a unique frame map
\[
 F^\flat:\Omega_{J_Q}(Q)\longrightarrow\Omega_{J_P}(P),
 \qquad F^\flat q_Q=q_PF^{-1}.
\]
Its value on a stage is
\[
 F^\flat([q]_{J_Q})
 =q_P\{p\in P\mid q\leq F(p)\}.
\]
Identity condition maps induce identity forcing maps. Composable condition maps satisfying the quotient equations induce composable forcing maps with the expected composite.
\end{proposition}
\begin{proof}
\leavevmode
\begin{enumerate}[label=\textup{(\roman*)},leftmargin=*,itemsep=3pt,topsep=4pt]
\item \textbf{Site-induced forcing comparison.} Apply \cref{lem:quotient-factor} to $q_PF^{-1}$ and $j_Q$.
This gives the unique factorization and proves that it is a frame map.
Substitution of $\uparrow q$ gives the stage formula.
\item \textbf{Unit and identity laws.} For $G:Q\to R$, compute
\[
 F^\flat G^\flat q_R
 =F^\flat q_QG^{-1}
 =q_PF^{-1}G^{-1}
 =q_P(GF)^{-1}.
\]
The right side therefore factors through $q_R$ and satisfies the required equation.
Uniqueness identifies this factor with $(GF)^\flat$.
The identity calculation is the same, and associativity follows from associativity of the original inverse-image maps.

Naturality is expressed by the commutative diagram
\[
\begin{tikzcd}[column sep=large,row sep=large]
\Up(Q)\ar[r,"F^{-1}"]\ar[d,"q_Q"']&\Up(P)\ar[d,"q_P"]\\ \Omega_{J_Q}(Q)\ar[r,"F^\flat"']&\Omega_{J_P}(P)
\end{tikzcd}
\]
\end{enumerate}
\end{proof}

This theorem gives a concrete route to the comparison square: verify the cover equations after $q_PF^{-1}$. It does not require a strict map between arbitrary chosen objects of covering sieves. In examples one verifies the equations on generating covers and then uses the least-localization property to obtain \eqref{eq:site-map-condition}.

\begin{proposition}[Operational equations under localization]\label{prop:operational-localization}
A functional answer relation remains functional after localization, and a subrelation of it is functional. Exhaustiveness is preserved by sheafification of the whole answer map. A descended oracle map or descended answer map remains descended after a pullback for which the localized left leg is surjective.
\end{proposition}
\begin{proof}
\leavevmode
\begin{enumerate}[label=\textup{(\roman*)},leftmargin=*,itemsep=3pt,topsep=4pt]
\item \textbf{Operational equations under localization.} Functionality of $R\incl Q\times A$ is the inclusion
\[
 R(q,a)\wedge R(q,a')\ \vdash\ a=a'.
\]
This is a geometric sequent and is preserved by inverse image.
Replacing $R$ by a subrelation preserves the inclusion directly.
Exhaustiveness is epimorphy of $R\to Q$, and sheafification preserves regular epimorphisms.
For descent, pull back the equation $t=\rho e$ to obtain the corresponding equation for the restricted maps; uniqueness of the descended result follows from the surjective left leg.
These statements concern the operational relation and its maps.
Effectivity of the restricted truth is the separate pullback calculation of \cref{thm:lax}.
\end{enumerate}
\end{proof}

\subsection{Localization of retained structures}
Suppose $j\leq k$ are nuclei on $H$ and $i:L_k\incl L_j$. The restricted realization is $\rho_k=\rho_ji$. Its frame map is the relative quotient applied to $\rho_j^*$. The following formulas assemble the constructions already made into a coherent tower.

\begin{theorem}[Localization transport]\label{thm:localization-transport}
For $\rho_j:L_j\to X$, a generalized oracle $x:T\to X$, and a local realization span $(e,t)$ over $L_j$, the following hold.
\begin{enumerate}
\item Canonical stages and histories restrict by pullback:
\[
 V_p^k=V_p^j\times_{L_j}L_k,\qquad
 \Hist_{\rho_k}(x)\cong\Hist_{\rho_j}(x)\times_{L_j}L_k.
\]
A displayed history basis pulls back to a join basis and hence gives the canonical relative site of the new history locale.
\item A section survives precisely when it factors through $L_k$. A least-history section which factors remains least.
\item Pull back the span to $\widetilde L_k=\widetilde L\times_{L_j}L_k$. Whenever its left leg $e_k$ is surjective, its recomputed candidate image has a canonical inclusion
\[
 \Cand(e_k,t_k)\incl L_k\times_{L_j}\Cand(e,t).
\]
This is an isomorphism precisely when the pulled-back image map
$\widetilde L_k\to L_k\times_{L_j}\Cand(e,t)$ is surjective.
\item When $e_k$ is surjective, every descended oracle observation for $(e,t)$ descends for $(e_k,t_k)$, by restriction of its descended value. If the whole oracle descends, its restricted oracle is the descended map of the pulled-back span.
\end{enumerate}
All comparisons commute with further localization. Open surjective left legs remain surjective, by \cref{thm:open-descent}.
\end{theorem}
\begin{proof}
\leavevmode
\begin{enumerate}[label=\textup{(\roman*)},leftmargin=*,itemsep=3pt,topsep=4pt]
\item \textbf{Pullback construction.} For (1), a canonical stage is an open sublocale determined by a stage generator.
Its inverse image is the open determined by the pulled-back generator.
The history equation is associativity of pullbacks.
Inverse image of locale opens preserves arbitrary joins and finite meets, so a displayed join basis remains a join basis of the pulled-back sublocale, internally in the relative base.
Apply \cref{cor:canonical-basis} there.
\item \textbf{Verification of the equations.} For (2), the section equation after factorization follows by composing with the sublocale inclusion.
Leastness is the inequality $\kappa\rho_j\leq\id$ of \cref{thm:least-history}.
Its restriction to $L_k$ gives $\kappa_k\rho_k\leq\id_{L_k}$: postcomposition with a locale embedding reflects the order, since its inverse image is surjective on frames.
Thus the factored section remains least.
\item \textbf{Pullback construction.} Write $C=\Cand(e,t)$ and factor the paired map as $\widetilde L\to C\incl L_j\times X$.
Its pullback through $L_k$ gives
\[
 \widetilde L_k\longrightarrow L_k\times_{L_j}C\incl L_k\times X.
\]
The image of the first arrow is precisely the new candidate image, because composing an image factorization with a sublocale inclusion preserves its image.
This gives the inclusion and the isomorphism criterion in (3).
\item \textbf{Coherence equations.} For (4), if $e^*V=t^*U$, pull back the equation to obtain
$e_k^*i^*V=t_k^*U$.
This supplies the descended value on $L_k$, and its uniqueness follows from surjectivity of $e_k$.
In particular $t=\sigma e$ pulls back to $t_k=(\sigma i)e_k$.
Coherence in all four clauses follows from the commuting projections of iterated pullbacks and the uniqueness of image and descent factorizations.
\end{enumerate}
\end{proof}

The first clause concerns the locale's relative basis, not the literal preservation of every old cover-object. The new semantic coverage is determined by its pulled-back basic opens. Likewise, (3) recomputes the image before comparing it; it does not presuppose stability of every localic surjection under every pullback.

For a displayed region formula $A_U(p)=\Must_p(U)$ from \cref{thm:regions}, the new realized truth is $k(A_U)$. Its raw formula is exact exactly when $A_U$ is $k$-fixed. Cover exhaustivity for the enlarged coverage proves this fixedness by the same argument as before. A stage section proves exactness at the new layer only when its map factors through the localized stage. These are equations on the existing construction, not independent new oracle axioms.

\subsection{Relative presentation and observational comparison}
After a change of base, a construction may use newly available oracle objects, predicates, or presentations. Its locale $\widehat L$ is then compared with the literal classifier base change $f^\#L$ by a specified map. To call it exact at the forcing level is to prove that map invertible. To call it exact for a realization, information domain, or effective observation family is to prove the corresponding commuting comparison and its claimed invertibility.

For selected observations, the necessary comparison is often smaller. Suppose there are maps
\[
 a:\widehat L\to Z,\qquad b:f^\#L\to Z
\]
or another common comparison context, and selected oracle families whose interpreted truths are to agree there. Form the equalizer of the two interpreted family maps into the common full truth object. Intersect it, by pullback, with their two admissibility objects. This constructs the maximal parameter subobject on which the comparison is both true and effective. It is a calculation with the chosen maps; it does not assert equality for unselected observations.

\begin{proposition}[Coherence of selected comparisons]\label{prop:selected-comparison}
Equalizers and admissibility pullbacks define selected observational comparisons stably under parameter substitution. Their composites have the domain requiring each intermediate comparison. The direct comparison has the larger domain on which only the final equation and marking conditions hold. The canonical inclusions satisfy the same coherence laws as \cref{thm:lax}.
\end{proposition}
\begin{proof}
\leavevmode
\begin{enumerate}[label=\textup{(\roman*)},leftmargin=*,itemsep=3pt,topsep=4pt]
\item \textbf{Kernel-pair descent.} An equalizer is the subobject where the two full maps agree.
Intersecting with pullbacks of the relevant markings expresses all the conditions simultaneously.
Pullback along a parameter map commutes with these finite limits.
For composition, conjunction of the two equalities implies the final equality by transitivity; conjunction of the intermediate marking conditions implies the final ones.
This gives the domain inclusion.
Two ways of composing these inclusions are the same mono factorization, proving coherence.
\end{enumerate}
\end{proof}

Finally, classifier base change preserves the geometric presentation of the constructions. It preserves finite-limit constructions of locales through the corresponding pullback of localic toposes. For a pushout or an image, the elementary universal property at least gives a canonical comparison from the newly constructed object to the pulled-back old one whenever the transported cocone or factorization is given. Such a comparison is used at exactly its proved strength. Isomorphisms of all named data, as in \cref{prop:comparison-levels}, do preserve images and pushouts; an arbitrary change of ambient topos is a separate preservation question.


\section{Intensional enhancement: coding fibrations, oracle realizers,
and certification}\label{sec:intensional-basechange}
We now connect the semantic base-change theorems with the genuine
name-presentation structures of \cref{sec:pres} and the finite operational certificates constructed below.
The comparisons concern specific covers, implementations, domains
and descriptions. The comparisons retain the denotational kernel pairs of the chosen name covers.

\subsection{Transport of the coding fibration}
For the ground topos $\E$, let $\Pres(\E)$ denote the category of regular
epimorphisms $q:C\onto X$ and their commuting implementation
squares, as in \cref{sec:pres}. Its projection to $\E$ is
a fibration, with cartesian arrows given by pullback.

\begin{theorem}[Geometric transport of name presentations and adequacy]
\label{thm:intensional-pres-fibration}
For every geometric morphism $f:\FF\to\E$, inverse image defines
a functor
\[
 f^{\Pres}:\Pres(\E)\longrightarrow\Pres(\FF),
 \qquad (C\xrightarrow{q}X)\longmapsto
 (f^*C\xrightarrow{f^*q}f^*X).
\]
It carries cartesian name reindexing to cartesian name
reindexing and satisfies the canonical composition coherence.
For every name cover $q$, it identifies the denotational kernel
pair by
\[
 f^*(C\times_XC)\cong f^*C\times_{f^*X}f^*C.
\]
Implementation equations, including equalities of the extent and
value maps of $\Sigma$-partial implemented operations, remain true
after transport. Their new extents are classified through the
imported truth object $\Theta_f$ of
\cref{def:imported-truth}; no preservation of native Rosolini
truth or native completion objects is asserted.
\end{theorem}
\begin{proof}
\leavevmode
\begin{enumerate}[label=\textup{(\roman*)},leftmargin=*,itemsep=3pt,topsep=4pt]
\item \textbf{Pullback construction.} A geometric inverse-image functor is left exact and preserves
regular epimorphisms.
Thus it sends presentation objects and
commuting squares to presentation objects and commuting squares.
The cartesian arrow over $h:Y\to X$ is the regular-epimorphic
pullback $Y\times_XC\to Y$.
Applying $f^*$ yields the pullback
of $f^*q$ along $f^*h$; hence the image arrow is cartesian.
The canonical comparison of iterated pullbacks preserves both
projections and is therefore associative by uniqueness.
\item \textbf{Inverse comparison.} Applying
left exactness to the pullback square defining $C\times_XC$
gives the displayed kernel-pair isomorphism.
\item \textbf{Pullback construction.} An implementation is a commuting square $q'\widehat u=uq$;
applying $f^*$ gives the same equation after transport. A partial
implementation is a domain--value span satisfying equality of its
domain mono and value with the semantic partial operation.
Inverse
image preserves its domain pullback, the value map and their
commuting triangles, by \cref{thm:partial-transport}.
If the
original classifier is $p:C\to\Sigma_\E$, the imported classifier
is $e_ff^*p:f^*C\to\Theta_f$ and its underlying subobject is the
pulled-back domain by \cref{eq:truth-comparison}.
Thus no
additional choice or code-extraction assumption enters the
transport theorem.
\end{enumerate}
\end{proof}

\subsection{The exact old-code-representable fragment}
For a searchable presentation $C$, let $N_C\onto\Pt D_C$
be a name cover constructed in \cref{thm:canonicalcoding}, and denote the effective
point inclusion by $j_C:\Pt D_C\incl D_C$. The comparison of
\cref{prop:base-approximation} gives
\[
 f^*N_C\longrightarrow f^*\Pt D_C\xrightarrow{f^*j_C}
 f^*D_C\xrightarrow{\beta_{f,C}}D_{f^\dagger C}.
\]

\begin{theorem}[Universal old-name representability image]
\label{thm:intensional-old-image}
Factor this map into its regular image:
\[
 f^*N_C\onto\operatorname{Rep}_{f,C}
 \incl D_{f^\dagger C}.
\]
Then
\[
 \operatorname{Rep}_{f,C}
 =\im\bigl(\beta_{f,C}f^*j_C:
    f^*\Pt D_C\to D_{f^\dagger C}\bigr).
\]
It is the smallest native subobject containing every transported
old effective denotation, and the transported names cover it.
The image construction is natural for explicitly implemented
maps and stable under substitution of object parameters. It need
not exhaust the native completion, nor does image formation by
itself establish a native $\Sigma$-classifier for the image.
\end{theorem}
\begin{proof}
\leavevmode
\begin{enumerate}[label=\textup{(\roman*)},leftmargin=*,itemsep=3pt,topsep=4pt]
\item \textbf{Universal old-name representability image.} The first arrow $f^*N_C\to f^*\Pt D_C$ is regular epic.
Precomposing a morphism with a regular epimorphism does not change
its regular image: factor the second morphism through its image,
then compose the first regular-epimorphic leg, and use uniqueness
of image factorizations.
This proves equality of the two images
and the regular-epimorphic covering by names.
\item \textbf{Pullback construction.} Every subobject
through which the composite factors contains its regular image,
which proves minimality. A commuting implemented square gives
commuting maps of the two composites; images inherit a comparison
by the universal factorization.
Pullback stability of regular
images in a topos gives parameter naturality.
The possible
failure of $\beta_{f,C}$ to be epic was exhibited in
\cref{prop:beta-not-epi}, so no unproved exhaustion follows.
\end{enumerate}
\end{proof}

\subsection{Realizer equations for oracle modalities}
\label{sec:intensional-realisers}
This calculation is made in the ground topos $\Eff$, using $K_1$ application and numerical realizers. Choose a partitioned-projective query presentation
whose carrier is $|Q|$ and whose naming map into the PCA is
$\varepsilon(q)$. Let the answer predicate have realizer fibre
$P(q)\subseteq\N$. Represent an arbitrary truth value by a
subset $S\subseteq\N$.

\begin{theorem}[Concrete tagged oracle computation realization]
\label{thm:intensional-realisers}
Define the monotone operation on realizer sets
\[
 \begin{split}
 \Phi_S(Z)={}&\{\langle0,s\rangle:s\in S\}\\
 &\cup\bigl\{\langle1,\langle\varepsilon(q),c\rangle\rangle:
 q\in|Q|,\ \forall a\in P(q),\ ca\downarrow\ \land\ ca\in Z\bigr\}.
 \end{split}
\]
Its least fixed point
$\mathcal O_P(S)=\bigcap\{Z\subseteq\N:\Phi_S(Z)\subseteq Z\}$
represents the Ahman--Bauer modality $o_P(S)$, up to the
realizability preorder, with uniform realizers for inflationarity,
finite meets, multiplication and the query rule. The construction
retains the actual codes of successive queries and answers; the
underlying propositional modality depends only on the support.
\end{theorem}
\begin{proof}
\leavevmode
\begin{enumerate}[label=\textup{(\roman*)},leftmargin=*,itemsep=3pt,topsep=4pt]
\item \textbf{Partiality and definedness.} The leaf code is a direct realizer of $S$. A query code is a name
$\varepsilon(q)$ and a partial-combinatory program $c$ realizing
$P(q)\Rightarrow Z$: it is defined and lands in $Z$ on every
valid answer realizer.
Thus $\Phi_S$ is the realizer interpretation
of $Z\mapsto S\vee\exists q\,(P(q)\Rightarrow Z)$.
Monotonicity gives the least-prefixed-point construction and its
induction principle.
The uniform return realizer sends $s$ to
$\langle0,s\rangle$.
\item \textbf{Oracle interpretation.} Given a tracked implication $S\to T$,
recursion on the tagged computation leaves query nodes unchanged
and replaces a leaf realizer by its translated result; this gives
monotonicity.
\item \textbf{Multiplication law.} For multiplication, replace every leaf whose payload
is a tagged computation by that computation, and compile every query continuation by recursion.
For finite meets, execute the
first tree; at each return execute the second tree and return the
pair of terminal realizers.
Projections of terminal pairs give
both converse maps.
These are finite constructor programs defined
by the recursion and delayed-continuation operations of \cref{rel:lem:abstraction,rel:thm:recursion}:
they do not inspect unavailable oracle answers.
\item \textbf{Partiality and definedness.} Their defining extent equations state $ca\downarrow$ only on
supplied answers.
The tracking of the resulting tagged transformations
in $K_1$ is the effective-topos realizability
calculation of \cite[Section~5]{ahmanbauer}, using its partitioned
query presentation and realizability implication.
It concerns
codes for well-founded, possibly unbounded, answer-dependent
interactions: it does not assert a uniform bound on the tree's
depth or branching, and it does not derive a tracker for an
unrelated object-level oracle map in the ground computational
signature.
\item \textbf{Tracked implementation.} The implemented evaluator of \cref{thm:canonicalcoding} supplies such a tracker when
its own implementation equations have been verified.
\item \textbf{Forcing construction.} The query code $\langle1,\langle\varepsilon(q),[a]\langle0,a\rangle\rangle\rangle$
realizes the imposed truth $P(q)$ after oracle closure: each valid answer is returned as a tagged leaf.
The leaf tag is essential to membership in the displayed closure.
If $k$ is any local operator
forcing all $P(q)$, interpret a tagged tree in $k(S)$ by structural
induction: at a leaf apply $S\to k(S)$; at a query use the
permitted implication realizer for $P(q)\Rightarrow k(S)$,
apply $k$, and use $kP(q)=\top$ and idempotence.
\item \textbf{Oracle interpretation.} The delayed
interpreter has the required realizability tracking;
its equation is checked on each valid answer before taking the
least closure.
Hence $\mathcal O_P(S)\leq k(S)$.
This is the
concrete realization of the least-fixed-point and leastness
calculations of \cref{lem:oracle-inductive-closure}, and matches
the tagged formula of \cite[Section~5]{ahmanbauer}.
\end{enumerate}
\end{proof}

\begin{corollary}[Extended-Weihrauch presentations and all local operators]
\label{cor:intensional-weihrauch}
For the query carrier above define the extended-Weihrauch predicate
\[
 \mathsf w_P(b)=\{P(q)\subseteq\N:\varepsilon(q)=b\}.
\]
Then the query clause in \cref{thm:intensional-realisers} equals
\[
 \bigl\{\langle1,b,c\rangle:\exists\theta\in\mathsf w_P(b)\;
   \forall a\in\theta\,[ca\downarrow\land ca\in Z]\bigr\}
\]
up to computably invertible pairing of codes. Conversely, for
$\mathsf w:\N\to\Pow(\Pow(\N))$, take
$Q_\mathsf w=\{(b,\theta):\theta\in\mathsf w(b)\}$ with name
$\varepsilon(b,\theta)=b$ and answer fibre $P(b,\theta)=\theta$.
The two least-fixed-point formulas are then equal by their indexed
unions. Every local operator in $\Eff$
has such a presentation, by the canonical query predicate of
\cref{thm:canonical-modal-oracle} and a partitioned-projective
cover, whose invariance is \cref{prop:query-cover-invariance}.
\end{corollary}
\begin{proof}
\leavevmode
\begin{enumerate}[label=\textup{(\roman*)},leftmargin=*,itemsep=3pt,topsep=4pt]
\item \textbf{Oracle interpretation.} Expand the quantifier $q\in|Q|$ in the query clause by its name
$b=\varepsilon(q)$ and the corresponding answer-realizer set
$\theta=P(q)$.
This gives the displayed union.
Conversely the
explicit carrier $Q_\mathsf w$ makes the same union its query
clause, so the resulting monotone operators and least fixed points
coincide.
For a general local operator $k$, use
$Q_k=\sum_{p:\Omega}k(p)$ and $P_k(p,h)=p$ from
\cref{thm:canonical-modal-oracle}.
Pull this query predicate to a
partitioned-projective cover.
\item \textbf{Verification of the equations.} Regular-epimorphic presentation
invariance preserves $o_{P_k}=k$, and the calculation just made
provides its extended-Weihrauch predicate.
No numerical
enumeration of arbitrary query objects and no selection of an
answer for every query has been used.
\end{enumerate}
\end{proof}

\subsection{Transport of finite execution evidence}
The accepted-run family is the one constructed in \cref{thm:operational-certificates}; its native oracle counterpart is \cref{prop:intensional-native-checker}. We transport the displayed checker, witness covers and value maps.

\begin{theorem}[Geometric transport of finite execution evidence]\label{thm:intensional-certificate-transport}
For a forcing geometric morphism $f:\FF\to\Eff$, finite numerical syntax, its constructor maps, and its accepted-run family have canonical comparisons
\[
 f^*\mathsf{Run}\cong\mathsf{Run}_{\FF},\qquad
 f^*\operatorname{Cert}_V\cong\operatorname{Cert}_{f^*V}.
\]
They preserve the finite run and wrapper equations and satisfy
\[
 f^*\bigl(\exists t\,[V(e,x,y,t)=1]\bigr)
 =\exists t\,[f^*V(e,x,y,t)=1]
\]
as subobjects, with the parameters transported by $f^*$. The transported search supports are $\Theta_f$-classified. The same conclusion holds for the actual finite-hit witness and value covers of presented operations, together with their kernel-pair equations.
\end{theorem}
\begin{proof}
\leavevmode
\begin{enumerate}[label=\textup{(\roman*)},leftmargin=*,itemsep=3pt,topsep=4pt]
\item \textbf{Kernel-pair descent.} Inverse image preserves the NNO, finite lists, finite products and sums, and the equalizers defining decidable grammar relations.
The primitive recursive construction of $V$ therefore transports with its defining equations.
Its accepted family is the equalizer of $V$ and the constant $1$, and existential support is a regular image.
Left exactness and preservation of regular images prove the two displayed comparisons.
The classifier of each transported positive support factors through $\Theta_f$ by \cref{prop:theta}.
\item \textbf{Partiality and definedness.} Apply the same finite-limit and image arguments to the displayed witness and value covers of \cref{thm:partialadequacy}; their denotational kernel pairs transport by \cref{thm:intensional-pres-fibration}.
The wrapper transformations are actual maps with verified checker equations, so applying $f^*$ preserves those equations and their composition.
\end{enumerate}
\end{proof}


The native checker and the transported ground checker have their own accepted families. A forcing assertion about a computation is related to finite execution evidence through the actual interpreter and its extent equation. Forcing a cover supplies local answers, whereas an execution certificate retains the particular finite answers and transitions used by that execution.


\part{External points, stabilization, and assembly}\label{part:external}

The categorical constructions acquire their forcing meaning when conditions are assembled into generic points and their finite information stabilizes. We use continuous ideals, requirement predicates, and fusion to make that passage explicit. The concluding synthesis then records which comparisons preserve full forcing truth, positive computation, and retained implementation witnesses.


\section{Continuous ideals and requirement constructions}\label{sec:external-points}


A forcing universe is determined by its presented locale. An external generic construction selects a point of that locale after the chosen change of base. The point is described by a continuous ideal, so its proof obligations can be read directly from the cover presentation.


Let $p_0\leq p_1\leq\cdots$ be a refinement chain and put
\[
 I_p=\{q\mid\exists n,\ q\leq p_n\}.
\]
It is inhabited at $p_0$ and lower. For witnesses $q\leq p_m$ and $r\leq p_n$, the term $p_{\max(m,n)}$ witnesses directedness.

\begin{proposition}[Cover-continuous chains]\label{prop:chain-point}
If for every $n$ and every $S\in J(p_n)$ there are $m\geq n$ and $s\in S$ with $s\leq p_m$, then $I_p$ is continuous and determines a point of $L_{P,J}$.
\end{proposition}
\begin{proof}
\leavevmode
\begin{enumerate}[label=\textup{(\roman*)},leftmargin=*,itemsep=3pt,topsep=4pt]
\item \textbf{Cover-continuous chains.} Let $q\in I_p$ and let $S$ cover $q$.
Choose a witness $q\leq p_n$.
Pullback stability gives the cover $S\cap\uparrow p_n$ at $p_n$.
The chain property gives $s$ in that sieve with $s\leq p_m$.
Thus $s\in S\cap I_p$.
Apply \cref{thm:points-ideals}.
\end{enumerate}
\end{proof}

Suppose $J_{\mathrm{Req}}$ is generated over a structural topology $J_0$ by upper requirement predicates $D_i$, with covers $D_i\cap\uparrow p$ at each $p$.

\begin{theorem}[Requirement-generic chains]\label{thm:requirement-point}
If the chain is $J_0$-continuous and for every $i,n$ there is $m\geq n$ with $D_i(p_m)$, then $I_p$ is $J_{\mathrm{Req}}$-continuous.
\end{theorem}
\begin{proof}
\leavevmode
\begin{enumerate}[label=\textup{(\roman*)},leftmargin=*,itemsep=3pt,topsep=4pt]
\item \textbf{Requirement-generic chains.} It meets structural covers by \cref{prop:chain-point}.
If $q\in I_p$, choose $n$ with $q\leq p_n$.
The stated requirement property gives $m\geq n$ with $D_i(p_m)$.
Then $p_m$ belongs to the requirement cover at $q$ and to $I_p$.
Thus it meets all generators.
Apply \cref{lem:ideal-topology}.
\end{enumerate}
\end{proof}

The localization and its points are separate constructions. The cover equations construct a locale whether or not a point-selection procedure is available. The last theorem checks a supplied chain. An actual chain constructor is obtained, for example, from a successor operation with a fair schedule.

\begin{theorem}[Canonical successor from a decidable certificate search]
\label{thm:decidable-certificate-successor}
Let $A=P\times I$ be a condition--requirement parameter object,
let $c:A\times\N\to2$ be a total decidable certificate verifier,
and define the supported search domain
\[
 D_c=\{u:A\mid\exists n:\N\;c(u,n)=1\}.
\]
There is a unique total least-certificate map
$\mu_c:D_c\to\N$ satisfying
\[
 c(u,\mu_c(u))=1,\qquad
 \forall m<\mu_c(u)\;c(u,m)=0.
\]
If an accepted certificate code $n$ has a constructed extension
$\operatorname{ext}(u,n)\in P$ with $p\leq\operatorname{ext}(u,n)$
and $D_i(\operatorname{ext}(u,n))$, then
\[
 s(u)=\operatorname{ext}(u,\mu_c(u))
\]
is an actual successor map on $D_c$. In particular it is total on
$P\times I$ whenever the displayed certificate-existence lemma
proves $D_c=A$. No axiom of choice or selector for arbitrary
regular epimorphisms is involved.
\end{theorem}
\begin{proof}
\leavevmode
\begin{enumerate}[label=\textup{(\roman*)},leftmargin=*,itemsep=3pt,topsep=4pt]
\item \textbf{Canonical successor from a decidable certificate search.} Over $D_c$ consider the relation
\[
 R(u,n)\iff c(u,n)=1\ \wedge\
          \forall m<n\,[c(u,m)=0].
\]
Given a witness $k$ of $c(u,k)=1$, finite decidable search through
$0,\ldots,k$ constructs a first successful index.
This proves
that $R\to D_c$ is a regular epimorphism, using only finite
recursion and existential elimination.
Two first successful
indices cannot differ: numerical comparison is decidable, and
the larger one's minimality clause contradicts success at the
smaller.
\item Thus $R\to D_c$ is also monic. A monic regular
epimorphism in a topos is an isomorphism; its inverse followed by
the code projection constructs $\mu_c$.
The two displayed
verification equations follow from its definition and are stable
under all parameter substitutions.
Substituting this least code
into the given extension map proves both successor equations.
If $D_c=A$, the domain mono is invertible, so the successor is
total on $A$.
\item \textbf{Observation preservation.} For numerically represented effective conditions
with a tracked verifier and extension interpreter, the same
least-number search is tracked on its proved domain, by the
computational interfaces of \cref{sec:numericalcodes,thm:partialadequacy}.
This last assertion does not
turn an arbitrary positive predicate into a decidable verifier:
positive leaf evidence must be included in the finite code when
needed to make the certificate check decidable.
\end{enumerate}
\end{proof}

\begin{proposition}[A scheduled construction by recursion]\label{prop:scheduled-chain}
Suppose $D_n$ are countably many upper requirements, and a map $s:P\times\N\to P$ satisfies
$p\leq s(p,n)$ and $D_n(s(p,n))$. Let $\nu:\N\to\N$ visit each number at arbitrarily late times. Recursion
\[
 p_{n+1}=s(p_n,\nu(n))
\]
constructs a refinement chain eventually in each $D_i$. If the structural generators have a displayed numerical indexing and the schedule also processes every such indexed obligation by its supplied refinement constructor, the resulting ideal is a point of the requirement localization.
\end{proposition}
\begin{proof}
\leavevmode
\begin{enumerate}[label=\textup{(\roman*)},leftmargin=*,itemsep=3pt,topsep=4pt]
\item \textbf{A scheduled construction by recursion.} The refinement equation follows by induction.
For any $i,n$, choose $k\geq n$ with $\nu(k)=i$.
Then $D_i(p_{k+1})$ by the equation for $s$, and upperness retains it at every later stage.
Hence the chain meets all requirement generators.
\item \textbf{Verification of the equations.} For the displayed numerically indexed structural generators, retain each obligation as a pair $(n,i)$ consisting of the already constructed stage $p_n$ and an index $i$ for one of its displayed generating covers. A fair enumeration of finite pairs schedules every such obligation after its base stage has been constructed; for example, dovetail the finite blocks of pairs whose two coordinates are at most the current stage.
When $(n,i)$ is processed at $p_k$ with $k\geq n$, pull the specified cover at $p_n$ back along $p_n\leq p_k$.
\item \textbf{Localization and coverage.} The structural successor operation returns a refinement $p_{k+1}$ lying in a member of this pulled-back sieve, hence in the original sieve.
Subsequent refinements retain that witness by upperness.
Thus for every $n$ and every displayed structural generator at $p_n$, a later chain stage lies above a member of the covering sieve.
This is exactly the condition of \cref{prop:chain-point} on the generating coverage.
\Cref{lem:ideal-topology} passes from the generators to their topology, and \cref{thm:requirement-point} gives continuity for the joint requirements.
\item \textbf{Functoriality and coherence.} This is a construction with the already supplied successor operations, not a selection from bare density.
\end{enumerate}
\end{proof}

The maps in this proposition are constructors, not choices inferred from bare existential density. When the displayed extension theorem includes a decidable finite-certificate verifier, \cref{thm:decidable-certificate-successor} constructs the successor on its exact certificate-existence domain. The front constructor of \cref{thm:uniform-front-recursion} and finite bushy certificate searches with their positive leaf witnesses recorded are instances; a general semidecision without a decidable accepted-certificate checker is not silently treated as such an instance. In an effective instance, a successor operation can often be built by dovetailing a positive search over coded extensions and returning its first successful witness. The proof that this search is total is then exactly the supplied extension lemma. This explains where the combinatorial proof enters an effective generic construction.

\subsection{Density and a basis extension lemma}
A stable truth $D\in\Omega_J(P)$ is dense when $\neg D=0$, with negation in the forcing frame. Its basis characterization is often a useful intermediate calculation.

\begin{proposition}[Density on basic stages]\label{prop:basis-density}
A stable truth $D$ is dense precisely when, for every condition $p$,
\[
 [p]_J\wedge D=0\quad\Longrightarrow\quad[p]_J=0.
\]
\end{proposition}
\begin{proof}
\leavevmode
\begin{enumerate}[label=\textup{(\roman*)},leftmargin=*,itemsep=3pt,topsep=4pt]
\item \textbf{Density on basic stages.} If $\neg D=0$, a stage with zero meet with $D$ lies below $\neg D$, hence is zero.
Conversely expand $\neg D$ as the join of the basic stages below it.
Each such stage has zero meet with $D$, and the stated implication makes each zero.
Their join is zero, proving density.
\end{enumerate}
\end{proof}

The conclusion is an equality test in a frame. A successor algorithm requires a witness of a suitable strengthening. These are related by an example-specific basis extension construction. For example, in the ordinary binary-string site, if an open $U$ is represented by a positive enumeration of cylinders and a cylinder $[\sigma]$ meets it positively, a witness names a compatible enumerated cylinder. The longer of its string and $\sigma$ is a strengthening with cylinder contained in $U$. The witness can be found by search over the supplied enumeration. This derives the extension operation for that presentation; frame density alone has not selected an enumerated cylinder.

When consistency is to be maintained, run the construction by recursion on the state object
$\{p\mid[p]_J\neq0\}$ with a constructor returning both the next condition and its consistency proof. The recursion then retains the proof component automatically. The later continuity argument uses the resulting underlying chain exactly as in \cref{prop:chain-point}.

\subsection{The external truth equation}
For a point with ideal $I$, a stable open $U$ satisfies
\[
 x^*(\widetilde U)
 =\exists p\,[I(p)\wedge U(p)]
 =\exists p\,[I(p)\wedge[p]_J\leq U].
\]
The first equality is the inverse-image formula of \cref{thm:points-ideals}, and the second is stable stage membership. Thus every geometric assertion true of the externalized generic has a condition in the selected ideal forcing its open interpretation, locally in the external parameter context. The assertion also applies to any fixed forced subobject interpreted before externalization. Identifying that pulled-back subobject with a fresh non-geometric interpretation requires the corresponding logical preservation, as explained next.

\subsection{What a point preserves}
Pulling a forced subobject back along a point gives the interpreted subobject in the selected model. This preserves geometric connectives and the domains and values of the already constructed partial maps. For a non-geometric assertion such as implication or universal quantification, the comparison is the actual inverse image of its forcing interpretation. A geometric inverse-image functor need not preserve that logical operation computed afresh in the target. The full internal forcing theorem of \cref{thm:kj} is therefore used in the forcing topos; external specialization is stated through its subobjects and the geometric comparisons just proved.


A refinement construction yields a generic point when its continuous ideal satisfies the presented covers. This provides the semantic target against which computational forcing predicates and their extension-density saturation can be compared.

\section{Computational forcing predicates and semantic saturation}\label{sec:forcing-saturation}


\subsection{Why saturation is a necessary construction}
There are computability forcing relations designed to keep the
complexity of the relation low. They need not already be saturated
semantic forcing. There is a concrete example in
\cite[Definition 3.19(1)]{patey-jumps}: to force a $\Sigma^0_2$ sentence
$\exists x\,\psi(x)$, the displayed witness on each promised part
must be smaller than the condition's tail bound $t$.
Fix $M\geq t$ and take
\[
 \theta_M=\exists x\,\forall z\,[x\geq M].
\]
This is true with witness $M$, and every ordinary internal interpretation
of arithmetic forces it at every stage. The cited raw clause fails at
this condition, because no $x<t$ can satisfy $x\geq M$, and its promise
has a nonempty part. Thus an unsaturated identification of this raw
relation with Kripke--Joyal forcing is false, even for a logical theorem.
Preserving it as a computational predicate and saturating it repairs
the identification without changing the theorem being interpreted.

\subsection{Arithmetic forcing from finite execution witnesses}
Work in the classical dense completion $\mathbb B=\RegOpen(P)$. A designated
binary oracle is a family of complementary values
$b_{q,0},b_{q,1}\in\mathbb B$. For fixed $e,x,y$, let
$\mathcal W_{e,x,y}$ be the numerical codes of candidate finite records
starting program $e$ on input $x$ and halting with output $y$.
For such a record $t$, put
\[
 w_t=\bigwedge_{(q,b)\text{ read in }t}b_{q,b}
\]
when its ordinary transitions and output are valid, and put $w_t=0$ otherwise.
Let $\beta_p=\Dens(\uparrow p)$ be the canonical basic truth value.

\begin{theorem}[Exact arithmetic semantics]\label{cal:thm:adequacy}
For each $e,x,y$,
\begin{equation}\label{cal:eq:executiontruth}
 \val{U^g(e,x)=y}_{\mathbb B}
   =\bigvee_{t\in\mathcal W_{e,x,y}}w_t
   =\Dens R_{e,x,y},
 \qquad
 R_{e,x,y}(p)\Longleftrightarrow
      \exists t\in\mathcal W_{e,x,y}\;(\beta_p\leq w_t).
\end{equation}
Here the last equality identifies the fixed upper set with its element of
$\mathbb B$. The sheaf interpretation of every first-order arithmetical formula
about this oracle is the ordinary Boolean-valued forcing interpretation.
\end{theorem}
\begin{proof}
\leavevmode
\begin{enumerate}[label=\textup{(\roman*)},leftmargin=*,itemsep=3pt,topsep=4pt]
\item \textbf{Tracked implementation.} A fixed finite record is accepted precisely on the intersection of the truth
values of its reads.
Existential quantification over its numerical code gives
the first join in \eqref{cal:eq:executiontruth}, by \cref{prop:intensional-native-checker,cal:thm:dialogues}.
The values $\beta_p$ form a basis: every $v\in\mathbb B$ satisfies
$v=\bigvee\{\beta_p:\beta_p\leq v\}$.
Applying this separately to each $w_t$
gives the second equality.
Persistence makes $R_{e,x,y}$ upper.
\item \textbf{Oracle interpretation.} For numerical equality and the oracle graph, the interpretations agree by
construction.
Finite conjunction and disjunction use Boolean meet and join;
negation uses Boolean complement.
For numerical quantification the constant
numerals cover the natural-number object locally.
Thus
\[
 \val{\exists n\,\varphi(n)}=\bigvee_n\val{\varphi(n)},\qquad
 \val{\forall n\,\varphi(n)}=\bigwedge_n\val{\varphi(n)}.
\]
Induction on formulas proves the assertion.
Computation atoms use
\eqref{cal:eq:executiontruth}; they require no additional logical clause.
\end{enumerate}
\end{proof}

The proof works without countability of $P$, and does not require completely
prime points of its Boolean locale. It directly covers the numerical generic
oracles of the Cohen, Mathias, perfect-tree, closed-class, and bushy examples,
and the characteristic function of the finite-axiom generic in functional
forcing. Sequence oracles with values in $\N$ use the corresponding countable
partition of truth values in place of the two Boolean values.

\subsection{A generic-filter separation proof}
We now perform the classical calibration in the ordinary
external Boolean numerical setting. It is not a use of excluded middle
inside the ground effective topos. For a numerically coded countable preorder $P$, put
$j=\Dens=\neg\neg$ on $\Up(P)$, using the classical identity of
\cref{cal:thm:density}. Its fixed points are the regular-open Boolean truth
values; $\neg U$ consists of conditions with no extension in $U$. The
extension-density nucleus is available internally, but its equality with
double negation is used here only in this classical calibration.

\begin{lemma}[Explicit generic chain and separation]\label{forcing:lem:separation}
Given a countable list of dense subsets of a numerically coded $P$,
there is a filter through each condition meeting every member of the
list. If two upper sets $U,V$ are met by exactly the same filters in
this class, then $jU=jV$.
\end{lemma}
\begin{proof}
\leavevmode
\begin{enumerate}[label=\textup{(\roman*)},leftmargin=*,itemsep=3pt,topsep=4pt]
\item \textbf{Forcing construction.} Starting at $p_0$, let $p_{n+1}$ be the condition with least numerical
code which strengthens $p_n$ and belongs to the $n$th dense set.
Density proves the candidate set is nonempty.
In the external Boolean
interpretation the least element exists, giving a defined recursion;
this asserts no ordinary computability of that selection.
If a
computational implementation is wanted, the corresponding predicates
are exactly decision queries of \cref{prop:decision}.
The downward
closure of the chain is the required filter.
\item \textbf{Verification of the equations.} Suppose $jU\not\subseteq jV$.
Choose $p\in jU\setminus jV$.
Since $p\notin jV$, there is $q\geq p$ with no extension in $V$.
Since $p\in jU$, there is $r\geq q$ in $U$.
The constructed filter
through $r$ meets all the specified dense sets and $U$.
It cannot meet
$V$, since a common stronger condition with a member of $V$ would
contradict the property of $q$.
This contradicts the stipulated equality
of meeting behaviour.
The reverse inclusion is symmetric.
\end{enumerate}
\end{proof}

\begin{lemma}[Arithmetic truth for the regular-open interpretation]\label{forcing:lem:arithmetic-truth}
For a countable family of numerical formulas about a presented generic
object, there is a countable family of dense sets such that, on filters
meeting them, ordinary arithmetic truth is equivalent to meeting the
regular-open Boolean value of each formula.
\end{lemma}
\begin{proof}
\leavevmode
\begin{enumerate}[label=\textup{(\roman*)},leftmargin=*,itemsep=3pt,topsep=4pt]
\item \textbf{Observation preservation.} Close the family under subformulas and substitution of natural-number
parameters; it remains countable.
Include the dense sets ensuring the
specified finite coordinates of the generic object are defined.
An
atomic finite observation has an upper witness set $U$ and Boolean
value $jU$.
Include $D_U=U\cup\neg U$, which is dense: from any
condition either an extension in $U$ exists or the condition is in
$\neg U$. A filter meeting $D_U$ meets $jU$ exactly when it meets
$U$, because $jU$ is incompatible with $\neg U$.
\item \textbf{Forcing construction.} For an existential formula with instance values $B_n$, the Boolean
value is $j(\bigcup_nB_n)$.
Include
$D_{\cup B_n}=(\bigcup_n B_n)\cup\neg(\bigcup_n B_n)$.
A filter meeting this dense set meets the existential value exactly
when it meets some $B_n$.
For negation include $D_B=B\cup\neg B$;
the filter meets $\neg B$ exactly when it does not meet $B$.
For conjunction, directedness identifies meeting $B\cap C$ with
meeting both.
\item \textbf{Existence and uniqueness.} These arguments inductively give truth for conjunction,
negation, and existential quantification; the other classical
connectives and universal quantifier are expressed through them.
Only countably many dense sets were added.
The chain construction of
\cref{forcing:lem:separation} proves they can all be met below every condition.
\end{enumerate}
\end{proof}

\begin{theorem}[Saturation of a computational forcing presentation]\label{forcing:thm:raw-saturation}
Let $U_\varphi=\{p:p\rawforces\varphi\}$ be a persistent numerical
forcing presentation for a fixed named generic object. Whenever its
proved truth lemma identifies meeting $U_\varphi$ with truth of
$\varphi$ on the countably generic filters, one has the exact identity
\[
 jU_\varphi=\val{\varphi}_{\RegOpen(P)}.
\]
In particular the monotonicity, density, and truth lemmas for Patey's
arithmetical relation, \cite[Lemmas 3.20--3.23]{patey-jumps}, identify its
unpointed $\Sigma^0_{n+2}$ and $\Pi^0_{n+2}$ predicates with semantic
truth after this saturation, for the generic real defined there.
\end{theorem}
\begin{proof}
\leavevmode
\begin{enumerate}[label=\textup{(\roman*)},leftmargin=*,itemsep=3pt,topsep=4pt]
\item \textbf{Forcing construction.} Combine the countable requirement families in the given proved truth
lemma and in \cref{forcing:lem:arithmetic-truth}.
Every filter generic for the combined
family meets $U_\varphi$ precisely when it meets the Boolean value
$B=\val{\varphi}$.
Apply \cref{forcing:lem:separation} to $U_\varphi$ and
$B$.
Since $jB=B$, the result is $jU_\varphi=B$.
For the stated application, the conditions and their oracle-computable
promise codes are countable.
\item \textbf{Verification of the equations.} The cited monotonicity lemma makes the
raw predicates upper, and the cited truth and genericity lemmas supply
the countable requirement family.
There is also an explicit Boolean
name for the same generic real, so no identification of names is left
implicit.
For each $n$, pad the finite predicate $n\in G$ to a
$\Sigma^0_2$ sentence $\psi_n$, for example
$\exists u\,\forall v\,[u=0\wedge n\in G]$.
Let
\[
 U_{n,1}=\{p:p\rawforces\psi_n\},\qquad
 U_{n,0}=\{p:p\rawforces\neg\psi_n\}.
\]
Their union is dense by the cited density lemma.
\item \textbf{Universal-property argument.} They are incompatible:
a common strengthening would extend to a filter satisfying the cited
truth lemma, making both a sentence and its negation true.
Thus
$jU_{n,0}$ and $jU_{n,1}$ are complementary Boolean values.
These
pairs, one for each $n$, define the Boolean name for a binary sequence.
A filter meeting the raw decision sets interprets its $n$th bit as
$b$ exactly when it meets $U_{n,b}$, and the cited truth lemma
identifies that bit with the source's generic real.
\item \textbf{Functoriality and coherence.} This proves the
needed name identification before applying the general saturation
argument to arbitrary higher-level formulas.
Its premises are
therefore supplied by constructions and the cited theorems, rather
than appended as an admissibility axiom.
At the first level, the
source's relation is pointed by a part; one retains that pointing
and retains the pointing as part of the relational presentation.
\end{enumerate}
\end{proof}

The filters in these proofs are tests for the specified countable
families. They are not being asserted to be completely prime points of
the entire complete Boolean frame; an atomless complete Boolean locale
need not have such points. The separation argument avoids that false
inference.

Combining \cref{forcing:thm:raw-saturation,thm:forcing-domain-oracle-identity}
identifies the saturated relation with sheaf semantics on the
corresponding localized ideal domain, with arithmetic interpreted by
\cref{arith:prop:closed-interpretation}. The following construction
retains the raw predicate, its complexity presentation, and its
semantic reflection together.

\subsection{Raw forcing predicates as propositions in the glued topos}
\label{sec:raw-forcing-gluing}
We apply the internal frame construction of \cref{glue:thm:frame}
in the Boolean calibration of the preceding truth lemma. Put
$H=\Up(P)$, $j=\neg\neg$, and $B_\varphi=\val{\varphi}_{\RegOpen(P)}$.
\Cref{forcing:thm:raw-saturation} has proved $jU_\varphi=B_\varphi$
for the cited forcing presentation.

\begin{corollary}[Modal interpretation retaining the raw forcing relation]
\label{glue:cor:raw-forcing}
In the glued frame $K_j$, the raw forcing predicate determines the
proposition
\[
 C_\varphi=\gamma_j(U_\varphi)
         =(U_\varphi,B_\varphi).
\]
Its raw projection is exactly $U_\varphi$, its semantic projection
is exactly $B_\varphi$, and
\begin{equation}\label{eq:raw-forcing-open-modality}
 o\Rightarrow C_\varphi=\gamma_j(B_\varphi)
                       =(B_\varphi,B_\varphi).
\end{equation}
The indexed code objects and their raw-predicate relations are
retained by \cref{glue:prop:codes}, with the contextual gluing of
\cref{glue:prop:slices} supplying parameter substitution.
\end{corollary}
\begin{proof}
\leavevmode
\begin{enumerate}[label=\textup{(\roman*)},leftmargin=*,itemsep=3pt,topsep=4pt]
\item \textbf{Modal interpretation retaining the raw forcing relation.} The saturation identity supplies the second coordinate of
$\gamma_j(U_\varphi)$.
Its two projections give the stated
predicates.
Equation \eqref{eq:glue-frame-modality} and
$jB_\varphi=B_\varphi$ prove
\eqref{eq:raw-forcing-open-modality}.
Apply the two cited propositions to the existing code-indexed relations, retaining their numerical indices.
\end{enumerate}
\end{proof}

For the sentence $\theta_M$ at the start of this section,
$U_{\theta_M}\ne1$ but $B_{\theta_M}=1$. Hence
\[
 C_{\theta_M}=(U_{\theta_M},1)\ne(1,1),\qquad
 o\Rightarrow C_{\theta_M}=(1,1).
\]
The proposition $C_\varphi$ records availability according to the
raw forcing clause; its semantic restriction interprets $\varphi$.
The displayed calculation gives their exact modal relationship,
including the bounded-witness discrepancy. It preserves the raw and interpreted components together with their exact modal comparison.

The gluing and slice theorems themselves are internal constructions
over $\Eff$. The application just made uses the Boolean calibration
in which the Patey saturation identity was proved. Arithmetic in a
localization is interpreted by
\cref{arith:prop:closed-interpretation}. These constructions retain
the actual raw complexity presentations; computing their semantic
reflection is the action of the specified modality, whose positive
effective domains were constructed in Part~\ref{part:effective}.


For a requirement layer $j_0\leq\Dens$, the flagged gluing of
\cref{cal:thm:gluing} retains the same predicate as
$(U_\varphi,j_0U_\varphi,\Dens U_\varphi)$. Its dense component is
$B_\varphi$ by \cref{forcing:thm:raw-saturation}; its middle component
records the selected requirements before full dense saturation.
The two marked open-adjunction units recover both comparison maps.

\section{Filtered stabilization and explicit fusion}\label{sec:stabilization}

Finite requirements can be solved individually without producing a coherent infinite object. Stabilization and fusion organise their solutions into compatible finite observations. We construct the resulting limit and verify the preservation of protected fronts, making clear which effective bounds and which generic existence conclusions the construction provides.


Let $D_n$ be finite-observation objects with transition maps $r_n:D_{n+1}\to D_n$. Let $\tau_n:P\to D_n$ satisfy $r_n\tau_{n+1}=\tau_n$. The compatible-sequence object is
\[
 D_\infty=
 \left\{(d_n)\in\prod_{n:\N}D_n\ \middle|\
             \forall n,\ r_n(d_{n+1})=d_n\right\}.
\]
This limit object carries its full $\Omega$-valued membership predicate.

\begin{definition}[Filtered stabilization]\label{def:stabilization}
A refinement chain is filtered-stabilizing when
$\tau_n(p_m)=\tau_n(p_n)$ for every $m\geq n$.
A limit-condition operation for such chains assigns a valid condition $p_\infty$ with
$p_n\leq p_\infty$ and $\tau_n(p_\infty)=\tau_n(p_n)$ for every $n$.
\end{definition}

\begin{proposition}[Canonical limit observation]\label{prop:limit-observation}
A filtered-stabilizing chain determines a unique compatible observation $d_\infty$ whose $n$th coordinate is $\tau_n(p_n)$. The construction is uniform in parameters.
\end{proposition}
\begin{proof}
\leavevmode
\begin{enumerate}[label=\textup{(\roman*)},leftmargin=*,itemsep=3pt,topsep=4pt]
\item \textbf{Canonical limit observation.} Put $d_n=\tau_n(p_n)$.
Then
\[
 r_nd_{n+1}=r_n\tau_{n+1}(p_{n+1})
 =\tau_n(p_{n+1})=\tau_n(p_n)=d_n.
\]
Thus the tuple factors through the defining subobject of $D_\infty$.
Uniqueness follows from its projections.
The same calculation in the slice over a parameter object gives the parameterized map.
\end{enumerate}
\end{proof}

This gives an observation, not yet a condition. We next construct a genuine limit condition for an explicit class of conditions, and verifies its required forcing properties.

\subsection{Finite fronts for perfect trees}
Work with nonempty decidable prefix-closed binary trees and specified finite splitting fronts. A front $F$ in a tree $T$ is a finite antichain of nodes such that every node of $T$ is comparable with some member of $F$. The stronger finite assertion here is useful: it gives the extension property needed in the construction without choosing a branch. A node on a front has two specified incompatible extensions in a subsequent front. Such data certify perfectness by iteration.

\begin{proposition}[Structural successor of a protected splitting front]\label{prop:front-successor}
For a decidable perfect binary tree $T$ and a finite front $F$ whose members have length at least $n$, there is a finite-search construction of a decidable perfect subtree $T'\subseteq T$ and a front $F'$ such that $T'$ agrees with $T$ through level $n$, every member of $F$ has two specified incompatible descendants in $F'$, and every member of $F'$ has length at least $n+1$.
\end{proposition}
\begin{proof}
\leavevmode
\begin{enumerate}[label=\textup{(\roman*)},leftmargin=*,itemsep=3pt,topsep=4pt]
\item \textbf{Structural successor of a protected splitting front.} For each $u\in F$, enumerate pairs of finite extensions of $u$ in $T$ and return the first pair $v_u^0,v_u^1$ that is incompatible and has both lengths at least $n+1$.
Tree membership, prefix order, and incompatibility are decidable.
Perfectness gives a successful pair; least-witness search on the numerical codes therefore constructs it.
Finite induction over the given front assembles all these pairs without a choice operation.
\item \textbf{Observation preservation.} Put
\[
 F'=\{v_u^i\mid u\in F,\ i<2\},\qquad
 T'=\{\sigma\in T\mid \exists v\in F'\;
             (\sigma\preceq v\ \vee\ v\preceq\sigma)\}.
\]
The finite search defining $T'$ is decidable.
The set $F'$ is an antichain and, by construction, a front of $T'$.
For any retained node $\sigma$, choose $v=v_u^i\in F'$ comparable with it.
If $v\preceq\sigma$, perfectness of $T$ above $\sigma$ gives two incompatible extensions retained in $T'$.
If $\sigma\preceq u$, the selected incompatible descendants $v_u^0,v_u^1$ extend $\sigma$.
\item \textbf{Forcing construction.} In the remaining case $u\preceq\sigma\preceq v$, and the two extensions of $v$ provided by perfectness of $T$ remain in $T'$.
Thus every node of $T'$ has incompatible extensions and $T'$ is perfect.
If $|\sigma|\leq n$ and $\sigma\in T$, the front property puts $\sigma$ below some $u\in F$, hence below a chosen descendant in $F'$.
Therefore $\sigma\in T'$, proving agreement through level $n$.
The remaining properties are built into $F'$.
\end{enumerate}
\end{proof}


\begin{theorem}[Uniform constructor for protected front states]
\label{thm:uniform-front-recursion}
Let $S$ be the object of triples $(n,T,F)$ in which $T$ is a
decidable perfect binary tree and $F$ is a displayed finite front
whose members all have length at least $n$. The finite search of
\cref{prop:front-successor} defines one total morphism
\[
 s:S\longrightarrow S,
 \qquad (n,T,F)\longmapsto(n+1,T',F').
\]
Starting with the full binary tree and the root front, NNO recursion
constructs a single family $(T_n,F_n)_{n:\N}$ satisfying the three
fusion equations below. Its membership map
$(n,\sigma)\mapsto[\sigma\in T_n]$ is jointly decidable; hence
$\sigma\mapsto T_{|\sigma|}(\sigma)$ is a decidable tree predicate.
In the effective numerical presentation, when tree codes and the
decidable membership operation are tracked, this recursion yields a
uniform tracked decision procedure for the fusion tree.
\end{theorem}
\begin{proof}
\leavevmode
\begin{enumerate}[label=\textup{(\roman*)},leftmargin=*,itemsep=3pt,topsep=4pt]
\item \textbf{Forcing construction.} Represent decidable tree membership by a map
$2^{<\N}\to2$, with the prefix, perfectness, and front equations
forming the subobject $S$.
For each front member, the relation
``two incompatible extensions of sufficient length belong to $T$''
is decidable on numerical codes.
Perfectness proves that this
numerical search succeeds.
The least-number principle for an
inhabited decidable numerical predicate constructs its least
successful code without countable choice.
Finite recursion over the
supplied front constructs all extension pairs.
\item \textbf{Verification of the equations.} The definitions of
$T'$ and $F'$ in \cref{prop:front-successor} are decidable finite
tests; its proof establishes exactly the properties needed for the
resulting triple to lie in $S$.
Therefore the construction is one
morphism $s:S\to S$ in the indicated parameter context.
\item \textbf{Evaluation and transposition.} Iterating $s$ by the NNO constructs a map $\N\to S$; evaluation of
its tree component is the stated jointly decidable relation.
At
stage $n$ its values satisfy front splitting and agreement with
the previous tree through level $n$, directly by the successor
proof.
Substitution $n=|\sigma|$ gives the decidable fusion
membership test.
\item \textbf{Observation preservation.} In the effective specialization, numerical least
search on the demonstrated total domain, primitive recursion, and
composition are tracked by the specified program operations of
\cref{sec:numericalcodes}; evaluating at $|\sigma|$ is therefore tracked uniformly.
\end{enumerate}
\end{proof}


Consider a decreasing sequence of trees $T_0\supseteq T_1\supseteq\cdots$ with fronts $F_n\subseteq T_n$ satisfying:
\begin{enumerate}
\item every node of $F_n$ has length at least $n$;
\item $T_{n+1}$ and $T_n$ agree on nodes of length at most $n$;
\item every member of $F_n$ has two specified incompatible extensions in $F_{n+1}$.
\end{enumerate}
The third clause retains the splitting certificates. \Cref{prop:front-successor} constructs the structural successor with all three properties by finite search. The theorem below uses the three resulting equations. \Cref{sec:diagonal-fusion} constructs front-preserving refinements satisfying all diagonal requirements in the halting decision extension. The second clause is exactly the finite-level preservation used in the proof.

\begin{theorem}[Constructed fusion tree]\label{thm:fusion-tree}
Let $(T_n,F_n)_{n:\N}$ be a single internal family of decidable perfect trees and displayed fronts with jointly decidable membership, satisfying the preceding three fusion equations. The predicate
\[
 T_\infty(\sigma)\iff T_{|\sigma|}(\sigma)
\]
defines a decidable prefix-closed perfect tree contained in every $T_n$. It has the stabilized finite-level observations and carries the retained splitting certificates.
\end{theorem}
\begin{proof}
\leavevmode
\begin{enumerate}[label=\textup{(\roman*)},leftmargin=*,itemsep=3pt,topsep=4pt]
\item \textbf{Constructed fusion tree.} First prove stabilization: for $m\geq k$, the trees $T_m$ and $T_k$ agree on every node of length at most $k$.
Induct on $m-k$.
At a step $l\geq k$, clause (2) gives agreement on length at most $l$, hence on length at most $k$.
\item \textbf{Observation preservation.} For a node $\sigma$ of length $k$, if $n\leq k$ then $T_k\subseteq T_n$; if $n\geq k$ the stabilization just proved identifies its membership in $T_n$ and $T_k$.
Hence $T_\infty\subseteq T_n$ for every $n$, and indeed $T_\infty=\bigcap_nT_n$ as subobjects.
Decidability is given by the single test $T_{|\sigma|}(\sigma)$ on the stipulated jointly decidable family.
Such a family is constructed explicitly in \cref{thm:uniform-front-recursion} for the structural successor; separate nonuniform external choices of individually decidable trees do not suffice.
\item \textbf{Universal-property argument.} If $\tau$ is a prefix of $\sigma\in T_\infty$ and $|\tau|=l\leq|\sigma|=k$, prefix closure gives $\tau\in T_k$.
Since $T_k\subseteq T_l$, we obtain $\tau\in T_\infty$.
Thus the new predicate is a tree.
\item \textbf{Functoriality and coherence.} Every node $u\in F_n$ remains in all later trees.
Indeed clause (3) gives an extension in $F_{n+1}$ and hence membership in $T_{n+1}$.
Iterating the specified extension gives, for every $m\geq n$, an extension in $F_m\subseteq T_m$; prefix closure gives $u\in T_m$.
Earlier trees contain $T_n$, so $u$ belongs to the intersection, hence to $T_\infty$.
\item \textbf{Observation preservation.} Now let $\sigma\in T_\infty$ and choose $n>|\sigma|$.
Since $\sigma\in T_n$, the front property makes it comparable with some $u\in F_n$.
Its shorter length forces $\sigma$ to be a prefix of $u$.
Clause (3) gives two incompatible extensions of $u$ in $F_{n+1}$, and both remain in $T_\infty$ by the previous paragraph.
Thus every node has two incompatible extensions in $T_\infty$.
This proves perfectness with the retained witnesses.
\item Finally, on nodes of length at most $n$, membership in $T_\infty$ agrees with membership in $T_n$ by stabilization.
Therefore its finite-level observations are exactly the stabilized observations, as required for a limit condition.
\end{enumerate}
\end{proof}

The finite-front organization is the finite-front method in perfect-set forcing \cite{sacks}. Here its effect is separated into the finite observation limit and the proof that the limit is still a condition. For the structural sequence constructed by \cref{thm:uniform-front-recursion}, the displayed recursion supplies a uniform decision program for $T_\infty$ in the effective numerical interpretation. The diagonal refinements of \cref{fusion:thm:diagonal} provide a requirement-bearing application with a uniform oracle-relative successor. Bare density by itself does not yield a tracked successor for an arbitrary different requirement.

A fusion limit need not itself generate a point of the structural locale by its principal ideal: one must still check covers. To obtain a point from the original chain, apply \cref{prop:chain-point}; to obtain a point below the fusion condition, continue a cover-continuous construction there. The role of fusion is to construct a common refinement while retaining finite decisions.

\subsection{Protected-front fusion meeting all diagonal requirements}\label{sec:diagonal-fusion}
\Cref{thm:uniform-front-recursion,thm:fusion-tree} construct the structural successor and limit. We now construct front-preserving refinements meeting every diagonal requirement.
Use the same ground enumeration $\varphi_e$ and its halting oracle $h$.
For each $e$, query $h(e)$. If it is positive, simulate until its value
$v_e$ is obtained; otherwise mark the coordinate as unrestricted. Define
a decidable predicate relative to $h$ by
\[
 A(e,b)\ \Longleftrightarrow\
 \text{the marked value is absent, or is not in $2$, or differs from $b$}.
\]
At least one of $A(e,0),A(e,1)$ always holds.

\begin{theorem}[Perfect diagonal fusion]\label{fusion:thm:diagonal}
There is a perfect binary tree $T_\infty$, decidable using $h$, every
branch $X$ of which satisfies
\[
 \varphi_e(e)\converges\quad\Longrightarrow\quad X(e)\ne\varphi_e(e).
\]
It is produced by a uniformly decidable protected-front fusion whose
stage operations and splitting certificates are explicit.
\end{theorem}
\begin{proof}
\leavevmode
\begin{enumerate}[label=\textup{(\roman*)},leftmargin=*,itemsep=3pt,topsep=4pt]
\item \textbf{Perfect diagonal fusion.} Construct a computable strictly increasing sequence $d_0<d_1<\cdots$
of indices of identically diverging programs.
For the finite register-program syntax of \cref{thm:operational-certificates}, append arbitrary distinct unused instructions
after an unconditional loop; these give infinitely many distinct codes.
Search these codes for the next larger one.
This constructs the padding
sequence and proves its divergence.
\item \textbf{Observation preservation.} Set $\ell_n=d_n+1$, and put
\[
 \begin{split}
 T_n&=\{\sigma:\forall e<\min(|\sigma|,\ell_n)\ A(e,\sigma(e))\},\\
 F_n&=T_n\cap2^{\ell_n}.
 \end{split}
\]
Membership is a finite Boolean test using $h$, jointly in $n,\sigma$.
Every $T_n$ is prefix closed and perfect: beyond its finitely many
restricted coordinates both bits are available.
Every shorter node
extends to length $\ell_n$ by choosing the least allowed bit at each
coordinate; every longer node has its prefix in $F_n$.
Hence $F_n$ is
the required finite front.
\item \textbf{Universal-property argument.} The trees decrease and $T_{n+1}$ agrees with $T_n$ through length
$\ell_n$.
Each $u\in F_n$ has two specified extensions in $F_{n+1}$:
at the unrestricted coordinate $d_{n+1}$ choose different bits, and
at the other new coordinates choose the least allowed bit.
These
extensions are incompatible and retain $u$.
Since $\ell_n\geq n+1$,
the finite-level fusion equations of \cref{thm:fusion-tree} all hold.
\item \textbf{Observation preservation.} Their limit has the direct description
\[
 T_\infty=\{\sigma:\forall e<|\sigma|\ A(e,\sigma(e))\}.
\]
It is decidable using $h$, and it equals $T_{|\sigma|}(\sigma)$ on each
input $\sigma$, since $\ell_{|\sigma|}>|\sigma|$.
The retained splitting
certificates prove perfectness, also directly seen at the arbitrarily
large free coordinates $d_n$.
Once $\ell_n>e$, requirement $e$ is
imposed at stage $n$ and preserved thereafter.
This proves the displayed
diagonal assertion for every branch.
\end{enumerate}
\end{proof}

No ground computable binary sequence is a branch: its own program index
would contradict the diagonal assertion. In particular the use of the
decision oracle is substantive. A ground decidable perfect tree has a
computable path by successive finite searches for splitting extensions,
so a claim that this same construction produced a ground computable
perfect tree would be false. This example proves a genuine
front-preserving requirement operation; it does not claim that every
dense subset of every perfect-tree order has a ground computable
refinement operator. The general perfect-set forcing theory remains
available by citation to \cite{sacks}.


Fusion turns compatible finite requirement solutions into a native infinite object while preserving the specified protected data. Its numerical programs remain those of the chosen interpreter, and the final implementation examples check the resulting generic values against their named finite execution witnesses.

\section{Intensional enhancement: concrete generic implementations}
\label{sec:intensional-external}

The forcing constructions now have both semantic points and numerical execution interfaces. We close the concrete development by comparing a one-query Cohen bit, a finite-axiom functional, and named kernels for restrained computations. These examples test the same distinction between a program's definedness and the effectiveness of its forcing observation.


\begin{theorem}[A one-query implementation of the Cohen generic bit]
\label{thm:intensional-cohen-read}
In the structural Cohen presentation, let $A_n$ be the local
object of valid coordinate answers $i\in2$, classified by
$U_{n,i}$. It is functional: the two answer opens are disjoint.
Its support is the coordinate-decision open
$U_{n,0}\vee U_{n,1}$, which becomes top in the Cohen sheaf topos.
The one-query tree
\[
 \operatorname{read}(n)=\operatorname{ask}
  \bigl(n,a\mapsto\operatorname{return}(\operatorname{bit}(a))\bigr)
\]
is denotationally coherent and its interpretation is the native
generic coordinate $\dot g(n):2$. For all $n$ and $i\in2$,
\[
 \llbracket\operatorname{read}(n)=i\rrbracket
  =q_J(U_{n,i}).
\]
This is an oracle-relative description of a native generic
computation, not a claim that the whole generic stream has a name
in the transported ground-model stream object.
\end{theorem}
\begin{proof}
\leavevmode
\begin{enumerate}[label=\textup{(\roman*)},leftmargin=*,itemsep=3pt,topsep=4pt]
\item \textbf{Oracle interpretation.} The answer relation is the coproduct of the two disjoint opens.
Binary splitting makes its projection onto the query object epic
after sheafification, while disjointness makes it monic.
Thus its
inverse is the coordinate function $\dot g$ by
\cref{prop:answer-reconstruction}.
Under a valid answer $a$,
$\operatorname{read}(n)$ returns $\operatorname{bit}(a)$, which is
also $\dot g(n)$ on that answer region.
Since the regions jointly
cover the query context, separatedness proves equality of the
interpreted maps.
\item \textbf{Verification of the equations.} The truth of the equality with $i$ is precisely
the corresponding coordinate-open interpretation, namely
$q_J(U_{n,i})$.
Two simultaneously valid answers have the same
bit because their validity overlap is empty when the bits differ,
so the protocol satisfies the coherence relation.
\end{enumerate}
\end{proof}

The stream computed one coordinate at a time above is the native
sequence of \cref{seq:prop:cohen}. Its finite observations agree with
the generic ideal by \eqref{eq:native-sequence-cylinder}, while
\cref{prop:intensional-new-streams} proves that the whole sequence
has no local lift to the imported ground sequence object.


\begin{corollary}[Named Mathias forcing has a proper effective scope]
\label{cor:intensional-mathias-scope}
For the named evaluator kernels of
\cref{thm:named-kernel-domain}, the uniformly bounded-use Mathias
family factors through $T_{\rho_M}^{\rm adm}$, while the
uniformly parametrized unbounded-use parity-of-first-member
family does not. Thus the effective named forcing observation
is neither determined by bare code existence nor a total
operation on every effective name.
\end{corollary}
\begin{proof}
\leavevmode
\begin{enumerate}[label=\textup{(\roman*)},leftmargin=*,itemsep=3pt,topsep=4pt]
\item \textbf{Observation preservation.} The classifier constructed in \cref{thm:mathias-bounded}
gives the required stable lift for bounded use.
The program
constructing the least selected integer has an effective
computation kernel, but \cref{thm:mathias-unbounded} proves that
its forced stage truth has no uniform $\Sigma$-classifier.
\Cref{thm:named-kernel-domain} identifies these two cases
exactly with factorization, respectively nonfactorization, through
the maximal admissibility domain.
\end{enumerate}
\end{proof}


\section{The assembled theory and its computational use}\label{sec:assembly}

The separate completion, forcing, oracle, and implementation comparisons can now be stated together. The purpose of this synthesis is not to conflate their objects, but to specify the maps between them and the equations those maps satisfy. It identifies which structures are native, which are imported, and which retain raw descriptions before localization.


\begin{theorem}[Site construction theorem]\label{thm:site-assembly}
A condition site $(P,J)$ gives canonically:
\begin{enumerate}
\item the nucleus $j_J$, frame $\Omega_J$, locale $L_{P,J}$, and localic sheaf topos $\FF_{P,J}$;
\item the generic continuous ideal, full Kripke--Joyal forcing semantics, and the free ideal-completion comparison;
\item the stable marking $S_J$ and its uniform stage-observation pairing;
\item for each atomic oracle interpretation satisfying its geometric relations, a unique localic realization;
\item the maximal admissible subobject of every specified effective oracle observation family.
\end{enumerate}
A cover normal-form calculation of the kind in \cref{thm:normal-form} constructs the effective saturation map on its exact displayed domain. For an effective evaluator kernel, a lift of that kernel through the admissibility object gives a uniform stage semidecision for its forced truth.
\end{theorem}
\begin{proof}
\leavevmode
\begin{enumerate}[label=\textup{(\roman*)},leftmargin=*,itemsep=3pt,topsep=4pt]
\item \textbf{Ideal and directedness calculation.} The nucleus and fixed-frame constructions are \cref{prop:site-nucleus,prop:fixed-frame}.
Their universal model and full semantics are \cref{prop:site-classifier,thm:generic-ideal,thm:kj}; \cref{thm:site-scott} supplies the ideal-completion comparison.
The stable marking and admissibility domain are the counits of \cref{thm:admissibility-coreflection} for the displayed marking maps.
An atomic interpretation extends uniquely by \cref{thm:rounded-classifier,cor:atomic-realization}, and stage evaluation gives the positive pairing.
\item Finally \cref{thm:normal-form} constructs the saturation classifier, while \cref{lem:kernel-effective,thm:kernel-forcing} identify the evaluator kernel and its stage trace.
Pullback uniqueness makes these constructions commute on each specified family.
\end{enumerate}
\end{proof}

\begin{theorem}[Forcing, domain sheaves, and oracle localization]\label{thm:forcing-domain-oracle-identity}
For a condition site $(P,J)$ put $D=\Idl(P)$. The inverse frame
isomorphisms
\[
 \Phi(U)(I)\iff\exists p\in I\,U(p),\qquad
 \Psi(V)(p)\iff V(\downarrow p)
\]
identify $\Up(P)$ with $\Sc(D)$. Put $\bar j=\Phi j_J\Psi$ and
let $L_{\bar j}$ have the fixed frame of $\bar j$. For the contextual
cover-answer interface $d_{\mathcal R}$ in the refinement presheaf
topos $\widehat P$, there are canonical equivalences
\[
 \FF_{P,J}\simeq\Sh_\E(L_{\bar j})
       \simeq\Sh_{o_{d_{\mathcal R}}}(\widehat P).
\]
They identify the represented condition opens, generic membership,
and associated sheaf reflectors. Transport of a retained name cover,
its denotational kernel pair, or a finite-run domain--value span gives
the same object or map in these three presentations.
The common sheafification adjunction constructs an Artin gluing
$\mathcal G_j\simeq\Sh_\E(K_j)$ whose semantic open is this forcing
topos. Its fully faithful embedding $\Gamma_j$ retains the raw
presentation and satisfies $s\Gamma_j=a$ and
$o_*s\Gamma_j\cong\Gamma_jia$. The condition-frame and Scott-frame
presentations identify this same glued frame and open modality.
\end{theorem}
\begin{proof}
\leavevmode
\begin{enumerate}[label=\textup{(\roman*)},leftmargin=*,itemsep=3pt,topsep=4pt]
\item \textbf{Scott continuity.} The frame maps and their inverse identities are
\cref{prop:scott-alex}; transporting the nucleus and its generic
generators gives \cref{thm:site-scott}.
The contextual covering-sieve
inclusions generate exactly the topology $j_J$ on the whole
presheaf topos by \cref{thm:full-cover-oracle}, which also identifies
the reflectors. A represented condition is sent to the coordinate
open $I\mapsto p\in I$, identifying generic membership.
The domain--value and certificate assertions are
\cref{thm:partial-transport,thm:intensional-certificate-transport}.
\item \textbf{Inverse comparison.} The retained presentation and kernel-pair equations are
\cref{thm:intensional-pres-fibration}.
All these equations are
preserved by the displayed equivalences.
\Cref{glue:thm:topos,glue:thm:frame} constructs the gluing and its
Scott-frame presentation from this adjunction; its action on the
retained diagrams is \cref{glue:prop:codes}.
The two semantic
presentations identify the right adjoints as well as the reflectors,
and hence identify the corresponding comma categories.

Naturality is expressed by the commutative diagram
\[
\begin{tikzcd}[column sep=large,row sep=large]
\Up(P)\ar[r,"\Phi"]\ar[d,"j_J"']&\Sc(D)\ar[d,"\bar j"]\\ \Up(P)\ar[r,"\Phi"']&\Sc(D)
\end{tikzcd}
\]
\end{enumerate}
\end{proof}

\begin{corollary}[Calibration of the constructed forcing sites]
\label{cal:cor:cases}
The structural Cohen, next-element Mathias, perfect-tree level,
closed-class jump-decision, enumerated bushy, and restrained-functional
requirement sites admit the dense localization of
\cref{cal:thm:localize} in the bases where their displayed successor
operations are defined. Classically, the dense completion of each specified
carrier is its regular-open forcing topos. The full Mathias, Sacks,
closed-class, basic bushy, and restrained-functional orders are identified
by their constructed carriers and code comparisons; the computable fragments
remain separate. Their designated numerical oracles
have the dialogue and Boolean-valued arithmetic semantics of
\cref{cal:thm:dialogues,cal:thm:adequacy}.
\end{corollary}
\begin{proof}
\leavevmode
\begin{enumerate}[label=\textup{(\roman*)},leftmargin=*,itemsep=3pt,topsep=4pt]
\item \textbf{Calibration of the constructed forcing sites.} Cohen conditions extend to decide any next coordinate.
Mathias conditions
append another reservoir element.
Inside any stronger perfect tree, a node
on a required finite level gives a cone restriction.
The closed-class
exhaustion decisions construct a nonempty subtree deciding the next jump
question and extending the stem.
Enumerated bushy conditions extend by
\cref{bush:thm:successor,cal:thm:bushy} in the halting decision base.
Restrained functionals extend by protection and the unused-prefix
use-growth operation of \cref{sec:restrained-forcing}.
\item \textbf{Localization and coverage.} Each construction
works at every stronger input condition, proving extension-density.
Apply \cref{cal:thm:localize,cal:cor:dense-oracle}.
\item \textbf{Tracked implementation.} The classical condition identifications and code comparisons are
\cref{prop:cohen-comparison,cal:prop:mathias,cal:prop:closed-codes,cal:prop:ks-codes,cal:prop:bushy-codes}
and the full perfect-tree calculation in \cref{sec:least-history-examples}.
The c.e. bushy carrier is retained separately from the unrestricted carrier,
as \cref{cal:prop:bush-notdense} requires; the same distinction for
computable perfect trees is \cref{cal:prop:sacks}.
Finite coordinates define
the indicated generic oracle in each site.
\item \textbf{Functoriality and coherence.} Its finite executions and
arithmetical interpretation are therefore exactly those of
\cref{cal:thm:dialogues,cal:thm:adequacy}.
\end{enumerate}
\end{proof}

\begin{theorem}[Information-domain construction theorem]\label{thm:domain-assembly}
For a searchable approximation presentation $C$, a nucleus on $\Sc(D_C)$, and a presented oracle map $F:D_C\to D_X$, the direct realization has frame map $q_jF^{-1}$. The native generic oracle is obtained by applying $F^\pi$ to the native generic rounded ideal; its full Scott observations and the native partial-evaluator kernel equations are computed unconditionally by \cref{thm:native-generic-computation}. The base effective observation family on which it has effective fixed truth is $\Adm(\rho_F)$. Its pairing with $\Pt D_C$ gives uniform information-state observation.

Whenever the completion comparison reconstructs an old-object generic representative, its relative Scott point is exactly the direct localic realization. All presented partial computations on that representative agree, on domains and values, with their compared native computations.
\end{theorem}
\begin{proof}
\leavevmode
\begin{enumerate}[label=\textup{(\roman*)},leftmargin=*,itemsep=3pt,topsep=4pt]
\item \textbf{Naturality and reindexing.} The inverse image is $q_jF^{-1}$ by \cref{prop:direct}; \cref{thm:native-generic-computation} computes its native ideal, observation and partial-domain equations using the naturality square of \cref{prop:base-approximation}.
Applying the marking coreflection gives $\Adm(\rho_F)$ and its effective-point pairing.
When the comparison reconstructs an old-object representative, \cref{thm:object-generic} identifies its frame map with $q_jF^{-1}$.
\item \textbf{Coherence equations.} Substitution of that representative into the domain--value comparison of \cref{thm:partial-transport} identifies both the partial domain and value map.
These comparisons preserve the defining projections, hence commute.
\end{enumerate}
\end{proof}

\begin{theorem}[Local realization construction theorem]\label{thm:span-assembly}
For a local realization span $L_j\leftarrow\widetilde L\to X$ with surjective left leg, the candidate image and the descended observation locale are canonically constructed. The effective descended fragment is
\[
 M_X\times_{\OO(\widetilde L)}S_j.
\]
Its positive join operations have the exact constructed scopes of \cref{thm:positive-descended-joins}.
The span descends globally precisely when its candidate projection is an isomorphism, equivalently when every oracle open descends. Open surjective left legs admit stable kernel-pair descent. Localization and isomorphism comparisons act as in \cref{thm:localization-transport,prop:comparison-levels}.
\end{theorem}
\begin{proof}
\leavevmode
\begin{enumerate}[label=\textup{(\roman*)},leftmargin=*,itemsep=3pt,topsep=4pt]
\item \textbf{Pullback construction.} Use the localic image nucleus for the candidate, the pullback frame for the pushout locale, and the pullback of markings for the effective fragment.
Their universal properties are \cref{prop:candidate,thm:descent-pushout,thm:marked-descent}.
The equivalent global descent criteria are \cref{prop:candidate,cor:full-descent}.
The open-map calculation proves stable descent in \cref{thm:open-descent}.
The localization and comparison statements follow from the actual pullback, image, and factorization maps already constructed in the cited results.
\end{enumerate}
\end{proof}

\begin{theorem}[Oracle-modality and effective-operation bridge]\label{thm:oracle-modal-assembly}
For a condition-site forcing package with a realization $\rho:L_j\to X$, the following constructions are compatible in their specified scopes.
\begin{enumerate}
\item An affirmative or multivalued response-open interface yields an internal answer relation with support \eqref{eq:oracle-support}. Its least answer-existence localization is the Ahman--Bauer oracle modality corresponding to the relative nucleus generated by that support.
\item The cover-answer family of the structural topology generates exactly the coverage nucleus, and the full oracle-sheaf comparison identifies the entire original forcing sheaf category and its associated sheaf reflector. Under subsequent localizations, the nucleus--topology correspondence and the tower isomorphism identify the resulting relative modal layers.
\item A selected effective answer family gives a uniform stage semidecision precisely on its pullback-defined admissibility domain. Effective restriction along localization is normal lax, and its preservation and reflection domains are the pullbacks of \cref{thm:lax,prop:effective-reflection}.
\item If the answer interface is used by a well-founded adaptive protocol, its transcript object has the constructed domain and value. Its coherent subobject is exactly the equifoliate-tree submonad, and its interpretation presents ordinary sheafification as the modal quotient by the explicitly calculated kernel pair; the finite-depth theory is its bounded fragment.
\item The canonical positive history domains supply the additional observation operation \eqref{eq:history-observation} on its maximal search domain, without assuming an overt history basis. For a local realization span, the effective descended marking $M_D=M_X\times_{\OO(\widetilde L)}S_j$ induces compatible oracle predicates and identical generated modalities in the common pointed base.
\end{enumerate}
These comparisons are internal and uniform in their stated parameter contexts. They distinguish equality of generated propositional modalities from equality of response objects, histories, effective presentations, or unrestricted sheafification constructions.
\end{theorem}
\begin{proof}
\leavevmode
\begin{enumerate}[label=\textup{(\roman*)},leftmargin=*,itemsep=3pt,topsep=4pt]
\item \textbf{Image and quotient construction.} For (1), apply \cref{prop:answer-reconstruction,prop:answer-support-quotient,thm:ab-comparison} to the image of the answer projection.
The cover-answer universal property and the exact comparison in \cref{cor:cover-ab} give (2), and the tower coherence is \cref{thm:tower}.
Item (3) is the admissibility pullback followed by \cref{thm:lax,prop:effective-reflection}.
The well-founded protocol theorem and \cref{thm:intensional-equifoliate,thm:intensional-kernel-pair,thm:intensional-monad-morphism} give (4), with the finite theorem obtained on the bounded images.
\item \textbf{Pullback construction.} The canonical positive history-output and search domains, with their overt specialization, give the first statement of (5), while the marked-pushout pullback and \cref{prop:modal-descent} give the second.
Every step is a construction already made in the indicated theorem, and their composites are equal because they preserve the defining pullback projections, actual response values, and modal truth predicates.
\end{enumerate}
\end{proof}

\begin{theorem}[Computationally structured sheaf realization]
\label{thm:intensional-assembly}
For the constructed condition-site localization, the displayed
computational signature determines a model $M$ of $\Comp$ and a
marked structure sheaf $\Ocomp$ on
$\mathcal G_j\simeq\Sh_\E(K_j)$. Its retained core is
$(o,\Gamma_{j,*}M)$, a model of $\Comp^\bullet$, and its native
extension consists of the constructed comparison sorts. It has the
following properties.
\begin{enumerate}
\item Closed restriction recovers $M$ faithfully, open restriction
recovers $a_*M$, and the open-adjunction unit recovers their comparison.
The complete structure and its exact implementation diagrams satisfy
effective descent and coherent parameter reindexing. For
$j_0\leq\Dens$, the flag of \cref{cal:thm:gluing,cal:cor:structure}
retains the raw, requirement, and dense layers in one structure and
recovers both localization maps.
\item Name covers, denotational kernel pairs, and exact partial
domain--value spans transport through forcing inverse image.
The description interpretation satisfies
\eqref{int:eq:transportalgebra} and \eqref{int:eq:nativetree}, including substitution
and common-refinement coherence. Native completion comparison retains
old names covering precisely $\operatorname{Rep}_{f,C}$ in the full
native completion, with both denotational relations.
\item The oracle W-type and its coherent name object present
sheafified denotations by the regular epimorphism with modal kernel
pair of \cref{thm:intensional-oracle-coding}. The native comparison
$\kappa$ respects supplied substitution and preserves and reflects
finite successful transcripts, including their supports and values.
Its coherent evaluation is $a_j\delta$.
\item Accepted executions and checked constructor transformations
transport with their supports. A native numerical Boolean oracle
defines an internal partial combinatory algebra with native
semidecidable application and its own finite execution witnesses. The
compiled dialogue translations preserve its exact domain and values;
classically its arithmetic is the Boolean-valued forcing interpretation.
\item The domain $T_\rho^{\rm adm}$ is the maximal domain of uniform
effective forcing observations on named computations. The marked
structure is compatible with the Scott-frame presentation of $K_j$;
its generating approximation and answer data recover the raw and
native completions and protocol monads on their respective
restrictions.
\end{enumerate}
\end{theorem}
\begin{proof}
\leavevmode
\begin{enumerate}[label=\textup{(\roman*)},leftmargin=*,itemsep=3pt,topsep=4pt]
\item \textbf{Partiality and definedness.} Construct $M$ from the actual name and witness covers by
\cref{struct:ex:numerical,thm:partialadequacy,cor:presented-execution-witnesses};
\cref{struct:prop:syntax} gives their single regular presentation.
For (1), apply
\cref{struct:thm:modelgluing,struct:prop:open,struct:prop:descent,glue:prop:slices,cal:thm:gluing,cal:cor:structure}.
Item (2) combines
\cref{thm:intensional-pres-fibration,thm:partial-transport,int:thm:transport,int:thm:refinement}
with the image calculation of
\cref{struct:prop:image,thm:intensional-old-image}.
\item \textbf{Image and quotient construction.} For (3), use the named oracle quotient and
\cref{struct:thm:Wcomparison,struct:thm:transcript,struct:cor:coherent};
\cref{struct:cor:nativeresult} also treats native result objects.
Item (4) is
\cref{thm:intensional-certificate-transport,prop:intensional-native-checker,struct:thm:pca,cal:thm:dialogues,cal:prop:transport,cal:thm:adequacy}.
\item Finally \cref{thm:named-kernel-domain,thm:kernel-forcing} gives the
maximal observation domain, and
\cref{glue:thm:frame,struct:prop:generators} gives the structural
recovery in (5).
The comparisons preserve their defining projections
and the actual transcript values; uniqueness of image, pullback,
initial-algebra, and descent maps gives their coherence.
\end{enumerate}
\end{proof}

\subsection{How the techniques fit a proof}
For a concrete oracle-forcing argument, one starts with conditions and cover relations or with an approximation matrix and a nucleus. The basis theorem translates between these forms when their displayed generators satisfy its equations. A small atomic interpretation then constructs the oracle realization. If the oracle is only available after a choice, a local realization span is constructed instead.

The computational question is a selected family of evaluator kernels. A region calculation may show that its forcing truth is a finite certificate formula, as in bounded-use Mathias forcing. Alternatively a structural saturation formula may provide its positive classifier, as in Cohen forcing. The output of this calculation is a map into the stable marking. It gives an actual uniform observation operation. The unbounded Mathias and Booleanization calculations explain why this operation cannot be extended automatically to all effective oracle predicates.

When a theorem adds requirements, their answer maps or cover equations generate the localization. The tower theorem gives its fixed frame, while \cref{thm:localization-transport} restricts regions, histories, and local candidates. A surviving least-history section remains least. A locally available oracle is usable on exactly the descended observation fragment computed by the pushout; effective use is selected by its marking.

Finally, a native generic model is always available from the classifier, and its native partial evaluation already has the unconditional domain-and-observation equations of \cref{thm:native-generic-computation}. An object representative in the pulled-back old completion is obtained through the actual comparison when it is exact; that additional representative then admits the old-object domain-and-value comparison. A continuous ideal selects an external point, while the diagonal fusion of \cref{fusion:thm:diagonal} supplies a requirement-bearing limit condition. The closed-class chain of \cref{closed:thm:low-basis} meets its explicitly generated covers, and \cref{forcing:thm:raw-saturation} gives the classical arithmetic truth identity for the cited raw forcing presentation. The native sequence and path identifications are
\cref{seq:thm:prefixes,seq:cor:paths}. The raw truth identity lifts
to the open-modal equation of \cref{glue:cor:raw-forcing}, with code
presentations organized by \cref{thm:intensional-assembly}. Their
marked computational structure retains the construction and recovers
its raw and semantic restrictions.

\paragraph{Scope of the forcing applications.}
The full carriers and the computational fragments remain distinct:
\cref{cal:prop:sacks,cal:prop:bush-notdense} gives explicit failures of
density. The next-element Mathias site and the Ellentuck site have different
coverages but the same classical dense completion by
\cref{cal:prop:mathias}. The exact comparisons use native full carriers in
the indicated bases; a change-of-base comparison from an imported carrier
retains its actual represented image.

The advanced bushy minimal-degree construction requires the tree, its two
growth-rate functions, and their splitting and fusion operations from
\cite[Section 5]{khan-miller}. Full Kumabe--Slaman degree coding requires
the essential-tree and requirement-dependent restraint operations of
\cite[Section 8.2]{slaman-woodin}. The basic successor and use-growth
constructions above identify their stated condition semantics and numerical
evaluation, while these stronger applications retain their specific
combinatorial proofs. Likewise the split Ramsey solution map of
\cref{cal:prop:ramsey} retains its histories; comparisons with a particular
Ramsey or partition forcing act on that forcing's full condition data.

\subsection{Position relative to the literature}\label{subsec:literature}
Enriched and KZ category theory supply the completion and splitting framework \cite{stubbe,hofmannsousa,hofmannwaszkiewicz}. Dominances and synthetic computability supply positive observations and partial elements \cite{rosolini,escardoknapp,bauer,bauerfixed}; their realizability interpretation uses \cite{hyland,vanoosten,frey}. The constructions here specify the positive pairing, general-carrier presentations and exact implementation comparisons. Restriction categories organize the broader partial-domain calculus \cite{cockettlack}.

The intensional organization uses P-categories and exposures
\cite{kavvosarticle}, with literal syntax and staged interpretation
distinguished as in \cite{daviespfenning,chenko}. Substitution is the
finite-tree algebraic construction, with the context-indexed account
of binding in \cite{fioreplotkinturi}. The temporal indexed theory
connects guarded domain theory \cite{birkedalguarded} with local
introspection \cite{ramesh}. The manuscript's compatibility results
are the exact interpretation, refinement and localization equations
\cref{int:thm:exactsyntax,int:thm:refinement,int:thm:transport}, realized
in the computational structure sheaf by \cref{struct:thm:modelgluing}.

Locales, sheaves and geometric theories provide the forcing and classifier framework \cite{johnstone,maclanemoerdijk,vickers,joyaltierney,vickersau,sterlingye}. Oracle sheafification, equifoliate computation and the extended-Weihrauch description are due to Ahman--Bauer \cite{ahmanbauer}; instance reductions and local-operator comparisons use \cite{bauerinstance,kihara}. The present comparisons retain effective observation domains, histories and program presentations alongside those constructions.

The classical dense-sheaf/Boolean-algebra correspondence is made explicit
in \cite[Theorem 4.22]{smykalla-mariano}; the distinction between mixing and
fullness, and the correspondence of mixing models with sheaves, is treated in
\cite[Section 4.2 and Proposition 5.3]{pierobon-viale}. Here
\cref{cal:thm:density,cal:thm:localize} separates constructive extension-density
from internal Booleanization, while \cref{cal:thm:adequacy} identifies the
arithmetic of the designated numerical interpreter. The Mathias/Ellentuck
comparison uses the basis description in \cite{dobrinen-ramsey}; the full
perfect-tree order is \cite[Definition 21]{geschke-quickert}.

Oracle adjunction and hierarchy--jump calibration are related to \cite{relative-recursion,faber-oosten,forster-kirst-mueck}. The forcing constructions retain their specific attribution: bushy methods to \cite{khan-miller} and the DNC separation to \cite[Theorem 1.8]{ambos-spies-etal}, arithmetical partition requirements to \cite{patey-jumps}, low-basis reductions to \cite{jockusch-soare}, and perfect-set and restrained-functional methods to \cite{sacks,slaman-woodin}. The displayed successor and fusion operations expose their actual oracle queries; no optimal jump bound is inferred merely from the Ramsey construction.

Artin gluing and its frame presentation use \cite{sga4-gluing,faul-manuell-frames,faul-manuell-siqueira}, with a related computational internal-language treatment in \cite{sterling-tait}. The regular syntactic presentation follows \cite{caramello}; \cref{sec:computational-structure} constructs its gluing, descent, and marked recovery. Its ordinary structured-topos formulation is distinguished there from the local-morphism convention of \cite{lurie}. The native-sequence proofs here use prefix domains and the Cohen construction; \cite{bhatt-scholze-proetale,scholze-condensed} provide a broader comparison between transported and native objects. Torsors, units and the square-zero deformation pattern are developed in \cite[Tags 03AG, 03AJ, 040D and 0C6R]{stacks-project}. The finite projective resolutions and obstruction classes are calculated explicitly here. The abelian gluing is the comma recollement of \cite[Sections~3.2 and~4.1]{franjou-pirashvili}, with its adjoints, defects and projection formulas computed for the displayed sheafification adjunction.


\appendix
\part*{Construction guide}

The preceding parts supplied several completions, restrictions, localizations, and comparison maps. This guide indexes them by the universal operation that constructs each one, rather than replaying the proofs. It records the interfaces needed to use a result without confusing the full semantic object with its effective or named fragment.


\section{Guide to the universal constructions}\label{app:guide}


The table records the principal construction dependencies. Detailed compatibility statements are collected in \cref{sec:assembly}.
\begin{longtable}{L{.25\textwidth}L{.37\textwidth}L{.29\textwidth}}
\toprule
Construction & Universal operation & Main comparison\\\midrule
\endhead
Computational foundation & Canonical covers, stream lifts and descent & \cref{thm:relativefoundation,lem:modestclosure}\\
Partiality and quantifiers & Universal domain span; subobject adjunctions & \cref{prop:partialmonad,lem:predicate-bc,prop:quantifier-closure}\\
Ordinary completion & Free directed completion and idempotent splitting & \cref{prop:freeextension,thm:rounded-matrix}\\
Admitted and observed completion & Positive-pairing pullback, KZ algebras and comparison & \cref{jd:thm:J,jd:thm:ccc,jd:thm:splitcomparison}\\
Effective presentations & Positive matrices and observed directed covers & \cref{thm:relations,thm:fulladequacy}\\
Spectra and intrinsic computation & Trace transpose and effective points & \cref{jd:thm:spectra,thm:observationfullness,thm:syntheticms}\\
Iteration and feedback & Distributive law and exact finite returns & \cref{jd:thm:distributive,exec:thm:elgot,exec:thm:trace,exec:thm:comparison}\\
Implementations & Name fibration, kernel pairs and exact domains & \cref{thm:canonicalcoding,exec:thm:codefeedback,thm:operational-certificates}\\
Intensional structure & Cartesian exposure and exact description algebras & \cref{int:thm:exposure,int:thm:exactsyntax,int:thm:refinement,int:thm:ifp}\\
Temporal interpretation & Prior--later adjunction and represented homs & \cref{int:thm:temporal,int:prop:temporaltransport}\\
Numerical finite information & Program-cover descent and witness transformations & \cref{thm:recursion,thm:finiteinformation,prop:finite-information-parameters,thm:ms}\\
Forcing presentations & Basis reflection and geometric ideal classifiers & \cref{thm:basis,thm:rounded-classifier,prop:site-classifier,thm:points-ideals}\\
Dense forcing calibration & Extension-density, code lifting and cover-answer localization & \cref{cal:thm:density,cal:lem:codes,cal:cor:dense-oracle,cal:cor:cases}\\
Effective observation & Indexed coreflection and quantified normal forms & \cref{thm:admissibility-coreflection,thm:normal-form,thm:kernel-forcing}\\
Geometric comparison & Universal membership and exact description transport & \cref{lem:ideal-beta,thm:classifier-basechange,thm:partial-transport,int:thm:transport}\\
Native data & Coherent finite prefixes and comparison maps & \cref{seq:thm:prefixes,seq:prop:transport,prop:intensional-new-streams}\\
Oracle localization & Answer support, full sheaf reflector and towers & \cref{thm:ab-comparison,thm:full-cover-oracle,thm:tower}\\
Adaptive protocols & W-recursion, modal descent and native transcript comparison & \cref{thm:well-founded-modal-adequacy,thm:intensional-kernel-pair,struct:thm:Wcomparison,struct:thm:transcript}\\
Numerical oracle adequacy & Dialogue compilers, exact domains and Boolean execution values & \cref{cal:thm:dialogues,cal:prop:transport,cal:thm:adequacy}\\
Arithmetic oracles & Presented predicates and decision localizations & \cref{arith:prop:first-level,arith:prop:jump-strict,arith:thm:post}\\
Regions and histories & Saturated containment, fibres and least-section adjunction & \cref{thm:regions,thm:history-site,thm:least-history,thm:history-search-domain}\\
Local candidates & Localic image and marked pushout & \cref{prop:candidate,thm:descent-pushout,thm:marked-descent,thm:positive-descended-joins}\\
Constructed requirements & Explicit witnesses, successor maps and fusion & \cref{bush:thm:successor,cal:thm:bushy,partition:prop:promises,ramsey:thm:homogeneous,cal:prop:ramsey,closed:thm:low-basis,ks:thm:extension,fusion:thm:diagonal}\\
Computational structure sheaf & Regular syntax, gluing and marked recovery & \cref{struct:prop:syntax,struct:thm:modelgluing,struct:prop:open,cal:thm:gluing,cal:cor:structure,thm:intensional-assembly}\\
Naming obstructions & Descent, torsors, recollement and finite resolutions & \cref{obs:prop:descent,obs:prop:recollement,obs:thm:crown,obs:thm:deformation}\\
Transported names and evidence & Preserved covers, regular images and accepted runs & \cref{thm:intensional-pres-fibration,thm:intensional-old-image,thm:intensional-certificate-transport}\\
\bottomrule
\end{longtable}

\paragraph{Notation.}
$\T$ is ordinary ideal completion, $\Th$ its observation-equipped lift, $\J$ the admitted-ideal monad, and $\F$ the semantic free-domain functor. The domain lift is $\Lift D=D_\bot$; $\LL$ is $\Sigma$-classified partiality. The full numerical function domain is $\FunDom$, with effective-point object $\PC=\Pt\FunDom$; the full subset domain is $\mathcal W$, with $W=\Pt\mathcal W$. The forcing topos is $\FF$. The realizer exposure is $\RealQ$, the free term monad is $\Syn$, and $\Prior\dashv\triangleright$ is the temporal adjunction.

For partiality, protocol substitution, and sheafification, Kleisli extension is written $f^\star$, with its monad determined by the displayed source and target. Binary Kleisli composition is $g\star f=g^\star f$; $f^\diamond$ denotes the positive predicate transformer. The notation $f^\#$ is retained for extension by the free-algebra universal property. Extension along a reflector unit into an arbitrary sheaf is written $\bar f$. The distinct solution operations are $F^\dagger$ for least fixed points, $f^\ddagger$ for finite-return iteration, and $h^\sharp$ for guarded solutions. For geometric change of base, $f^*$ acts on objects, $f^\#$ on locales, and $f^\dagger$ on approximation presentations; the latter is a presentation transformation, not a fixed-point operation.

The selected domain observations are $\OO_D$, whereas $\OO(L)$ is the full frame of a locale. The Scott frame is $\Sc(D)$ and its locale is $\ScLoc(D)$; effective points $\Pt D$ are distinct from locale points. The category of posets is $\Pos$, whereas $\LPos_K$ is locale positivity. The object $\RO(B,r)$ consists of rounded positive traces; $\RegOpen(P)$ denotes regular opens. Presented arithmetic uses $\SigA_n^A,\PiA_n^A$, separately from the truth object $\Sigma$ and the classical complexity classes $\Sigma^0_n,\Pi^0_n$. Substitution in an object slice uses the indexed constructions; comparison with objects formed in another topos uses the explicit maps of Part~\ref{part:basechange}.

\phantomsection
\addcontentsline{toc}{section}{References}

\begin{thebibliography}{99}
\bibitem{soare-terminology}
Robert I. Soare.
\newblock Computability and recursion.
\newblock \emph{Bulletin of Symbolic Logic}, 2(3):284--321, 1996.
\newblock \url{https://doi.org/10.2307/420992}.

\bibitem{bauer}
Andrej Bauer.
\newblock First steps in synthetic computability theory.
\newblock \emph{Electronic Notes in Theoretical Computer Science}, 155:5--31, 2006.
\newblock \url{https://doi.org/10.1016/j.entcs.2005.11.049}.

\bibitem{bauerfixed}
Andrej Bauer.
\newblock On fixed-point theorems in synthetic computability.
\newblock \textbf{Tbilisi Mathematical Journal}, 10(3):167--181, 2017.
\newblock \url{https://doi.org/10.1515/tmj-2017-0107}.

\bibitem{birkedaloosten}
Lars Birkedal and Jaap van Oosten.
\newblock Relative and modified relative realizability.
\newblock \textbf{Annals of Pure and Applied Logic}, 118(1--2):115--132, 2002.
\newblock \url{https://doi.org/10.1016/S0168-0072(01)00122-1}.

\bibitem{cockettlack}
J. R. B. Cockett and Stephen Lack.
\newblock Restriction categories I: categories of partial maps.
\newblock \emph{Theoretical Computer Science}, 270(1--2):223--259, 2002.
\newblock \url{https://doi.org/10.1016/S0304-3975(00)00382-0}.

\bibitem{escardoknapp}
Mart\'in H. Escard\'o and Cory M. Knapp.
\newblock Partial elements and recursion via dominances in univalent type theory.
\newblock In \emph{CSL 2017}, LIPIcs 82, 21:1--21:16, 2017.
\newblock \url{https://doi.org/10.4230/LIPIcs.CSL.2017.21}.

\bibitem{frey}
Jonas Frey.
\newblock \textbf{A fibrational study of realizability toposes}.
\newblock PhD thesis, Universit\'e Paris Diderot, 2013; revised preprint, 2014.
\newblock \url{https://arxiv.org/abs/1403.3672}.

\bibitem{freyassemblies}
Jonas Frey.
\newblock Characterizing partitioned assemblies and realizability toposes.
\newblock \emph{Journal of Pure and Applied Algebra}, 223(5):2000--2014, 2019.
\newblock \url{https://doi.org/10.1016/j.jpaa.2018.08.012}.
\newblock Author version: \url{https://arxiv.org/abs/1404.6997}.


\bibitem{hasegawa}
Masahito Hasegawa.
\newblock The uniformity principle on traced monoidal categories.
\newblock \textbf{Publications of the Research Institute for Mathematical Sciences},
40(3):991--1014, 2004.
\newblock Author version: \url{https://www.kurims.kyoto-u.ac.jp/~hassei/papers/prims04.pdf}.

\bibitem{makkai}
Michael Makkai.
\newblock Avoiding the axiom of choice in general category theory.
\newblock \textbf{Journal of Pure and Applied Algebra}, 108(2):109--173, 1996.
\newblock \url{https://doi.org/10.1016/0022-4049(95)00029-1}.

\bibitem{hofstra}
Pieter Jan Wouter Hofstra.
\newblock \emph{Completions in Realizability}.
\newblock PhD thesis, Utrecht University, 2003.
\newblock \url{https://webspace.science.uu.nl/~ooste110/realizability/Thesis_Hofstra_A4.pdf}.

\bibitem{hofmannsousa}
Dirk Hofmann and Lurdes Sousa.
\newblock Aspects of algebraic algebras.
\newblock \textbf{Logical Methods in Computer Science}, 13(3), article 4, 2017.
\newblock \url{https://doi.org/10.23638/LMCS-13(3:4)2017}.

\bibitem{hofmannwaszkiewicz}
Dirk Hofmann and Pawe\l{} Waszkiewicz.
\newblock Approximation in quantale-enriched categories.
\newblock \textbf{Topology and its Applications}, 158(8):963--977, 2011.
\newblock Preprint: \url{https://arxiv.org/abs/1004.2228}.

\bibitem{hyland}
J. M. E. Hyland.
\newblock The effective topos.
\newblock In A. S. Troelstra and D. van Dalen, editors,
\emph{The L. E. J. Brouwer Centenary Symposium}, Studies in Logic and the
Foundations of Mathematics 110, pages 165--216. North-Holland, 1982.

\bibitem{jsv}
Andr\'e Joyal, Ross Street, and Dominic Verity.
\newblock Traced monoidal categories.
\newblock \textbf{Mathematical Proceedings of the Cambridge Philosophical Society},
119(3):447--468, 1996.
\newblock \url{https://doi.org/10.1017/S0305004100074338}.

\bibitem{pratt}
Vaughan Pratt.
\newblock Chu spaces as a semantic bridge between linear logic and mathematics.
\newblock \textbf{Theoretical Computer Science}, 294(3):439--471, 2003.
\newblock \url{https://doi.org/10.1016/S0304-3975(01)00169-4}.

\bibitem{rosolini}
Giuseppe Rosolini.
\newblock \emph{Continuity and Effectiveness in Topoi}.
\newblock DPhil thesis, University of Oxford, 1986.

\bibitem{stekelenburg}
Wouter Pieter Stekelenburg.
\newblock Regular functors and relative realisability categories.
\newblock \textbf{Mathematical Structures in Computer Science}, 23(5):1082--1110, 2013.
\newblock \url{https://doi.org/10.1017/S0960129512000990}.

\bibitem{sterlingye}
Jonathan Sterling and Lingyuan Ye.
\newblock Domains and classifying topoi.
\newblock Preprint, 2025; version 3, 2026.
\newblock \url{https://arxiv.org/abs/2505.13096v3}.

\bibitem{vanoosten}
Jaap van Oosten.
\newblock \emph{Realizability: An Introduction to its Categorical Side}.
\newblock Studies in Logic and the Foundations of Mathematics 152.
Elsevier, 2008.

\bibitem{jungmoshiervickers}
Achim Jung, M. Andrew Moshier, and Steven Vickers.
\newblock Presenting dcpos and dcpo algebras.
\newblock \textbf{Electronic Notes in Theoretical Computer Science}, 218:209--229, 2008.
\newblock \href{https://achimjungbham.github.io/pub/papers/Jung-Moshier-Vickers-2008-Presenting-dcpos-and-dcpo-algebras.pdf}{Author version}.

\bibitem{lawvere}
F. William Lawvere.
\newblock Diagonal arguments and cartesian closed categories.
\newblock In \textbf{Category Theory, Homology Theory and their Applications II}, Lecture Notes in Mathematics 92, pages 134--145. Springer, 1969.
\newblock Reprinted in \textbf{Reprints in Theory and Applications of Categories}, 15:1--13, 2006.
\newblock \url{https://www.tac.mta.ca/tac/reprints/articles/15/tr15.pdf}.

\bibitem{stubbe}
Isar Stubbe.
\newblock ``Hausdorff distance'' via conical cocompletion.
\newblock \textbf{Cahiers de Topologie et G\'eom\'etrie Diff\'erentielle Cat\'egoriques}, 51(1):51--76, 2010.
\newblock \url{https://arxiv.org/abs/0903.2722}.

\bibitem{goncharovmiliusrauch}
Sergey Goncharov, Stefan Milius, and Christoph Rauch.
\newblock Complete Elgot monads and coalgebraic resumptions.
\newblock \textbf{Electronic Notes in Theoretical Computer Science}, 325:147--168, 2016.
\newblock \url{https://doi.org/10.1016/j.entcs.2016.09.036}.

\bibitem{goncharovetal}
Sergey Goncharov, Lutz Schr\"oder, Christoph Rauch, and Maciej Pir\'og.
\newblock Guarded and unguarded iteration for generalized processes.
\newblock \textbf{Logical Methods in Computer Science}, 15(3), article 1, 2019.
\newblock \url{https://doi.org/10.23638/LMCS-15(3:1)2019}.

\bibitem{goncharovfoundations}
Sergey Goncharov.
\newblock Uniform Elgot iteration in foundations.
\newblock In \textbf{48th International Colloquium on Automata, Languages, and Programming}, LIPIcs 198, article 131, 2021.
\newblock \url{https://doi.org/10.4230/LIPIcs.ICALP.2021.131}.

\bibitem{miliuslitak}
Stefan Milius and Tadeusz Litak.
\newblock Guard your daggers and traces: properties of guarded (co-)recursion.
\newblock \textbf{Fundamenta Informaticae}, 150(3--4):407--449, 2017.
\newblock \url{https://doi.org/10.3233/FI-2017-1475}.

\bibitem{abramskyhaghverdisc}
Samson Abramsky, Esfandiar Haghverdi, and Philip Scott.
\newblock Geometry of Interaction and linear combinatory algebras.
\newblock \textbf{Mathematical Structures in Computer Science}, 12(5):625--665, 2002.
\newblock \url{https://doi.org/10.1017/S0960129502003730}.

\bibitem{hasegawaclosed}
Masahito Hasegawa.
\newblock On traced monoidal closed categories.
\newblock \textbf{Mathematical Structures in Computer Science}, 19(2):217--244, 2009.
\newblock \url{https://doi.org/10.1017/S0960129508007184}.

\bibitem{ahmanbauer}
Danel Ahman and Andrej Bauer.
\newblock Sheaves as oracle computations.
\newblock Preprint, 2026, version 2.
\newblock \url{https://arxiv.org/abs/2602.22135v2}.
\newblock Invited-talk version in \emph{FSCD 2026}, LIPIcs 378, 1:1--1:17.
\newblock \url{https://doi.org/10.4230/LIPIcs.FSCD.2026.1}.


\bibitem{caramello}
Olivia Caramello.
\newblock Topos-theoretic background.
\newblock Preprint, 2014, Sections~5--6.
\newblock \url{https://preprints.ihes.fr/storage/ToposTheoreticPreliminariesOliviaCaramelloRev.pdf}.

\bibitem{lurie}
Jacob Lurie.
\newblock Derived Algebraic Geometry V: Structured Spaces.
\newblock 2011, Definition~1.2.8.
\newblock \url{https://www.math.ias.edu/~lurie/papers/DAG-V.pdf}.

\bibitem{vandenbergmoerdijk}
Benno van den Berg and Ieke Moerdijk.
\newblock W-types in sheaves.
\newblock Preprint, 2008.
\newblock \url{https://arxiv.org/abs/0810.2398}.

\bibitem{moerdijkpalmgren}
Ieke Moerdijk and Erik Palmgren.
\newblock Wellfounded trees in categories.
\newblock \emph{Annals of Pure and Applied Logic},
104(1--3):189--218, 2000.
\newblock \url{https://doi.org/10.1016/S0168-0072(00)00012-9}.

\bibitem{gambinokock}
Nicola Gambino and Joachim Kock.
\newblock Polynomial functors and polynomial monads.
\newblock \emph{Mathematical Proceedings of the Cambridge
Philosophical Society}, 154(1):153--192, 2013.
\newblock \url{https://doi.org/10.1017/S0305004112000394}.

\bibitem{bauerinstance}
Andrej Bauer.
\newblock Instance reducibility and Weihrauch degrees.
\newblock \emph{Logical Methods in Computer Science}, 18(3), article 20, 2022.
\newblock \url{https://doi.org/10.46298/LMCS-18(3:20)2022}.

\bibitem{kihara}
Takayuki Kihara.
\newblock Lawvere--Tierney topologies for computability theorists.
\newblock \emph{Transactions of the American Mathematical Society, Series B}, 10:48--85, 2023.
\newblock \url{https://doi.org/10.1090/btran/134}.

\bibitem{johnstone}
Peter T. Johnstone.
\newblock \emph{Sketches of an Elephant: A Topos Theory Compendium}.
\newblock Volumes 1 and 2, Oxford Logic Guides 43 and 44.
Oxford University Press, 2002.

\bibitem{joyaltierney}
Andr\'e Joyal and Myles Tierney.
\newblock \emph{An Extension of the Galois Theory of Grothendieck}.
\newblock Memoirs of the American Mathematical Society 51, no.~309, 1984.
\newblock \url{https://doi.org/10.1090/memo/0309}.

\bibitem{maclanemoerdijk}
Saunders Mac Lane and Ieke Moerdijk.
\newblock \emph{Sheaves in Geometry and Logic: A First Introduction to Topos Theory}.
\newblock Universitext. Springer, 1992.
\newblock \url{https://doi.org/10.1007/978-1-4612-0927-0}.

\bibitem{mathias}
A. R. D. Mathias.
\newblock Happy families.
\newblock \emph{Annals of Mathematical Logic}, 12(1):59--111, 1977.

\bibitem{sacks}
Gerald E. Sacks.
\newblock Forcing with perfect closed sets.
\newblock In \emph{Axiomatic Set Theory}, Proceedings of Symposia in Pure
Mathematics 13, Part I, pages 331--355. American Mathematical Society, 1971.

\bibitem{vickers}
Steven Vickers.
\newblock \emph{Topology via Logic}.
\newblock Cambridge Tracts in Theoretical Computer Science 5.
Cambridge University Press, 1989.

\bibitem{vickersau}
Steven Vickers.
\newblock Arithmetic universes and classifying toposes.
\newblock Preprint, 2017.
\newblock \url{https://arxiv.org/abs/1701.04611}.


\bibitem{relative-recursion} Jaap van Oosten.
\emph{A general form of relative recursion}. Manuscript, revised
9 December 2005.
\url{https://webspace.science.uu.nl/~ooste110/realizability/relrec.pdf}.

\bibitem{faber-oosten} Eric Faber and Jaap van Oosten.
Effective operations of type 2 in PCAs.
\emph{Computability} 5 (2016), 127--146.
\url{https://doi.org/10.3233/COM-150048}.
Author version:
\url{https://webspace.science.uu.nl/~ooste110/realizability/type2inpcas.pdf}.

\bibitem{forster-kirst-mueck} Yannick Forster, Dominik Kirst, and Niklas M\"uck.
The Kleene--Post and Post's theorem in the calculus of inductive
constructions. \emph{CSL 2024}, LIPIcs 288, article 37, 37:1--37:20.
\url{https://doi.org/10.4230/LIPIcs.CSL.2024.37}.

\bibitem{khan-miller} Mushfeq Khan and Joseph S. Miller.
\emph{Forcing with bushy trees}. Author manuscript, 23 April 2019.
\url{https://people.math.wisc.edu/~jsmiller8/Papers/bushy_trees.pdf}.

\bibitem{patey-jumps} Ludovic Patey.
\emph{Controlling iterated jumps of solutions to combinatorial
problems}. arXiv:1509.05340v2, 19 January 2016.
\url{https://arxiv.org/abs/1509.05340v2}.
Section and lemma numbering refers to this 32-page version.

\bibitem{jockusch-soare} Carl G. Jockusch, Jr. and Robert I. Soare.
$\Pi^0_1$ classes and degrees of theories.
\emph{Transactions of the American Mathematical Society} 173 (1972),
33--56. \url{https://doi.org/10.1090/S0002-9947-1972-0316227-0}.

\bibitem{slaman-woodin} Theodore A. Slaman and W. Hugh Woodin.
\emph{Definability in degree structures}. Manuscript, 21 July 2005, Chapter 8,
especially Section 8.2 and Definitions 1.1.3 and 8.2.2.
\url{https://math.hawaii.edu/~bjoern/pdf/sw.pdf}.

\bibitem{sga4-gluing}
Michael Artin, Alexander Grothendieck, and Jean-Louis Verdier, editors.
\newblock \emph{Th\'eorie des topos et cohomologie \'etale des sch\'emas
(SGA~4), Tome~1: Th\'eorie des topos}.
\newblock Lecture Notes in Mathematics 269. Springer, 1972.
\newblock Expos\'e~IV, Section~9, especially Th\'eor\`eme~9.5.4.
\newblock \url{https://doi.org/10.1007/BFb0081551}.

\bibitem{faul-manuell-frames}
Peter F. Faul and Graham R. Manuell.
\newblock Artin glueings of frames as semidirect products.
\newblock Author manuscript, arXiv:1907.05104v2, 2019.
\newblock \url{https://arxiv.org/abs/1907.05104v2}.

\bibitem{faul-manuell-siqueira}
Peter F. Faul, Graham Manuell, and Jos\'e Siqueira.
\newblock Artin glueings of toposes as adjoint split extensions.
\newblock \emph{Journal of Pure and Applied Algebra}, 227(5):107273, 2023.
\newblock \url{https://doi.org/10.1016/j.jpaa.2022.107273}.
\newblock Author version: \url{https://arxiv.org/abs/2012.04963v2}.

\bibitem{sterling-tait}
Jonathan Sterling.
\newblock \emph{First Steps in Synthetic Tait Computability:
The Objective Metatheory of Cubical Type Theory}.
\newblock PhD thesis, Carnegie Mellon University, CMU-CS-21-142, 2021.
\newblock \url{https://csd.cmu.edu/sites/default/files/phd-thesis/CMU-CS-21-142.pdf}.

\bibitem{bhatt-scholze-proetale}
Bhargav Bhatt and Peter Scholze.
\newblock The pro-\'etale topology for schemes.
\newblock \emph{Ast\'erisque}, 369:99--201, 2015.
\newblock \url{https://doi.org/10.24033/ast.960}.
\newblock Author version: \url{https://www.math.uni-bonn.de/people/scholze/proetale.pdf}.

\bibitem{scholze-condensed}
Peter Scholze, reporting joint work with Dustin Clausen.
\newblock \emph{Lectures on Condensed Mathematics}.
\newblock University of Bonn lecture notes, 2019.
\newblock \url{https://www.math.uni-bonn.de/people/scholze/Condensed.pdf}.

\bibitem{franjou-pirashvili}
Vincent Franjou and Teimuraz Pirashvili.
\newblock Comparison of abelian categories recollements.
\newblock \textbf{Documenta Mathematica}, 9:41--56, 2004.
\newblock \url{https://doi.org/10.4171/DM/156}.
\newblock Author version: \url{https://arxiv.org/abs/math/0403081}.

\bibitem{stacks-project}
The Stacks Project Authors.
\newblock \emph{The Stacks Project}.
\newblock First cohomology and torsors, Tags 03AG and 03AJ;
first cohomology and invertible sheaves, Tag 040D;
Picard groups of thickenings, Tag 0C6R.
\newblock \url{https://stacks.math.columbia.edu/tag/03AG};
\url{https://stacks.math.columbia.edu/tag/03AJ};
\url{https://stacks.math.columbia.edu/tag/040D};
\url{https://stacks.math.columbia.edu/tag/0C6R}.


\bibitem{ambos-spies-etal}
K.~Ambos-Spies, B.~Kjos-Hanssen, S.~Lempp, and T.~A.~Slaman,
\emph{Comparing DNR and WWKL},
Journal of Symbolic Logic \textbf{69} (2004), no.~4, 1089--1104.
Theorem 1.8 and Section 2.
\url{https://math.berkeley.edu/~slaman/papers/DNR.pdf}.


\bibitem{dobrinen-ramsey}
N.~Dobrinen,
\emph{Tutorial: Ramsey theory in Forcing}, RIMS lectures, 2016.
Slides on Ellentuck space and Mathias forcing.
\url{https://www.sci.shizuoka.ac.jp/~math/yorioka/RIMS2016/161128-3_Dobrinen.pdf}.


\bibitem{geschke-quickert}
S.~Geschke and S.~Quickert,
\emph{On Sacks Forcing and the Sacks Property}, author PDF, Section 2.
\url{https://www.math.uni-hamburg.de/home/geschke/papers/OnSacks_main.pdf}.


\bibitem{pierobon-viale}
M.~Pierobon and M.~Viale,
\emph{Boolean valued models, presheaves, and \texorpdfstring{\'etal\'e}{etale} spaces},
arXiv:2006.14852v6, 23 July 2026.
\url{https://arxiv.org/abs/2006.14852v6}.


\bibitem{smykalla-mariano}
M.~V.~Smykalla and H.~L.~Mariano,
\emph{Relating forcing relations}, arXiv:2605.26896v1, 26 May 2026.
Theorem 4.22.
\url{https://arxiv.org/abs/2605.26896v1}.

\bibitem{kavvosarticle}
G. A. Kavvos.
\newblock On the semantics of intensionality.
\newblock In \emph{Foundations of Software Science and Computation
Structures (FoSSaCS 2017)}, Lecture Notes in Computer Science 10203,
pages 550--566. Springer, 2017.
\newblock \url{https://doi.org/10.1007/978-3-662-54458-7_32}.
\newblock Author version: \url{https://arxiv.org/abs/1602.01365v6}.

\bibitem{chenko}
Liang-Ting Chen and Hsiang-Shang Ko.
\newblock Realising intensional S4 and GL modalities.
\newblock In \emph{30th EACSL Annual Conference on Computer Science
Logic (CSL 2022)}, LIPIcs 216, pages 14:1--14:17, 2022.
\newblock \url{https://doi.org/10.4230/LIPIcs.CSL.2022.14}.

\bibitem{fioreplotkinturi}
Marcelo Fiore, Gordon Plotkin, and Daniele Turi.
\newblock Abstract syntax and variable binding.
\newblock In \emph{Proceedings of the 14th Annual IEEE Symposium on
Logic in Computer Science (LICS 1999)}, pages 193--202, 1999.
\newblock \url{https://doi.org/10.1109/LICS.1999.782615}.

\bibitem{daviespfenning}
Rowan Davies and Frank Pfenning.
\newblock A modal analysis of staged computation.
\newblock \emph{Journal of the ACM}, 48(3):555--604, 2001.
\newblock \url{https://doi.org/10.1145/382780.382785}.

\bibitem{birkedalguarded}
Lars Birkedal, Rasmus Ejlers M{\o}gelberg, Jan Schwinghammer,
and Kristian St{\o}vring.
\newblock First steps in synthetic guarded domain theory:
step-indexing in the topos of trees.
\newblock \emph{Logical Methods in Computer Science}, 8(4), article 1, 2012.
\newblock \url{https://doi.org/10.2168/LMCS-8(4:1)2012}.

\bibitem{ramesh}
Sridhar Ramesh.
\newblock \emph{Introspective Theories and Geminal Categories}.
\newblock PhD dissertation, University of California, Berkeley, 2023.
\newblock \url{https://escholarship.org/uc/item/3mn0c475}.


\end{thebibliography}

\end{document}
